\documentclass{article}

\usepackage[utf8]{inputenc}
\usepackage[T1]{fontenc}
\usepackage[catalan]{babel}
\usepackage{cmbright}
\usepackage{lmodern}
\usepackage{geometry}
\usepackage{tikz}
\usepackage{hyperref}
\usepackage[bf,sf,pagestyles]{titlesec}
\usepackage{titling}
\usepackage[runin]{abstract}
\usepackage[font={footnotesize, sf}, labelfont=bf]{caption} 
\usepackage{siunitx}
\usepackage{graphicx}
\usepackage{booktabs}
\usepackage{listings}
\usepackage{textcomp}
\usepackage{amsmath,amssymb}
\usepackage[sort]{cleveref}
\usepackage{enumitem}

\renewcommand{\familydefault}{\rmdefault}
\renewcommand{\sfdefault}{cmbr}
\renewcommand{\ttdefault}{cmtl}


\geometry{
	a4paper,
	right = 2.5cm,
	left = 2.5cm,
	bottom = 3cm,
	top = 3cm
}

\hypersetup{
	colorlinks,
	linkcolor = {red!50!blue},
	citecolor = {red!50!blue},
	linktoc = page
}

\numberwithin{table}{section}
\numberwithin{figure}{section}
\numberwithin{equation}{section}

\graphicspath{{./figs/}}

% Unitats
\sisetup{
	inter-unit-product = \ensuremath{ \, },
	allow-number-unit-breaks = true,
	math-celsius = {}^{\circ}\kern-\scriptspace C,
	detect-family = true,
	detect-shape = true,
	list-final-separator = { y },
	list-pair-separator = { y },
	list-units = single,
	separate-uncertainty = true
}

\lstdefinestyle{customc}{
  belowcaptionskip=1\baselineskip,
	backgroundcolor = \color{gray!20!white},
  breaklines=true,
	breakatwhitespace = true,
	postbreak = \textrightarrow\,,
  xleftmargin=\parindent,
	frame = l,
  language=C,
  showstringspaces=false,
  basicstyle=\small\ttfamily,
  keywordstyle=\bfseries\color{green!40!black},
  commentstyle=\sffamily\itshape\color{purple!40!black},
  identifierstyle=\color{blue},
  stringstyle=\color{orange},
	numbers = left, 
	numberstyle = \tiny\ttfamily,
}
\lstset{texcl=true,style=customc}
\renewcommand{\lstlistingname}{Programa}

\newpagestyle{pagina}{
	\headrule
	\sethead*{\sffamily \bfseries Práctica 4}{}{\theauthor}
	\footrule
	\setfoot*{}{}{\sffamily \thepage}
}
\pagestyle{pagina}


\title{\sffamily \bfseries Modelització i anàlisi financera de l'empresa Caterpillar \vspace{1em} \hrule}
\author{\sffamily Andreu Arderiu, Arnau Mas}
\date{\sffamily 14 de juny de 2019}

\begin{document}

\begin{titlepage}
	\clearpage
	\thispagestyle{empty}
	\maketitle	

	\renewcommand{\abstractname}{\bfseries \sffamily Resum:}
	\begin{abstract}
		Abstract en català
	\end{abstract}

	\renewcommand{\abstractname}{\bfseries \sffamily Abstract:}
	\begin{abstract}
		Abstract en anglès
	\end{abstract}
	
	\begin{figure}[b]
		\thispagestyle{empty}
		\center \includegraphics[scale = 0.3]{uab}
	\end{figure}
\end{titlepage}

\textbf{En aquest estudi de l'empresa Caterpillar amb \textit{Symbol} bursàtil a la borsa de Nova York CAT, ens centrarem en els anys $2015-2018$. Mitjançant la comanda $getSymbols$ de R prenem un total de $1005$ dades corresponents als preus de tancament diari. A partir d'aquestes dades, podem calcular els estadístics bàsics de les rendibilitats, és a dir, els increments dels logaritmes dels preus. Trobem així una mitjana anual de $\hat{\mu}_a=7.786318$, una desviació estàndard anual de $\hat{\sigma}_a=26.7607$, i finalment una curtosi diària de $\hat{ku}_d=2.806694$. }

\section{Context}
\begin{figure}[htb]
	\centering \small \sffamily
	\includegraphics[width=0.4\textwidth]{cat-logo}
	\caption{Logotip de Caterpillar}
	\label{fig:logo}
\end{figure}

Caterpillar Inc. (NYSE: CAT) és una empresa multinacional de maquinària industrial i de construcció amb seu a Illinois. Actualment és l'empresa líder en l'àmbit de la mquinària de construcció i industrial \cite{cat} i es troba en la posició 65 al rànquing \emph{Fortune 500} i en la posició 238 del rànquing \emph{Fortune Global 500} de l'any 2019 \cite{fortune,global}. El preu actual de l'acció de Caterpillar és d'uns \$127. Actualment té \num{104000} treballadors. 

Caterpillar va néixer l'any 1925 de la fusió de les empreses Holt Manufacturing Company i C. L: Best Tractor Co. El fundador de la primera, Benjamin Holt, va patentar un dels primers vehicles que funcionaven amb tracció per eruga (d'aquí el nom Caterpillar), un sistema de tracció que consisteix en una sèrie de plaques rígides unides per una cinta que gira al voltant de les rodes del vehicle. Amb aquest sistema, la superfície de contacte del vehicle amb el sòl s'incrementa de manera que s'exerceix menys pressió i és possible circular per a terrenys més complicats. Originalment, els vehicles de Holt es feien servir a l'oest americà, on la combinació del terreny àrid i irregular del desert californià i la manca d'infrastructura dificultaven les tasques agràries i de transport. Molt aviat, però, l'empresa va desenvolupar vehicles militars durant la Primera Guerra mundial, sent un dels precursors del tanc. Durant el període d'entreguerra i sobretot després de la Segona Guerra mundial, Caterpillar va evolucionar cap a l'àmbit de la construcció, primer al mercat americà i després a escala global, sobretot als països en desenvolupament. 

\begin{figure}[htb]
	\centering \sffamily \small
	% Created by tikzDevice version 0.12 on 2019-06-12 19:07:25
% !TEX encoding = UTF-8 Unicode
\begin{tikzpicture}[x=1pt,y=1pt]
\definecolor{fillColor}{RGB}{255,255,255}
\path[use as bounding box,fill=fillColor,fill opacity=0.00] (0,0) rectangle (433.62,122.86);
\begin{scope}
\path[clip] ( 21.60, 32.40) rectangle (412.02,122.86);
\definecolor{drawColor}{RGB}{169,169,169}

\path[draw=drawColor,line width= 0.4pt,line join=round,line cap=round] ( 36.06, 35.75) -- ( 36.06,107.54);

\path[draw=drawColor,line width= 0.4pt,line join=round,line cap=round] ( 56.23, 35.75) -- ( 56.23,107.54);

\path[draw=drawColor,line width= 0.4pt,line join=round,line cap=round] ( 76.30, 35.75) -- ( 76.30,107.54);

\path[draw=drawColor,line width= 0.4pt,line join=round,line cap=round] ( 96.36, 35.75) -- ( 96.36,107.54);

\path[draw=drawColor,line width= 0.4pt,line join=round,line cap=round] (116.48, 35.75) -- (116.48,107.54);

\path[draw=drawColor,line width= 0.4pt,line join=round,line cap=round] (136.49, 35.75) -- (136.49,107.54);

\path[draw=drawColor,line width= 0.4pt,line join=round,line cap=round] (156.51, 35.75) -- (156.51,107.54);

\path[draw=drawColor,line width= 0.4pt,line join=round,line cap=round] (176.68, 35.75) -- (176.68,107.54);

\path[draw=drawColor,line width= 0.4pt,line join=round,line cap=round] (196.80, 35.75) -- (196.80,107.54);

\path[draw=drawColor,line width= 0.4pt,line join=round,line cap=round] (216.86, 35.75) -- (216.86,107.54);

\path[draw=drawColor,line width= 0.4pt,line join=round,line cap=round] (236.93, 35.75) -- (236.93,107.54);

\path[draw=drawColor,line width= 0.4pt,line join=round,line cap=round] (256.94, 35.75) -- (256.94,107.54);

\path[draw=drawColor,line width= 0.4pt,line join=round,line cap=round] (277.11, 35.75) -- (277.11,107.54);

\path[draw=drawColor,line width= 0.4pt,line join=round,line cap=round] (297.18, 35.75) -- (297.18,107.54);

\path[draw=drawColor,line width= 0.4pt,line join=round,line cap=round] (317.24, 35.75) -- (317.24,107.54);

\path[draw=drawColor,line width= 0.4pt,line join=round,line cap=round] (337.31, 35.75) -- (337.31,107.54);

\path[draw=drawColor,line width= 0.4pt,line join=round,line cap=round] (357.38, 35.75) -- (357.38,107.54);

\path[draw=drawColor,line width= 0.4pt,line join=round,line cap=round] (377.39, 35.75) -- (377.39,107.54);

\path[draw=drawColor,line width= 0.4pt,line join=round,line cap=round] (397.56, 35.75) -- (397.56,107.54);

\path[draw=drawColor,line width= 0.4pt,line join=round,line cap=round] (397.56, 35.75) -- (397.56,107.54);
\end{scope}
\begin{scope}
\path[clip] (  0.00,  0.00) rectangle (433.62,122.86);
\definecolor{drawColor}{gray}{0.96}

\path[draw=drawColor,line width= 0.4pt,line join=round,line cap=round] ( 36.06, 32.40) -- (397.56, 32.40);

\path[draw=drawColor,line width= 0.4pt,line join=round,line cap=round] ( 36.06, 32.40) -- ( 36.06, 35.64);

\path[draw=drawColor,line width= 0.4pt,line join=round,line cap=round] ( 56.23, 32.40) -- ( 56.23, 35.64);

\path[draw=drawColor,line width= 0.4pt,line join=round,line cap=round] ( 76.30, 32.40) -- ( 76.30, 35.64);

\path[draw=drawColor,line width= 0.4pt,line join=round,line cap=round] ( 96.36, 32.40) -- ( 96.36, 35.64);

\path[draw=drawColor,line width= 0.4pt,line join=round,line cap=round] (116.48, 32.40) -- (116.48, 35.64);

\path[draw=drawColor,line width= 0.4pt,line join=round,line cap=round] (136.49, 32.40) -- (136.49, 35.64);

\path[draw=drawColor,line width= 0.4pt,line join=round,line cap=round] (156.51, 32.40) -- (156.51, 35.64);

\path[draw=drawColor,line width= 0.4pt,line join=round,line cap=round] (176.68, 32.40) -- (176.68, 35.64);

\path[draw=drawColor,line width= 0.4pt,line join=round,line cap=round] (196.80, 32.40) -- (196.80, 35.64);

\path[draw=drawColor,line width= 0.4pt,line join=round,line cap=round] (216.86, 32.40) -- (216.86, 35.64);

\path[draw=drawColor,line width= 0.4pt,line join=round,line cap=round] (236.93, 32.40) -- (236.93, 35.64);

\path[draw=drawColor,line width= 0.4pt,line join=round,line cap=round] (256.94, 32.40) -- (256.94, 35.64);

\path[draw=drawColor,line width= 0.4pt,line join=round,line cap=round] (277.11, 32.40) -- (277.11, 35.64);

\path[draw=drawColor,line width= 0.4pt,line join=round,line cap=round] (297.18, 32.40) -- (297.18, 35.64);

\path[draw=drawColor,line width= 0.4pt,line join=round,line cap=round] (317.24, 32.40) -- (317.24, 35.64);

\path[draw=drawColor,line width= 0.4pt,line join=round,line cap=round] (337.31, 32.40) -- (337.31, 35.64);

\path[draw=drawColor,line width= 0.4pt,line join=round,line cap=round] (357.38, 32.40) -- (357.38, 35.64);

\path[draw=drawColor,line width= 0.4pt,line join=round,line cap=round] (377.39, 32.40) -- (377.39, 35.64);

\path[draw=drawColor,line width= 0.4pt,line join=round,line cap=round] (397.56, 32.40) -- (397.56, 35.64);
\end{scope}
\begin{scope}
\path[clip] (  0.00,  0.00) rectangle (433.62,122.86);
\definecolor{drawColor}{gray}{0.20}

\path[draw=drawColor,line width= 0.4pt,line join=round,line cap=round] ( 36.06, 32.40) -- (397.56, 32.40);

\path[draw=drawColor,line width= 0.6pt,line join=round,line cap=round] ( 36.06, 32.40) -- ( 36.06, 28.08);

\path[draw=drawColor,line width= 0.6pt,line join=round,line cap=round] ( 56.23, 32.40) -- ( 56.23, 28.08);

\path[draw=drawColor,line width= 0.6pt,line join=round,line cap=round] ( 76.30, 32.40) -- ( 76.30, 28.08);

\path[draw=drawColor,line width= 0.6pt,line join=round,line cap=round] ( 96.36, 32.40) -- ( 96.36, 28.08);

\path[draw=drawColor,line width= 0.6pt,line join=round,line cap=round] (116.48, 32.40) -- (116.48, 28.08);

\path[draw=drawColor,line width= 0.6pt,line join=round,line cap=round] (136.49, 32.40) -- (136.49, 28.08);

\path[draw=drawColor,line width= 0.6pt,line join=round,line cap=round] (156.51, 32.40) -- (156.51, 28.08);

\path[draw=drawColor,line width= 0.6pt,line join=round,line cap=round] (176.68, 32.40) -- (176.68, 28.08);

\path[draw=drawColor,line width= 0.6pt,line join=round,line cap=round] (196.80, 32.40) -- (196.80, 28.08);

\path[draw=drawColor,line width= 0.6pt,line join=round,line cap=round] (216.86, 32.40) -- (216.86, 28.08);

\path[draw=drawColor,line width= 0.6pt,line join=round,line cap=round] (236.93, 32.40) -- (236.93, 28.08);

\path[draw=drawColor,line width= 0.6pt,line join=round,line cap=round] (256.94, 32.40) -- (256.94, 28.08);

\path[draw=drawColor,line width= 0.6pt,line join=round,line cap=round] (277.11, 32.40) -- (277.11, 28.08);

\path[draw=drawColor,line width= 0.6pt,line join=round,line cap=round] (297.18, 32.40) -- (297.18, 28.08);

\path[draw=drawColor,line width= 0.6pt,line join=round,line cap=round] (317.24, 32.40) -- (317.24, 28.08);

\path[draw=drawColor,line width= 0.6pt,line join=round,line cap=round] (337.31, 32.40) -- (337.31, 28.08);

\path[draw=drawColor,line width= 0.6pt,line join=round,line cap=round] (357.38, 32.40) -- (357.38, 28.08);

\path[draw=drawColor,line width= 0.6pt,line join=round,line cap=round] (377.39, 32.40) -- (377.39, 28.08);

\path[draw=drawColor,line width= 0.6pt,line join=round,line cap=round] (397.56, 32.40) -- (397.56, 28.08);

\path[draw=drawColor,line width= 0.6pt,line join=round,line cap=round] (397.56, 32.40) -- (397.56, 28.08);

\node[text=drawColor,anchor=base,inner sep=0pt, outer sep=0pt, scale=  0.81] at ( 36.06, 27.00) {};

\node[text=drawColor,anchor=base,inner sep=0pt, outer sep=0pt, scale=  0.81] at ( 36.06, 17.28) {Dec};

\node[text=drawColor,anchor=base,inner sep=0pt, outer sep=0pt, scale=  0.81] at ( 36.06,  7.56) {2000};

\node[text=drawColor,anchor=base,inner sep=0pt, outer sep=0pt, scale=  0.81] at ( 76.30, 27.00) {};

\node[text=drawColor,anchor=base,inner sep=0pt, outer sep=0pt, scale=  0.81] at ( 76.30, 17.28) {Dec};

\node[text=drawColor,anchor=base,inner sep=0pt, outer sep=0pt, scale=  0.81] at ( 76.30,  7.56) {2002};

\node[text=drawColor,anchor=base,inner sep=0pt, outer sep=0pt, scale=  0.81] at (116.48, 27.00) {};

\node[text=drawColor,anchor=base,inner sep=0pt, outer sep=0pt, scale=  0.81] at (116.48, 17.28) {Dec};

\node[text=drawColor,anchor=base,inner sep=0pt, outer sep=0pt, scale=  0.81] at (116.48,  7.56) {2004};

\node[text=drawColor,anchor=base,inner sep=0pt, outer sep=0pt, scale=  0.81] at (156.51, 27.00) {};

\node[text=drawColor,anchor=base,inner sep=0pt, outer sep=0pt, scale=  0.81] at (156.51, 17.28) {Dec};

\node[text=drawColor,anchor=base,inner sep=0pt, outer sep=0pt, scale=  0.81] at (156.51,  7.56) {2006};

\node[text=drawColor,anchor=base,inner sep=0pt, outer sep=0pt, scale=  0.81] at (196.80, 27.00) {};

\node[text=drawColor,anchor=base,inner sep=0pt, outer sep=0pt, scale=  0.81] at (196.80, 17.28) {Dec};

\node[text=drawColor,anchor=base,inner sep=0pt, outer sep=0pt, scale=  0.81] at (196.80,  7.56) {2008};

\node[text=drawColor,anchor=base,inner sep=0pt, outer sep=0pt, scale=  0.81] at (236.93, 27.00) {};

\node[text=drawColor,anchor=base,inner sep=0pt, outer sep=0pt, scale=  0.81] at (236.93, 17.28) {Dec};

\node[text=drawColor,anchor=base,inner sep=0pt, outer sep=0pt, scale=  0.81] at (236.93,  7.56) {2010};

\node[text=drawColor,anchor=base,inner sep=0pt, outer sep=0pt, scale=  0.81] at (277.11, 27.00) {};

\node[text=drawColor,anchor=base,inner sep=0pt, outer sep=0pt, scale=  0.81] at (277.11, 17.28) {Dec};

\node[text=drawColor,anchor=base,inner sep=0pt, outer sep=0pt, scale=  0.81] at (277.11,  7.56) {2012};

\node[text=drawColor,anchor=base,inner sep=0pt, outer sep=0pt, scale=  0.81] at (317.24, 27.00) {};

\node[text=drawColor,anchor=base,inner sep=0pt, outer sep=0pt, scale=  0.81] at (317.24, 17.28) {Dec};

\node[text=drawColor,anchor=base,inner sep=0pt, outer sep=0pt, scale=  0.81] at (317.24,  7.56) {2014};

\node[text=drawColor,anchor=base,inner sep=0pt, outer sep=0pt, scale=  0.81] at (357.38, 27.00) {};

\node[text=drawColor,anchor=base,inner sep=0pt, outer sep=0pt, scale=  0.81] at (357.38, 17.28) {Dec};

\node[text=drawColor,anchor=base,inner sep=0pt, outer sep=0pt, scale=  0.81] at (357.38,  7.56) {2016};

\node[text=drawColor,anchor=base,inner sep=0pt, outer sep=0pt, scale=  0.81] at (397.56, 27.00) {};

\node[text=drawColor,anchor=base,inner sep=0pt, outer sep=0pt, scale=  0.81] at (397.56, 17.28) {Dec};

\node[text=drawColor,anchor=base,inner sep=0pt, outer sep=0pt, scale=  0.81] at (397.56,  7.56) {2018};
\end{scope}
\begin{scope}
\path[clip] ( 21.60,107.54) rectangle (412.02,119.51);
\definecolor{drawColor}{gray}{0.20}

\node[text=drawColor,anchor=base west,inner sep=0pt, outer sep=0pt, scale=  0.99] at ( 36.06,111.25) {\bfseries Tancaments anuals entre 2000 i 2018};

\node[text=drawColor,anchor=base east,inner sep=0pt, outer sep=0pt, scale=  0.90] at (392.16,111.46) {2000-12-29 / 2018-12-31};
\end{scope}
\begin{scope}
\path[clip] ( 21.60, 35.75) rectangle (412.02,107.54);
\definecolor{drawColor}{RGB}{169,169,169}

\path[draw=drawColor,line width= 0.4pt,line join=round,line cap=round] ( 36.06, 44.88) -- (397.56, 44.88);

\path[draw=drawColor,line width= 0.4pt,line join=round,line cap=round] ( 36.06, 55.54) -- (397.56, 55.54);

\path[draw=drawColor,line width= 0.4pt,line join=round,line cap=round] ( 36.06, 66.20) -- (397.56, 66.20);

\path[draw=drawColor,line width= 0.4pt,line join=round,line cap=round] ( 36.06, 76.86) -- (397.56, 76.86);

\path[draw=drawColor,line width= 0.4pt,line join=round,line cap=round] ( 36.06, 87.52) -- (397.56, 87.52);

\path[draw=drawColor,line width= 0.4pt,line join=round,line cap=round] ( 36.06, 98.17) -- (397.56, 98.17);
\end{scope}
\begin{scope}
\path[clip] (  0.00,  0.00) rectangle (433.62,122.86);
\definecolor{drawColor}{gray}{0.20}

\node[text=drawColor,anchor=base east,inner sep=0pt, outer sep=0pt, scale=  0.81] at ( 25.26, 43.02) { 40};

\node[text=drawColor,anchor=base east,inner sep=0pt, outer sep=0pt, scale=  0.81] at ( 25.26, 53.68) { 60};

\node[text=drawColor,anchor=base east,inner sep=0pt, outer sep=0pt, scale=  0.81] at ( 25.26, 64.34) { 80};

\node[text=drawColor,anchor=base east,inner sep=0pt, outer sep=0pt, scale=  0.81] at ( 25.26, 75.00) {100};

\node[text=drawColor,anchor=base east,inner sep=0pt, outer sep=0pt, scale=  0.81] at ( 25.26, 85.66) {120};

\node[text=drawColor,anchor=base east,inner sep=0pt, outer sep=0pt, scale=  0.81] at ( 25.26, 96.32) {140};

\node[text=drawColor,anchor=base west,inner sep=0pt, outer sep=0pt, scale=  0.81] at (408.36, 43.02) { 40};

\node[text=drawColor,anchor=base west,inner sep=0pt, outer sep=0pt, scale=  0.81] at (408.36, 53.68) { 60};

\node[text=drawColor,anchor=base west,inner sep=0pt, outer sep=0pt, scale=  0.81] at (408.36, 64.34) { 80};

\node[text=drawColor,anchor=base west,inner sep=0pt, outer sep=0pt, scale=  0.81] at (408.36, 75.00) {100};

\node[text=drawColor,anchor=base west,inner sep=0pt, outer sep=0pt, scale=  0.81] at (408.36, 85.66) {120};

\node[text=drawColor,anchor=base west,inner sep=0pt, outer sep=0pt, scale=  0.81] at (408.36, 96.32) {140};
\end{scope}
\begin{scope}
\path[clip] ( 21.60, 35.75) rectangle (412.02,107.54);
\definecolor{drawColor}{RGB}{0,0,255}

\path[draw=drawColor,line width= 0.8pt,line join=round] ( 36.06, 36.17) --
	( 56.23, 37.49) --
	( 76.30, 35.75) --
	( 96.36, 45.69) --
	(116.48, 49.55) --
	(136.49, 54.35) --
	(156.51, 56.25) --
	(176.68, 62.24) --
	(196.80, 47.37) --
	(216.86, 53.94) --
	(236.93, 73.48) --
	(256.94, 71.85) --
	(277.11, 71.32) --
	(297.18, 71.96) --
	(317.24, 72.34) --
	(337.31, 59.78) --
	(357.38, 72.99) --
	(377.39,107.54) --
	(397.56, 91.28);
\end{scope}
\end{tikzpicture}

	\caption{Preus de tancament anual entre 2000 i 2018 de Caterpillar}
	\label{fig:tancament}
\end{figure}

A la \cref{fig:tancament} podem veure els preus de tancament anuals entre l'any 2000 i l'actualitat. La tendència és en general a l'alça. Tot i així, cal remarcar una caiguda de prop del \SI{50}{\percent} durant l'any 2008, naturalment deguda a la crisi financera del mateix any. Aquesta crisi va afectar fortament el sector de la construcció, que és una de les principals àrees de negoci de Caterpillar.

\section{Previsió futura}
\begin{figure}[htb]
	\centering \sffamily \small
	% Created by tikzDevice version 0.12 on 2019-06-12 19:23:20
% !TEX encoding = UTF-8 Unicode
\begin{tikzpicture}[x=1pt,y=1pt]
\definecolor{fillColor}{RGB}{255,255,255}
\path[use as bounding box,fill=fillColor,fill opacity=0.00] (0,0) rectangle (433.62,216.81);
\begin{scope}
\path[clip] ( 21.60, 32.40) rectangle (412.02,216.81);
\definecolor{drawColor}{RGB}{169,169,169}

\path[draw=drawColor,line width= 0.4pt,line join=round,line cap=round] ( 36.06, 84.16) -- ( 36.06,192.01);

\path[draw=drawColor,line width= 0.4pt,line join=round,line cap=round] ( 80.75, 84.16) -- ( 80.75,192.01);

\path[draw=drawColor,line width= 0.4pt,line join=round,line cap=round] (127.18, 84.16) -- (127.18,192.01);

\path[draw=drawColor,line width= 0.4pt,line join=round,line cap=round] (171.62, 84.16) -- (171.62,192.01);

\path[draw=drawColor,line width= 0.4pt,line join=round,line cap=round] (217.80, 84.16) -- (217.80,192.01);

\path[draw=drawColor,line width= 0.4pt,line join=round,line cap=round] (262.74, 84.16) -- (262.74,192.01);

\path[draw=drawColor,line width= 0.4pt,line join=round,line cap=round] (308.18, 84.16) -- (308.18,192.01);

\path[draw=drawColor,line width= 0.4pt,line join=round,line cap=round] (353.12, 84.16) -- (353.12,192.01);

\path[draw=drawColor,line width= 0.4pt,line join=round,line cap=round] (397.56, 84.16) -- (397.56,192.01);
\end{scope}
\begin{scope}
\path[clip] (  0.00,  0.00) rectangle (433.62,216.81);
\definecolor{drawColor}{gray}{0.20}

\path[draw=drawColor,line width= 0.4pt,line join=round,line cap=round] ( 36.06, 32.40) -- (397.56, 32.40);

\path[draw=drawColor,line width= 0.6pt,line join=round,line cap=round] ( 36.06, 32.40) -- ( 36.06, 28.08);

\path[draw=drawColor,line width= 0.6pt,line join=round,line cap=round] ( 80.75, 32.40) -- ( 80.75, 28.08);

\path[draw=drawColor,line width= 0.6pt,line join=round,line cap=round] (127.18, 32.40) -- (127.18, 28.08);

\path[draw=drawColor,line width= 0.6pt,line join=round,line cap=round] (171.62, 32.40) -- (171.62, 28.08);

\path[draw=drawColor,line width= 0.6pt,line join=round,line cap=round] (217.80, 32.40) -- (217.80, 28.08);

\path[draw=drawColor,line width= 0.6pt,line join=round,line cap=round] (262.74, 32.40) -- (262.74, 28.08);

\path[draw=drawColor,line width= 0.6pt,line join=round,line cap=round] (308.18, 32.40) -- (308.18, 28.08);

\path[draw=drawColor,line width= 0.6pt,line join=round,line cap=round] (353.12, 32.40) -- (353.12, 28.08);

\path[draw=drawColor,line width= 0.6pt,line join=round,line cap=round] (397.56, 32.40) -- (397.56, 28.08);

\node[text=drawColor,anchor=base,inner sep=0pt, outer sep=0pt, scale=  0.81] at ( 36.06, 17.28) {Jan 02};

\node[text=drawColor,anchor=base,inner sep=0pt, outer sep=0pt, scale=  0.81] at ( 36.06,  7.56) {2015};

\node[text=drawColor,anchor=base,inner sep=0pt, outer sep=0pt, scale=  0.81] at ( 80.75, 17.28) {Jul 01};

\node[text=drawColor,anchor=base,inner sep=0pt, outer sep=0pt, scale=  0.81] at ( 80.75,  7.56) {2015};

\node[text=drawColor,anchor=base,inner sep=0pt, outer sep=0pt, scale=  0.81] at (127.18, 17.28) {Jan 04};

\node[text=drawColor,anchor=base,inner sep=0pt, outer sep=0pt, scale=  0.81] at (127.18,  7.56) {2016};

\node[text=drawColor,anchor=base,inner sep=0pt, outer sep=0pt, scale=  0.81] at (171.62, 17.28) {Jul 01};

\node[text=drawColor,anchor=base,inner sep=0pt, outer sep=0pt, scale=  0.81] at (171.62,  7.56) {2016};

\node[text=drawColor,anchor=base,inner sep=0pt, outer sep=0pt, scale=  0.81] at (217.80, 17.28) {Jan 03};

\node[text=drawColor,anchor=base,inner sep=0pt, outer sep=0pt, scale=  0.81] at (217.80,  7.56) {2017};

\node[text=drawColor,anchor=base,inner sep=0pt, outer sep=0pt, scale=  0.81] at (262.74, 17.28) {Jul 03};

\node[text=drawColor,anchor=base,inner sep=0pt, outer sep=0pt, scale=  0.81] at (262.74,  7.56) {2017};

\node[text=drawColor,anchor=base,inner sep=0pt, outer sep=0pt, scale=  0.81] at (308.18, 17.28) {Jan 02};

\node[text=drawColor,anchor=base,inner sep=0pt, outer sep=0pt, scale=  0.81] at (308.18,  7.56) {2018};

\node[text=drawColor,anchor=base,inner sep=0pt, outer sep=0pt, scale=  0.81] at (353.12, 17.28) {Jul 02};

\node[text=drawColor,anchor=base,inner sep=0pt, outer sep=0pt, scale=  0.81] at (353.12,  7.56) {2018};

\node[text=drawColor,anchor=base,inner sep=0pt, outer sep=0pt, scale=  0.81] at (397.56, 17.28) {Dec 28};

\node[text=drawColor,anchor=base,inner sep=0pt, outer sep=0pt, scale=  0.81] at (397.56,  7.56) {2018};
\end{scope}
\begin{scope}
\path[clip] ( 21.60,192.01) rectangle (412.02,209.98);
\definecolor{drawColor}{gray}{0.20}

\node[text=drawColor,anchor=base west,inner sep=0pt, outer sep=0pt, scale=  0.99] at ( 36.06,198.72) {\bfseries Evolució del preu de tancament};

\node[text=drawColor,anchor=base east,inner sep=0pt, outer sep=0pt, scale=  0.90] at (392.16,198.93) {2015-01-02 / 2018-12-28};
\end{scope}
\begin{scope}
\path[clip] ( 21.60, 84.16) rectangle (412.02,192.01);
\definecolor{drawColor}{RGB}{169,169,169}

\path[draw=drawColor,line width= 0.4pt,line join=round,line cap=round] ( 36.06, 86.16) -- (397.56, 86.16);

\path[draw=drawColor,line width= 0.4pt,line join=round,line cap=round] ( 36.06,105.25) -- (397.56,105.25);

\path[draw=drawColor,line width= 0.4pt,line join=round,line cap=round] ( 36.06,124.34) -- (397.56,124.34);

\path[draw=drawColor,line width= 0.4pt,line join=round,line cap=round] ( 36.06,143.43) -- (397.56,143.43);

\path[draw=drawColor,line width= 0.4pt,line join=round,line cap=round] ( 36.06,162.52) -- (397.56,162.52);

\path[draw=drawColor,line width= 0.4pt,line join=round,line cap=round] ( 36.06,181.61) -- (397.56,181.61);
\end{scope}
\begin{scope}
\path[clip] (  0.00,  0.00) rectangle (433.62,216.81);
\definecolor{drawColor}{gray}{0.20}

\node[text=drawColor,anchor=base east,inner sep=0pt, outer sep=0pt, scale=  0.81] at ( 25.26, 84.30) { 60};

\node[text=drawColor,anchor=base east,inner sep=0pt, outer sep=0pt, scale=  0.81] at ( 25.26,103.39) { 80};

\node[text=drawColor,anchor=base east,inner sep=0pt, outer sep=0pt, scale=  0.81] at ( 25.26,122.48) {100};

\node[text=drawColor,anchor=base east,inner sep=0pt, outer sep=0pt, scale=  0.81] at ( 25.26,141.57) {120};

\node[text=drawColor,anchor=base east,inner sep=0pt, outer sep=0pt, scale=  0.81] at ( 25.26,160.66) {140};

\node[text=drawColor,anchor=base east,inner sep=0pt, outer sep=0pt, scale=  0.81] at ( 25.26,179.75) {160};

\node[text=drawColor,anchor=base west,inner sep=0pt, outer sep=0pt, scale=  0.81] at (408.36, 84.30) { 60};

\node[text=drawColor,anchor=base west,inner sep=0pt, outer sep=0pt, scale=  0.81] at (408.36,103.39) { 80};

\node[text=drawColor,anchor=base west,inner sep=0pt, outer sep=0pt, scale=  0.81] at (408.36,122.48) {100};

\node[text=drawColor,anchor=base west,inner sep=0pt, outer sep=0pt, scale=  0.81] at (408.36,141.57) {120};

\node[text=drawColor,anchor=base west,inner sep=0pt, outer sep=0pt, scale=  0.81] at (408.36,160.66) {140};

\node[text=drawColor,anchor=base west,inner sep=0pt, outer sep=0pt, scale=  0.81] at (408.36,179.75) {160};
\end{scope}
\begin{scope}
\path[clip] ( 21.60, 84.16) rectangle (412.02,192.01);
\definecolor{drawColor}{RGB}{0,0,0}

\path[draw=drawColor,line width= 0.8pt,line join=round] ( 36.06,116.59) --
	( 36.80,111.96) --
	( 37.05,111.43) --
	( 37.30,112.70) --
	( 37.55,113.56) --
	( 37.80,112.55) --
	( 38.54,111.83) --
	( 38.79,111.33) --
	( 39.04,111.04) --
	( 39.29,109.37) --
	( 39.54,108.93) --
	( 40.53,109.04) --
	( 40.78,110.38) --
	( 41.03,111.76) --
	( 41.27,110.60) --
	( 42.02,111.01) --
	( 42.27,105.11) --
	( 42.52,105.26) --
	( 42.76,105.25) --
	( 43.01,105.22) --
	( 43.76,106.05) --
	( 44.01,108.99) --
	( 44.25,107.11) --
	( 44.50,108.66) --
	( 44.75,108.31) --
	( 45.49,109.71) --
	( 45.74,108.97) --
	( 45.99,108.50) --
	( 46.24,108.62) --
	( 46.49,110.15) --
	( 47.48,110.38) --
	( 47.73,109.83) --
	( 47.98,108.40) --
	( 48.23,109.49) --
	( 48.97,108.58) --
	( 49.22,109.55) --
	( 49.47,110.17) --
	( 49.72,108.88) --
	( 49.96,108.02) --
	( 50.71,108.11) --
	( 50.96,107.96) --
	( 51.21,106.90) --
	( 51.45,106.69) --
	( 51.70,105.31) --
	( 52.45,106.02) --
	( 52.69,105.16) --
	( 52.94,104.96) --
	( 53.19,105.43) --
	( 53.44,104.51) --
	( 54.18,105.22) --
	( 54.43,103.77) --
	( 54.68,106.52) --
	( 54.93,105.34) --
	( 55.18,106.39) --
	( 55.92,106.29) --
	( 56.17,105.71) --
	( 56.42,105.12) --
	( 56.67,105.14) --
	( 56.92,104.93) --
	( 57.66,106.56) --
	( 57.91,105.28) --
	( 58.16,104.91) --
	( 58.41,105.48) --
	( 59.40,106.67) --
	( 59.65,105.77) --
	( 59.90,105.67) --
	( 60.14,105.90) --
	( 60.39,107.73) --
	( 61.14,107.29) --
	( 61.38,108.11) --
	( 61.63,110.17) --
	( 61.88,109.39) --
	( 62.13,108.38) --
	( 62.87,109.64) --
	( 63.12,108.99) --
	( 63.37,109.90) --
	( 63.62,109.82) --
	( 63.87,109.64) --
	( 64.61,110.34) --
	( 64.86,111.17) --
	( 65.11,112.41) --
	( 65.36,111.82) --
	( 65.61,112.28) --
	( 66.35,112.22) --
	( 66.60,111.93) --
	( 66.85,111.92) --
	( 67.10,111.39) --
	( 67.34,112.23) --
	( 68.09,113.62) --
	( 68.34,113.04) --
	( 68.59,113.31) --
	( 68.83,113.59) --
	( 69.08,113.30) --
	( 69.83,113.16) --
	( 70.07,112.16) --
	( 70.32,112.82) --
	( 70.57,114.16) --
	( 70.82,113.49) --
	( 71.81,112.73) --
	( 72.06,112.81) --
	( 72.31,110.99) --
	( 72.56,110.33) --
	( 73.30,110.53) --
	( 73.55,111.14) --
	( 73.80,111.20) --
	( 74.05,110.47) --
	( 74.30,111.02) --
	( 75.04,111.12) --
	( 75.29,111.67) --
	( 75.54,113.34) --
	( 75.79,113.05) --
	( 76.03,112.74) --
	( 76.78,112.06) --
	( 77.03,111.54) --
	( 77.27,112.17) --
	( 77.52,112.35) --
	( 77.77,112.43) --
	( 78.52,113.14) --
	( 78.76,113.32) --
	( 79.01,113.05) --
	( 79.26,111.66) --
	( 79.51,111.76) --
	( 80.25,110.31) --
	( 80.50,109.85) --
	( 80.75,109.02) --
	( 81.00,109.43) --
	( 81.99,108.26) --
	( 82.24,109.06) --
	( 82.49,107.42) --
	( 82.74,106.86) --
	( 82.99,107.30) --
	( 83.73,108.72) --
	( 83.98,109.51) --
	( 84.23,109.22) --
	( 84.48,108.84) --
	( 84.72,108.27) --
	( 85.47,107.40) --
	( 85.72,107.37) --
	( 85.96,105.02) --
	( 86.21,102.27) --
	( 86.46,101.53) --
	( 87.21,100.78) --
	( 87.45,103.13) --
	( 87.70,102.70) --
	( 87.95,103.67) --
	( 88.20,103.94) --
	( 88.94,102.63) --
	( 89.19,102.00) --
	( 89.44,102.27) --
	( 89.69,102.83) --
	( 89.94,102.66) --
	( 90.68,105.40) --
	( 90.93,103.38) --
	( 91.18,104.22) --
	( 91.43,103.11) --
	( 91.68,103.81) --
	( 92.42,103.86) --
	( 92.67,104.12) --
	( 92.92,102.78) --
	( 93.17,101.82) --
	( 93.41,100.48) --
	( 94.16, 98.40) --
	( 94.41, 97.67) --
	( 94.65, 99.40) --
	( 94.90,101.11) --
	( 95.15,101.38) --
	( 95.90,101.85) --
	( 96.14,100.38) --
	( 96.39,101.53) --
	( 96.64, 99.95) --
	( 96.89, 98.66) --
	( 97.88, 99.81) --
	( 98.13, 98.53) --
	( 98.38, 98.01) --
	( 98.63, 98.21) --
	( 99.37, 98.35) --
	( 99.62,100.08) --
	( 99.87,101.09) --
	(100.12, 99.58) --
	(100.37, 97.48) --
	(101.11, 97.77) --
	(101.36, 97.31) --
	(101.61, 95.90) --
	(101.85, 91.70) --
	(102.10, 90.91) --
	(102.85, 89.78) --
	(103.10, 90.27) --
	(103.34, 91.28) --
	(103.59, 90.35) --
	(103.84, 91.60) --
	(104.59, 94.92) --
	(104.83, 96.49) --
	(105.08, 96.07) --
	(105.33, 97.46) --
	(105.58, 96.95) --
	(106.32, 96.18) --
	(106.57, 95.87) --
	(106.82, 96.38) --
	(107.07, 96.50) --
	(107.32, 95.40) --
	(108.06, 95.01) --
	(108.31, 95.96) --
	(108.56, 95.39) --
	(108.81, 96.54) --
	(109.06, 97.37) --
	(109.80, 97.16) --
	(110.05, 96.08) --
	(110.30, 97.58) --
	(110.54, 97.58) --
	(110.79, 98.56) --
	(111.54, 99.85) --
	(111.79,100.24) --
	(112.03,100.05) --
	(112.28, 99.73) --
	(112.53, 99.37) --
	(113.28, 97.51) --
	(113.52, 98.02) --
	(113.77, 97.53) --
	(114.02, 94.43) --
	(114.27, 95.35) --
	(115.01, 96.08) --
	(115.26, 95.12) --
	(115.51, 96.01) --
	(115.76, 95.72) --
	(116.01, 96.79) --
	(116.75, 96.68) --
	(117.00, 97.04) --
	(117.25, 97.13) --
	(117.75, 96.87) --
	(118.49, 98.23) --
	(118.74, 97.19) --
	(118.99, 96.69) --
	(119.23, 95.36) --
	(119.48, 95.69) --
	(120.23, 94.14) --
	(120.48, 92.40) --
	(120.72, 91.80) --
	(120.97, 92.25) --
	(121.22, 91.34) --
	(121.97, 92.00) --
	(122.21, 92.60) --
	(122.46, 93.53) --
	(122.71, 90.84) --
	(122.96, 91.04) --
	(123.70, 91.16) --
	(123.95, 94.19) --
	(124.20, 95.60) --
	(124.45, 95.11) --
	(125.44, 94.35) --
	(125.69, 94.93) --
	(125.94, 94.45) --
	(126.19, 93.76) --
	(127.18, 93.79) --
	(127.43, 93.11) --
	(127.68, 92.10) --
	(127.92, 89.92) --
	(128.17, 89.30) --
	(128.92, 87.56) --
	(129.17, 87.69) --
	(129.41, 87.01) --
	(129.66, 88.33) --
	(129.91, 86.04) --
	(130.90, 85.23) --
	(131.15, 85.02) --
	(131.40, 85.86) --
	(131.65, 87.09) --
	(132.39, 84.16) --
	(132.64, 85.36) --
	(132.89, 84.56) --
	(133.14, 87.19) --
	(133.39, 88.30) --
	(134.13, 88.74) --
	(134.38, 86.80) --
	(134.63, 89.28) --
	(134.88, 91.85) --
	(135.12, 92.00) --
	(135.87, 90.53) --
	(136.12, 89.91) --
	(136.37, 88.20) --
	(136.61, 87.51) --
	(136.86, 89.17) --
	(137.86, 91.13) --
	(138.10, 93.09) --
	(138.35, 92.00) --
	(138.60, 91.33) --
	(139.35, 93.14) --
	(139.59, 91.68) --
	(139.84, 91.77) --
	(140.09, 92.33) --
	(140.34, 92.72) --
	(141.08, 93.51) --
	(141.33, 94.84) --
	(141.58, 95.11) --
	(141.83, 97.37) --
	(142.08, 98.42) --
	(142.82,100.26) --
	(143.07, 97.36) --
	(143.32, 97.49) --
	(143.57, 97.00) --
	(143.81, 98.38) --
	(144.56, 98.31) --
	(144.81, 98.03) --
	(145.06, 99.85) --
	(145.30,101.34) --
	(145.55,100.93) --
	(146.30,101.33) --
	(146.55,101.00) --
	(146.79, 99.22) --
	(147.04,100.75) --
	(148.04,100.78) --
	(148.28,101.44) --
	(148.53,101.92) --
	(148.78,101.95) --
	(149.03,102.19) --
	(149.77,101.16) --
	(150.02,100.71) --
	(150.27,100.69) --
	(150.52, 99.68) --
	(150.77, 99.86) --
	(151.51,100.12) --
	(151.76,101.53) --
	(152.01,104.42) --
	(152.26,104.34) --
	(152.50,104.46) --
	(153.25,104.53) --
	(153.50,105.62) --
	(153.75,105.00) --
	(153.99,103.97) --
	(154.24,103.65) --
	(154.99,102.19) --
	(155.24,103.01) --
	(155.48,103.99) --
	(155.73,103.10) --
	(155.98,103.07) --
	(156.73,103.19) --
	(156.97,101.78) --
	(157.22, 99.75) --
	(157.47, 98.37) --
	(157.72, 98.91) --
	(158.46, 96.45) --
	(158.71, 98.10) --
	(158.96, 98.02) --
	(159.21, 97.33) --
	(159.46, 95.77) --
	(160.20, 96.37) --
	(160.45, 96.71) --
	(160.70, 96.27) --
	(160.95, 95.16) --
	(161.19, 95.58) --
	(161.94, 96.09) --
	(162.19, 96.74) --
	(162.44, 98.16) --
	(162.68, 97.69) --
	(162.93, 97.58) --
	(163.93, 98.10) --
	(164.17, 97.87) --
	(164.42, 99.16) --
	(164.67,100.52) --
	(165.42,101.83) --
	(165.66,102.20) --
	(165.91,103.45) --
	(166.16,102.54) --
	(166.41,101.46) --
	(167.15,100.70) --
	(167.40,100.34) --
	(167.65,100.54) --
	(167.90,100.47) --
	(168.15,101.36) --
	(168.89,101.84) --
	(169.14,101.90) --
	(169.39,101.84) --
	(169.64,103.55) --
	(169.88, 98.60) --
	(170.63, 97.02) --
	(170.88, 98.11) --
	(171.13, 99.74) --
	(171.37,101.25) --
	(171.62,101.86) --
	(172.62, 99.89) --
	(172.86,100.31) --
	(173.11,100.52) --
	(173.36,102.74) --
	(174.11,103.15) --
	(174.35,105.06) --
	(174.60,104.95) --
	(174.85,105.31) --
	(175.10,105.92) --
	(175.84,105.08) --
	(176.09,105.00) --
	(176.34,104.95) --
	(176.59,105.26) --
	(176.84,104.66) --
	(177.58,104.00) --
	(177.83,107.87) --
	(178.08,109.16) --
	(178.33,108.04) --
	(178.57,107.88) --
	(179.32,107.11) --
	(179.57,106.43) --
	(179.82,107.19) --
	(180.06,106.94) --
	(180.31,107.69) --
	(181.06,108.53) --
	(181.31,107.95) --
	(181.55,107.78) --
	(181.80,108.34) --
	(182.05,108.11) --
	(182.80,109.21) --
	(183.04,109.34) --
	(183.29,109.46) --
	(183.54,108.48) --
	(183.79,108.92) --
	(184.53,108.68) --
	(184.78,109.08) --
	(185.03,108.26) --
	(185.28,107.96) --
	(185.53,107.79) --
	(186.27,108.21) --
	(186.52,107.62) --
	(186.77,107.11) --
	(187.02,106.63) --
	(187.26,106.86) --
	(188.26,107.24) --
	(188.51,108.80) --
	(188.75,108.63) --
	(189.00,106.00) --
	(189.75,107.08) --
	(190.00,105.47) --
	(190.24,106.22) --
	(190.49,107.19) --
	(190.74,107.21) --
	(191.49,107.04) --
	(191.73,106.87) --
	(191.98,108.56) --
	(192.23,108.59) --
	(192.48,107.58) --
	(193.22,107.51) --
	(193.47,108.00) --
	(193.72,111.54) --
	(193.97,112.38) --
	(194.22,113.62) --
	(194.96,113.15) --
	(195.21,112.42) --
	(195.46,114.24) --
	(195.71,114.12) --
	(195.95,113.33) --
	(196.70,113.10) --
	(196.95,112.38) --
	(197.20,112.48) --
	(197.44,111.90) --
	(197.69,112.57) --
	(198.44,112.21) --
	(198.69,112.14) --
	(198.93,112.89) --
	(199.18,111.58) --
	(199.43,111.29) --
	(200.18,110.97) --
	(200.42,109.53) --
	(200.67,109.19) --
	(200.92,108.12) --
	(201.17,108.95) --
	(201.91,108.55) --
	(202.16,107.39) --
	(202.41,106.31) --
	(202.66,106.46) --
	(202.91,107.45) --
	(203.65,108.83) --
	(203.90,109.72) --
	(204.15,115.94) --
	(204.40,118.09) --
	(204.64,117.67) --
	(205.39,118.78) --
	(205.64,119.03) --
	(205.89,117.99) --
	(206.13,117.45) --
	(206.38,117.03) --
	(207.13,117.56) --
	(207.38,118.25) --
	(207.62,120.69) --
	(208.12,120.34) --
	(208.86,119.47) --
	(209.11,118.65) --
	(209.36,120.10) --
	(209.61,120.75) --
	(209.86,119.70) --
	(210.60,119.04) --
	(210.85,119.78) --
	(211.10,121.79) --
	(211.35,120.78) --
	(211.60,120.07) --
	(212.34,119.64) --
	(212.59,120.54) --
	(212.84,118.36) --
	(213.09,119.12) --
	(213.33,117.26) --
	(214.08,117.38) --
	(214.33,118.94) --
	(214.58,118.44) --
	(214.82,118.73) --
	(215.07,118.92) --
	(216.07,119.14) --
	(216.31,117.80) --
	(216.56,117.94) --
	(216.81,117.41) --
	(217.80,118.60) --
	(218.05,118.20) --
	(218.30,117.66) --
	(218.55,117.70) --
	(219.29,117.06) --
	(219.54,118.45) --
	(219.79,119.23) --
	(220.04,118.60) --
	(220.29,119.07) --
	(221.28,118.20) --
	(221.53,117.97) --
	(221.78,118.02) --
	(222.02,119.17) --
	(222.77,119.05) --
	(223.02,120.75) --
	(223.27,122.57) --
	(223.51,121.69) --
	(223.76,123.38) --
	(224.51,121.28) --
	(224.76,120.20) --
	(225.00,119.67) --
	(225.25,118.39) --
	(225.50,117.93) --
	(226.24,117.53) --
	(226.49,117.95) --
	(226.74,117.57) --
	(226.99,118.57) --
	(227.24,120.82) --
	(227.98,122.91) --
	(228.23,122.52) --
	(228.48,123.40) --
	(228.73,122.68) --
	(228.98,123.24) --
	(229.97,122.53) --
	(230.22,122.62) --
	(230.47,120.09) --
	(230.71,120.03) --
	(231.46,121.90) --
	(231.71,121.15) --
	(231.96,122.98) --
	(232.20,118.96) --
	(232.45,119.68) --
	(233.20,120.20) --
	(233.44,120.46) --
	(233.69,117.88) --
	(233.94,116.12) --
	(234.19,117.00) --
	(234.93,117.31) --
	(235.18,116.57) --
	(235.43,118.00) --
	(235.68,117.52) --
	(235.93,117.57) --
	(236.67,119.95) --
	(236.92,117.11) --
	(237.17,117.11) --
	(237.42,116.92) --
	(237.67,116.85) --
	(238.41,116.24) --
	(238.66,117.61) --
	(238.91,117.52) --
	(239.16,117.87) --
	(239.40,117.43) --
	(240.15,116.96) --
	(240.40,118.74) --
	(240.65,118.84) --
	(240.89,120.35) --
	(241.14,120.06) --
	(241.89,121.61) --
	(242.13,121.57) --
	(242.38,119.43) --
	(242.63,117.75) --
	(243.62,118.75) --
	(243.87,118.99) --
	(244.12,118.11) --
	(244.37,119.24) --
	(244.62,118.92) --
	(245.36,121.30) --
	(245.61,128.56) --
	(245.86,128.79) --
	(246.11,126.90) --
	(246.36,126.50) --
	(247.10,126.25) --
	(247.35,125.78) --
	(247.60,125.81) --
	(247.85,123.76) --
	(248.09,124.00) --
	(248.84,123.74) --
	(249.09,123.66) --
	(249.34,124.68) --
	(249.58,125.29) --
	(249.83,125.02) --
	(250.58,126.65) --
	(250.82,126.57) --
	(251.07,124.47) --
	(251.32,124.55) --
	(251.57,126.66) --
	(252.31,126.53) --
	(252.56,127.81) --
	(252.81,128.11) --
	(253.06,129.13) --
	(253.31,129.74) --
	(254.30,129.56) --
	(254.55,129.52) --
	(254.80,129.73) --
	(255.05,130.02) --
	(255.79,129.30) --
	(256.04,128.68) --
	(256.29,127.69) --
	(256.54,129.12) --
	(256.78,129.99) --
	(257.53,130.13) --
	(257.78,129.83) --
	(258.02,128.84) --
	(258.27,130.45) --
	(258.52,131.59) --
	(259.27,131.49) --
	(259.51,131.06) --
	(259.76,127.65) --
	(260.01,128.01) --
	(260.26,128.26) --
	(261.00,128.40) --
	(261.25,128.09) --
	(261.50,130.50) --
	(261.75,129.72) --
	(262.00,131.46) --
	(262.74,130.99) --
	(263.24,131.49) --
	(263.49,130.55) --
	(263.74,130.95) --
	(264.48,131.82) --
	(264.73,132.21) --
	(264.98,132.89) --
	(265.23,132.42) --
	(265.47,132.71) --
	(266.22,132.78) --
	(266.47,131.09) --
	(266.71,131.83) --
	(266.96,131.33) --
	(267.21,130.63) --
	(267.96,132.15) --
	(268.20,138.22) --
	(268.45,137.25) --
	(268.70,138.06) --
	(268.95,137.80) --
	(269.69,137.66) --
	(269.94,136.84) --
	(270.19,136.83) --
	(270.44,137.28) --
	(270.69,138.04) --
	(271.43,138.38) --
	(271.68,138.09) --
	(271.93,138.03) --
	(272.18,136.89) --
	(272.43,136.61) --
	(273.17,137.43) --
	(273.42,137.37) --
	(273.67,137.82) --
	(273.92,136.88) --
	(274.16,137.63) --
	(274.91,137.32) --
	(275.16,139.44) --
	(275.40,138.42) --
	(275.65,138.45) --
	(275.90,138.99) --
	(276.65,138.72) --
	(276.89,139.62) --
	(277.14,141.09) --
	(277.39,141.03) --
	(277.64,141.79) --
	(278.63,141.81) --
	(278.88,140.55) --
	(279.13,141.30) --
	(279.38,141.35) --
	(280.12,142.35) --
	(280.37,144.33) --
	(280.62,143.08) --
	(280.87,143.89) --
	(281.12,144.74) --
	(281.86,147.09) --
	(282.11,147.96) --
	(282.36,148.06) --
	(282.60,148.00) --
	(282.85,147.66) --
	(283.60,147.55) --
	(283.85,147.74) --
	(284.09,147.75) --
	(284.34,148.42) --
	(284.59,147.93) --
	(285.34,147.94) --
	(285.58,148.69) --
	(285.83,149.86) --
	(286.08,149.54) --
	(286.33,150.05) --
	(287.07,150.00) --
	(287.32,151.25) --
	(287.57,151.64) --
	(287.82,152.97) --
	(288.07,153.65) --
	(288.81,154.38) --
	(289.06,153.49) --
	(289.31,154.21) --
	(289.56,154.46) --
	(289.81,154.27) --
	(290.55,154.58) --
	(290.80,160.84) --
	(291.05,159.50) --
	(291.29,159.60) --
	(291.54,160.43) --
	(292.29,159.17) --
	(292.54,158.51) --
	(292.78,158.98) --
	(293.03,159.15) --
	(293.28,159.30) --
	(294.03,160.34) --
	(294.27,161.39) --
	(294.52,159.93) --
	(294.77,158.07) --
	(295.02,159.16) --
	(295.76,159.21) --
	(296.01,160.17) --
	(296.26,156.89) --
	(296.51,159.05) --
	(296.76,158.83) --
	(297.50,159.57) --
	(297.75,160.23) --
	(298.00,160.62) --
	(298.50,160.03) --
	(299.24,159.84) --
	(299.49,161.56) --
	(299.74,160.67) --
	(299.98,163.62) --
	(300.23,163.97) --
	(300.98,163.95) --
	(301.23,162.65) --
	(301.47,162.85) --
	(301.72,165.28) --
	(301.97,166.21) --
	(302.72,166.12) --
	(302.96,165.79) --
	(303.21,170.70) --
	(303.46,168.52) --
	(303.71,168.91) --
	(304.45,171.85) --
	(304.70,172.93) --
	(304.95,174.47) --
	(305.20,176.50) --
	(305.45,177.55) --
	(306.44,178.21) --
	(306.69,179.24) --
	(306.94,180.10) --
	(307.19,179.30) --
	(308.18,178.79) --
	(308.43,179.02) --
	(308.67,181.08) --
	(308.92,183.48) --
	(309.67,187.37) --
	(309.92,187.75) --
	(310.16,187.21) --
	(310.41,190.39) --
	(310.66,191.44) --
	(311.65,190.50) --
	(311.90,189.73) --
	(312.15,190.04) --
	(312.40,191.55) --
	(313.14,192.01) --
	(313.39,190.61) --
	(313.64,189.57) --
	(313.89,190.56) --
	(314.14,188.35) --
	(314.88,184.07) --
	(315.13,185.20) --
	(315.38,184.27) --
	(315.63,183.75) --
	(315.87,179.22) --
	(316.62,173.10) --
	(316.87,178.18) --
	(317.12,176.21) --
	(317.36,168.24) --
	(317.61,171.31) --
	(318.36,174.25) --
	(318.61,175.86) --
	(318.85,179.76) --
	(319.10,181.59) --
	(319.35,178.07) --
	(320.34,177.37) --
	(320.59,177.06) --
	(320.84,180.52) --
	(321.09,183.91) --
	(321.83,185.13) --
	(322.08,182.81) --
	(322.33,176.49) --
	(322.58,172.29) --
	(322.83,168.61) --
	(323.57,173.14) --
	(323.82,175.65) --
	(324.07,173.51) --
	(324.32,175.49) --
	(324.56,179.94) --
	(325.31,176.36) --
	(325.56,175.59) --
	(325.81,174.49) --
	(326.05,176.43) --
	(326.30,178.23) --
	(327.05,174.12) --
	(327.30,175.94) --
	(327.54,177.60) --
	(327.79,169.11) --
	(328.04,166.62) --
	(328.79,171.29) --
	(329.03,169.19) --
	(329.28,167.45) --
	(329.53,169.57) --
	(330.52,166.23) --
	(330.77,167.35) --
	(331.02,167.47) --
	(331.27,170.28) --
	(331.52,165.38) --
	(332.26,165.46) --
	(332.51,170.24) --
	(332.76,169.00) --
	(333.01,172.20) --
	(333.25,172.29) --
	(334.00,174.11) --
	(334.25,175.23) --
	(334.50,177.75) --
	(334.74,177.51) --
	(334.99,175.17) --
	(335.74,175.87) --
	(335.99,166.76) --
	(336.23,167.48) --
	(336.48,168.54) --
	(336.73,167.03) --
	(337.48,166.68) --
	(337.72,166.74) --
	(337.97,166.21) --
	(338.22,166.13) --
	(338.47,168.60) --
	(339.21,171.89) --
	(339.46,172.00) --
	(339.71,174.56) --
	(339.96,175.86) --
	(340.21,176.71) --
	(340.95,177.11) --
	(341.20,174.54) --
	(341.45,175.69) --
	(341.70,175.59) --
	(341.94,177.52) --
	(342.69,180.58) --
	(342.94,177.98) --
	(343.19,179.45) --
	(343.43,179.42) --
	(343.68,177.65) --
	(344.68,175.52) --
	(344.92,177.28) --
	(345.17,173.89) --
	(345.42,175.43) --
	(346.17,174.45) --
	(346.41,174.94) --
	(346.66,178.28) --
	(346.91,178.07) --
	(347.16,177.14) --
	(347.90,178.08) --
	(348.15,179.22) --
	(348.40,176.56) --
	(348.65,175.06) --
	(348.90,172.09) --
	(349.64,170.81) --
	(349.89,165.67) --
	(350.14,165.51) --
	(350.39,162.07) --
	(350.63,162.46) --
	(351.38,159.26) --
	(351.63,158.26) --
	(351.88,157.38) --
	(352.12,158.73) --
	(352.37,158.39) --
	(353.12,158.98) --
	(353.37,156.63) --
	(353.86,158.52) --
	(354.11,158.14) --
	(354.86,163.45) --
	(355.10,163.71) --
	(355.35,159.43) --
	(355.60,161.97) --
	(355.85,163.24) --
	(356.59,160.69) --
	(356.84,161.52) --
	(357.09,163.59) --
	(357.34,161.97) --
	(357.59,159.49) --
	(358.33,158.97) --
	(358.58,160.58) --
	(358.83,162.98) --
	(359.08,164.98) --
	(359.32,164.96) --
	(360.07,162.28) --
	(360.32,166.15) --
	(360.57,161.13) --
	(360.81,160.62) --
	(361.06,161.10) --
	(361.81,162.02) --
	(362.06,164.79) --
	(362.30,162.24) --
	(362.55,161.11) --
	(362.80,158.63) --
	(363.55,157.77) --
	(363.79,157.67) --
	(364.04,154.90) --
	(364.29,158.95) --
	(364.54,161.89) --
	(365.28,161.25) --
	(365.53,162.51) --
	(365.78,162.17) --
	(366.03,159.46) --
	(366.28,160.81) --
	(367.02,164.47) --
	(367.27,164.43) --
	(367.52,164.30) --
	(367.77,161.62) --
	(368.01,161.42) --
	(369.01,160.84) --
	(369.26,163.08) --
	(369.50,164.02) --
	(369.75,163.78) --
	(370.50,164.48) --
	(370.75,164.46) --
	(370.99,166.61) --
	(371.24,167.81) --
	(371.49,167.20) --
	(372.24,168.46) --
	(372.48,171.18) --
	(372.73,174.70) --
	(372.98,177.79) --
	(373.23,178.16) --
	(373.97,175.88) --
	(374.22,175.97) --
	(374.47,175.09) --
	(374.72,174.42) --
	(374.97,174.44) --
	(375.71,174.24) --
	(375.96,176.67) --
	(376.21,179.91) --
	(376.46,178.51) --
	(376.70,175.23) --
	(377.45,175.26) --
	(377.70,171.55) --
	(377.95,166.07) --
	(378.19,163.45) --
	(378.44,164.50) --
	(379.19,164.24) --
	(379.44,166.06) --
	(379.68,163.80) --
	(379.93,158.51) --
	(380.18,154.24) --
	(380.93,151.74) --
	(381.17,142.46) --
	(381.42,136.12) --
	(381.67,139.26) --
	(381.92,138.71) --
	(382.66,137.68) --
	(382.91,141.01) --
	(383.16,144.69) --
	(383.41,147.92) --
	(383.66,148.96) --
	(384.40,149.55) --
	(384.65,152.34) --
	(384.90,157.84) --
	(385.15,154.30) --
	(385.39,150.03) --
	(386.14,148.20) --
	(386.39,149.12) --
	(386.64,148.30) --
	(386.88,152.42) --
	(387.13,152.94) --
	(387.88,149.14) --
	(388.13,145.60) --
	(388.37,147.12) --
	(388.87,145.65) --
	(389.61,148.01) --
	(389.86,147.82) --
	(390.11,153.71) --
	(390.36,153.20) --
	(390.61,158.39) --
	(391.35,161.52) --
	(391.60,152.33) --
	(392.10,151.42) --
	(392.35,146.82) --
	(393.09,146.67) --
	(393.34,146.52) --
	(393.59,148.56) --
	(393.84,149.06) --
	(394.08,149.89) --
	(394.83,147.70) --
	(395.08,147.51) --
	(395.33,145.65) --
	(395.57,144.91) --
	(395.82,143.50) --
	(396.57,140.52) --
	(397.06,147.97) --
	(397.31,149.80) --
	(397.56,148.79);
\end{scope}
\begin{scope}
\path[clip] ( 21.60, 39.23) rectangle (412.02, 75.18);
\definecolor{drawColor}{RGB}{169,169,169}

\path[draw=drawColor,line width= 0.4pt,line join=round,line cap=round] ( 36.06, 32.40) -- ( 36.06,216.81);

\path[draw=drawColor,line width= 0.4pt,line join=round,line cap=round] ( 80.75, 32.40) -- ( 80.75,216.81);

\path[draw=drawColor,line width= 0.4pt,line join=round,line cap=round] (127.18, 32.40) -- (127.18,216.81);

\path[draw=drawColor,line width= 0.4pt,line join=round,line cap=round] (171.62, 32.40) -- (171.62,216.81);

\path[draw=drawColor,line width= 0.4pt,line join=round,line cap=round] (217.80, 32.40) -- (217.80,216.81);

\path[draw=drawColor,line width= 0.4pt,line join=round,line cap=round] (262.74, 32.40) -- (262.74,216.81);

\path[draw=drawColor,line width= 0.4pt,line join=round,line cap=round] (308.18, 32.40) -- (308.18,216.81);

\path[draw=drawColor,line width= 0.4pt,line join=round,line cap=round] (353.12, 32.40) -- (353.12,216.81);

\path[draw=drawColor,line width= 0.4pt,line join=round,line cap=round] (397.56, 32.40) -- (397.56,216.81);
\definecolor{drawColor}{RGB}{255,0,0}

\path[draw=drawColor,line width= 0.8pt,line join=round] ( 36.80, 57.55) -- ( 36.80, 44.91);

\path[draw=drawColor,line width= 0.8pt,line join=round] ( 37.05, 57.55) -- ( 37.05, 56.04);

\path[draw=drawColor,line width= 0.8pt,line join=round] ( 37.30, 57.55) -- ( 37.30, 61.13);

\path[draw=drawColor,line width= 0.8pt,line join=round] ( 37.55, 57.55) -- ( 37.55, 59.92);

\path[draw=drawColor,line width= 0.8pt,line join=round] ( 37.80, 57.55) -- ( 37.80, 54.74);

\path[draw=drawColor,line width= 0.8pt,line join=round] ( 38.54, 57.55) -- ( 38.54, 55.52);

\path[draw=drawColor,line width= 0.8pt,line join=round] ( 38.79, 57.55) -- ( 38.79, 56.15);

\path[draw=drawColor,line width= 0.8pt,line join=round] ( 39.04, 57.55) -- ( 39.04, 56.73);

\path[draw=drawColor,line width= 0.8pt,line join=round] ( 39.29, 57.55) -- ( 39.29, 52.76);

\path[draw=drawColor,line width= 0.8pt,line join=round] ( 39.54, 57.55) -- ( 39.54, 56.27);

\path[draw=drawColor,line width= 0.8pt,line join=round] ( 40.53, 57.55) -- ( 40.53, 57.85);

\path[draw=drawColor,line width= 0.8pt,line join=round] ( 40.78, 57.55) -- ( 40.78, 61.40);

\path[draw=drawColor,line width= 0.8pt,line join=round] ( 41.03, 57.55) -- ( 41.03, 61.47);

\path[draw=drawColor,line width= 0.8pt,line join=round] ( 41.27, 57.55) -- ( 41.27, 54.28);

\path[draw=drawColor,line width= 0.8pt,line join=round] ( 42.02, 57.55) -- ( 42.02, 58.69);

\path[draw=drawColor,line width= 0.8pt,line join=round] ( 42.27, 57.55) -- ( 42.27, 40.18);

\path[draw=drawColor,line width= 0.8pt,line join=round] ( 42.52, 57.55) -- ( 42.52, 58.01);

\path[draw=drawColor,line width= 0.8pt,line join=round] ( 42.76, 57.55) -- ( 42.76, 57.52);

\path[draw=drawColor,line width= 0.8pt,line join=round] ( 43.01, 57.55) -- ( 43.01, 57.46);

\path[draw=drawColor,line width= 0.8pt,line join=round] ( 43.76, 57.55) -- ( 43.76, 60.07);

\path[draw=drawColor,line width= 0.8pt,line join=round] ( 44.01, 57.55) -- ( 44.01, 66.26);

\path[draw=drawColor,line width= 0.8pt,line join=round] ( 44.25, 57.55) -- ( 44.25, 52.01);

\path[draw=drawColor,line width= 0.8pt,line join=round] ( 44.50, 57.55) -- ( 44.50, 62.11);

\path[draw=drawColor,line width= 0.8pt,line join=round] ( 44.75, 57.55) -- ( 44.75, 56.54);

\path[draw=drawColor,line width= 0.8pt,line join=round] ( 45.49, 57.55) -- ( 45.49, 61.60);

\path[draw=drawColor,line width= 0.8pt,line join=round] ( 45.74, 57.55) -- ( 45.74, 55.42);

\path[draw=drawColor,line width= 0.8pt,line join=round] ( 45.99, 57.55) -- ( 45.99, 56.18);

\path[draw=drawColor,line width= 0.8pt,line join=round] ( 46.24, 57.55) -- ( 46.24, 57.88);

\path[draw=drawColor,line width= 0.8pt,line join=round] ( 46.49, 57.55) -- ( 46.49, 61.97);

\path[draw=drawColor,line width= 0.8pt,line join=round] ( 47.48, 57.55) -- ( 47.48, 58.23);

\path[draw=drawColor,line width= 0.8pt,line join=round] ( 47.73, 57.55) -- ( 47.73, 55.96);

\path[draw=drawColor,line width= 0.8pt,line join=round] ( 47.98, 57.55) -- ( 47.98, 53.39);

\path[draw=drawColor,line width= 0.8pt,line join=round] ( 48.23, 57.55) -- ( 48.23, 60.71);

\path[draw=drawColor,line width= 0.8pt,line join=round] ( 48.97, 57.55) -- ( 48.97, 54.91);

\path[draw=drawColor,line width= 0.8pt,line join=round] ( 49.22, 57.55) -- ( 49.22, 60.35);

\path[draw=drawColor,line width= 0.8pt,line join=round] ( 49.47, 57.55) -- ( 49.47, 59.33);

\path[draw=drawColor,line width= 0.8pt,line join=round] ( 49.72, 57.55) -- ( 49.72, 53.82);

\path[draw=drawColor,line width= 0.8pt,line join=round] ( 49.96, 57.55) -- ( 49.96, 55.03);

\path[draw=drawColor,line width= 0.8pt,line join=round] ( 50.71, 57.55) -- ( 50.71, 57.83);

\path[draw=drawColor,line width= 0.8pt,line join=round] ( 50.96, 57.55) -- ( 50.96, 57.10);

\path[draw=drawColor,line width= 0.8pt,line join=round] ( 51.21, 57.55) -- ( 51.21, 54.40);

\path[draw=drawColor,line width= 0.8pt,line join=round] ( 51.45, 57.55) -- ( 51.45, 56.92);

\path[draw=drawColor,line width= 0.8pt,line join=round] ( 51.70, 57.55) -- ( 51.70, 53.36);

\path[draw=drawColor,line width= 0.8pt,line join=round] ( 52.45, 57.55) -- ( 52.45, 59.72);

\path[draw=drawColor,line width= 0.8pt,line join=round] ( 52.69, 57.55) -- ( 52.69, 54.94);

\path[draw=drawColor,line width= 0.8pt,line join=round] ( 52.94, 57.55) -- ( 52.94, 56.93);

\path[draw=drawColor,line width= 0.8pt,line join=round] ( 53.19, 57.55) -- ( 53.19, 58.97);

\path[draw=drawColor,line width= 0.8pt,line join=round] ( 53.44, 57.55) -- ( 53.44, 54.74);

\path[draw=drawColor,line width= 0.8pt,line join=round] ( 54.18, 57.55) -- ( 54.18, 59.71);

\path[draw=drawColor,line width= 0.8pt,line join=round] ( 54.43, 57.55) -- ( 54.43, 53.07);

\path[draw=drawColor,line width= 0.8pt,line join=round] ( 54.68, 57.55) -- ( 54.68, 65.95);

\path[draw=drawColor,line width= 0.8pt,line join=round] ( 54.93, 57.55) -- ( 54.93, 53.97);

\path[draw=drawColor,line width= 0.8pt,line join=round] ( 55.18, 57.55) -- ( 55.18, 60.72);

\path[draw=drawColor,line width= 0.8pt,line join=round] ( 55.92, 57.55) -- ( 55.92, 57.26);

\path[draw=drawColor,line width= 0.8pt,line join=round] ( 56.17, 57.55) -- ( 56.17, 55.79);

\path[draw=drawColor,line width= 0.8pt,line join=round] ( 56.42, 57.55) -- ( 56.42, 55.74);

\path[draw=drawColor,line width= 0.8pt,line join=round] ( 56.67, 57.55) -- ( 56.67, 57.60);

\path[draw=drawColor,line width= 0.8pt,line join=round] ( 56.92, 57.55) -- ( 56.92, 56.93);

\path[draw=drawColor,line width= 0.8pt,line join=round] ( 57.66, 57.55) -- ( 57.66, 62.47);

\path[draw=drawColor,line width= 0.8pt,line join=round] ( 57.91, 57.55) -- ( 57.91, 53.68);

\path[draw=drawColor,line width= 0.8pt,line join=round] ( 58.16, 57.55) -- ( 58.16, 56.41);

\path[draw=drawColor,line width= 0.8pt,line join=round] ( 58.41, 57.55) -- ( 58.41, 59.29);

\path[draw=drawColor,line width= 0.8pt,line join=round] ( 59.40, 57.55) -- ( 59.40, 61.15);

\path[draw=drawColor,line width= 0.8pt,line join=round] ( 59.65, 57.55) -- ( 59.65, 54.81);

\path[draw=drawColor,line width= 0.8pt,line join=round] ( 59.90, 57.55) -- ( 59.90, 57.26);

\path[draw=drawColor,line width= 0.8pt,line join=round] ( 60.14, 57.55) -- ( 60.14, 58.24);

\path[draw=drawColor,line width= 0.8pt,line join=round] ( 60.39, 57.55) -- ( 60.39, 63.03);

\path[draw=drawColor,line width= 0.8pt,line join=round] ( 61.14, 57.55) -- ( 61.14, 56.24);

\path[draw=drawColor,line width= 0.8pt,line join=round] ( 61.38, 57.55) -- ( 61.38, 59.97);

\path[draw=drawColor,line width= 0.8pt,line join=round] ( 61.63, 57.55) -- ( 61.63, 63.53);

\path[draw=drawColor,line width= 0.8pt,line join=round] ( 61.88, 57.55) -- ( 61.88, 55.29);

\path[draw=drawColor,line width= 0.8pt,line join=round] ( 62.13, 57.55) -- ( 62.13, 54.60);

\path[draw=drawColor,line width= 0.8pt,line join=round] ( 62.87, 57.55) -- ( 62.87, 61.21);

\path[draw=drawColor,line width= 0.8pt,line join=round] ( 63.12, 57.55) -- ( 63.12, 55.67);

\path[draw=drawColor,line width= 0.8pt,line join=round] ( 63.37, 57.55) -- ( 63.37, 60.17);

\path[draw=drawColor,line width= 0.8pt,line join=round] ( 63.62, 57.55) -- ( 63.62, 57.33);

\path[draw=drawColor,line width= 0.8pt,line join=round] ( 63.87, 57.55) -- ( 63.87, 57.02);

\path[draw=drawColor,line width= 0.8pt,line join=round] ( 64.61, 57.55) -- ( 64.61, 59.55);

\path[draw=drawColor,line width= 0.8pt,line join=round] ( 64.86, 57.55) -- ( 64.86, 59.91);

\path[draw=drawColor,line width= 0.8pt,line join=round] ( 65.11, 57.55) -- ( 65.11, 61.03);

\path[draw=drawColor,line width= 0.8pt,line join=round] ( 65.36, 57.55) -- ( 65.36, 55.89);

\path[draw=drawColor,line width= 0.8pt,line join=round] ( 65.61, 57.55) -- ( 65.61, 58.86);

\path[draw=drawColor,line width= 0.8pt,line join=round] ( 66.35, 57.55) -- ( 66.35, 57.36);

\path[draw=drawColor,line width= 0.8pt,line join=round] ( 66.60, 57.55) -- ( 66.60, 56.74);

\path[draw=drawColor,line width= 0.8pt,line join=round] ( 66.85, 57.55) -- ( 66.85, 57.52);

\path[draw=drawColor,line width= 0.8pt,line join=round] ( 67.10, 57.55) -- ( 67.10, 56.04);

\path[draw=drawColor,line width= 0.8pt,line join=round] ( 67.34, 57.55) -- ( 67.34, 59.91);

\path[draw=drawColor,line width= 0.8pt,line join=round] ( 68.09, 57.55) -- ( 68.09, 61.41);

\path[draw=drawColor,line width= 0.8pt,line join=round] ( 68.34, 57.55) -- ( 68.34, 55.94);

\path[draw=drawColor,line width= 0.8pt,line join=round] ( 68.59, 57.55) -- ( 68.59, 58.28);

\path[draw=drawColor,line width= 0.8pt,line join=round] ( 68.83, 57.55) -- ( 68.83, 58.33);

\path[draw=drawColor,line width= 0.8pt,line join=round] ( 69.08, 57.55) -- ( 69.08, 56.73);

\path[draw=drawColor,line width= 0.8pt,line join=round] ( 69.83, 57.55) -- ( 69.83, 57.18);

\path[draw=drawColor,line width= 0.8pt,line join=round] ( 70.07, 57.55) -- ( 70.07, 54.76);

\path[draw=drawColor,line width= 0.8pt,line join=round] ( 70.32, 57.55) -- ( 70.32, 59.38);

\path[draw=drawColor,line width= 0.8pt,line join=round] ( 70.57, 57.55) -- ( 70.57, 61.23);

\path[draw=drawColor,line width= 0.8pt,line join=round] ( 70.82, 57.55) -- ( 70.82, 55.71);

\path[draw=drawColor,line width= 0.8pt,line join=round] ( 71.81, 57.55) -- ( 71.81, 55.46);

\path[draw=drawColor,line width= 0.8pt,line join=round] ( 72.06, 57.55) -- ( 72.06, 57.76);

\path[draw=drawColor,line width= 0.8pt,line join=round] ( 72.31, 57.55) -- ( 72.31, 52.43);

\path[draw=drawColor,line width= 0.8pt,line join=round] ( 72.56, 57.55) -- ( 72.56, 55.67);

\path[draw=drawColor,line width= 0.8pt,line join=round] ( 73.30, 57.55) -- ( 73.30, 58.12);

\path[draw=drawColor,line width= 0.8pt,line join=round] ( 73.55, 57.55) -- ( 73.55, 59.28);

\path[draw=drawColor,line width= 0.8pt,line join=round] ( 73.80, 57.55) -- ( 73.80, 57.71);

\path[draw=drawColor,line width= 0.8pt,line join=round] ( 74.05, 57.55) -- ( 74.05, 55.48);

\path[draw=drawColor,line width= 0.8pt,line join=round] ( 74.30, 57.55) -- ( 74.30, 59.12);

\path[draw=drawColor,line width= 0.8pt,line join=round] ( 75.04, 57.55) -- ( 75.04, 57.82);

\path[draw=drawColor,line width= 0.8pt,line join=round] ( 75.29, 57.55) -- ( 75.29, 59.11);

\path[draw=drawColor,line width= 0.8pt,line join=round] ( 75.54, 57.55) -- ( 75.54, 62.20);

\path[draw=drawColor,line width= 0.8pt,line join=round] ( 75.79, 57.55) -- ( 75.79, 56.73);

\path[draw=drawColor,line width= 0.8pt,line join=round] ( 76.03, 57.55) -- ( 76.03, 56.70);

\path[draw=drawColor,line width= 0.8pt,line join=round] ( 76.78, 57.55) -- ( 76.78, 55.63);

\path[draw=drawColor,line width= 0.8pt,line join=round] ( 77.03, 57.55) -- ( 77.03, 56.10);

\path[draw=drawColor,line width= 0.8pt,line join=round] ( 77.27, 57.55) -- ( 77.27, 59.31);

\path[draw=drawColor,line width= 0.8pt,line join=round] ( 77.52, 57.55) -- ( 77.52, 58.05);

\path[draw=drawColor,line width= 0.8pt,line join=round] ( 77.77, 57.55) -- ( 77.77, 57.76);

\path[draw=drawColor,line width= 0.8pt,line join=round] ( 78.52, 57.55) -- ( 78.52, 59.53);

\path[draw=drawColor,line width= 0.8pt,line join=round] ( 78.76, 57.55) -- ( 78.76, 58.02);

\path[draw=drawColor,line width= 0.8pt,line join=round] ( 79.01, 57.55) -- ( 79.01, 56.81);

\path[draw=drawColor,line width= 0.8pt,line join=round] ( 79.26, 57.55) -- ( 79.26, 53.68);

\path[draw=drawColor,line width= 0.8pt,line join=round] ( 79.51, 57.55) -- ( 79.51, 57.81);

\path[draw=drawColor,line width= 0.8pt,line join=round] ( 80.25, 57.55) -- ( 80.25, 53.43);

\path[draw=drawColor,line width= 0.8pt,line join=round] ( 80.50, 57.55) -- ( 80.50, 56.23);

\path[draw=drawColor,line width= 0.8pt,line join=round] ( 80.75, 57.55) -- ( 80.75, 55.14);

\path[draw=drawColor,line width= 0.8pt,line join=round] ( 81.00, 57.55) -- ( 81.00, 58.74);

\path[draw=drawColor,line width= 0.8pt,line join=round] ( 81.99, 57.55) -- ( 81.99, 54.12);

\path[draw=drawColor,line width= 0.8pt,line join=round] ( 82.24, 57.55) -- ( 82.24, 59.89);

\path[draw=drawColor,line width= 0.8pt,line join=round] ( 82.49, 57.55) -- ( 82.49, 52.72);

\path[draw=drawColor,line width= 0.8pt,line join=round] ( 82.74, 57.55) -- ( 82.74, 55.90);

\path[draw=drawColor,line width= 0.8pt,line join=round] ( 82.99, 57.55) -- ( 82.99, 58.85);

\path[draw=drawColor,line width= 0.8pt,line join=round] ( 83.73, 57.55) -- ( 83.73, 61.73);

\path[draw=drawColor,line width= 0.8pt,line join=round] ( 83.98, 57.55) -- ( 83.98, 59.82);

\path[draw=drawColor,line width= 0.8pt,line join=round] ( 84.23, 57.55) -- ( 84.23, 56.72);

\path[draw=drawColor,line width= 0.8pt,line join=round] ( 84.48, 57.55) -- ( 84.48, 56.44);

\path[draw=drawColor,line width= 0.8pt,line join=round] ( 84.72, 57.55) -- ( 84.72, 55.87);

\path[draw=drawColor,line width= 0.8pt,line join=round] ( 85.47, 57.55) -- ( 85.47, 54.98);

\path[draw=drawColor,line width= 0.8pt,line join=round] ( 85.72, 57.55) -- ( 85.72, 57.46);

\path[draw=drawColor,line width= 0.8pt,line join=round] ( 85.96, 57.55) -- ( 85.96, 50.47);

\path[draw=drawColor,line width= 0.8pt,line join=round] ( 86.21, 57.55) -- ( 86.21, 48.98);

\path[draw=drawColor,line width= 0.8pt,line join=round] ( 86.46, 57.55) -- ( 86.46, 55.17);

\path[draw=drawColor,line width= 0.8pt,line join=round] ( 87.21, 57.55) -- ( 87.21, 55.15);

\path[draw=drawColor,line width= 0.8pt,line join=round] ( 87.45, 57.55) -- ( 87.45, 65.03);

\path[draw=drawColor,line width= 0.8pt,line join=round] ( 87.70, 57.55) -- ( 87.70, 56.19);

\path[draw=drawColor,line width= 0.8pt,line join=round] ( 87.95, 57.55) -- ( 87.95, 60.60);

\path[draw=drawColor,line width= 0.8pt,line join=round] ( 88.20, 57.55) -- ( 88.20, 58.38);

\path[draw=drawColor,line width= 0.8pt,line join=round] ( 88.94, 57.55) -- ( 88.94, 53.45);

\path[draw=drawColor,line width= 0.8pt,line join=round] ( 89.19, 57.55) -- ( 89.19, 55.55);

\path[draw=drawColor,line width= 0.8pt,line join=round] ( 89.44, 57.55) -- ( 89.44, 58.40);

\path[draw=drawColor,line width= 0.8pt,line join=round] ( 89.69, 57.55) -- ( 89.69, 59.30);

\path[draw=drawColor,line width= 0.8pt,line join=round] ( 89.94, 57.55) -- ( 89.94, 57.03);

\path[draw=drawColor,line width= 0.8pt,line join=round] ( 90.68, 57.55) -- ( 90.68, 66.04);

\path[draw=drawColor,line width= 0.8pt,line join=round] ( 90.93, 57.55) -- ( 90.93, 51.30);

\path[draw=drawColor,line width= 0.8pt,line join=round] ( 91.18, 57.55) -- ( 91.18, 60.16);

\path[draw=drawColor,line width= 0.8pt,line join=round] ( 91.43, 57.55) -- ( 91.43, 54.10);

\path[draw=drawColor,line width= 0.8pt,line join=round] ( 91.68, 57.55) -- ( 91.68, 59.72);

\path[draw=drawColor,line width= 0.8pt,line join=round] ( 92.42, 57.55) -- ( 92.42, 57.69);

\path[draw=drawColor,line width= 0.8pt,line join=round] ( 92.67, 57.55) -- ( 92.67, 58.37);

\path[draw=drawColor,line width= 0.8pt,line join=round] ( 92.92, 57.55) -- ( 92.92, 53.34);

\path[draw=drawColor,line width= 0.8pt,line join=round] ( 93.17, 57.55) -- ( 93.17, 54.52);

\path[draw=drawColor,line width= 0.8pt,line join=round] ( 93.41, 57.55) -- ( 93.41, 53.21);

\path[draw=drawColor,line width= 0.8pt,line join=round] ( 94.16, 57.55) -- ( 94.16, 50.67);

\path[draw=drawColor,line width= 0.8pt,line join=round] ( 94.41, 57.55) -- ( 94.41, 55.10);

\path[draw=drawColor,line width= 0.8pt,line join=round] ( 94.65, 57.55) -- ( 94.65, 63.33);

\path[draw=drawColor,line width= 0.8pt,line join=round] ( 94.90, 57.55) -- ( 94.90, 63.12);

\path[draw=drawColor,line width= 0.8pt,line join=round] ( 95.15, 57.55) -- ( 95.15, 58.44);

\path[draw=drawColor,line width= 0.8pt,line join=round] ( 95.90, 57.55) -- ( 95.90, 59.04);

\path[draw=drawColor,line width= 0.8pt,line join=round] ( 96.14, 57.55) -- ( 96.14, 52.80);

\path[draw=drawColor,line width= 0.8pt,line join=round] ( 96.39, 57.55) -- ( 96.39, 61.25);

\path[draw=drawColor,line width= 0.8pt,line join=round] ( 96.64, 57.55) -- ( 96.64, 52.44);

\path[draw=drawColor,line width= 0.8pt,line join=round] ( 96.89, 57.55) -- ( 96.89, 53.28);

\path[draw=drawColor,line width= 0.8pt,line join=round] ( 97.88, 57.55) -- ( 97.88, 61.34);

\path[draw=drawColor,line width= 0.8pt,line join=round] ( 98.13, 57.55) -- ( 98.13, 53.31);

\path[draw=drawColor,line width= 0.8pt,line join=round] ( 98.38, 57.55) -- ( 98.38, 55.81);

\path[draw=drawColor,line width= 0.8pt,line join=round] ( 98.63, 57.55) -- ( 98.63, 58.22);

\path[draw=drawColor,line width= 0.8pt,line join=round] ( 99.37, 57.55) -- ( 99.37, 57.99);

\path[draw=drawColor,line width= 0.8pt,line join=round] ( 99.62, 57.55) -- ( 99.62, 63.27);

\path[draw=drawColor,line width= 0.8pt,line join=round] ( 99.87, 57.55) -- ( 99.87, 60.83);

\path[draw=drawColor,line width= 0.8pt,line join=round] (100.12, 57.55) -- (100.12, 52.63);

\path[draw=drawColor,line width= 0.8pt,line join=round] (100.37, 57.55) -- (100.37, 50.52);

\path[draw=drawColor,line width= 0.8pt,line join=round] (101.11, 57.55) -- (101.11, 58.52);

\path[draw=drawColor,line width= 0.8pt,line join=round] (101.36, 57.55) -- (101.36, 55.99);

\path[draw=drawColor,line width= 0.8pt,line join=round] (101.61, 57.55) -- (101.61, 52.68);

\path[draw=drawColor,line width= 0.8pt,line join=round] (101.85, 57.55) -- (101.85, 42.46);

\path[draw=drawColor,line width= 0.8pt,line join=round] (102.10, 57.55) -- (102.10, 54.62);

\path[draw=drawColor,line width= 0.8pt,line join=round] (102.85, 57.55) -- (102.85, 53.24);

\path[draw=drawColor,line width= 0.8pt,line join=round] (103.10, 57.55) -- (103.10, 59.44);

\path[draw=drawColor,line width= 0.8pt,line join=round] (103.34, 57.55) -- (103.34, 61.32);

\path[draw=drawColor,line width= 0.8pt,line join=round] (103.59, 57.55) -- (103.59, 54.06);

\path[draw=drawColor,line width= 0.8pt,line join=round] (103.84, 57.55) -- (103.84, 62.24);

\path[draw=drawColor,line width= 0.8pt,line join=round] (104.59, 57.55) -- (104.59, 69.57);

\path[draw=drawColor,line width= 0.8pt,line join=round] (104.83, 57.55) -- (104.83, 63.00);

\path[draw=drawColor,line width= 0.8pt,line join=round] (105.08, 57.55) -- (105.08, 56.09);

\path[draw=drawColor,line width= 0.8pt,line join=round] (105.33, 57.55) -- (105.33, 62.33);

\path[draw=drawColor,line width= 0.8pt,line join=round] (105.58, 57.55) -- (105.58, 55.79);

\path[draw=drawColor,line width= 0.8pt,line join=round] (106.32, 57.55) -- (106.32, 54.92);

\path[draw=drawColor,line width= 0.8pt,line join=round] (106.57, 57.55) -- (106.57, 56.45);

\path[draw=drawColor,line width= 0.8pt,line join=round] (106.82, 57.55) -- (106.82, 59.33);

\path[draw=drawColor,line width= 0.8pt,line join=round] (107.07, 57.55) -- (107.07, 57.94);

\path[draw=drawColor,line width= 0.8pt,line join=round] (107.32, 57.55) -- (107.32, 53.73);

\path[draw=drawColor,line width= 0.8pt,line join=round] (108.06, 57.55) -- (108.06, 56.17);

\path[draw=drawColor,line width= 0.8pt,line join=round] (108.31, 57.55) -- (108.31, 60.89);

\path[draw=drawColor,line width= 0.8pt,line join=round] (108.56, 57.55) -- (108.56, 55.55);

\path[draw=drawColor,line width= 0.8pt,line join=round] (108.81, 57.55) -- (108.81, 61.56);

\path[draw=drawColor,line width= 0.8pt,line join=round] (109.06, 57.55) -- (109.06, 60.39);

\path[draw=drawColor,line width= 0.8pt,line join=round] (109.80, 57.55) -- (109.80, 56.80);

\path[draw=drawColor,line width= 0.8pt,line join=round] (110.05, 57.55) -- (110.05, 53.83);

\path[draw=drawColor,line width= 0.8pt,line join=round] (110.30, 57.55) -- (110.30, 62.72);

\path[draw=drawColor,line width= 0.8pt,line join=round] (110.54, 57.55) -- (110.54, 57.55);

\path[draw=drawColor,line width= 0.8pt,line join=round] (110.79, 57.55) -- (110.79, 60.82);

\path[draw=drawColor,line width= 0.8pt,line join=round] (111.54, 57.55) -- (111.54, 61.82);

\path[draw=drawColor,line width= 0.8pt,line join=round] (111.79, 57.55) -- (111.79, 58.83);

\path[draw=drawColor,line width= 0.8pt,line join=round] (112.03, 57.55) -- (112.03, 56.92);

\path[draw=drawColor,line width= 0.8pt,line join=round] (112.28, 57.55) -- (112.28, 56.51);

\path[draw=drawColor,line width= 0.8pt,line join=round] (112.53, 57.55) -- (112.53, 56.35);

\path[draw=drawColor,line width= 0.8pt,line join=round] (113.28, 57.55) -- (113.28, 51.31);

\path[draw=drawColor,line width= 0.8pt,line join=round] (113.52, 57.55) -- (113.52, 59.29);

\path[draw=drawColor,line width= 0.8pt,line join=round] (113.77, 57.55) -- (113.77, 55.87);

\path[draw=drawColor,line width= 0.8pt,line join=round] (114.02, 57.55) -- (114.02, 46.77);

\path[draw=drawColor,line width= 0.8pt,line join=round] (114.27, 57.55) -- (114.27, 60.81);

\path[draw=drawColor,line width= 0.8pt,line join=round] (115.01, 57.55) -- (115.01, 60.08);

\path[draw=drawColor,line width= 0.8pt,line join=round] (115.26, 57.55) -- (115.26, 54.21);

\path[draw=drawColor,line width= 0.8pt,line join=round] (115.51, 57.55) -- (115.51, 60.65);

\path[draw=drawColor,line width= 0.8pt,line join=round] (115.76, 57.55) -- (115.76, 56.55);

\path[draw=drawColor,line width= 0.8pt,line join=round] (116.01, 57.55) -- (116.01, 61.24);

\path[draw=drawColor,line width= 0.8pt,line join=round] (116.75, 57.55) -- (116.75, 57.15);

\path[draw=drawColor,line width= 0.8pt,line join=round] (117.00, 57.55) -- (117.00, 58.79);

\path[draw=drawColor,line width= 0.8pt,line join=round] (117.25, 57.55) -- (117.25, 57.84);

\path[draw=drawColor,line width= 0.8pt,line join=round] (117.75, 57.55) -- (117.75, 56.66);

\path[draw=drawColor,line width= 0.8pt,line join=round] (118.49, 57.55) -- (118.49, 62.18);

\path[draw=drawColor,line width= 0.8pt,line join=round] (118.74, 57.55) -- (118.74, 54.02);

\path[draw=drawColor,line width= 0.8pt,line join=round] (118.99, 57.55) -- (118.99, 55.81);

\path[draw=drawColor,line width= 0.8pt,line join=round] (119.23, 57.55) -- (119.23, 52.94);

\path[draw=drawColor,line width= 0.8pt,line join=round] (119.48, 57.55) -- (119.48, 58.68);

\path[draw=drawColor,line width= 0.8pt,line join=round] (120.23, 57.55) -- (120.23, 52.09);

\path[draw=drawColor,line width= 0.8pt,line join=round] (120.48, 57.55) -- (120.48, 51.26);

\path[draw=drawColor,line width= 0.8pt,line join=round] (120.72, 57.55) -- (120.72, 55.33);

\path[draw=drawColor,line width= 0.8pt,line join=round] (120.97, 57.55) -- (120.97, 59.20);

\path[draw=drawColor,line width= 0.8pt,line join=round] (121.22, 57.55) -- (121.22, 54.19);

\path[draw=drawColor,line width= 0.8pt,line join=round] (121.97, 57.55) -- (121.97, 59.99);

\path[draw=drawColor,line width= 0.8pt,line join=round] (122.21, 57.55) -- (122.21, 59.76);

\path[draw=drawColor,line width= 0.8pt,line join=round] (122.46, 57.55) -- (122.46, 60.91);

\path[draw=drawColor,line width= 0.8pt,line join=round] (122.71, 57.55) -- (122.71, 47.64);

\path[draw=drawColor,line width= 0.8pt,line join=round] (122.96, 57.55) -- (122.96, 58.30);

\path[draw=drawColor,line width= 0.8pt,line join=round] (123.70, 57.55) -- (123.70, 58.01);

\path[draw=drawColor,line width= 0.8pt,line join=round] (123.95, 57.55) -- (123.95, 68.60);

\path[draw=drawColor,line width= 0.8pt,line join=round] (124.20, 57.55) -- (124.20, 62.53);

\path[draw=drawColor,line width= 0.8pt,line join=round] (124.45, 57.55) -- (124.45, 55.84);

\path[draw=drawColor,line width= 0.8pt,line join=round] (125.44, 57.55) -- (125.44, 54.84);

\path[draw=drawColor,line width= 0.8pt,line join=round] (125.69, 57.55) -- (125.69, 59.61);

\path[draw=drawColor,line width= 0.8pt,line join=round] (125.94, 57.55) -- (125.94, 55.86);

\path[draw=drawColor,line width= 0.8pt,line join=round] (126.19, 57.55) -- (126.19, 55.06);

\path[draw=drawColor,line width= 0.8pt,line join=round] (127.18, 57.55) -- (127.18, 57.65);

\path[draw=drawColor,line width= 0.8pt,line join=round] (127.43, 57.55) -- (127.43, 55.10);

\path[draw=drawColor,line width= 0.8pt,line join=round] (127.68, 57.55) -- (127.68, 53.85);

\path[draw=drawColor,line width= 0.8pt,line join=round] (127.92, 57.55) -- (127.92, 49.38);

\path[draw=drawColor,line width= 0.8pt,line join=round] (128.17, 57.55) -- (128.17, 55.16);

\path[draw=drawColor,line width= 0.8pt,line join=round] (128.92, 57.55) -- (128.92, 50.75);

\path[draw=drawColor,line width= 0.8pt,line join=round] (129.17, 57.55) -- (129.17, 58.04);

\path[draw=drawColor,line width= 0.8pt,line join=round] (129.41, 57.55) -- (129.41, 54.84);

\path[draw=drawColor,line width= 0.8pt,line join=round] (129.66, 57.55) -- (129.66, 62.77);

\path[draw=drawColor,line width= 0.8pt,line join=round] (129.91, 57.55) -- (129.91, 48.39);

\path[draw=drawColor,line width= 0.8pt,line join=round] (130.90, 57.55) -- (130.90, 54.25);

\path[draw=drawColor,line width= 0.8pt,line join=round] (131.15, 57.55) -- (131.15, 56.68);

\path[draw=drawColor,line width= 0.8pt,line join=round] (131.40, 57.55) -- (131.40, 61.01);

\path[draw=drawColor,line width= 0.8pt,line join=round] (131.65, 57.55) -- (131.65, 62.53);

\path[draw=drawColor,line width= 0.8pt,line join=round] (132.39, 57.55) -- (132.39, 45.51);

\path[draw=drawColor,line width= 0.8pt,line join=round] (132.64, 57.55) -- (132.64, 62.52);

\path[draw=drawColor,line width= 0.8pt,line join=round] (132.89, 57.55) -- (132.89, 54.21);

\path[draw=drawColor,line width= 0.8pt,line join=round] (133.14, 57.55) -- (133.14, 68.32);

\path[draw=drawColor,line width= 0.8pt,line join=round] (133.39, 57.55) -- (133.39, 61.93);

\path[draw=drawColor,line width= 0.8pt,line join=round] (134.13, 57.55) -- (134.13, 59.26);

\path[draw=drawColor,line width= 0.8pt,line join=round] (134.38, 57.55) -- (134.38, 49.88);

\path[draw=drawColor,line width= 0.8pt,line join=round] (134.63, 57.55) -- (134.63, 67.32);

\path[draw=drawColor,line width= 0.8pt,line join=round] (134.88, 57.55) -- (134.88, 67.25);

\path[draw=drawColor,line width= 0.8pt,line join=round] (135.12, 57.55) -- (135.12, 58.11);

\path[draw=drawColor,line width= 0.8pt,line join=round] (135.87, 57.55) -- (135.87, 52.05);

\path[draw=drawColor,line width= 0.8pt,line join=round] (136.12, 57.55) -- (136.12, 55.19);

\path[draw=drawColor,line width= 0.8pt,line join=round] (136.37, 57.55) -- (136.37, 50.93);

\path[draw=drawColor,line width= 0.8pt,line join=round] (136.61, 57.55) -- (136.61, 54.79);

\path[draw=drawColor,line width= 0.8pt,line join=round] (136.86, 57.55) -- (136.86, 64.06);

\path[draw=drawColor,line width= 0.8pt,line join=round] (137.86, 57.55) -- (137.86, 65.03);

\path[draw=drawColor,line width= 0.8pt,line join=round] (138.10, 57.55) -- (138.10, 64.76);

\path[draw=drawColor,line width= 0.8pt,line join=round] (138.35, 57.55) -- (138.35, 53.56);

\path[draw=drawColor,line width= 0.8pt,line join=round] (138.60, 57.55) -- (138.60, 55.07);

\path[draw=drawColor,line width= 0.8pt,line join=round] (139.35, 57.55) -- (139.35, 64.18);

\path[draw=drawColor,line width= 0.8pt,line join=round] (139.59, 57.55) -- (139.59, 52.19);

\path[draw=drawColor,line width= 0.8pt,line join=round] (139.84, 57.55) -- (139.84, 57.90);

\path[draw=drawColor,line width= 0.8pt,line join=round] (140.09, 57.55) -- (140.09, 59.62);

\path[draw=drawColor,line width= 0.8pt,line join=round] (140.34, 57.55) -- (140.34, 58.94);

\path[draw=drawColor,line width= 0.8pt,line join=round] (141.08, 57.55) -- (141.08, 60.42);

\path[draw=drawColor,line width= 0.8pt,line join=round] (141.33, 57.55) -- (141.33, 62.28);

\path[draw=drawColor,line width= 0.8pt,line join=round] (141.58, 57.55) -- (141.58, 58.52);

\path[draw=drawColor,line width= 0.8pt,line join=round] (141.83, 57.55) -- (141.83, 65.37);

\path[draw=drawColor,line width= 0.8pt,line join=round] (142.08, 57.55) -- (142.08, 61.06);

\path[draw=drawColor,line width= 0.8pt,line join=round] (142.82, 57.55) -- (142.82, 63.64);

\path[draw=drawColor,line width= 0.8pt,line join=round] (143.07, 57.55) -- (143.07, 47.87);

\path[draw=drawColor,line width= 0.8pt,line join=round] (143.32, 57.55) -- (143.32, 58.00);

\path[draw=drawColor,line width= 0.8pt,line join=round] (143.57, 57.55) -- (143.57, 55.89);

\path[draw=drawColor,line width= 0.8pt,line join=round] (143.81, 57.55) -- (143.81, 62.20);

\path[draw=drawColor,line width= 0.8pt,line join=round] (144.56, 57.55) -- (144.56, 57.32);

\path[draw=drawColor,line width= 0.8pt,line join=round] (144.81, 57.55) -- (144.81, 56.61);

\path[draw=drawColor,line width= 0.8pt,line join=round] (145.06, 57.55) -- (145.06, 63.58);

\path[draw=drawColor,line width= 0.8pt,line join=round] (145.30, 57.55) -- (145.30, 62.38);

\path[draw=drawColor,line width= 0.8pt,line join=round] (145.55, 57.55) -- (145.55, 56.22);

\path[draw=drawColor,line width= 0.8pt,line join=round] (146.30, 57.55) -- (146.30, 58.84);

\path[draw=drawColor,line width= 0.8pt,line join=round] (146.55, 57.55) -- (146.55, 56.50);

\path[draw=drawColor,line width= 0.8pt,line join=round] (146.79, 57.55) -- (146.79, 51.71);

\path[draw=drawColor,line width= 0.8pt,line join=round] (147.04, 57.55) -- (147.04, 62.58);

\path[draw=drawColor,line width= 0.8pt,line join=round] (148.04, 57.55) -- (148.04, 57.64);

\path[draw=drawColor,line width= 0.8pt,line join=round] (148.28, 57.55) -- (148.28, 59.67);

\path[draw=drawColor,line width= 0.8pt,line join=round] (148.53, 57.55) -- (148.53, 59.07);

\path[draw=drawColor,line width= 0.8pt,line join=round] (148.78, 57.55) -- (148.78, 57.64);

\path[draw=drawColor,line width= 0.8pt,line join=round] (149.03, 57.55) -- (149.03, 58.31);

\path[draw=drawColor,line width= 0.8pt,line join=round] (149.77, 57.55) -- (149.77, 54.28);

\path[draw=drawColor,line width= 0.8pt,line join=round] (150.02, 57.55) -- (150.02, 56.06);

\path[draw=drawColor,line width= 0.8pt,line join=round] (150.27, 57.55) -- (150.27, 57.48);

\path[draw=drawColor,line width= 0.8pt,line join=round] (150.52, 57.55) -- (150.52, 54.27);

\path[draw=drawColor,line width= 0.8pt,line join=round] (150.77, 57.55) -- (150.77, 58.11);

\path[draw=drawColor,line width= 0.8pt,line join=round] (151.51, 57.55) -- (151.51, 58.42);

\path[draw=drawColor,line width= 0.8pt,line join=round] (151.76, 57.55) -- (151.76, 62.09);

\path[draw=drawColor,line width= 0.8pt,line join=round] (152.01, 57.55) -- (152.01, 66.64);

\path[draw=drawColor,line width= 0.8pt,line join=round] (152.26, 57.55) -- (152.26, 57.31);

\path[draw=drawColor,line width= 0.8pt,line join=round] (152.50, 57.55) -- (152.50, 57.90);

\path[draw=drawColor,line width= 0.8pt,line join=round] (153.25, 57.55) -- (153.25, 57.78);

\path[draw=drawColor,line width= 0.8pt,line join=round] (153.50, 57.55) -- (153.50, 60.87);

\path[draw=drawColor,line width= 0.8pt,line join=round] (153.75, 57.55) -- (153.75, 55.65);

\path[draw=drawColor,line width= 0.8pt,line join=round] (153.99, 57.55) -- (153.99, 54.37);

\path[draw=drawColor,line width= 0.8pt,line join=round] (154.24, 57.55) -- (154.24, 56.54);

\path[draw=drawColor,line width= 0.8pt,line join=round] (154.99, 57.55) -- (154.99, 52.95);

\path[draw=drawColor,line width= 0.8pt,line join=round] (155.24, 57.55) -- (155.24, 60.14);

\path[draw=drawColor,line width= 0.8pt,line join=round] (155.48, 57.55) -- (155.48, 60.62);

\path[draw=drawColor,line width= 0.8pt,line join=round] (155.73, 57.55) -- (155.73, 54.78);

\path[draw=drawColor,line width= 0.8pt,line join=round] (155.98, 57.55) -- (155.98, 57.46);

\path[draw=drawColor,line width= 0.8pt,line join=round] (156.73, 57.55) -- (156.73, 57.91);

\path[draw=drawColor,line width= 0.8pt,line join=round] (156.97, 57.55) -- (156.97, 53.07);

\path[draw=drawColor,line width= 0.8pt,line join=round] (157.22, 57.55) -- (157.22, 50.99);

\path[draw=drawColor,line width= 0.8pt,line join=round] (157.47, 57.55) -- (157.47, 52.95);

\path[draw=drawColor,line width= 0.8pt,line join=round] (157.72, 57.55) -- (157.72, 59.36);

\path[draw=drawColor,line width= 0.8pt,line join=round] (158.46, 57.55) -- (158.46, 49.20);

\path[draw=drawColor,line width= 0.8pt,line join=round] (158.71, 57.55) -- (158.71, 63.17);

\path[draw=drawColor,line width= 0.8pt,line join=round] (158.96, 57.55) -- (158.96, 57.29);

\path[draw=drawColor,line width= 0.8pt,line join=round] (159.21, 57.55) -- (159.21, 55.19);

\path[draw=drawColor,line width= 0.8pt,line join=round] (159.46, 57.55) -- (159.46, 52.19);

\path[draw=drawColor,line width= 0.8pt,line join=round] (160.20, 57.55) -- (160.20, 59.63);

\path[draw=drawColor,line width= 0.8pt,line join=round] (160.45, 57.55) -- (160.45, 58.70);

\path[draw=drawColor,line width= 0.8pt,line join=round] (160.70, 57.55) -- (160.70, 56.03);

\path[draw=drawColor,line width= 0.8pt,line join=round] (160.95, 57.55) -- (160.95, 53.68);

\path[draw=drawColor,line width= 0.8pt,line join=round] (161.19, 57.55) -- (161.19, 59.02);

\path[draw=drawColor,line width= 0.8pt,line join=round] (161.94, 57.55) -- (161.94, 59.31);

\path[draw=drawColor,line width= 0.8pt,line join=round] (162.19, 57.55) -- (162.19, 59.82);

\path[draw=drawColor,line width= 0.8pt,line join=round] (162.44, 57.55) -- (162.44, 62.35);

\path[draw=drawColor,line width= 0.8pt,line join=round] (162.68, 57.55) -- (162.68, 55.97);

\path[draw=drawColor,line width= 0.8pt,line join=round] (162.93, 57.55) -- (162.93, 57.16);

\path[draw=drawColor,line width= 0.8pt,line join=round] (163.93, 57.55) -- (163.93, 59.32);

\path[draw=drawColor,line width= 0.8pt,line join=round] (164.17, 57.55) -- (164.17, 56.77);

\path[draw=drawColor,line width= 0.8pt,line join=round] (164.42, 57.55) -- (164.42, 61.86);

\path[draw=drawColor,line width= 0.8pt,line join=round] (164.67, 57.55) -- (164.67, 62.00);

\path[draw=drawColor,line width= 0.8pt,line join=round] (165.42, 57.55) -- (165.42, 61.79);

\path[draw=drawColor,line width= 0.8pt,line join=round] (165.66, 57.55) -- (165.66, 58.73);

\path[draw=drawColor,line width= 0.8pt,line join=round] (165.91, 57.55) -- (165.91, 61.46);

\path[draw=drawColor,line width= 0.8pt,line join=round] (166.16, 57.55) -- (166.16, 54.69);

\path[draw=drawColor,line width= 0.8pt,line join=round] (166.41, 57.55) -- (166.41, 54.11);

\path[draw=drawColor,line width= 0.8pt,line join=round] (167.15, 57.55) -- (167.15, 55.08);

\path[draw=drawColor,line width= 0.8pt,line join=round] (167.40, 57.55) -- (167.40, 56.40);

\path[draw=drawColor,line width= 0.8pt,line join=round] (167.65, 57.55) -- (167.65, 58.20);

\path[draw=drawColor,line width= 0.8pt,line join=round] (167.90, 57.55) -- (167.90, 57.30);

\path[draw=drawColor,line width= 0.8pt,line join=round] (168.15, 57.55) -- (168.15, 60.45);

\path[draw=drawColor,line width= 0.8pt,line join=round] (168.89, 57.55) -- (168.89, 59.07);

\path[draw=drawColor,line width= 0.8pt,line join=round] (169.14, 57.55) -- (169.14, 57.73);

\path[draw=drawColor,line width= 0.8pt,line join=round] (169.39, 57.55) -- (169.39, 57.36);

\path[draw=drawColor,line width= 0.8pt,line join=round] (169.64, 57.55) -- (169.64, 62.94);

\path[draw=drawColor,line width= 0.8pt,line join=round] (169.88, 57.55) -- (169.88, 41.55);

\path[draw=drawColor,line width= 0.8pt,line join=round] (170.63, 57.55) -- (170.63, 52.22);

\path[draw=drawColor,line width= 0.8pt,line join=round] (170.88, 57.55) -- (170.88, 61.24);

\path[draw=drawColor,line width= 0.8pt,line join=round] (171.13, 57.55) -- (171.13, 62.98);

\path[draw=drawColor,line width= 0.8pt,line join=round] (171.37, 57.55) -- (171.37, 62.45);

\path[draw=drawColor,line width= 0.8pt,line join=round] (171.62, 57.55) -- (171.62, 59.50);

\path[draw=drawColor,line width= 0.8pt,line join=round] (172.62, 57.55) -- (172.62, 51.15);

\path[draw=drawColor,line width= 0.8pt,line join=round] (172.86, 57.55) -- (172.86, 58.95);

\path[draw=drawColor,line width= 0.8pt,line join=round] (173.11, 57.55) -- (173.11, 58.23);

\path[draw=drawColor,line width= 0.8pt,line join=round] (173.36, 57.55) -- (173.36, 64.64);

\path[draw=drawColor,line width= 0.8pt,line join=round] (174.11, 57.55) -- (174.11, 58.84);

\path[draw=drawColor,line width= 0.8pt,line join=round] (174.35, 57.55) -- (174.35, 63.46);

\path[draw=drawColor,line width= 0.8pt,line join=round] (174.60, 57.55) -- (174.60, 57.22);

\path[draw=drawColor,line width= 0.8pt,line join=round] (174.85, 57.55) -- (174.85, 58.63);

\path[draw=drawColor,line width= 0.8pt,line join=round] (175.10, 57.55) -- (175.10, 59.40);

\path[draw=drawColor,line width= 0.8pt,line join=round] (175.84, 57.55) -- (175.84, 54.99);

\path[draw=drawColor,line width= 0.8pt,line join=round] (176.09, 57.55) -- (176.09, 57.31);

\path[draw=drawColor,line width= 0.8pt,line join=round] (176.34, 57.55) -- (176.34, 57.40);

\path[draw=drawColor,line width= 0.8pt,line join=round] (176.59, 57.55) -- (176.59, 58.48);

\path[draw=drawColor,line width= 0.8pt,line join=round] (176.84, 57.55) -- (176.84, 55.70);

\path[draw=drawColor,line width= 0.8pt,line join=round] (177.58, 57.55) -- (177.58, 55.51);

\path[draw=drawColor,line width= 0.8pt,line join=round] (177.83, 57.55) -- (177.83, 69.27);

\path[draw=drawColor,line width= 0.8pt,line join=round] (178.08, 57.55) -- (178.08, 61.32);

\path[draw=drawColor,line width= 0.8pt,line join=round] (178.33, 57.55) -- (178.33, 54.25);

\path[draw=drawColor,line width= 0.8pt,line join=round] (178.57, 57.55) -- (178.57, 57.10);

\path[draw=drawColor,line width= 0.8pt,line join=round] (179.32, 57.55) -- (179.32, 55.25);

\path[draw=drawColor,line width= 0.8pt,line join=round] (179.57, 57.55) -- (179.57, 55.52);

\path[draw=drawColor,line width= 0.8pt,line join=round] (179.82, 57.55) -- (179.82, 59.80);

\path[draw=drawColor,line width= 0.8pt,line join=round] (180.06, 57.55) -- (180.06, 56.81);

\path[draw=drawColor,line width= 0.8pt,line join=round] (180.31, 57.55) -- (180.31, 59.79);

\path[draw=drawColor,line width= 0.8pt,line join=round] (181.06, 57.55) -- (181.06, 60.02);

\path[draw=drawColor,line width= 0.8pt,line join=round] (181.31, 57.55) -- (181.31, 55.84);

\path[draw=drawColor,line width= 0.8pt,line join=round] (181.55, 57.55) -- (181.55, 57.04);

\path[draw=drawColor,line width= 0.8pt,line join=round] (181.80, 57.55) -- (181.80, 59.20);

\path[draw=drawColor,line width= 0.8pt,line join=round] (182.05, 57.55) -- (182.05, 56.87);

\path[draw=drawColor,line width= 0.8pt,line join=round] (182.80, 57.55) -- (182.80, 60.75);

\path[draw=drawColor,line width= 0.8pt,line join=round] (183.04, 57.55) -- (183.04, 57.93);

\path[draw=drawColor,line width= 0.8pt,line join=round] (183.29, 57.55) -- (183.29, 57.88);

\path[draw=drawColor,line width= 0.8pt,line join=round] (183.54, 57.55) -- (183.54, 54.68);

\path[draw=drawColor,line width= 0.8pt,line join=round] (183.79, 57.55) -- (183.79, 58.83);

\path[draw=drawColor,line width= 0.8pt,line join=round] (184.53, 57.55) -- (184.53, 56.85);

\path[draw=drawColor,line width= 0.8pt,line join=round] (184.78, 57.55) -- (184.78, 58.71);

\path[draw=drawColor,line width= 0.8pt,line join=round] (185.03, 57.55) -- (185.03, 55.15);

\path[draw=drawColor,line width= 0.8pt,line join=round] (185.28, 57.55) -- (185.28, 56.68);

\path[draw=drawColor,line width= 0.8pt,line join=round] (185.53, 57.55) -- (185.53, 57.04);

\path[draw=drawColor,line width= 0.8pt,line join=round] (186.27, 57.55) -- (186.27, 58.78);

\path[draw=drawColor,line width= 0.8pt,line join=round] (186.52, 57.55) -- (186.52, 55.80);

\path[draw=drawColor,line width= 0.8pt,line join=round] (186.77, 57.55) -- (186.77, 56.04);

\path[draw=drawColor,line width= 0.8pt,line join=round] (187.02, 57.55) -- (187.02, 56.12);

\path[draw=drawColor,line width= 0.8pt,line join=round] (187.26, 57.55) -- (187.26, 58.23);

\path[draw=drawColor,line width= 0.8pt,line join=round] (188.26, 57.55) -- (188.26, 58.66);

\path[draw=drawColor,line width= 0.8pt,line join=round] (188.51, 57.55) -- (188.51, 62.16);

\path[draw=drawColor,line width= 0.8pt,line join=round] (188.75, 57.55) -- (188.75, 57.04);

\path[draw=drawColor,line width= 0.8pt,line join=round] (189.00, 57.55) -- (189.00, 49.75);

\path[draw=drawColor,line width= 0.8pt,line join=round] (189.75, 57.55) -- (189.75, 60.78);

\path[draw=drawColor,line width= 0.8pt,line join=round] (190.00, 57.55) -- (190.00, 52.69);

\path[draw=drawColor,line width= 0.8pt,line join=round] (190.24, 57.55) -- (190.24, 59.83);

\path[draw=drawColor,line width= 0.8pt,line join=round] (190.49, 57.55) -- (190.49, 60.43);

\path[draw=drawColor,line width= 0.8pt,line join=round] (190.74, 57.55) -- (190.74, 57.60);

\path[draw=drawColor,line width= 0.8pt,line join=round] (191.49, 57.55) -- (191.49, 57.06);

\path[draw=drawColor,line width= 0.8pt,line join=round] (191.73, 57.55) -- (191.73, 57.03);

\path[draw=drawColor,line width= 0.8pt,line join=round] (191.98, 57.55) -- (191.98, 62.54);

\path[draw=drawColor,line width= 0.8pt,line join=round] (192.23, 57.55) -- (192.23, 57.63);

\path[draw=drawColor,line width= 0.8pt,line join=round] (192.48, 57.55) -- (192.48, 54.57);

\path[draw=drawColor,line width= 0.8pt,line join=round] (193.22, 57.55) -- (193.22, 57.35);

\path[draw=drawColor,line width= 0.8pt,line join=round] (193.47, 57.55) -- (193.47, 58.98);

\path[draw=drawColor,line width= 0.8pt,line join=round] (193.72, 57.55) -- (193.72, 67.75);

\path[draw=drawColor,line width= 0.8pt,line join=round] (193.97, 57.55) -- (193.97, 59.90);

\path[draw=drawColor,line width= 0.8pt,line join=round] (194.22, 57.55) -- (194.22, 60.98);

\path[draw=drawColor,line width= 0.8pt,line join=round] (194.96, 57.55) -- (194.96, 56.26);

\path[draw=drawColor,line width= 0.8pt,line join=round] (195.21, 57.55) -- (195.21, 55.50);

\path[draw=drawColor,line width= 0.8pt,line join=round] (195.46, 57.55) -- (195.46, 62.58);

\path[draw=drawColor,line width= 0.8pt,line join=round] (195.71, 57.55) -- (195.71, 57.21);

\path[draw=drawColor,line width= 0.8pt,line join=round] (195.95, 57.55) -- (195.95, 55.40);

\path[draw=drawColor,line width= 0.8pt,line join=round] (196.70, 57.55) -- (196.70, 56.89);

\path[draw=drawColor,line width= 0.8pt,line join=round] (196.95, 57.55) -- (196.95, 55.56);

\path[draw=drawColor,line width= 0.8pt,line join=round] (197.20, 57.55) -- (197.20, 57.84);

\path[draw=drawColor,line width= 0.8pt,line join=round] (197.44, 57.55) -- (197.44, 55.92);

\path[draw=drawColor,line width= 0.8pt,line join=round] (197.69, 57.55) -- (197.69, 59.41);

\path[draw=drawColor,line width= 0.8pt,line join=round] (198.44, 57.55) -- (198.44, 56.53);

\path[draw=drawColor,line width= 0.8pt,line join=round] (198.69, 57.55) -- (198.69, 57.36);

\path[draw=drawColor,line width= 0.8pt,line join=round] (198.93, 57.55) -- (198.93, 59.62);

\path[draw=drawColor,line width= 0.8pt,line join=round] (199.18, 57.55) -- (199.18, 53.89);

\path[draw=drawColor,line width= 0.8pt,line join=round] (199.43, 57.55) -- (199.43, 56.74);

\path[draw=drawColor,line width= 0.8pt,line join=round] (200.18, 57.55) -- (200.18, 56.63);

\path[draw=drawColor,line width= 0.8pt,line join=round] (200.42, 57.55) -- (200.42, 53.42);

\path[draw=drawColor,line width= 0.8pt,line join=round] (200.67, 57.55) -- (200.67, 56.58);

\path[draw=drawColor,line width= 0.8pt,line join=round] (200.92, 57.55) -- (200.92, 54.42);

\path[draw=drawColor,line width= 0.8pt,line join=round] (201.17, 57.55) -- (201.17, 59.97);

\path[draw=drawColor,line width= 0.8pt,line join=round] (201.91, 57.55) -- (201.91, 56.38);

\path[draw=drawColor,line width= 0.8pt,line join=round] (202.16, 57.55) -- (202.16, 54.11);

\path[draw=drawColor,line width= 0.8pt,line join=round] (202.41, 57.55) -- (202.41, 54.32);

\path[draw=drawColor,line width= 0.8pt,line join=round] (202.66, 57.55) -- (202.66, 58.00);

\path[draw=drawColor,line width= 0.8pt,line join=round] (202.91, 57.55) -- (202.91, 60.51);

\path[draw=drawColor,line width= 0.8pt,line join=round] (203.65, 57.55) -- (203.65, 61.59);

\path[draw=drawColor,line width= 0.8pt,line join=round] (203.90, 57.55) -- (203.90, 60.12);

\path[draw=drawColor,line width= 0.8pt,line join=round] (204.15, 57.55) -- (204.15, 74.83);

\path[draw=drawColor,line width= 0.8pt,line join=round] (204.40, 57.55) -- (204.40, 63.22);

\path[draw=drawColor,line width= 0.8pt,line join=round] (204.64, 57.55) -- (204.64, 56.45);

\path[draw=drawColor,line width= 0.8pt,line join=round] (205.39, 57.55) -- (205.39, 60.43);

\path[draw=drawColor,line width= 0.8pt,line join=round] (205.64, 57.55) -- (205.64, 58.21);

\path[draw=drawColor,line width= 0.8pt,line join=round] (205.89, 57.55) -- (205.89, 54.84);

\path[draw=drawColor,line width= 0.8pt,line join=round] (206.13, 57.55) -- (206.13, 56.12);

\path[draw=drawColor,line width= 0.8pt,line join=round] (206.38, 57.55) -- (206.38, 56.44);

\path[draw=drawColor,line width= 0.8pt,line join=round] (207.13, 57.55) -- (207.13, 58.95);

\path[draw=drawColor,line width= 0.8pt,line join=round] (207.38, 57.55) -- (207.38, 59.34);

\path[draw=drawColor,line width= 0.8pt,line join=round] (207.62, 57.55) -- (207.62, 63.83);

\path[draw=drawColor,line width= 0.8pt,line join=round] (208.12, 57.55) -- (208.12, 56.65);

\path[draw=drawColor,line width= 0.8pt,line join=round] (208.86, 57.55) -- (208.86, 55.32);

\path[draw=drawColor,line width= 0.8pt,line join=round] (209.11, 57.55) -- (209.11, 55.42);

\path[draw=drawColor,line width= 0.8pt,line join=round] (209.36, 57.55) -- (209.36, 61.28);

\path[draw=drawColor,line width= 0.8pt,line join=round] (209.61, 57.55) -- (209.61, 59.20);

\path[draw=drawColor,line width= 0.8pt,line join=round] (209.86, 57.55) -- (209.86, 54.87);

\path[draw=drawColor,line width= 0.8pt,line join=round] (210.60, 57.55) -- (210.60, 55.85);

\path[draw=drawColor,line width= 0.8pt,line join=round] (210.85, 57.55) -- (210.85, 59.44);

\path[draw=drawColor,line width= 0.8pt,line join=round] (211.10, 57.55) -- (211.10, 62.65);

\path[draw=drawColor,line width= 0.8pt,line join=round] (211.35, 57.55) -- (211.35, 54.99);

\path[draw=drawColor,line width= 0.8pt,line join=round] (211.60, 57.55) -- (211.60, 55.75);

\path[draw=drawColor,line width= 0.8pt,line join=round] (212.34, 57.55) -- (212.34, 56.45);

\path[draw=drawColor,line width= 0.8pt,line join=round] (212.59, 57.55) -- (212.59, 59.84);

\path[draw=drawColor,line width= 0.8pt,line join=round] (212.84, 57.55) -- (212.84, 51.95);

\path[draw=drawColor,line width= 0.8pt,line join=round] (213.09, 57.55) -- (213.09, 59.50);

\path[draw=drawColor,line width= 0.8pt,line join=round] (213.33, 57.55) -- (213.33, 52.69);

\path[draw=drawColor,line width= 0.8pt,line join=round] (214.08, 57.55) -- (214.08, 57.87);

\path[draw=drawColor,line width= 0.8pt,line join=round] (214.33, 57.55) -- (214.33, 61.61);

\path[draw=drawColor,line width= 0.8pt,line join=round] (214.58, 57.55) -- (214.58, 56.26);

\path[draw=drawColor,line width= 0.8pt,line join=round] (214.82, 57.55) -- (214.82, 58.29);

\path[draw=drawColor,line width= 0.8pt,line join=round] (215.07, 57.55) -- (215.07, 58.04);

\path[draw=drawColor,line width= 0.8pt,line join=round] (216.07, 57.55) -- (216.07, 58.11);

\path[draw=drawColor,line width= 0.8pt,line join=round] (216.31, 57.55) -- (216.31, 54.07);

\path[draw=drawColor,line width= 0.8pt,line join=round] (216.56, 57.55) -- (216.56, 57.90);

\path[draw=drawColor,line width= 0.8pt,line join=round] (216.81, 57.55) -- (216.81, 56.17);

\path[draw=drawColor,line width= 0.8pt,line join=round] (217.80, 57.55) -- (217.80, 60.67);

\path[draw=drawColor,line width= 0.8pt,line join=round] (218.05, 57.55) -- (218.05, 56.50);

\path[draw=drawColor,line width= 0.8pt,line join=round] (218.30, 57.55) -- (218.30, 56.12);

\path[draw=drawColor,line width= 0.8pt,line join=round] (218.55, 57.55) -- (218.55, 57.65);

\path[draw=drawColor,line width= 0.8pt,line join=round] (219.29, 57.55) -- (219.29, 55.86);

\path[draw=drawColor,line width= 0.8pt,line join=round] (219.54, 57.55) -- (219.54, 61.20);

\path[draw=drawColor,line width= 0.8pt,line join=round] (219.79, 57.55) -- (219.79, 59.57);

\path[draw=drawColor,line width= 0.8pt,line join=round] (220.04, 57.55) -- (220.04, 55.92);

\path[draw=drawColor,line width= 0.8pt,line join=round] (220.29, 57.55) -- (220.29, 58.76);

\path[draw=drawColor,line width= 0.8pt,line join=round] (221.28, 57.55) -- (221.28, 55.29);

\path[draw=drawColor,line width= 0.8pt,line join=round] (221.53, 57.55) -- (221.53, 56.95);

\path[draw=drawColor,line width= 0.8pt,line join=round] (221.78, 57.55) -- (221.78, 57.67);

\path[draw=drawColor,line width= 0.8pt,line join=round] (222.02, 57.55) -- (222.02, 60.52);

\path[draw=drawColor,line width= 0.8pt,line join=round] (222.77, 57.55) -- (222.77, 57.25);

\path[draw=drawColor,line width= 0.8pt,line join=round] (223.02, 57.55) -- (223.02, 61.90);

\path[draw=drawColor,line width= 0.8pt,line join=round] (223.27, 57.55) -- (223.27, 62.12);

\path[draw=drawColor,line width= 0.8pt,line join=round] (223.51, 57.55) -- (223.51, 55.33);

\path[draw=drawColor,line width= 0.8pt,line join=round] (223.76, 57.55) -- (223.76, 61.75);

\path[draw=drawColor,line width= 0.8pt,line join=round] (224.51, 57.55) -- (224.51, 52.31);

\path[draw=drawColor,line width= 0.8pt,line join=round] (224.76, 57.55) -- (224.76, 54.81);

\path[draw=drawColor,line width= 0.8pt,line join=round] (225.00, 57.55) -- (225.00, 56.20);

\path[draw=drawColor,line width= 0.8pt,line join=round] (225.25, 57.55) -- (225.25, 54.24);

\path[draw=drawColor,line width= 0.8pt,line join=round] (225.50, 57.55) -- (225.50, 56.32);

\path[draw=drawColor,line width= 0.8pt,line join=round] (226.24, 57.55) -- (226.24, 56.52);

\path[draw=drawColor,line width= 0.8pt,line join=round] (226.49, 57.55) -- (226.49, 58.65);

\path[draw=drawColor,line width= 0.8pt,line join=round] (226.74, 57.55) -- (226.74, 56.54);

\path[draw=drawColor,line width= 0.8pt,line join=round] (226.99, 57.55) -- (226.99, 60.16);

\path[draw=drawColor,line width= 0.8pt,line join=round] (227.24, 57.55) -- (227.24, 63.30);

\path[draw=drawColor,line width= 0.8pt,line join=round] (227.98, 57.55) -- (227.98, 62.78);

\path[draw=drawColor,line width= 0.8pt,line join=round] (228.23, 57.55) -- (228.23, 56.57);

\path[draw=drawColor,line width= 0.8pt,line join=round] (228.48, 57.55) -- (228.48, 59.74);

\path[draw=drawColor,line width= 0.8pt,line join=round] (228.73, 57.55) -- (228.73, 55.75);

\path[draw=drawColor,line width= 0.8pt,line join=round] (228.98, 57.55) -- (228.98, 58.94);

\path[draw=drawColor,line width= 0.8pt,line join=round] (229.97, 57.55) -- (229.97, 55.77);

\path[draw=drawColor,line width= 0.8pt,line join=round] (230.22, 57.55) -- (230.22, 57.78);

\path[draw=drawColor,line width= 0.8pt,line join=round] (230.47, 57.55) -- (230.47, 51.17);

\path[draw=drawColor,line width= 0.8pt,line join=round] (230.71, 57.55) -- (230.71, 57.37);

\path[draw=drawColor,line width= 0.8pt,line join=round] (231.46, 57.55) -- (231.46, 62.28);

\path[draw=drawColor,line width= 0.8pt,line join=round] (231.71, 57.55) -- (231.71, 55.67);

\path[draw=drawColor,line width= 0.8pt,line join=round] (231.96, 57.55) -- (231.96, 62.13);

\path[draw=drawColor,line width= 0.8pt,line join=round] (232.20, 57.55) -- (232.20, 47.35);

\path[draw=drawColor,line width= 0.8pt,line join=round] (232.45, 57.55) -- (232.45, 59.41);

\path[draw=drawColor,line width= 0.8pt,line join=round] (233.20, 57.55) -- (233.20, 58.86);

\path[draw=drawColor,line width= 0.8pt,line join=round] (233.44, 57.55) -- (233.44, 58.20);

\path[draw=drawColor,line width= 0.8pt,line join=round] (233.69, 57.55) -- (233.69, 50.89);

\path[draw=drawColor,line width= 0.8pt,line join=round] (233.94, 57.55) -- (233.94, 52.90);

\path[draw=drawColor,line width= 0.8pt,line join=round] (234.19, 57.55) -- (234.19, 59.88);

\path[draw=drawColor,line width= 0.8pt,line join=round] (234.93, 57.55) -- (234.93, 58.38);

\path[draw=drawColor,line width= 0.8pt,line join=round] (235.18, 57.55) -- (235.18, 55.58);

\path[draw=drawColor,line width= 0.8pt,line join=round] (235.43, 57.55) -- (235.43, 61.32);

\path[draw=drawColor,line width= 0.8pt,line join=round] (235.68, 57.55) -- (235.68, 56.27);

\path[draw=drawColor,line width= 0.8pt,line join=round] (235.93, 57.55) -- (235.93, 57.70);

\path[draw=drawColor,line width= 0.8pt,line join=round] (236.67, 57.55) -- (236.67, 63.71);

\path[draw=drawColor,line width= 0.8pt,line join=round] (236.92, 57.55) -- (236.92, 50.18);

\path[draw=drawColor,line width= 0.8pt,line join=round] (237.17, 57.55) -- (237.17, 57.55);

\path[draw=drawColor,line width= 0.8pt,line join=round] (237.42, 57.55) -- (237.42, 57.04);

\path[draw=drawColor,line width= 0.8pt,line join=round] (237.67, 57.55) -- (237.67, 57.34);

\path[draw=drawColor,line width= 0.8pt,line join=round] (238.41, 57.55) -- (238.41, 55.92);

\path[draw=drawColor,line width= 0.8pt,line join=round] (238.66, 57.55) -- (238.66, 61.18);

\path[draw=drawColor,line width= 0.8pt,line join=round] (238.91, 57.55) -- (238.91, 57.32);

\path[draw=drawColor,line width= 0.8pt,line join=round] (239.16, 57.55) -- (239.16, 58.45);

\path[draw=drawColor,line width= 0.8pt,line join=round] (239.40, 57.55) -- (239.40, 56.39);

\path[draw=drawColor,line width= 0.8pt,line join=round] (240.15, 57.55) -- (240.15, 56.31);

\path[draw=drawColor,line width= 0.8pt,line join=round] (240.40, 57.55) -- (240.40, 62.20);

\path[draw=drawColor,line width= 0.8pt,line join=round] (240.65, 57.55) -- (240.65, 57.82);

\path[draw=drawColor,line width= 0.8pt,line join=round] (240.89, 57.55) -- (240.89, 61.42);

\path[draw=drawColor,line width= 0.8pt,line join=round] (241.14, 57.55) -- (241.14, 56.82);

\path[draw=drawColor,line width= 0.8pt,line join=round] (241.89, 57.55) -- (241.89, 61.46);

\path[draw=drawColor,line width= 0.8pt,line join=round] (242.13, 57.55) -- (242.13, 57.45);

\path[draw=drawColor,line width= 0.8pt,line join=round] (242.38, 57.55) -- (242.38, 52.11);

\path[draw=drawColor,line width= 0.8pt,line join=round] (242.63, 57.55) -- (242.63, 53.18);

\path[draw=drawColor,line width= 0.8pt,line join=round] (243.62, 57.55) -- (243.62, 60.13);

\path[draw=drawColor,line width= 0.8pt,line join=round] (243.87, 57.55) -- (243.87, 58.16);

\path[draw=drawColor,line width= 0.8pt,line join=round] (244.12, 57.55) -- (244.12, 55.26);

\path[draw=drawColor,line width= 0.8pt,line join=round] (244.37, 57.55) -- (244.37, 60.49);

\path[draw=drawColor,line width= 0.8pt,line join=round] (244.62, 57.55) -- (244.62, 56.71);

\path[draw=drawColor,line width= 0.8pt,line join=round] (245.36, 57.55) -- (245.36, 63.62);

\path[draw=drawColor,line width= 0.8pt,line join=round] (245.61, 57.55) -- (245.61, 75.18);

\path[draw=drawColor,line width= 0.8pt,line join=round] (245.86, 57.55) -- (245.86, 58.08);

\path[draw=drawColor,line width= 0.8pt,line join=round] (246.11, 57.55) -- (246.11, 53.10);

\path[draw=drawColor,line width= 0.8pt,line join=round] (246.36, 57.55) -- (246.36, 56.59);

\path[draw=drawColor,line width= 0.8pt,line join=round] (247.10, 57.55) -- (247.10, 56.95);

\path[draw=drawColor,line width= 0.8pt,line join=round] (247.35, 57.55) -- (247.35, 56.42);

\path[draw=drawColor,line width= 0.8pt,line join=round] (247.60, 57.55) -- (247.60, 57.61);

\path[draw=drawColor,line width= 0.8pt,line join=round] (247.85, 57.55) -- (247.85, 52.56);

\path[draw=drawColor,line width= 0.8pt,line join=round] (248.09, 57.55) -- (248.09, 58.13);

\path[draw=drawColor,line width= 0.8pt,line join=round] (248.84, 57.55) -- (248.84, 56.91);

\path[draw=drawColor,line width= 0.8pt,line join=round] (249.09, 57.55) -- (249.09, 57.36);

\path[draw=drawColor,line width= 0.8pt,line join=round] (249.34, 57.55) -- (249.34, 60.04);

\path[draw=drawColor,line width= 0.8pt,line join=round] (249.58, 57.55) -- (249.58, 59.00);

\path[draw=drawColor,line width= 0.8pt,line join=round] (249.83, 57.55) -- (249.83, 56.90);

\path[draw=drawColor,line width= 0.8pt,line join=round] (250.58, 57.55) -- (250.58, 61.47);

\path[draw=drawColor,line width= 0.8pt,line join=round] (250.82, 57.55) -- (250.82, 57.36);

\path[draw=drawColor,line width= 0.8pt,line join=round] (251.07, 57.55) -- (251.07, 52.48);

\path[draw=drawColor,line width= 0.8pt,line join=round] (251.32, 57.55) -- (251.32, 57.73);

\path[draw=drawColor,line width= 0.8pt,line join=round] (251.57, 57.55) -- (251.57, 62.63);

\path[draw=drawColor,line width= 0.8pt,line join=round] (252.31, 57.55) -- (252.31, 57.23);

\path[draw=drawColor,line width= 0.8pt,line join=round] (252.56, 57.55) -- (252.56, 60.58);

\path[draw=drawColor,line width= 0.8pt,line join=round] (252.81, 57.55) -- (252.81, 58.26);

\path[draw=drawColor,line width= 0.8pt,line join=round] (253.06, 57.55) -- (253.06, 59.93);

\path[draw=drawColor,line width= 0.8pt,line join=round] (253.31, 57.55) -- (253.31, 58.96);

\path[draw=drawColor,line width= 0.8pt,line join=round] (254.30, 57.55) -- (254.30, 57.13);

\path[draw=drawColor,line width= 0.8pt,line join=round] (254.55, 57.55) -- (254.55, 57.46);

\path[draw=drawColor,line width= 0.8pt,line join=round] (254.80, 57.55) -- (254.80, 58.03);

\path[draw=drawColor,line width= 0.8pt,line join=round] (255.05, 57.55) -- (255.05, 58.21);

\path[draw=drawColor,line width= 0.8pt,line join=round] (255.79, 57.55) -- (255.79, 55.89);

\path[draw=drawColor,line width= 0.8pt,line join=round] (256.04, 57.55) -- (256.04, 56.10);

\path[draw=drawColor,line width= 0.8pt,line join=round] (256.29, 57.55) -- (256.29, 55.22);

\path[draw=drawColor,line width= 0.8pt,line join=round] (256.54, 57.55) -- (256.54, 60.90);

\path[draw=drawColor,line width= 0.8pt,line join=round] (256.78, 57.55) -- (256.78, 59.56);

\path[draw=drawColor,line width= 0.8pt,line join=round] (257.53, 57.55) -- (257.53, 57.88);

\path[draw=drawColor,line width= 0.8pt,line join=round] (257.78, 57.55) -- (257.78, 56.84);

\path[draw=drawColor,line width= 0.8pt,line join=round] (258.02, 57.55) -- (258.02, 55.24);

\path[draw=drawColor,line width= 0.8pt,line join=round] (258.27, 57.55) -- (258.27, 61.28);

\path[draw=drawColor,line width= 0.8pt,line join=round] (258.52, 57.55) -- (258.52, 60.16);

\path[draw=drawColor,line width= 0.8pt,line join=round] (259.27, 57.55) -- (259.27, 57.31);

\path[draw=drawColor,line width= 0.8pt,line join=round] (259.51, 57.55) -- (259.51, 56.57);

\path[draw=drawColor,line width= 0.8pt,line join=round] (259.76, 57.55) -- (259.76, 49.64);

\path[draw=drawColor,line width= 0.8pt,line join=round] (260.01, 57.55) -- (260.01, 58.38);

\path[draw=drawColor,line width= 0.8pt,line join=round] (260.26, 57.55) -- (260.26, 58.15);

\path[draw=drawColor,line width= 0.8pt,line join=round] (261.00, 57.55) -- (261.00, 57.86);

\path[draw=drawColor,line width= 0.8pt,line join=round] (261.25, 57.55) -- (261.25, 56.83);

\path[draw=drawColor,line width= 0.8pt,line join=round] (261.50, 57.55) -- (261.50, 63.13);

\path[draw=drawColor,line width= 0.8pt,line join=round] (261.75, 57.55) -- (261.75, 55.77);

\path[draw=drawColor,line width= 0.8pt,line join=round] (262.00, 57.55) -- (262.00, 61.53);

\path[draw=drawColor,line width= 0.8pt,line join=round] (262.74, 57.55) -- (262.74, 56.48);

\path[draw=drawColor,line width= 0.8pt,line join=round] (263.24, 57.55) -- (263.24, 58.68);

\path[draw=drawColor,line width= 0.8pt,line join=round] (263.49, 57.55) -- (263.49, 55.41);

\path[draw=drawColor,line width= 0.8pt,line join=round] (263.74, 57.55) -- (263.74, 58.44);

\path[draw=drawColor,line width= 0.8pt,line join=round] (264.48, 57.55) -- (264.48, 59.54);

\path[draw=drawColor,line width= 0.8pt,line join=round] (264.73, 57.55) -- (264.73, 58.41);

\path[draw=drawColor,line width= 0.8pt,line join=round] (264.98, 57.55) -- (264.98, 59.09);

\path[draw=drawColor,line width= 0.8pt,line join=round] (265.23, 57.55) -- (265.23, 56.50);

\path[draw=drawColor,line width= 0.8pt,line join=round] (265.47, 57.55) -- (265.47, 58.19);

\path[draw=drawColor,line width= 0.8pt,line join=round] (266.22, 57.55) -- (266.22, 57.70);

\path[draw=drawColor,line width= 0.8pt,line join=round] (266.47, 57.55) -- (266.47, 53.73);

\path[draw=drawColor,line width= 0.8pt,line join=round] (266.71, 57.55) -- (266.71, 59.24);

\path[draw=drawColor,line width= 0.8pt,line join=round] (266.96, 57.55) -- (266.96, 56.40);

\path[draw=drawColor,line width= 0.8pt,line join=round] (267.21, 57.55) -- (267.21, 55.96);

\path[draw=drawColor,line width= 0.8pt,line join=round] (267.96, 57.55) -- (267.96, 61.00);

\path[draw=drawColor,line width= 0.8pt,line join=round] (268.20, 57.55) -- (268.20, 70.86);

\path[draw=drawColor,line width= 0.8pt,line join=round] (268.45, 57.55) -- (268.45, 55.46);

\path[draw=drawColor,line width= 0.8pt,line join=round] (268.70, 57.55) -- (268.70, 59.28);

\path[draw=drawColor,line width= 0.8pt,line join=round] (268.95, 57.55) -- (268.95, 56.99);

\path[draw=drawColor,line width= 0.8pt,line join=round] (269.69, 57.55) -- (269.69, 57.24);

\path[draw=drawColor,line width= 0.8pt,line join=round] (269.94, 57.55) -- (269.94, 55.80);

\path[draw=drawColor,line width= 0.8pt,line join=round] (270.19, 57.55) -- (270.19, 57.53);

\path[draw=drawColor,line width= 0.8pt,line join=round] (270.44, 57.55) -- (270.44, 58.51);

\path[draw=drawColor,line width= 0.8pt,line join=round] (270.69, 57.55) -- (270.69, 59.16);

\path[draw=drawColor,line width= 0.8pt,line join=round] (271.43, 57.55) -- (271.43, 58.28);

\path[draw=drawColor,line width= 0.8pt,line join=round] (271.68, 57.55) -- (271.68, 56.94);

\path[draw=drawColor,line width= 0.8pt,line join=round] (271.93, 57.55) -- (271.93, 57.40);

\path[draw=drawColor,line width= 0.8pt,line join=round] (272.18, 57.55) -- (272.18, 55.11);

\path[draw=drawColor,line width= 0.8pt,line join=round] (272.43, 57.55) -- (272.43, 56.93);

\path[draw=drawColor,line width= 0.8pt,line join=round] (273.17, 57.55) -- (273.17, 59.31);

\path[draw=drawColor,line width= 0.8pt,line join=round] (273.42, 57.55) -- (273.42, 57.42);

\path[draw=drawColor,line width= 0.8pt,line join=round] (273.67, 57.55) -- (273.67, 58.51);

\path[draw=drawColor,line width= 0.8pt,line join=round] (273.92, 57.55) -- (273.92, 55.54);

\path[draw=drawColor,line width= 0.8pt,line join=round] (274.16, 57.55) -- (274.16, 59.15);

\path[draw=drawColor,line width= 0.8pt,line join=round] (274.91, 57.55) -- (274.91, 56.89);

\path[draw=drawColor,line width= 0.8pt,line join=round] (275.16, 57.55) -- (275.16, 62.06);

\path[draw=drawColor,line width= 0.8pt,line join=round] (275.40, 57.55) -- (275.40, 55.38);

\path[draw=drawColor,line width= 0.8pt,line join=round] (275.65, 57.55) -- (275.65, 57.61);

\path[draw=drawColor,line width= 0.8pt,line join=round] (275.90, 57.55) -- (275.90, 58.70);

\path[draw=drawColor,line width= 0.8pt,line join=round] (276.65, 57.55) -- (276.65, 56.98);

\path[draw=drawColor,line width= 0.8pt,line join=round] (276.89, 57.55) -- (276.89, 59.44);

\path[draw=drawColor,line width= 0.8pt,line join=round] (277.14, 57.55) -- (277.14, 60.62);

\path[draw=drawColor,line width= 0.8pt,line join=round] (277.39, 57.55) -- (277.39, 57.43);

\path[draw=drawColor,line width= 0.8pt,line join=round] (277.64, 57.55) -- (277.64, 59.11);

\path[draw=drawColor,line width= 0.8pt,line join=round] (278.63, 57.55) -- (278.63, 57.59);

\path[draw=drawColor,line width= 0.8pt,line join=round] (278.88, 57.55) -- (278.88, 54.93);

\path[draw=drawColor,line width= 0.8pt,line join=round] (279.13, 57.55) -- (279.13, 59.11);

\path[draw=drawColor,line width= 0.8pt,line join=round] (279.38, 57.55) -- (279.38, 57.64);

\path[draw=drawColor,line width= 0.8pt,line join=round] (280.12, 57.55) -- (280.12, 59.61);

\path[draw=drawColor,line width= 0.8pt,line join=round] (280.37, 57.55) -- (280.37, 61.57);

\path[draw=drawColor,line width= 0.8pt,line join=round] (280.62, 57.55) -- (280.62, 55.01);

\path[draw=drawColor,line width= 0.8pt,line join=round] (280.87, 57.55) -- (280.87, 59.20);

\path[draw=drawColor,line width= 0.8pt,line join=round] (281.12, 57.55) -- (281.12, 59.26);

\path[draw=drawColor,line width= 0.8pt,line join=round] (281.86, 57.55) -- (281.86, 62.22);

\path[draw=drawColor,line width= 0.8pt,line join=round] (282.11, 57.55) -- (282.11, 59.25);

\path[draw=drawColor,line width= 0.8pt,line join=round] (282.36, 57.55) -- (282.36, 57.75);

\path[draw=drawColor,line width= 0.8pt,line join=round] (282.60, 57.55) -- (282.60, 57.43);

\path[draw=drawColor,line width= 0.8pt,line join=round] (282.85, 57.55) -- (282.85, 56.87);

\path[draw=drawColor,line width= 0.8pt,line join=round] (283.60, 57.55) -- (283.60, 57.34);

\path[draw=drawColor,line width= 0.8pt,line join=round] (283.85, 57.55) -- (283.85, 57.90);

\path[draw=drawColor,line width= 0.8pt,line join=round] (284.09, 57.55) -- (284.09, 57.56);

\path[draw=drawColor,line width= 0.8pt,line join=round] (284.34, 57.55) -- (284.34, 58.87);

\path[draw=drawColor,line width= 0.8pt,line join=round] (284.59, 57.55) -- (284.59, 56.58);

\path[draw=drawColor,line width= 0.8pt,line join=round] (285.34, 57.55) -- (285.34, 57.56);

\path[draw=drawColor,line width= 0.8pt,line join=round] (285.58, 57.55) -- (285.58, 59.02);

\path[draw=drawColor,line width= 0.8pt,line join=round] (285.83, 57.55) -- (285.83, 59.82);

\path[draw=drawColor,line width= 0.8pt,line join=round] (286.08, 57.55) -- (286.08, 56.92);

\path[draw=drawColor,line width= 0.8pt,line join=round] (286.33, 57.55) -- (286.33, 58.52);

\path[draw=drawColor,line width= 0.8pt,line join=round] (287.07, 57.55) -- (287.07, 57.45);

\path[draw=drawColor,line width= 0.8pt,line join=round] (287.32, 57.55) -- (287.32, 59.94);

\path[draw=drawColor,line width= 0.8pt,line join=round] (287.57, 57.55) -- (287.57, 58.29);

\path[draw=drawColor,line width= 0.8pt,line join=round] (287.82, 57.55) -- (287.82, 60.05);

\path[draw=drawColor,line width= 0.8pt,line join=round] (288.07, 57.55) -- (288.07, 58.83);

\path[draw=drawColor,line width= 0.8pt,line join=round] (288.81, 57.55) -- (288.81, 58.90);

\path[draw=drawColor,line width= 0.8pt,line join=round] (289.06, 57.55) -- (289.06, 55.89);

\path[draw=drawColor,line width= 0.8pt,line join=round] (289.31, 57.55) -- (289.31, 58.88);

\path[draw=drawColor,line width= 0.8pt,line join=round] (289.56, 57.55) -- (289.56, 58.01);

\path[draw=drawColor,line width= 0.8pt,line join=round] (289.81, 57.55) -- (289.81, 57.21);

\path[draw=drawColor,line width= 0.8pt,line join=round] (290.55, 57.55) -- (290.55, 58.11);

\path[draw=drawColor,line width= 0.8pt,line join=round] (290.80, 57.55) -- (290.80, 68.87);

\path[draw=drawColor,line width= 0.8pt,line join=round] (291.05, 57.55) -- (291.05, 55.17);

\path[draw=drawColor,line width= 0.8pt,line join=round] (291.29, 57.55) -- (291.29, 57.72);

\path[draw=drawColor,line width= 0.8pt,line join=round] (291.54, 57.55) -- (291.54, 59.02);

\path[draw=drawColor,line width= 0.8pt,line join=round] (292.29, 57.55) -- (292.29, 55.30);

\path[draw=drawColor,line width= 0.8pt,line join=round] (292.54, 57.55) -- (292.54, 56.36);

\path[draw=drawColor,line width= 0.8pt,line join=round] (292.78, 57.55) -- (292.78, 58.38);

\path[draw=drawColor,line width= 0.8pt,line join=round] (293.03, 57.55) -- (293.03, 57.85);

\path[draw=drawColor,line width= 0.8pt,line join=round] (293.28, 57.55) -- (293.28, 57.82);

\path[draw=drawColor,line width= 0.8pt,line join=round] (294.03, 57.55) -- (294.03, 59.38);

\path[draw=drawColor,line width= 0.8pt,line join=round] (294.27, 57.55) -- (294.27, 59.40);

\path[draw=drawColor,line width= 0.8pt,line join=round] (294.52, 57.55) -- (294.52, 54.98);

\path[draw=drawColor,line width= 0.8pt,line join=round] (294.77, 57.55) -- (294.77, 54.21);

\path[draw=drawColor,line width= 0.8pt,line join=round] (295.02, 57.55) -- (295.02, 59.50);

\path[draw=drawColor,line width= 0.8pt,line join=round] (295.76, 57.55) -- (295.76, 57.63);

\path[draw=drawColor,line width= 0.8pt,line join=round] (296.01, 57.55) -- (296.01, 59.26);

\path[draw=drawColor,line width= 0.8pt,line join=round] (296.26, 57.55) -- (296.26, 51.64);

\path[draw=drawColor,line width= 0.8pt,line join=round] (296.51, 57.55) -- (296.51, 61.44);

\path[draw=drawColor,line width= 0.8pt,line join=round] (296.76, 57.55) -- (296.76, 57.15);

\path[draw=drawColor,line width= 0.8pt,line join=round] (297.50, 57.55) -- (297.50, 58.88);

\path[draw=drawColor,line width= 0.8pt,line join=round] (297.75, 57.55) -- (297.75, 58.72);

\path[draw=drawColor,line width= 0.8pt,line join=round] (298.00, 57.55) -- (298.00, 58.24);

\path[draw=drawColor,line width= 0.8pt,line join=round] (298.50, 57.55) -- (298.50, 56.50);

\path[draw=drawColor,line width= 0.8pt,line join=round] (299.24, 57.55) -- (299.24, 57.21);

\path[draw=drawColor,line width= 0.8pt,line join=round] (299.49, 57.55) -- (299.49, 60.58);

\path[draw=drawColor,line width= 0.8pt,line join=round] (299.74, 57.55) -- (299.74, 55.98);

\path[draw=drawColor,line width= 0.8pt,line join=round] (299.98, 57.55) -- (299.98, 62.70);

\path[draw=drawColor,line width= 0.8pt,line join=round] (300.23, 57.55) -- (300.23, 58.16);

\path[draw=drawColor,line width= 0.8pt,line join=round] (300.98, 57.55) -- (300.98, 57.51);

\path[draw=drawColor,line width= 0.8pt,line join=round] (301.23, 57.55) -- (301.23, 55.30);

\path[draw=drawColor,line width= 0.8pt,line join=round] (301.47, 57.55) -- (301.47, 57.88);

\path[draw=drawColor,line width= 0.8pt,line join=round] (301.72, 57.55) -- (301.72, 61.74);

\path[draw=drawColor,line width= 0.8pt,line join=round] (301.97, 57.55) -- (301.97, 59.12);

\path[draw=drawColor,line width= 0.8pt,line join=round] (302.72, 57.55) -- (302.72, 57.40);

\path[draw=drawColor,line width= 0.8pt,line join=round] (302.96, 57.55) -- (302.96, 56.98);

\path[draw=drawColor,line width= 0.8pt,line join=round] (303.21, 57.55) -- (303.21, 65.77);

\path[draw=drawColor,line width= 0.8pt,line join=round] (303.46, 57.55) -- (303.46, 53.93);

\path[draw=drawColor,line width= 0.8pt,line join=round] (303.71, 57.55) -- (303.71, 58.20);

\path[draw=drawColor,line width= 0.8pt,line join=round] (304.45, 57.55) -- (304.45, 62.39);

\path[draw=drawColor,line width= 0.8pt,line join=round] (304.70, 57.55) -- (304.70, 59.31);

\path[draw=drawColor,line width= 0.8pt,line join=round] (304.95, 57.55) -- (304.95, 60.02);

\path[draw=drawColor,line width= 0.8pt,line join=round] (305.20, 57.55) -- (305.20, 60.76);

\path[draw=drawColor,line width= 0.8pt,line join=round] (305.45, 57.55) -- (305.45, 59.21);

\path[draw=drawColor,line width= 0.8pt,line join=round] (306.44, 57.55) -- (306.44, 58.58);

\path[draw=drawColor,line width= 0.8pt,line join=round] (306.69, 57.55) -- (306.69, 59.15);

\path[draw=drawColor,line width= 0.8pt,line join=round] (306.94, 57.55) -- (306.94, 58.87);

\path[draw=drawColor,line width= 0.8pt,line join=round] (307.19, 57.55) -- (307.19, 56.31);

\path[draw=drawColor,line width= 0.8pt,line join=round] (308.18, 57.55) -- (308.18, 56.75);

\path[draw=drawColor,line width= 0.8pt,line join=round] (308.43, 57.55) -- (308.43, 57.90);

\path[draw=drawColor,line width= 0.8pt,line join=round] (308.67, 57.55) -- (308.67, 60.72);

\path[draw=drawColor,line width= 0.8pt,line join=round] (308.92, 57.55) -- (308.92, 61.20);

\path[draw=drawColor,line width= 0.8pt,line join=round] (309.67, 57.55) -- (309.67, 63.33);

\path[draw=drawColor,line width= 0.8pt,line join=round] (309.92, 57.55) -- (309.92, 58.11);

\path[draw=drawColor,line width= 0.8pt,line join=round] (310.16, 57.55) -- (310.16, 56.76);

\path[draw=drawColor,line width= 0.8pt,line join=round] (310.41, 57.55) -- (310.41, 62.18);

\path[draw=drawColor,line width= 0.8pt,line join=round] (310.66, 57.55) -- (310.66, 59.06);

\path[draw=drawColor,line width= 0.8pt,line join=round] (311.65, 57.55) -- (311.65, 56.19);

\path[draw=drawColor,line width= 0.8pt,line join=round] (311.90, 57.55) -- (311.90, 56.43);

\path[draw=drawColor,line width= 0.8pt,line join=round] (312.15, 57.55) -- (312.15, 58.00);

\path[draw=drawColor,line width= 0.8pt,line join=round] (312.40, 57.55) -- (312.40, 59.72);

\path[draw=drawColor,line width= 0.8pt,line join=round] (313.14, 57.55) -- (313.14, 58.20);

\path[draw=drawColor,line width= 0.8pt,line join=round] (313.39, 57.55) -- (313.39, 55.55);

\path[draw=drawColor,line width= 0.8pt,line join=round] (313.64, 57.55) -- (313.64, 56.04);

\path[draw=drawColor,line width= 0.8pt,line join=round] (313.89, 57.55) -- (313.89, 58.97);

\path[draw=drawColor,line width= 0.8pt,line join=round] (314.14, 57.55) -- (314.14, 54.35);

\path[draw=drawColor,line width= 0.8pt,line join=round] (314.88, 57.55) -- (314.88, 51.21);

\path[draw=drawColor,line width= 0.8pt,line join=round] (315.13, 57.55) -- (315.13, 59.23);

\path[draw=drawColor,line width= 0.8pt,line join=round] (315.38, 57.55) -- (315.38, 56.15);

\path[draw=drawColor,line width= 0.8pt,line join=round] (315.63, 57.55) -- (315.63, 56.77);

\path[draw=drawColor,line width= 0.8pt,line join=round] (315.87, 57.55) -- (315.87, 50.62);

\path[draw=drawColor,line width= 0.8pt,line join=round] (316.62, 57.55) -- (316.62, 47.86);

\path[draw=drawColor,line width= 0.8pt,line join=round] (316.87, 57.55) -- (316.87, 65.62);

\path[draw=drawColor,line width= 0.8pt,line join=round] (317.12, 57.55) -- (317.12, 54.44);

\path[draw=drawColor,line width= 0.8pt,line join=round] (317.36, 57.55) -- (317.36, 44.59);

\path[draw=drawColor,line width= 0.8pt,line join=round] (317.61, 57.55) -- (317.61, 62.63);

\path[draw=drawColor,line width= 0.8pt,line join=round] (318.36, 57.55) -- (318.36, 62.31);

\path[draw=drawColor,line width= 0.8pt,line join=round] (318.61, 57.55) -- (318.61, 60.10);

\path[draw=drawColor,line width= 0.8pt,line join=round] (318.85, 57.55) -- (318.85, 63.65);

\path[draw=drawColor,line width= 0.8pt,line join=round] (319.10, 57.55) -- (319.10, 60.36);

\path[draw=drawColor,line width= 0.8pt,line join=round] (319.35, 57.55) -- (319.35, 52.11);

\path[draw=drawColor,line width= 0.8pt,line join=round] (320.34, 57.55) -- (320.34, 56.45);

\path[draw=drawColor,line width= 0.8pt,line join=round] (320.59, 57.55) -- (320.59, 57.05);

\path[draw=drawColor,line width= 0.8pt,line join=round] (320.84, 57.55) -- (320.84, 62.93);

\path[draw=drawColor,line width= 0.8pt,line join=round] (321.09, 57.55) -- (321.09, 62.70);

\path[draw=drawColor,line width= 0.8pt,line join=round] (321.83, 57.55) -- (321.83, 59.37);

\path[draw=drawColor,line width= 0.8pt,line join=round] (322.08, 57.55) -- (322.08, 54.06);

\path[draw=drawColor,line width= 0.8pt,line join=round] (322.33, 57.55) -- (322.33, 47.76);

\path[draw=drawColor,line width= 0.8pt,line join=round] (322.58, 57.55) -- (322.58, 50.82);

\path[draw=drawColor,line width= 0.8pt,line join=round] (322.83, 57.55) -- (322.83, 51.50);

\path[draw=drawColor,line width= 0.8pt,line join=round] (323.57, 57.55) -- (323.57, 64.97);

\path[draw=drawColor,line width= 0.8pt,line join=round] (323.82, 57.55) -- (323.82, 61.57);

\path[draw=drawColor,line width= 0.8pt,line join=round] (324.07, 57.55) -- (324.07, 54.13);

\path[draw=drawColor,line width= 0.8pt,line join=round] (324.32, 57.55) -- (324.32, 60.72);

\path[draw=drawColor,line width= 0.8pt,line join=round] (324.56, 57.55) -- (324.56, 64.51);

\path[draw=drawColor,line width= 0.8pt,line join=round] (325.31, 57.55) -- (325.31, 51.96);

\path[draw=drawColor,line width= 0.8pt,line join=round] (325.56, 57.55) -- (325.56, 56.32);

\path[draw=drawColor,line width= 0.8pt,line join=round] (325.81, 57.55) -- (325.81, 55.80);

\path[draw=drawColor,line width= 0.8pt,line join=round] (326.05, 57.55) -- (326.05, 60.63);

\path[draw=drawColor,line width= 0.8pt,line join=round] (326.30, 57.55) -- (326.30, 60.38);

\path[draw=drawColor,line width= 0.8pt,line join=round] (327.05, 57.55) -- (327.05, 51.04);

\path[draw=drawColor,line width= 0.8pt,line join=round] (327.30, 57.55) -- (327.30, 60.45);

\path[draw=drawColor,line width= 0.8pt,line join=round] (327.54, 57.55) -- (327.54, 60.16);

\path[draw=drawColor,line width= 0.8pt,line join=round] (327.79, 57.55) -- (327.79, 43.84);

\path[draw=drawColor,line width= 0.8pt,line join=round] (328.04, 57.55) -- (328.04, 53.37);

\path[draw=drawColor,line width= 0.8pt,line join=round] (328.79, 57.55) -- (328.79, 65.33);

\path[draw=drawColor,line width= 0.8pt,line join=round] (329.03, 57.55) -- (329.03, 54.08);

\path[draw=drawColor,line width= 0.8pt,line join=round] (329.28, 57.55) -- (329.28, 54.63);

\path[draw=drawColor,line width= 0.8pt,line join=round] (329.53, 57.55) -- (329.53, 61.08);

\path[draw=drawColor,line width= 0.8pt,line join=round] (330.52, 57.55) -- (330.52, 51.96);

\path[draw=drawColor,line width= 0.8pt,line join=round] (330.77, 57.55) -- (330.77, 59.43);

\path[draw=drawColor,line width= 0.8pt,line join=round] (331.02, 57.55) -- (331.02, 57.74);

\path[draw=drawColor,line width= 0.8pt,line join=round] (331.27, 57.55) -- (331.27, 62.23);

\path[draw=drawColor,line width= 0.8pt,line join=round] (331.52, 57.55) -- (331.52, 49.32);

\path[draw=drawColor,line width= 0.8pt,line join=round] (332.26, 57.55) -- (332.26, 57.69);

\path[draw=drawColor,line width= 0.8pt,line join=round] (332.51, 57.55) -- (332.51, 65.56);

\path[draw=drawColor,line width= 0.8pt,line join=round] (332.76, 57.55) -- (332.76, 55.49);

\path[draw=drawColor,line width= 0.8pt,line join=round] (333.01, 57.55) -- (333.01, 62.80);

\path[draw=drawColor,line width= 0.8pt,line join=round] (333.25, 57.55) -- (333.25, 57.69);

\path[draw=drawColor,line width= 0.8pt,line join=round] (334.00, 57.55) -- (334.00, 60.49);

\path[draw=drawColor,line width= 0.8pt,line join=round] (334.25, 57.55) -- (334.25, 59.33);

\path[draw=drawColor,line width= 0.8pt,line join=round] (334.50, 57.55) -- (334.50, 61.52);

\path[draw=drawColor,line width= 0.8pt,line join=round] (334.74, 57.55) -- (334.74, 57.17);

\path[draw=drawColor,line width= 0.8pt,line join=round] (334.99, 57.55) -- (334.99, 53.85);

\path[draw=drawColor,line width= 0.8pt,line join=round] (335.74, 57.55) -- (335.74, 58.67);

\path[draw=drawColor,line width= 0.8pt,line join=round] (335.99, 57.55) -- (335.99, 42.63);

\path[draw=drawColor,line width= 0.8pt,line join=round] (336.23, 57.55) -- (336.23, 58.75);

\path[draw=drawColor,line width= 0.8pt,line join=round] (336.48, 57.55) -- (336.48, 59.34);

\path[draw=drawColor,line width= 0.8pt,line join=round] (336.73, 57.55) -- (336.73, 55.00);

\path[draw=drawColor,line width= 0.8pt,line join=round] (337.48, 57.55) -- (337.48, 56.97);

\path[draw=drawColor,line width= 0.8pt,line join=round] (337.72, 57.55) -- (337.72, 57.64);

\path[draw=drawColor,line width= 0.8pt,line join=round] (337.97, 57.55) -- (337.97, 56.64);

\path[draw=drawColor,line width= 0.8pt,line join=round] (338.22, 57.55) -- (338.22, 57.42);

\path[draw=drawColor,line width= 0.8pt,line join=round] (338.47, 57.55) -- (338.47, 61.71);

\path[draw=drawColor,line width= 0.8pt,line join=round] (339.21, 57.55) -- (339.21, 62.97);

\path[draw=drawColor,line width= 0.8pt,line join=round] (339.46, 57.55) -- (339.46, 57.72);

\path[draw=drawColor,line width= 0.8pt,line join=round] (339.71, 57.55) -- (339.71, 61.67);

\path[draw=drawColor,line width= 0.8pt,line join=round] (339.96, 57.55) -- (339.96, 59.61);

\path[draw=drawColor,line width= 0.8pt,line join=round] (340.21, 57.55) -- (340.21, 58.89);

\path[draw=drawColor,line width= 0.8pt,line join=round] (340.95, 57.55) -- (340.95, 58.18);

\path[draw=drawColor,line width= 0.8pt,line join=round] (341.20, 57.55) -- (341.20, 53.47);

\path[draw=drawColor,line width= 0.8pt,line join=round] (341.45, 57.55) -- (341.45, 59.39);

\path[draw=drawColor,line width= 0.8pt,line join=round] (341.70, 57.55) -- (341.70, 57.38);

\path[draw=drawColor,line width= 0.8pt,line join=round] (341.94, 57.55) -- (341.94, 60.59);

\path[draw=drawColor,line width= 0.8pt,line join=round] (342.69, 57.55) -- (342.69, 62.30);

\path[draw=drawColor,line width= 0.8pt,line join=round] (342.94, 57.55) -- (342.94, 53.52);

\path[draw=drawColor,line width= 0.8pt,line join=round] (343.19, 57.55) -- (343.19, 59.83);

\path[draw=drawColor,line width= 0.8pt,line join=round] (343.43, 57.55) -- (343.43, 57.49);

\path[draw=drawColor,line width= 0.8pt,line join=round] (343.68, 57.55) -- (343.68, 54.80);

\path[draw=drawColor,line width= 0.8pt,line join=round] (344.68, 57.55) -- (344.68, 54.19);

\path[draw=drawColor,line width= 0.8pt,line join=round] (344.92, 57.55) -- (344.92, 60.32);

\path[draw=drawColor,line width= 0.8pt,line join=round] (345.17, 57.55) -- (345.17, 52.16);

\path[draw=drawColor,line width= 0.8pt,line join=round] (345.42, 57.55) -- (345.42, 60.00);

\path[draw=drawColor,line width= 0.8pt,line join=round] (346.17, 57.55) -- (346.17, 55.99);

\path[draw=drawColor,line width= 0.8pt,line join=round] (346.41, 57.55) -- (346.41, 58.32);

\path[draw=drawColor,line width= 0.8pt,line join=round] (346.66, 57.55) -- (346.66, 62.82);

\path[draw=drawColor,line width= 0.8pt,line join=round] (346.91, 57.55) -- (346.91, 57.22);

\path[draw=drawColor,line width= 0.8pt,line join=round] (347.16, 57.55) -- (347.16, 56.10);

\path[draw=drawColor,line width= 0.8pt,line join=round] (347.90, 57.55) -- (347.90, 59.01);

\path[draw=drawColor,line width= 0.8pt,line join=round] (348.15, 57.55) -- (348.15, 59.31);

\path[draw=drawColor,line width= 0.8pt,line join=round] (348.40, 57.55) -- (348.40, 53.40);

\path[draw=drawColor,line width= 0.8pt,line join=round] (348.65, 57.55) -- (348.65, 55.17);

\path[draw=drawColor,line width= 0.8pt,line join=round] (348.90, 57.55) -- (348.90, 52.75);

\path[draw=drawColor,line width= 0.8pt,line join=round] (349.64, 57.55) -- (349.64, 55.46);

\path[draw=drawColor,line width= 0.8pt,line join=round] (349.89, 57.55) -- (349.89, 48.96);

\path[draw=drawColor,line width= 0.8pt,line join=round] (350.14, 57.55) -- (350.14, 57.27);

\path[draw=drawColor,line width= 0.8pt,line join=round] (350.39, 57.55) -- (350.39, 51.61);

\path[draw=drawColor,line width= 0.8pt,line join=round] (350.63, 57.55) -- (350.63, 58.23);

\path[draw=drawColor,line width= 0.8pt,line join=round] (351.38, 57.55) -- (351.38, 51.88);

\path[draw=drawColor,line width= 0.8pt,line join=round] (351.63, 57.55) -- (351.63, 55.76);

\path[draw=drawColor,line width= 0.8pt,line join=round] (351.88, 57.55) -- (351.88, 55.94);

\path[draw=drawColor,line width= 0.8pt,line join=round] (352.12, 57.55) -- (352.12, 59.99);

\path[draw=drawColor,line width= 0.8pt,line join=round] (352.37, 57.55) -- (352.37, 56.93);

\path[draw=drawColor,line width= 0.8pt,line join=round] (353.12, 57.55) -- (353.12, 58.61);

\path[draw=drawColor,line width= 0.8pt,line join=round] (353.37, 57.55) -- (353.37, 53.30);

\path[draw=drawColor,line width= 0.8pt,line join=round] (353.86, 57.55) -- (353.86, 60.97);

\path[draw=drawColor,line width= 0.8pt,line join=round] (354.11, 57.55) -- (354.11, 56.86);

\path[draw=drawColor,line width= 0.8pt,line join=round] (354.86, 57.55) -- (354.86, 66.92);

\path[draw=drawColor,line width= 0.8pt,line join=round] (355.10, 57.55) -- (355.10, 58.01);

\path[draw=drawColor,line width= 0.8pt,line join=round] (355.35, 57.55) -- (355.35, 50.02);

\path[draw=drawColor,line width= 0.8pt,line join=round] (355.60, 57.55) -- (355.60, 62.03);

\path[draw=drawColor,line width= 0.8pt,line join=round] (355.85, 57.55) -- (355.85, 59.76);

\path[draw=drawColor,line width= 0.8pt,line join=round] (356.59, 57.55) -- (356.59, 53.08);

\path[draw=drawColor,line width= 0.8pt,line join=round] (356.84, 57.55) -- (356.84, 59.01);

\path[draw=drawColor,line width= 0.8pt,line join=round] (357.09, 57.55) -- (357.09, 61.16);

\path[draw=drawColor,line width= 0.8pt,line join=round] (357.34, 57.55) -- (357.34, 54.72);

\path[draw=drawColor,line width= 0.8pt,line join=round] (357.59, 57.55) -- (357.59, 53.16);

\path[draw=drawColor,line width= 0.8pt,line join=round] (358.33, 57.55) -- (358.33, 56.62);

\path[draw=drawColor,line width= 0.8pt,line join=round] (358.58, 57.55) -- (358.58, 60.42);

\path[draw=drawColor,line width= 0.8pt,line join=round] (358.83, 57.55) -- (358.83, 61.75);

\path[draw=drawColor,line width= 0.8pt,line join=round] (359.08, 57.55) -- (359.08, 61.00);

\path[draw=drawColor,line width= 0.8pt,line join=round] (359.32, 57.55) -- (359.32, 57.51);

\path[draw=drawColor,line width= 0.8pt,line join=round] (360.07, 57.55) -- (360.07, 52.91);

\path[draw=drawColor,line width= 0.8pt,line join=round] (360.32, 57.55) -- (360.32, 64.20);

\path[draw=drawColor,line width= 0.8pt,line join=round] (360.57, 57.55) -- (360.57, 48.86);

\path[draw=drawColor,line width= 0.8pt,line join=round] (360.81, 57.55) -- (360.81, 56.65);

\path[draw=drawColor,line width= 0.8pt,line join=round] (361.06, 57.55) -- (361.06, 58.39);

\path[draw=drawColor,line width= 0.8pt,line join=round] (361.81, 57.55) -- (361.81, 59.17);

\path[draw=drawColor,line width= 0.8pt,line join=round] (362.06, 57.55) -- (362.06, 62.34);

\path[draw=drawColor,line width= 0.8pt,line join=round] (362.30, 57.55) -- (362.30, 53.13);

\path[draw=drawColor,line width= 0.8pt,line join=round] (362.55, 57.55) -- (362.55, 55.55);

\path[draw=drawColor,line width= 0.8pt,line join=round] (362.80, 57.55) -- (362.80, 53.13);

\path[draw=drawColor,line width= 0.8pt,line join=round] (363.55, 57.55) -- (363.55, 56.00);

\path[draw=drawColor,line width= 0.8pt,line join=round] (363.79, 57.55) -- (363.79, 57.37);

\path[draw=drawColor,line width= 0.8pt,line join=round] (364.04, 57.55) -- (364.04, 52.48);

\path[draw=drawColor,line width= 0.8pt,line join=round] (364.29, 57.55) -- (364.29, 64.91);

\path[draw=drawColor,line width= 0.8pt,line join=round] (364.54, 57.55) -- (364.54, 62.75);

\path[draw=drawColor,line width= 0.8pt,line join=round] (365.28, 57.55) -- (365.28, 56.42);

\path[draw=drawColor,line width= 0.8pt,line join=round] (365.53, 57.55) -- (365.53, 59.75);

\path[draw=drawColor,line width= 0.8pt,line join=round] (365.78, 57.55) -- (365.78, 56.95);

\path[draw=drawColor,line width= 0.8pt,line join=round] (366.03, 57.55) -- (366.03, 52.76);

\path[draw=drawColor,line width= 0.8pt,line join=round] (366.28, 57.55) -- (366.28, 59.95);

\path[draw=drawColor,line width= 0.8pt,line join=round] (367.02, 57.55) -- (367.02, 63.91);

\path[draw=drawColor,line width= 0.8pt,line join=round] (367.27, 57.55) -- (367.27, 57.48);

\path[draw=drawColor,line width= 0.8pt,line join=round] (367.52, 57.55) -- (367.52, 57.32);

\path[draw=drawColor,line width= 0.8pt,line join=round] (367.77, 57.55) -- (367.77, 52.90);

\path[draw=drawColor,line width= 0.8pt,line join=round] (368.01, 57.55) -- (368.01, 57.19);

\path[draw=drawColor,line width= 0.8pt,line join=round] (369.01, 57.55) -- (369.01, 56.52);

\path[draw=drawColor,line width= 0.8pt,line join=round] (369.26, 57.55) -- (369.26, 61.47);

\path[draw=drawColor,line width= 0.8pt,line join=round] (369.50, 57.55) -- (369.50, 59.16);

\path[draw=drawColor,line width= 0.8pt,line join=round] (369.75, 57.55) -- (369.75, 57.13);

\path[draw=drawColor,line width= 0.8pt,line join=round] (370.50, 57.55) -- (370.50, 58.75);

\path[draw=drawColor,line width= 0.8pt,line join=round] (370.75, 57.55) -- (370.75, 57.51);

\path[draw=drawColor,line width= 0.8pt,line join=round] (370.99, 57.55) -- (370.99, 61.21);

\path[draw=drawColor,line width= 0.8pt,line join=round] (371.24, 57.55) -- (371.24, 59.57);

\path[draw=drawColor,line width= 0.8pt,line join=round] (371.49, 57.55) -- (371.49, 56.52);

\path[draw=drawColor,line width= 0.8pt,line join=round] (372.24, 57.55) -- (372.24, 59.66);

\path[draw=drawColor,line width= 0.8pt,line join=round] (372.48, 57.55) -- (372.48, 62.04);

\path[draw=drawColor,line width= 0.8pt,line join=round] (372.73, 57.55) -- (372.73, 63.24);

\path[draw=drawColor,line width= 0.8pt,line join=round] (372.98, 57.55) -- (372.98, 62.44);

\path[draw=drawColor,line width= 0.8pt,line join=round] (373.23, 57.55) -- (373.23, 58.11);

\path[draw=drawColor,line width= 0.8pt,line join=round] (373.97, 57.55) -- (373.97, 53.97);

\path[draw=drawColor,line width= 0.8pt,line join=round] (374.22, 57.55) -- (374.22, 57.68);

\path[draw=drawColor,line width= 0.8pt,line join=round] (374.47, 57.55) -- (374.47, 56.15);

\path[draw=drawColor,line width= 0.8pt,line join=round] (374.72, 57.55) -- (374.72, 56.48);

\path[draw=drawColor,line width= 0.8pt,line join=round] (374.97, 57.55) -- (374.97, 57.58);

\path[draw=drawColor,line width= 0.8pt,line join=round] (375.71, 57.55) -- (375.71, 57.22);

\path[draw=drawColor,line width= 0.8pt,line join=round] (375.96, 57.55) -- (375.96, 61.40);

\path[draw=drawColor,line width= 0.8pt,line join=round] (376.21, 57.55) -- (376.21, 62.61);

\path[draw=drawColor,line width= 0.8pt,line join=round] (376.46, 57.55) -- (376.46, 55.37);

\path[draw=drawColor,line width= 0.8pt,line join=round] (376.70, 57.55) -- (376.70, 52.38);

\path[draw=drawColor,line width= 0.8pt,line join=round] (377.45, 57.55) -- (377.45, 57.61);

\path[draw=drawColor,line width= 0.8pt,line join=round] (377.70, 57.55) -- (377.70, 51.56);

\path[draw=drawColor,line width= 0.8pt,line join=round] (377.95, 57.55) -- (377.95, 48.42);

\path[draw=drawColor,line width= 0.8pt,line join=round] (378.19, 57.55) -- (378.19, 53.04);

\path[draw=drawColor,line width= 0.8pt,line join=round] (378.44, 57.55) -- (378.44, 59.36);

\path[draw=drawColor,line width= 0.8pt,line join=round] (379.19, 57.55) -- (379.19, 57.10);

\path[draw=drawColor,line width= 0.8pt,line join=round] (379.44, 57.55) -- (379.44, 60.66);

\path[draw=drawColor,line width= 0.8pt,line join=round] (379.68, 57.55) -- (379.68, 53.67);

\path[draw=drawColor,line width= 0.8pt,line join=round] (379.93, 57.55) -- (379.93, 48.23);

\path[draw=drawColor,line width= 0.8pt,line join=round] (380.18, 57.55) -- (380.18, 49.73);

\path[draw=drawColor,line width= 0.8pt,line join=round] (380.93, 57.55) -- (380.93, 52.87);

\path[draw=drawColor,line width= 0.8pt,line join=round] (381.17, 57.55) -- (381.17, 39.23);

\path[draw=drawColor,line width= 0.8pt,line join=round] (381.42, 57.55) -- (381.42, 44.17);

\path[draw=drawColor,line width= 0.8pt,line join=round] (381.67, 57.55) -- (381.67, 64.27);

\path[draw=drawColor,line width= 0.8pt,line join=round] (381.92, 57.55) -- (381.92, 56.37);

\path[draw=drawColor,line width= 0.8pt,line join=round] (382.66, 57.55) -- (382.66, 55.37);

\path[draw=drawColor,line width= 0.8pt,line join=round] (382.91, 57.55) -- (382.91, 64.55);

\path[draw=drawColor,line width= 0.8pt,line join=round] (383.16, 57.55) -- (383.16, 65.08);

\path[draw=drawColor,line width= 0.8pt,line join=round] (383.41, 57.55) -- (383.41, 63.95);

\path[draw=drawColor,line width= 0.8pt,line join=round] (383.66, 57.55) -- (383.66, 59.57);

\path[draw=drawColor,line width= 0.8pt,line join=round] (384.40, 57.55) -- (384.40, 58.69);

\path[draw=drawColor,line width= 0.8pt,line join=round] (384.65, 57.55) -- (384.65, 62.87);

\path[draw=drawColor,line width= 0.8pt,line join=round] (384.90, 57.55) -- (384.90, 67.72);

\path[draw=drawColor,line width= 0.8pt,line join=round] (385.15, 57.55) -- (385.15, 51.06);

\path[draw=drawColor,line width= 0.8pt,line join=round] (385.39, 57.55) -- (385.39, 49.46);

\path[draw=drawColor,line width= 0.8pt,line join=round] (386.14, 57.55) -- (386.14, 54.01);

\path[draw=drawColor,line width= 0.8pt,line join=round] (386.39, 57.55) -- (386.39, 59.33);

\path[draw=drawColor,line width= 0.8pt,line join=round] (386.64, 57.55) -- (386.64, 55.95);

\path[draw=drawColor,line width= 0.8pt,line join=round] (386.88, 57.55) -- (386.88, 65.46);

\path[draw=drawColor,line width= 0.8pt,line join=round] (387.13, 57.55) -- (387.13, 58.52);

\path[draw=drawColor,line width= 0.8pt,line join=round] (387.88, 57.55) -- (387.88, 50.30);

\path[draw=drawColor,line width= 0.8pt,line join=round] (388.13, 57.55) -- (388.13, 50.58);

\path[draw=drawColor,line width= 0.8pt,line join=round] (388.37, 57.55) -- (388.37, 60.57);

\path[draw=drawColor,line width= 0.8pt,line join=round] (388.87, 57.55) -- (388.87, 54.61);

\path[draw=drawColor,line width= 0.8pt,line join=round] (389.61, 57.55) -- (389.61, 62.22);

\path[draw=drawColor,line width= 0.8pt,line join=round] (389.86, 57.55) -- (389.86, 57.17);

\path[draw=drawColor,line width= 0.8pt,line join=round] (390.11, 57.55) -- (390.11, 68.81);

\path[draw=drawColor,line width= 0.8pt,line join=round] (390.36, 57.55) -- (390.36, 56.58);

\path[draw=drawColor,line width= 0.8pt,line join=round] (390.61, 57.55) -- (390.61, 67.08);

\path[draw=drawColor,line width= 0.8pt,line join=round] (391.35, 57.55) -- (391.35, 63.11);

\path[draw=drawColor,line width= 0.8pt,line join=round] (391.60, 57.55) -- (391.60, 40.81);

\path[draw=drawColor,line width= 0.8pt,line join=round] (392.10, 57.55) -- (392.10, 55.83);

\path[draw=drawColor,line width= 0.8pt,line join=round] (392.35, 57.55) -- (392.35, 48.63);

\path[draw=drawColor,line width= 0.8pt,line join=round] (393.09, 57.55) -- (393.09, 57.24);

\path[draw=drawColor,line width= 0.8pt,line join=round] (393.34, 57.55) -- (393.34, 57.26);

\path[draw=drawColor,line width= 0.8pt,line join=round] (393.59, 57.55) -- (393.59, 61.54);

\path[draw=drawColor,line width= 0.8pt,line join=round] (393.84, 57.55) -- (393.84, 58.53);

\path[draw=drawColor,line width= 0.8pt,line join=round] (394.08, 57.55) -- (394.08, 59.15);

\path[draw=drawColor,line width= 0.8pt,line join=round] (394.83, 57.55) -- (394.83, 53.28);

\path[draw=drawColor,line width= 0.8pt,line join=round] (395.08, 57.55) -- (395.08, 57.17);

\path[draw=drawColor,line width= 0.8pt,line join=round] (395.33, 57.55) -- (395.33, 53.88);

\path[draw=drawColor,line width= 0.8pt,line join=round] (395.57, 57.55) -- (395.57, 56.06);

\path[draw=drawColor,line width= 0.8pt,line join=round] (395.82, 57.55) -- (395.82, 54.69);

\path[draw=drawColor,line width= 0.8pt,line join=round] (396.57, 57.55) -- (396.57, 51.41);

\path[draw=drawColor,line width= 0.8pt,line join=round] (397.06, 57.55) -- (397.06, 72.61);

\path[draw=drawColor,line width= 0.8pt,line join=round] (397.31, 57.55) -- (397.31, 61.09);

\path[draw=drawColor,line width= 0.8pt,line join=round] (397.56, 57.55) -- (397.56, 55.59);
\definecolor{drawColor}{RGB}{169,169,169}

\path[draw=drawColor,line width= 0.4pt,line join=round,line cap=round] ( 36.06, 45.90) -- (397.56, 45.90);

\path[draw=drawColor,line width= 0.4pt,line join=round,line cap=round] ( 36.06, 57.55) -- (397.56, 57.55);

\path[draw=drawColor,line width= 0.4pt,line join=round,line cap=round] ( 36.06, 69.20) -- (397.56, 69.20);
\end{scope}
\begin{scope}
\path[clip] (  0.00,  0.00) rectangle (433.62,216.81);
\definecolor{drawColor}{gray}{0.20}

\node[text=drawColor,anchor=base west,inner sep=0pt, outer sep=0pt, scale=  0.81] at ( 28.07, 44.04) {-5};

\node[text=drawColor,anchor=base west,inner sep=0pt, outer sep=0pt, scale=  0.81] at ( 28.07, 55.69) { 0};

\node[text=drawColor,anchor=base west,inner sep=0pt, outer sep=0pt, scale=  0.81] at ( 28.07, 67.34) { 5};

\node[text=drawColor,anchor=base west,inner sep=0pt, outer sep=0pt, scale=  0.81] at (398.06, 44.04) {-5};

\node[text=drawColor,anchor=base west,inner sep=0pt, outer sep=0pt, scale=  0.81] at (398.06, 55.69) { 0};

\node[text=drawColor,anchor=base west,inner sep=0pt, outer sep=0pt, scale=  0.81] at (398.06, 67.34) { 5};
\end{scope}
\end{tikzpicture}

	\caption{Evolució del preu de tancament i rendibilitat diàries de Caterpillar entre 2015 i 2018}
	\label{fig:serie preus}
\end{figure}
El model més senzill per a l'evolució dels actius financers d'una empresa és la hipòtesi que les rendibilitats diàries són independents normalment distribuïdes amb paràmetres independents del temps. La rendibilitat (a vegades anomenada renbilitat neta per a distingir-la de la rendibilitat bruta) diària es defineix com el logaritme del quocient de preus. És a dir, si \( P_t \) és el preu el dia \( t \) aleshores la rendibilitat d'aquest dia, \( r_t \) és
\begin{equation*}
	r_t = \log{\left(\frac{P_t}{P_{t-1}}\right)}.
\end{equation*}
La hipòtesi, doncs, és que \( r_t \sim N(\mu, \sigma^2) \). Aquest model s'anomena el model de Samuelson-Bachelier. Observem que si \( \sigma = 0 \) aleshores la rendibilitat és constant i els preus evolucionen exponencialment segons,
\begin{equation*}
	P_t = e^{r}P_{t-1}
\end{equation*}
que és la fórmula de l'interès compost. Així doncs, \( \sigma \) governa el risc en l'evolució i s'anomena la volatilitat. 

Si introduïm els log-preus, \( p_t = \log{P_t} \), tenim que \( p_t = p_{t-1} + r_t \), i per tant que l'evolució dels log-preus és un passeig aleatori. Seguint el model de Samuelson-Bachieler podem estimar els paràmetres de la distribució de les rendibilitats i amb aquests fer una previsió de l'evolució dels preus dels actius.  

A la \cref{fig:serie preus} hi ha l'històric dels preus de tancament i rendibilitats diàries de Caterpillar al llarg dels últims 4 anys (2015-2018), en total 1004 dades. La mitjana i desviació estàndard mostrals són
\begin{equation} \label{eq:parametres}
	\begin{gathered}
		\hat{\mu}_d = 0.0311 \\
		\hat{\sigma}_d	= 1.70 \\
	\end{gathered}
\end{equation}

\section{Resta}

\begin{figure}[htb]
	\centering \sffamily \small
	% Created by tikzDevice version 0.12 on 2019-06-12 00:01:05
% !TEX encoding = UTF-8 Unicode
\begin{tikzpicture}[x=1pt,y=1pt]
\definecolor{fillColor}{RGB}{255,255,255}
\path[use as bounding box,fill=fillColor,fill opacity=0.00] (0,0) rectangle (361.35,289.08);
\begin{scope}
\path[clip] ( 21.60, 32.40) rectangle (339.75,289.08);
\definecolor{drawColor}{RGB}{169,169,169}

\path[draw=drawColor,line width= 0.4pt,line join=round,line cap=round] ( 33.38, 41.91) -- ( 33.38,245.62);

\path[draw=drawColor,line width= 0.4pt,line join=round,line cap=round] ( 69.80, 41.91) -- ( 69.80,245.62);

\path[draw=drawColor,line width= 0.4pt,line join=round,line cap=round] (107.64, 41.91) -- (107.64,245.62);

\path[draw=drawColor,line width= 0.4pt,line join=round,line cap=round] (143.85, 41.91) -- (143.85,245.62);

\path[draw=drawColor,line width= 0.4pt,line join=round,line cap=round] (181.48, 41.91) -- (181.48,245.62);

\path[draw=drawColor,line width= 0.4pt,line join=round,line cap=round] (218.10, 41.91) -- (218.10,245.62);

\path[draw=drawColor,line width= 0.4pt,line join=round,line cap=round] (255.13, 41.91) -- (255.13,245.62);

\path[draw=drawColor,line width= 0.4pt,line join=round,line cap=round] (291.75, 41.91) -- (291.75,245.62);

\path[draw=drawColor,line width= 0.4pt,line join=round,line cap=round] (327.97, 41.91) -- (327.97,245.62);
\end{scope}
\begin{scope}
\path[clip] (  0.00,  0.00) rectangle (361.35,289.08);
\definecolor{drawColor}{gray}{0.20}

\path[draw=drawColor,line width= 0.4pt,line join=round,line cap=round] ( 33.38, 32.40) -- (327.97, 32.40);

\path[draw=drawColor,line width= 0.6pt,line join=round,line cap=round] ( 33.38, 32.40) -- ( 33.38, 28.08);

\path[draw=drawColor,line width= 0.6pt,line join=round,line cap=round] ( 69.80, 32.40) -- ( 69.80, 28.08);

\path[draw=drawColor,line width= 0.6pt,line join=round,line cap=round] (107.64, 32.40) -- (107.64, 28.08);

\path[draw=drawColor,line width= 0.6pt,line join=round,line cap=round] (143.85, 32.40) -- (143.85, 28.08);

\path[draw=drawColor,line width= 0.6pt,line join=round,line cap=round] (181.48, 32.40) -- (181.48, 28.08);

\path[draw=drawColor,line width= 0.6pt,line join=round,line cap=round] (218.10, 32.40) -- (218.10, 28.08);

\path[draw=drawColor,line width= 0.6pt,line join=round,line cap=round] (255.13, 32.40) -- (255.13, 28.08);

\path[draw=drawColor,line width= 0.6pt,line join=round,line cap=round] (291.75, 32.40) -- (291.75, 28.08);

\path[draw=drawColor,line width= 0.6pt,line join=round,line cap=round] (327.97, 32.40) -- (327.97, 28.08);

\node[text=drawColor,anchor=base,inner sep=0pt, outer sep=0pt, scale=  0.81] at ( 33.38, 17.28) {Jan 02};

\node[text=drawColor,anchor=base,inner sep=0pt, outer sep=0pt, scale=  0.81] at ( 33.38,  7.56) {2015};

\node[text=drawColor,anchor=base,inner sep=0pt, outer sep=0pt, scale=  0.81] at ( 69.80, 17.28) {Jul 01};

\node[text=drawColor,anchor=base,inner sep=0pt, outer sep=0pt, scale=  0.81] at ( 69.80,  7.56) {2015};

\node[text=drawColor,anchor=base,inner sep=0pt, outer sep=0pt, scale=  0.81] at (107.64, 17.28) {Jan 04};

\node[text=drawColor,anchor=base,inner sep=0pt, outer sep=0pt, scale=  0.81] at (107.64,  7.56) {2016};

\node[text=drawColor,anchor=base,inner sep=0pt, outer sep=0pt, scale=  0.81] at (143.85, 17.28) {Jul 01};

\node[text=drawColor,anchor=base,inner sep=0pt, outer sep=0pt, scale=  0.81] at (143.85,  7.56) {2016};

\node[text=drawColor,anchor=base,inner sep=0pt, outer sep=0pt, scale=  0.81] at (181.48, 17.28) {Jan 03};

\node[text=drawColor,anchor=base,inner sep=0pt, outer sep=0pt, scale=  0.81] at (181.48,  7.56) {2017};

\node[text=drawColor,anchor=base,inner sep=0pt, outer sep=0pt, scale=  0.81] at (218.10, 17.28) {Jul 03};

\node[text=drawColor,anchor=base,inner sep=0pt, outer sep=0pt, scale=  0.81] at (218.10,  7.56) {2017};

\node[text=drawColor,anchor=base,inner sep=0pt, outer sep=0pt, scale=  0.81] at (255.13, 17.28) {Jan 02};

\node[text=drawColor,anchor=base,inner sep=0pt, outer sep=0pt, scale=  0.81] at (255.13,  7.56) {2018};

\node[text=drawColor,anchor=base,inner sep=0pt, outer sep=0pt, scale=  0.81] at (291.75, 17.28) {Jul 02};

\node[text=drawColor,anchor=base,inner sep=0pt, outer sep=0pt, scale=  0.81] at (291.75,  7.56) {2018};

\node[text=drawColor,anchor=base,inner sep=0pt, outer sep=0pt, scale=  0.81] at (327.97, 17.28) {Dec 28};

\node[text=drawColor,anchor=base,inner sep=0pt, outer sep=0pt, scale=  0.81] at (327.97,  7.56) {2018};
\end{scope}
\begin{scope}
\path[clip] ( 21.60,245.62) rectangle (339.75,279.57);
\definecolor{drawColor}{gray}{0.20}

\node[text=drawColor,anchor=base west,inner sep=0pt, outer sep=0pt, scale=  0.99] at ( 33.38,260.32) {\bfseries Evolució d'una inversió de \$100};

\node[text=drawColor,anchor=base east,inner sep=0pt, outer sep=0pt, scale=  0.90] at (322.57,260.53) {2015-01-02 / 2018-12-28};
\end{scope}
\begin{scope}
\path[clip] ( 21.60, 41.91) rectangle (339.75,245.62);
\definecolor{drawColor}{RGB}{169,169,169}

\path[draw=drawColor,line width= 0.4pt,line join=round,line cap=round] ( 33.38, 70.02) -- (327.97, 70.02);

\path[draw=drawColor,line width= 0.4pt,line join=round,line cap=round] ( 33.38,103.16) -- (327.97,103.16);

\path[draw=drawColor,line width= 0.4pt,line join=round,line cap=round] ( 33.38,136.29) -- (327.97,136.29);

\path[draw=drawColor,line width= 0.4pt,line join=round,line cap=round] ( 33.38,169.43) -- (327.97,169.43);

\path[draw=drawColor,line width= 0.4pt,line join=round,line cap=round] ( 33.38,202.56) -- (327.97,202.56);

\path[draw=drawColor,line width= 0.4pt,line join=round,line cap=round] ( 33.38,235.69) -- (327.97,235.69);
\end{scope}
\begin{scope}
\path[clip] (  0.00,  0.00) rectangle (361.35,289.08);
\definecolor{drawColor}{gray}{0.20}

\node[text=drawColor,anchor=base east,inner sep=0pt, outer sep=0pt, scale=  0.81] at ( 22.58, 68.16) { 80};

\node[text=drawColor,anchor=base east,inner sep=0pt, outer sep=0pt, scale=  0.81] at ( 22.58,101.30) {100};

\node[text=drawColor,anchor=base east,inner sep=0pt, outer sep=0pt, scale=  0.81] at ( 22.58,134.43) {120};

\node[text=drawColor,anchor=base east,inner sep=0pt, outer sep=0pt, scale=  0.81] at ( 22.58,167.57) {140};

\node[text=drawColor,anchor=base east,inner sep=0pt, outer sep=0pt, scale=  0.81] at ( 22.58,200.70) {160};

\node[text=drawColor,anchor=base east,inner sep=0pt, outer sep=0pt, scale=  0.81] at ( 22.58,233.83) {180};

\node[text=drawColor,anchor=base west,inner sep=0pt, outer sep=0pt, scale=  0.81] at (338.77, 68.16) { 80};

\node[text=drawColor,anchor=base west,inner sep=0pt, outer sep=0pt, scale=  0.81] at (338.77,101.30) {100};

\node[text=drawColor,anchor=base west,inner sep=0pt, outer sep=0pt, scale=  0.81] at (338.77,134.43) {120};

\node[text=drawColor,anchor=base west,inner sep=0pt, outer sep=0pt, scale=  0.81] at (338.77,167.57) {140};

\node[text=drawColor,anchor=base west,inner sep=0pt, outer sep=0pt, scale=  0.81] at (338.77,200.70) {160};

\node[text=drawColor,anchor=base west,inner sep=0pt, outer sep=0pt, scale=  0.81] at (338.77,233.83) {180};
\end{scope}
\begin{scope}
\path[clip] ( 21.60, 41.91) rectangle (339.75,245.62);
\definecolor{drawColor}{RGB}{0,0,255}

\path[draw=drawColor,line width= 0.8pt,line join=round] ( 33.38,103.16) --
	( 33.99, 94.41) --
	( 34.19, 93.40) --
	( 34.39, 95.82) --
	( 34.60, 97.44) --
	( 34.80, 95.53) --
	( 35.41, 94.16) --
	( 35.61, 93.22) --
	( 35.81, 92.68) --
	( 36.01, 89.53) --
	( 36.22, 88.70) --
	( 37.03, 88.90) --
	( 37.23, 91.42) --
	( 37.43, 94.03) --
	( 37.63, 91.85) --
	( 38.24, 92.61) --
	( 38.44, 81.47) --
	( 38.64, 81.76) --
	( 38.85, 81.74) --
	( 39.05, 81.68) --
	( 39.66, 83.25) --
	( 39.86, 88.81) --
	( 40.06, 85.25) --
	( 40.26, 88.17) --
	( 40.46, 87.53) --
	( 41.07, 90.16) --
	( 41.27, 88.77) --
	( 41.48, 87.89) --
	( 41.68, 88.10) --
	( 41.88, 90.99) --
	( 42.69, 91.44) --
	( 42.89, 90.39) --
	( 43.09, 87.69) --
	( 43.30, 89.74) --
	( 43.90, 88.03) --
	( 44.11, 89.85) --
	( 44.31, 91.02) --
	( 44.51, 88.59) --
	( 44.71, 86.97) --
	( 45.32, 87.15) --
	( 45.52, 86.86) --
	( 45.73, 84.86) --
	( 45.93, 84.46) --
	( 46.13, 81.85) --
	( 46.74, 83.20) --
	( 46.94, 81.57) --
	( 47.14, 81.20) --
	( 47.34, 82.08) --
	( 47.55, 80.35) --
	( 48.15, 81.68) --
	( 48.36, 78.94) --
	( 48.56, 84.14) --
	( 48.76, 81.90) --
	( 48.96, 83.88) --
	( 49.57, 83.70) --
	( 49.77, 82.60) --
	( 49.97, 81.48) --
	( 50.18, 81.52) --
	( 50.38, 81.14) --
	( 50.99, 84.21) --
	( 51.19, 81.79) --
	( 51.39, 81.09) --
	( 51.59, 82.17) --
	( 52.40, 84.42) --
	( 52.60, 82.71) --
	( 52.81, 82.53) --
	( 53.01, 82.96) --
	( 53.21, 86.43) --
	( 53.82, 85.60) --
	( 54.02, 87.15) --
	( 54.22, 91.04) --
	( 54.42, 89.56) --
	( 54.63, 87.65) --
	( 55.23, 90.03) --
	( 55.44, 88.81) --
	( 55.64, 90.52) --
	( 55.84, 90.37) --
	( 56.04, 90.03) --
	( 56.65, 91.35) --
	( 56.85, 92.92) --
	( 57.06, 95.26) --
	( 57.26, 94.14) --
	( 57.46, 95.03) --
	( 58.07, 94.90) --
	( 58.27, 94.36) --
	( 58.47, 94.34) --
	( 58.67, 93.33) --
	( 58.88, 94.92) --
	( 59.48, 97.55) --
	( 59.69, 96.45) --
	( 59.89, 96.96) --
	( 60.09, 97.50) --
	( 60.29, 96.94) --
	( 60.90, 96.68) --
	( 61.10, 94.79) --
	( 61.30, 96.04) --
	( 61.51, 98.56) --
	( 61.71, 97.30) --
	( 62.52, 95.87) --
	( 62.72, 96.02) --
	( 62.92, 92.57) --
	( 63.12, 91.33) --
	( 63.73, 91.71) --
	( 63.93, 92.86) --
	( 64.14, 92.97) --
	( 64.34, 91.60) --
	( 64.54, 92.65) --
	( 65.15, 92.83) --
	( 65.35, 93.87) --
	( 65.55, 97.03) --
	( 65.76, 96.47) --
	( 65.96, 95.89) --
	( 66.56, 94.59) --
	( 66.77, 93.62) --
	( 66.97, 94.81) --
	( 67.17, 95.15) --
	( 67.37, 95.30) --
	( 67.98, 96.65) --
	( 68.18, 96.97) --
	( 68.39, 96.47) --
	( 68.59, 93.85) --
	( 68.79, 94.03) --
	( 69.40, 91.29) --
	( 69.60, 90.43) --
	( 69.80, 88.86) --
	( 70.00, 89.63) --
	( 70.81, 87.42) --
	( 71.02, 88.93) --
	( 71.22, 85.83) --
	( 71.42, 84.78) --
	( 71.62, 85.61) --
	( 72.23, 88.30) --
	( 72.43, 89.78) --
	( 72.63, 89.24) --
	( 72.84, 88.52) --
	( 73.04, 87.43) --
	( 73.65, 85.79) --
	( 73.85, 85.74) --
	( 74.05, 81.30) --
	( 74.25, 76.11) --
	( 74.46, 74.71) --
	( 75.06, 73.30) --
	( 75.26, 77.73) --
	( 75.47, 76.92) --
	( 75.67, 78.76) --
	( 75.87, 79.27) --
	( 76.48, 76.80) --
	( 76.68, 75.61) --
	( 76.88, 76.11) --
	( 77.09, 77.16) --
	( 77.29, 76.85) --
	( 77.89, 82.03) --
	( 78.10, 78.20) --
	( 78.30, 79.79) --
	( 78.50, 77.70) --
	( 78.70, 79.01) --
	( 79.31, 79.10) --
	( 79.51, 79.61) --
	( 79.72, 77.07) --
	( 79.92, 75.26) --
	( 80.12, 72.72) --
	( 80.73, 68.79) --
	( 80.93, 67.42) --
	( 81.13, 70.68) --
	( 81.33, 73.91) --
	( 81.54, 74.43) --
	( 82.14, 75.32) --
	( 82.35, 72.54) --
	( 82.55, 74.71) --
	( 82.75, 71.73) --
	( 82.95, 69.30) --
	( 83.76, 71.46) --
	( 83.96, 69.04) --
	( 84.17, 68.07) --
	( 84.37, 68.45) --
	( 84.98, 68.70) --
	( 85.18, 71.96) --
	( 85.38, 73.88) --
	( 85.58, 71.03) --
	( 85.79, 67.06) --
	( 86.39, 67.60) --
	( 86.59, 66.74) --
	( 86.80, 64.07) --
	( 87.00, 56.13) --
	( 87.20, 54.65) --
	( 87.81, 52.51) --
	( 88.01, 53.45) --
	( 88.21, 55.34) --
	( 88.42, 53.59) --
	( 88.62, 55.95) --
	( 89.22, 62.23) --
	( 89.43, 65.18) --
	( 89.63, 64.39) --
	( 89.83, 67.02) --
	( 90.03, 66.05) --
	( 90.64, 64.61) --
	( 90.84, 64.01) --
	( 91.05, 64.99) --
	( 91.25, 65.20) --
	( 91.45, 63.13) --
	( 92.06, 62.39) --
	( 92.26, 64.19) --
	( 92.46, 63.11) --
	( 92.66, 65.29) --
	( 92.87, 66.86) --
	( 93.47, 66.45) --
	( 93.68, 64.41) --
	( 93.88, 67.26) --
	( 94.08, 67.26) --
	( 94.28, 69.10) --
	( 94.89, 71.53) --
	( 95.09, 72.27) --
	( 95.29, 71.91) --
	( 95.50, 71.32) --
	( 95.70, 70.63) --
	( 96.31, 67.11) --
	( 96.51, 68.09) --
	( 96.71, 67.15) --
	( 96.91, 61.29) --
	( 97.12, 63.04) --
	( 97.72, 64.41) --
	( 97.92, 62.61) --
	( 98.13, 64.28) --
	( 98.33, 63.74) --
	( 98.53, 65.76) --
	( 99.14, 65.55) --
	( 99.34, 66.23) --
	( 99.54, 66.39) --
	( 99.95, 65.91) --
	(100.55, 68.48) --
	(100.76, 66.52) --
	(100.96, 65.56) --
	(101.16, 63.06) --
	(101.36, 63.67) --
	(101.97, 60.75) --
	(102.17, 57.47) --
	(102.38, 56.33) --
	(102.58, 57.18) --
	(102.78, 55.47) --
	(103.39, 56.71) --
	(103.59, 57.85) --
	(103.79, 59.60) --
	(103.99, 54.51) --
	(104.20, 54.89) --
	(104.80, 55.12) --
	(105.01, 60.84) --
	(105.21, 63.51) --
	(105.41, 62.59) --
	(106.22, 61.15) --
	(106.42, 62.25) --
	(106.62, 61.34) --
	(106.83, 60.03) --
	(107.64, 60.08) --
	(107.84, 58.80) --
	(108.04, 56.89) --
	(108.24, 52.78) --
	(108.45, 51.61) --
	(109.05, 48.33) --
	(109.25, 48.56) --
	(109.46, 47.28) --
	(109.66, 49.77) --
	(109.86, 45.44) --
	(110.67, 43.93) --
	(110.87, 43.53) --
	(111.08, 45.12) --
	(111.28, 47.44) --
	(111.88, 41.91) --
	(112.09, 44.16) --
	(112.29, 42.65) --
	(112.49, 47.62) --
	(112.69, 49.71) --
	(113.30, 50.54) --
	(113.50, 46.88) --
	(113.71, 51.57) --
	(113.91, 56.42) --
	(114.11, 56.71) --
	(114.72, 53.93) --
	(114.92, 52.76) --
	(115.12, 49.53) --
	(115.32, 48.22) --
	(115.53, 51.35) --
	(116.34, 55.07) --
	(116.54, 58.77) --
	(116.74, 56.71) --
	(116.94, 55.45) --
	(117.55, 58.86) --
	(117.75, 56.10) --
	(117.95, 56.28) --
	(118.16, 57.34) --
	(118.36, 58.06) --
	(118.97, 59.56) --
	(119.17, 62.07) --
	(119.37, 62.59) --
	(119.57, 66.86) --
	(119.78, 68.83) --
	(120.38, 72.31) --
	(120.58, 66.83) --
	(120.79, 67.08) --
	(120.99, 66.16) --
	(121.19, 68.75) --
	(121.80, 68.63) --
	(122.00, 68.11) --
	(122.20, 71.53) --
	(122.41, 74.34) --
	(122.61, 73.57) --
	(123.22, 74.33) --
	(123.42, 73.71) --
	(123.62, 70.34) --
	(123.82, 73.24) --
	(124.63, 73.30) --
	(124.83, 74.54) --
	(125.04, 75.44) --
	(125.24, 75.50) --
	(125.44, 75.95) --
	(126.05, 74.02) --
	(126.25, 73.15) --
	(126.45, 73.12) --
	(126.65, 71.23) --
	(126.86, 71.55) --
	(127.46, 72.05) --
	(127.67, 74.71) --
	(127.87, 80.17) --
	(128.07, 80.02) --
	(128.27, 80.24) --
	(128.88, 80.38) --
	(129.08, 82.44) --
	(129.28, 81.27) --
	(129.49, 79.32) --
	(129.69, 78.71) --
	(130.30, 75.95) --
	(130.50, 77.50) --
	(130.70, 79.36) --
	(130.90, 77.68) --
	(131.11, 77.63) --
	(131.71, 77.84) --
	(131.91, 75.17) --
	(132.12, 71.35) --
	(132.32, 68.74) --
	(132.52, 69.76) --
	(133.13, 65.11) --
	(133.33, 68.23) --
	(133.53, 68.09) --
	(133.74, 66.77) --
	(133.94, 63.83) --
	(134.55, 64.97) --
	(134.75, 65.60) --
	(134.95, 64.77) --
	(135.15, 62.68) --
	(135.35, 63.47) --
	(135.96, 64.43) --
	(136.16, 65.67) --
	(136.37, 68.34) --
	(136.57, 67.46) --
	(136.77, 67.24) --
	(137.58, 68.23) --
	(137.78, 67.80) --
	(137.98, 70.23) --
	(138.19, 72.79) --
	(138.79, 75.28) --
	(139.00, 75.99) --
	(139.20, 78.33) --
	(139.40, 76.62) --
	(139.60, 74.58) --
	(140.21, 73.14) --
	(140.41, 72.47) --
	(140.61, 72.85) --
	(140.82, 72.70) --
	(141.02, 74.40) --
	(141.63, 75.30) --
	(141.83, 75.41) --
	(142.03, 75.30) --
	(142.23, 78.53) --
	(142.44, 69.17) --
	(143.04, 66.19) --
	(143.25, 68.25) --
	(143.45, 71.33) --
	(143.65, 74.18) --
	(143.85, 75.34) --
	(144.66, 71.60) --
	(144.86, 72.42) --
	(145.07, 72.81) --
	(145.27, 77.00) --
	(145.88, 77.77) --
	(146.08, 81.38) --
	(146.28, 81.18) --
	(146.48, 81.85) --
	(146.68, 83.00) --
	(147.29, 81.41) --
	(147.49, 81.27) --
	(147.70, 81.18) --
	(147.90, 81.76) --
	(148.10, 80.62) --
	(148.71, 79.38) --
	(148.91, 86.70) --
	(149.11, 89.13) --
	(149.31, 87.00) --
	(149.52, 86.71) --
	(150.12, 85.25) --
	(150.33, 83.97) --
	(150.53, 85.40) --
	(150.73, 84.93) --
	(150.93, 86.35) --
	(151.54, 87.94) --
	(151.74, 86.84) --
	(151.95, 86.52) --
	(152.15, 87.58) --
	(152.35, 87.15) --
	(152.96, 89.22) --
	(153.16, 89.47) --
	(153.36, 89.69) --
	(153.56, 87.83) --
	(153.77, 88.66) --
	(154.37, 88.21) --
	(154.58, 88.97) --
	(154.78, 87.42) --
	(154.98, 86.86) --
	(155.18, 86.53) --
	(155.79, 87.33) --
	(155.99, 86.21) --
	(156.19, 85.25) --
	(156.40, 84.35) --
	(156.60, 84.78) --
	(157.41, 85.49) --
	(157.61, 88.44) --
	(157.81, 88.12) --
	(158.01, 83.16) --
	(158.62, 85.20) --
	(158.82, 82.15) --
	(159.03, 83.58) --
	(159.23, 85.40) --
	(159.43, 85.43) --
	(160.04, 85.13) --
	(160.24, 84.80) --
	(160.44, 87.99) --
	(160.64, 88.05) --
	(160.85, 86.14) --
	(161.45, 86.01) --
	(161.66, 86.93) --
	(161.86, 93.62) --
	(162.06, 95.21) --
	(162.26, 97.55) --
	(162.87, 96.67) --
	(163.07, 95.28) --
	(163.28, 98.72) --
	(163.48, 98.49) --
	(163.68, 97.01) --
	(164.29, 96.56) --
	(164.49, 95.21) --
	(164.69, 95.40) --
	(164.89, 94.30) --
	(165.10, 95.57) --
	(165.70, 94.88) --
	(165.91, 94.76) --
	(166.11, 96.16) --
	(166.31, 93.69) --
	(166.51, 93.15) --
	(167.12, 92.54) --
	(167.32, 89.82) --
	(167.52, 89.18) --
	(167.73, 87.16) --
	(167.93, 88.73) --
	(168.54, 87.98) --
	(168.74, 85.78) --
	(168.94, 83.74) --
	(169.14, 84.03) --
	(169.34, 85.90) --
	(169.95, 88.50) --
	(170.15, 90.18) --
	(170.36,101.93) --
	(170.56,105.99) --
	(170.76,105.20) --
	(171.37,107.29) --
	(171.57,107.77) --
	(171.77,105.81) --
	(171.98,104.78) --
	(172.18,103.99) --
	(172.78,105.00) --
	(172.99,106.30) --
	(173.19,110.91) --
	(173.59,110.24) --
	(174.20,108.60) --
	(174.40,107.05) --
	(174.61,109.79) --
	(174.81,111.02) --
	(175.01,109.04) --
	(175.62,107.79) --
	(175.82,109.18) --
	(176.02,112.98) --
	(176.22,111.07) --
	(176.43,109.74) --
	(177.03,108.93) --
	(177.24,110.62) --
	(177.44,106.51) --
	(177.64,107.94) --
	(177.84,104.42) --
	(178.45,104.65) --
	(178.65,107.59) --
	(178.85,106.66) --
	(179.06,107.20) --
	(179.26,107.56) --
	(180.07,107.97) --
	(180.27,105.45) --
	(180.47,105.70) --
	(180.67,104.71) --
	(181.48,106.96) --
	(181.69,106.21) --
	(181.89,105.18) --
	(182.09,105.25) --
	(182.70,104.04) --
	(182.90,106.67) --
	(183.10,108.15) --
	(183.31,106.96) --
	(183.51,107.85) --
	(184.32,106.21) --
	(184.52,105.77) --
	(184.72,105.86) --
	(184.92,108.03) --
	(185.53,107.81) --
	(185.73,111.02) --
	(185.94,114.46) --
	(186.14,112.79) --
	(186.34,115.98) --
	(186.95,112.01) --
	(187.15,109.97) --
	(187.35,108.98) --
	(187.55,106.57) --
	(187.76,105.68) --
	(188.36,104.94) --
	(188.57,105.74) --
	(188.77,105.02) --
	(188.97,106.91) --
	(189.17,111.15) --
	(189.78,115.09) --
	(189.98,114.36) --
	(190.18,116.03) --
	(190.39,114.66) --
	(190.59,115.73) --
	(191.40,114.37) --
	(191.60,114.55) --
	(191.80,109.78) --
	(192.01,109.65) --
	(192.61,113.18) --
	(192.81,111.78) --
	(193.02,115.24) --
	(193.22,107.63) --
	(193.42,109.00) --
	(194.03,109.97) --
	(194.23,110.46) --
	(194.43,105.59) --
	(194.64,102.27) --
	(194.84,103.93) --
	(195.44,104.53) --
	(195.65,103.12) --
	(195.85,105.83) --
	(196.05,104.91) --
	(196.25,105.02) --
	(196.86,109.50) --
	(197.06,104.15) --
	(197.27,104.15) --
	(197.47,103.79) --
	(197.67,103.64) --
	(198.28,102.49) --
	(198.48,105.09) --
	(198.68,104.93) --
	(198.88,105.57) --
	(199.09,104.74) --
	(199.69,103.86) --
	(199.90,107.21) --
	(200.10,107.41) --
	(200.30,110.26) --
	(200.50,109.72) --
	(201.11,112.64) --
	(201.31,112.57) --
	(201.51,108.53) --
	(201.72,105.36) --
	(202.53,107.23) --
	(202.73,107.68) --
	(202.93,106.02) --
	(203.13,108.17) --
	(203.34,107.56) --
	(203.94,112.05) --
	(204.14,125.77) --
	(204.35,126.20) --
	(204.55,122.63) --
	(204.75,121.87) --
	(205.36,121.41) --
	(205.56,120.52) --
	(205.76,120.58) --
	(205.97,116.70) --
	(206.17,117.15) --
	(206.77,116.66) --
	(206.98,116.52) --
	(207.18,118.45) --
	(207.38,119.58) --
	(207.58,119.08) --
	(208.19,122.16) --
	(208.39,122.02) --
	(208.60,118.05) --
	(208.80,118.20) --
	(209.00,122.18) --
	(209.61,121.93) --
	(209.81,124.34) --
	(210.01,124.92) --
	(210.21,126.85) --
	(210.42,128.00) --
	(211.23,127.66) --
	(211.43,127.59) --
	(211.63,127.99) --
	(211.83,128.53) --
	(212.44,127.18) --
	(212.64,126.00) --
	(212.84,124.13) --
	(213.05,126.83) --
	(213.25,128.47) --
	(213.86,128.74) --
	(214.06,128.17) --
	(214.26,126.29) --
	(214.46,129.34) --
	(214.67,131.50) --
	(215.27,131.30) --
	(215.47,130.49) --
	(215.68,124.06) --
	(215.88,124.72) --
	(216.08,125.21) --
	(216.69,125.46) --
	(216.89,124.89) --
	(217.09,129.43) --
	(217.30,127.97) --
	(217.50,131.25) --
	(218.10,130.37) --
	(218.51,131.30) --
	(218.71,129.54) --
	(218.91,130.28) --
	(219.52,131.94) --
	(219.72,132.66) --
	(219.93,133.95) --
	(220.13,133.07) --
	(220.33,133.61) --
	(220.94,133.74) --
	(221.14,130.55) --
	(221.34,131.95) --
	(221.54,131.00) --
	(221.75,129.68) --
	(222.35,132.55) --
	(222.56,144.02) --
	(222.76,142.18) --
	(222.96,143.71) --
	(223.16,143.22) --
	(223.77,142.95) --
	(223.97,141.42) --
	(224.17,141.40) --
	(224.38,142.25) --
	(224.58,143.67) --
	(225.19,144.32) --
	(225.39,143.78) --
	(225.59,143.66) --
	(225.79,141.51) --
	(226.00,140.97) --
	(226.60,142.52) --
	(226.80,142.41) --
	(227.01,143.26) --
	(227.21,141.49) --
	(227.41,142.90) --
	(228.02,142.32) --
	(228.22,146.32) --
	(228.42,144.39) --
	(228.63,144.45) --
	(228.83,145.48) --
	(229.44,144.97) --
	(229.64,146.67) --
	(229.84,149.44) --
	(230.04,149.34) --
	(230.24,150.76) --
	(231.05,150.80) --
	(231.26,148.42) --
	(231.46,149.84) --
	(231.66,149.93) --
	(232.27,151.82) --
	(232.47,155.56) --
	(232.67,153.19) --
	(232.87,154.73) --
	(233.08,156.33) --
	(233.68,160.77) --
	(233.89,162.41) --
	(234.09,162.61) --
	(234.29,162.50) --
	(234.49,161.85) --
	(235.10,161.65) --
	(235.30,161.99) --
	(235.50,162.01) --
	(235.71,163.29) --
	(235.91,162.35) --
	(236.52,162.37) --
	(236.72,163.80) --
	(236.92,166.01) --
	(237.12,165.40) --
	(237.33,166.36) --
	(237.93,166.27) --
	(238.13,168.63) --
	(238.34,169.37) --
	(238.54,171.87) --
	(238.74,173.17) --
	(239.35,174.54) --
	(239.55,172.87) --
	(239.75,174.22) --
	(239.96,174.69) --
	(240.16,174.34) --
	(240.77,174.92) --
	(240.97,186.75) --
	(241.17,184.23) --
	(241.37,184.41) --
	(241.57,185.97) --
	(242.18,183.59) --
	(242.38,182.35) --
	(242.59,183.23) --
	(242.79,183.56) --
	(242.99,183.85) --
	(243.60,185.79) --
	(243.80,187.78) --
	(244.00,185.04) --
	(244.20,181.52) --
	(244.41,183.58) --
	(245.01,183.67) --
	(245.22,185.49) --
	(245.42,179.28) --
	(245.62,183.36) --
	(245.82,182.95) --
	(246.43,184.35) --
	(246.63,185.60) --
	(246.83,186.34) --
	(247.24,185.22) --
	(247.85,184.86) --
	(248.05,188.10) --
	(248.25,186.43) --
	(248.45,192.00) --
	(248.66,192.66) --
	(249.26,192.63) --
	(249.47,190.18) --
	(249.67,190.54) --
	(249.87,195.13) --
	(250.07,196.88) --
	(250.68,196.72) --
	(250.88,196.09) --
	(251.08,205.38) --
	(251.29,201.25) --
	(251.49,201.99) --
	(252.10,207.54) --
	(252.30,209.60) --
	(252.50,212.50) --
	(252.70,216.32) --
	(252.90,218.32) --
	(253.71,219.57) --
	(253.92,221.51) --
	(254.12,223.14) --
	(254.32,221.62) --
	(255.13,220.65) --
	(255.33,221.08) --
	(255.53,224.98) --
	(255.74,229.52) --
	(256.34,236.86) --
	(256.55,237.58) --
	(256.75,236.57) --
	(256.95,242.57) --
	(257.15,244.56) --
	(257.96,242.77) --
	(258.16,241.31) --
	(258.37,241.91) --
	(258.57,244.76) --
	(259.18,245.62) --
	(259.38,242.99) --
	(259.58,241.02) --
	(259.78,242.88) --
	(259.99,238.72) --
	(260.59,230.64) --
	(260.80,232.76) --
	(261.00,231.00) --
	(261.20,230.02) --
	(261.40,221.46) --
	(262.01,209.90) --
	(262.21,219.51) --
	(262.41,215.78) --
	(262.62,200.72) --
	(262.82,206.53) --
	(263.43,212.08) --
	(263.63,215.11) --
	(263.83,222.49) --
	(264.03,225.95) --
	(264.23,219.30) --
	(265.04,217.98) --
	(265.25,217.38) --
	(265.45,223.93) --
	(265.65,230.33) --
	(266.26,232.64) --
	(266.46,228.26) --
	(266.66,216.30) --
	(266.86,208.37) --
	(267.07,201.43) --
	(267.67,209.97) --
	(267.88,214.72) --
	(268.08,210.68) --
	(268.28,214.43) --
	(268.48,222.83) --
	(269.09,216.07) --
	(269.29,214.61) --
	(269.50,212.53) --
	(269.70,216.19) --
	(269.90,219.60) --
	(270.51,211.83) --
	(270.71,215.27) --
	(270.91,218.41) --
	(271.11,202.36) --
	(271.32,197.66) --
	(271.92,206.49) --
	(272.13,202.53) --
	(272.33,199.23) --
	(272.53,203.23) --
	(273.34,196.94) --
	(273.54,199.05) --
	(273.74,199.26) --
	(273.95,204.58) --
	(274.15,195.31) --
	(274.76,195.48) --
	(274.96,204.51) --
	(275.16,202.17) --
	(275.36,208.21) --
	(275.56,208.37) --
	(276.17,211.81) --
	(276.37,213.92) --
	(276.58,218.68) --
	(276.78,218.23) --
	(276.98,213.81) --
	(277.59,215.15) --
	(277.79,197.93) --
	(277.99,199.28) --
	(278.20,201.30) --
	(278.40,198.43) --
	(279.00,197.78) --
	(279.21,197.89) --
	(279.41,196.88) --
	(279.61,196.74) --
	(279.81,201.41) --
	(280.42,207.63) --
	(280.62,207.83) --
	(280.83,212.66) --
	(281.03,215.11) --
	(281.23,216.72) --
	(281.84,217.47) --
	(282.04,212.62) --
	(282.24,214.81) --
	(282.44,214.61) --
	(282.65,218.25) --
	(283.25,224.04) --
	(283.46,219.13) --
	(283.66,221.91) --
	(283.86,221.84) --
	(284.06,218.50) --
	(284.87,214.48) --
	(285.07,217.80) --
	(285.28,211.40) --
	(285.48,214.30) --
	(286.09,212.46) --
	(286.29,213.38) --
	(286.49,219.69) --
	(286.69,219.30) --
	(286.89,217.55) --
	(287.50,219.31) --
	(287.70,221.46) --
	(287.91,216.45) --
	(288.11,213.62) --
	(288.31,207.99) --
	(288.92,205.57) --
	(289.12,195.87) --
	(289.32,195.57) --
	(289.53,189.08) --
	(289.73,189.81) --
	(290.33,183.76) --
	(290.54,181.88) --
	(290.74,180.20) --
	(290.94,182.76) --
	(291.14,182.12) --
	(291.75,183.23) --
	(291.95,178.80) --
	(292.36,182.37) --
	(292.56,181.65) --
	(293.17,191.67) --
	(293.37,192.18) --
	(293.57,184.08) --
	(293.77,188.88) --
	(293.98,191.28) --
	(294.58,186.46) --
	(294.79,188.03) --
	(294.99,191.94) --
	(295.19,188.88) --
	(295.39,184.19) --
	(296.00,183.22) --
	(296.20,186.26) --
	(296.40,190.79) --
	(296.61,194.58) --
	(296.81,194.54) --
	(297.42,189.47) --
	(297.62,196.77) --
	(297.82,187.29) --
	(298.02,186.34) --
	(298.23,187.24) --
	(298.83,188.99) --
	(299.03,194.21) --
	(299.24,189.40) --
	(299.44,187.25) --
	(299.64,182.57) --
	(300.25,180.94) --
	(300.45,180.76) --
	(300.65,175.53) --
	(300.86,183.18) --
	(301.06,188.73) --
	(301.66,187.53) --
	(301.87,189.91) --
	(302.07,189.26) --
	(302.27,184.14) --
	(302.47,186.70) --
	(303.08,193.60) --
	(303.28,193.53) --
	(303.49,193.28) --
	(303.69,188.23) --
	(303.89,187.85) --
	(304.70,186.75) --
	(304.90,190.99) --
	(305.10,192.75) --
	(305.31,192.30) --
	(305.91,193.62) --
	(306.12,193.58) --
	(306.32,197.64) --
	(306.52,199.91) --
	(306.72,198.76) --
	(307.33,201.14) --
	(307.53,206.28) --
	(307.73,212.93) --
	(307.94,218.77) --
	(308.14,219.46) --
	(308.75,215.17) --
	(308.95,215.33) --
	(309.15,213.67) --
	(309.35,212.41) --
	(309.56,212.44) --
	(310.16,212.07) --
	(310.36,216.65) --
	(310.57,222.78) --
	(310.77,220.13) --
	(310.97,213.92) --
	(311.58,213.99) --
	(311.78,206.98) --
	(311.98,196.63) --
	(312.19,191.67) --
	(312.39,193.66) --
	(312.99,193.17) --
	(313.20,196.61) --
	(313.40,192.34) --
	(313.60,182.35) --
	(313.80,174.27) --
	(314.41,169.57) --
	(314.61,152.02) --
	(314.82,140.05) --
	(315.02,145.98) --
	(315.22,144.94) --
	(315.83,143.01) --
	(316.03,149.28) --
	(316.23,156.24) --
	(316.43,162.34) --
	(316.64,164.30) --
	(317.24,165.42) --
	(317.45,170.68) --
	(317.65,181.09) --
	(317.85,174.40) --
	(318.05,166.32) --
	(318.66,162.88) --
	(318.86,164.61) --
	(319.06,163.06) --
	(319.27,170.85) --
	(319.47,171.82) --
	(320.08,164.64) --
	(320.28,157.95) --
	(320.48,160.84) --
	(320.89,158.04) --
	(321.49,162.52) --
	(321.69,162.16) --
	(321.90,173.28) --
	(322.10,172.31) --
	(322.30,182.12) --
	(322.91,188.03) --
	(323.11,170.67) --
	(323.52,168.95) --
	(323.72,160.26) --
	(324.32,159.97) --
	(324.53,159.70) --
	(324.73,163.54) --
	(324.93,164.50) --
	(325.13,166.07) --
	(325.74,161.92) --
	(325.94,161.56) --
	(326.15,158.06) --
	(326.35,156.66) --
	(326.55,153.99) --
	(327.16,148.36) --
	(327.56,162.44) --
	(327.76,165.89) --
	(327.97,163.98);
\end{scope}
\begin{scope}
\path[clip] ( 21.60, 41.91) rectangle (339.75,245.62);
\definecolor{drawColor}{RGB}{255,0,0}

\path[draw=drawColor,line width= 0.8pt,line join=round] ( 33.38,103.16) --
	( 33.99,100.08) --
	( 34.19, 98.87) --
	( 34.39,100.85) --
	( 34.60,103.85) --
	( 34.80,102.27) --
	( 35.41,101.37) --
	( 35.61,101.12) --
	( 35.81, 99.39) --
	( 36.01, 98.40) --
	( 36.22,100.17) --
	( 37.03,100.21) --
	( 37.23,100.57) --
	( 37.43,102.98) --
	( 37.63,101.67) --
	( 38.24,101.72) --
	( 38.44, 99.02) --
	( 38.64, 97.20) --
	( 38.85, 99.29) --
	( 39.05, 96.95) --
	( 39.66, 98.77) --
	( 39.86,101.61) --
	( 40.06,101.67) --
	( 40.26,103.64) --
	( 40.46,103.08) --
	( 41.07,102.19) --
	( 41.27,103.49) --
	( 41.48,103.43) --
	( 41.68,104.45) --
	( 41.88,104.89) --
	( 42.69,105.15) --
	( 42.89,104.99) --
	( 43.09,104.58) --
	( 43.30,106.01) --
	( 43.90,105.79) --
	( 44.11,106.65) --
	( 44.31,106.80) --
	( 44.51,106.70) --
	( 44.71,105.94) --
	( 45.32,107.39) --
	( 45.52,106.60) --
	( 45.73,105.61) --
	( 45.93,105.97) --
	( 46.13,103.38) --
	( 46.74,104.67) --
	( 46.94,101.58) --
	( 47.14,101.32) --
	( 47.34,103.74) --
	( 47.55,102.38) --
	( 48.15,104.50) --
	( 48.36,103.31) --
	( 48.56,105.42) --
	( 48.76,104.33) --
	( 48.96,105.90) --
	( 49.57,105.79) --
	( 49.77,104.81) --
	( 49.97,102.09) --
	( 50.18,101.72) --
	( 50.38,102.04) --
	( 50.99,104.49) --
	( 51.19,102.63) --
	( 51.39,101.91) --
	( 51.59,102.51) --
	( 52.40,103.60) --
	( 52.60,103.55) --
	( 52.81,103.80) --
	( 53.01,104.33) --
	( 53.21,105.25) --
	( 53.82,104.50) --
	( 54.02,105.05) --
	( 54.22,105.76) --
	( 54.42,105.69) --
	( 54.63,103.10) --
	( 55.23,105.03) --
	( 55.44,104.24) --
	( 55.64,105.07) --
	( 55.84,105.25) --
	( 56.04,105.45) --
	( 56.65,105.06) --
	( 56.85,105.73) --
	( 57.06,105.04) --
	( 57.26,103.23) --
	( 57.46,104.93) --
	( 58.07,105.36) --
	( 58.27,104.04) --
	( 58.47,103.24) --
	( 58.67,104.00) --
	( 58.88,106.48) --
	( 59.48,105.69) --
	( 59.69,105.34) --
	( 59.89,105.27) --
	( 60.09,107.05) --
	( 60.29,107.24) --
	( 60.90,107.49) --
	( 61.10,107.61) --
	( 61.30,107.36) --
	( 61.51,107.36) --
	( 61.71,106.86) --
	( 62.52,105.10) --
	( 62.72,106.22) --
	( 62.92,105.88) --
	( 63.12,104.81) --
	( 63.73,105.08) --
	( 63.93,104.82) --
	( 64.14,105.42) --
	( 64.34,103.83) --
	( 64.54,103.31) --
	( 65.15,102.54) --
	( 65.35,102.52) --
	( 65.55,104.71) --
	( 65.76,105.08) --
	( 65.96,103.77) --
	( 66.56,102.77) --
	( 66.77,103.82) --
	( 66.97,104.11) --
	( 67.17,105.79) --
	( 67.37,104.86) --
	( 67.98,105.82) --
	( 68.18,106.05) --
	( 68.39,104.39) --
	( 68.59,103.69) --
	( 68.79,104.21) --
	( 69.40,100.96) --
	( 69.60,101.17) --
	( 69.80,102.46) --
	( 70.00,102.20) --
	( 70.81,101.77) --
	( 71.02,102.64) --
	( 71.22,100.21) --
	( 71.42,100.52) --
	( 71.62,102.48) --
	( 72.23,104.50) --
	( 72.43,105.21) --
	( 72.63,105.18) --
	( 72.84,105.83) --
	( 73.04,105.51) --
	( 73.65,105.64) --
	( 73.85,103.96) --
	( 74.05,103.33) --
	( 74.25,102.22) --
	( 74.46,100.70) --
	( 75.06, 99.51) --
	( 75.26,101.27) --
	( 75.47,102.40) --
	( 75.67,102.35) --
	( 75.87,101.83) --
	( 76.48,100.98) --
	( 76.68,100.54) --
	( 76.88,100.44) --
	( 77.09, 99.32) --
	( 77.29, 98.89) --
	( 77.89,101.13) --
	( 78.10, 99.16) --
	( 78.30, 99.16) --
	( 78.50, 99.21) --
	( 78.70, 99.85) --
	( 79.31,100.48) --
	( 79.51,100.17) --
	( 79.72, 98.66) --
	( 79.92, 95.33) --
	( 80.12, 90.40) --
	( 80.73, 84.93) --
	( 80.93, 83.03) --
	( 81.13, 88.78) --
	( 81.33, 92.21) --
	( 81.54, 92.10) --
	( 82.14, 91.03) --
	( 82.35, 86.67) --
	( 82.55, 89.39) --
	( 82.75, 89.61) --
	( 82.95, 87.08) --
	( 83.76, 90.71) --
	( 83.96, 88.49) --
	( 84.17, 89.20) --
	( 84.37, 90.15) --
	( 84.98, 89.58) --
	( 85.18, 91.70) --
	( 85.38, 93.00) --
	( 85.58, 92.40) --
	( 85.79, 89.70) --
	( 86.39, 90.87) --
	( 86.59, 89.20) --
	( 86.80, 88.73) --
	( 87.00, 88.00) --
	( 87.20, 89.05) --
	( 87.81, 86.15) --
	( 88.01, 86.59) --
	( 88.21, 88.77) --
	( 88.42, 88.66) --
	( 88.62, 90.52) --
	( 89.22, 93.34) --
	( 89.43, 93.47) --
	( 89.63, 94.60) --
	( 89.83, 95.89) --
	( 90.03, 96.20) --
	( 90.64, 96.64) --
	( 90.84, 96.18) --
	( 91.05, 94.72) --
	( 91.25, 96.74) --
	( 91.45, 97.43) --
	( 92.06, 97.56) --
	( 92.26, 97.44) --
	( 92.46, 96.99) --
	( 92.66, 99.96) --
	( 92.87,101.43) --
	( 93.47,101.21) --
	( 93.68,100.82) --
	( 93.88,102.66) --
	( 94.08,102.44) --
	( 94.28,101.58) --
	( 94.89,103.12) --
	( 95.09,103.95) --
	( 95.29,103.48) --
	( 95.50,103.44) --
	( 95.70,103.88) --
	( 96.31,102.21) --
	( 96.51,102.46) --
	( 96.71,101.94) --
	( 96.91, 99.58) --
	( 97.12, 97.70) --
	( 97.72, 99.91) --
	( 97.92, 99.97) --
	( 98.13,102.27) --
	( 98.33,102.23) --
	( 98.53,103.07) --
	( 99.14,102.78) --
	( 99.34,102.96) --
	( 99.54,102.98) --
	( 99.95,102.98) --
	(100.55,102.11) --
	(100.76,103.67) --
	(100.96,102.20) --
	(101.16, 99.86) --
	(101.36,103.29) --
	(101.97,102.21) --
	(102.17,100.70) --
	(102.38, 99.99) --
	(102.58,100.76) --
	(102.78, 97.88) --
	(103.39, 98.84) --
	(103.59,100.30) --
	(103.79,102.38) --
	(103.99,100.03) --
	(104.20, 96.61) --
	(104.80, 97.76) --
	(105.01, 99.30) --
	(105.21,101.02) --
	(105.41,100.55) --
	(106.22,100.33) --
	(106.42,102.12) --
	(106.62,101.03) --
	(106.83, 99.37) --
	(107.64, 96.80) --
	(107.84, 96.89) --
	(108.04, 94.55) --
	(108.24, 90.91) --
	(108.45, 89.35) --
	(109.05, 89.83) --
	(109.25, 90.93) --
	(109.46, 87.54) --
	(109.66, 89.65) --
	(109.86, 86.02) --
	(110.67, 86.28) --
	(110.87, 83.96) --
	(111.08, 85.04) --
	(111.28, 87.00) --
	(111.88, 85.06) --
	(112.09, 87.68) --
	(112.29, 85.61) --
	(112.49, 86.78) --
	(112.69, 90.46) --
	(113.30, 90.30) --
	(113.50, 87.56) --
	(113.71, 89.26) --
	(113.91, 90.00) --
	(114.11, 88.03) --
	(114.72, 86.38) --
	(114.92, 86.26) --
	(115.12, 85.34) --
	(115.32, 82.97) --
	(115.53, 85.89) --
	(116.34, 87.95) --
	(116.54, 90.35) --
	(116.74, 89.97) --
	(116.94, 89.77) --
	(117.55, 91.90) --
	(117.75, 90.14) --
	(117.95, 90.64) --
	(118.16, 92.61) --
	(118.36, 92.07) --
	(118.97, 90.93) --
	(119.17, 94.17) --
	(119.37, 94.48) --
	(119.57, 94.90) --
	(119.78, 95.48) --
	(120.38, 96.11) --
	(120.58, 95.09) --
	(120.79, 95.42) --
	(120.99, 95.37) --
	(121.19, 97.40) --
	(121.80, 97.55) --
	(122.00, 97.76) --
	(122.20, 98.45) --
	(122.41, 99.89) --
	(122.61,101.01) --
	(123.22,101.22) --
	(123.42,100.83) --
	(123.62,100.09) --
	(123.82,100.21) --
	(124.63,100.39) --
	(124.83,101.30) --
	(125.04,102.08) --
	(125.24,101.78) --
	(125.44,102.78) --
	(126.05,102.27) --
	(126.25,101.02) --
	(126.45,102.07) --
	(126.65,100.45) --
	(126.86,100.78) --
	(127.46,100.59) --
	(127.67,102.12) --
	(127.87,103.86) --
	(128.07,104.03) --
	(128.27,103.76) --
	(128.88,104.75) --
	(129.08,105.21) --
	(129.28,105.60) --
	(129.49,104.55) --
	(129.69,104.74) --
	(130.30,104.50) --
	(130.50,104.62) --
	(130.70,105.10) --
	(130.90,103.14) --
	(131.11,102.61) --
	(131.71,103.70) --
	(131.91,102.40) --
	(132.12,101.47) --
	(132.32,101.56) --
	(132.52,102.30) --
	(133.13,101.98) --
	(133.33,104.04) --
	(133.53,102.03) --
	(133.74,102.11) --
	(133.94,100.39) --
	(134.55,102.02) --
	(134.75,100.34) --
	(134.95,100.31) --
	(135.15, 99.46) --
	(135.35,100.07) --
	(135.96,100.00) --
	(136.16,101.98) --
	(136.37,103.33) --
	(136.57,103.11) --
	(136.77,103.53) --
	(137.58,102.73) --
	(137.78,102.76) --
	(137.98,103.21) --
	(138.19,102.92) --
	(138.79,103.97) --
	(139.00,104.14) --
	(139.20,104.76) --
	(139.40,104.57) --
	(139.60,103.46) --
	(140.21,102.22) --
	(140.41,101.69) --
	(140.61,101.37) --
	(140.82,102.23) --
	(141.02,101.69) --
	(141.63,102.90) --
	(141.83,103.13) --
	(142.03,102.67) --
	(142.23,104.81) --
	(142.44, 99.14) --
	(143.04, 96.72) --
	(143.25, 99.23) --
	(143.45,101.87) --
	(143.65,104.06) --
	(143.85,104.24) --
	(144.66,103.23) --
	(144.86,103.95) --
	(145.07,103.74) --
	(145.27,106.07) --
	(145.88,106.82) --
	(146.08,107.94) --
	(146.28,108.17) --
	(146.48,109.41) --
	(146.68,109.51) --
	(147.29,109.66) --
	(147.49,109.90) --
	(147.70,110.24) --
	(147.90,109.51) --
	(148.10,110.01) --
	(148.71,109.29) --
	(148.91,109.11) --
	(149.11,109.10) --
	(149.31,108.95) --
	(149.52,108.73) --
	(150.12,108.47) --
	(150.33,107.62) --
	(150.53,108.01) --
	(150.73,107.98) --
	(150.93,109.76) --
	(151.54,109.63) --
	(151.74,109.66) --
	(151.95,109.31) --
	(152.15,110.41) --
	(152.35,110.06) --
	(152.96,110.62) --
	(153.16,109.84) --
	(153.36,110.04) --
	(153.56,110.26) --
	(153.77,109.84) --
	(154.37,109.63) --
	(154.58,109.79) --
	(154.78,109.18) --
	(154.98,108.88) --
	(155.18,108.38) --
	(155.79,109.38) --
	(155.99,108.93) --
	(156.19,108.43) --
	(156.40,108.60) --
	(156.60,109.28) --
	(157.41,109.71) --
	(157.61,109.60) --
	(157.81,109.17) --
	(158.01,105.50) --
	(158.62,107.73) --
	(158.82,105.33) --
	(159.03,105.03) --
	(159.23,106.68) --
	(159.43,105.86) --
	(160.04,105.83) --
	(160.24,105.92) --
	(160.44,107.44) --
	(160.64,108.36) --
	(160.85,107.14) --
	(161.45,105.59) --
	(161.66,106.83) --
	(161.86,107.86) --
	(162.06,106.04) --
	(162.26,107.57) --
	(162.87,107.07) --
	(163.07,106.27) --
	(163.28,107.32) --
	(163.48,107.20) --
	(163.68,106.94) --
	(164.29,107.77) --
	(164.49,105.90) --
	(164.69,106.05) --
	(164.89,105.63) --
	(165.10,106.00) --
	(165.70,105.51) --
	(165.91,106.21) --
	(166.11,106.59) --
	(166.31,106.22) --
	(166.51,106.06) --
	(167.12,106.78) --
	(167.32,106.28) --
	(167.52,106.56) --
	(167.73,106.29) --
	(167.93,106.21) --
	(168.54,106.03) --
	(168.74,105.05) --
	(168.94,104.33) --
	(169.14,104.07) --
	(169.34,103.67) --
	(169.95,107.12) --
	(170.15,107.80) --
	(170.36,110.19) --
	(170.56,112.21) --
	(170.76,112.58) --
	(171.37,112.78) --
	(171.57,113.28) --
	(171.77,112.77) --
	(171.98,113.11) --
	(172.18,112.77) --
	(172.78,113.60) --
	(172.99,114.22) --
	(173.19,114.77) --
	(173.59,115.41) --
	(174.20,114.91) --
	(174.40,115.13) --
	(174.61,115.15) --
	(174.81,115.78) --
	(175.01,115.58) --
	(175.62,116.01) --
	(175.82,116.34) --
	(176.02,119.11) --
	(176.22,119.71) --
	(176.43,121.03) --
	(177.03,121.40) --
	(177.24,122.46) --
	(177.44,121.36) --
	(177.64,121.92) --
	(177.84,121.83) --
	(178.45,122.20) --
	(178.65,123.05) --
	(178.85,122.75) --
	(179.06,122.54) --
	(179.26,122.67) --
	(180.07,122.78) --
	(180.27,121.74) --
	(180.47,121.62) --
	(180.67,121.08) --
	(181.48,122.19) --
	(181.69,122.75) --
	(181.89,122.35) --
	(182.09,122.95) --
	(182.70,122.24) --
	(182.90,121.95) --
	(183.10,122.86) --
	(183.31,122.28) --
	(183.51,122.23) --
	(184.32,121.68) --
	(184.52,121.48) --
	(184.72,120.80) --
	(184.92,121.68) --
	(185.53,121.43) --
	(185.73,122.48) --
	(185.94,123.93) --
	(186.14,124.23) --
	(186.34,124.16) --
	(186.95,123.02) --
	(187.15,122.03) --
	(187.35,122.28) --
	(187.55,122.22) --
	(187.76,123.95) --
	(188.36,123.78) --
	(188.57,124.13) --
	(188.77,123.79) --
	(188.97,124.89) --
	(189.17,125.79) --
	(189.78,127.12) --
	(189.98,127.98) --
	(190.18,128.97) --
	(190.39,129.05) --
	(190.59,129.09) --
	(191.40,130.19) --
	(191.60,130.49) --
	(191.80,130.82) --
	(192.01,130.92) --
	(192.61,131.07) --
	(192.81,130.84) --
	(193.02,133.65) --
	(193.22,132.61) --
	(193.42,132.63) --
	(194.03,132.16) --
	(194.23,131.88) --
	(194.43,131.24) --
	(194.64,131.26) --
	(194.84,131.68) --
	(195.44,131.48) --
	(195.65,131.07) --
	(195.85,132.12) --
	(196.05,131.97) --
	(196.25,131.79) --
	(196.86,131.71) --
	(197.06,129.50) --
	(197.27,129.43) --
	(197.47,129.39) --
	(197.67,128.83) --
	(198.28,128.41) --
	(198.48,129.81) --
	(198.68,129.41) --
	(198.88,130.06) --
	(199.09,129.45) --
	(199.69,129.33) --
	(199.90,129.69) --
	(200.10,129.31) --
	(200.30,129.45) --
	(200.50,129.38) --
	(201.11,129.40) --
	(201.31,129.34) --
	(201.51,128.79) --
	(201.72,127.50) --
	(202.53,129.21) --
	(202.73,128.15) --
	(202.93,127.05) --
	(203.13,128.67) --
	(203.34,128.38) --
	(203.94,130.39) --
	(204.14,132.54) --
	(204.35,132.35) --
	(204.55,132.41) --
	(204.75,132.03) --
	(205.36,131.78) --
	(205.56,132.11) --
	(205.76,132.19) --
	(205.97,132.13) --
	(206.17,132.64) --
	(206.77,132.69) --
	(206.98,132.35) --
	(207.18,132.05) --
	(207.38,131.83) --
	(207.58,131.62) --
	(208.19,132.41) --
	(208.39,132.39) --
	(208.60,128.93) --
	(208.80,129.45) --
	(209.00,130.77) --
	(209.61,131.60) --
	(209.81,132.00) --
	(210.01,132.69) --
	(210.21,133.35) --
	(210.42,133.33) --
	(211.23,132.85) --
	(211.43,132.66) --
	(211.63,133.92) --
	(211.83,134.50) --
	(212.44,134.29) --
	(212.64,133.85) --
	(212.84,134.19) --
	(213.05,134.28) --
	(213.25,135.11) --
	(213.86,134.77) --
	(214.06,135.63) --
	(214.26,136.06) --
	(214.46,135.92) --
	(214.67,136.15) --
	(215.27,137.49) --
	(215.47,136.92) --
	(215.68,136.39) --
	(215.88,136.27) --
	(216.08,136.25) --
	(216.69,136.38) --
	(216.89,135.47) --
	(217.09,136.80) --
	(217.30,135.25) --
	(217.50,135.83) --
	(218.10,137.03) --
	(218.51,137.02) --
	(218.71,135.55) --
	(218.91,136.43) --
	(219.52,136.37) --
	(219.72,136.38) --
	(219.93,137.52) --
	(220.13,137.72) --
	(220.33,138.50) --
	(220.94,138.43) --
	(221.14,137.92) --
	(221.34,138.53) --
	(221.54,138.26) --
	(221.75,137.97) --
	(222.35,137.35) --
	(222.56,138.28) --
	(222.76,139.18) --
	(222.96,139.98) --
	(223.16,140.29) --
	(223.77,140.86) --
	(223.97,141.53) --
	(224.17,142.02) --
	(224.38,142.11) --
	(224.58,142.73) --
	(225.19,142.97) --
	(225.39,142.66) --
	(225.59,142.32) --
	(225.79,140.42) --
	(226.00,140.55) --
	(226.60,141.81) --
	(226.80,141.86) --
	(227.01,142.10) --
	(227.21,139.55) --
	(227.41,138.85) --
	(228.02,139.12) --
	(228.22,140.94) --
	(228.42,140.12) --
	(228.63,139.86) --
	(228.83,140.14) --
	(229.44,140.09) --
	(229.64,140.62) --
	(229.84,140.87) --
	(230.04,141.39) --
	(230.24,141.75) --
	(231.05,139.58) --
	(231.26,140.08) --
	(231.46,139.87) --
	(231.66,139.99) --
	(232.27,142.40) --
	(232.47,142.97) --
	(232.67,143.34) --
	(232.87,143.76) --
	(233.08,144.36) --
	(233.68,144.95) --
	(233.89,145.31) --
	(234.09,145.70) --
	(234.29,145.21) --
	(234.49,145.12) --
	(235.10,144.62) --
	(235.30,144.51) --
	(235.50,145.03) --
	(235.71,145.41) --
	(235.91,145.63) --
	(236.52,147.05) --
	(236.72,147.83) --
	(236.92,148.02) --
	(237.12,149.07) --
	(237.33,149.06) --
	(237.93,148.94) --
	(238.13,149.59) --
	(238.34,149.98) --
	(238.54,149.68) --
	(238.74,149.97) --
	(239.35,150.76) --
	(239.55,151.14) --
	(239.75,152.62) --
	(239.96,152.67) --
	(240.16,154.21) --
	(240.77,153.70) --
	(240.97,155.26) --
	(241.17,154.22) --
	(241.37,154.88) --
	(241.57,155.19) --
	(242.18,154.40) --
	(242.38,154.66) --
	(242.59,155.20) --
	(242.79,155.96) --
	(242.99,156.17) --
	(243.60,156.25) --
	(243.80,156.34) --
	(244.00,156.39) --
	(244.20,155.45) --
	(244.41,155.08) --
	(245.01,155.24) --
	(245.22,154.96) --
	(245.42,153.68) --
	(245.62,155.42) --
	(245.82,154.49) --
	(246.43,155.16) --
	(246.63,156.65) --
	(246.83,156.05) --
	(247.24,156.34) --
	(247.85,156.56) --
	(248.05,158.93) --
	(248.25,159.90) --
	(248.45,162.98) --
	(248.66,162.60) --
	(249.26,163.14) --
	(249.47,162.13) --
	(249.67,161.76) --
	(249.87,162.41) --
	(250.07,163.51) --
	(250.68,164.04) --
	(250.88,165.14) --
	(251.08,165.89) --
	(251.29,165.18) --
	(251.49,166.50) --
	(252.10,167.81) --
	(252.30,167.46) --
	(252.50,167.20) --
	(252.70,167.72) --
	(252.90,167.45) --
	(253.71,167.38) --
	(253.92,167.64) --
	(254.12,168.23) --
	(254.32,167.13) --
	(255.13,168.10) --
	(255.33,169.02) --
	(255.53,170.44) --
	(255.74,172.49) --
	(256.34,172.37) --
	(256.55,173.32) --
	(256.75,173.17) --
	(256.95,175.08) --
	(257.15,177.20) --
	(257.96,177.11) --
	(258.16,180.10) --
	(258.37,179.20) --
	(258.57,179.70) --
	(259.18,181.02) --
	(259.38,180.99) --
	(259.58,181.37) --
	(259.78,182.68) --
	(259.99,184.76) --
	(260.59,183.11) --
	(260.80,179.74) --
	(261.00,180.42) --
	(261.20,180.76) --
	(261.40,174.58) --
	(262.01,163.66) --
	(262.21,168.93) --
	(262.41,168.75) --
	(262.62,159.15) --
	(262.82,162.22) --
	(263.43,166.04) --
	(263.63,166.40) --
	(263.83,168.75) --
	(264.03,171.60) --
	(264.23,171.78) --
	(265.04,169.41) --
	(265.25,167.86) --
	(265.45,169.39) --
	(265.65,172.62) --
	(266.26,176.33) --
	(266.46,173.55) --
	(266.66,170.01) --
	(266.86,166.11) --
	(267.07,165.45) --
	(267.67,168.58) --
	(267.88,168.66) --
	(268.08,167.89) --
	(268.28,168.77) --
	(268.48,172.86) --
	(269.09,171.40) --
	(269.29,169.80) --
	(269.50,167.49) --
	(269.70,168.57) --
	(269.90,169.24) --
	(270.51,166.13) --
	(270.71,167.21) --
	(270.91,166.79) --
	(271.11,160.06) --
	(271.32,156.11) --
	(271.92,162.33) --
	(272.13,159.13) --
	(272.33,159.04) --
	(272.53,161.41) --
	(273.34,157.14) --
	(273.54,160.76) --
	(273.74,162.90) --
	(273.95,165.14) --
	(274.15,159.82) --
	(274.76,160.26) --
	(274.96,164.24) --
	(275.16,162.21) --
	(275.36,164.94) --
	(275.56,163.80) --
	(276.17,165.77) --
	(276.37,167.76) --
	(276.58,167.40) --
	(276.78,166.63) --
	(276.98,164.75) --
	(277.59,164.62) --
	(277.79,160.67) --
	(277.99,161.23) --
	(278.20,163.44) --
	(278.40,163.34) --
	(279.00,161.97) --
	(279.21,161.37) --
	(279.41,159.75) --
	(279.61,159.80) --
	(279.81,162.89) --
	(280.42,163.77) --
	(280.62,163.80) --
	(280.83,165.49) --
	(281.03,167.32) --
	(281.23,168.17) --
	(281.84,168.81) --
	(282.04,167.01) --
	(282.24,167.59) --
	(282.44,167.08) --
	(282.65,167.09) --
	(283.25,169.86) --
	(283.46,168.20) --
	(283.66,168.69) --
	(283.86,167.99) --
	(284.06,167.45) --
	(284.87,163.81) --
	(285.07,166.65) --
	(285.28,164.31) --
	(285.48,166.35) --
	(286.09,168.01) --
	(286.29,167.88) --
	(286.49,171.10) --
	(286.69,171.98) --
	(286.89,172.68) --
	(287.50,172.73) --
	(287.70,172.72) --
	(287.91,171.61) --
	(288.11,171.37) --
	(288.31,170.58) --
	(288.92,169.62) --
	(289.12,166.95) --
	(289.32,166.56) --
	(289.53,164.74) --
	(289.73,165.85) --
	(290.33,162.80) --
	(290.54,163.08) --
	(290.74,161.54) --
	(290.94,162.46) --
	(291.14,162.97) --
	(291.75,163.30) --
	(291.95,162.07) --
	(292.36,163.76) --
	(292.56,164.69) --
	(293.17,167.66) --
	(293.37,168.99) --
	(293.57,166.96) --
	(293.77,169.04) --
	(293.98,169.92) --
	(294.58,170.34) --
	(294.79,170.85) --
	(294.99,171.59) --
	(295.19,170.34) --
	(295.39,170.28) --
	(296.00,170.15) --
	(296.20,171.99) --
	(296.40,173.59) --
	(296.61,174.64) --
	(296.81,173.93) --
	(297.42,172.59) --
	(297.62,173.60) --
	(297.82,172.84) --
	(298.02,172.77) --
	(298.23,174.04) --
	(298.83,174.40) --
	(299.03,175.58) --
	(299.24,175.16) --
	(299.44,174.47) --
	(299.64,172.65) --
	(300.25,171.48) --
	(300.45,172.53) --
	(300.65,171.25) --
	(300.86,174.93) --
	(301.06,175.96) --
	(301.66,176.79) --
	(301.87,177.38) --
	(302.07,176.55) --
	(302.27,175.84) --
	(302.47,177.08) --
	(303.08,179.49) --
	(303.28,179.62) --
	(303.49,180.19) --
	(303.69,178.91) --
	(303.89,178.70) --
	(304.70,178.59) --
	(304.90,178.80) --
	(305.10,178.99) --
	(305.31,178.25) --
	(305.91,177.70) --
	(306.12,178.76) --
	(306.32,179.02) --
	(306.52,180.39) --
	(306.72,180.47) --
	(307.33,179.61) --
	(307.53,181.32) --
	(307.73,182.80) --
	(307.94,185.13) --
	(308.14,185.94) --
	(308.75,184.25) --
	(308.95,183.60) --
	(309.15,182.61) --
	(309.35,183.12) --
	(309.56,183.29) --
	(310.16,185.08) --
	(310.36,186.22) --
	(310.57,186.73) --
	(310.77,184.86) --
	(310.97,183.18) --
	(311.58,183.55) --
	(311.78,183.03) --
	(311.98,175.30) --
	(312.19,170.23) --
	(312.39,172.90) --
	(312.99,172.07) --
	(313.20,177.16) --
	(313.40,176.30) --
	(313.60,173.26) --
	(313.80,173.87) --
	(314.41,172.69) --
	(314.61,171.52) --
	(314.82,165.87) --
	(315.02,169.60) --
	(315.22,166.84) --
	(315.83,164.56) --
	(316.03,168.58) --
	(316.23,170.82) --
	(316.43,173.28) --
	(316.64,172.26) --
	(317.24,174.03) --
	(317.45,175.64) --
	(317.65,180.70) --
	(317.85,180.81) --
	(318.05,178.93) --
	(318.66,173.34) --
	(318.86,172.40) --
	(319.06,170.49) --
	(319.27,172.43) --
	(319.47,173.58) --
	(320.08,169.90) --
	(320.28,164.78) --
	(320.48,164.77) --
	(320.89,163.11) --
	(321.49,166.40) --
	(321.69,167.41) --
	(321.90,173.14) --
	(322.10,172.89) --
	(322.30,174.74) --
	(322.91,177.42) --
	(323.11,169.99) --
	(323.52,169.25) --
	(323.72,164.06) --
	(324.32,164.38) --
	(324.53,163.89) --
	(324.73,165.35) --
	(324.93,166.00) --
	(325.13,161.38) --
	(325.74,156.67) --
	(325.94,157.44) --
	(326.15,154.17) --
	(326.35,149.86) --
	(326.55,146.01) --
	(327.16,139.94) --
	(327.56,150.03) --
	(327.76,152.45) --
	(327.97,151.74);
\end{scope}
\begin{scope}
\path[clip] ( 21.60, 41.91) rectangle (339.75,245.62);
\definecolor{drawColor}{RGB}{0,0,255}

\path[draw=drawColor,line width= 0.8pt,line join=round,line cap=round] (280.62,234.82) -- (296.82,234.82);
\definecolor{drawColor}{RGB}{255,0,0}

\path[draw=drawColor,line width= 0.8pt,line join=round,line cap=round] (280.62,224.02) -- (296.82,224.02);
\definecolor{drawColor}{gray}{0.20}

\node[text=drawColor,anchor=base west,inner sep=0pt, outer sep=0pt, scale=  0.90] at (304.92,231.72) {CAT};

\node[text=drawColor,anchor=base west,inner sep=0pt, outer sep=0pt, scale=  0.90] at (304.92,220.92) {DJI};
\end{scope}
\end{tikzpicture}

	\caption{Evolució d'una inversió de \$100 en Caterpillar i l'índex DJI entre l'any 2000 i el 2018}
	\label{fig:inversio}
\end{figure}

\begin{figure}[htb]
	\centering \sffamily \small
	% Created by tikzDevice version 0.12 on 2019-06-13 23:40:41
% !TEX encoding = UTF-8 Unicode
\begin{tikzpicture}[x=1pt,y=1pt]
\definecolor{fillColor}{RGB}{255,255,255}
\path[use as bounding box,fill=fillColor,fill opacity=0.00] (0,0) rectangle (252.94,252.94);
\begin{scope}
\path[clip] ( 48.00,  0.00) rectangle (252.94, 50.59);
\definecolor{fillColor}{RGB}{173,216,230}

\path[fill=fillColor] (142.38, 15.93) --
	(142.38, 34.66) --
	(163.94, 34.66) --
	(163.94, 15.93) --
	cycle;
\definecolor{drawColor}{RGB}{0,0,0}

\path[draw=drawColor,line width= 0.4pt,line join=round] (152.60, 15.93) -- (152.60, 34.66);

\path[draw=drawColor,line width= 0.4pt,dash pattern=on 4pt off 4pt ,line join=round,line cap=round] (110.55, 25.29) -- (142.38, 25.29);

\path[draw=drawColor,line width= 0.4pt,dash pattern=on 4pt off 4pt ,line join=round,line cap=round] (195.67, 25.29) -- (163.94, 25.29);

\path[draw=drawColor,line width= 0.4pt,line join=round,line cap=round] (110.55, 20.61) -- (110.55, 29.98);

\path[draw=drawColor,line width= 0.4pt,line join=round,line cap=round] (195.67, 20.61) -- (195.67, 29.98);
\definecolor{drawColor}{RGB}{255,255,255}

\path[draw=drawColor,line width= 0.4pt,line join=round,line cap=round] (142.38, 15.93) --
	(142.38, 34.66) --
	(163.94, 34.66) --
	(163.94, 15.93) --
	(142.38, 15.93);
\definecolor{drawColor}{gray}{0.40}

\path[draw=drawColor,line width= 0.4pt,line join=round,line cap=round] ( 83.32, 23.04) -- ( 87.82, 27.54);

\path[draw=drawColor,line width= 0.4pt,line join=round,line cap=round] ( 83.32, 27.54) -- ( 87.82, 23.04);

\path[draw=drawColor,line width= 0.4pt,line join=round,line cap=round] ( 58.33, 23.04) -- ( 62.83, 27.54);

\path[draw=drawColor,line width= 0.4pt,line join=round,line cap=round] ( 58.33, 27.54) -- ( 62.83, 23.04);

\path[draw=drawColor,line width= 0.4pt,line join=round,line cap=round] (196.02, 23.04) -- (200.52, 27.54);

\path[draw=drawColor,line width= 0.4pt,line join=round,line cap=round] (196.02, 27.54) -- (200.52, 23.04);

\path[draw=drawColor,line width= 0.4pt,line join=round,line cap=round] (194.37, 23.04) -- (198.87, 27.54);

\path[draw=drawColor,line width= 0.4pt,line join=round,line cap=round] (194.37, 27.54) -- (198.87, 23.04);

\path[draw=drawColor,line width= 0.4pt,line join=round,line cap=round] (104.79, 23.04) -- (109.29, 27.54);

\path[draw=drawColor,line width= 0.4pt,line join=round,line cap=round] (104.79, 27.54) -- (109.29, 23.04);

\path[draw=drawColor,line width= 0.4pt,line join=round,line cap=round] (194.87, 23.04) -- (199.37, 27.54);

\path[draw=drawColor,line width= 0.4pt,line join=round,line cap=round] (194.87, 27.54) -- (199.37, 23.04);

\path[draw=drawColor,line width= 0.4pt,line join=round,line cap=round] ( 70.41, 23.04) -- ( 74.91, 27.54);

\path[draw=drawColor,line width= 0.4pt,line join=round,line cap=round] ( 70.41, 27.54) -- ( 74.91, 23.04);

\path[draw=drawColor,line width= 0.4pt,line join=round,line cap=round] (213.51, 23.04) -- (218.01, 27.54);

\path[draw=drawColor,line width= 0.4pt,line join=round,line cap=round] (213.51, 27.54) -- (218.01, 23.04);

\path[draw=drawColor,line width= 0.4pt,line join=round,line cap=round] ( 93.14, 23.04) -- ( 97.64, 27.54);

\path[draw=drawColor,line width= 0.4pt,line join=round,line cap=round] ( 93.14, 27.54) -- ( 97.64, 23.04);

\path[draw=drawColor,line width= 0.4pt,line join=round,line cap=round] ( 97.71, 23.04) -- (102.21, 27.54);

\path[draw=drawColor,line width= 0.4pt,line join=round,line cap=round] ( 97.71, 27.54) -- (102.21, 23.04);

\path[draw=drawColor,line width= 0.4pt,line join=round,line cap=round] (208.39, 23.04) -- (212.89, 27.54);

\path[draw=drawColor,line width= 0.4pt,line join=round,line cap=round] (208.39, 27.54) -- (212.89, 23.04);

\path[draw=drawColor,line width= 0.4pt,line join=round,line cap=round] (106.93, 23.04) -- (111.43, 27.54);

\path[draw=drawColor,line width= 0.4pt,line join=round,line cap=round] (106.93, 27.54) -- (111.43, 23.04);

\path[draw=drawColor,line width= 0.4pt,line join=round,line cap=round] (101.68, 23.04) -- (106.18, 27.54);

\path[draw=drawColor,line width= 0.4pt,line join=round,line cap=round] (101.68, 27.54) -- (106.18, 23.04);

\path[draw=drawColor,line width= 0.4pt,line join=round,line cap=round] ( 86.49, 23.04) -- ( 90.99, 27.54);

\path[draw=drawColor,line width= 0.4pt,line join=round,line cap=round] ( 86.49, 27.54) -- ( 90.99, 23.04);

\path[draw=drawColor,line width= 0.4pt,line join=round,line cap=round] (206.90, 23.04) -- (211.40, 27.54);

\path[draw=drawColor,line width= 0.4pt,line join=round,line cap=round] (206.90, 27.54) -- (211.40, 23.04);

\path[draw=drawColor,line width= 0.4pt,line join=round,line cap=round] (201.64, 23.04) -- (206.14, 27.54);

\path[draw=drawColor,line width= 0.4pt,line join=round,line cap=round] (201.64, 27.54) -- (206.14, 23.04);

\path[draw=drawColor,line width= 0.4pt,line join=round,line cap=round] (201.24, 23.04) -- (205.74, 27.54);

\path[draw=drawColor,line width= 0.4pt,line join=round,line cap=round] (201.24, 27.54) -- (205.74, 23.04);

\path[draw=drawColor,line width= 0.4pt,line join=round,line cap=round] ( 98.97, 23.04) -- (103.47, 27.54);

\path[draw=drawColor,line width= 0.4pt,line join=round,line cap=round] ( 98.97, 27.54) -- (103.47, 23.04);

\path[draw=drawColor,line width= 0.4pt,line join=round,line cap=round] (198.05, 23.04) -- (202.55, 27.54);

\path[draw=drawColor,line width= 0.4pt,line join=round,line cap=round] (198.05, 27.54) -- (202.55, 23.04);

\path[draw=drawColor,line width= 0.4pt,line join=round,line cap=round] (105.99, 23.04) -- (110.49, 27.54);

\path[draw=drawColor,line width= 0.4pt,line join=round,line cap=round] (105.99, 27.54) -- (110.49, 23.04);

\path[draw=drawColor,line width= 0.4pt,line join=round,line cap=round] ( 65.58, 23.04) -- ( 70.08, 27.54);

\path[draw=drawColor,line width= 0.4pt,line join=round,line cap=round] ( 65.58, 27.54) -- ( 70.08, 23.04);

\path[draw=drawColor,line width= 0.4pt,line join=round,line cap=round] (211.91, 23.04) -- (216.41, 27.54);

\path[draw=drawColor,line width= 0.4pt,line join=round,line cap=round] (211.91, 27.54) -- (216.41, 23.04);

\path[draw=drawColor,line width= 0.4pt,line join=round,line cap=round] (203.89, 23.04) -- (208.39, 27.54);

\path[draw=drawColor,line width= 0.4pt,line join=round,line cap=round] (203.89, 27.54) -- (208.39, 23.04);

\path[draw=drawColor,line width= 0.4pt,line join=round,line cap=round] (241.26, 23.04) -- (245.76, 27.54);

\path[draw=drawColor,line width= 0.4pt,line join=round,line cap=round] (241.26, 27.54) -- (245.76, 23.04);

\path[draw=drawColor,line width= 0.4pt,line join=round,line cap=round] ( 96.21, 23.04) -- (100.71, 27.54);

\path[draw=drawColor,line width= 0.4pt,line join=round,line cap=round] ( 96.21, 27.54) -- (100.71, 23.04);

\path[draw=drawColor,line width= 0.4pt,line join=round,line cap=round] (243.10, 23.04) -- (247.60, 27.54);

\path[draw=drawColor,line width= 0.4pt,line join=round,line cap=round] (243.10, 27.54) -- (247.60, 23.04);

\path[draw=drawColor,line width= 0.4pt,line join=round,line cap=round] (220.30, 23.04) -- (224.80, 27.54);

\path[draw=drawColor,line width= 0.4pt,line join=round,line cap=round] (220.30, 27.54) -- (224.80, 23.04);

\path[draw=drawColor,line width= 0.4pt,line join=round,line cap=round] (209.83, 23.04) -- (214.33, 27.54);

\path[draw=drawColor,line width= 0.4pt,line join=round,line cap=round] (209.83, 27.54) -- (214.33, 23.04);

\path[draw=drawColor,line width= 0.4pt,line join=round,line cap=round] ( 98.92, 23.04) -- (103.42, 27.54);

\path[draw=drawColor,line width= 0.4pt,line join=round,line cap=round] ( 98.92, 27.54) -- (103.42, 23.04);

\path[draw=drawColor,line width= 0.4pt,line join=round,line cap=round] ( 81.61, 23.04) -- ( 86.11, 27.54);

\path[draw=drawColor,line width= 0.4pt,line join=round,line cap=round] ( 81.61, 27.54) -- ( 86.11, 23.04);

\path[draw=drawColor,line width= 0.4pt,line join=round,line cap=round] ( 98.39, 23.04) -- (102.89, 27.54);

\path[draw=drawColor,line width= 0.4pt,line join=round,line cap=round] ( 98.39, 27.54) -- (102.89, 23.04);

\path[draw=drawColor,line width= 0.4pt,line join=round,line cap=round] ( 77.68, 23.04) -- ( 82.18, 27.54);

\path[draw=drawColor,line width= 0.4pt,line join=round,line cap=round] ( 77.68, 27.54) -- ( 82.18, 23.04);

\path[draw=drawColor,line width= 0.4pt,line join=round,line cap=round] (106.59, 23.04) -- (111.09, 27.54);

\path[draw=drawColor,line width= 0.4pt,line join=round,line cap=round] (106.59, 27.54) -- (111.09, 23.04);

\path[draw=drawColor,line width= 0.4pt,line join=round,line cap=round] ( 71.28, 23.04) -- ( 75.78, 27.54);

\path[draw=drawColor,line width= 0.4pt,line join=round,line cap=round] ( 71.28, 27.54) -- ( 75.78, 23.04);

\path[draw=drawColor,line width= 0.4pt,line join=round,line cap=round] (104.69, 23.04) -- (109.19, 27.54);

\path[draw=drawColor,line width= 0.4pt,line join=round,line cap=round] (104.69, 27.54) -- (109.19, 23.04);

\path[draw=drawColor,line width= 0.4pt,line join=round,line cap=round] (199.52, 23.04) -- (204.02, 27.54);

\path[draw=drawColor,line width= 0.4pt,line join=round,line cap=round] (199.52, 27.54) -- (204.02, 23.04);

\path[draw=drawColor,line width= 0.4pt,line join=round,line cap=round] (104.19, 23.04) -- (108.69, 27.54);

\path[draw=drawColor,line width= 0.4pt,line join=round,line cap=round] (104.19, 27.54) -- (108.69, 23.04);

\path[draw=drawColor,line width= 0.4pt,line join=round,line cap=round] (101.86, 23.04) -- (106.36, 27.54);

\path[draw=drawColor,line width= 0.4pt,line join=round,line cap=round] (101.86, 27.54) -- (106.36, 23.04);

\path[draw=drawColor,line width= 0.4pt,line join=round,line cap=round] (100.85, 23.04) -- (105.35, 27.54);

\path[draw=drawColor,line width= 0.4pt,line join=round,line cap=round] (100.85, 27.54) -- (105.35, 23.04);

\path[draw=drawColor,line width= 0.4pt,line join=round,line cap=round] ( 53.34, 23.04) -- ( 57.84, 27.54);

\path[draw=drawColor,line width= 0.4pt,line join=round,line cap=round] ( 53.34, 27.54) -- ( 57.84, 23.04);

\path[draw=drawColor,line width= 0.4pt,line join=round,line cap=round] ( 79.39, 23.04) -- ( 83.89, 27.54);

\path[draw=drawColor,line width= 0.4pt,line join=round,line cap=round] ( 79.39, 27.54) -- ( 83.89, 23.04);

\path[draw=drawColor,line width= 0.4pt,line join=round,line cap=round] (203.72, 23.04) -- (208.22, 27.54);

\path[draw=drawColor,line width= 0.4pt,line join=round,line cap=round] (203.72, 27.54) -- (208.22, 23.04);

\path[draw=drawColor,line width= 0.4pt,line join=round,line cap=round] (107.36, 23.04) -- (111.86, 27.54);

\path[draw=drawColor,line width= 0.4pt,line join=round,line cap=round] (107.36, 27.54) -- (111.86, 23.04);

\path[draw=drawColor,line width= 0.4pt,line join=round,line cap=round] (209.48, 23.04) -- (213.98, 27.54);

\path[draw=drawColor,line width= 0.4pt,line join=round,line cap=round] (209.48, 27.54) -- (213.98, 23.04);

\path[draw=drawColor,line width= 0.4pt,line join=round,line cap=round] (200.36, 23.04) -- (204.86, 27.54);

\path[draw=drawColor,line width= 0.4pt,line join=round,line cap=round] (200.36, 27.54) -- (204.86, 23.04);

\path[draw=drawColor,line width= 0.4pt,line join=round,line cap=round] ( 61.68, 23.04) -- ( 66.18, 27.54);

\path[draw=drawColor,line width= 0.4pt,line join=round,line cap=round] ( 61.68, 27.54) -- ( 66.18, 23.04);

\path[draw=drawColor,line width= 0.4pt,line join=round,line cap=round] (102.95, 23.04) -- (107.45, 27.54);

\path[draw=drawColor,line width= 0.4pt,line join=round,line cap=round] (102.95, 27.54) -- (107.45, 23.04);

\path[draw=drawColor,line width= 0.4pt,line join=round,line cap=round] (229.54, 23.04) -- (234.04, 27.54);

\path[draw=drawColor,line width= 0.4pt,line join=round,line cap=round] (229.54, 27.54) -- (234.04, 23.04);
\end{scope}
\begin{scope}
\path[clip] (  0.00, 50.59) rectangle (252.94,252.94);
\definecolor{drawColor}{RGB}{0,0,0}

\node[text=drawColor,anchor=base,inner sep=0pt, outer sep=0pt, scale=  1.20] at (150.47,230.80) {\bfseries Distribució de les rendibilitats diàries};

\node[text=drawColor,rotate= 90.00,anchor=base,inner sep=0pt, outer sep=0pt, scale=  1.00] at (  9.60,139.77) {Densitat};
\end{scope}
\begin{scope}
\path[clip] (  0.00,  0.00) rectangle (252.94,252.94);
\definecolor{drawColor}{RGB}{0,0,0}

\path[draw=drawColor,line width= 0.4pt,line join=round,line cap=round] ( 55.59, 62.59) -- (245.35, 62.59);

\path[draw=drawColor,line width= 0.4pt,line join=round,line cap=round] ( 55.59, 62.59) -- ( 55.59, 56.59);

\path[draw=drawColor,line width= 0.4pt,line join=round,line cap=round] (103.03, 62.59) -- (103.03, 56.59);

\path[draw=drawColor,line width= 0.4pt,line join=round,line cap=round] (150.47, 62.59) -- (150.47, 56.59);

\path[draw=drawColor,line width= 0.4pt,line join=round,line cap=round] (197.91, 62.59) -- (197.91, 56.59);

\path[draw=drawColor,line width= 0.4pt,line join=round,line cap=round] (245.35, 62.59) -- (245.35, 56.59);

\node[text=drawColor,anchor=base,inner sep=0pt, outer sep=0pt, scale=  1.00] at ( 55.59, 40.99) {-0.10};

\node[text=drawColor,anchor=base,inner sep=0pt, outer sep=0pt, scale=  1.00] at (103.03, 40.99) {-0.05};

\node[text=drawColor,anchor=base,inner sep=0pt, outer sep=0pt, scale=  1.00] at (150.47, 40.99) {0.00};

\node[text=drawColor,anchor=base,inner sep=0pt, outer sep=0pt, scale=  1.00] at (197.91, 40.99) {0.05};

\node[text=drawColor,anchor=base,inner sep=0pt, outer sep=0pt, scale=  1.00] at (245.35, 40.99) {0.10};

\path[draw=drawColor,line width= 0.4pt,line join=round,line cap=round] ( 48.00, 68.31) -- ( 48.00,211.23);

\path[draw=drawColor,line width= 0.4pt,line join=round,line cap=round] ( 48.00, 68.31) -- ( 42.00, 68.31);

\path[draw=drawColor,line width= 0.4pt,line join=round,line cap=round] ( 48.00, 92.13) -- ( 42.00, 92.13);

\path[draw=drawColor,line width= 0.4pt,line join=round,line cap=round] ( 48.00,115.95) -- ( 42.00,115.95);

\path[draw=drawColor,line width= 0.4pt,line join=round,line cap=round] ( 48.00,139.77) -- ( 42.00,139.77);

\path[draw=drawColor,line width= 0.4pt,line join=round,line cap=round] ( 48.00,163.59) -- ( 42.00,163.59);

\path[draw=drawColor,line width= 0.4pt,line join=round,line cap=round] ( 48.00,187.41) -- ( 42.00,187.41);

\path[draw=drawColor,line width= 0.4pt,line join=round,line cap=round] ( 48.00,211.23) -- ( 42.00,211.23);

\node[text=drawColor,rotate= 90.00,anchor=base,inner sep=0pt, outer sep=0pt, scale=  1.00] at ( 33.60, 68.31) {0};

\node[text=drawColor,rotate= 90.00,anchor=base,inner sep=0pt, outer sep=0pt, scale=  1.00] at ( 33.60, 92.13) {5};

\node[text=drawColor,rotate= 90.00,anchor=base,inner sep=0pt, outer sep=0pt, scale=  1.00] at ( 33.60,115.95) {10};

\node[text=drawColor,rotate= 90.00,anchor=base,inner sep=0pt, outer sep=0pt, scale=  1.00] at ( 33.60,139.77) {15};

\node[text=drawColor,rotate= 90.00,anchor=base,inner sep=0pt, outer sep=0pt, scale=  1.00] at ( 33.60,163.59) {20};

\node[text=drawColor,rotate= 90.00,anchor=base,inner sep=0pt, outer sep=0pt, scale=  1.00] at ( 33.60,187.41) {25};

\node[text=drawColor,rotate= 90.00,anchor=base,inner sep=0pt, outer sep=0pt, scale=  1.00] at ( 33.60,211.23) {30};
\end{scope}
\begin{scope}
\path[clip] ( 48.00, 62.59) rectangle (252.94,216.94);
\definecolor{fillColor}{RGB}{173,216,230}

\path[fill=fillColor] ( 74.57, 68.31) rectangle ( 84.06, 69.73);

\path[fill=fillColor] ( 84.06, 68.31) rectangle ( 93.54, 69.73);

\path[fill=fillColor] ( 93.54, 68.31) rectangle (103.03, 70.68);

\path[fill=fillColor] (103.03, 68.31) rectangle (112.52, 71.15);

\path[fill=fillColor] (112.52, 68.31) rectangle (122.01, 77.80);

\path[fill=fillColor] (122.01, 68.31) rectangle (131.50, 92.98);

\path[fill=fillColor] (131.50, 68.31) rectangle (140.98,126.67);

\path[fill=fillColor] (140.98, 68.31) rectangle (150.47,202.59);

\path[fill=fillColor] (150.47, 68.31) rectangle (159.96,195.47);

\path[fill=fillColor] (159.96, 68.31) rectangle (169.45,135.21);

\path[fill=fillColor] (169.45, 68.31) rectangle (178.94, 99.62);

\path[fill=fillColor] (178.94, 68.31) rectangle (188.43, 76.85);

\path[fill=fillColor] (188.43, 68.31) rectangle (197.91, 73.05);

\path[fill=fillColor] (197.91, 68.31) rectangle (207.40, 69.73);

\path[fill=fillColor] (207.40, 68.31) rectangle (216.89, 68.78);

\path[fill=fillColor] (216.89, 68.31) rectangle (226.38, 69.25);
\definecolor{drawColor}{RGB}{218,112,214}

\path[draw=drawColor,line width= 1.2pt,dash pattern=on 1pt off 3pt ,line join=round,line cap=round] ( 55.59, 68.31) --
	( 57.49, 68.31) --
	( 59.39, 68.31) --
	( 61.28, 68.31) --
	( 63.18, 68.31) --
	( 65.08, 68.31) --
	( 66.98, 68.31) --
	( 68.87, 68.31) --
	( 70.77, 68.31) --
	( 72.67, 68.31) --
	( 74.57, 68.31) --
	( 76.46, 68.31) --
	( 78.36, 68.31) --
	( 80.26, 68.31) --
	( 82.16, 68.32) --
	( 84.06, 68.33) --
	( 85.95, 68.34) --
	( 87.85, 68.36) --
	( 89.75, 68.39) --
	( 91.65, 68.43) --
	( 93.54, 68.50) --
	( 95.44, 68.60) --
	( 97.34, 68.75) --
	( 99.24, 68.96) --
	(101.13, 69.25) --
	(103.03, 69.66) --
	(104.93, 70.22) --
	(106.83, 70.96) --
	(108.72, 71.95) --
	(110.62, 73.24) --
	(112.52, 74.89) --
	(114.42, 76.97) --
	(116.31, 79.55) --
	(118.21, 82.69) --
	(120.11, 86.46) --
	(122.01, 90.89) --
	(123.91, 96.02) --
	(125.80,101.85) --
	(127.70,108.33) --
	(129.60,115.40) --
	(131.50,122.96) --
	(133.39,130.85) --
	(135.29,138.88) --
	(137.19,146.85) --
	(139.09,154.50) --
	(140.98,161.58) --
	(142.88,167.85) --
	(144.78,173.06) --
	(146.68,177.02) --
	(148.57,179.56) --
	(150.47,180.58) --
	(152.37,180.04) --
	(154.27,177.97) --
	(156.17,174.43) --
	(158.06,169.59) --
	(159.96,163.63) --
	(161.86,156.78) --
	(163.76,149.28) --
	(165.65,141.38) --
	(167.55,133.35) --
	(169.45,125.39) --
	(171.35,117.71) --
	(173.24,110.47) --
	(175.14,103.80) --
	(177.04, 97.76) --
	(178.94, 92.42) --
	(180.83, 87.77) --
	(182.73, 83.80) --
	(184.63, 80.46) --
	(186.53, 77.72) --
	(188.43, 75.49) --
	(190.32, 73.71) --
	(192.22, 72.32) --
	(194.12, 71.24) --
	(196.02, 70.43) --
	(197.91, 69.82) --
	(199.81, 69.37) --
	(201.71, 69.04) --
	(203.61, 68.81) --
	(205.50, 68.64) --
	(207.40, 68.53) --
	(209.30, 68.45) --
	(211.20, 68.40) --
	(213.09, 68.37) --
	(214.99, 68.34) --
	(216.89, 68.33) --
	(218.79, 68.32) --
	(220.69, 68.31) --
	(222.58, 68.31) --
	(224.48, 68.31) --
	(226.38, 68.31) --
	(228.28, 68.31) --
	(230.17, 68.31) --
	(232.07, 68.31) --
	(233.97, 68.31) --
	(235.87, 68.31) --
	(237.76, 68.31) --
	(239.66, 68.31) --
	(241.56, 68.31) --
	(243.46, 68.31) --
	(245.35, 68.31);
\definecolor{drawColor}{RGB}{50,205,50}

\path[draw=drawColor,line width= 1.2pt,dash pattern=on 1pt off 3pt ,line join=round,line cap=round] ( 55.59, 68.35) --
	( 57.49, 68.36) --
	( 59.39, 68.37) --
	( 61.28, 68.38) --
	( 63.18, 68.40) --
	( 65.08, 68.41) --
	( 66.98, 68.43) --
	( 68.87, 68.45) --
	( 70.77, 68.47) --
	( 72.67, 68.50) --
	( 74.57, 68.53) --
	( 76.46, 68.57) --
	( 78.36, 68.62) --
	( 80.26, 68.67) --
	( 82.16, 68.73) --
	( 84.06, 68.81) --
	( 85.95, 68.89) --
	( 87.85, 69.00) --
	( 89.75, 69.12) --
	( 91.65, 69.27) --
	( 93.54, 69.44) --
	( 95.44, 69.64) --
	( 97.34, 69.89) --
	( 99.24, 70.18) --
	(101.13, 70.52) --
	(103.03, 70.94) --
	(104.93, 71.43) --
	(106.83, 72.03) --
	(108.72, 72.74) --
	(110.62, 73.60) --
	(112.52, 74.63) --
	(114.42, 75.87) --
	(116.31, 77.38) --
	(118.21, 79.20) --
	(120.11, 81.40) --
	(122.01, 84.07) --
	(123.91, 87.30) --
	(125.80, 91.22) --
	(127.70, 95.95) --
	(129.60,101.66) --
	(131.50,108.50) --
	(133.39,116.63) --
	(135.29,126.19) --
	(137.19,137.22) --
	(139.09,149.64) --
	(140.98,163.12) --
	(142.88,177.01) --
	(144.78,190.30) --
	(146.68,201.64) --
	(148.57,209.58) --
	(150.47,212.94) --
	(152.37,211.16) --
	(154.27,204.54) --
	(156.17,194.11) --
	(158.06,181.27) --
	(159.96,167.44) --
	(161.86,153.75) --
	(163.76,140.95) --
	(165.65,129.47) --
	(167.55,119.45) --
	(169.45,110.88) --
	(171.35,103.66) --
	(173.24, 97.62) --
	(175.14, 92.60) --
	(177.04, 88.44) --
	(178.94, 85.01) --
	(180.83, 82.17) --
	(182.73, 79.84) --
	(184.63, 77.91) --
	(186.53, 76.31) --
	(188.43, 74.99) --
	(190.32, 73.90) --
	(192.22, 72.99) --
	(194.12, 72.23) --
	(196.02, 71.61) --
	(197.91, 71.08) --
	(199.81, 70.64) --
	(201.71, 70.28) --
	(203.61, 69.97) --
	(205.50, 69.71) --
	(207.40, 69.50) --
	(209.30, 69.32) --
	(211.20, 69.16) --
	(213.09, 69.03) --
	(214.99, 68.93) --
	(216.89, 68.83) --
	(218.79, 68.75) --
	(220.69, 68.69) --
	(222.58, 68.63) --
	(224.48, 68.58) --
	(226.38, 68.54) --
	(228.28, 68.51) --
	(230.17, 68.48) --
	(232.07, 68.46) --
	(233.97, 68.43) --
	(235.87, 68.42) --
	(237.76, 68.40) --
	(239.66, 68.39) --
	(241.56, 68.38) --
	(243.46, 68.37) --
	(245.35, 68.36);

\path[] ( 54.15,212.31) rectangle (127.17,176.31);
\definecolor{drawColor}{RGB}{218,112,214}

\path[draw=drawColor,line width= 1.6pt,dash pattern=on 1pt off 3pt ,line join=round,line cap=round] ( 63.15,200.31) -- ( 81.15,200.31);
\definecolor{drawColor}{RGB}{50,205,50}

\path[draw=drawColor,line width= 1.6pt,dash pattern=on 1pt off 3pt ,line join=round,line cap=round] ( 63.15,188.31) -- ( 81.15,188.31);
\definecolor{drawColor}{RGB}{0,0,0}

\node[text=drawColor,anchor=base west,inner sep=0pt, outer sep=0pt, scale=  1.00] at ( 90.15,196.87) {Normal};

\node[text=drawColor,anchor=base west,inner sep=0pt, outer sep=0pt, scale=  1.00] at ( 90.15,184.87) {NIG};
\end{scope}
\end{tikzpicture}

	\caption{Evolució del preu de tancament de Caterpillar entre 2015 i 2018}
	\label{fig:distribucio rendibilitats}
\end{figure}
\begin{figure}[htb]
	\centering \sffamily \small
	% Created by tikzDevice version 0.12 on 2019-06-13 19:48:56
% !TEX encoding = UTF-8 Unicode
\begin{tikzpicture}[x=1pt,y=1pt]
\definecolor{fillColor}{RGB}{255,255,255}
\path[use as bounding box,fill=fillColor,fill opacity=0.00] (0,0) rectangle (361.35,180.67);
\begin{scope}
\path[clip] ( 48.00, 48.00) rectangle (174.67,156.67);
\definecolor{drawColor}{gray}{0.40}

\path[draw=drawColor,line width= 0.4pt,line join=round,line cap=round] ( 50.44, 49.77) -- ( 54.94, 54.27);

\path[draw=drawColor,line width= 0.4pt,line join=round,line cap=round] ( 50.44, 54.27) -- ( 54.94, 49.77);

\path[draw=drawColor,line width= 0.4pt,line join=round,line cap=round] ( 50.56, 49.80) -- ( 55.06, 54.30);

\path[draw=drawColor,line width= 0.4pt,line join=round,line cap=round] ( 50.56, 54.30) -- ( 55.06, 49.80);

\path[draw=drawColor,line width= 0.4pt,line join=round,line cap=round] ( 50.68, 49.82) -- ( 55.18, 54.32);

\path[draw=drawColor,line width= 0.4pt,line join=round,line cap=round] ( 50.68, 54.32) -- ( 55.18, 49.82);

\path[draw=drawColor,line width= 0.4pt,line join=round,line cap=round] ( 50.79, 49.86) -- ( 55.29, 54.36);

\path[draw=drawColor,line width= 0.4pt,line join=round,line cap=round] ( 50.79, 54.36) -- ( 55.29, 49.86);

\path[draw=drawColor,line width= 0.4pt,line join=round,line cap=round] ( 50.91, 49.91) -- ( 55.41, 54.41);

\path[draw=drawColor,line width= 0.4pt,line join=round,line cap=round] ( 50.91, 54.41) -- ( 55.41, 49.91);

\path[draw=drawColor,line width= 0.4pt,line join=round,line cap=round] ( 51.03, 49.92) -- ( 55.53, 54.42);

\path[draw=drawColor,line width= 0.4pt,line join=round,line cap=round] ( 51.03, 54.42) -- ( 55.53, 49.92);

\path[draw=drawColor,line width= 0.4pt,line join=round,line cap=round] ( 51.14, 50.03) -- ( 55.64, 54.53);

\path[draw=drawColor,line width= 0.4pt,line join=round,line cap=round] ( 51.14, 54.53) -- ( 55.64, 50.03);

\path[draw=drawColor,line width= 0.4pt,line join=round,line cap=round] ( 51.26, 50.07) -- ( 55.76, 54.57);

\path[draw=drawColor,line width= 0.4pt,line join=round,line cap=round] ( 51.26, 54.57) -- ( 55.76, 50.07);

\path[draw=drawColor,line width= 0.4pt,line join=round,line cap=round] ( 51.38, 50.13) -- ( 55.88, 54.63);

\path[draw=drawColor,line width= 0.4pt,line join=round,line cap=round] ( 51.38, 54.63) -- ( 55.88, 50.13);

\path[draw=drawColor,line width= 0.4pt,line join=round,line cap=round] ( 51.49, 50.18) -- ( 55.99, 54.68);

\path[draw=drawColor,line width= 0.4pt,line join=round,line cap=round] ( 51.49, 54.68) -- ( 55.99, 50.18);

\path[draw=drawColor,line width= 0.4pt,line join=round,line cap=round] ( 51.61, 50.29) -- ( 56.11, 54.79);

\path[draw=drawColor,line width= 0.4pt,line join=round,line cap=round] ( 51.61, 54.79) -- ( 56.11, 50.29);

\path[draw=drawColor,line width= 0.4pt,line join=round,line cap=round] ( 51.73, 50.63) -- ( 56.23, 55.13);

\path[draw=drawColor,line width= 0.4pt,line join=round,line cap=round] ( 51.73, 55.13) -- ( 56.23, 50.63);

\path[draw=drawColor,line width= 0.4pt,line join=round,line cap=round] ( 51.84, 50.84) -- ( 56.34, 55.34);

\path[draw=drawColor,line width= 0.4pt,line join=round,line cap=round] ( 51.84, 55.34) -- ( 56.34, 50.84);

\path[draw=drawColor,line width= 0.4pt,line join=round,line cap=round] ( 51.96, 50.96) -- ( 56.46, 55.46);

\path[draw=drawColor,line width= 0.4pt,line join=round,line cap=round] ( 51.96, 55.46) -- ( 56.46, 50.96);

\path[draw=drawColor,line width= 0.4pt,line join=round,line cap=round] ( 52.08, 51.02) -- ( 56.58, 55.52);

\path[draw=drawColor,line width= 0.4pt,line join=round,line cap=round] ( 52.08, 55.52) -- ( 56.58, 51.02);

\path[draw=drawColor,line width= 0.4pt,line join=round,line cap=round] ( 52.20, 51.07) -- ( 56.70, 55.57);

\path[draw=drawColor,line width= 0.4pt,line join=round,line cap=round] ( 52.20, 55.57) -- ( 56.70, 51.07);

\path[draw=drawColor,line width= 0.4pt,line join=round,line cap=round] ( 52.31, 51.08) -- ( 56.81, 55.58);

\path[draw=drawColor,line width= 0.4pt,line join=round,line cap=round] ( 52.31, 55.58) -- ( 56.81, 51.08);

\path[draw=drawColor,line width= 0.4pt,line join=round,line cap=round] ( 52.43, 51.27) -- ( 56.93, 55.77);

\path[draw=drawColor,line width= 0.4pt,line join=round,line cap=round] ( 52.43, 55.77) -- ( 56.93, 51.27);

\path[draw=drawColor,line width= 0.4pt,line join=round,line cap=round] ( 52.55, 51.36) -- ( 57.05, 55.86);

\path[draw=drawColor,line width= 0.4pt,line join=round,line cap=round] ( 52.55, 55.86) -- ( 57.05, 51.36);

\path[draw=drawColor,line width= 0.4pt,line join=round,line cap=round] ( 52.66, 51.38) -- ( 57.16, 55.88);

\path[draw=drawColor,line width= 0.4pt,line join=round,line cap=round] ( 52.66, 55.88) -- ( 57.16, 51.38);

\path[draw=drawColor,line width= 0.4pt,line join=round,line cap=round] ( 52.78, 51.52) -- ( 57.28, 56.02);

\path[draw=drawColor,line width= 0.4pt,line join=round,line cap=round] ( 52.78, 56.02) -- ( 57.28, 51.52);

\path[draw=drawColor,line width= 0.4pt,line join=round,line cap=round] ( 52.90, 51.68) -- ( 57.40, 56.18);

\path[draw=drawColor,line width= 0.4pt,line join=round,line cap=round] ( 52.90, 56.18) -- ( 57.40, 51.68);

\path[draw=drawColor,line width= 0.4pt,line join=round,line cap=round] ( 53.01, 51.75) -- ( 57.51, 56.25);

\path[draw=drawColor,line width= 0.4pt,line join=round,line cap=round] ( 53.01, 56.25) -- ( 57.51, 51.75);

\path[draw=drawColor,line width= 0.4pt,line join=round,line cap=round] ( 53.13, 51.77) -- ( 57.63, 56.27);

\path[draw=drawColor,line width= 0.4pt,line join=round,line cap=round] ( 53.13, 56.27) -- ( 57.63, 51.77);

\path[draw=drawColor,line width= 0.4pt,line join=round,line cap=round] ( 53.25, 51.95) -- ( 57.75, 56.45);

\path[draw=drawColor,line width= 0.4pt,line join=round,line cap=round] ( 53.25, 56.45) -- ( 57.75, 51.95);

\path[draw=drawColor,line width= 0.4pt,line join=round,line cap=round] ( 53.37, 52.05) -- ( 57.87, 56.55);

\path[draw=drawColor,line width= 0.4pt,line join=round,line cap=round] ( 53.37, 56.55) -- ( 57.87, 52.05);

\path[draw=drawColor,line width= 0.4pt,line join=round,line cap=round] ( 53.48, 52.10) -- ( 57.98, 56.60);

\path[draw=drawColor,line width= 0.4pt,line join=round,line cap=round] ( 53.48, 56.60) -- ( 57.98, 52.10);

\path[draw=drawColor,line width= 0.4pt,line join=round,line cap=round] ( 53.60, 52.18) -- ( 58.10, 56.68);

\path[draw=drawColor,line width= 0.4pt,line join=round,line cap=round] ( 53.60, 56.68) -- ( 58.10, 52.18);

\path[draw=drawColor,line width= 0.4pt,line join=round,line cap=round] ( 53.72, 52.35) -- ( 58.22, 56.85);

\path[draw=drawColor,line width= 0.4pt,line join=round,line cap=round] ( 53.72, 56.85) -- ( 58.22, 52.35);

\path[draw=drawColor,line width= 0.4pt,line join=round,line cap=round] ( 53.83, 52.44) -- ( 58.33, 56.94);

\path[draw=drawColor,line width= 0.4pt,line join=round,line cap=round] ( 53.83, 56.94) -- ( 58.33, 52.44);

\path[draw=drawColor,line width= 0.4pt,line join=round,line cap=round] ( 53.95, 52.46) -- ( 58.45, 56.96);

\path[draw=drawColor,line width= 0.4pt,line join=round,line cap=round] ( 53.95, 56.96) -- ( 58.45, 52.46);

\path[draw=drawColor,line width= 0.4pt,line join=round,line cap=round] ( 54.07, 52.60) -- ( 58.57, 57.10);

\path[draw=drawColor,line width= 0.4pt,line join=round,line cap=round] ( 54.07, 57.10) -- ( 58.57, 52.60);

\path[draw=drawColor,line width= 0.4pt,line join=round,line cap=round] ( 54.18, 52.76) -- ( 58.68, 57.26);

\path[draw=drawColor,line width= 0.4pt,line join=round,line cap=round] ( 54.18, 57.26) -- ( 58.68, 52.76);

\path[draw=drawColor,line width= 0.4pt,line join=round,line cap=round] ( 54.30, 52.94) -- ( 58.80, 57.44);

\path[draw=drawColor,line width= 0.4pt,line join=round,line cap=round] ( 54.30, 57.44) -- ( 58.80, 52.94);

\path[draw=drawColor,line width= 0.4pt,line join=round,line cap=round] ( 54.42, 53.10) -- ( 58.92, 57.60);

\path[draw=drawColor,line width= 0.4pt,line join=round,line cap=round] ( 54.42, 57.60) -- ( 58.92, 53.10);

\path[draw=drawColor,line width= 0.4pt,line join=round,line cap=round] ( 54.53, 53.32) -- ( 59.03, 57.82);

\path[draw=drawColor,line width= 0.4pt,line join=round,line cap=round] ( 54.53, 57.82) -- ( 59.03, 53.32);

\path[draw=drawColor,line width= 0.4pt,line join=round,line cap=round] ( 54.65, 53.40) -- ( 59.15, 57.90);

\path[draw=drawColor,line width= 0.4pt,line join=round,line cap=round] ( 54.65, 57.90) -- ( 59.15, 53.40);

\path[draw=drawColor,line width= 0.4pt,line join=round,line cap=round] ( 54.77, 53.48) -- ( 59.27, 57.98);

\path[draw=drawColor,line width= 0.4pt,line join=round,line cap=round] ( 54.77, 57.98) -- ( 59.27, 53.48);

\path[draw=drawColor,line width= 0.4pt,line join=round,line cap=round] ( 54.89, 53.54) -- ( 59.39, 58.04);

\path[draw=drawColor,line width= 0.4pt,line join=round,line cap=round] ( 54.89, 58.04) -- ( 59.39, 53.54);

\path[draw=drawColor,line width= 0.4pt,line join=round,line cap=round] ( 55.00, 53.62) -- ( 59.50, 58.12);

\path[draw=drawColor,line width= 0.4pt,line join=round,line cap=round] ( 55.00, 58.12) -- ( 59.50, 53.62);

\path[draw=drawColor,line width= 0.4pt,line join=round,line cap=round] ( 55.12, 53.73) -- ( 59.62, 58.23);

\path[draw=drawColor,line width= 0.4pt,line join=round,line cap=round] ( 55.12, 58.23) -- ( 59.62, 53.73);

\path[draw=drawColor,line width= 0.4pt,line join=round,line cap=round] ( 55.24, 53.85) -- ( 59.74, 58.35);

\path[draw=drawColor,line width= 0.4pt,line join=round,line cap=round] ( 55.24, 58.35) -- ( 59.74, 53.85);

\path[draw=drawColor,line width= 0.4pt,line join=round,line cap=round] ( 55.35, 53.97) -- ( 59.85, 58.47);

\path[draw=drawColor,line width= 0.4pt,line join=round,line cap=round] ( 55.35, 58.47) -- ( 59.85, 53.97);

\path[draw=drawColor,line width= 0.4pt,line join=round,line cap=round] ( 55.47, 54.02) -- ( 59.97, 58.52);

\path[draw=drawColor,line width= 0.4pt,line join=round,line cap=round] ( 55.47, 58.52) -- ( 59.97, 54.02);

\path[draw=drawColor,line width= 0.4pt,line join=round,line cap=round] ( 55.59, 54.12) -- ( 60.09, 58.62);

\path[draw=drawColor,line width= 0.4pt,line join=round,line cap=round] ( 55.59, 58.62) -- ( 60.09, 54.12);

\path[draw=drawColor,line width= 0.4pt,line join=round,line cap=round] ( 55.70, 54.21) -- ( 60.20, 58.71);

\path[draw=drawColor,line width= 0.4pt,line join=round,line cap=round] ( 55.70, 58.71) -- ( 60.20, 54.21);

\path[draw=drawColor,line width= 0.4pt,line join=round,line cap=round] ( 55.82, 54.24) -- ( 60.32, 58.74);

\path[draw=drawColor,line width= 0.4pt,line join=round,line cap=round] ( 55.82, 58.74) -- ( 60.32, 54.24);

\path[draw=drawColor,line width= 0.4pt,line join=round,line cap=round] ( 55.94, 54.41) -- ( 60.44, 58.91);

\path[draw=drawColor,line width= 0.4pt,line join=round,line cap=round] ( 55.94, 58.91) -- ( 60.44, 54.41);

\path[draw=drawColor,line width= 0.4pt,line join=round,line cap=round] ( 56.05, 54.45) -- ( 60.55, 58.95);

\path[draw=drawColor,line width= 0.4pt,line join=round,line cap=round] ( 56.05, 58.95) -- ( 60.55, 54.45);

\path[draw=drawColor,line width= 0.4pt,line join=round,line cap=round] ( 56.17, 54.52) -- ( 60.67, 59.02);

\path[draw=drawColor,line width= 0.4pt,line join=round,line cap=round] ( 56.17, 59.02) -- ( 60.67, 54.52);

\path[draw=drawColor,line width= 0.4pt,line join=round,line cap=round] ( 56.29, 54.61) -- ( 60.79, 59.11);

\path[draw=drawColor,line width= 0.4pt,line join=round,line cap=round] ( 56.29, 59.11) -- ( 60.79, 54.61);

\path[draw=drawColor,line width= 0.4pt,line join=round,line cap=round] ( 56.41, 54.69) -- ( 60.91, 59.19);

\path[draw=drawColor,line width= 0.4pt,line join=round,line cap=round] ( 56.41, 59.19) -- ( 60.91, 54.69);

\path[draw=drawColor,line width= 0.4pt,line join=round,line cap=round] ( 56.52, 54.71) -- ( 61.02, 59.21);

\path[draw=drawColor,line width= 0.4pt,line join=round,line cap=round] ( 56.52, 59.21) -- ( 61.02, 54.71);

\path[draw=drawColor,line width= 0.4pt,line join=round,line cap=round] ( 56.64, 54.91) -- ( 61.14, 59.41);

\path[draw=drawColor,line width= 0.4pt,line join=round,line cap=round] ( 56.64, 59.41) -- ( 61.14, 54.91);

\path[draw=drawColor,line width= 0.4pt,line join=round,line cap=round] ( 56.76, 55.08) -- ( 61.26, 59.58);

\path[draw=drawColor,line width= 0.4pt,line join=round,line cap=round] ( 56.76, 59.58) -- ( 61.26, 55.08);

\path[draw=drawColor,line width= 0.4pt,line join=round,line cap=round] ( 56.87, 55.21) -- ( 61.37, 59.71);

\path[draw=drawColor,line width= 0.4pt,line join=round,line cap=round] ( 56.87, 59.71) -- ( 61.37, 55.21);

\path[draw=drawColor,line width= 0.4pt,line join=round,line cap=round] ( 56.99, 55.33) -- ( 61.49, 59.83);

\path[draw=drawColor,line width= 0.4pt,line join=round,line cap=round] ( 56.99, 59.83) -- ( 61.49, 55.33);

\path[draw=drawColor,line width= 0.4pt,line join=round,line cap=round] ( 57.11, 55.40) -- ( 61.61, 59.90);

\path[draw=drawColor,line width= 0.4pt,line join=round,line cap=round] ( 57.11, 59.90) -- ( 61.61, 55.40);

\path[draw=drawColor,line width= 0.4pt,line join=round,line cap=round] ( 57.22, 55.54) -- ( 61.72, 60.04);

\path[draw=drawColor,line width= 0.4pt,line join=round,line cap=round] ( 57.22, 60.04) -- ( 61.72, 55.54);

\path[draw=drawColor,line width= 0.4pt,line join=round,line cap=round] ( 57.34, 55.96) -- ( 61.84, 60.46);

\path[draw=drawColor,line width= 0.4pt,line join=round,line cap=round] ( 57.34, 60.46) -- ( 61.84, 55.96);

\path[draw=drawColor,line width= 0.4pt,line join=round,line cap=round] ( 57.46, 56.11) -- ( 61.96, 60.61);

\path[draw=drawColor,line width= 0.4pt,line join=round,line cap=round] ( 57.46, 60.61) -- ( 61.96, 56.11);

\path[draw=drawColor,line width= 0.4pt,line join=round,line cap=round] ( 57.58, 56.14) -- ( 62.08, 60.64);

\path[draw=drawColor,line width= 0.4pt,line join=round,line cap=round] ( 57.58, 60.64) -- ( 62.08, 56.14);

\path[draw=drawColor,line width= 0.4pt,line join=round,line cap=round] ( 57.69, 56.15) -- ( 62.19, 60.65);

\path[draw=drawColor,line width= 0.4pt,line join=round,line cap=round] ( 57.69, 60.65) -- ( 62.19, 56.15);

\path[draw=drawColor,line width= 0.4pt,line join=round,line cap=round] ( 57.81, 56.28) -- ( 62.31, 60.78);

\path[draw=drawColor,line width= 0.4pt,line join=round,line cap=round] ( 57.81, 60.78) -- ( 62.31, 56.28);

\path[draw=drawColor,line width= 0.4pt,line join=round,line cap=round] ( 57.93, 56.39) -- ( 62.43, 60.89);

\path[draw=drawColor,line width= 0.4pt,line join=round,line cap=round] ( 57.93, 60.89) -- ( 62.43, 56.39);

\path[draw=drawColor,line width= 0.4pt,line join=round,line cap=round] ( 58.04, 56.48) -- ( 62.54, 60.98);

\path[draw=drawColor,line width= 0.4pt,line join=round,line cap=round] ( 58.04, 60.98) -- ( 62.54, 56.48);

\path[draw=drawColor,line width= 0.4pt,line join=round,line cap=round] ( 58.16, 56.53) -- ( 62.66, 61.03);

\path[draw=drawColor,line width= 0.4pt,line join=round,line cap=round] ( 58.16, 61.03) -- ( 62.66, 56.53);

\path[draw=drawColor,line width= 0.4pt,line join=round,line cap=round] ( 58.28, 56.53) -- ( 62.78, 61.03);

\path[draw=drawColor,line width= 0.4pt,line join=round,line cap=round] ( 58.28, 61.03) -- ( 62.78, 56.53);

\path[draw=drawColor,line width= 0.4pt,line join=round,line cap=round] ( 58.39, 56.68) -- ( 62.89, 61.18);

\path[draw=drawColor,line width= 0.4pt,line join=round,line cap=round] ( 58.39, 61.18) -- ( 62.89, 56.68);

\path[draw=drawColor,line width= 0.4pt,line join=round,line cap=round] ( 58.51, 56.75) -- ( 63.01, 61.25);

\path[draw=drawColor,line width= 0.4pt,line join=round,line cap=round] ( 58.51, 61.25) -- ( 63.01, 56.75);

\path[draw=drawColor,line width= 0.4pt,line join=round,line cap=round] ( 58.63, 56.75) -- ( 63.13, 61.25);

\path[draw=drawColor,line width= 0.4pt,line join=round,line cap=round] ( 58.63, 61.25) -- ( 63.13, 56.75);

\path[draw=drawColor,line width= 0.4pt,line join=round,line cap=round] ( 58.74, 56.84) -- ( 63.24, 61.34);

\path[draw=drawColor,line width= 0.4pt,line join=round,line cap=round] ( 58.74, 61.34) -- ( 63.24, 56.84);

\path[draw=drawColor,line width= 0.4pt,line join=round,line cap=round] ( 58.86, 57.09) -- ( 63.36, 61.59);

\path[draw=drawColor,line width= 0.4pt,line join=round,line cap=round] ( 58.86, 61.59) -- ( 63.36, 57.09);

\path[draw=drawColor,line width= 0.4pt,line join=round,line cap=round] ( 58.98, 57.28) -- ( 63.48, 61.78);

\path[draw=drawColor,line width= 0.4pt,line join=round,line cap=round] ( 58.98, 61.78) -- ( 63.48, 57.28);

\path[draw=drawColor,line width= 0.4pt,line join=round,line cap=round] ( 59.10, 57.44) -- ( 63.60, 61.94);

\path[draw=drawColor,line width= 0.4pt,line join=round,line cap=round] ( 59.10, 61.94) -- ( 63.60, 57.44);

\path[draw=drawColor,line width= 0.4pt,line join=round,line cap=round] ( 59.21, 57.47) -- ( 63.71, 61.97);

\path[draw=drawColor,line width= 0.4pt,line join=round,line cap=round] ( 59.21, 61.97) -- ( 63.71, 57.47);

\path[draw=drawColor,line width= 0.4pt,line join=round,line cap=round] ( 59.33, 57.61) -- ( 63.83, 62.11);

\path[draw=drawColor,line width= 0.4pt,line join=round,line cap=round] ( 59.33, 62.11) -- ( 63.83, 57.61);

\path[draw=drawColor,line width= 0.4pt,line join=round,line cap=round] ( 59.45, 57.61) -- ( 63.95, 62.11);

\path[draw=drawColor,line width= 0.4pt,line join=round,line cap=round] ( 59.45, 62.11) -- ( 63.95, 57.61);

\path[draw=drawColor,line width= 0.4pt,line join=round,line cap=round] ( 59.56, 57.85) -- ( 64.06, 62.35);

\path[draw=drawColor,line width= 0.4pt,line join=round,line cap=round] ( 59.56, 62.35) -- ( 64.06, 57.85);

\path[draw=drawColor,line width= 0.4pt,line join=round,line cap=round] ( 59.68, 58.07) -- ( 64.18, 62.57);

\path[draw=drawColor,line width= 0.4pt,line join=round,line cap=round] ( 59.68, 62.57) -- ( 64.18, 58.07);

\path[draw=drawColor,line width= 0.4pt,line join=round,line cap=round] ( 59.80, 58.26) -- ( 64.30, 62.76);

\path[draw=drawColor,line width= 0.4pt,line join=round,line cap=round] ( 59.80, 62.76) -- ( 64.30, 58.26);

\path[draw=drawColor,line width= 0.4pt,line join=round,line cap=round] ( 59.91, 58.28) -- ( 64.41, 62.78);

\path[draw=drawColor,line width= 0.4pt,line join=round,line cap=round] ( 59.91, 62.78) -- ( 64.41, 58.28);

\path[draw=drawColor,line width= 0.4pt,line join=round,line cap=round] ( 60.03, 58.28) -- ( 64.53, 62.78);

\path[draw=drawColor,line width= 0.4pt,line join=round,line cap=round] ( 60.03, 62.78) -- ( 64.53, 58.28);

\path[draw=drawColor,line width= 0.4pt,line join=round,line cap=round] ( 60.15, 58.40) -- ( 64.65, 62.90);

\path[draw=drawColor,line width= 0.4pt,line join=round,line cap=round] ( 60.15, 62.90) -- ( 64.65, 58.40);

\path[draw=drawColor,line width= 0.4pt,line join=round,line cap=round] ( 60.26, 58.49) -- ( 64.76, 62.99);

\path[draw=drawColor,line width= 0.4pt,line join=round,line cap=round] ( 60.26, 62.99) -- ( 64.76, 58.49);

\path[draw=drawColor,line width= 0.4pt,line join=round,line cap=round] ( 60.38, 58.51) -- ( 64.88, 63.01);

\path[draw=drawColor,line width= 0.4pt,line join=round,line cap=round] ( 60.38, 63.01) -- ( 64.88, 58.51);

\path[draw=drawColor,line width= 0.4pt,line join=round,line cap=round] ( 60.50, 58.52) -- ( 65.00, 63.02);

\path[draw=drawColor,line width= 0.4pt,line join=round,line cap=round] ( 60.50, 63.02) -- ( 65.00, 58.52);

\path[draw=drawColor,line width= 0.4pt,line join=round,line cap=round] ( 60.62, 58.67) -- ( 65.12, 63.17);

\path[draw=drawColor,line width= 0.4pt,line join=round,line cap=round] ( 60.62, 63.17) -- ( 65.12, 58.67);

\path[draw=drawColor,line width= 0.4pt,line join=round,line cap=round] ( 60.73, 58.90) -- ( 65.23, 63.40);

\path[draw=drawColor,line width= 0.4pt,line join=round,line cap=round] ( 60.73, 63.40) -- ( 65.23, 58.90);

\path[draw=drawColor,line width= 0.4pt,line join=round,line cap=round] ( 60.85, 59.02) -- ( 65.35, 63.52);

\path[draw=drawColor,line width= 0.4pt,line join=round,line cap=round] ( 60.85, 63.52) -- ( 65.35, 59.02);

\path[draw=drawColor,line width= 0.4pt,line join=round,line cap=round] ( 60.97, 59.03) -- ( 65.47, 63.53);

\path[draw=drawColor,line width= 0.4pt,line join=round,line cap=round] ( 60.97, 63.53) -- ( 65.47, 59.03);

\path[draw=drawColor,line width= 0.4pt,line join=round,line cap=round] ( 61.08, 59.05) -- ( 65.58, 63.55);

\path[draw=drawColor,line width= 0.4pt,line join=round,line cap=round] ( 61.08, 63.55) -- ( 65.58, 59.05);

\path[draw=drawColor,line width= 0.4pt,line join=round,line cap=round] ( 61.20, 59.17) -- ( 65.70, 63.67);

\path[draw=drawColor,line width= 0.4pt,line join=round,line cap=round] ( 61.20, 63.67) -- ( 65.70, 59.17);

\path[draw=drawColor,line width= 0.4pt,line join=round,line cap=round] ( 61.32, 59.20) -- ( 65.82, 63.70);

\path[draw=drawColor,line width= 0.4pt,line join=round,line cap=round] ( 61.32, 63.70) -- ( 65.82, 59.20);

\path[draw=drawColor,line width= 0.4pt,line join=round,line cap=round] ( 61.43, 59.21) -- ( 65.93, 63.71);

\path[draw=drawColor,line width= 0.4pt,line join=round,line cap=round] ( 61.43, 63.71) -- ( 65.93, 59.21);

\path[draw=drawColor,line width= 0.4pt,line join=round,line cap=round] ( 61.55, 59.57) -- ( 66.05, 64.07);

\path[draw=drawColor,line width= 0.4pt,line join=round,line cap=round] ( 61.55, 64.07) -- ( 66.05, 59.57);

\path[draw=drawColor,line width= 0.4pt,line join=round,line cap=round] ( 61.67, 59.68) -- ( 66.17, 64.18);

\path[draw=drawColor,line width= 0.4pt,line join=round,line cap=round] ( 61.67, 64.18) -- ( 66.17, 59.68);

\path[draw=drawColor,line width= 0.4pt,line join=round,line cap=round] ( 61.78, 59.68) -- ( 66.28, 64.18);

\path[draw=drawColor,line width= 0.4pt,line join=round,line cap=round] ( 61.78, 64.18) -- ( 66.28, 59.68);

\path[draw=drawColor,line width= 0.4pt,line join=round,line cap=round] ( 61.90, 59.72) -- ( 66.40, 64.22);

\path[draw=drawColor,line width= 0.4pt,line join=round,line cap=round] ( 61.90, 64.22) -- ( 66.40, 59.72);

\path[draw=drawColor,line width= 0.4pt,line join=round,line cap=round] ( 62.02, 59.77) -- ( 66.52, 64.27);

\path[draw=drawColor,line width= 0.4pt,line join=round,line cap=round] ( 62.02, 64.27) -- ( 66.52, 59.77);

\path[draw=drawColor,line width= 0.4pt,line join=round,line cap=round] ( 62.14, 59.91) -- ( 66.64, 64.41);

\path[draw=drawColor,line width= 0.4pt,line join=round,line cap=round] ( 62.14, 64.41) -- ( 66.64, 59.91);

\path[draw=drawColor,line width= 0.4pt,line join=round,line cap=round] ( 62.25, 59.92) -- ( 66.75, 64.42);

\path[draw=drawColor,line width= 0.4pt,line join=round,line cap=round] ( 62.25, 64.42) -- ( 66.75, 59.92);

\path[draw=drawColor,line width= 0.4pt,line join=round,line cap=round] ( 62.37, 60.02) -- ( 66.87, 64.52);

\path[draw=drawColor,line width= 0.4pt,line join=round,line cap=round] ( 62.37, 64.52) -- ( 66.87, 60.02);

\path[draw=drawColor,line width= 0.4pt,line join=round,line cap=round] ( 62.49, 60.12) -- ( 66.99, 64.62);

\path[draw=drawColor,line width= 0.4pt,line join=round,line cap=round] ( 62.49, 64.62) -- ( 66.99, 60.12);

\path[draw=drawColor,line width= 0.4pt,line join=round,line cap=round] ( 62.60, 60.21) -- ( 67.10, 64.71);

\path[draw=drawColor,line width= 0.4pt,line join=round,line cap=round] ( 62.60, 64.71) -- ( 67.10, 60.21);

\path[draw=drawColor,line width= 0.4pt,line join=round,line cap=round] ( 62.72, 60.35) -- ( 67.22, 64.85);

\path[draw=drawColor,line width= 0.4pt,line join=round,line cap=round] ( 62.72, 64.85) -- ( 67.22, 60.35);

\path[draw=drawColor,line width= 0.4pt,line join=round,line cap=round] ( 62.84, 60.52) -- ( 67.34, 65.02);

\path[draw=drawColor,line width= 0.4pt,line join=round,line cap=round] ( 62.84, 65.02) -- ( 67.34, 60.52);

\path[draw=drawColor,line width= 0.4pt,line join=round,line cap=round] ( 62.95, 60.53) -- ( 67.45, 65.03);

\path[draw=drawColor,line width= 0.4pt,line join=round,line cap=round] ( 62.95, 65.03) -- ( 67.45, 60.53);

\path[draw=drawColor,line width= 0.4pt,line join=round,line cap=round] ( 63.07, 60.62) -- ( 67.57, 65.12);

\path[draw=drawColor,line width= 0.4pt,line join=round,line cap=round] ( 63.07, 65.12) -- ( 67.57, 60.62);

\path[draw=drawColor,line width= 0.4pt,line join=round,line cap=round] ( 63.19, 60.63) -- ( 67.69, 65.13);

\path[draw=drawColor,line width= 0.4pt,line join=round,line cap=round] ( 63.19, 65.13) -- ( 67.69, 60.63);

\path[draw=drawColor,line width= 0.4pt,line join=round,line cap=round] ( 63.31, 60.78) -- ( 67.81, 65.28);

\path[draw=drawColor,line width= 0.4pt,line join=round,line cap=round] ( 63.31, 65.28) -- ( 67.81, 60.78);

\path[draw=drawColor,line width= 0.4pt,line join=round,line cap=round] ( 63.42, 60.89) -- ( 67.92, 65.39);

\path[draw=drawColor,line width= 0.4pt,line join=round,line cap=round] ( 63.42, 65.39) -- ( 67.92, 60.89);

\path[draw=drawColor,line width= 0.4pt,line join=round,line cap=round] ( 63.54, 60.91) -- ( 68.04, 65.41);

\path[draw=drawColor,line width= 0.4pt,line join=round,line cap=round] ( 63.54, 65.41) -- ( 68.04, 60.91);

\path[draw=drawColor,line width= 0.4pt,line join=round,line cap=round] ( 63.66, 60.99) -- ( 68.16, 65.49);

\path[draw=drawColor,line width= 0.4pt,line join=round,line cap=round] ( 63.66, 65.49) -- ( 68.16, 60.99);

\path[draw=drawColor,line width= 0.4pt,line join=round,line cap=round] ( 63.77, 61.03) -- ( 68.27, 65.53);

\path[draw=drawColor,line width= 0.4pt,line join=round,line cap=round] ( 63.77, 65.53) -- ( 68.27, 61.03);

\path[draw=drawColor,line width= 0.4pt,line join=round,line cap=round] ( 63.89, 61.13) -- ( 68.39, 65.63);

\path[draw=drawColor,line width= 0.4pt,line join=round,line cap=round] ( 63.89, 65.63) -- ( 68.39, 61.13);

\path[draw=drawColor,line width= 0.4pt,line join=round,line cap=round] ( 64.01, 61.19) -- ( 68.51, 65.69);

\path[draw=drawColor,line width= 0.4pt,line join=round,line cap=round] ( 64.01, 65.69) -- ( 68.51, 61.19);

\path[draw=drawColor,line width= 0.4pt,line join=round,line cap=round] ( 64.12, 61.28) -- ( 68.62, 65.78);

\path[draw=drawColor,line width= 0.4pt,line join=round,line cap=round] ( 64.12, 65.78) -- ( 68.62, 61.28);

\path[draw=drawColor,line width= 0.4pt,line join=round,line cap=round] ( 64.24, 61.39) -- ( 68.74, 65.89);

\path[draw=drawColor,line width= 0.4pt,line join=round,line cap=round] ( 64.24, 65.89) -- ( 68.74, 61.39);

\path[draw=drawColor,line width= 0.4pt,line join=round,line cap=round] ( 64.36, 61.62) -- ( 68.86, 66.12);

\path[draw=drawColor,line width= 0.4pt,line join=round,line cap=round] ( 64.36, 66.12) -- ( 68.86, 61.62);

\path[draw=drawColor,line width= 0.4pt,line join=round,line cap=round] ( 64.47, 61.80) -- ( 68.97, 66.30);

\path[draw=drawColor,line width= 0.4pt,line join=round,line cap=round] ( 64.47, 66.30) -- ( 68.97, 61.80);

\path[draw=drawColor,line width= 0.4pt,line join=round,line cap=round] ( 64.59, 62.33) -- ( 69.09, 66.83);

\path[draw=drawColor,line width= 0.4pt,line join=round,line cap=round] ( 64.59, 66.83) -- ( 69.09, 62.33);

\path[draw=drawColor,line width= 0.4pt,line join=round,line cap=round] ( 64.71, 62.36) -- ( 69.21, 66.86);

\path[draw=drawColor,line width= 0.4pt,line join=round,line cap=round] ( 64.71, 66.86) -- ( 69.21, 62.36);

\path[draw=drawColor,line width= 0.4pt,line join=round,line cap=round] ( 64.83, 62.38) -- ( 69.33, 66.88);

\path[draw=drawColor,line width= 0.4pt,line join=round,line cap=round] ( 64.83, 66.88) -- ( 69.33, 62.38);

\path[draw=drawColor,line width= 0.4pt,line join=round,line cap=round] ( 64.94, 62.40) -- ( 69.44, 66.90);

\path[draw=drawColor,line width= 0.4pt,line join=round,line cap=round] ( 64.94, 66.90) -- ( 69.44, 62.40);

\path[draw=drawColor,line width= 0.4pt,line join=round,line cap=round] ( 65.06, 62.60) -- ( 69.56, 67.10);

\path[draw=drawColor,line width= 0.4pt,line join=round,line cap=round] ( 65.06, 67.10) -- ( 69.56, 62.60);

\path[draw=drawColor,line width= 0.4pt,line join=round,line cap=round] ( 65.18, 62.63) -- ( 69.68, 67.13);

\path[draw=drawColor,line width= 0.4pt,line join=round,line cap=round] ( 65.18, 67.13) -- ( 69.68, 62.63);

\path[draw=drawColor,line width= 0.4pt,line join=round,line cap=round] ( 65.29, 63.10) -- ( 69.79, 67.60);

\path[draw=drawColor,line width= 0.4pt,line join=round,line cap=round] ( 65.29, 67.60) -- ( 69.79, 63.10);

\path[draw=drawColor,line width= 0.4pt,line join=round,line cap=round] ( 65.41, 63.17) -- ( 69.91, 67.67);

\path[draw=drawColor,line width= 0.4pt,line join=round,line cap=round] ( 65.41, 67.67) -- ( 69.91, 63.17);

\path[draw=drawColor,line width= 0.4pt,line join=round,line cap=round] ( 65.53, 63.23) -- ( 70.03, 67.73);

\path[draw=drawColor,line width= 0.4pt,line join=round,line cap=round] ( 65.53, 67.73) -- ( 70.03, 63.23);

\path[draw=drawColor,line width= 0.4pt,line join=round,line cap=round] ( 65.64, 63.25) -- ( 70.14, 67.75);

\path[draw=drawColor,line width= 0.4pt,line join=round,line cap=round] ( 65.64, 67.75) -- ( 70.14, 63.25);

\path[draw=drawColor,line width= 0.4pt,line join=round,line cap=round] ( 65.76, 63.41) -- ( 70.26, 67.91);

\path[draw=drawColor,line width= 0.4pt,line join=round,line cap=round] ( 65.76, 67.91) -- ( 70.26, 63.41);

\path[draw=drawColor,line width= 0.4pt,line join=round,line cap=round] ( 65.88, 63.46) -- ( 70.38, 67.96);

\path[draw=drawColor,line width= 0.4pt,line join=round,line cap=round] ( 65.88, 67.96) -- ( 70.38, 63.46);

\path[draw=drawColor,line width= 0.4pt,line join=round,line cap=round] ( 65.99, 63.66) -- ( 70.49, 68.16);

\path[draw=drawColor,line width= 0.4pt,line join=round,line cap=round] ( 65.99, 68.16) -- ( 70.49, 63.66);

\path[draw=drawColor,line width= 0.4pt,line join=round,line cap=round] ( 66.11, 63.88) -- ( 70.61, 68.38);

\path[draw=drawColor,line width= 0.4pt,line join=round,line cap=round] ( 66.11, 68.38) -- ( 70.61, 63.88);

\path[draw=drawColor,line width= 0.4pt,line join=round,line cap=round] ( 66.23, 63.92) -- ( 70.73, 68.42);

\path[draw=drawColor,line width= 0.4pt,line join=round,line cap=round] ( 66.23, 68.42) -- ( 70.73, 63.92);

\path[draw=drawColor,line width= 0.4pt,line join=round,line cap=round] ( 66.35, 64.14) -- ( 70.85, 68.64);

\path[draw=drawColor,line width= 0.4pt,line join=round,line cap=round] ( 66.35, 68.64) -- ( 70.85, 64.14);

\path[draw=drawColor,line width= 0.4pt,line join=round,line cap=round] ( 66.46, 64.20) -- ( 70.96, 68.70);

\path[draw=drawColor,line width= 0.4pt,line join=round,line cap=round] ( 66.46, 68.70) -- ( 70.96, 64.20);

\path[draw=drawColor,line width= 0.4pt,line join=round,line cap=round] ( 66.58, 64.42) -- ( 71.08, 68.92);

\path[draw=drawColor,line width= 0.4pt,line join=round,line cap=round] ( 66.58, 68.92) -- ( 71.08, 64.42);

\path[draw=drawColor,line width= 0.4pt,line join=round,line cap=round] ( 66.70, 64.42) -- ( 71.20, 68.92);

\path[draw=drawColor,line width= 0.4pt,line join=round,line cap=round] ( 66.70, 68.92) -- ( 71.20, 64.42);

\path[draw=drawColor,line width= 0.4pt,line join=round,line cap=round] ( 66.81, 64.47) -- ( 71.31, 68.97);

\path[draw=drawColor,line width= 0.4pt,line join=round,line cap=round] ( 66.81, 68.97) -- ( 71.31, 64.47);

\path[draw=drawColor,line width= 0.4pt,line join=round,line cap=round] ( 66.93, 64.55) -- ( 71.43, 69.05);

\path[draw=drawColor,line width= 0.4pt,line join=round,line cap=round] ( 66.93, 69.05) -- ( 71.43, 64.55);

\path[draw=drawColor,line width= 0.4pt,line join=round,line cap=round] ( 67.05, 64.62) -- ( 71.55, 69.12);

\path[draw=drawColor,line width= 0.4pt,line join=round,line cap=round] ( 67.05, 69.12) -- ( 71.55, 64.62);

\path[draw=drawColor,line width= 0.4pt,line join=round,line cap=round] ( 67.16, 64.69) -- ( 71.66, 69.19);

\path[draw=drawColor,line width= 0.4pt,line join=round,line cap=round] ( 67.16, 69.19) -- ( 71.66, 64.69);

\path[draw=drawColor,line width= 0.4pt,line join=round,line cap=round] ( 67.28, 64.73) -- ( 71.78, 69.23);

\path[draw=drawColor,line width= 0.4pt,line join=round,line cap=round] ( 67.28, 69.23) -- ( 71.78, 64.73);

\path[draw=drawColor,line width= 0.4pt,line join=round,line cap=round] ( 67.40, 64.79) -- ( 71.90, 69.29);

\path[draw=drawColor,line width= 0.4pt,line join=round,line cap=round] ( 67.40, 69.29) -- ( 71.90, 64.79);

\path[draw=drawColor,line width= 0.4pt,line join=round,line cap=round] ( 67.52, 64.80) -- ( 72.02, 69.30);

\path[draw=drawColor,line width= 0.4pt,line join=round,line cap=round] ( 67.52, 69.30) -- ( 72.02, 64.80);

\path[draw=drawColor,line width= 0.4pt,line join=round,line cap=round] ( 67.63, 65.16) -- ( 72.13, 69.66);

\path[draw=drawColor,line width= 0.4pt,line join=round,line cap=round] ( 67.63, 69.66) -- ( 72.13, 65.16);

\path[draw=drawColor,line width= 0.4pt,line join=round,line cap=round] ( 67.75, 65.17) -- ( 72.25, 69.67);

\path[draw=drawColor,line width= 0.4pt,line join=round,line cap=round] ( 67.75, 69.67) -- ( 72.25, 65.17);

\path[draw=drawColor,line width= 0.4pt,line join=round,line cap=round] ( 67.87, 65.31) -- ( 72.37, 69.81);

\path[draw=drawColor,line width= 0.4pt,line join=round,line cap=round] ( 67.87, 69.81) -- ( 72.37, 65.31);

\path[draw=drawColor,line width= 0.4pt,line join=round,line cap=round] ( 67.98, 65.32) -- ( 72.48, 69.82);

\path[draw=drawColor,line width= 0.4pt,line join=round,line cap=round] ( 67.98, 69.82) -- ( 72.48, 65.32);

\path[draw=drawColor,line width= 0.4pt,line join=round,line cap=round] ( 68.10, 65.32) -- ( 72.60, 69.82);

\path[draw=drawColor,line width= 0.4pt,line join=round,line cap=round] ( 68.10, 69.82) -- ( 72.60, 65.32);

\path[draw=drawColor,line width= 0.4pt,line join=round,line cap=round] ( 68.22, 65.48) -- ( 72.72, 69.98);

\path[draw=drawColor,line width= 0.4pt,line join=round,line cap=round] ( 68.22, 69.98) -- ( 72.72, 65.48);

\path[draw=drawColor,line width= 0.4pt,line join=round,line cap=round] ( 68.33, 65.57) -- ( 72.83, 70.07);

\path[draw=drawColor,line width= 0.4pt,line join=round,line cap=round] ( 68.33, 70.07) -- ( 72.83, 65.57);

\path[draw=drawColor,line width= 0.4pt,line join=round,line cap=round] ( 68.45, 65.57) -- ( 72.95, 70.07);

\path[draw=drawColor,line width= 0.4pt,line join=round,line cap=round] ( 68.45, 70.07) -- ( 72.95, 65.57);

\path[draw=drawColor,line width= 0.4pt,line join=round,line cap=round] ( 68.57, 65.67) -- ( 73.07, 70.17);

\path[draw=drawColor,line width= 0.4pt,line join=round,line cap=round] ( 68.57, 70.17) -- ( 73.07, 65.67);

\path[draw=drawColor,line width= 0.4pt,line join=round,line cap=round] ( 68.68, 65.70) -- ( 73.18, 70.20);

\path[draw=drawColor,line width= 0.4pt,line join=round,line cap=round] ( 68.68, 70.20) -- ( 73.18, 65.70);

\path[draw=drawColor,line width= 0.4pt,line join=round,line cap=round] ( 68.80, 65.71) -- ( 73.30, 70.21);

\path[draw=drawColor,line width= 0.4pt,line join=round,line cap=round] ( 68.80, 70.21) -- ( 73.30, 65.71);

\path[draw=drawColor,line width= 0.4pt,line join=round,line cap=round] ( 68.92, 66.00) -- ( 73.42, 70.50);

\path[draw=drawColor,line width= 0.4pt,line join=round,line cap=round] ( 68.92, 70.50) -- ( 73.42, 66.00);

\path[draw=drawColor,line width= 0.4pt,line join=round,line cap=round] ( 69.04, 66.15) -- ( 73.54, 70.65);

\path[draw=drawColor,line width= 0.4pt,line join=round,line cap=round] ( 69.04, 70.65) -- ( 73.54, 66.15);

\path[draw=drawColor,line width= 0.4pt,line join=round,line cap=round] ( 69.15, 66.29) -- ( 73.65, 70.79);

\path[draw=drawColor,line width= 0.4pt,line join=round,line cap=round] ( 69.15, 70.79) -- ( 73.65, 66.29);

\path[draw=drawColor,line width= 0.4pt,line join=round,line cap=round] ( 69.27, 66.52) -- ( 73.77, 71.02);

\path[draw=drawColor,line width= 0.4pt,line join=round,line cap=round] ( 69.27, 71.02) -- ( 73.77, 66.52);

\path[draw=drawColor,line width= 0.4pt,line join=round,line cap=round] ( 69.39, 66.65) -- ( 73.89, 71.15);

\path[draw=drawColor,line width= 0.4pt,line join=round,line cap=round] ( 69.39, 71.15) -- ( 73.89, 66.65);

\path[draw=drawColor,line width= 0.4pt,line join=round,line cap=round] ( 69.50, 66.77) -- ( 74.00, 71.27);

\path[draw=drawColor,line width= 0.4pt,line join=round,line cap=round] ( 69.50, 71.27) -- ( 74.00, 66.77);

\path[draw=drawColor,line width= 0.4pt,line join=round,line cap=round] ( 69.62, 67.27) -- ( 74.12, 71.77);

\path[draw=drawColor,line width= 0.4pt,line join=round,line cap=round] ( 69.62, 71.77) -- ( 74.12, 67.27);

\path[draw=drawColor,line width= 0.4pt,line join=round,line cap=round] ( 69.74, 67.64) -- ( 74.24, 72.14);

\path[draw=drawColor,line width= 0.4pt,line join=round,line cap=round] ( 69.74, 72.14) -- ( 74.24, 67.64);

\path[draw=drawColor,line width= 0.4pt,line join=round,line cap=round] ( 69.85, 67.84) -- ( 74.35, 72.34);

\path[draw=drawColor,line width= 0.4pt,line join=round,line cap=round] ( 69.85, 72.34) -- ( 74.35, 67.84);

\path[draw=drawColor,line width= 0.4pt,line join=round,line cap=round] ( 69.97, 67.93) -- ( 74.47, 72.43);

\path[draw=drawColor,line width= 0.4pt,line join=round,line cap=round] ( 69.97, 72.43) -- ( 74.47, 67.93);

\path[draw=drawColor,line width= 0.4pt,line join=round,line cap=round] ( 70.09, 68.02) -- ( 74.59, 72.52);

\path[draw=drawColor,line width= 0.4pt,line join=round,line cap=round] ( 70.09, 72.52) -- ( 74.59, 68.02);

\path[draw=drawColor,line width= 0.4pt,line join=round,line cap=round] ( 70.20, 68.04) -- ( 74.70, 72.54);

\path[draw=drawColor,line width= 0.4pt,line join=round,line cap=round] ( 70.20, 72.54) -- ( 74.70, 68.04);

\path[draw=drawColor,line width= 0.4pt,line join=round,line cap=round] ( 70.32, 68.46) -- ( 74.82, 72.96);

\path[draw=drawColor,line width= 0.4pt,line join=round,line cap=round] ( 70.32, 72.96) -- ( 74.82, 68.46);

\path[draw=drawColor,line width= 0.4pt,line join=round,line cap=round] ( 70.44, 68.50) -- ( 74.94, 73.00);

\path[draw=drawColor,line width= 0.4pt,line join=round,line cap=round] ( 70.44, 73.00) -- ( 74.94, 68.50);

\path[draw=drawColor,line width= 0.4pt,line join=round,line cap=round] ( 70.56, 68.53) -- ( 75.06, 73.03);

\path[draw=drawColor,line width= 0.4pt,line join=round,line cap=round] ( 70.56, 73.03) -- ( 75.06, 68.53);

\path[draw=drawColor,line width= 0.4pt,line join=round,line cap=round] ( 70.67, 68.72) -- ( 75.17, 73.22);

\path[draw=drawColor,line width= 0.4pt,line join=round,line cap=round] ( 70.67, 73.22) -- ( 75.17, 68.72);

\path[draw=drawColor,line width= 0.4pt,line join=round,line cap=round] ( 70.79, 68.85) -- ( 75.29, 73.35);

\path[draw=drawColor,line width= 0.4pt,line join=round,line cap=round] ( 70.79, 73.35) -- ( 75.29, 68.85);

\path[draw=drawColor,line width= 0.4pt,line join=round,line cap=round] ( 70.91, 68.89) -- ( 75.41, 73.39);

\path[draw=drawColor,line width= 0.4pt,line join=round,line cap=round] ( 70.91, 73.39) -- ( 75.41, 68.89);

\path[draw=drawColor,line width= 0.4pt,line join=round,line cap=round] ( 71.02, 68.99) -- ( 75.52, 73.49);

\path[draw=drawColor,line width= 0.4pt,line join=round,line cap=round] ( 71.02, 73.49) -- ( 75.52, 68.99);

\path[draw=drawColor,line width= 0.4pt,line join=round,line cap=round] ( 71.14, 69.12) -- ( 75.64, 73.62);

\path[draw=drawColor,line width= 0.4pt,line join=round,line cap=round] ( 71.14, 73.62) -- ( 75.64, 69.12);

\path[draw=drawColor,line width= 0.4pt,line join=round,line cap=round] ( 71.26, 69.25) -- ( 75.76, 73.75);

\path[draw=drawColor,line width= 0.4pt,line join=round,line cap=round] ( 71.26, 73.75) -- ( 75.76, 69.25);

\path[draw=drawColor,line width= 0.4pt,line join=round,line cap=round] ( 71.37, 69.28) -- ( 75.87, 73.78);

\path[draw=drawColor,line width= 0.4pt,line join=round,line cap=round] ( 71.37, 73.78) -- ( 75.87, 69.28);

\path[draw=drawColor,line width= 0.4pt,line join=round,line cap=round] ( 71.49, 69.38) -- ( 75.99, 73.88);

\path[draw=drawColor,line width= 0.4pt,line join=round,line cap=round] ( 71.49, 73.88) -- ( 75.99, 69.38);

\path[draw=drawColor,line width= 0.4pt,line join=round,line cap=round] ( 71.61, 69.41) -- ( 76.11, 73.91);

\path[draw=drawColor,line width= 0.4pt,line join=round,line cap=round] ( 71.61, 73.91) -- ( 76.11, 69.41);

\path[draw=drawColor,line width= 0.4pt,line join=round,line cap=round] ( 71.72, 69.61) -- ( 76.22, 74.11);

\path[draw=drawColor,line width= 0.4pt,line join=round,line cap=round] ( 71.72, 74.11) -- ( 76.22, 69.61);

\path[draw=drawColor,line width= 0.4pt,line join=round,line cap=round] ( 71.84, 69.63) -- ( 76.34, 74.13);

\path[draw=drawColor,line width= 0.4pt,line join=round,line cap=round] ( 71.84, 74.13) -- ( 76.34, 69.63);

\path[draw=drawColor,line width= 0.4pt,line join=round,line cap=round] ( 71.96, 69.64) -- ( 76.46, 74.14);

\path[draw=drawColor,line width= 0.4pt,line join=round,line cap=round] ( 71.96, 74.14) -- ( 76.46, 69.64);

\path[draw=drawColor,line width= 0.4pt,line join=round,line cap=round] ( 72.08, 69.82) -- ( 76.58, 74.32);

\path[draw=drawColor,line width= 0.4pt,line join=round,line cap=round] ( 72.08, 74.32) -- ( 76.58, 69.82);

\path[draw=drawColor,line width= 0.4pt,line join=round,line cap=round] ( 72.19, 70.14) -- ( 76.69, 74.64);

\path[draw=drawColor,line width= 0.4pt,line join=round,line cap=round] ( 72.19, 74.64) -- ( 76.69, 70.14);

\path[draw=drawColor,line width= 0.4pt,line join=round,line cap=round] ( 72.31, 70.20) -- ( 76.81, 74.70);

\path[draw=drawColor,line width= 0.4pt,line join=round,line cap=round] ( 72.31, 74.70) -- ( 76.81, 70.20);

\path[draw=drawColor,line width= 0.4pt,line join=round,line cap=round] ( 72.43, 70.31) -- ( 76.93, 74.81);

\path[draw=drawColor,line width= 0.4pt,line join=round,line cap=round] ( 72.43, 74.81) -- ( 76.93, 70.31);

\path[draw=drawColor,line width= 0.4pt,line join=round,line cap=round] ( 72.54, 70.35) -- ( 77.04, 74.85);

\path[draw=drawColor,line width= 0.4pt,line join=round,line cap=round] ( 72.54, 74.85) -- ( 77.04, 70.35);

\path[draw=drawColor,line width= 0.4pt,line join=round,line cap=round] ( 72.66, 70.70) -- ( 77.16, 75.20);

\path[draw=drawColor,line width= 0.4pt,line join=round,line cap=round] ( 72.66, 75.20) -- ( 77.16, 70.70);

\path[draw=drawColor,line width= 0.4pt,line join=round,line cap=round] ( 72.78, 70.71) -- ( 77.28, 75.21);

\path[draw=drawColor,line width= 0.4pt,line join=round,line cap=round] ( 72.78, 75.21) -- ( 77.28, 70.71);

\path[draw=drawColor,line width= 0.4pt,line join=round,line cap=round] ( 72.89, 70.79) -- ( 77.39, 75.29);

\path[draw=drawColor,line width= 0.4pt,line join=round,line cap=round] ( 72.89, 75.29) -- ( 77.39, 70.79);

\path[draw=drawColor,line width= 0.4pt,line join=round,line cap=round] ( 73.01, 70.81) -- ( 77.51, 75.31);

\path[draw=drawColor,line width= 0.4pt,line join=round,line cap=round] ( 73.01, 75.31) -- ( 77.51, 70.81);

\path[draw=drawColor,line width= 0.4pt,line join=round,line cap=round] ( 73.13, 70.86) -- ( 77.63, 75.36);

\path[draw=drawColor,line width= 0.4pt,line join=round,line cap=round] ( 73.13, 75.36) -- ( 77.63, 70.86);

\path[draw=drawColor,line width= 0.4pt,line join=round,line cap=round] ( 73.25, 70.92) -- ( 77.75, 75.42);

\path[draw=drawColor,line width= 0.4pt,line join=round,line cap=round] ( 73.25, 75.42) -- ( 77.75, 70.92);

\path[draw=drawColor,line width= 0.4pt,line join=round,line cap=round] ( 73.36, 71.10) -- ( 77.86, 75.60);

\path[draw=drawColor,line width= 0.4pt,line join=round,line cap=round] ( 73.36, 75.60) -- ( 77.86, 71.10);

\path[draw=drawColor,line width= 0.4pt,line join=round,line cap=round] ( 73.48, 71.31) -- ( 77.98, 75.81);

\path[draw=drawColor,line width= 0.4pt,line join=round,line cap=round] ( 73.48, 75.81) -- ( 77.98, 71.31);

\path[draw=drawColor,line width= 0.4pt,line join=round,line cap=round] ( 73.60, 71.39) -- ( 78.10, 75.89);

\path[draw=drawColor,line width= 0.4pt,line join=round,line cap=round] ( 73.60, 75.89) -- ( 78.10, 71.39);

\path[draw=drawColor,line width= 0.4pt,line join=round,line cap=round] ( 73.71, 71.52) -- ( 78.21, 76.02);

\path[draw=drawColor,line width= 0.4pt,line join=round,line cap=round] ( 73.71, 76.02) -- ( 78.21, 71.52);

\path[draw=drawColor,line width= 0.4pt,line join=round,line cap=round] ( 73.83, 71.67) -- ( 78.33, 76.17);

\path[draw=drawColor,line width= 0.4pt,line join=round,line cap=round] ( 73.83, 76.17) -- ( 78.33, 71.67);

\path[draw=drawColor,line width= 0.4pt,line join=round,line cap=round] ( 73.95, 71.68) -- ( 78.45, 76.18);

\path[draw=drawColor,line width= 0.4pt,line join=round,line cap=round] ( 73.95, 76.18) -- ( 78.45, 71.68);

\path[draw=drawColor,line width= 0.4pt,line join=round,line cap=round] ( 74.06, 71.74) -- ( 78.56, 76.24);

\path[draw=drawColor,line width= 0.4pt,line join=round,line cap=round] ( 74.06, 76.24) -- ( 78.56, 71.74);

\path[draw=drawColor,line width= 0.4pt,line join=round,line cap=round] ( 74.18, 72.04) -- ( 78.68, 76.54);

\path[draw=drawColor,line width= 0.4pt,line join=round,line cap=round] ( 74.18, 76.54) -- ( 78.68, 72.04);

\path[draw=drawColor,line width= 0.4pt,line join=round,line cap=round] ( 74.30, 72.05) -- ( 78.80, 76.55);

\path[draw=drawColor,line width= 0.4pt,line join=round,line cap=round] ( 74.30, 76.55) -- ( 78.80, 72.05);

\path[draw=drawColor,line width= 0.4pt,line join=round,line cap=round] ( 74.41, 72.08) -- ( 78.91, 76.58);

\path[draw=drawColor,line width= 0.4pt,line join=round,line cap=round] ( 74.41, 76.58) -- ( 78.91, 72.08);

\path[draw=drawColor,line width= 0.4pt,line join=round,line cap=round] ( 74.53, 72.22) -- ( 79.03, 76.72);

\path[draw=drawColor,line width= 0.4pt,line join=round,line cap=round] ( 74.53, 76.72) -- ( 79.03, 72.22);

\path[draw=drawColor,line width= 0.4pt,line join=round,line cap=round] ( 74.65, 72.26) -- ( 79.15, 76.76);

\path[draw=drawColor,line width= 0.4pt,line join=round,line cap=round] ( 74.65, 76.76) -- ( 79.15, 72.26);

\path[draw=drawColor,line width= 0.4pt,line join=round,line cap=round] ( 74.77, 72.26) -- ( 79.27, 76.76);

\path[draw=drawColor,line width= 0.4pt,line join=round,line cap=round] ( 74.77, 76.76) -- ( 79.27, 72.26);

\path[draw=drawColor,line width= 0.4pt,line join=round,line cap=round] ( 74.88, 72.30) -- ( 79.38, 76.80);

\path[draw=drawColor,line width= 0.4pt,line join=round,line cap=round] ( 74.88, 76.80) -- ( 79.38, 72.30);

\path[draw=drawColor,line width= 0.4pt,line join=round,line cap=round] ( 75.00, 72.39) -- ( 79.50, 76.89);

\path[draw=drawColor,line width= 0.4pt,line join=round,line cap=round] ( 75.00, 76.89) -- ( 79.50, 72.39);

\path[draw=drawColor,line width= 0.4pt,line join=round,line cap=round] ( 75.12, 72.42) -- ( 79.62, 76.92);

\path[draw=drawColor,line width= 0.4pt,line join=round,line cap=round] ( 75.12, 76.92) -- ( 79.62, 72.42);

\path[draw=drawColor,line width= 0.4pt,line join=round,line cap=round] ( 75.23, 72.66) -- ( 79.73, 77.16);

\path[draw=drawColor,line width= 0.4pt,line join=round,line cap=round] ( 75.23, 77.16) -- ( 79.73, 72.66);

\path[draw=drawColor,line width= 0.4pt,line join=round,line cap=round] ( 75.35, 72.89) -- ( 79.85, 77.39);

\path[draw=drawColor,line width= 0.4pt,line join=round,line cap=round] ( 75.35, 77.39) -- ( 79.85, 72.89);

\path[draw=drawColor,line width= 0.4pt,line join=round,line cap=round] ( 75.47, 72.99) -- ( 79.97, 77.49);

\path[draw=drawColor,line width= 0.4pt,line join=round,line cap=round] ( 75.47, 77.49) -- ( 79.97, 72.99);

\path[draw=drawColor,line width= 0.4pt,line join=round,line cap=round] ( 75.58, 73.07) -- ( 80.08, 77.57);

\path[draw=drawColor,line width= 0.4pt,line join=round,line cap=round] ( 75.58, 77.57) -- ( 80.08, 73.07);

\path[draw=drawColor,line width= 0.4pt,line join=round,line cap=round] ( 75.70, 73.31) -- ( 80.20, 77.81);

\path[draw=drawColor,line width= 0.4pt,line join=round,line cap=round] ( 75.70, 77.81) -- ( 80.20, 73.31);

\path[draw=drawColor,line width= 0.4pt,line join=round,line cap=round] ( 75.82, 73.31) -- ( 80.32, 77.81);

\path[draw=drawColor,line width= 0.4pt,line join=round,line cap=round] ( 75.82, 77.81) -- ( 80.32, 73.31);

\path[draw=drawColor,line width= 0.4pt,line join=round,line cap=round] ( 75.93, 73.35) -- ( 80.43, 77.85);

\path[draw=drawColor,line width= 0.4pt,line join=round,line cap=round] ( 75.93, 77.85) -- ( 80.43, 73.35);

\path[draw=drawColor,line width= 0.4pt,line join=round,line cap=round] ( 76.05, 73.42) -- ( 80.55, 77.92);

\path[draw=drawColor,line width= 0.4pt,line join=round,line cap=round] ( 76.05, 77.92) -- ( 80.55, 73.42);

\path[draw=drawColor,line width= 0.4pt,line join=round,line cap=round] ( 76.17, 73.58) -- ( 80.67, 78.08);

\path[draw=drawColor,line width= 0.4pt,line join=round,line cap=round] ( 76.17, 78.08) -- ( 80.67, 73.58);

\path[draw=drawColor,line width= 0.4pt,line join=round,line cap=round] ( 76.29, 73.63) -- ( 80.79, 78.13);

\path[draw=drawColor,line width= 0.4pt,line join=round,line cap=round] ( 76.29, 78.13) -- ( 80.79, 73.63);

\path[draw=drawColor,line width= 0.4pt,line join=round,line cap=round] ( 76.40, 73.65) -- ( 80.90, 78.15);

\path[draw=drawColor,line width= 0.4pt,line join=round,line cap=round] ( 76.40, 78.15) -- ( 80.90, 73.65);

\path[draw=drawColor,line width= 0.4pt,line join=round,line cap=round] ( 76.52, 74.00) -- ( 81.02, 78.50);

\path[draw=drawColor,line width= 0.4pt,line join=round,line cap=round] ( 76.52, 78.50) -- ( 81.02, 74.00);

\path[draw=drawColor,line width= 0.4pt,line join=round,line cap=round] ( 76.64, 74.02) -- ( 81.14, 78.52);

\path[draw=drawColor,line width= 0.4pt,line join=round,line cap=round] ( 76.64, 78.52) -- ( 81.14, 74.02);

\path[draw=drawColor,line width= 0.4pt,line join=round,line cap=round] ( 76.75, 74.13) -- ( 81.25, 78.63);

\path[draw=drawColor,line width= 0.4pt,line join=round,line cap=round] ( 76.75, 78.63) -- ( 81.25, 74.13);

\path[draw=drawColor,line width= 0.4pt,line join=round,line cap=round] ( 76.87, 74.25) -- ( 81.37, 78.75);

\path[draw=drawColor,line width= 0.4pt,line join=round,line cap=round] ( 76.87, 78.75) -- ( 81.37, 74.25);

\path[draw=drawColor,line width= 0.4pt,line join=round,line cap=round] ( 76.99, 74.39) -- ( 81.49, 78.89);

\path[draw=drawColor,line width= 0.4pt,line join=round,line cap=round] ( 76.99, 78.89) -- ( 81.49, 74.39);

\path[draw=drawColor,line width= 0.4pt,line join=round,line cap=round] ( 77.10, 74.44) -- ( 81.60, 78.94);

\path[draw=drawColor,line width= 0.4pt,line join=round,line cap=round] ( 77.10, 78.94) -- ( 81.60, 74.44);

\path[draw=drawColor,line width= 0.4pt,line join=round,line cap=round] ( 77.22, 74.51) -- ( 81.72, 79.01);

\path[draw=drawColor,line width= 0.4pt,line join=round,line cap=round] ( 77.22, 79.01) -- ( 81.72, 74.51);

\path[draw=drawColor,line width= 0.4pt,line join=round,line cap=round] ( 77.34, 74.79) -- ( 81.84, 79.29);

\path[draw=drawColor,line width= 0.4pt,line join=round,line cap=round] ( 77.34, 79.29) -- ( 81.84, 74.79);

\path[draw=drawColor,line width= 0.4pt,line join=round,line cap=round] ( 77.46, 74.83) -- ( 81.96, 79.33);

\path[draw=drawColor,line width= 0.4pt,line join=round,line cap=round] ( 77.46, 79.33) -- ( 81.96, 74.83);

\path[draw=drawColor,line width= 0.4pt,line join=round,line cap=round] ( 77.57, 74.85) -- ( 82.07, 79.35);

\path[draw=drawColor,line width= 0.4pt,line join=round,line cap=round] ( 77.57, 79.35) -- ( 82.07, 74.85);

\path[draw=drawColor,line width= 0.4pt,line join=round,line cap=round] ( 77.69, 75.05) -- ( 82.19, 79.55);

\path[draw=drawColor,line width= 0.4pt,line join=round,line cap=round] ( 77.69, 79.55) -- ( 82.19, 75.05);

\path[draw=drawColor,line width= 0.4pt,line join=round,line cap=round] ( 77.81, 75.13) -- ( 82.31, 79.63);

\path[draw=drawColor,line width= 0.4pt,line join=round,line cap=round] ( 77.81, 79.63) -- ( 82.31, 75.13);

\path[draw=drawColor,line width= 0.4pt,line join=round,line cap=round] ( 77.92, 75.25) -- ( 82.42, 79.75);

\path[draw=drawColor,line width= 0.4pt,line join=round,line cap=round] ( 77.92, 79.75) -- ( 82.42, 75.25);

\path[draw=drawColor,line width= 0.4pt,line join=round,line cap=round] ( 78.04, 75.32) -- ( 82.54, 79.82);

\path[draw=drawColor,line width= 0.4pt,line join=round,line cap=round] ( 78.04, 79.82) -- ( 82.54, 75.32);

\path[draw=drawColor,line width= 0.4pt,line join=round,line cap=round] ( 78.16, 75.37) -- ( 82.66, 79.87);

\path[draw=drawColor,line width= 0.4pt,line join=round,line cap=round] ( 78.16, 79.87) -- ( 82.66, 75.37);

\path[draw=drawColor,line width= 0.4pt,line join=round,line cap=round] ( 78.27, 75.38) -- ( 82.77, 79.88);

\path[draw=drawColor,line width= 0.4pt,line join=round,line cap=round] ( 78.27, 79.88) -- ( 82.77, 75.38);

\path[draw=drawColor,line width= 0.4pt,line join=round,line cap=round] ( 78.39, 75.55) -- ( 82.89, 80.05);

\path[draw=drawColor,line width= 0.4pt,line join=round,line cap=round] ( 78.39, 80.05) -- ( 82.89, 75.55);

\path[draw=drawColor,line width= 0.4pt,line join=round,line cap=round] ( 78.51, 75.65) -- ( 83.01, 80.15);

\path[draw=drawColor,line width= 0.4pt,line join=round,line cap=round] ( 78.51, 80.15) -- ( 83.01, 75.65);

\path[draw=drawColor,line width= 0.4pt,line join=round,line cap=round] ( 78.62, 75.65) -- ( 83.12, 80.15);

\path[draw=drawColor,line width= 0.4pt,line join=round,line cap=round] ( 78.62, 80.15) -- ( 83.12, 75.65);

\path[draw=drawColor,line width= 0.4pt,line join=round,line cap=round] ( 78.74, 75.66) -- ( 83.24, 80.16);

\path[draw=drawColor,line width= 0.4pt,line join=round,line cap=round] ( 78.74, 80.16) -- ( 83.24, 75.66);

\path[draw=drawColor,line width= 0.4pt,line join=round,line cap=round] ( 78.86, 75.71) -- ( 83.36, 80.21);

\path[draw=drawColor,line width= 0.4pt,line join=round,line cap=round] ( 78.86, 80.21) -- ( 83.36, 75.71);

\path[draw=drawColor,line width= 0.4pt,line join=round,line cap=round] ( 78.98, 75.74) -- ( 83.48, 80.24);

\path[draw=drawColor,line width= 0.4pt,line join=round,line cap=round] ( 78.98, 80.24) -- ( 83.48, 75.74);

\path[draw=drawColor,line width= 0.4pt,line join=round,line cap=round] ( 79.09, 75.93) -- ( 83.59, 80.43);

\path[draw=drawColor,line width= 0.4pt,line join=round,line cap=round] ( 79.09, 80.43) -- ( 83.59, 75.93);

\path[draw=drawColor,line width= 0.4pt,line join=round,line cap=round] ( 79.21, 76.04) -- ( 83.71, 80.54);

\path[draw=drawColor,line width= 0.4pt,line join=round,line cap=round] ( 79.21, 80.54) -- ( 83.71, 76.04);

\path[draw=drawColor,line width= 0.4pt,line join=round,line cap=round] ( 79.33, 76.43) -- ( 83.83, 80.93);

\path[draw=drawColor,line width= 0.4pt,line join=round,line cap=round] ( 79.33, 80.93) -- ( 83.83, 76.43);

\path[draw=drawColor,line width= 0.4pt,line join=round,line cap=round] ( 79.44, 76.69) -- ( 83.94, 81.19);

\path[draw=drawColor,line width= 0.4pt,line join=round,line cap=round] ( 79.44, 81.19) -- ( 83.94, 76.69);

\path[draw=drawColor,line width= 0.4pt,line join=round,line cap=round] ( 79.56, 76.80) -- ( 84.06, 81.30);

\path[draw=drawColor,line width= 0.4pt,line join=round,line cap=round] ( 79.56, 81.30) -- ( 84.06, 76.80);

\path[draw=drawColor,line width= 0.4pt,line join=round,line cap=round] ( 79.68, 76.83) -- ( 84.18, 81.33);

\path[draw=drawColor,line width= 0.4pt,line join=round,line cap=round] ( 79.68, 81.33) -- ( 84.18, 76.83);

\path[draw=drawColor,line width= 0.4pt,line join=round,line cap=round] ( 79.79, 76.87) -- ( 84.29, 81.37);

\path[draw=drawColor,line width= 0.4pt,line join=round,line cap=round] ( 79.79, 81.37) -- ( 84.29, 76.87);

\path[draw=drawColor,line width= 0.4pt,line join=round,line cap=round] ( 79.91, 77.18) -- ( 84.41, 81.68);

\path[draw=drawColor,line width= 0.4pt,line join=round,line cap=round] ( 79.91, 81.68) -- ( 84.41, 77.18);

\path[draw=drawColor,line width= 0.4pt,line join=round,line cap=round] ( 80.03, 77.27) -- ( 84.53, 81.77);

\path[draw=drawColor,line width= 0.4pt,line join=round,line cap=round] ( 80.03, 81.77) -- ( 84.53, 77.27);

\path[draw=drawColor,line width= 0.4pt,line join=round,line cap=round] ( 80.14, 77.58) -- ( 84.64, 82.08);

\path[draw=drawColor,line width= 0.4pt,line join=round,line cap=round] ( 80.14, 82.08) -- ( 84.64, 77.58);

\path[draw=drawColor,line width= 0.4pt,line join=round,line cap=round] ( 80.26, 77.62) -- ( 84.76, 82.12);

\path[draw=drawColor,line width= 0.4pt,line join=round,line cap=round] ( 80.26, 82.12) -- ( 84.76, 77.62);

\path[draw=drawColor,line width= 0.4pt,line join=round,line cap=round] ( 80.38, 77.64) -- ( 84.88, 82.14);

\path[draw=drawColor,line width= 0.4pt,line join=round,line cap=round] ( 80.38, 82.14) -- ( 84.88, 77.64);

\path[draw=drawColor,line width= 0.4pt,line join=round,line cap=round] ( 80.50, 77.79) -- ( 85.00, 82.29);

\path[draw=drawColor,line width= 0.4pt,line join=round,line cap=round] ( 80.50, 82.29) -- ( 85.00, 77.79);

\path[draw=drawColor,line width= 0.4pt,line join=round,line cap=round] ( 80.61, 77.80) -- ( 85.11, 82.30);

\path[draw=drawColor,line width= 0.4pt,line join=round,line cap=round] ( 80.61, 82.30) -- ( 85.11, 77.80);

\path[draw=drawColor,line width= 0.4pt,line join=round,line cap=round] ( 80.73, 77.85) -- ( 85.23, 82.35);

\path[draw=drawColor,line width= 0.4pt,line join=round,line cap=round] ( 80.73, 82.35) -- ( 85.23, 77.85);

\path[draw=drawColor,line width= 0.4pt,line join=round,line cap=round] ( 80.85, 78.01) -- ( 85.35, 82.51);

\path[draw=drawColor,line width= 0.4pt,line join=round,line cap=round] ( 80.85, 82.51) -- ( 85.35, 78.01);

\path[draw=drawColor,line width= 0.4pt,line join=round,line cap=round] ( 80.96, 78.02) -- ( 85.46, 82.52);

\path[draw=drawColor,line width= 0.4pt,line join=round,line cap=round] ( 80.96, 82.52) -- ( 85.46, 78.02);

\path[draw=drawColor,line width= 0.4pt,line join=round,line cap=round] ( 81.08, 78.10) -- ( 85.58, 82.60);

\path[draw=drawColor,line width= 0.4pt,line join=round,line cap=round] ( 81.08, 82.60) -- ( 85.58, 78.10);

\path[draw=drawColor,line width= 0.4pt,line join=round,line cap=round] ( 81.20, 78.15) -- ( 85.70, 82.65);

\path[draw=drawColor,line width= 0.4pt,line join=round,line cap=round] ( 81.20, 82.65) -- ( 85.70, 78.15);

\path[draw=drawColor,line width= 0.4pt,line join=round,line cap=round] ( 81.31, 78.16) -- ( 85.81, 82.66);

\path[draw=drawColor,line width= 0.4pt,line join=round,line cap=round] ( 81.31, 82.66) -- ( 85.81, 78.16);

\path[draw=drawColor,line width= 0.4pt,line join=round,line cap=round] ( 81.43, 78.28) -- ( 85.93, 82.78);

\path[draw=drawColor,line width= 0.4pt,line join=round,line cap=round] ( 81.43, 82.78) -- ( 85.93, 78.28);

\path[draw=drawColor,line width= 0.4pt,line join=round,line cap=round] ( 81.55, 78.30) -- ( 86.05, 82.80);

\path[draw=drawColor,line width= 0.4pt,line join=round,line cap=round] ( 81.55, 82.80) -- ( 86.05, 78.30);

\path[draw=drawColor,line width= 0.4pt,line join=round,line cap=round] ( 81.66, 78.43) -- ( 86.16, 82.93);

\path[draw=drawColor,line width= 0.4pt,line join=round,line cap=round] ( 81.66, 82.93) -- ( 86.16, 78.43);

\path[draw=drawColor,line width= 0.4pt,line join=round,line cap=round] ( 81.78, 78.51) -- ( 86.28, 83.01);

\path[draw=drawColor,line width= 0.4pt,line join=round,line cap=round] ( 81.78, 83.01) -- ( 86.28, 78.51);

\path[draw=drawColor,line width= 0.4pt,line join=round,line cap=round] ( 81.90, 78.55) -- ( 86.40, 83.05);

\path[draw=drawColor,line width= 0.4pt,line join=round,line cap=round] ( 81.90, 83.05) -- ( 86.40, 78.55);

\path[draw=drawColor,line width= 0.4pt,line join=round,line cap=round] ( 82.02, 78.65) -- ( 86.52, 83.15);

\path[draw=drawColor,line width= 0.4pt,line join=round,line cap=round] ( 82.02, 83.15) -- ( 86.52, 78.65);

\path[draw=drawColor,line width= 0.4pt,line join=round,line cap=round] ( 82.13, 78.72) -- ( 86.63, 83.22);

\path[draw=drawColor,line width= 0.4pt,line join=round,line cap=round] ( 82.13, 83.22) -- ( 86.63, 78.72);

\path[draw=drawColor,line width= 0.4pt,line join=round,line cap=round] ( 82.25, 78.78) -- ( 86.75, 83.28);

\path[draw=drawColor,line width= 0.4pt,line join=round,line cap=round] ( 82.25, 83.28) -- ( 86.75, 78.78);

\path[draw=drawColor,line width= 0.4pt,line join=round,line cap=round] ( 82.37, 78.83) -- ( 86.87, 83.33);

\path[draw=drawColor,line width= 0.4pt,line join=round,line cap=round] ( 82.37, 83.33) -- ( 86.87, 78.83);

\path[draw=drawColor,line width= 0.4pt,line join=round,line cap=round] ( 82.48, 78.87) -- ( 86.98, 83.37);

\path[draw=drawColor,line width= 0.4pt,line join=round,line cap=round] ( 82.48, 83.37) -- ( 86.98, 78.87);

\path[draw=drawColor,line width= 0.4pt,line join=round,line cap=round] ( 82.60, 79.04) -- ( 87.10, 83.54);

\path[draw=drawColor,line width= 0.4pt,line join=round,line cap=round] ( 82.60, 83.54) -- ( 87.10, 79.04);

\path[draw=drawColor,line width= 0.4pt,line join=round,line cap=round] ( 82.72, 79.06) -- ( 87.22, 83.56);

\path[draw=drawColor,line width= 0.4pt,line join=round,line cap=round] ( 82.72, 83.56) -- ( 87.22, 79.06);

\path[draw=drawColor,line width= 0.4pt,line join=round,line cap=round] ( 82.83, 79.08) -- ( 87.33, 83.58);

\path[draw=drawColor,line width= 0.4pt,line join=round,line cap=round] ( 82.83, 83.58) -- ( 87.33, 79.08);

\path[draw=drawColor,line width= 0.4pt,line join=round,line cap=round] ( 82.95, 79.09) -- ( 87.45, 83.59);

\path[draw=drawColor,line width= 0.4pt,line join=round,line cap=round] ( 82.95, 83.59) -- ( 87.45, 79.09);

\path[draw=drawColor,line width= 0.4pt,line join=round,line cap=round] ( 83.07, 79.15) -- ( 87.57, 83.65);

\path[draw=drawColor,line width= 0.4pt,line join=round,line cap=round] ( 83.07, 83.65) -- ( 87.57, 79.15);

\path[draw=drawColor,line width= 0.4pt,line join=round,line cap=round] ( 83.19, 79.34) -- ( 87.69, 83.84);

\path[draw=drawColor,line width= 0.4pt,line join=round,line cap=round] ( 83.19, 83.84) -- ( 87.69, 79.34);

\path[draw=drawColor,line width= 0.4pt,line join=round,line cap=round] ( 83.30, 79.36) -- ( 87.80, 83.86);

\path[draw=drawColor,line width= 0.4pt,line join=round,line cap=round] ( 83.30, 83.86) -- ( 87.80, 79.36);

\path[draw=drawColor,line width= 0.4pt,line join=round,line cap=round] ( 83.42, 79.41) -- ( 87.92, 83.91);

\path[draw=drawColor,line width= 0.4pt,line join=round,line cap=round] ( 83.42, 83.91) -- ( 87.92, 79.41);

\path[draw=drawColor,line width= 0.4pt,line join=round,line cap=round] ( 83.54, 79.59) -- ( 88.04, 84.09);

\path[draw=drawColor,line width= 0.4pt,line join=round,line cap=round] ( 83.54, 84.09) -- ( 88.04, 79.59);

\path[draw=drawColor,line width= 0.4pt,line join=round,line cap=round] ( 83.65, 79.61) -- ( 88.15, 84.11);

\path[draw=drawColor,line width= 0.4pt,line join=round,line cap=round] ( 83.65, 84.11) -- ( 88.15, 79.61);

\path[draw=drawColor,line width= 0.4pt,line join=round,line cap=round] ( 83.77, 79.70) -- ( 88.27, 84.20);

\path[draw=drawColor,line width= 0.4pt,line join=round,line cap=round] ( 83.77, 84.20) -- ( 88.27, 79.70);

\path[draw=drawColor,line width= 0.4pt,line join=round,line cap=round] ( 83.89, 79.76) -- ( 88.39, 84.26);

\path[draw=drawColor,line width= 0.4pt,line join=round,line cap=round] ( 83.89, 84.26) -- ( 88.39, 79.76);

\path[draw=drawColor,line width= 0.4pt,line join=round,line cap=round] ( 84.00, 79.78) -- ( 88.50, 84.28);

\path[draw=drawColor,line width= 0.4pt,line join=round,line cap=round] ( 84.00, 84.28) -- ( 88.50, 79.78);

\path[draw=drawColor,line width= 0.4pt,line join=round,line cap=round] ( 84.12, 79.88) -- ( 88.62, 84.38);

\path[draw=drawColor,line width= 0.4pt,line join=round,line cap=round] ( 84.12, 84.38) -- ( 88.62, 79.88);

\path[draw=drawColor,line width= 0.4pt,line join=round,line cap=round] ( 84.24, 80.04) -- ( 88.74, 84.54);

\path[draw=drawColor,line width= 0.4pt,line join=round,line cap=round] ( 84.24, 84.54) -- ( 88.74, 80.04);

\path[draw=drawColor,line width= 0.4pt,line join=round,line cap=round] ( 84.35, 80.13) -- ( 88.85, 84.63);

\path[draw=drawColor,line width= 0.4pt,line join=round,line cap=round] ( 84.35, 84.63) -- ( 88.85, 80.13);

\path[draw=drawColor,line width= 0.4pt,line join=round,line cap=round] ( 84.47, 80.15) -- ( 88.97, 84.65);

\path[draw=drawColor,line width= 0.4pt,line join=round,line cap=round] ( 84.47, 84.65) -- ( 88.97, 80.15);

\path[draw=drawColor,line width= 0.4pt,line join=round,line cap=round] ( 84.59, 80.21) -- ( 89.09, 84.71);

\path[draw=drawColor,line width= 0.4pt,line join=round,line cap=round] ( 84.59, 84.71) -- ( 89.09, 80.21);

\path[draw=drawColor,line width= 0.4pt,line join=round,line cap=round] ( 84.71, 80.58) -- ( 89.21, 85.08);

\path[draw=drawColor,line width= 0.4pt,line join=round,line cap=round] ( 84.71, 85.08) -- ( 89.21, 80.58);

\path[draw=drawColor,line width= 0.4pt,line join=round,line cap=round] ( 84.82, 80.68) -- ( 89.32, 85.18);

\path[draw=drawColor,line width= 0.4pt,line join=round,line cap=round] ( 84.82, 85.18) -- ( 89.32, 80.68);

\path[draw=drawColor,line width= 0.4pt,line join=round,line cap=round] ( 84.94, 80.68) -- ( 89.44, 85.18);

\path[draw=drawColor,line width= 0.4pt,line join=round,line cap=round] ( 84.94, 85.18) -- ( 89.44, 80.68);

\path[draw=drawColor,line width= 0.4pt,line join=round,line cap=round] ( 85.06, 80.69) -- ( 89.56, 85.19);

\path[draw=drawColor,line width= 0.4pt,line join=round,line cap=round] ( 85.06, 85.19) -- ( 89.56, 80.69);

\path[draw=drawColor,line width= 0.4pt,line join=round,line cap=round] ( 85.17, 80.70) -- ( 89.67, 85.20);

\path[draw=drawColor,line width= 0.4pt,line join=round,line cap=round] ( 85.17, 85.20) -- ( 89.67, 80.70);

\path[draw=drawColor,line width= 0.4pt,line join=round,line cap=round] ( 85.29, 80.83) -- ( 89.79, 85.33);

\path[draw=drawColor,line width= 0.4pt,line join=round,line cap=round] ( 85.29, 85.33) -- ( 89.79, 80.83);

\path[draw=drawColor,line width= 0.4pt,line join=round,line cap=round] ( 85.41, 80.93) -- ( 89.91, 85.43);

\path[draw=drawColor,line width= 0.4pt,line join=round,line cap=round] ( 85.41, 85.43) -- ( 89.91, 80.93);

\path[draw=drawColor,line width= 0.4pt,line join=round,line cap=round] ( 85.52, 81.25) -- ( 90.02, 85.75);

\path[draw=drawColor,line width= 0.4pt,line join=round,line cap=round] ( 85.52, 85.75) -- ( 90.02, 81.25);

\path[draw=drawColor,line width= 0.4pt,line join=round,line cap=round] ( 85.64, 81.26) -- ( 90.14, 85.76);

\path[draw=drawColor,line width= 0.4pt,line join=round,line cap=round] ( 85.64, 85.76) -- ( 90.14, 81.26);

\path[draw=drawColor,line width= 0.4pt,line join=round,line cap=round] ( 85.76, 81.29) -- ( 90.26, 85.79);

\path[draw=drawColor,line width= 0.4pt,line join=round,line cap=round] ( 85.76, 85.79) -- ( 90.26, 81.29);

\path[draw=drawColor,line width= 0.4pt,line join=round,line cap=round] ( 85.87, 81.34) -- ( 90.37, 85.84);

\path[draw=drawColor,line width= 0.4pt,line join=round,line cap=round] ( 85.87, 85.84) -- ( 90.37, 81.34);

\path[draw=drawColor,line width= 0.4pt,line join=round,line cap=round] ( 85.99, 81.52) -- ( 90.49, 86.02);

\path[draw=drawColor,line width= 0.4pt,line join=round,line cap=round] ( 85.99, 86.02) -- ( 90.49, 81.52);

\path[draw=drawColor,line width= 0.4pt,line join=round,line cap=round] ( 86.11, 81.54) -- ( 90.61, 86.04);

\path[draw=drawColor,line width= 0.4pt,line join=round,line cap=round] ( 86.11, 86.04) -- ( 90.61, 81.54);

\path[draw=drawColor,line width= 0.4pt,line join=round,line cap=round] ( 86.23, 81.56) -- ( 90.73, 86.06);

\path[draw=drawColor,line width= 0.4pt,line join=round,line cap=round] ( 86.23, 86.06) -- ( 90.73, 81.56);

\path[draw=drawColor,line width= 0.4pt,line join=round,line cap=round] ( 86.34, 81.84) -- ( 90.84, 86.34);

\path[draw=drawColor,line width= 0.4pt,line join=round,line cap=round] ( 86.34, 86.34) -- ( 90.84, 81.84);

\path[draw=drawColor,line width= 0.4pt,line join=round,line cap=round] ( 86.46, 81.84) -- ( 90.96, 86.34);

\path[draw=drawColor,line width= 0.4pt,line join=round,line cap=round] ( 86.46, 86.34) -- ( 90.96, 81.84);

\path[draw=drawColor,line width= 0.4pt,line join=round,line cap=round] ( 86.58, 81.87) -- ( 91.08, 86.37);

\path[draw=drawColor,line width= 0.4pt,line join=round,line cap=round] ( 86.58, 86.37) -- ( 91.08, 81.87);

\path[draw=drawColor,line width= 0.4pt,line join=round,line cap=round] ( 86.69, 82.07) -- ( 91.19, 86.57);

\path[draw=drawColor,line width= 0.4pt,line join=round,line cap=round] ( 86.69, 86.57) -- ( 91.19, 82.07);

\path[draw=drawColor,line width= 0.4pt,line join=round,line cap=round] ( 86.81, 82.10) -- ( 91.31, 86.60);

\path[draw=drawColor,line width= 0.4pt,line join=round,line cap=round] ( 86.81, 86.60) -- ( 91.31, 82.10);

\path[draw=drawColor,line width= 0.4pt,line join=round,line cap=round] ( 86.93, 82.21) -- ( 91.43, 86.71);

\path[draw=drawColor,line width= 0.4pt,line join=round,line cap=round] ( 86.93, 86.71) -- ( 91.43, 82.21);

\path[draw=drawColor,line width= 0.4pt,line join=round,line cap=round] ( 87.04, 82.28) -- ( 91.54, 86.78);

\path[draw=drawColor,line width= 0.4pt,line join=round,line cap=round] ( 87.04, 86.78) -- ( 91.54, 82.28);

\path[draw=drawColor,line width= 0.4pt,line join=round,line cap=round] ( 87.16, 82.36) -- ( 91.66, 86.86);

\path[draw=drawColor,line width= 0.4pt,line join=round,line cap=round] ( 87.16, 86.86) -- ( 91.66, 82.36);

\path[draw=drawColor,line width= 0.4pt,line join=round,line cap=round] ( 87.28, 82.45) -- ( 91.78, 86.95);

\path[draw=drawColor,line width= 0.4pt,line join=round,line cap=round] ( 87.28, 86.95) -- ( 91.78, 82.45);

\path[draw=drawColor,line width= 0.4pt,line join=round,line cap=round] ( 87.39, 82.67) -- ( 91.89, 87.17);

\path[draw=drawColor,line width= 0.4pt,line join=round,line cap=round] ( 87.39, 87.17) -- ( 91.89, 82.67);

\path[draw=drawColor,line width= 0.4pt,line join=round,line cap=round] ( 87.51, 82.77) -- ( 92.01, 87.27);

\path[draw=drawColor,line width= 0.4pt,line join=round,line cap=round] ( 87.51, 87.27) -- ( 92.01, 82.77);

\path[draw=drawColor,line width= 0.4pt,line join=round,line cap=round] ( 87.63, 82.93) -- ( 92.13, 87.43);

\path[draw=drawColor,line width= 0.4pt,line join=round,line cap=round] ( 87.63, 87.43) -- ( 92.13, 82.93);

\path[draw=drawColor,line width= 0.4pt,line join=round,line cap=round] ( 87.75, 83.06) -- ( 92.25, 87.56);

\path[draw=drawColor,line width= 0.4pt,line join=round,line cap=round] ( 87.75, 87.56) -- ( 92.25, 83.06);

\path[draw=drawColor,line width= 0.4pt,line join=round,line cap=round] ( 87.86, 83.08) -- ( 92.36, 87.58);

\path[draw=drawColor,line width= 0.4pt,line join=round,line cap=round] ( 87.86, 87.58) -- ( 92.36, 83.08);

\path[draw=drawColor,line width= 0.4pt,line join=round,line cap=round] ( 87.98, 83.21) -- ( 92.48, 87.71);

\path[draw=drawColor,line width= 0.4pt,line join=round,line cap=round] ( 87.98, 87.71) -- ( 92.48, 83.21);

\path[draw=drawColor,line width= 0.4pt,line join=round,line cap=round] ( 88.10, 83.23) -- ( 92.60, 87.73);

\path[draw=drawColor,line width= 0.4pt,line join=round,line cap=round] ( 88.10, 87.73) -- ( 92.60, 83.23);

\path[draw=drawColor,line width= 0.4pt,line join=round,line cap=round] ( 88.21, 83.64) -- ( 92.71, 88.14);

\path[draw=drawColor,line width= 0.4pt,line join=round,line cap=round] ( 88.21, 88.14) -- ( 92.71, 83.64);

\path[draw=drawColor,line width= 0.4pt,line join=round,line cap=round] ( 88.33, 83.69) -- ( 92.83, 88.19);

\path[draw=drawColor,line width= 0.4pt,line join=round,line cap=round] ( 88.33, 88.19) -- ( 92.83, 83.69);

\path[draw=drawColor,line width= 0.4pt,line join=round,line cap=round] ( 88.45, 83.80) -- ( 92.95, 88.30);

\path[draw=drawColor,line width= 0.4pt,line join=round,line cap=round] ( 88.45, 88.30) -- ( 92.95, 83.80);

\path[draw=drawColor,line width= 0.4pt,line join=round,line cap=round] ( 88.56, 83.84) -- ( 93.06, 88.34);

\path[draw=drawColor,line width= 0.4pt,line join=round,line cap=round] ( 88.56, 88.34) -- ( 93.06, 83.84);

\path[draw=drawColor,line width= 0.4pt,line join=round,line cap=round] ( 88.68, 84.13) -- ( 93.18, 88.63);

\path[draw=drawColor,line width= 0.4pt,line join=round,line cap=round] ( 88.68, 88.63) -- ( 93.18, 84.13);

\path[draw=drawColor,line width= 0.4pt,line join=round,line cap=round] ( 88.80, 84.30) -- ( 93.30, 88.80);

\path[draw=drawColor,line width= 0.4pt,line join=round,line cap=round] ( 88.80, 88.80) -- ( 93.30, 84.30);

\path[draw=drawColor,line width= 0.4pt,line join=round,line cap=round] ( 88.92, 84.41) -- ( 93.42, 88.91);

\path[draw=drawColor,line width= 0.4pt,line join=round,line cap=round] ( 88.92, 88.91) -- ( 93.42, 84.41);

\path[draw=drawColor,line width= 0.4pt,line join=round,line cap=round] ( 89.03, 84.44) -- ( 93.53, 88.94);

\path[draw=drawColor,line width= 0.4pt,line join=round,line cap=round] ( 89.03, 88.94) -- ( 93.53, 84.44);

\path[draw=drawColor,line width= 0.4pt,line join=round,line cap=round] ( 89.15, 84.63) -- ( 93.65, 89.13);

\path[draw=drawColor,line width= 0.4pt,line join=round,line cap=round] ( 89.15, 89.13) -- ( 93.65, 84.63);

\path[draw=drawColor,line width= 0.4pt,line join=round,line cap=round] ( 89.27, 84.68) -- ( 93.77, 89.18);

\path[draw=drawColor,line width= 0.4pt,line join=round,line cap=round] ( 89.27, 89.18) -- ( 93.77, 84.68);

\path[draw=drawColor,line width= 0.4pt,line join=round,line cap=round] ( 89.38, 84.69) -- ( 93.88, 89.19);

\path[draw=drawColor,line width= 0.4pt,line join=round,line cap=round] ( 89.38, 89.19) -- ( 93.88, 84.69);

\path[draw=drawColor,line width= 0.4pt,line join=round,line cap=round] ( 89.50, 84.80) -- ( 94.00, 89.30);

\path[draw=drawColor,line width= 0.4pt,line join=round,line cap=round] ( 89.50, 89.30) -- ( 94.00, 84.80);

\path[draw=drawColor,line width= 0.4pt,line join=round,line cap=round] ( 89.62, 84.98) -- ( 94.12, 89.48);

\path[draw=drawColor,line width= 0.4pt,line join=round,line cap=round] ( 89.62, 89.48) -- ( 94.12, 84.98);

\path[draw=drawColor,line width= 0.4pt,line join=round,line cap=round] ( 89.73, 84.99) -- ( 94.23, 89.49);

\path[draw=drawColor,line width= 0.4pt,line join=round,line cap=round] ( 89.73, 89.49) -- ( 94.23, 84.99);

\path[draw=drawColor,line width= 0.4pt,line join=round,line cap=round] ( 89.85, 85.05) -- ( 94.35, 89.55);

\path[draw=drawColor,line width= 0.4pt,line join=round,line cap=round] ( 89.85, 89.55) -- ( 94.35, 85.05);

\path[draw=drawColor,line width= 0.4pt,line join=round,line cap=round] ( 89.97, 85.13) -- ( 94.47, 89.63);

\path[draw=drawColor,line width= 0.4pt,line join=round,line cap=round] ( 89.97, 89.63) -- ( 94.47, 85.13);

\path[draw=drawColor,line width= 0.4pt,line join=round,line cap=round] ( 90.08, 85.16) -- ( 94.58, 89.66);

\path[draw=drawColor,line width= 0.4pt,line join=round,line cap=round] ( 90.08, 89.66) -- ( 94.58, 85.16);

\path[draw=drawColor,line width= 0.4pt,line join=round,line cap=round] ( 90.20, 85.25) -- ( 94.70, 89.75);

\path[draw=drawColor,line width= 0.4pt,line join=round,line cap=round] ( 90.20, 89.75) -- ( 94.70, 85.25);

\path[draw=drawColor,line width= 0.4pt,line join=round,line cap=round] ( 90.32, 85.25) -- ( 94.82, 89.75);

\path[draw=drawColor,line width= 0.4pt,line join=round,line cap=round] ( 90.32, 89.75) -- ( 94.82, 85.25);

\path[draw=drawColor,line width= 0.4pt,line join=round,line cap=round] ( 90.44, 85.33) -- ( 94.94, 89.83);

\path[draw=drawColor,line width= 0.4pt,line join=round,line cap=round] ( 90.44, 89.83) -- ( 94.94, 85.33);

\path[draw=drawColor,line width= 0.4pt,line join=round,line cap=round] ( 90.55, 85.36) -- ( 95.05, 89.86);

\path[draw=drawColor,line width= 0.4pt,line join=round,line cap=round] ( 90.55, 89.86) -- ( 95.05, 85.36);

\path[draw=drawColor,line width= 0.4pt,line join=round,line cap=round] ( 90.67, 85.64) -- ( 95.17, 90.14);

\path[draw=drawColor,line width= 0.4pt,line join=round,line cap=round] ( 90.67, 90.14) -- ( 95.17, 85.64);

\path[draw=drawColor,line width= 0.4pt,line join=round,line cap=round] ( 90.79, 85.66) -- ( 95.29, 90.16);

\path[draw=drawColor,line width= 0.4pt,line join=round,line cap=round] ( 90.79, 90.16) -- ( 95.29, 85.66);

\path[draw=drawColor,line width= 0.4pt,line join=round,line cap=round] ( 90.90, 85.84) -- ( 95.40, 90.34);

\path[draw=drawColor,line width= 0.4pt,line join=round,line cap=round] ( 90.90, 90.34) -- ( 95.40, 85.84);

\path[draw=drawColor,line width= 0.4pt,line join=round,line cap=round] ( 91.02, 85.85) -- ( 95.52, 90.35);

\path[draw=drawColor,line width= 0.4pt,line join=round,line cap=round] ( 91.02, 90.35) -- ( 95.52, 85.85);

\path[draw=drawColor,line width= 0.4pt,line join=round,line cap=round] ( 91.14, 85.89) -- ( 95.64, 90.39);

\path[draw=drawColor,line width= 0.4pt,line join=round,line cap=round] ( 91.14, 90.39) -- ( 95.64, 85.89);

\path[draw=drawColor,line width= 0.4pt,line join=round,line cap=round] ( 91.25, 85.92) -- ( 95.75, 90.42);

\path[draw=drawColor,line width= 0.4pt,line join=round,line cap=round] ( 91.25, 90.42) -- ( 95.75, 85.92);

\path[draw=drawColor,line width= 0.4pt,line join=round,line cap=round] ( 91.37, 86.03) -- ( 95.87, 90.53);

\path[draw=drawColor,line width= 0.4pt,line join=round,line cap=round] ( 91.37, 90.53) -- ( 95.87, 86.03);

\path[draw=drawColor,line width= 0.4pt,line join=round,line cap=round] ( 91.49, 86.12) -- ( 95.99, 90.62);

\path[draw=drawColor,line width= 0.4pt,line join=round,line cap=round] ( 91.49, 90.62) -- ( 95.99, 86.12);

\path[draw=drawColor,line width= 0.4pt,line join=round,line cap=round] ( 91.60, 86.12) -- ( 96.10, 90.62);

\path[draw=drawColor,line width= 0.4pt,line join=round,line cap=round] ( 91.60, 90.62) -- ( 96.10, 86.12);

\path[draw=drawColor,line width= 0.4pt,line join=round,line cap=round] ( 91.72, 86.13) -- ( 96.22, 90.63);

\path[draw=drawColor,line width= 0.4pt,line join=round,line cap=round] ( 91.72, 90.63) -- ( 96.22, 86.13);

\path[draw=drawColor,line width= 0.4pt,line join=round,line cap=round] ( 91.84, 86.29) -- ( 96.34, 90.79);

\path[draw=drawColor,line width= 0.4pt,line join=round,line cap=round] ( 91.84, 90.79) -- ( 96.34, 86.29);

\path[draw=drawColor,line width= 0.4pt,line join=round,line cap=round] ( 91.96, 86.32) -- ( 96.46, 90.82);

\path[draw=drawColor,line width= 0.4pt,line join=round,line cap=round] ( 91.96, 90.82) -- ( 96.46, 86.32);

\path[draw=drawColor,line width= 0.4pt,line join=round,line cap=round] ( 92.07, 86.37) -- ( 96.57, 90.87);

\path[draw=drawColor,line width= 0.4pt,line join=round,line cap=round] ( 92.07, 90.87) -- ( 96.57, 86.37);

\path[draw=drawColor,line width= 0.4pt,line join=round,line cap=round] ( 92.19, 86.42) -- ( 96.69, 90.92);

\path[draw=drawColor,line width= 0.4pt,line join=round,line cap=round] ( 92.19, 90.92) -- ( 96.69, 86.42);

\path[draw=drawColor,line width= 0.4pt,line join=round,line cap=round] ( 92.31, 86.48) -- ( 96.81, 90.98);

\path[draw=drawColor,line width= 0.4pt,line join=round,line cap=round] ( 92.31, 90.98) -- ( 96.81, 86.48);

\path[draw=drawColor,line width= 0.4pt,line join=round,line cap=round] ( 92.42, 86.71) -- ( 96.92, 91.21);

\path[draw=drawColor,line width= 0.4pt,line join=round,line cap=round] ( 92.42, 91.21) -- ( 96.92, 86.71);

\path[draw=drawColor,line width= 0.4pt,line join=round,line cap=round] ( 92.54, 86.77) -- ( 97.04, 91.27);

\path[draw=drawColor,line width= 0.4pt,line join=round,line cap=round] ( 92.54, 91.27) -- ( 97.04, 86.77);

\path[draw=drawColor,line width= 0.4pt,line join=round,line cap=round] ( 92.66, 86.80) -- ( 97.16, 91.30);

\path[draw=drawColor,line width= 0.4pt,line join=round,line cap=round] ( 92.66, 91.30) -- ( 97.16, 86.80);

\path[draw=drawColor,line width= 0.4pt,line join=round,line cap=round] ( 92.77, 86.83) -- ( 97.27, 91.33);

\path[draw=drawColor,line width= 0.4pt,line join=round,line cap=round] ( 92.77, 91.33) -- ( 97.27, 86.83);

\path[draw=drawColor,line width= 0.4pt,line join=round,line cap=round] ( 92.89, 86.87) -- ( 97.39, 91.37);

\path[draw=drawColor,line width= 0.4pt,line join=round,line cap=round] ( 92.89, 91.37) -- ( 97.39, 86.87);

\path[draw=drawColor,line width= 0.4pt,line join=round,line cap=round] ( 93.01, 86.98) -- ( 97.51, 91.48);

\path[draw=drawColor,line width= 0.4pt,line join=round,line cap=round] ( 93.01, 91.48) -- ( 97.51, 86.98);

\path[draw=drawColor,line width= 0.4pt,line join=round,line cap=round] ( 93.13, 87.27) -- ( 97.63, 91.77);

\path[draw=drawColor,line width= 0.4pt,line join=round,line cap=round] ( 93.13, 91.77) -- ( 97.63, 87.27);

\path[draw=drawColor,line width= 0.4pt,line join=round,line cap=round] ( 93.24, 87.38) -- ( 97.74, 91.88);

\path[draw=drawColor,line width= 0.4pt,line join=round,line cap=round] ( 93.24, 91.88) -- ( 97.74, 87.38);

\path[draw=drawColor,line width= 0.4pt,line join=round,line cap=round] ( 93.36, 87.41) -- ( 97.86, 91.91);

\path[draw=drawColor,line width= 0.4pt,line join=round,line cap=round] ( 93.36, 91.91) -- ( 97.86, 87.41);

\path[draw=drawColor,line width= 0.4pt,line join=round,line cap=round] ( 93.48, 87.58) -- ( 97.98, 92.08);

\path[draw=drawColor,line width= 0.4pt,line join=round,line cap=round] ( 93.48, 92.08) -- ( 97.98, 87.58);

\path[draw=drawColor,line width= 0.4pt,line join=round,line cap=round] ( 93.59, 87.67) -- ( 98.09, 92.17);

\path[draw=drawColor,line width= 0.4pt,line join=round,line cap=round] ( 93.59, 92.17) -- ( 98.09, 87.67);

\path[draw=drawColor,line width= 0.4pt,line join=round,line cap=round] ( 93.71, 87.73) -- ( 98.21, 92.23);

\path[draw=drawColor,line width= 0.4pt,line join=round,line cap=round] ( 93.71, 92.23) -- ( 98.21, 87.73);

\path[draw=drawColor,line width= 0.4pt,line join=round,line cap=round] ( 93.83, 87.87) -- ( 98.33, 92.37);

\path[draw=drawColor,line width= 0.4pt,line join=round,line cap=round] ( 93.83, 92.37) -- ( 98.33, 87.87);

\path[draw=drawColor,line width= 0.4pt,line join=round,line cap=round] ( 93.94, 88.01) -- ( 98.44, 92.51);

\path[draw=drawColor,line width= 0.4pt,line join=round,line cap=round] ( 93.94, 92.51) -- ( 98.44, 88.01);

\path[draw=drawColor,line width= 0.4pt,line join=round,line cap=round] ( 94.06, 88.01) -- ( 98.56, 92.51);

\path[draw=drawColor,line width= 0.4pt,line join=round,line cap=round] ( 94.06, 92.51) -- ( 98.56, 88.01);

\path[draw=drawColor,line width= 0.4pt,line join=round,line cap=round] ( 94.18, 88.29) -- ( 98.68, 92.79);

\path[draw=drawColor,line width= 0.4pt,line join=round,line cap=round] ( 94.18, 92.79) -- ( 98.68, 88.29);

\path[draw=drawColor,line width= 0.4pt,line join=round,line cap=round] ( 94.29, 88.40) -- ( 98.79, 92.90);

\path[draw=drawColor,line width= 0.4pt,line join=round,line cap=round] ( 94.29, 92.90) -- ( 98.79, 88.40);

\path[draw=drawColor,line width= 0.4pt,line join=round,line cap=round] ( 94.41, 88.51) -- ( 98.91, 93.01);

\path[draw=drawColor,line width= 0.4pt,line join=round,line cap=round] ( 94.41, 93.01) -- ( 98.91, 88.51);

\path[draw=drawColor,line width= 0.4pt,line join=round,line cap=round] ( 94.53, 88.65) -- ( 99.03, 93.15);

\path[draw=drawColor,line width= 0.4pt,line join=round,line cap=round] ( 94.53, 93.15) -- ( 99.03, 88.65);

\path[draw=drawColor,line width= 0.4pt,line join=round,line cap=round] ( 94.65, 88.68) -- ( 99.15, 93.18);

\path[draw=drawColor,line width= 0.4pt,line join=round,line cap=round] ( 94.65, 93.18) -- ( 99.15, 88.68);

\path[draw=drawColor,line width= 0.4pt,line join=round,line cap=round] ( 94.76, 88.74) -- ( 99.26, 93.24);

\path[draw=drawColor,line width= 0.4pt,line join=round,line cap=round] ( 94.76, 93.24) -- ( 99.26, 88.74);

\path[draw=drawColor,line width= 0.4pt,line join=round,line cap=round] ( 94.88, 88.77) -- ( 99.38, 93.27);

\path[draw=drawColor,line width= 0.4pt,line join=round,line cap=round] ( 94.88, 93.27) -- ( 99.38, 88.77);

\path[draw=drawColor,line width= 0.4pt,line join=round,line cap=round] ( 95.00, 88.84) -- ( 99.50, 93.34);

\path[draw=drawColor,line width= 0.4pt,line join=round,line cap=round] ( 95.00, 93.34) -- ( 99.50, 88.84);

\path[draw=drawColor,line width= 0.4pt,line join=round,line cap=round] ( 95.11, 88.87) -- ( 99.61, 93.37);

\path[draw=drawColor,line width= 0.4pt,line join=round,line cap=round] ( 95.11, 93.37) -- ( 99.61, 88.87);

\path[draw=drawColor,line width= 0.4pt,line join=round,line cap=round] ( 95.23, 89.05) -- ( 99.73, 93.55);

\path[draw=drawColor,line width= 0.4pt,line join=round,line cap=round] ( 95.23, 93.55) -- ( 99.73, 89.05);

\path[draw=drawColor,line width= 0.4pt,line join=round,line cap=round] ( 95.35, 89.19) -- ( 99.85, 93.69);

\path[draw=drawColor,line width= 0.4pt,line join=round,line cap=round] ( 95.35, 93.69) -- ( 99.85, 89.19);

\path[draw=drawColor,line width= 0.4pt,line join=round,line cap=round] ( 95.46, 89.21) -- ( 99.96, 93.71);

\path[draw=drawColor,line width= 0.4pt,line join=round,line cap=round] ( 95.46, 93.71) -- ( 99.96, 89.21);

\path[draw=drawColor,line width= 0.4pt,line join=round,line cap=round] ( 95.58, 89.51) -- (100.08, 94.01);

\path[draw=drawColor,line width= 0.4pt,line join=round,line cap=round] ( 95.58, 94.01) -- (100.08, 89.51);

\path[draw=drawColor,line width= 0.4pt,line join=round,line cap=round] ( 95.70, 89.62) -- (100.20, 94.12);

\path[draw=drawColor,line width= 0.4pt,line join=round,line cap=round] ( 95.70, 94.12) -- (100.20, 89.62);

\path[draw=drawColor,line width= 0.4pt,line join=round,line cap=round] ( 95.81, 89.63) -- (100.31, 94.13);

\path[draw=drawColor,line width= 0.4pt,line join=round,line cap=round] ( 95.81, 94.13) -- (100.31, 89.63);

\path[draw=drawColor,line width= 0.4pt,line join=round,line cap=round] ( 95.93, 89.73) -- (100.43, 94.23);

\path[draw=drawColor,line width= 0.4pt,line join=round,line cap=round] ( 95.93, 94.23) -- (100.43, 89.73);

\path[draw=drawColor,line width= 0.4pt,line join=round,line cap=round] ( 96.05, 89.91) -- (100.55, 94.41);

\path[draw=drawColor,line width= 0.4pt,line join=round,line cap=round] ( 96.05, 94.41) -- (100.55, 89.91);

\path[draw=drawColor,line width= 0.4pt,line join=round,line cap=round] ( 96.17, 90.06) -- (100.67, 94.56);

\path[draw=drawColor,line width= 0.4pt,line join=round,line cap=round] ( 96.17, 94.56) -- (100.67, 90.06);

\path[draw=drawColor,line width= 0.4pt,line join=round,line cap=round] ( 96.28, 90.16) -- (100.78, 94.66);

\path[draw=drawColor,line width= 0.4pt,line join=round,line cap=round] ( 96.28, 94.66) -- (100.78, 90.16);

\path[draw=drawColor,line width= 0.4pt,line join=round,line cap=round] ( 96.40, 90.27) -- (100.90, 94.77);

\path[draw=drawColor,line width= 0.4pt,line join=round,line cap=round] ( 96.40, 94.77) -- (100.90, 90.27);

\path[draw=drawColor,line width= 0.4pt,line join=round,line cap=round] ( 96.52, 90.45) -- (101.02, 94.95);

\path[draw=drawColor,line width= 0.4pt,line join=round,line cap=round] ( 96.52, 94.95) -- (101.02, 90.45);

\path[draw=drawColor,line width= 0.4pt,line join=round,line cap=round] ( 96.63, 90.45) -- (101.13, 94.95);

\path[draw=drawColor,line width= 0.4pt,line join=round,line cap=round] ( 96.63, 94.95) -- (101.13, 90.45);

\path[draw=drawColor,line width= 0.4pt,line join=round,line cap=round] ( 96.75, 90.62) -- (101.25, 95.12);

\path[draw=drawColor,line width= 0.4pt,line join=round,line cap=round] ( 96.75, 95.12) -- (101.25, 90.62);

\path[draw=drawColor,line width= 0.4pt,line join=round,line cap=round] ( 96.87, 90.67) -- (101.37, 95.17);

\path[draw=drawColor,line width= 0.4pt,line join=round,line cap=round] ( 96.87, 95.17) -- (101.37, 90.67);

\path[draw=drawColor,line width= 0.4pt,line join=round,line cap=round] ( 96.98, 90.78) -- (101.48, 95.28);

\path[draw=drawColor,line width= 0.4pt,line join=round,line cap=round] ( 96.98, 95.28) -- (101.48, 90.78);

\path[draw=drawColor,line width= 0.4pt,line join=round,line cap=round] ( 97.10, 90.96) -- (101.60, 95.46);

\path[draw=drawColor,line width= 0.4pt,line join=round,line cap=round] ( 97.10, 95.46) -- (101.60, 90.96);

\path[draw=drawColor,line width= 0.4pt,line join=round,line cap=round] ( 97.22, 91.02) -- (101.72, 95.52);

\path[draw=drawColor,line width= 0.4pt,line join=round,line cap=round] ( 97.22, 95.52) -- (101.72, 91.02);

\path[draw=drawColor,line width= 0.4pt,line join=round,line cap=round] ( 97.33, 91.04) -- (101.83, 95.54);

\path[draw=drawColor,line width= 0.4pt,line join=round,line cap=round] ( 97.33, 95.54) -- (101.83, 91.04);

\path[draw=drawColor,line width= 0.4pt,line join=round,line cap=round] ( 97.45, 91.06) -- (101.95, 95.56);

\path[draw=drawColor,line width= 0.4pt,line join=round,line cap=round] ( 97.45, 95.56) -- (101.95, 91.06);

\path[draw=drawColor,line width= 0.4pt,line join=round,line cap=round] ( 97.57, 91.14) -- (102.07, 95.64);

\path[draw=drawColor,line width= 0.4pt,line join=round,line cap=round] ( 97.57, 95.64) -- (102.07, 91.14);

\path[draw=drawColor,line width= 0.4pt,line join=round,line cap=round] ( 97.69, 91.15) -- (102.19, 95.65);

\path[draw=drawColor,line width= 0.4pt,line join=round,line cap=round] ( 97.69, 95.65) -- (102.19, 91.15);

\path[draw=drawColor,line width= 0.4pt,line join=round,line cap=round] ( 97.80, 91.20) -- (102.30, 95.70);

\path[draw=drawColor,line width= 0.4pt,line join=round,line cap=round] ( 97.80, 95.70) -- (102.30, 91.20);

\path[draw=drawColor,line width= 0.4pt,line join=round,line cap=round] ( 97.92, 91.20) -- (102.42, 95.70);

\path[draw=drawColor,line width= 0.4pt,line join=round,line cap=round] ( 97.92, 95.70) -- (102.42, 91.20);

\path[draw=drawColor,line width= 0.4pt,line join=round,line cap=round] ( 98.04, 91.24) -- (102.54, 95.74);

\path[draw=drawColor,line width= 0.4pt,line join=round,line cap=round] ( 98.04, 95.74) -- (102.54, 91.24);

\path[draw=drawColor,line width= 0.4pt,line join=round,line cap=round] ( 98.15, 91.37) -- (102.65, 95.87);

\path[draw=drawColor,line width= 0.4pt,line join=round,line cap=round] ( 98.15, 95.87) -- (102.65, 91.37);

\path[draw=drawColor,line width= 0.4pt,line join=round,line cap=round] ( 98.27, 91.39) -- (102.77, 95.89);

\path[draw=drawColor,line width= 0.4pt,line join=round,line cap=round] ( 98.27, 95.89) -- (102.77, 91.39);

\path[draw=drawColor,line width= 0.4pt,line join=round,line cap=round] ( 98.39, 91.46) -- (102.89, 95.96);

\path[draw=drawColor,line width= 0.4pt,line join=round,line cap=round] ( 98.39, 95.96) -- (102.89, 91.46);

\path[draw=drawColor,line width= 0.4pt,line join=round,line cap=round] ( 98.50, 91.62) -- (103.00, 96.12);

\path[draw=drawColor,line width= 0.4pt,line join=round,line cap=round] ( 98.50, 96.12) -- (103.00, 91.62);

\path[draw=drawColor,line width= 0.4pt,line join=round,line cap=round] ( 98.62, 91.78) -- (103.12, 96.28);

\path[draw=drawColor,line width= 0.4pt,line join=round,line cap=round] ( 98.62, 96.28) -- (103.12, 91.78);

\path[draw=drawColor,line width= 0.4pt,line join=round,line cap=round] ( 98.74, 91.80) -- (103.24, 96.30);

\path[draw=drawColor,line width= 0.4pt,line join=round,line cap=round] ( 98.74, 96.30) -- (103.24, 91.80);

\path[draw=drawColor,line width= 0.4pt,line join=round,line cap=round] ( 98.86, 92.00) -- (103.36, 96.50);

\path[draw=drawColor,line width= 0.4pt,line join=round,line cap=round] ( 98.86, 96.50) -- (103.36, 92.00);

\path[draw=drawColor,line width= 0.4pt,line join=round,line cap=round] ( 98.97, 92.35) -- (103.47, 96.85);

\path[draw=drawColor,line width= 0.4pt,line join=round,line cap=round] ( 98.97, 96.85) -- (103.47, 92.35);

\path[draw=drawColor,line width= 0.4pt,line join=round,line cap=round] ( 99.09, 92.48) -- (103.59, 96.98);

\path[draw=drawColor,line width= 0.4pt,line join=round,line cap=round] ( 99.09, 96.98) -- (103.59, 92.48);

\path[draw=drawColor,line width= 0.4pt,line join=round,line cap=round] ( 99.21, 92.49) -- (103.71, 96.99);

\path[draw=drawColor,line width= 0.4pt,line join=round,line cap=round] ( 99.21, 96.99) -- (103.71, 92.49);

\path[draw=drawColor,line width= 0.4pt,line join=round,line cap=round] ( 99.32, 92.56) -- (103.82, 97.06);

\path[draw=drawColor,line width= 0.4pt,line join=round,line cap=round] ( 99.32, 97.06) -- (103.82, 92.56);

\path[draw=drawColor,line width= 0.4pt,line join=round,line cap=round] ( 99.44, 92.56) -- (103.94, 97.06);

\path[draw=drawColor,line width= 0.4pt,line join=round,line cap=round] ( 99.44, 97.06) -- (103.94, 92.56);

\path[draw=drawColor,line width= 0.4pt,line join=round,line cap=round] ( 99.56, 92.58) -- (104.06, 97.08);

\path[draw=drawColor,line width= 0.4pt,line join=round,line cap=round] ( 99.56, 97.08) -- (104.06, 92.58);

\path[draw=drawColor,line width= 0.4pt,line join=round,line cap=round] ( 99.67, 92.62) -- (104.17, 97.12);

\path[draw=drawColor,line width= 0.4pt,line join=round,line cap=round] ( 99.67, 97.12) -- (104.17, 92.62);

\path[draw=drawColor,line width= 0.4pt,line join=round,line cap=round] ( 99.79, 92.71) -- (104.29, 97.21);

\path[draw=drawColor,line width= 0.4pt,line join=round,line cap=round] ( 99.79, 97.21) -- (104.29, 92.71);

\path[draw=drawColor,line width= 0.4pt,line join=round,line cap=round] ( 99.91, 92.86) -- (104.41, 97.36);

\path[draw=drawColor,line width= 0.4pt,line join=round,line cap=round] ( 99.91, 97.36) -- (104.41, 92.86);

\path[draw=drawColor,line width= 0.4pt,line join=round,line cap=round] (100.02, 93.28) -- (104.52, 97.78);

\path[draw=drawColor,line width= 0.4pt,line join=round,line cap=round] (100.02, 97.78) -- (104.52, 93.28);

\path[draw=drawColor,line width= 0.4pt,line join=round,line cap=round] (100.14, 93.29) -- (104.64, 97.79);

\path[draw=drawColor,line width= 0.4pt,line join=round,line cap=round] (100.14, 97.79) -- (104.64, 93.29);

\path[draw=drawColor,line width= 0.4pt,line join=round,line cap=round] (100.26, 93.37) -- (104.76, 97.87);

\path[draw=drawColor,line width= 0.4pt,line join=round,line cap=round] (100.26, 97.87) -- (104.76, 93.37);

\path[draw=drawColor,line width= 0.4pt,line join=round,line cap=round] (100.38, 93.68) -- (104.88, 98.18);

\path[draw=drawColor,line width= 0.4pt,line join=round,line cap=round] (100.38, 98.18) -- (104.88, 93.68);

\path[draw=drawColor,line width= 0.4pt,line join=round,line cap=round] (100.49, 93.78) -- (104.99, 98.28);

\path[draw=drawColor,line width= 0.4pt,line join=round,line cap=round] (100.49, 98.28) -- (104.99, 93.78);

\path[draw=drawColor,line width= 0.4pt,line join=round,line cap=round] (100.61, 94.01) -- (105.11, 98.51);

\path[draw=drawColor,line width= 0.4pt,line join=round,line cap=round] (100.61, 98.51) -- (105.11, 94.01);

\path[draw=drawColor,line width= 0.4pt,line join=round,line cap=round] (100.73, 94.01) -- (105.23, 98.51);

\path[draw=drawColor,line width= 0.4pt,line join=round,line cap=round] (100.73, 98.51) -- (105.23, 94.01);

\path[draw=drawColor,line width= 0.4pt,line join=round,line cap=round] (100.84, 94.08) -- (105.34, 98.58);

\path[draw=drawColor,line width= 0.4pt,line join=round,line cap=round] (100.84, 98.58) -- (105.34, 94.08);

\path[draw=drawColor,line width= 0.4pt,line join=round,line cap=round] (100.96, 94.25) -- (105.46, 98.75);

\path[draw=drawColor,line width= 0.4pt,line join=round,line cap=round] (100.96, 98.75) -- (105.46, 94.25);

\path[draw=drawColor,line width= 0.4pt,line join=round,line cap=round] (101.08, 94.27) -- (105.58, 98.77);

\path[draw=drawColor,line width= 0.4pt,line join=round,line cap=round] (101.08, 98.77) -- (105.58, 94.27);

\path[draw=drawColor,line width= 0.4pt,line join=round,line cap=round] (101.19, 94.27) -- (105.69, 98.77);

\path[draw=drawColor,line width= 0.4pt,line join=round,line cap=round] (101.19, 98.77) -- (105.69, 94.27);

\path[draw=drawColor,line width= 0.4pt,line join=round,line cap=round] (101.31, 94.33) -- (105.81, 98.83);

\path[draw=drawColor,line width= 0.4pt,line join=round,line cap=round] (101.31, 98.83) -- (105.81, 94.33);

\path[draw=drawColor,line width= 0.4pt,line join=round,line cap=round] (101.43, 94.55) -- (105.93, 99.05);

\path[draw=drawColor,line width= 0.4pt,line join=round,line cap=round] (101.43, 99.05) -- (105.93, 94.55);

\path[draw=drawColor,line width= 0.4pt,line join=round,line cap=round] (101.54, 94.71) -- (106.04, 99.21);

\path[draw=drawColor,line width= 0.4pt,line join=round,line cap=round] (101.54, 99.21) -- (106.04, 94.71);

\path[draw=drawColor,line width= 0.4pt,line join=round,line cap=round] (101.66, 94.72) -- (106.16, 99.22);

\path[draw=drawColor,line width= 0.4pt,line join=round,line cap=round] (101.66, 99.22) -- (106.16, 94.72);

\path[draw=drawColor,line width= 0.4pt,line join=round,line cap=round] (101.78, 94.75) -- (106.28, 99.25);

\path[draw=drawColor,line width= 0.4pt,line join=round,line cap=round] (101.78, 99.25) -- (106.28, 94.75);

\path[draw=drawColor,line width= 0.4pt,line join=round,line cap=round] (101.90, 94.86) -- (106.40, 99.36);

\path[draw=drawColor,line width= 0.4pt,line join=round,line cap=round] (101.90, 99.36) -- (106.40, 94.86);

\path[draw=drawColor,line width= 0.4pt,line join=round,line cap=round] (102.01, 94.94) -- (106.51, 99.44);

\path[draw=drawColor,line width= 0.4pt,line join=round,line cap=round] (102.01, 99.44) -- (106.51, 94.94);

\path[draw=drawColor,line width= 0.4pt,line join=round,line cap=round] (102.13, 94.95) -- (106.63, 99.45);

\path[draw=drawColor,line width= 0.4pt,line join=round,line cap=round] (102.13, 99.45) -- (106.63, 94.95);

\path[draw=drawColor,line width= 0.4pt,line join=round,line cap=round] (102.25, 94.98) -- (106.75, 99.48);

\path[draw=drawColor,line width= 0.4pt,line join=round,line cap=round] (102.25, 99.48) -- (106.75, 94.98);

\path[draw=drawColor,line width= 0.4pt,line join=round,line cap=round] (102.36, 95.14) -- (106.86, 99.64);

\path[draw=drawColor,line width= 0.4pt,line join=round,line cap=round] (102.36, 99.64) -- (106.86, 95.14);

\path[draw=drawColor,line width= 0.4pt,line join=round,line cap=round] (102.48, 95.20) -- (106.98, 99.70);

\path[draw=drawColor,line width= 0.4pt,line join=round,line cap=round] (102.48, 99.70) -- (106.98, 95.20);

\path[draw=drawColor,line width= 0.4pt,line join=round,line cap=round] (102.60, 95.28) -- (107.10, 99.78);

\path[draw=drawColor,line width= 0.4pt,line join=round,line cap=round] (102.60, 99.78) -- (107.10, 95.28);

\path[draw=drawColor,line width= 0.4pt,line join=round,line cap=round] (102.71, 95.36) -- (107.21, 99.86);

\path[draw=drawColor,line width= 0.4pt,line join=round,line cap=round] (102.71, 99.86) -- (107.21, 95.36);

\path[draw=drawColor,line width= 0.4pt,line join=round,line cap=round] (102.83, 95.39) -- (107.33, 99.89);

\path[draw=drawColor,line width= 0.4pt,line join=round,line cap=round] (102.83, 99.89) -- (107.33, 95.39);

\path[draw=drawColor,line width= 0.4pt,line join=round,line cap=round] (102.95, 95.44) -- (107.45, 99.94);

\path[draw=drawColor,line width= 0.4pt,line join=round,line cap=round] (102.95, 99.94) -- (107.45, 95.44);

\path[draw=drawColor,line width= 0.4pt,line join=round,line cap=round] (103.07, 95.53) -- (107.57,100.03);

\path[draw=drawColor,line width= 0.4pt,line join=round,line cap=round] (103.07,100.03) -- (107.57, 95.53);

\path[draw=drawColor,line width= 0.4pt,line join=round,line cap=round] (103.18, 95.78) -- (107.68,100.28);

\path[draw=drawColor,line width= 0.4pt,line join=round,line cap=round] (103.18,100.28) -- (107.68, 95.78);

\path[draw=drawColor,line width= 0.4pt,line join=round,line cap=round] (103.30, 95.89) -- (107.80,100.39);

\path[draw=drawColor,line width= 0.4pt,line join=round,line cap=round] (103.30,100.39) -- (107.80, 95.89);

\path[draw=drawColor,line width= 0.4pt,line join=round,line cap=round] (103.42, 96.02) -- (107.92,100.52);

\path[draw=drawColor,line width= 0.4pt,line join=round,line cap=round] (103.42,100.52) -- (107.92, 96.02);

\path[draw=drawColor,line width= 0.4pt,line join=round,line cap=round] (103.53, 96.06) -- (108.03,100.56);

\path[draw=drawColor,line width= 0.4pt,line join=round,line cap=round] (103.53,100.56) -- (108.03, 96.06);

\path[draw=drawColor,line width= 0.4pt,line join=round,line cap=round] (103.65, 96.09) -- (108.15,100.59);

\path[draw=drawColor,line width= 0.4pt,line join=round,line cap=round] (103.65,100.59) -- (108.15, 96.09);

\path[draw=drawColor,line width= 0.4pt,line join=round,line cap=round] (103.77, 96.14) -- (108.27,100.64);

\path[draw=drawColor,line width= 0.4pt,line join=round,line cap=round] (103.77,100.64) -- (108.27, 96.14);

\path[draw=drawColor,line width= 0.4pt,line join=round,line cap=round] (103.88, 96.19) -- (108.38,100.69);

\path[draw=drawColor,line width= 0.4pt,line join=round,line cap=round] (103.88,100.69) -- (108.38, 96.19);

\path[draw=drawColor,line width= 0.4pt,line join=round,line cap=round] (104.00, 96.21) -- (108.50,100.71);

\path[draw=drawColor,line width= 0.4pt,line join=round,line cap=round] (104.00,100.71) -- (108.50, 96.21);

\path[draw=drawColor,line width= 0.4pt,line join=round,line cap=round] (104.12, 96.27) -- (108.62,100.77);

\path[draw=drawColor,line width= 0.4pt,line join=round,line cap=round] (104.12,100.77) -- (108.62, 96.27);

\path[draw=drawColor,line width= 0.4pt,line join=round,line cap=round] (104.23, 96.45) -- (108.73,100.95);

\path[draw=drawColor,line width= 0.4pt,line join=round,line cap=round] (104.23,100.95) -- (108.73, 96.45);

\path[draw=drawColor,line width= 0.4pt,line join=round,line cap=round] (104.35, 96.50) -- (108.85,101.00);

\path[draw=drawColor,line width= 0.4pt,line join=round,line cap=round] (104.35,101.00) -- (108.85, 96.50);

\path[draw=drawColor,line width= 0.4pt,line join=round,line cap=round] (104.47, 96.55) -- (108.97,101.05);

\path[draw=drawColor,line width= 0.4pt,line join=round,line cap=round] (104.47,101.05) -- (108.97, 96.55);

\path[draw=drawColor,line width= 0.4pt,line join=round,line cap=round] (104.59, 96.69) -- (109.09,101.19);

\path[draw=drawColor,line width= 0.4pt,line join=round,line cap=round] (104.59,101.19) -- (109.09, 96.69);

\path[draw=drawColor,line width= 0.4pt,line join=round,line cap=round] (104.70, 96.70) -- (109.20,101.20);

\path[draw=drawColor,line width= 0.4pt,line join=round,line cap=round] (104.70,101.20) -- (109.20, 96.70);

\path[draw=drawColor,line width= 0.4pt,line join=round,line cap=round] (104.82, 96.70) -- (109.32,101.20);

\path[draw=drawColor,line width= 0.4pt,line join=round,line cap=round] (104.82,101.20) -- (109.32, 96.70);

\path[draw=drawColor,line width= 0.4pt,line join=round,line cap=round] (104.94, 96.75) -- (109.44,101.25);

\path[draw=drawColor,line width= 0.4pt,line join=round,line cap=round] (104.94,101.25) -- (109.44, 96.75);

\path[draw=drawColor,line width= 0.4pt,line join=round,line cap=round] (105.05, 96.76) -- (109.55,101.26);

\path[draw=drawColor,line width= 0.4pt,line join=round,line cap=round] (105.05,101.26) -- (109.55, 96.76);

\path[draw=drawColor,line width= 0.4pt,line join=round,line cap=round] (105.17, 96.88) -- (109.67,101.38);

\path[draw=drawColor,line width= 0.4pt,line join=round,line cap=round] (105.17,101.38) -- (109.67, 96.88);

\path[draw=drawColor,line width= 0.4pt,line join=round,line cap=round] (105.29, 96.91) -- (109.79,101.41);

\path[draw=drawColor,line width= 0.4pt,line join=round,line cap=round] (105.29,101.41) -- (109.79, 96.91);

\path[draw=drawColor,line width= 0.4pt,line join=round,line cap=round] (105.40, 96.96) -- (109.90,101.46);

\path[draw=drawColor,line width= 0.4pt,line join=round,line cap=round] (105.40,101.46) -- (109.90, 96.96);

\path[draw=drawColor,line width= 0.4pt,line join=round,line cap=round] (105.52, 97.23) -- (110.02,101.73);

\path[draw=drawColor,line width= 0.4pt,line join=round,line cap=round] (105.52,101.73) -- (110.02, 97.23);

\path[draw=drawColor,line width= 0.4pt,line join=round,line cap=round] (105.64, 97.23) -- (110.14,101.73);

\path[draw=drawColor,line width= 0.4pt,line join=round,line cap=round] (105.64,101.73) -- (110.14, 97.23);

\path[draw=drawColor,line width= 0.4pt,line join=round,line cap=round] (105.75, 97.27) -- (110.25,101.77);

\path[draw=drawColor,line width= 0.4pt,line join=round,line cap=round] (105.75,101.77) -- (110.25, 97.27);

\path[draw=drawColor,line width= 0.4pt,line join=round,line cap=round] (105.87, 97.45) -- (110.37,101.95);

\path[draw=drawColor,line width= 0.4pt,line join=round,line cap=round] (105.87,101.95) -- (110.37, 97.45);

\path[draw=drawColor,line width= 0.4pt,line join=round,line cap=round] (105.99, 97.53) -- (110.49,102.03);

\path[draw=drawColor,line width= 0.4pt,line join=round,line cap=round] (105.99,102.03) -- (110.49, 97.53);

\path[draw=drawColor,line width= 0.4pt,line join=round,line cap=round] (106.11, 97.58) -- (110.61,102.08);

\path[draw=drawColor,line width= 0.4pt,line join=round,line cap=round] (106.11,102.08) -- (110.61, 97.58);

\path[draw=drawColor,line width= 0.4pt,line join=round,line cap=round] (106.22, 97.68) -- (110.72,102.18);

\path[draw=drawColor,line width= 0.4pt,line join=round,line cap=round] (106.22,102.18) -- (110.72, 97.68);

\path[draw=drawColor,line width= 0.4pt,line join=round,line cap=round] (106.34, 97.89) -- (110.84,102.39);

\path[draw=drawColor,line width= 0.4pt,line join=round,line cap=round] (106.34,102.39) -- (110.84, 97.89);

\path[draw=drawColor,line width= 0.4pt,line join=round,line cap=round] (106.46, 97.94) -- (110.96,102.44);

\path[draw=drawColor,line width= 0.4pt,line join=round,line cap=round] (106.46,102.44) -- (110.96, 97.94);

\path[draw=drawColor,line width= 0.4pt,line join=round,line cap=round] (106.57, 97.97) -- (111.07,102.47);

\path[draw=drawColor,line width= 0.4pt,line join=round,line cap=round] (106.57,102.47) -- (111.07, 97.97);

\path[draw=drawColor,line width= 0.4pt,line join=round,line cap=round] (106.69, 97.99) -- (111.19,102.49);

\path[draw=drawColor,line width= 0.4pt,line join=round,line cap=round] (106.69,102.49) -- (111.19, 97.99);

\path[draw=drawColor,line width= 0.4pt,line join=round,line cap=round] (106.81, 98.00) -- (111.31,102.50);

\path[draw=drawColor,line width= 0.4pt,line join=round,line cap=round] (106.81,102.50) -- (111.31, 98.00);

\path[draw=drawColor,line width= 0.4pt,line join=round,line cap=round] (106.92, 98.03) -- (111.42,102.53);

\path[draw=drawColor,line width= 0.4pt,line join=round,line cap=round] (106.92,102.53) -- (111.42, 98.03);

\path[draw=drawColor,line width= 0.4pt,line join=round,line cap=round] (107.04, 98.28) -- (111.54,102.78);

\path[draw=drawColor,line width= 0.4pt,line join=round,line cap=round] (107.04,102.78) -- (111.54, 98.28);

\path[draw=drawColor,line width= 0.4pt,line join=round,line cap=round] (107.16, 98.33) -- (111.66,102.83);

\path[draw=drawColor,line width= 0.4pt,line join=round,line cap=round] (107.16,102.83) -- (111.66, 98.33);

\path[draw=drawColor,line width= 0.4pt,line join=round,line cap=round] (107.27, 98.37) -- (111.77,102.87);

\path[draw=drawColor,line width= 0.4pt,line join=round,line cap=round] (107.27,102.87) -- (111.77, 98.37);

\path[draw=drawColor,line width= 0.4pt,line join=round,line cap=round] (107.39, 98.71) -- (111.89,103.21);

\path[draw=drawColor,line width= 0.4pt,line join=round,line cap=round] (107.39,103.21) -- (111.89, 98.71);

\path[draw=drawColor,line width= 0.4pt,line join=round,line cap=round] (107.51, 98.71) -- (112.01,103.21);

\path[draw=drawColor,line width= 0.4pt,line join=round,line cap=round] (107.51,103.21) -- (112.01, 98.71);

\path[draw=drawColor,line width= 0.4pt,line join=round,line cap=round] (107.63, 98.72) -- (112.13,103.22);

\path[draw=drawColor,line width= 0.4pt,line join=round,line cap=round] (107.63,103.22) -- (112.13, 98.72);

\path[draw=drawColor,line width= 0.4pt,line join=round,line cap=round] (107.74, 98.76) -- (112.24,103.26);

\path[draw=drawColor,line width= 0.4pt,line join=round,line cap=round] (107.74,103.26) -- (112.24, 98.76);

\path[draw=drawColor,line width= 0.4pt,line join=round,line cap=round] (107.86, 98.79) -- (112.36,103.29);

\path[draw=drawColor,line width= 0.4pt,line join=round,line cap=round] (107.86,103.29) -- (112.36, 98.79);

\path[draw=drawColor,line width= 0.4pt,line join=round,line cap=round] (107.98, 98.87) -- (112.48,103.37);

\path[draw=drawColor,line width= 0.4pt,line join=round,line cap=round] (107.98,103.37) -- (112.48, 98.87);

\path[draw=drawColor,line width= 0.4pt,line join=round,line cap=round] (108.09, 99.15) -- (112.59,103.65);

\path[draw=drawColor,line width= 0.4pt,line join=round,line cap=round] (108.09,103.65) -- (112.59, 99.15);

\path[draw=drawColor,line width= 0.4pt,line join=round,line cap=round] (108.21, 99.15) -- (112.71,103.65);

\path[draw=drawColor,line width= 0.4pt,line join=round,line cap=round] (108.21,103.65) -- (112.71, 99.15);

\path[draw=drawColor,line width= 0.4pt,line join=round,line cap=round] (108.33, 99.39) -- (112.83,103.89);

\path[draw=drawColor,line width= 0.4pt,line join=round,line cap=round] (108.33,103.89) -- (112.83, 99.39);

\path[draw=drawColor,line width= 0.4pt,line join=round,line cap=round] (108.44, 99.39) -- (112.94,103.89);

\path[draw=drawColor,line width= 0.4pt,line join=round,line cap=round] (108.44,103.89) -- (112.94, 99.39);

\path[draw=drawColor,line width= 0.4pt,line join=round,line cap=round] (108.56, 99.55) -- (113.06,104.05);

\path[draw=drawColor,line width= 0.4pt,line join=round,line cap=round] (108.56,104.05) -- (113.06, 99.55);

\path[draw=drawColor,line width= 0.4pt,line join=round,line cap=round] (108.68, 99.66) -- (113.18,104.16);

\path[draw=drawColor,line width= 0.4pt,line join=round,line cap=round] (108.68,104.16) -- (113.18, 99.66);

\path[draw=drawColor,line width= 0.4pt,line join=round,line cap=round] (108.80, 99.89) -- (113.30,104.39);

\path[draw=drawColor,line width= 0.4pt,line join=round,line cap=round] (108.80,104.39) -- (113.30, 99.89);

\path[draw=drawColor,line width= 0.4pt,line join=round,line cap=round] (108.91, 99.91) -- (113.41,104.41);

\path[draw=drawColor,line width= 0.4pt,line join=round,line cap=round] (108.91,104.41) -- (113.41, 99.91);

\path[draw=drawColor,line width= 0.4pt,line join=round,line cap=round] (109.03, 99.94) -- (113.53,104.44);

\path[draw=drawColor,line width= 0.4pt,line join=round,line cap=round] (109.03,104.44) -- (113.53, 99.94);

\path[draw=drawColor,line width= 0.4pt,line join=round,line cap=round] (109.15, 99.95) -- (113.65,104.45);

\path[draw=drawColor,line width= 0.4pt,line join=round,line cap=round] (109.15,104.45) -- (113.65, 99.95);

\path[draw=drawColor,line width= 0.4pt,line join=round,line cap=round] (109.26,100.05) -- (113.76,104.55);

\path[draw=drawColor,line width= 0.4pt,line join=round,line cap=round] (109.26,104.55) -- (113.76,100.05);

\path[draw=drawColor,line width= 0.4pt,line join=round,line cap=round] (109.38,100.25) -- (113.88,104.75);

\path[draw=drawColor,line width= 0.4pt,line join=round,line cap=round] (109.38,104.75) -- (113.88,100.25);

\path[draw=drawColor,line width= 0.4pt,line join=round,line cap=round] (109.50,100.27) -- (114.00,104.77);

\path[draw=drawColor,line width= 0.4pt,line join=round,line cap=round] (109.50,104.77) -- (114.00,100.27);

\path[draw=drawColor,line width= 0.4pt,line join=round,line cap=round] (109.61,100.35) -- (114.11,104.85);

\path[draw=drawColor,line width= 0.4pt,line join=round,line cap=round] (109.61,104.85) -- (114.11,100.35);

\path[draw=drawColor,line width= 0.4pt,line join=round,line cap=round] (109.73,100.36) -- (114.23,104.86);

\path[draw=drawColor,line width= 0.4pt,line join=round,line cap=round] (109.73,104.86) -- (114.23,100.36);

\path[draw=drawColor,line width= 0.4pt,line join=round,line cap=round] (109.85,100.42) -- (114.35,104.92);

\path[draw=drawColor,line width= 0.4pt,line join=round,line cap=round] (109.85,104.92) -- (114.35,100.42);

\path[draw=drawColor,line width= 0.4pt,line join=round,line cap=round] (109.96,100.44) -- (114.46,104.94);

\path[draw=drawColor,line width= 0.4pt,line join=round,line cap=round] (109.96,104.94) -- (114.46,100.44);

\path[draw=drawColor,line width= 0.4pt,line join=round,line cap=round] (110.08,100.46) -- (114.58,104.96);

\path[draw=drawColor,line width= 0.4pt,line join=round,line cap=round] (110.08,104.96) -- (114.58,100.46);

\path[draw=drawColor,line width= 0.4pt,line join=round,line cap=round] (110.20,100.50) -- (114.70,105.00);

\path[draw=drawColor,line width= 0.4pt,line join=round,line cap=round] (110.20,105.00) -- (114.70,100.50);

\path[draw=drawColor,line width= 0.4pt,line join=round,line cap=round] (110.32,100.78) -- (114.82,105.28);

\path[draw=drawColor,line width= 0.4pt,line join=round,line cap=round] (110.32,105.28) -- (114.82,100.78);

\path[draw=drawColor,line width= 0.4pt,line join=round,line cap=round] (110.43,100.93) -- (114.93,105.43);

\path[draw=drawColor,line width= 0.4pt,line join=round,line cap=round] (110.43,105.43) -- (114.93,100.93);

\path[draw=drawColor,line width= 0.4pt,line join=round,line cap=round] (110.55,100.98) -- (115.05,105.48);

\path[draw=drawColor,line width= 0.4pt,line join=round,line cap=round] (110.55,105.48) -- (115.05,100.98);

\path[draw=drawColor,line width= 0.4pt,line join=round,line cap=round] (110.67,101.07) -- (115.17,105.57);

\path[draw=drawColor,line width= 0.4pt,line join=round,line cap=round] (110.67,105.57) -- (115.17,101.07);

\path[draw=drawColor,line width= 0.4pt,line join=round,line cap=round] (110.78,101.09) -- (115.28,105.59);

\path[draw=drawColor,line width= 0.4pt,line join=round,line cap=round] (110.78,105.59) -- (115.28,101.09);

\path[draw=drawColor,line width= 0.4pt,line join=round,line cap=round] (110.90,101.11) -- (115.40,105.61);

\path[draw=drawColor,line width= 0.4pt,line join=round,line cap=round] (110.90,105.61) -- (115.40,101.11);

\path[draw=drawColor,line width= 0.4pt,line join=round,line cap=round] (111.02,101.12) -- (115.52,105.62);

\path[draw=drawColor,line width= 0.4pt,line join=round,line cap=round] (111.02,105.62) -- (115.52,101.12);

\path[draw=drawColor,line width= 0.4pt,line join=round,line cap=round] (111.13,101.28) -- (115.63,105.78);

\path[draw=drawColor,line width= 0.4pt,line join=round,line cap=round] (111.13,105.78) -- (115.63,101.28);

\path[draw=drawColor,line width= 0.4pt,line join=round,line cap=round] (111.25,101.38) -- (115.75,105.88);

\path[draw=drawColor,line width= 0.4pt,line join=round,line cap=round] (111.25,105.88) -- (115.75,101.38);

\path[draw=drawColor,line width= 0.4pt,line join=round,line cap=round] (111.37,101.39) -- (115.87,105.89);

\path[draw=drawColor,line width= 0.4pt,line join=round,line cap=round] (111.37,105.89) -- (115.87,101.39);

\path[draw=drawColor,line width= 0.4pt,line join=round,line cap=round] (111.48,101.55) -- (115.98,106.05);

\path[draw=drawColor,line width= 0.4pt,line join=round,line cap=round] (111.48,106.05) -- (115.98,101.55);

\path[draw=drawColor,line width= 0.4pt,line join=round,line cap=round] (111.60,101.59) -- (116.10,106.09);

\path[draw=drawColor,line width= 0.4pt,line join=round,line cap=round] (111.60,106.09) -- (116.10,101.59);

\path[draw=drawColor,line width= 0.4pt,line join=round,line cap=round] (111.72,101.68) -- (116.22,106.18);

\path[draw=drawColor,line width= 0.4pt,line join=round,line cap=round] (111.72,106.18) -- (116.22,101.68);

\path[draw=drawColor,line width= 0.4pt,line join=round,line cap=round] (111.84,101.84) -- (116.34,106.34);

\path[draw=drawColor,line width= 0.4pt,line join=round,line cap=round] (111.84,106.34) -- (116.34,101.84);

\path[draw=drawColor,line width= 0.4pt,line join=round,line cap=round] (111.95,101.93) -- (116.45,106.43);

\path[draw=drawColor,line width= 0.4pt,line join=round,line cap=round] (111.95,106.43) -- (116.45,101.93);

\path[draw=drawColor,line width= 0.4pt,line join=round,line cap=round] (112.07,101.94) -- (116.57,106.44);

\path[draw=drawColor,line width= 0.4pt,line join=round,line cap=round] (112.07,106.44) -- (116.57,101.94);

\path[draw=drawColor,line width= 0.4pt,line join=round,line cap=round] (112.19,102.24) -- (116.69,106.74);

\path[draw=drawColor,line width= 0.4pt,line join=round,line cap=round] (112.19,106.74) -- (116.69,102.24);

\path[draw=drawColor,line width= 0.4pt,line join=round,line cap=round] (112.30,102.26) -- (116.80,106.76);

\path[draw=drawColor,line width= 0.4pt,line join=round,line cap=round] (112.30,106.76) -- (116.80,102.26);

\path[draw=drawColor,line width= 0.4pt,line join=round,line cap=round] (112.42,102.67) -- (116.92,107.17);

\path[draw=drawColor,line width= 0.4pt,line join=round,line cap=round] (112.42,107.17) -- (116.92,102.67);

\path[draw=drawColor,line width= 0.4pt,line join=round,line cap=round] (112.54,102.70) -- (117.04,107.20);

\path[draw=drawColor,line width= 0.4pt,line join=round,line cap=round] (112.54,107.20) -- (117.04,102.70);

\path[draw=drawColor,line width= 0.4pt,line join=round,line cap=round] (112.65,102.72) -- (117.15,107.22);

\path[draw=drawColor,line width= 0.4pt,line join=round,line cap=round] (112.65,107.22) -- (117.15,102.72);

\path[draw=drawColor,line width= 0.4pt,line join=round,line cap=round] (112.77,102.73) -- (117.27,107.23);

\path[draw=drawColor,line width= 0.4pt,line join=round,line cap=round] (112.77,107.23) -- (117.27,102.73);

\path[draw=drawColor,line width= 0.4pt,line join=round,line cap=round] (112.89,102.83) -- (117.39,107.33);

\path[draw=drawColor,line width= 0.4pt,line join=round,line cap=round] (112.89,107.33) -- (117.39,102.83);

\path[draw=drawColor,line width= 0.4pt,line join=round,line cap=round] (113.01,102.99) -- (117.51,107.49);

\path[draw=drawColor,line width= 0.4pt,line join=round,line cap=round] (113.01,107.49) -- (117.51,102.99);

\path[draw=drawColor,line width= 0.4pt,line join=round,line cap=round] (113.12,103.00) -- (117.62,107.50);

\path[draw=drawColor,line width= 0.4pt,line join=round,line cap=round] (113.12,107.50) -- (117.62,103.00);

\path[draw=drawColor,line width= 0.4pt,line join=round,line cap=round] (113.24,103.15) -- (117.74,107.65);

\path[draw=drawColor,line width= 0.4pt,line join=round,line cap=round] (113.24,107.65) -- (117.74,103.15);

\path[draw=drawColor,line width= 0.4pt,line join=round,line cap=round] (113.36,103.18) -- (117.86,107.68);

\path[draw=drawColor,line width= 0.4pt,line join=round,line cap=round] (113.36,107.68) -- (117.86,103.18);

\path[draw=drawColor,line width= 0.4pt,line join=round,line cap=round] (113.47,103.25) -- (117.97,107.75);

\path[draw=drawColor,line width= 0.4pt,line join=round,line cap=round] (113.47,107.75) -- (117.97,103.25);

\path[draw=drawColor,line width= 0.4pt,line join=round,line cap=round] (113.59,103.43) -- (118.09,107.93);

\path[draw=drawColor,line width= 0.4pt,line join=round,line cap=round] (113.59,107.93) -- (118.09,103.43);

\path[draw=drawColor,line width= 0.4pt,line join=round,line cap=round] (113.71,103.47) -- (118.21,107.97);

\path[draw=drawColor,line width= 0.4pt,line join=round,line cap=round] (113.71,107.97) -- (118.21,103.47);

\path[draw=drawColor,line width= 0.4pt,line join=round,line cap=round] (113.82,103.49) -- (118.32,107.99);

\path[draw=drawColor,line width= 0.4pt,line join=round,line cap=round] (113.82,107.99) -- (118.32,103.49);

\path[draw=drawColor,line width= 0.4pt,line join=round,line cap=round] (113.94,103.50) -- (118.44,108.00);

\path[draw=drawColor,line width= 0.4pt,line join=round,line cap=round] (113.94,108.00) -- (118.44,103.50);

\path[draw=drawColor,line width= 0.4pt,line join=round,line cap=round] (114.06,103.54) -- (118.56,108.04);

\path[draw=drawColor,line width= 0.4pt,line join=round,line cap=round] (114.06,108.04) -- (118.56,103.54);

\path[draw=drawColor,line width= 0.4pt,line join=round,line cap=round] (114.17,103.73) -- (118.67,108.23);

\path[draw=drawColor,line width= 0.4pt,line join=round,line cap=round] (114.17,108.23) -- (118.67,103.73);

\path[draw=drawColor,line width= 0.4pt,line join=round,line cap=round] (114.29,103.78) -- (118.79,108.28);

\path[draw=drawColor,line width= 0.4pt,line join=round,line cap=round] (114.29,108.28) -- (118.79,103.78);

\path[draw=drawColor,line width= 0.4pt,line join=round,line cap=round] (114.41,103.79) -- (118.91,108.29);

\path[draw=drawColor,line width= 0.4pt,line join=round,line cap=round] (114.41,108.29) -- (118.91,103.79);

\path[draw=drawColor,line width= 0.4pt,line join=round,line cap=round] (114.53,103.81) -- (119.03,108.31);

\path[draw=drawColor,line width= 0.4pt,line join=round,line cap=round] (114.53,108.31) -- (119.03,103.81);

\path[draw=drawColor,line width= 0.4pt,line join=round,line cap=round] (114.64,103.81) -- (119.14,108.31);

\path[draw=drawColor,line width= 0.4pt,line join=round,line cap=round] (114.64,108.31) -- (119.14,103.81);

\path[draw=drawColor,line width= 0.4pt,line join=round,line cap=round] (114.76,103.86) -- (119.26,108.36);

\path[draw=drawColor,line width= 0.4pt,line join=round,line cap=round] (114.76,108.36) -- (119.26,103.86);

\path[draw=drawColor,line width= 0.4pt,line join=round,line cap=round] (114.88,104.22) -- (119.38,108.72);

\path[draw=drawColor,line width= 0.4pt,line join=round,line cap=round] (114.88,108.72) -- (119.38,104.22);

\path[draw=drawColor,line width= 0.4pt,line join=round,line cap=round] (114.99,104.32) -- (119.49,108.82);

\path[draw=drawColor,line width= 0.4pt,line join=round,line cap=round] (114.99,108.82) -- (119.49,104.32);

\path[draw=drawColor,line width= 0.4pt,line join=round,line cap=round] (115.11,105.02) -- (119.61,109.52);

\path[draw=drawColor,line width= 0.4pt,line join=round,line cap=round] (115.11,109.52) -- (119.61,105.02);

\path[draw=drawColor,line width= 0.4pt,line join=round,line cap=round] (115.23,105.09) -- (119.73,109.59);

\path[draw=drawColor,line width= 0.4pt,line join=round,line cap=round] (115.23,109.59) -- (119.73,105.09);

\path[draw=drawColor,line width= 0.4pt,line join=round,line cap=round] (115.34,105.11) -- (119.84,109.61);

\path[draw=drawColor,line width= 0.4pt,line join=round,line cap=round] (115.34,109.61) -- (119.84,105.11);

\path[draw=drawColor,line width= 0.4pt,line join=round,line cap=round] (115.46,105.16) -- (119.96,109.66);

\path[draw=drawColor,line width= 0.4pt,line join=round,line cap=round] (115.46,109.66) -- (119.96,105.16);

\path[draw=drawColor,line width= 0.4pt,line join=round,line cap=round] (115.58,105.18) -- (120.08,109.68);

\path[draw=drawColor,line width= 0.4pt,line join=round,line cap=round] (115.58,109.68) -- (120.08,105.18);

\path[draw=drawColor,line width= 0.4pt,line join=round,line cap=round] (115.69,105.20) -- (120.19,109.70);

\path[draw=drawColor,line width= 0.4pt,line join=round,line cap=round] (115.69,109.70) -- (120.19,105.20);

\path[draw=drawColor,line width= 0.4pt,line join=round,line cap=round] (115.81,105.23) -- (120.31,109.73);

\path[draw=drawColor,line width= 0.4pt,line join=round,line cap=round] (115.81,109.73) -- (120.31,105.23);

\path[draw=drawColor,line width= 0.4pt,line join=round,line cap=round] (115.93,105.25) -- (120.43,109.75);

\path[draw=drawColor,line width= 0.4pt,line join=round,line cap=round] (115.93,109.75) -- (120.43,105.25);

\path[draw=drawColor,line width= 0.4pt,line join=round,line cap=round] (116.05,105.36) -- (120.55,109.86);

\path[draw=drawColor,line width= 0.4pt,line join=round,line cap=round] (116.05,109.86) -- (120.55,105.36);

\path[draw=drawColor,line width= 0.4pt,line join=round,line cap=round] (116.16,105.50) -- (120.66,110.00);

\path[draw=drawColor,line width= 0.4pt,line join=round,line cap=round] (116.16,110.00) -- (120.66,105.50);

\path[draw=drawColor,line width= 0.4pt,line join=round,line cap=round] (116.28,105.59) -- (120.78,110.09);

\path[draw=drawColor,line width= 0.4pt,line join=round,line cap=round] (116.28,110.09) -- (120.78,105.59);

\path[draw=drawColor,line width= 0.4pt,line join=round,line cap=round] (116.40,105.62) -- (120.90,110.12);

\path[draw=drawColor,line width= 0.4pt,line join=round,line cap=round] (116.40,110.12) -- (120.90,105.62);

\path[draw=drawColor,line width= 0.4pt,line join=round,line cap=round] (116.51,105.77) -- (121.01,110.27);

\path[draw=drawColor,line width= 0.4pt,line join=round,line cap=round] (116.51,110.27) -- (121.01,105.77);

\path[draw=drawColor,line width= 0.4pt,line join=round,line cap=round] (116.63,106.14) -- (121.13,110.64);

\path[draw=drawColor,line width= 0.4pt,line join=round,line cap=round] (116.63,110.64) -- (121.13,106.14);

\path[draw=drawColor,line width= 0.4pt,line join=round,line cap=round] (116.75,106.47) -- (121.25,110.97);

\path[draw=drawColor,line width= 0.4pt,line join=round,line cap=round] (116.75,110.97) -- (121.25,106.47);

\path[draw=drawColor,line width= 0.4pt,line join=round,line cap=round] (116.86,106.52) -- (121.36,111.02);

\path[draw=drawColor,line width= 0.4pt,line join=round,line cap=round] (116.86,111.02) -- (121.36,106.52);

\path[draw=drawColor,line width= 0.4pt,line join=round,line cap=round] (116.98,106.52) -- (121.48,111.02);

\path[draw=drawColor,line width= 0.4pt,line join=round,line cap=round] (116.98,111.02) -- (121.48,106.52);

\path[draw=drawColor,line width= 0.4pt,line join=round,line cap=round] (117.10,106.55) -- (121.60,111.05);

\path[draw=drawColor,line width= 0.4pt,line join=round,line cap=round] (117.10,111.05) -- (121.60,106.55);

\path[draw=drawColor,line width= 0.4pt,line join=round,line cap=round] (117.21,106.55) -- (121.71,111.05);

\path[draw=drawColor,line width= 0.4pt,line join=round,line cap=round] (117.21,111.05) -- (121.71,106.55);

\path[draw=drawColor,line width= 0.4pt,line join=round,line cap=round] (117.33,106.56) -- (121.83,111.06);

\path[draw=drawColor,line width= 0.4pt,line join=round,line cap=round] (117.33,111.06) -- (121.83,106.56);

\path[draw=drawColor,line width= 0.4pt,line join=round,line cap=round] (117.45,106.63) -- (121.95,111.13);

\path[draw=drawColor,line width= 0.4pt,line join=round,line cap=round] (117.45,111.13) -- (121.95,106.63);

\path[draw=drawColor,line width= 0.4pt,line join=round,line cap=round] (117.57,106.79) -- (122.07,111.29);

\path[draw=drawColor,line width= 0.4pt,line join=round,line cap=round] (117.57,111.29) -- (122.07,106.79);

\path[draw=drawColor,line width= 0.4pt,line join=round,line cap=round] (117.68,107.04) -- (122.18,111.54);

\path[draw=drawColor,line width= 0.4pt,line join=round,line cap=round] (117.68,111.54) -- (122.18,107.04);

\path[draw=drawColor,line width= 0.4pt,line join=round,line cap=round] (117.80,107.10) -- (122.30,111.60);

\path[draw=drawColor,line width= 0.4pt,line join=round,line cap=round] (117.80,111.60) -- (122.30,107.10);

\path[draw=drawColor,line width= 0.4pt,line join=round,line cap=round] (117.92,107.21) -- (122.42,111.71);

\path[draw=drawColor,line width= 0.4pt,line join=round,line cap=round] (117.92,111.71) -- (122.42,107.21);

\path[draw=drawColor,line width= 0.4pt,line join=round,line cap=round] (118.03,107.37) -- (122.53,111.87);

\path[draw=drawColor,line width= 0.4pt,line join=round,line cap=round] (118.03,111.87) -- (122.53,107.37);

\path[draw=drawColor,line width= 0.4pt,line join=round,line cap=round] (118.15,107.53) -- (122.65,112.03);

\path[draw=drawColor,line width= 0.4pt,line join=round,line cap=round] (118.15,112.03) -- (122.65,107.53);

\path[draw=drawColor,line width= 0.4pt,line join=round,line cap=round] (118.27,107.64) -- (122.77,112.14);

\path[draw=drawColor,line width= 0.4pt,line join=round,line cap=round] (118.27,112.14) -- (122.77,107.64);

\path[draw=drawColor,line width= 0.4pt,line join=round,line cap=round] (118.38,107.65) -- (122.88,112.15);

\path[draw=drawColor,line width= 0.4pt,line join=round,line cap=round] (118.38,112.15) -- (122.88,107.65);

\path[draw=drawColor,line width= 0.4pt,line join=round,line cap=round] (118.50,107.68) -- (123.00,112.18);

\path[draw=drawColor,line width= 0.4pt,line join=round,line cap=round] (118.50,112.18) -- (123.00,107.68);

\path[draw=drawColor,line width= 0.4pt,line join=round,line cap=round] (118.62,107.70) -- (123.12,112.20);

\path[draw=drawColor,line width= 0.4pt,line join=round,line cap=round] (118.62,112.20) -- (123.12,107.70);

\path[draw=drawColor,line width= 0.4pt,line join=round,line cap=round] (118.74,107.75) -- (123.24,112.25);

\path[draw=drawColor,line width= 0.4pt,line join=round,line cap=round] (118.74,112.25) -- (123.24,107.75);

\path[draw=drawColor,line width= 0.4pt,line join=round,line cap=round] (118.85,107.83) -- (123.35,112.33);

\path[draw=drawColor,line width= 0.4pt,line join=round,line cap=round] (118.85,112.33) -- (123.35,107.83);

\path[draw=drawColor,line width= 0.4pt,line join=round,line cap=round] (118.97,107.93) -- (123.47,112.43);

\path[draw=drawColor,line width= 0.4pt,line join=round,line cap=round] (118.97,112.43) -- (123.47,107.93);

\path[draw=drawColor,line width= 0.4pt,line join=round,line cap=round] (119.09,108.04) -- (123.59,112.54);

\path[draw=drawColor,line width= 0.4pt,line join=round,line cap=round] (119.09,112.54) -- (123.59,108.04);

\path[draw=drawColor,line width= 0.4pt,line join=round,line cap=round] (119.20,108.05) -- (123.70,112.55);

\path[draw=drawColor,line width= 0.4pt,line join=round,line cap=round] (119.20,112.55) -- (123.70,108.05);

\path[draw=drawColor,line width= 0.4pt,line join=round,line cap=round] (119.32,108.05) -- (123.82,112.55);

\path[draw=drawColor,line width= 0.4pt,line join=round,line cap=round] (119.32,112.55) -- (123.82,108.05);

\path[draw=drawColor,line width= 0.4pt,line join=round,line cap=round] (119.44,108.07) -- (123.94,112.57);

\path[draw=drawColor,line width= 0.4pt,line join=round,line cap=round] (119.44,112.57) -- (123.94,108.07);

\path[draw=drawColor,line width= 0.4pt,line join=round,line cap=round] (119.55,108.17) -- (124.05,112.67);

\path[draw=drawColor,line width= 0.4pt,line join=round,line cap=round] (119.55,112.67) -- (124.05,108.17);

\path[draw=drawColor,line width= 0.4pt,line join=round,line cap=round] (119.67,108.18) -- (124.17,112.68);

\path[draw=drawColor,line width= 0.4pt,line join=round,line cap=round] (119.67,112.68) -- (124.17,108.18);

\path[draw=drawColor,line width= 0.4pt,line join=round,line cap=round] (119.79,108.49) -- (124.29,112.99);

\path[draw=drawColor,line width= 0.4pt,line join=round,line cap=round] (119.79,112.99) -- (124.29,108.49);

\path[draw=drawColor,line width= 0.4pt,line join=round,line cap=round] (119.90,108.66) -- (124.40,113.16);

\path[draw=drawColor,line width= 0.4pt,line join=round,line cap=round] (119.90,113.16) -- (124.40,108.66);

\path[draw=drawColor,line width= 0.4pt,line join=round,line cap=round] (120.02,108.75) -- (124.52,113.25);

\path[draw=drawColor,line width= 0.4pt,line join=round,line cap=round] (120.02,113.25) -- (124.52,108.75);

\path[draw=drawColor,line width= 0.4pt,line join=round,line cap=round] (120.14,108.81) -- (124.64,113.31);

\path[draw=drawColor,line width= 0.4pt,line join=round,line cap=round] (120.14,113.31) -- (124.64,108.81);

\path[draw=drawColor,line width= 0.4pt,line join=round,line cap=round] (120.26,108.81) -- (124.76,113.31);

\path[draw=drawColor,line width= 0.4pt,line join=round,line cap=round] (120.26,113.31) -- (124.76,108.81);

\path[draw=drawColor,line width= 0.4pt,line join=round,line cap=round] (120.37,108.92) -- (124.87,113.42);

\path[draw=drawColor,line width= 0.4pt,line join=round,line cap=round] (120.37,113.42) -- (124.87,108.92);

\path[draw=drawColor,line width= 0.4pt,line join=round,line cap=round] (120.49,109.01) -- (124.99,113.51);

\path[draw=drawColor,line width= 0.4pt,line join=round,line cap=round] (120.49,113.51) -- (124.99,109.01);

\path[draw=drawColor,line width= 0.4pt,line join=round,line cap=round] (120.61,109.24) -- (125.11,113.74);

\path[draw=drawColor,line width= 0.4pt,line join=round,line cap=round] (120.61,113.74) -- (125.11,109.24);

\path[draw=drawColor,line width= 0.4pt,line join=round,line cap=round] (120.72,109.38) -- (125.22,113.88);

\path[draw=drawColor,line width= 0.4pt,line join=round,line cap=round] (120.72,113.88) -- (125.22,109.38);

\path[draw=drawColor,line width= 0.4pt,line join=round,line cap=round] (120.84,109.88) -- (125.34,114.38);

\path[draw=drawColor,line width= 0.4pt,line join=round,line cap=round] (120.84,114.38) -- (125.34,109.88);

\path[draw=drawColor,line width= 0.4pt,line join=round,line cap=round] (120.96,109.91) -- (125.46,114.41);

\path[draw=drawColor,line width= 0.4pt,line join=round,line cap=round] (120.96,114.41) -- (125.46,109.91);

\path[draw=drawColor,line width= 0.4pt,line join=round,line cap=round] (121.07,109.91) -- (125.57,114.41);

\path[draw=drawColor,line width= 0.4pt,line join=round,line cap=round] (121.07,114.41) -- (125.57,109.91);

\path[draw=drawColor,line width= 0.4pt,line join=round,line cap=round] (121.19,109.92) -- (125.69,114.42);

\path[draw=drawColor,line width= 0.4pt,line join=round,line cap=round] (121.19,114.42) -- (125.69,109.92);

\path[draw=drawColor,line width= 0.4pt,line join=round,line cap=round] (121.31,110.01) -- (125.81,114.51);

\path[draw=drawColor,line width= 0.4pt,line join=round,line cap=round] (121.31,114.51) -- (125.81,110.01);

\path[draw=drawColor,line width= 0.4pt,line join=round,line cap=round] (121.42,110.05) -- (125.92,114.55);

\path[draw=drawColor,line width= 0.4pt,line join=round,line cap=round] (121.42,114.55) -- (125.92,110.05);

\path[draw=drawColor,line width= 0.4pt,line join=round,line cap=round] (121.54,110.15) -- (126.04,114.65);

\path[draw=drawColor,line width= 0.4pt,line join=round,line cap=round] (121.54,114.65) -- (126.04,110.15);

\path[draw=drawColor,line width= 0.4pt,line join=round,line cap=round] (121.66,110.30) -- (126.16,114.80);

\path[draw=drawColor,line width= 0.4pt,line join=round,line cap=round] (121.66,114.80) -- (126.16,110.30);

\path[draw=drawColor,line width= 0.4pt,line join=round,line cap=round] (121.78,110.47) -- (126.28,114.97);

\path[draw=drawColor,line width= 0.4pt,line join=round,line cap=round] (121.78,114.97) -- (126.28,110.47);

\path[draw=drawColor,line width= 0.4pt,line join=round,line cap=round] (121.89,110.66) -- (126.39,115.16);

\path[draw=drawColor,line width= 0.4pt,line join=round,line cap=round] (121.89,115.16) -- (126.39,110.66);

\path[draw=drawColor,line width= 0.4pt,line join=round,line cap=round] (122.01,110.71) -- (126.51,115.21);

\path[draw=drawColor,line width= 0.4pt,line join=round,line cap=round] (122.01,115.21) -- (126.51,110.71);

\path[draw=drawColor,line width= 0.4pt,line join=round,line cap=round] (122.13,110.79) -- (126.63,115.29);

\path[draw=drawColor,line width= 0.4pt,line join=round,line cap=round] (122.13,115.29) -- (126.63,110.79);

\path[draw=drawColor,line width= 0.4pt,line join=round,line cap=round] (122.24,111.19) -- (126.74,115.69);

\path[draw=drawColor,line width= 0.4pt,line join=round,line cap=round] (122.24,115.69) -- (126.74,111.19);

\path[draw=drawColor,line width= 0.4pt,line join=round,line cap=round] (122.36,111.53) -- (126.86,116.03);

\path[draw=drawColor,line width= 0.4pt,line join=round,line cap=round] (122.36,116.03) -- (126.86,111.53);

\path[draw=drawColor,line width= 0.4pt,line join=round,line cap=round] (122.48,111.59) -- (126.98,116.09);

\path[draw=drawColor,line width= 0.4pt,line join=round,line cap=round] (122.48,116.09) -- (126.98,111.59);

\path[draw=drawColor,line width= 0.4pt,line join=round,line cap=round] (122.59,111.64) -- (127.09,116.14);

\path[draw=drawColor,line width= 0.4pt,line join=round,line cap=round] (122.59,116.14) -- (127.09,111.64);

\path[draw=drawColor,line width= 0.4pt,line join=round,line cap=round] (122.71,111.64) -- (127.21,116.14);

\path[draw=drawColor,line width= 0.4pt,line join=round,line cap=round] (122.71,116.14) -- (127.21,111.64);

\path[draw=drawColor,line width= 0.4pt,line join=round,line cap=round] (122.83,111.69) -- (127.33,116.19);

\path[draw=drawColor,line width= 0.4pt,line join=round,line cap=round] (122.83,116.19) -- (127.33,111.69);

\path[draw=drawColor,line width= 0.4pt,line join=round,line cap=round] (122.94,111.71) -- (127.44,116.21);

\path[draw=drawColor,line width= 0.4pt,line join=round,line cap=round] (122.94,116.21) -- (127.44,111.71);

\path[draw=drawColor,line width= 0.4pt,line join=round,line cap=round] (123.06,111.79) -- (127.56,116.29);

\path[draw=drawColor,line width= 0.4pt,line join=round,line cap=round] (123.06,116.29) -- (127.56,111.79);

\path[draw=drawColor,line width= 0.4pt,line join=round,line cap=round] (123.18,112.37) -- (127.68,116.87);

\path[draw=drawColor,line width= 0.4pt,line join=round,line cap=round] (123.18,116.87) -- (127.68,112.37);

\path[draw=drawColor,line width= 0.4pt,line join=round,line cap=round] (123.30,112.76) -- (127.80,117.26);

\path[draw=drawColor,line width= 0.4pt,line join=round,line cap=round] (123.30,117.26) -- (127.80,112.76);

\path[draw=drawColor,line width= 0.4pt,line join=round,line cap=round] (123.41,112.97) -- (127.91,117.47);

\path[draw=drawColor,line width= 0.4pt,line join=round,line cap=round] (123.41,117.47) -- (127.91,112.97);

\path[draw=drawColor,line width= 0.4pt,line join=round,line cap=round] (123.53,113.23) -- (128.03,117.73);

\path[draw=drawColor,line width= 0.4pt,line join=round,line cap=round] (123.53,117.73) -- (128.03,113.23);

\path[draw=drawColor,line width= 0.4pt,line join=round,line cap=round] (123.65,113.34) -- (128.15,117.84);

\path[draw=drawColor,line width= 0.4pt,line join=round,line cap=round] (123.65,117.84) -- (128.15,113.34);

\path[draw=drawColor,line width= 0.4pt,line join=round,line cap=round] (123.76,113.49) -- (128.26,117.99);

\path[draw=drawColor,line width= 0.4pt,line join=round,line cap=round] (123.76,117.99) -- (128.26,113.49);

\path[draw=drawColor,line width= 0.4pt,line join=round,line cap=round] (123.88,113.58) -- (128.38,118.08);

\path[draw=drawColor,line width= 0.4pt,line join=round,line cap=round] (123.88,118.08) -- (128.38,113.58);

\path[draw=drawColor,line width= 0.4pt,line join=round,line cap=round] (124.00,113.64) -- (128.50,118.14);

\path[draw=drawColor,line width= 0.4pt,line join=round,line cap=round] (124.00,118.14) -- (128.50,113.64);

\path[draw=drawColor,line width= 0.4pt,line join=round,line cap=round] (124.11,113.70) -- (128.61,118.20);

\path[draw=drawColor,line width= 0.4pt,line join=round,line cap=round] (124.11,118.20) -- (128.61,113.70);

\path[draw=drawColor,line width= 0.4pt,line join=round,line cap=round] (124.23,113.77) -- (128.73,118.27);

\path[draw=drawColor,line width= 0.4pt,line join=round,line cap=round] (124.23,118.27) -- (128.73,113.77);

\path[draw=drawColor,line width= 0.4pt,line join=round,line cap=round] (124.35,113.83) -- (128.85,118.33);

\path[draw=drawColor,line width= 0.4pt,line join=round,line cap=round] (124.35,118.33) -- (128.85,113.83);

\path[draw=drawColor,line width= 0.4pt,line join=round,line cap=round] (124.47,113.87) -- (128.97,118.37);

\path[draw=drawColor,line width= 0.4pt,line join=round,line cap=round] (124.47,118.37) -- (128.97,113.87);

\path[draw=drawColor,line width= 0.4pt,line join=round,line cap=round] (124.58,114.03) -- (129.08,118.53);

\path[draw=drawColor,line width= 0.4pt,line join=round,line cap=round] (124.58,118.53) -- (129.08,114.03);

\path[draw=drawColor,line width= 0.4pt,line join=round,line cap=round] (124.70,114.07) -- (129.20,118.57);

\path[draw=drawColor,line width= 0.4pt,line join=round,line cap=round] (124.70,118.57) -- (129.20,114.07);

\path[draw=drawColor,line width= 0.4pt,line join=round,line cap=round] (124.82,114.25) -- (129.32,118.75);

\path[draw=drawColor,line width= 0.4pt,line join=round,line cap=round] (124.82,118.75) -- (129.32,114.25);

\path[draw=drawColor,line width= 0.4pt,line join=round,line cap=round] (124.93,114.30) -- (129.43,118.80);

\path[draw=drawColor,line width= 0.4pt,line join=round,line cap=round] (124.93,118.80) -- (129.43,114.30);

\path[draw=drawColor,line width= 0.4pt,line join=round,line cap=round] (125.05,114.42) -- (129.55,118.92);

\path[draw=drawColor,line width= 0.4pt,line join=round,line cap=round] (125.05,118.92) -- (129.55,114.42);

\path[draw=drawColor,line width= 0.4pt,line join=round,line cap=round] (125.17,114.49) -- (129.67,118.99);

\path[draw=drawColor,line width= 0.4pt,line join=round,line cap=round] (125.17,118.99) -- (129.67,114.49);

\path[draw=drawColor,line width= 0.4pt,line join=round,line cap=round] (125.28,114.55) -- (129.78,119.05);

\path[draw=drawColor,line width= 0.4pt,line join=round,line cap=round] (125.28,119.05) -- (129.78,114.55);

\path[draw=drawColor,line width= 0.4pt,line join=round,line cap=round] (125.40,114.85) -- (129.90,119.35);

\path[draw=drawColor,line width= 0.4pt,line join=round,line cap=round] (125.40,119.35) -- (129.90,114.85);

\path[draw=drawColor,line width= 0.4pt,line join=round,line cap=round] (125.52,114.93) -- (130.02,119.43);

\path[draw=drawColor,line width= 0.4pt,line join=round,line cap=round] (125.52,119.43) -- (130.02,114.93);

\path[draw=drawColor,line width= 0.4pt,line join=round,line cap=round] (125.63,115.37) -- (130.13,119.87);

\path[draw=drawColor,line width= 0.4pt,line join=round,line cap=round] (125.63,119.87) -- (130.13,115.37);

\path[draw=drawColor,line width= 0.4pt,line join=round,line cap=round] (125.75,115.38) -- (130.25,119.88);

\path[draw=drawColor,line width= 0.4pt,line join=round,line cap=round] (125.75,119.88) -- (130.25,115.38);

\path[draw=drawColor,line width= 0.4pt,line join=round,line cap=round] (125.87,115.44) -- (130.37,119.94);

\path[draw=drawColor,line width= 0.4pt,line join=round,line cap=round] (125.87,119.94) -- (130.37,115.44);

\path[draw=drawColor,line width= 0.4pt,line join=round,line cap=round] (125.99,115.49) -- (130.49,119.99);

\path[draw=drawColor,line width= 0.4pt,line join=round,line cap=round] (125.99,119.99) -- (130.49,115.49);

\path[draw=drawColor,line width= 0.4pt,line join=round,line cap=round] (126.10,115.51) -- (130.60,120.01);

\path[draw=drawColor,line width= 0.4pt,line join=round,line cap=round] (126.10,120.01) -- (130.60,115.51);

\path[draw=drawColor,line width= 0.4pt,line join=round,line cap=round] (126.22,115.69) -- (130.72,120.19);

\path[draw=drawColor,line width= 0.4pt,line join=round,line cap=round] (126.22,120.19) -- (130.72,115.69);

\path[draw=drawColor,line width= 0.4pt,line join=round,line cap=round] (126.34,115.71) -- (130.84,120.21);

\path[draw=drawColor,line width= 0.4pt,line join=round,line cap=round] (126.34,120.21) -- (130.84,115.71);

\path[draw=drawColor,line width= 0.4pt,line join=round,line cap=round] (126.45,115.81) -- (130.95,120.31);

\path[draw=drawColor,line width= 0.4pt,line join=round,line cap=round] (126.45,120.31) -- (130.95,115.81);

\path[draw=drawColor,line width= 0.4pt,line join=round,line cap=round] (126.57,115.88) -- (131.07,120.38);

\path[draw=drawColor,line width= 0.4pt,line join=round,line cap=round] (126.57,120.38) -- (131.07,115.88);

\path[draw=drawColor,line width= 0.4pt,line join=round,line cap=round] (126.69,115.91) -- (131.19,120.41);

\path[draw=drawColor,line width= 0.4pt,line join=round,line cap=round] (126.69,120.41) -- (131.19,115.91);

\path[draw=drawColor,line width= 0.4pt,line join=round,line cap=round] (126.80,115.95) -- (131.30,120.45);

\path[draw=drawColor,line width= 0.4pt,line join=round,line cap=round] (126.80,120.45) -- (131.30,115.95);

\path[draw=drawColor,line width= 0.4pt,line join=round,line cap=round] (126.92,115.99) -- (131.42,120.49);

\path[draw=drawColor,line width= 0.4pt,line join=round,line cap=round] (126.92,120.49) -- (131.42,115.99);

\path[draw=drawColor,line width= 0.4pt,line join=round,line cap=round] (127.04,116.09) -- (131.54,120.59);

\path[draw=drawColor,line width= 0.4pt,line join=round,line cap=round] (127.04,120.59) -- (131.54,116.09);

\path[draw=drawColor,line width= 0.4pt,line join=round,line cap=round] (127.15,116.18) -- (131.65,120.68);

\path[draw=drawColor,line width= 0.4pt,line join=round,line cap=round] (127.15,120.68) -- (131.65,116.18);

\path[draw=drawColor,line width= 0.4pt,line join=round,line cap=round] (127.27,116.68) -- (131.77,121.18);

\path[draw=drawColor,line width= 0.4pt,line join=round,line cap=round] (127.27,121.18) -- (131.77,116.68);

\path[draw=drawColor,line width= 0.4pt,line join=round,line cap=round] (127.39,116.72) -- (131.89,121.22);

\path[draw=drawColor,line width= 0.4pt,line join=round,line cap=round] (127.39,121.22) -- (131.89,116.72);

\path[draw=drawColor,line width= 0.4pt,line join=round,line cap=round] (127.51,116.80) -- (132.01,121.30);

\path[draw=drawColor,line width= 0.4pt,line join=round,line cap=round] (127.51,121.30) -- (132.01,116.80);

\path[draw=drawColor,line width= 0.4pt,line join=round,line cap=round] (127.62,116.84) -- (132.12,121.34);

\path[draw=drawColor,line width= 0.4pt,line join=round,line cap=round] (127.62,121.34) -- (132.12,116.84);

\path[draw=drawColor,line width= 0.4pt,line join=round,line cap=round] (127.74,116.92) -- (132.24,121.42);

\path[draw=drawColor,line width= 0.4pt,line join=round,line cap=round] (127.74,121.42) -- (132.24,116.92);

\path[draw=drawColor,line width= 0.4pt,line join=round,line cap=round] (127.86,116.99) -- (132.36,121.49);

\path[draw=drawColor,line width= 0.4pt,line join=round,line cap=round] (127.86,121.49) -- (132.36,116.99);

\path[draw=drawColor,line width= 0.4pt,line join=round,line cap=round] (127.97,117.06) -- (132.47,121.56);

\path[draw=drawColor,line width= 0.4pt,line join=round,line cap=round] (127.97,121.56) -- (132.47,117.06);

\path[draw=drawColor,line width= 0.4pt,line join=round,line cap=round] (128.09,117.18) -- (132.59,121.68);

\path[draw=drawColor,line width= 0.4pt,line join=round,line cap=round] (128.09,121.68) -- (132.59,117.18);

\path[draw=drawColor,line width= 0.4pt,line join=round,line cap=round] (128.21,117.40) -- (132.71,121.90);

\path[draw=drawColor,line width= 0.4pt,line join=round,line cap=round] (128.21,121.90) -- (132.71,117.40);

\path[draw=drawColor,line width= 0.4pt,line join=round,line cap=round] (128.32,117.46) -- (132.82,121.96);

\path[draw=drawColor,line width= 0.4pt,line join=round,line cap=round] (128.32,121.96) -- (132.82,117.46);

\path[draw=drawColor,line width= 0.4pt,line join=round,line cap=round] (128.44,117.49) -- (132.94,121.99);

\path[draw=drawColor,line width= 0.4pt,line join=round,line cap=round] (128.44,121.99) -- (132.94,117.49);

\path[draw=drawColor,line width= 0.4pt,line join=round,line cap=round] (128.56,117.55) -- (133.06,122.05);

\path[draw=drawColor,line width= 0.4pt,line join=round,line cap=round] (128.56,122.05) -- (133.06,117.55);

\path[draw=drawColor,line width= 0.4pt,line join=round,line cap=round] (128.68,117.56) -- (133.18,122.06);

\path[draw=drawColor,line width= 0.4pt,line join=round,line cap=round] (128.68,122.06) -- (133.18,117.56);

\path[draw=drawColor,line width= 0.4pt,line join=round,line cap=round] (128.79,117.60) -- (133.29,122.10);

\path[draw=drawColor,line width= 0.4pt,line join=round,line cap=round] (128.79,122.10) -- (133.29,117.60);

\path[draw=drawColor,line width= 0.4pt,line join=round,line cap=round] (128.91,117.86) -- (133.41,122.36);

\path[draw=drawColor,line width= 0.4pt,line join=round,line cap=round] (128.91,122.36) -- (133.41,117.86);

\path[draw=drawColor,line width= 0.4pt,line join=round,line cap=round] (129.03,117.98) -- (133.53,122.48);

\path[draw=drawColor,line width= 0.4pt,line join=round,line cap=round] (129.03,122.48) -- (133.53,117.98);

\path[draw=drawColor,line width= 0.4pt,line join=round,line cap=round] (129.14,118.18) -- (133.64,122.68);

\path[draw=drawColor,line width= 0.4pt,line join=round,line cap=round] (129.14,122.68) -- (133.64,118.18);

\path[draw=drawColor,line width= 0.4pt,line join=round,line cap=round] (129.26,118.20) -- (133.76,122.70);

\path[draw=drawColor,line width= 0.4pt,line join=round,line cap=round] (129.26,122.70) -- (133.76,118.20);

\path[draw=drawColor,line width= 0.4pt,line join=round,line cap=round] (129.38,118.37) -- (133.88,122.87);

\path[draw=drawColor,line width= 0.4pt,line join=round,line cap=round] (129.38,122.87) -- (133.88,118.37);

\path[draw=drawColor,line width= 0.4pt,line join=round,line cap=round] (129.49,118.56) -- (133.99,123.06);

\path[draw=drawColor,line width= 0.4pt,line join=round,line cap=round] (129.49,123.06) -- (133.99,118.56);

\path[draw=drawColor,line width= 0.4pt,line join=round,line cap=round] (129.61,118.58) -- (134.11,123.08);

\path[draw=drawColor,line width= 0.4pt,line join=round,line cap=round] (129.61,123.08) -- (134.11,118.58);

\path[draw=drawColor,line width= 0.4pt,line join=round,line cap=round] (129.73,118.63) -- (134.23,123.13);

\path[draw=drawColor,line width= 0.4pt,line join=round,line cap=round] (129.73,123.13) -- (134.23,118.63);

\path[draw=drawColor,line width= 0.4pt,line join=round,line cap=round] (129.84,118.71) -- (134.34,123.21);

\path[draw=drawColor,line width= 0.4pt,line join=round,line cap=round] (129.84,123.21) -- (134.34,118.71);

\path[draw=drawColor,line width= 0.4pt,line join=round,line cap=round] (129.96,118.72) -- (134.46,123.22);

\path[draw=drawColor,line width= 0.4pt,line join=round,line cap=round] (129.96,123.22) -- (134.46,118.72);

\path[draw=drawColor,line width= 0.4pt,line join=round,line cap=round] (130.08,118.99) -- (134.58,123.49);

\path[draw=drawColor,line width= 0.4pt,line join=round,line cap=round] (130.08,123.49) -- (134.58,118.99);

\path[draw=drawColor,line width= 0.4pt,line join=round,line cap=round] (130.20,119.01) -- (134.70,123.51);

\path[draw=drawColor,line width= 0.4pt,line join=round,line cap=round] (130.20,123.51) -- (134.70,119.01);

\path[draw=drawColor,line width= 0.4pt,line join=round,line cap=round] (130.31,119.03) -- (134.81,123.53);

\path[draw=drawColor,line width= 0.4pt,line join=round,line cap=round] (130.31,123.53) -- (134.81,119.03);

\path[draw=drawColor,line width= 0.4pt,line join=round,line cap=round] (130.43,119.15) -- (134.93,123.65);

\path[draw=drawColor,line width= 0.4pt,line join=round,line cap=round] (130.43,123.65) -- (134.93,119.15);

\path[draw=drawColor,line width= 0.4pt,line join=round,line cap=round] (130.55,119.18) -- (135.05,123.68);

\path[draw=drawColor,line width= 0.4pt,line join=round,line cap=round] (130.55,123.68) -- (135.05,119.18);

\path[draw=drawColor,line width= 0.4pt,line join=round,line cap=round] (130.66,119.26) -- (135.16,123.76);

\path[draw=drawColor,line width= 0.4pt,line join=round,line cap=round] (130.66,123.76) -- (135.16,119.26);

\path[draw=drawColor,line width= 0.4pt,line join=round,line cap=round] (130.78,119.52) -- (135.28,124.02);

\path[draw=drawColor,line width= 0.4pt,line join=round,line cap=round] (130.78,124.02) -- (135.28,119.52);

\path[draw=drawColor,line width= 0.4pt,line join=round,line cap=round] (130.90,119.55) -- (135.40,124.05);

\path[draw=drawColor,line width= 0.4pt,line join=round,line cap=round] (130.90,124.05) -- (135.40,119.55);

\path[draw=drawColor,line width= 0.4pt,line join=round,line cap=round] (131.01,119.58) -- (135.51,124.08);

\path[draw=drawColor,line width= 0.4pt,line join=round,line cap=round] (131.01,124.08) -- (135.51,119.58);

\path[draw=drawColor,line width= 0.4pt,line join=round,line cap=round] (131.13,119.60) -- (135.63,124.10);

\path[draw=drawColor,line width= 0.4pt,line join=round,line cap=round] (131.13,124.10) -- (135.63,119.60);

\path[draw=drawColor,line width= 0.4pt,line join=round,line cap=round] (131.25,119.70) -- (135.75,124.20);

\path[draw=drawColor,line width= 0.4pt,line join=round,line cap=round] (131.25,124.20) -- (135.75,119.70);

\path[draw=drawColor,line width= 0.4pt,line join=round,line cap=round] (131.36,119.90) -- (135.86,124.40);

\path[draw=drawColor,line width= 0.4pt,line join=round,line cap=round] (131.36,124.40) -- (135.86,119.90);

\path[draw=drawColor,line width= 0.4pt,line join=round,line cap=round] (131.48,119.97) -- (135.98,124.47);

\path[draw=drawColor,line width= 0.4pt,line join=round,line cap=round] (131.48,124.47) -- (135.98,119.97);

\path[draw=drawColor,line width= 0.4pt,line join=round,line cap=round] (131.60,120.12) -- (136.10,124.62);

\path[draw=drawColor,line width= 0.4pt,line join=round,line cap=round] (131.60,124.62) -- (136.10,120.12);

\path[draw=drawColor,line width= 0.4pt,line join=round,line cap=round] (131.72,120.22) -- (136.22,124.72);

\path[draw=drawColor,line width= 0.4pt,line join=round,line cap=round] (131.72,124.72) -- (136.22,120.22);

\path[draw=drawColor,line width= 0.4pt,line join=round,line cap=round] (131.83,120.23) -- (136.33,124.73);

\path[draw=drawColor,line width= 0.4pt,line join=round,line cap=round] (131.83,124.73) -- (136.33,120.23);

\path[draw=drawColor,line width= 0.4pt,line join=round,line cap=round] (131.95,120.24) -- (136.45,124.74);

\path[draw=drawColor,line width= 0.4pt,line join=round,line cap=round] (131.95,124.74) -- (136.45,120.24);

\path[draw=drawColor,line width= 0.4pt,line join=round,line cap=round] (132.07,120.45) -- (136.57,124.95);

\path[draw=drawColor,line width= 0.4pt,line join=round,line cap=round] (132.07,124.95) -- (136.57,120.45);

\path[draw=drawColor,line width= 0.4pt,line join=round,line cap=round] (132.18,120.47) -- (136.68,124.97);

\path[draw=drawColor,line width= 0.4pt,line join=round,line cap=round] (132.18,124.97) -- (136.68,120.47);

\path[draw=drawColor,line width= 0.4pt,line join=round,line cap=round] (132.30,120.52) -- (136.80,125.02);

\path[draw=drawColor,line width= 0.4pt,line join=round,line cap=round] (132.30,125.02) -- (136.80,120.52);

\path[draw=drawColor,line width= 0.4pt,line join=round,line cap=round] (132.42,120.58) -- (136.92,125.08);

\path[draw=drawColor,line width= 0.4pt,line join=round,line cap=round] (132.42,125.08) -- (136.92,120.58);

\path[draw=drawColor,line width= 0.4pt,line join=round,line cap=round] (132.53,120.60) -- (137.03,125.10);

\path[draw=drawColor,line width= 0.4pt,line join=round,line cap=round] (132.53,125.10) -- (137.03,120.60);

\path[draw=drawColor,line width= 0.4pt,line join=round,line cap=round] (132.65,120.70) -- (137.15,125.20);

\path[draw=drawColor,line width= 0.4pt,line join=round,line cap=round] (132.65,125.20) -- (137.15,120.70);

\path[draw=drawColor,line width= 0.4pt,line join=round,line cap=round] (132.77,120.77) -- (137.27,125.27);

\path[draw=drawColor,line width= 0.4pt,line join=round,line cap=round] (132.77,125.27) -- (137.27,120.77);

\path[draw=drawColor,line width= 0.4pt,line join=round,line cap=round] (132.88,120.77) -- (137.38,125.27);

\path[draw=drawColor,line width= 0.4pt,line join=round,line cap=round] (132.88,125.27) -- (137.38,120.77);

\path[draw=drawColor,line width= 0.4pt,line join=round,line cap=round] (133.00,120.79) -- (137.50,125.29);

\path[draw=drawColor,line width= 0.4pt,line join=round,line cap=round] (133.00,125.29) -- (137.50,120.79);

\path[draw=drawColor,line width= 0.4pt,line join=round,line cap=round] (133.12,120.79) -- (137.62,125.29);

\path[draw=drawColor,line width= 0.4pt,line join=round,line cap=round] (133.12,125.29) -- (137.62,120.79);

\path[draw=drawColor,line width= 0.4pt,line join=round,line cap=round] (133.24,120.84) -- (137.74,125.34);

\path[draw=drawColor,line width= 0.4pt,line join=round,line cap=round] (133.24,125.34) -- (137.74,120.84);

\path[draw=drawColor,line width= 0.4pt,line join=round,line cap=round] (133.35,120.93) -- (137.85,125.43);

\path[draw=drawColor,line width= 0.4pt,line join=round,line cap=round] (133.35,125.43) -- (137.85,120.93);

\path[draw=drawColor,line width= 0.4pt,line join=round,line cap=round] (133.47,120.96) -- (137.97,125.46);

\path[draw=drawColor,line width= 0.4pt,line join=round,line cap=round] (133.47,125.46) -- (137.97,120.96);

\path[draw=drawColor,line width= 0.4pt,line join=round,line cap=round] (133.59,120.96) -- (138.09,125.46);

\path[draw=drawColor,line width= 0.4pt,line join=round,line cap=round] (133.59,125.46) -- (138.09,120.96);

\path[draw=drawColor,line width= 0.4pt,line join=round,line cap=round] (133.70,120.97) -- (138.20,125.47);

\path[draw=drawColor,line width= 0.4pt,line join=round,line cap=round] (133.70,125.47) -- (138.20,120.97);

\path[draw=drawColor,line width= 0.4pt,line join=round,line cap=round] (133.82,121.02) -- (138.32,125.52);

\path[draw=drawColor,line width= 0.4pt,line join=round,line cap=round] (133.82,125.52) -- (138.32,121.02);

\path[draw=drawColor,line width= 0.4pt,line join=round,line cap=round] (133.94,121.10) -- (138.44,125.60);

\path[draw=drawColor,line width= 0.4pt,line join=round,line cap=round] (133.94,125.60) -- (138.44,121.10);

\path[draw=drawColor,line width= 0.4pt,line join=round,line cap=round] (134.05,121.23) -- (138.55,125.73);

\path[draw=drawColor,line width= 0.4pt,line join=round,line cap=round] (134.05,125.73) -- (138.55,121.23);

\path[draw=drawColor,line width= 0.4pt,line join=round,line cap=round] (134.17,121.30) -- (138.67,125.80);

\path[draw=drawColor,line width= 0.4pt,line join=round,line cap=round] (134.17,125.80) -- (138.67,121.30);

\path[draw=drawColor,line width= 0.4pt,line join=round,line cap=round] (134.29,121.42) -- (138.79,125.92);

\path[draw=drawColor,line width= 0.4pt,line join=round,line cap=round] (134.29,125.92) -- (138.79,121.42);

\path[draw=drawColor,line width= 0.4pt,line join=round,line cap=round] (134.41,121.47) -- (138.91,125.97);

\path[draw=drawColor,line width= 0.4pt,line join=round,line cap=round] (134.41,125.97) -- (138.91,121.47);

\path[draw=drawColor,line width= 0.4pt,line join=round,line cap=round] (134.52,121.49) -- (139.02,125.99);

\path[draw=drawColor,line width= 0.4pt,line join=round,line cap=round] (134.52,125.99) -- (139.02,121.49);

\path[draw=drawColor,line width= 0.4pt,line join=round,line cap=round] (134.64,121.53) -- (139.14,126.03);

\path[draw=drawColor,line width= 0.4pt,line join=round,line cap=round] (134.64,126.03) -- (139.14,121.53);

\path[draw=drawColor,line width= 0.4pt,line join=round,line cap=round] (134.76,121.67) -- (139.26,126.17);

\path[draw=drawColor,line width= 0.4pt,line join=round,line cap=round] (134.76,126.17) -- (139.26,121.67);

\path[draw=drawColor,line width= 0.4pt,line join=round,line cap=round] (134.87,121.69) -- (139.37,126.19);

\path[draw=drawColor,line width= 0.4pt,line join=round,line cap=round] (134.87,126.19) -- (139.37,121.69);

\path[draw=drawColor,line width= 0.4pt,line join=round,line cap=round] (134.99,121.82) -- (139.49,126.32);

\path[draw=drawColor,line width= 0.4pt,line join=round,line cap=round] (134.99,126.32) -- (139.49,121.82);

\path[draw=drawColor,line width= 0.4pt,line join=round,line cap=round] (135.11,121.83) -- (139.61,126.33);

\path[draw=drawColor,line width= 0.4pt,line join=round,line cap=round] (135.11,126.33) -- (139.61,121.83);

\path[draw=drawColor,line width= 0.4pt,line join=round,line cap=round] (135.22,122.01) -- (139.72,126.51);

\path[draw=drawColor,line width= 0.4pt,line join=round,line cap=round] (135.22,126.51) -- (139.72,122.01);

\path[draw=drawColor,line width= 0.4pt,line join=round,line cap=round] (135.34,122.06) -- (139.84,126.56);

\path[draw=drawColor,line width= 0.4pt,line join=round,line cap=round] (135.34,126.56) -- (139.84,122.06);

\path[draw=drawColor,line width= 0.4pt,line join=round,line cap=round] (135.46,122.06) -- (139.96,126.56);

\path[draw=drawColor,line width= 0.4pt,line join=round,line cap=round] (135.46,126.56) -- (139.96,122.06);

\path[draw=drawColor,line width= 0.4pt,line join=round,line cap=round] (135.57,122.10) -- (140.07,126.60);

\path[draw=drawColor,line width= 0.4pt,line join=round,line cap=round] (135.57,126.60) -- (140.07,122.10);

\path[draw=drawColor,line width= 0.4pt,line join=round,line cap=round] (135.69,122.69) -- (140.19,127.19);

\path[draw=drawColor,line width= 0.4pt,line join=round,line cap=round] (135.69,127.19) -- (140.19,122.69);

\path[draw=drawColor,line width= 0.4pt,line join=round,line cap=round] (135.81,122.70) -- (140.31,127.20);

\path[draw=drawColor,line width= 0.4pt,line join=round,line cap=round] (135.81,127.20) -- (140.31,122.70);

\path[draw=drawColor,line width= 0.4pt,line join=round,line cap=round] (135.93,122.74) -- (140.43,127.24);

\path[draw=drawColor,line width= 0.4pt,line join=round,line cap=round] (135.93,127.24) -- (140.43,122.74);

\path[draw=drawColor,line width= 0.4pt,line join=round,line cap=round] (136.04,123.03) -- (140.54,127.53);

\path[draw=drawColor,line width= 0.4pt,line join=round,line cap=round] (136.04,127.53) -- (140.54,123.03);

\path[draw=drawColor,line width= 0.4pt,line join=round,line cap=round] (136.16,123.11) -- (140.66,127.61);

\path[draw=drawColor,line width= 0.4pt,line join=round,line cap=round] (136.16,127.61) -- (140.66,123.11);

\path[draw=drawColor,line width= 0.4pt,line join=round,line cap=round] (136.28,123.17) -- (140.78,127.67);

\path[draw=drawColor,line width= 0.4pt,line join=round,line cap=round] (136.28,127.67) -- (140.78,123.17);

\path[draw=drawColor,line width= 0.4pt,line join=round,line cap=round] (136.39,123.25) -- (140.89,127.75);

\path[draw=drawColor,line width= 0.4pt,line join=round,line cap=round] (136.39,127.75) -- (140.89,123.25);

\path[draw=drawColor,line width= 0.4pt,line join=round,line cap=round] (136.51,123.40) -- (141.01,127.90);

\path[draw=drawColor,line width= 0.4pt,line join=round,line cap=round] (136.51,127.90) -- (141.01,123.40);

\path[draw=drawColor,line width= 0.4pt,line join=round,line cap=round] (136.63,123.42) -- (141.13,127.92);

\path[draw=drawColor,line width= 0.4pt,line join=round,line cap=round] (136.63,127.92) -- (141.13,123.42);

\path[draw=drawColor,line width= 0.4pt,line join=round,line cap=round] (136.74,123.42) -- (141.24,127.92);

\path[draw=drawColor,line width= 0.4pt,line join=round,line cap=round] (136.74,127.92) -- (141.24,123.42);

\path[draw=drawColor,line width= 0.4pt,line join=round,line cap=round] (136.86,123.77) -- (141.36,128.27);

\path[draw=drawColor,line width= 0.4pt,line join=round,line cap=round] (136.86,128.27) -- (141.36,123.77);

\path[draw=drawColor,line width= 0.4pt,line join=round,line cap=round] (136.98,123.81) -- (141.48,128.31);

\path[draw=drawColor,line width= 0.4pt,line join=round,line cap=round] (136.98,128.31) -- (141.48,123.81);

\path[draw=drawColor,line width= 0.4pt,line join=round,line cap=round] (137.09,123.81) -- (141.59,128.31);

\path[draw=drawColor,line width= 0.4pt,line join=round,line cap=round] (137.09,128.31) -- (141.59,123.81);

\path[draw=drawColor,line width= 0.4pt,line join=round,line cap=round] (137.21,123.88) -- (141.71,128.38);

\path[draw=drawColor,line width= 0.4pt,line join=round,line cap=round] (137.21,128.38) -- (141.71,123.88);

\path[draw=drawColor,line width= 0.4pt,line join=round,line cap=round] (137.33,123.91) -- (141.83,128.41);

\path[draw=drawColor,line width= 0.4pt,line join=round,line cap=round] (137.33,128.41) -- (141.83,123.91);

\path[draw=drawColor,line width= 0.4pt,line join=round,line cap=round] (137.45,123.98) -- (141.95,128.48);

\path[draw=drawColor,line width= 0.4pt,line join=round,line cap=round] (137.45,128.48) -- (141.95,123.98);

\path[draw=drawColor,line width= 0.4pt,line join=round,line cap=round] (137.56,124.25) -- (142.06,128.75);

\path[draw=drawColor,line width= 0.4pt,line join=round,line cap=round] (137.56,128.75) -- (142.06,124.25);

\path[draw=drawColor,line width= 0.4pt,line join=round,line cap=round] (137.68,124.36) -- (142.18,128.86);

\path[draw=drawColor,line width= 0.4pt,line join=round,line cap=round] (137.68,128.86) -- (142.18,124.36);

\path[draw=drawColor,line width= 0.4pt,line join=round,line cap=round] (137.80,124.75) -- (142.30,129.25);

\path[draw=drawColor,line width= 0.4pt,line join=round,line cap=round] (137.80,129.25) -- (142.30,124.75);

\path[draw=drawColor,line width= 0.4pt,line join=round,line cap=round] (137.91,124.79) -- (142.41,129.29);

\path[draw=drawColor,line width= 0.4pt,line join=round,line cap=round] (137.91,129.29) -- (142.41,124.79);

\path[draw=drawColor,line width= 0.4pt,line join=round,line cap=round] (138.03,124.82) -- (142.53,129.32);

\path[draw=drawColor,line width= 0.4pt,line join=round,line cap=round] (138.03,129.32) -- (142.53,124.82);

\path[draw=drawColor,line width= 0.4pt,line join=round,line cap=round] (138.15,124.86) -- (142.65,129.36);

\path[draw=drawColor,line width= 0.4pt,line join=round,line cap=round] (138.15,129.36) -- (142.65,124.86);

\path[draw=drawColor,line width= 0.4pt,line join=round,line cap=round] (138.26,125.06) -- (142.76,129.56);

\path[draw=drawColor,line width= 0.4pt,line join=round,line cap=round] (138.26,129.56) -- (142.76,125.06);

\path[draw=drawColor,line width= 0.4pt,line join=round,line cap=round] (138.38,125.14) -- (142.88,129.64);

\path[draw=drawColor,line width= 0.4pt,line join=round,line cap=round] (138.38,129.64) -- (142.88,125.14);

\path[draw=drawColor,line width= 0.4pt,line join=round,line cap=round] (138.50,125.16) -- (143.00,129.66);

\path[draw=drawColor,line width= 0.4pt,line join=round,line cap=round] (138.50,129.66) -- (143.00,125.16);

\path[draw=drawColor,line width= 0.4pt,line join=round,line cap=round] (138.62,125.19) -- (143.12,129.69);

\path[draw=drawColor,line width= 0.4pt,line join=round,line cap=round] (138.62,129.69) -- (143.12,125.19);

\path[draw=drawColor,line width= 0.4pt,line join=round,line cap=round] (138.73,125.45) -- (143.23,129.95);

\path[draw=drawColor,line width= 0.4pt,line join=round,line cap=round] (138.73,129.95) -- (143.23,125.45);

\path[draw=drawColor,line width= 0.4pt,line join=round,line cap=round] (138.85,125.58) -- (143.35,130.08);

\path[draw=drawColor,line width= 0.4pt,line join=round,line cap=round] (138.85,130.08) -- (143.35,125.58);

\path[draw=drawColor,line width= 0.4pt,line join=round,line cap=round] (138.97,125.74) -- (143.47,130.24);

\path[draw=drawColor,line width= 0.4pt,line join=round,line cap=round] (138.97,130.24) -- (143.47,125.74);

\path[draw=drawColor,line width= 0.4pt,line join=round,line cap=round] (139.08,125.74) -- (143.58,130.24);

\path[draw=drawColor,line width= 0.4pt,line join=round,line cap=round] (139.08,130.24) -- (143.58,125.74);

\path[draw=drawColor,line width= 0.4pt,line join=round,line cap=round] (139.20,125.75) -- (143.70,130.25);

\path[draw=drawColor,line width= 0.4pt,line join=round,line cap=round] (139.20,130.25) -- (143.70,125.75);

\path[draw=drawColor,line width= 0.4pt,line join=round,line cap=round] (139.32,125.84) -- (143.82,130.34);

\path[draw=drawColor,line width= 0.4pt,line join=round,line cap=round] (139.32,130.34) -- (143.82,125.84);

\path[draw=drawColor,line width= 0.4pt,line join=round,line cap=round] (139.43,125.87) -- (143.93,130.37);

\path[draw=drawColor,line width= 0.4pt,line join=round,line cap=round] (139.43,130.37) -- (143.93,125.87);

\path[draw=drawColor,line width= 0.4pt,line join=round,line cap=round] (139.55,125.92) -- (144.05,130.42);

\path[draw=drawColor,line width= 0.4pt,line join=round,line cap=round] (139.55,130.42) -- (144.05,125.92);

\path[draw=drawColor,line width= 0.4pt,line join=round,line cap=round] (139.67,126.30) -- (144.17,130.80);

\path[draw=drawColor,line width= 0.4pt,line join=round,line cap=round] (139.67,130.80) -- (144.17,126.30);

\path[draw=drawColor,line width= 0.4pt,line join=round,line cap=round] (139.78,126.37) -- (144.28,130.87);

\path[draw=drawColor,line width= 0.4pt,line join=round,line cap=round] (139.78,130.87) -- (144.28,126.37);

\path[draw=drawColor,line width= 0.4pt,line join=round,line cap=round] (139.90,126.50) -- (144.40,131.00);

\path[draw=drawColor,line width= 0.4pt,line join=round,line cap=round] (139.90,131.00) -- (144.40,126.50);

\path[draw=drawColor,line width= 0.4pt,line join=round,line cap=round] (140.02,126.54) -- (144.52,131.04);

\path[draw=drawColor,line width= 0.4pt,line join=round,line cap=round] (140.02,131.04) -- (144.52,126.54);

\path[draw=drawColor,line width= 0.4pt,line join=round,line cap=round] (140.14,126.57) -- (144.64,131.07);

\path[draw=drawColor,line width= 0.4pt,line join=round,line cap=round] (140.14,131.07) -- (144.64,126.57);

\path[draw=drawColor,line width= 0.4pt,line join=round,line cap=round] (140.25,126.68) -- (144.75,131.18);

\path[draw=drawColor,line width= 0.4pt,line join=round,line cap=round] (140.25,131.18) -- (144.75,126.68);

\path[draw=drawColor,line width= 0.4pt,line join=round,line cap=round] (140.37,126.77) -- (144.87,131.27);

\path[draw=drawColor,line width= 0.4pt,line join=round,line cap=round] (140.37,131.27) -- (144.87,126.77);

\path[draw=drawColor,line width= 0.4pt,line join=round,line cap=round] (140.49,126.83) -- (144.99,131.33);

\path[draw=drawColor,line width= 0.4pt,line join=round,line cap=round] (140.49,131.33) -- (144.99,126.83);

\path[draw=drawColor,line width= 0.4pt,line join=round,line cap=round] (140.60,126.94) -- (145.10,131.44);

\path[draw=drawColor,line width= 0.4pt,line join=round,line cap=round] (140.60,131.44) -- (145.10,126.94);

\path[draw=drawColor,line width= 0.4pt,line join=round,line cap=round] (140.72,127.12) -- (145.22,131.62);

\path[draw=drawColor,line width= 0.4pt,line join=round,line cap=round] (140.72,131.62) -- (145.22,127.12);

\path[draw=drawColor,line width= 0.4pt,line join=round,line cap=round] (140.84,127.15) -- (145.34,131.65);

\path[draw=drawColor,line width= 0.4pt,line join=round,line cap=round] (140.84,131.65) -- (145.34,127.15);

\path[draw=drawColor,line width= 0.4pt,line join=round,line cap=round] (140.95,127.28) -- (145.45,131.78);

\path[draw=drawColor,line width= 0.4pt,line join=round,line cap=round] (140.95,131.78) -- (145.45,127.28);

\path[draw=drawColor,line width= 0.4pt,line join=round,line cap=round] (141.07,127.28) -- (145.57,131.78);

\path[draw=drawColor,line width= 0.4pt,line join=round,line cap=round] (141.07,131.78) -- (145.57,127.28);

\path[draw=drawColor,line width= 0.4pt,line join=round,line cap=round] (141.19,127.38) -- (145.69,131.88);

\path[draw=drawColor,line width= 0.4pt,line join=round,line cap=round] (141.19,131.88) -- (145.69,127.38);

\path[draw=drawColor,line width= 0.4pt,line join=round,line cap=round] (141.30,127.50) -- (145.80,132.00);

\path[draw=drawColor,line width= 0.4pt,line join=round,line cap=round] (141.30,132.00) -- (145.80,127.50);

\path[draw=drawColor,line width= 0.4pt,line join=round,line cap=round] (141.42,127.52) -- (145.92,132.02);

\path[draw=drawColor,line width= 0.4pt,line join=round,line cap=round] (141.42,132.02) -- (145.92,127.52);

\path[draw=drawColor,line width= 0.4pt,line join=round,line cap=round] (141.54,127.74) -- (146.04,132.24);

\path[draw=drawColor,line width= 0.4pt,line join=round,line cap=round] (141.54,132.24) -- (146.04,127.74);

\path[draw=drawColor,line width= 0.4pt,line join=round,line cap=round] (141.66,127.80) -- (146.16,132.30);

\path[draw=drawColor,line width= 0.4pt,line join=round,line cap=round] (141.66,132.30) -- (146.16,127.80);

\path[draw=drawColor,line width= 0.4pt,line join=round,line cap=round] (141.77,127.94) -- (146.27,132.44);

\path[draw=drawColor,line width= 0.4pt,line join=round,line cap=round] (141.77,132.44) -- (146.27,127.94);

\path[draw=drawColor,line width= 0.4pt,line join=round,line cap=round] (141.89,128.01) -- (146.39,132.51);

\path[draw=drawColor,line width= 0.4pt,line join=round,line cap=round] (141.89,132.51) -- (146.39,128.01);

\path[draw=drawColor,line width= 0.4pt,line join=round,line cap=round] (142.01,128.23) -- (146.51,132.73);

\path[draw=drawColor,line width= 0.4pt,line join=round,line cap=round] (142.01,132.73) -- (146.51,128.23);

\path[draw=drawColor,line width= 0.4pt,line join=round,line cap=round] (142.12,128.37) -- (146.62,132.87);

\path[draw=drawColor,line width= 0.4pt,line join=round,line cap=round] (142.12,132.87) -- (146.62,128.37);

\path[draw=drawColor,line width= 0.4pt,line join=round,line cap=round] (142.24,128.50) -- (146.74,133.00);

\path[draw=drawColor,line width= 0.4pt,line join=round,line cap=round] (142.24,133.00) -- (146.74,128.50);

\path[draw=drawColor,line width= 0.4pt,line join=round,line cap=round] (142.36,128.56) -- (146.86,133.06);

\path[draw=drawColor,line width= 0.4pt,line join=round,line cap=round] (142.36,133.06) -- (146.86,128.56);

\path[draw=drawColor,line width= 0.4pt,line join=round,line cap=round] (142.47,128.70) -- (146.97,133.20);

\path[draw=drawColor,line width= 0.4pt,line join=round,line cap=round] (142.47,133.20) -- (146.97,128.70);

\path[draw=drawColor,line width= 0.4pt,line join=round,line cap=round] (142.59,128.71) -- (147.09,133.21);

\path[draw=drawColor,line width= 0.4pt,line join=round,line cap=round] (142.59,133.21) -- (147.09,128.71);

\path[draw=drawColor,line width= 0.4pt,line join=round,line cap=round] (142.71,128.74) -- (147.21,133.24);

\path[draw=drawColor,line width= 0.4pt,line join=round,line cap=round] (142.71,133.24) -- (147.21,128.74);

\path[draw=drawColor,line width= 0.4pt,line join=round,line cap=round] (142.82,128.75) -- (147.32,133.25);

\path[draw=drawColor,line width= 0.4pt,line join=round,line cap=round] (142.82,133.25) -- (147.32,128.75);

\path[draw=drawColor,line width= 0.4pt,line join=round,line cap=round] (142.94,128.79) -- (147.44,133.29);

\path[draw=drawColor,line width= 0.4pt,line join=round,line cap=round] (142.94,133.29) -- (147.44,128.79);

\path[draw=drawColor,line width= 0.4pt,line join=round,line cap=round] (143.06,130.00) -- (147.56,134.50);

\path[draw=drawColor,line width= 0.4pt,line join=round,line cap=round] (143.06,134.50) -- (147.56,130.00);

\path[draw=drawColor,line width= 0.4pt,line join=round,line cap=round] (143.18,130.21) -- (147.68,134.71);

\path[draw=drawColor,line width= 0.4pt,line join=round,line cap=round] (143.18,134.71) -- (147.68,130.21);

\path[draw=drawColor,line width= 0.4pt,line join=round,line cap=round] (143.29,130.30) -- (147.79,134.80);

\path[draw=drawColor,line width= 0.4pt,line join=round,line cap=round] (143.29,134.80) -- (147.79,130.30);

\path[draw=drawColor,line width= 0.4pt,line join=round,line cap=round] (143.41,130.44) -- (147.91,134.94);

\path[draw=drawColor,line width= 0.4pt,line join=round,line cap=round] (143.41,134.94) -- (147.91,130.44);

\path[draw=drawColor,line width= 0.4pt,line join=round,line cap=round] (143.53,130.53) -- (148.03,135.03);

\path[draw=drawColor,line width= 0.4pt,line join=round,line cap=round] (143.53,135.03) -- (148.03,130.53);

\path[draw=drawColor,line width= 0.4pt,line join=round,line cap=round] (143.64,130.75) -- (148.14,135.25);

\path[draw=drawColor,line width= 0.4pt,line join=round,line cap=round] (143.64,135.25) -- (148.14,130.75);

\path[draw=drawColor,line width= 0.4pt,line join=round,line cap=round] (143.76,130.77) -- (148.26,135.27);

\path[draw=drawColor,line width= 0.4pt,line join=round,line cap=round] (143.76,135.27) -- (148.26,130.77);

\path[draw=drawColor,line width= 0.4pt,line join=round,line cap=round] (143.88,130.86) -- (148.38,135.36);

\path[draw=drawColor,line width= 0.4pt,line join=round,line cap=round] (143.88,135.36) -- (148.38,130.86);

\path[draw=drawColor,line width= 0.4pt,line join=round,line cap=round] (143.99,130.87) -- (148.49,135.37);

\path[draw=drawColor,line width= 0.4pt,line join=round,line cap=round] (143.99,135.37) -- (148.49,130.87);

\path[draw=drawColor,line width= 0.4pt,line join=round,line cap=round] (144.11,130.98) -- (148.61,135.48);

\path[draw=drawColor,line width= 0.4pt,line join=round,line cap=round] (144.11,135.48) -- (148.61,130.98);

\path[draw=drawColor,line width= 0.4pt,line join=round,line cap=round] (144.23,131.01) -- (148.73,135.51);

\path[draw=drawColor,line width= 0.4pt,line join=round,line cap=round] (144.23,135.51) -- (148.73,131.01);

\path[draw=drawColor,line width= 0.4pt,line join=round,line cap=round] (144.35,131.29) -- (148.85,135.79);

\path[draw=drawColor,line width= 0.4pt,line join=round,line cap=round] (144.35,135.79) -- (148.85,131.29);

\path[draw=drawColor,line width= 0.4pt,line join=round,line cap=round] (144.46,131.32) -- (148.96,135.82);

\path[draw=drawColor,line width= 0.4pt,line join=round,line cap=round] (144.46,135.82) -- (148.96,131.32);

\path[draw=drawColor,line width= 0.4pt,line join=round,line cap=round] (144.58,131.43) -- (149.08,135.93);

\path[draw=drawColor,line width= 0.4pt,line join=round,line cap=round] (144.58,135.93) -- (149.08,131.43);

\path[draw=drawColor,line width= 0.4pt,line join=round,line cap=round] (144.70,131.52) -- (149.20,136.02);

\path[draw=drawColor,line width= 0.4pt,line join=round,line cap=round] (144.70,136.02) -- (149.20,131.52);

\path[draw=drawColor,line width= 0.4pt,line join=round,line cap=round] (144.81,131.91) -- (149.31,136.41);

\path[draw=drawColor,line width= 0.4pt,line join=round,line cap=round] (144.81,136.41) -- (149.31,131.91);

\path[draw=drawColor,line width= 0.4pt,line join=round,line cap=round] (144.93,131.93) -- (149.43,136.43);

\path[draw=drawColor,line width= 0.4pt,line join=round,line cap=round] (144.93,136.43) -- (149.43,131.93);

\path[draw=drawColor,line width= 0.4pt,line join=round,line cap=round] (145.05,131.97) -- (149.55,136.47);

\path[draw=drawColor,line width= 0.4pt,line join=round,line cap=round] (145.05,136.47) -- (149.55,131.97);

\path[draw=drawColor,line width= 0.4pt,line join=round,line cap=round] (145.16,132.00) -- (149.66,136.50);

\path[draw=drawColor,line width= 0.4pt,line join=round,line cap=round] (145.16,136.50) -- (149.66,132.00);

\path[draw=drawColor,line width= 0.4pt,line join=round,line cap=round] (145.28,132.08) -- (149.78,136.58);

\path[draw=drawColor,line width= 0.4pt,line join=round,line cap=round] (145.28,136.58) -- (149.78,132.08);

\path[draw=drawColor,line width= 0.4pt,line join=round,line cap=round] (145.40,132.20) -- (149.90,136.70);

\path[draw=drawColor,line width= 0.4pt,line join=round,line cap=round] (145.40,136.70) -- (149.90,132.20);

\path[draw=drawColor,line width= 0.4pt,line join=round,line cap=round] (145.51,132.22) -- (150.01,136.72);

\path[draw=drawColor,line width= 0.4pt,line join=round,line cap=round] (145.51,136.72) -- (150.01,132.22);

\path[draw=drawColor,line width= 0.4pt,line join=round,line cap=round] (145.63,132.27) -- (150.13,136.77);

\path[draw=drawColor,line width= 0.4pt,line join=round,line cap=round] (145.63,136.77) -- (150.13,132.27);

\path[draw=drawColor,line width= 0.4pt,line join=round,line cap=round] (145.75,132.42) -- (150.25,136.92);

\path[draw=drawColor,line width= 0.4pt,line join=round,line cap=round] (145.75,136.92) -- (150.25,132.42);

\path[draw=drawColor,line width= 0.4pt,line join=round,line cap=round] (145.87,132.53) -- (150.37,137.03);

\path[draw=drawColor,line width= 0.4pt,line join=round,line cap=round] (145.87,137.03) -- (150.37,132.53);

\path[draw=drawColor,line width= 0.4pt,line join=round,line cap=round] (145.98,132.54) -- (150.48,137.04);

\path[draw=drawColor,line width= 0.4pt,line join=round,line cap=round] (145.98,137.04) -- (150.48,132.54);

\path[draw=drawColor,line width= 0.4pt,line join=round,line cap=round] (146.10,132.87) -- (150.60,137.37);

\path[draw=drawColor,line width= 0.4pt,line join=round,line cap=round] (146.10,137.37) -- (150.60,132.87);

\path[draw=drawColor,line width= 0.4pt,line join=round,line cap=round] (146.22,132.94) -- (150.72,137.44);

\path[draw=drawColor,line width= 0.4pt,line join=round,line cap=round] (146.22,137.44) -- (150.72,132.94);

\path[draw=drawColor,line width= 0.4pt,line join=round,line cap=round] (146.33,132.95) -- (150.83,137.45);

\path[draw=drawColor,line width= 0.4pt,line join=round,line cap=round] (146.33,137.45) -- (150.83,132.95);

\path[draw=drawColor,line width= 0.4pt,line join=round,line cap=round] (146.45,132.95) -- (150.95,137.45);

\path[draw=drawColor,line width= 0.4pt,line join=round,line cap=round] (146.45,137.45) -- (150.95,132.95);

\path[draw=drawColor,line width= 0.4pt,line join=round,line cap=round] (146.57,133.14) -- (151.07,137.64);

\path[draw=drawColor,line width= 0.4pt,line join=round,line cap=round] (146.57,137.64) -- (151.07,133.14);

\path[draw=drawColor,line width= 0.4pt,line join=round,line cap=round] (146.68,133.21) -- (151.18,137.71);

\path[draw=drawColor,line width= 0.4pt,line join=round,line cap=round] (146.68,137.71) -- (151.18,133.21);

\path[draw=drawColor,line width= 0.4pt,line join=round,line cap=round] (146.80,133.34) -- (151.30,137.84);

\path[draw=drawColor,line width= 0.4pt,line join=round,line cap=round] (146.80,137.84) -- (151.30,133.34);

\path[draw=drawColor,line width= 0.4pt,line join=round,line cap=round] (146.92,133.55) -- (151.42,138.05);

\path[draw=drawColor,line width= 0.4pt,line join=round,line cap=round] (146.92,138.05) -- (151.42,133.55);

\path[draw=drawColor,line width= 0.4pt,line join=round,line cap=round] (147.03,133.62) -- (151.53,138.12);

\path[draw=drawColor,line width= 0.4pt,line join=round,line cap=round] (147.03,138.12) -- (151.53,133.62);

\path[draw=drawColor,line width= 0.4pt,line join=round,line cap=round] (147.15,133.68) -- (151.65,138.18);

\path[draw=drawColor,line width= 0.4pt,line join=round,line cap=round] (147.15,138.18) -- (151.65,133.68);

\path[draw=drawColor,line width= 0.4pt,line join=round,line cap=round] (147.27,133.94) -- (151.77,138.44);

\path[draw=drawColor,line width= 0.4pt,line join=round,line cap=round] (147.27,138.44) -- (151.77,133.94);

\path[draw=drawColor,line width= 0.4pt,line join=round,line cap=round] (147.39,134.01) -- (151.89,138.51);

\path[draw=drawColor,line width= 0.4pt,line join=round,line cap=round] (147.39,138.51) -- (151.89,134.01);

\path[draw=drawColor,line width= 0.4pt,line join=round,line cap=round] (147.50,134.09) -- (152.00,138.59);

\path[draw=drawColor,line width= 0.4pt,line join=round,line cap=round] (147.50,138.59) -- (152.00,134.09);

\path[draw=drawColor,line width= 0.4pt,line join=round,line cap=round] (147.62,134.15) -- (152.12,138.65);

\path[draw=drawColor,line width= 0.4pt,line join=round,line cap=round] (147.62,138.65) -- (152.12,134.15);

\path[draw=drawColor,line width= 0.4pt,line join=round,line cap=round] (147.74,134.53) -- (152.24,139.03);

\path[draw=drawColor,line width= 0.4pt,line join=round,line cap=round] (147.74,139.03) -- (152.24,134.53);

\path[draw=drawColor,line width= 0.4pt,line join=round,line cap=round] (147.85,134.63) -- (152.35,139.13);

\path[draw=drawColor,line width= 0.4pt,line join=round,line cap=round] (147.85,139.13) -- (152.35,134.63);

\path[draw=drawColor,line width= 0.4pt,line join=round,line cap=round] (147.97,134.71) -- (152.47,139.21);

\path[draw=drawColor,line width= 0.4pt,line join=round,line cap=round] (147.97,139.21) -- (152.47,134.71);

\path[draw=drawColor,line width= 0.4pt,line join=round,line cap=round] (148.09,134.75) -- (152.59,139.25);

\path[draw=drawColor,line width= 0.4pt,line join=round,line cap=round] (148.09,139.25) -- (152.59,134.75);

\path[draw=drawColor,line width= 0.4pt,line join=round,line cap=round] (148.20,134.77) -- (152.70,139.27);

\path[draw=drawColor,line width= 0.4pt,line join=round,line cap=round] (148.20,139.27) -- (152.70,134.77);

\path[draw=drawColor,line width= 0.4pt,line join=round,line cap=round] (148.32,134.94) -- (152.82,139.44);

\path[draw=drawColor,line width= 0.4pt,line join=round,line cap=round] (148.32,139.44) -- (152.82,134.94);

\path[draw=drawColor,line width= 0.4pt,line join=round,line cap=round] (148.44,135.09) -- (152.94,139.59);

\path[draw=drawColor,line width= 0.4pt,line join=round,line cap=round] (148.44,139.59) -- (152.94,135.09);

\path[draw=drawColor,line width= 0.4pt,line join=round,line cap=round] (148.56,135.23) -- (153.06,139.73);

\path[draw=drawColor,line width= 0.4pt,line join=round,line cap=round] (148.56,139.73) -- (153.06,135.23);

\path[draw=drawColor,line width= 0.4pt,line join=round,line cap=round] (148.67,135.25) -- (153.17,139.75);

\path[draw=drawColor,line width= 0.4pt,line join=round,line cap=round] (148.67,139.75) -- (153.17,135.25);

\path[draw=drawColor,line width= 0.4pt,line join=round,line cap=round] (148.79,135.51) -- (153.29,140.01);

\path[draw=drawColor,line width= 0.4pt,line join=round,line cap=round] (148.79,140.01) -- (153.29,135.51);

\path[draw=drawColor,line width= 0.4pt,line join=round,line cap=round] (148.91,135.62) -- (153.41,140.12);

\path[draw=drawColor,line width= 0.4pt,line join=round,line cap=round] (148.91,140.12) -- (153.41,135.62);

\path[draw=drawColor,line width= 0.4pt,line join=round,line cap=round] (149.02,135.67) -- (153.52,140.17);

\path[draw=drawColor,line width= 0.4pt,line join=round,line cap=round] (149.02,140.17) -- (153.52,135.67);

\path[draw=drawColor,line width= 0.4pt,line join=round,line cap=round] (149.14,135.83) -- (153.64,140.33);

\path[draw=drawColor,line width= 0.4pt,line join=round,line cap=round] (149.14,140.33) -- (153.64,135.83);

\path[draw=drawColor,line width= 0.4pt,line join=round,line cap=round] (149.26,135.92) -- (153.76,140.42);

\path[draw=drawColor,line width= 0.4pt,line join=round,line cap=round] (149.26,140.42) -- (153.76,135.92);

\path[draw=drawColor,line width= 0.4pt,line join=round,line cap=round] (149.37,135.92) -- (153.87,140.42);

\path[draw=drawColor,line width= 0.4pt,line join=round,line cap=round] (149.37,140.42) -- (153.87,135.92);

\path[draw=drawColor,line width= 0.4pt,line join=round,line cap=round] (149.49,135.96) -- (153.99,140.46);

\path[draw=drawColor,line width= 0.4pt,line join=round,line cap=round] (149.49,140.46) -- (153.99,135.96);

\path[draw=drawColor,line width= 0.4pt,line join=round,line cap=round] (149.61,135.98) -- (154.11,140.48);

\path[draw=drawColor,line width= 0.4pt,line join=round,line cap=round] (149.61,140.48) -- (154.11,135.98);

\path[draw=drawColor,line width= 0.4pt,line join=round,line cap=round] (149.72,136.07) -- (154.22,140.57);

\path[draw=drawColor,line width= 0.4pt,line join=round,line cap=round] (149.72,140.57) -- (154.22,136.07);

\path[draw=drawColor,line width= 0.4pt,line join=round,line cap=round] (149.84,136.13) -- (154.34,140.63);

\path[draw=drawColor,line width= 0.4pt,line join=round,line cap=round] (149.84,140.63) -- (154.34,136.13);

\path[draw=drawColor,line width= 0.4pt,line join=round,line cap=round] (149.96,136.16) -- (154.46,140.66);

\path[draw=drawColor,line width= 0.4pt,line join=round,line cap=round] (149.96,140.66) -- (154.46,136.16);

\path[draw=drawColor,line width= 0.4pt,line join=round,line cap=round] (150.08,136.20) -- (154.58,140.70);

\path[draw=drawColor,line width= 0.4pt,line join=round,line cap=round] (150.08,140.70) -- (154.58,136.20);

\path[draw=drawColor,line width= 0.4pt,line join=round,line cap=round] (150.19,136.35) -- (154.69,140.85);

\path[draw=drawColor,line width= 0.4pt,line join=round,line cap=round] (150.19,140.85) -- (154.69,136.35);

\path[draw=drawColor,line width= 0.4pt,line join=round,line cap=round] (150.31,136.38) -- (154.81,140.88);

\path[draw=drawColor,line width= 0.4pt,line join=round,line cap=round] (150.31,140.88) -- (154.81,136.38);

\path[draw=drawColor,line width= 0.4pt,line join=round,line cap=round] (150.43,136.57) -- (154.93,141.07);

\path[draw=drawColor,line width= 0.4pt,line join=round,line cap=round] (150.43,141.07) -- (154.93,136.57);

\path[draw=drawColor,line width= 0.4pt,line join=round,line cap=round] (150.54,136.59) -- (155.04,141.09);

\path[draw=drawColor,line width= 0.4pt,line join=round,line cap=round] (150.54,141.09) -- (155.04,136.59);

\path[draw=drawColor,line width= 0.4pt,line join=round,line cap=round] (150.66,136.59) -- (155.16,141.09);

\path[draw=drawColor,line width= 0.4pt,line join=round,line cap=round] (150.66,141.09) -- (155.16,136.59);

\path[draw=drawColor,line width= 0.4pt,line join=round,line cap=round] (150.78,136.70) -- (155.28,141.20);

\path[draw=drawColor,line width= 0.4pt,line join=round,line cap=round] (150.78,141.20) -- (155.28,136.70);

\path[draw=drawColor,line width= 0.4pt,line join=round,line cap=round] (150.89,137.01) -- (155.39,141.51);

\path[draw=drawColor,line width= 0.4pt,line join=round,line cap=round] (150.89,141.51) -- (155.39,137.01);

\path[draw=drawColor,line width= 0.4pt,line join=round,line cap=round] (151.01,137.02) -- (155.51,141.52);

\path[draw=drawColor,line width= 0.4pt,line join=round,line cap=round] (151.01,141.52) -- (155.51,137.02);

\path[draw=drawColor,line width= 0.4pt,line join=round,line cap=round] (151.13,137.07) -- (155.63,141.57);

\path[draw=drawColor,line width= 0.4pt,line join=round,line cap=round] (151.13,141.57) -- (155.63,137.07);

\path[draw=drawColor,line width= 0.4pt,line join=round,line cap=round] (151.24,137.12) -- (155.74,141.62);

\path[draw=drawColor,line width= 0.4pt,line join=round,line cap=round] (151.24,141.62) -- (155.74,137.12);

\path[draw=drawColor,line width= 0.4pt,line join=round,line cap=round] (151.36,137.23) -- (155.86,141.73);

\path[draw=drawColor,line width= 0.4pt,line join=round,line cap=round] (151.36,141.73) -- (155.86,137.23);

\path[draw=drawColor,line width= 0.4pt,line join=round,line cap=round] (151.48,137.31) -- (155.98,141.81);

\path[draw=drawColor,line width= 0.4pt,line join=round,line cap=round] (151.48,141.81) -- (155.98,137.31);

\path[draw=drawColor,line width= 0.4pt,line join=round,line cap=round] (151.60,137.36) -- (156.10,141.86);

\path[draw=drawColor,line width= 0.4pt,line join=round,line cap=round] (151.60,141.86) -- (156.10,137.36);

\path[draw=drawColor,line width= 0.4pt,line join=round,line cap=round] (151.71,137.38) -- (156.21,141.88);

\path[draw=drawColor,line width= 0.4pt,line join=round,line cap=round] (151.71,141.88) -- (156.21,137.38);

\path[draw=drawColor,line width= 0.4pt,line join=round,line cap=round] (151.83,137.38) -- (156.33,141.88);

\path[draw=drawColor,line width= 0.4pt,line join=round,line cap=round] (151.83,141.88) -- (156.33,137.38);

\path[draw=drawColor,line width= 0.4pt,line join=round,line cap=round] (151.95,137.40) -- (156.45,141.90);

\path[draw=drawColor,line width= 0.4pt,line join=round,line cap=round] (151.95,141.90) -- (156.45,137.40);

\path[draw=drawColor,line width= 0.4pt,line join=round,line cap=round] (152.06,137.66) -- (156.56,142.16);

\path[draw=drawColor,line width= 0.4pt,line join=round,line cap=round] (152.06,142.16) -- (156.56,137.66);

\path[draw=drawColor,line width= 0.4pt,line join=round,line cap=round] (152.18,137.67) -- (156.68,142.17);

\path[draw=drawColor,line width= 0.4pt,line join=round,line cap=round] (152.18,142.17) -- (156.68,137.67);

\path[draw=drawColor,line width= 0.4pt,line join=round,line cap=round] (152.30,137.73) -- (156.80,142.23);

\path[draw=drawColor,line width= 0.4pt,line join=round,line cap=round] (152.30,142.23) -- (156.80,137.73);

\path[draw=drawColor,line width= 0.4pt,line join=round,line cap=round] (152.41,137.83) -- (156.91,142.33);

\path[draw=drawColor,line width= 0.4pt,line join=round,line cap=round] (152.41,142.33) -- (156.91,137.83);

\path[draw=drawColor,line width= 0.4pt,line join=round,line cap=round] (152.53,137.87) -- (157.03,142.37);

\path[draw=drawColor,line width= 0.4pt,line join=round,line cap=round] (152.53,142.37) -- (157.03,137.87);

\path[draw=drawColor,line width= 0.4pt,line join=round,line cap=round] (152.65,137.88) -- (157.15,142.38);

\path[draw=drawColor,line width= 0.4pt,line join=round,line cap=round] (152.65,142.38) -- (157.15,137.88);

\path[draw=drawColor,line width= 0.4pt,line join=round,line cap=round] (152.76,137.97) -- (157.26,142.47);

\path[draw=drawColor,line width= 0.4pt,line join=round,line cap=round] (152.76,142.47) -- (157.26,137.97);

\path[draw=drawColor,line width= 0.4pt,line join=round,line cap=round] (152.88,138.03) -- (157.38,142.53);

\path[draw=drawColor,line width= 0.4pt,line join=round,line cap=round] (152.88,142.53) -- (157.38,138.03);

\path[draw=drawColor,line width= 0.4pt,line join=round,line cap=round] (153.00,138.07) -- (157.50,142.57);

\path[draw=drawColor,line width= 0.4pt,line join=round,line cap=round] (153.00,142.57) -- (157.50,138.07);

\path[draw=drawColor,line width= 0.4pt,line join=round,line cap=round] (153.12,138.40) -- (157.62,142.90);

\path[draw=drawColor,line width= 0.4pt,line join=round,line cap=round] (153.12,142.90) -- (157.62,138.40);

\path[draw=drawColor,line width= 0.4pt,line join=round,line cap=round] (153.23,138.55) -- (157.73,143.05);

\path[draw=drawColor,line width= 0.4pt,line join=round,line cap=round] (153.23,143.05) -- (157.73,138.55);

\path[draw=drawColor,line width= 0.4pt,line join=round,line cap=round] (153.35,138.68) -- (157.85,143.18);

\path[draw=drawColor,line width= 0.4pt,line join=round,line cap=round] (153.35,143.18) -- (157.85,138.68);

\path[draw=drawColor,line width= 0.4pt,line join=round,line cap=round] (153.47,138.71) -- (157.97,143.21);

\path[draw=drawColor,line width= 0.4pt,line join=round,line cap=round] (153.47,143.21) -- (157.97,138.71);

\path[draw=drawColor,line width= 0.4pt,line join=round,line cap=round] (153.58,138.74) -- (158.08,143.24);

\path[draw=drawColor,line width= 0.4pt,line join=round,line cap=round] (153.58,143.24) -- (158.08,138.74);

\path[draw=drawColor,line width= 0.4pt,line join=round,line cap=round] (153.70,138.75) -- (158.20,143.25);

\path[draw=drawColor,line width= 0.4pt,line join=round,line cap=round] (153.70,143.25) -- (158.20,138.75);

\path[draw=drawColor,line width= 0.4pt,line join=round,line cap=round] (153.82,138.94) -- (158.32,143.44);

\path[draw=drawColor,line width= 0.4pt,line join=round,line cap=round] (153.82,143.44) -- (158.32,138.94);

\path[draw=drawColor,line width= 0.4pt,line join=round,line cap=round] (153.93,139.05) -- (158.43,143.55);

\path[draw=drawColor,line width= 0.4pt,line join=round,line cap=round] (153.93,143.55) -- (158.43,139.05);

\path[draw=drawColor,line width= 0.4pt,line join=round,line cap=round] (154.05,139.24) -- (158.55,143.74);

\path[draw=drawColor,line width= 0.4pt,line join=round,line cap=round] (154.05,143.74) -- (158.55,139.24);

\path[draw=drawColor,line width= 0.4pt,line join=round,line cap=round] (154.17,139.41) -- (158.67,143.91);

\path[draw=drawColor,line width= 0.4pt,line join=round,line cap=round] (154.17,143.91) -- (158.67,139.41);

\path[draw=drawColor,line width= 0.4pt,line join=round,line cap=round] (154.29,139.55) -- (158.79,144.05);

\path[draw=drawColor,line width= 0.4pt,line join=round,line cap=round] (154.29,144.05) -- (158.79,139.55);

\path[draw=drawColor,line width= 0.4pt,line join=round,line cap=round] (154.40,139.71) -- (158.90,144.21);

\path[draw=drawColor,line width= 0.4pt,line join=round,line cap=round] (154.40,144.21) -- (158.90,139.71);

\path[draw=drawColor,line width= 0.4pt,line join=round,line cap=round] (154.52,139.84) -- (159.02,144.34);

\path[draw=drawColor,line width= 0.4pt,line join=round,line cap=round] (154.52,144.34) -- (159.02,139.84);

\path[draw=drawColor,line width= 0.4pt,line join=round,line cap=round] (154.64,140.00) -- (159.14,144.50);

\path[draw=drawColor,line width= 0.4pt,line join=round,line cap=round] (154.64,144.50) -- (159.14,140.00);

\path[draw=drawColor,line width= 0.4pt,line join=round,line cap=round] (154.75,140.04) -- (159.25,144.54);

\path[draw=drawColor,line width= 0.4pt,line join=round,line cap=round] (154.75,144.54) -- (159.25,140.04);

\path[draw=drawColor,line width= 0.4pt,line join=round,line cap=round] (154.87,140.08) -- (159.37,144.58);

\path[draw=drawColor,line width= 0.4pt,line join=round,line cap=round] (154.87,144.58) -- (159.37,140.08);

\path[draw=drawColor,line width= 0.4pt,line join=round,line cap=round] (154.99,140.23) -- (159.49,144.73);

\path[draw=drawColor,line width= 0.4pt,line join=round,line cap=round] (154.99,144.73) -- (159.49,140.23);

\path[draw=drawColor,line width= 0.4pt,line join=round,line cap=round] (155.10,140.29) -- (159.60,144.79);

\path[draw=drawColor,line width= 0.4pt,line join=round,line cap=round] (155.10,144.79) -- (159.60,140.29);

\path[draw=drawColor,line width= 0.4pt,line join=round,line cap=round] (155.22,140.36) -- (159.72,144.86);

\path[draw=drawColor,line width= 0.4pt,line join=round,line cap=round] (155.22,144.86) -- (159.72,140.36);

\path[draw=drawColor,line width= 0.4pt,line join=round,line cap=round] (155.34,140.38) -- (159.84,144.88);

\path[draw=drawColor,line width= 0.4pt,line join=round,line cap=round] (155.34,144.88) -- (159.84,140.38);

\path[draw=drawColor,line width= 0.4pt,line join=round,line cap=round] (155.45,140.49) -- (159.95,144.99);

\path[draw=drawColor,line width= 0.4pt,line join=round,line cap=round] (155.45,144.99) -- (159.95,140.49);

\path[draw=drawColor,line width= 0.4pt,line join=round,line cap=round] (155.57,140.57) -- (160.07,145.07);

\path[draw=drawColor,line width= 0.4pt,line join=round,line cap=round] (155.57,145.07) -- (160.07,140.57);

\path[draw=drawColor,line width= 0.4pt,line join=round,line cap=round] (155.69,140.57) -- (160.19,145.07);

\path[draw=drawColor,line width= 0.4pt,line join=round,line cap=round] (155.69,145.07) -- (160.19,140.57);

\path[draw=drawColor,line width= 0.4pt,line join=round,line cap=round] (155.81,140.64) -- (160.31,145.14);

\path[draw=drawColor,line width= 0.4pt,line join=round,line cap=round] (155.81,145.14) -- (160.31,140.64);

\path[draw=drawColor,line width= 0.4pt,line join=round,line cap=round] (155.92,140.66) -- (160.42,145.16);

\path[draw=drawColor,line width= 0.4pt,line join=round,line cap=round] (155.92,145.16) -- (160.42,140.66);

\path[draw=drawColor,line width= 0.4pt,line join=round,line cap=round] (156.04,140.66) -- (160.54,145.16);

\path[draw=drawColor,line width= 0.4pt,line join=round,line cap=round] (156.04,145.16) -- (160.54,140.66);

\path[draw=drawColor,line width= 0.4pt,line join=round,line cap=round] (156.16,140.74) -- (160.66,145.24);

\path[draw=drawColor,line width= 0.4pt,line join=round,line cap=round] (156.16,145.24) -- (160.66,140.74);

\path[draw=drawColor,line width= 0.4pt,line join=round,line cap=round] (156.27,140.75) -- (160.77,145.25);

\path[draw=drawColor,line width= 0.4pt,line join=round,line cap=round] (156.27,145.25) -- (160.77,140.75);

\path[draw=drawColor,line width= 0.4pt,line join=round,line cap=round] (156.39,140.79) -- (160.89,145.29);

\path[draw=drawColor,line width= 0.4pt,line join=round,line cap=round] (156.39,145.29) -- (160.89,140.79);

\path[draw=drawColor,line width= 0.4pt,line join=round,line cap=round] (156.51,140.81) -- (161.01,145.31);

\path[draw=drawColor,line width= 0.4pt,line join=round,line cap=round] (156.51,145.31) -- (161.01,140.81);

\path[draw=drawColor,line width= 0.4pt,line join=round,line cap=round] (156.62,140.97) -- (161.12,145.47);

\path[draw=drawColor,line width= 0.4pt,line join=round,line cap=round] (156.62,145.47) -- (161.12,140.97);

\path[draw=drawColor,line width= 0.4pt,line join=round,line cap=round] (156.74,140.97) -- (161.24,145.47);

\path[draw=drawColor,line width= 0.4pt,line join=round,line cap=round] (156.74,145.47) -- (161.24,140.97);

\path[draw=drawColor,line width= 0.4pt,line join=round,line cap=round] (156.86,141.04) -- (161.36,145.54);

\path[draw=drawColor,line width= 0.4pt,line join=round,line cap=round] (156.86,145.54) -- (161.36,141.04);

\path[draw=drawColor,line width= 0.4pt,line join=round,line cap=round] (156.97,141.06) -- (161.47,145.56);

\path[draw=drawColor,line width= 0.4pt,line join=round,line cap=round] (156.97,145.56) -- (161.47,141.06);

\path[draw=drawColor,line width= 0.4pt,line join=round,line cap=round] (157.09,141.15) -- (161.59,145.65);

\path[draw=drawColor,line width= 0.4pt,line join=round,line cap=round] (157.09,145.65) -- (161.59,141.15);

\path[draw=drawColor,line width= 0.4pt,line join=round,line cap=round] (157.21,141.19) -- (161.71,145.69);

\path[draw=drawColor,line width= 0.4pt,line join=round,line cap=round] (157.21,145.69) -- (161.71,141.19);

\path[draw=drawColor,line width= 0.4pt,line join=round,line cap=round] (157.33,141.21) -- (161.83,145.71);

\path[draw=drawColor,line width= 0.4pt,line join=round,line cap=round] (157.33,145.71) -- (161.83,141.21);

\path[draw=drawColor,line width= 0.4pt,line join=round,line cap=round] (157.44,141.35) -- (161.94,145.85);

\path[draw=drawColor,line width= 0.4pt,line join=round,line cap=round] (157.44,145.85) -- (161.94,141.35);

\path[draw=drawColor,line width= 0.4pt,line join=round,line cap=round] (157.56,141.36) -- (162.06,145.86);

\path[draw=drawColor,line width= 0.4pt,line join=round,line cap=round] (157.56,145.86) -- (162.06,141.36);

\path[draw=drawColor,line width= 0.4pt,line join=round,line cap=round] (157.68,141.53) -- (162.18,146.03);

\path[draw=drawColor,line width= 0.4pt,line join=round,line cap=round] (157.68,146.03) -- (162.18,141.53);

\path[draw=drawColor,line width= 0.4pt,line join=round,line cap=round] (157.79,141.59) -- (162.29,146.09);

\path[draw=drawColor,line width= 0.4pt,line join=round,line cap=round] (157.79,146.09) -- (162.29,141.59);

\path[draw=drawColor,line width= 0.4pt,line join=round,line cap=round] (157.91,141.64) -- (162.41,146.14);

\path[draw=drawColor,line width= 0.4pt,line join=round,line cap=round] (157.91,146.14) -- (162.41,141.64);

\path[draw=drawColor,line width= 0.4pt,line join=round,line cap=round] (158.03,141.83) -- (162.53,146.33);

\path[draw=drawColor,line width= 0.4pt,line join=round,line cap=round] (158.03,146.33) -- (162.53,141.83);

\path[draw=drawColor,line width= 0.4pt,line join=round,line cap=round] (158.14,141.85) -- (162.64,146.35);

\path[draw=drawColor,line width= 0.4pt,line join=round,line cap=round] (158.14,146.35) -- (162.64,141.85);

\path[draw=drawColor,line width= 0.4pt,line join=round,line cap=round] (158.26,141.87) -- (162.76,146.37);

\path[draw=drawColor,line width= 0.4pt,line join=round,line cap=round] (158.26,146.37) -- (162.76,141.87);

\path[draw=drawColor,line width= 0.4pt,line join=round,line cap=round] (158.38,141.89) -- (162.88,146.39);

\path[draw=drawColor,line width= 0.4pt,line join=round,line cap=round] (158.38,146.39) -- (162.88,141.89);

\path[draw=drawColor,line width= 0.4pt,line join=round,line cap=round] (158.50,142.02) -- (163.00,146.52);

\path[draw=drawColor,line width= 0.4pt,line join=round,line cap=round] (158.50,146.52) -- (163.00,142.02);

\path[draw=drawColor,line width= 0.4pt,line join=round,line cap=round] (158.61,142.03) -- (163.11,146.53);

\path[draw=drawColor,line width= 0.4pt,line join=round,line cap=round] (158.61,146.53) -- (163.11,142.03);

\path[draw=drawColor,line width= 0.4pt,line join=round,line cap=round] (158.73,142.12) -- (163.23,146.62);

\path[draw=drawColor,line width= 0.4pt,line join=round,line cap=round] (158.73,146.62) -- (163.23,142.12);

\path[draw=drawColor,line width= 0.4pt,line join=round,line cap=round] (158.85,142.19) -- (163.35,146.69);

\path[draw=drawColor,line width= 0.4pt,line join=round,line cap=round] (158.85,146.69) -- (163.35,142.19);

\path[draw=drawColor,line width= 0.4pt,line join=round,line cap=round] (158.96,142.19) -- (163.46,146.69);

\path[draw=drawColor,line width= 0.4pt,line join=round,line cap=round] (158.96,146.69) -- (163.46,142.19);

\path[draw=drawColor,line width= 0.4pt,line join=round,line cap=round] (159.08,142.27) -- (163.58,146.77);

\path[draw=drawColor,line width= 0.4pt,line join=round,line cap=round] (159.08,146.77) -- (163.58,142.27);

\path[draw=drawColor,line width= 0.4pt,line join=round,line cap=round] (159.20,142.40) -- (163.70,146.90);

\path[draw=drawColor,line width= 0.4pt,line join=round,line cap=round] (159.20,146.90) -- (163.70,142.40);

\path[draw=drawColor,line width= 0.4pt,line join=round,line cap=round] (159.31,142.43) -- (163.81,146.93);

\path[draw=drawColor,line width= 0.4pt,line join=round,line cap=round] (159.31,146.93) -- (163.81,142.43);

\path[draw=drawColor,line width= 0.4pt,line join=round,line cap=round] (159.43,142.47) -- (163.93,146.97);

\path[draw=drawColor,line width= 0.4pt,line join=round,line cap=round] (159.43,146.97) -- (163.93,142.47);

\path[draw=drawColor,line width= 0.4pt,line join=round,line cap=round] (159.55,142.59) -- (164.05,147.09);

\path[draw=drawColor,line width= 0.4pt,line join=round,line cap=round] (159.55,147.09) -- (164.05,142.59);

\path[draw=drawColor,line width= 0.4pt,line join=round,line cap=round] (159.66,142.63) -- (164.16,147.13);

\path[draw=drawColor,line width= 0.4pt,line join=round,line cap=round] (159.66,147.13) -- (164.16,142.63);

\path[draw=drawColor,line width= 0.4pt,line join=round,line cap=round] (159.78,142.68) -- (164.28,147.18);

\path[draw=drawColor,line width= 0.4pt,line join=round,line cap=round] (159.78,147.18) -- (164.28,142.68);

\path[draw=drawColor,line width= 0.4pt,line join=round,line cap=round] (159.90,142.74) -- (164.40,147.24);

\path[draw=drawColor,line width= 0.4pt,line join=round,line cap=round] (159.90,147.24) -- (164.40,142.74);

\path[draw=drawColor,line width= 0.4pt,line join=round,line cap=round] (160.02,142.78) -- (164.52,147.28);

\path[draw=drawColor,line width= 0.4pt,line join=round,line cap=round] (160.02,147.28) -- (164.52,142.78);

\path[draw=drawColor,line width= 0.4pt,line join=round,line cap=round] (160.13,142.93) -- (164.63,147.43);

\path[draw=drawColor,line width= 0.4pt,line join=round,line cap=round] (160.13,147.43) -- (164.63,142.93);

\path[draw=drawColor,line width= 0.4pt,line join=round,line cap=round] (160.25,143.12) -- (164.75,147.62);

\path[draw=drawColor,line width= 0.4pt,line join=round,line cap=round] (160.25,147.62) -- (164.75,143.12);

\path[draw=drawColor,line width= 0.4pt,line join=round,line cap=round] (160.37,143.14) -- (164.87,147.64);

\path[draw=drawColor,line width= 0.4pt,line join=round,line cap=round] (160.37,147.64) -- (164.87,143.14);

\path[draw=drawColor,line width= 0.4pt,line join=round,line cap=round] (160.48,143.24) -- (164.98,147.74);

\path[draw=drawColor,line width= 0.4pt,line join=round,line cap=round] (160.48,147.74) -- (164.98,143.24);

\path[draw=drawColor,line width= 0.4pt,line join=round,line cap=round] (160.60,143.25) -- (165.10,147.75);

\path[draw=drawColor,line width= 0.4pt,line join=round,line cap=round] (160.60,147.75) -- (165.10,143.25);

\path[draw=drawColor,line width= 0.4pt,line join=round,line cap=round] (160.72,143.33) -- (165.22,147.83);

\path[draw=drawColor,line width= 0.4pt,line join=round,line cap=round] (160.72,147.83) -- (165.22,143.33);

\path[draw=drawColor,line width= 0.4pt,line join=round,line cap=round] (160.83,143.39) -- (165.33,147.89);

\path[draw=drawColor,line width= 0.4pt,line join=round,line cap=round] (160.83,147.89) -- (165.33,143.39);

\path[draw=drawColor,line width= 0.4pt,line join=round,line cap=round] (160.95,143.62) -- (165.45,148.12);

\path[draw=drawColor,line width= 0.4pt,line join=round,line cap=round] (160.95,148.12) -- (165.45,143.62);

\path[draw=drawColor,line width= 0.4pt,line join=round,line cap=round] (161.07,143.66) -- (165.57,148.16);

\path[draw=drawColor,line width= 0.4pt,line join=round,line cap=round] (161.07,148.16) -- (165.57,143.66);

\path[draw=drawColor,line width= 0.4pt,line join=round,line cap=round] (161.18,143.67) -- (165.68,148.17);

\path[draw=drawColor,line width= 0.4pt,line join=round,line cap=round] (161.18,148.17) -- (165.68,143.67);

\path[draw=drawColor,line width= 0.4pt,line join=round,line cap=round] (161.30,143.78) -- (165.80,148.28);

\path[draw=drawColor,line width= 0.4pt,line join=round,line cap=round] (161.30,148.28) -- (165.80,143.78);

\path[draw=drawColor,line width= 0.4pt,line join=round,line cap=round] (161.42,143.92) -- (165.92,148.42);

\path[draw=drawColor,line width= 0.4pt,line join=round,line cap=round] (161.42,148.42) -- (165.92,143.92);

\path[draw=drawColor,line width= 0.4pt,line join=round,line cap=round] (161.54,143.97) -- (166.04,148.47);

\path[draw=drawColor,line width= 0.4pt,line join=round,line cap=round] (161.54,148.47) -- (166.04,143.97);

\path[draw=drawColor,line width= 0.4pt,line join=round,line cap=round] (161.65,144.03) -- (166.15,148.53);

\path[draw=drawColor,line width= 0.4pt,line join=round,line cap=round] (161.65,148.53) -- (166.15,144.03);

\path[draw=drawColor,line width= 0.4pt,line join=round,line cap=round] (161.77,144.11) -- (166.27,148.61);

\path[draw=drawColor,line width= 0.4pt,line join=round,line cap=round] (161.77,148.61) -- (166.27,144.11);

\path[draw=drawColor,line width= 0.4pt,line join=round,line cap=round] (161.89,144.17) -- (166.39,148.67);

\path[draw=drawColor,line width= 0.4pt,line join=round,line cap=round] (161.89,148.67) -- (166.39,144.17);

\path[draw=drawColor,line width= 0.4pt,line join=round,line cap=round] (162.00,144.18) -- (166.50,148.68);

\path[draw=drawColor,line width= 0.4pt,line join=round,line cap=round] (162.00,148.68) -- (166.50,144.18);

\path[draw=drawColor,line width= 0.4pt,line join=round,line cap=round] (162.12,144.50) -- (166.62,149.00);

\path[draw=drawColor,line width= 0.4pt,line join=round,line cap=round] (162.12,149.00) -- (166.62,144.50);

\path[draw=drawColor,line width= 0.4pt,line join=round,line cap=round] (162.24,144.66) -- (166.74,149.16);

\path[draw=drawColor,line width= 0.4pt,line join=round,line cap=round] (162.24,149.16) -- (166.74,144.66);

\path[draw=drawColor,line width= 0.4pt,line join=round,line cap=round] (162.35,144.77) -- (166.85,149.27);

\path[draw=drawColor,line width= 0.4pt,line join=round,line cap=round] (162.35,149.27) -- (166.85,144.77);

\path[draw=drawColor,line width= 0.4pt,line join=round,line cap=round] (162.47,144.85) -- (166.97,149.35);

\path[draw=drawColor,line width= 0.4pt,line join=round,line cap=round] (162.47,149.35) -- (166.97,144.85);

\path[draw=drawColor,line width= 0.4pt,line join=round,line cap=round] (162.59,144.90) -- (167.09,149.40);

\path[draw=drawColor,line width= 0.4pt,line join=round,line cap=round] (162.59,149.40) -- (167.09,144.90);

\path[draw=drawColor,line width= 0.4pt,line join=round,line cap=round] (162.70,144.93) -- (167.20,149.43);

\path[draw=drawColor,line width= 0.4pt,line join=round,line cap=round] (162.70,149.43) -- (167.20,144.93);

\path[draw=drawColor,line width= 0.4pt,line join=round,line cap=round] (162.82,145.05) -- (167.32,149.55);

\path[draw=drawColor,line width= 0.4pt,line join=round,line cap=round] (162.82,149.55) -- (167.32,145.05);

\path[draw=drawColor,line width= 0.4pt,line join=round,line cap=round] (162.94,145.30) -- (167.44,149.80);

\path[draw=drawColor,line width= 0.4pt,line join=round,line cap=round] (162.94,149.80) -- (167.44,145.30);

\path[draw=drawColor,line width= 0.4pt,line join=round,line cap=round] (163.06,145.47) -- (167.56,149.97);

\path[draw=drawColor,line width= 0.4pt,line join=round,line cap=round] (163.06,149.97) -- (167.56,145.47);

\path[draw=drawColor,line width= 0.4pt,line join=round,line cap=round] (163.17,145.53) -- (167.67,150.03);

\path[draw=drawColor,line width= 0.4pt,line join=round,line cap=round] (163.17,150.03) -- (167.67,145.53);

\path[draw=drawColor,line width= 0.4pt,line join=round,line cap=round] (163.29,145.73) -- (167.79,150.23);

\path[draw=drawColor,line width= 0.4pt,line join=round,line cap=round] (163.29,150.23) -- (167.79,145.73);

\path[draw=drawColor,line width= 0.4pt,line join=round,line cap=round] (163.41,145.96) -- (167.91,150.46);

\path[draw=drawColor,line width= 0.4pt,line join=round,line cap=round] (163.41,150.46) -- (167.91,145.96);

\path[draw=drawColor,line width= 0.4pt,line join=round,line cap=round] (163.52,145.99) -- (168.02,150.49);

\path[draw=drawColor,line width= 0.4pt,line join=round,line cap=round] (163.52,150.49) -- (168.02,145.99);

\path[draw=drawColor,line width= 0.4pt,line join=round,line cap=round] (163.64,146.11) -- (168.14,150.61);

\path[draw=drawColor,line width= 0.4pt,line join=round,line cap=round] (163.64,150.61) -- (168.14,146.11);

\path[draw=drawColor,line width= 0.4pt,line join=round,line cap=round] (163.76,146.50) -- (168.26,151.00);

\path[draw=drawColor,line width= 0.4pt,line join=round,line cap=round] (163.76,151.00) -- (168.26,146.50);

\path[draw=drawColor,line width= 0.4pt,line join=round,line cap=round] (163.87,146.56) -- (168.37,151.06);

\path[draw=drawColor,line width= 0.4pt,line join=round,line cap=round] (163.87,151.06) -- (168.37,146.56);

\path[draw=drawColor,line width= 0.4pt,line join=round,line cap=round] (163.99,146.69) -- (168.49,151.19);

\path[draw=drawColor,line width= 0.4pt,line join=round,line cap=round] (163.99,151.19) -- (168.49,146.69);

\path[draw=drawColor,line width= 0.4pt,line join=round,line cap=round] (164.11,146.86) -- (168.61,151.36);

\path[draw=drawColor,line width= 0.4pt,line join=round,line cap=round] (164.11,151.36) -- (168.61,146.86);

\path[draw=drawColor,line width= 0.4pt,line join=round,line cap=round] (164.23,147.06) -- (168.73,151.56);

\path[draw=drawColor,line width= 0.4pt,line join=round,line cap=round] (164.23,151.56) -- (168.73,147.06);

\path[draw=drawColor,line width= 0.4pt,line join=round,line cap=round] (164.34,147.14) -- (168.84,151.64);

\path[draw=drawColor,line width= 0.4pt,line join=round,line cap=round] (164.34,151.64) -- (168.84,147.14);

\path[draw=drawColor,line width= 0.4pt,line join=round,line cap=round] (164.46,147.21) -- (168.96,151.71);

\path[draw=drawColor,line width= 0.4pt,line join=round,line cap=round] (164.46,151.71) -- (168.96,147.21);

\path[draw=drawColor,line width= 0.4pt,line join=round,line cap=round] (164.58,147.22) -- (169.08,151.72);

\path[draw=drawColor,line width= 0.4pt,line join=round,line cap=round] (164.58,151.72) -- (169.08,147.22);

\path[draw=drawColor,line width= 0.4pt,line join=round,line cap=round] (164.69,147.28) -- (169.19,151.78);

\path[draw=drawColor,line width= 0.4pt,line join=round,line cap=round] (164.69,151.78) -- (169.19,147.28);

\path[draw=drawColor,line width= 0.4pt,line join=round,line cap=round] (164.81,147.57) -- (169.31,152.07);

\path[draw=drawColor,line width= 0.4pt,line join=round,line cap=round] (164.81,152.07) -- (169.31,147.57);

\path[draw=drawColor,line width= 0.4pt,line join=round,line cap=round] (164.93,147.61) -- (169.43,152.11);

\path[draw=drawColor,line width= 0.4pt,line join=round,line cap=round] (164.93,152.11) -- (169.43,147.61);

\path[draw=drawColor,line width= 0.4pt,line join=round,line cap=round] (165.04,147.70) -- (169.54,152.20);

\path[draw=drawColor,line width= 0.4pt,line join=round,line cap=round] (165.04,152.20) -- (169.54,147.70);

\path[draw=drawColor,line width= 0.4pt,line join=round,line cap=round] (165.16,147.82) -- (169.66,152.32);

\path[draw=drawColor,line width= 0.4pt,line join=round,line cap=round] (165.16,152.32) -- (169.66,147.82);

\path[draw=drawColor,line width= 0.4pt,line join=round,line cap=round] (165.28,147.88) -- (169.78,152.38);

\path[draw=drawColor,line width= 0.4pt,line join=round,line cap=round] (165.28,152.38) -- (169.78,147.88);

\path[draw=drawColor,line width= 0.4pt,line join=round,line cap=round] (165.39,148.01) -- (169.89,152.51);

\path[draw=drawColor,line width= 0.4pt,line join=round,line cap=round] (165.39,152.51) -- (169.89,148.01);

\path[draw=drawColor,line width= 0.4pt,line join=round,line cap=round] (165.51,148.17) -- (170.01,152.67);

\path[draw=drawColor,line width= 0.4pt,line join=round,line cap=round] (165.51,152.67) -- (170.01,148.17);

\path[draw=drawColor,line width= 0.4pt,line join=round,line cap=round] (165.63,148.25) -- (170.13,152.75);

\path[draw=drawColor,line width= 0.4pt,line join=round,line cap=round] (165.63,152.75) -- (170.13,148.25);

\path[draw=drawColor,line width= 0.4pt,line join=round,line cap=round] (165.75,148.43) -- (170.25,152.93);

\path[draw=drawColor,line width= 0.4pt,line join=round,line cap=round] (165.75,152.93) -- (170.25,148.43);

\path[draw=drawColor,line width= 0.4pt,line join=round,line cap=round] (165.86,148.70) -- (170.36,153.20);

\path[draw=drawColor,line width= 0.4pt,line join=round,line cap=round] (165.86,153.20) -- (170.36,148.70);

\path[draw=drawColor,line width= 0.4pt,line join=round,line cap=round] (165.98,148.88) -- (170.48,153.38);

\path[draw=drawColor,line width= 0.4pt,line join=round,line cap=round] (165.98,153.38) -- (170.48,148.88);

\path[draw=drawColor,line width= 0.4pt,line join=round,line cap=round] (166.10,148.97) -- (170.60,153.47);

\path[draw=drawColor,line width= 0.4pt,line join=round,line cap=round] (166.10,153.47) -- (170.60,148.97);

\path[draw=drawColor,line width= 0.4pt,line join=round,line cap=round] (166.21,149.06) -- (170.71,153.56);

\path[draw=drawColor,line width= 0.4pt,line join=round,line cap=round] (166.21,153.56) -- (170.71,149.06);

\path[draw=drawColor,line width= 0.4pt,line join=round,line cap=round] (166.33,149.10) -- (170.83,153.60);

\path[draw=drawColor,line width= 0.4pt,line join=round,line cap=round] (166.33,153.60) -- (170.83,149.10);

\path[draw=drawColor,line width= 0.4pt,line join=round,line cap=round] (166.45,149.28) -- (170.95,153.78);

\path[draw=drawColor,line width= 0.4pt,line join=round,line cap=round] (166.45,153.78) -- (170.95,149.28);

\path[draw=drawColor,line width= 0.4pt,line join=round,line cap=round] (166.56,149.30) -- (171.06,153.80);

\path[draw=drawColor,line width= 0.4pt,line join=round,line cap=round] (166.56,153.80) -- (171.06,149.30);

\path[draw=drawColor,line width= 0.4pt,line join=round,line cap=round] (166.68,149.52) -- (171.18,154.02);

\path[draw=drawColor,line width= 0.4pt,line join=round,line cap=round] (166.68,154.02) -- (171.18,149.52);

\path[draw=drawColor,line width= 0.4pt,line join=round,line cap=round] (166.80,149.61) -- (171.30,154.11);

\path[draw=drawColor,line width= 0.4pt,line join=round,line cap=round] (166.80,154.11) -- (171.30,149.61);

\path[draw=drawColor,line width= 0.4pt,line join=round,line cap=round] (166.91,149.68) -- (171.41,154.18);

\path[draw=drawColor,line width= 0.4pt,line join=round,line cap=round] (166.91,154.18) -- (171.41,149.68);

\path[draw=drawColor,line width= 0.4pt,line join=round,line cap=round] (167.03,149.70) -- (171.53,154.20);

\path[draw=drawColor,line width= 0.4pt,line join=round,line cap=round] (167.03,154.20) -- (171.53,149.70);

\path[draw=drawColor,line width= 0.4pt,line join=round,line cap=round] (167.15,149.80) -- (171.65,154.30);

\path[draw=drawColor,line width= 0.4pt,line join=round,line cap=round] (167.15,154.30) -- (171.65,149.80);

\path[draw=drawColor,line width= 0.4pt,line join=round,line cap=round] (167.27,149.87) -- (171.77,154.37);

\path[draw=drawColor,line width= 0.4pt,line join=round,line cap=round] (167.27,154.37) -- (171.77,149.87);

\path[draw=drawColor,line width= 0.4pt,line join=round,line cap=round] (167.38,150.10) -- (171.88,154.60);

\path[draw=drawColor,line width= 0.4pt,line join=round,line cap=round] (167.38,154.60) -- (171.88,150.10);

\path[draw=drawColor,line width= 0.4pt,line join=round,line cap=round] (167.50,150.27) -- (172.00,154.77);

\path[draw=drawColor,line width= 0.4pt,line join=round,line cap=round] (167.50,154.77) -- (172.00,150.27);

\path[draw=drawColor,line width= 0.4pt,line join=round,line cap=round] (167.62,150.39) -- (172.12,154.89);

\path[draw=drawColor,line width= 0.4pt,line join=round,line cap=round] (167.62,154.89) -- (172.12,150.39);

\path[draw=drawColor,line width= 0.4pt,line join=round,line cap=round] (167.73,150.40) -- (172.23,154.90);

\path[draw=drawColor,line width= 0.4pt,line join=round,line cap=round] (167.73,154.90) -- (172.23,150.40);
\end{scope}
\begin{scope}
\path[clip] (  0.00,  0.00) rectangle (361.35,180.67);
\definecolor{drawColor}{RGB}{0,0,0}

\path[draw=drawColor,line width= 0.4pt,line join=round,line cap=round] ( 52.57, 48.00) -- (169.98, 48.00);

\path[draw=drawColor,line width= 0.4pt,line join=round,line cap=round] ( 52.57, 48.00) -- ( 52.57, 42.00);

\path[draw=drawColor,line width= 0.4pt,line join=round,line cap=round] ( 76.06, 48.00) -- ( 76.06, 42.00);

\path[draw=drawColor,line width= 0.4pt,line join=round,line cap=round] ( 99.54, 48.00) -- ( 99.54, 42.00);

\path[draw=drawColor,line width= 0.4pt,line join=round,line cap=round] (123.02, 48.00) -- (123.02, 42.00);

\path[draw=drawColor,line width= 0.4pt,line join=round,line cap=round] (146.50, 48.00) -- (146.50, 42.00);

\path[draw=drawColor,line width= 0.4pt,line join=round,line cap=round] (169.98, 48.00) -- (169.98, 42.00);

\node[text=drawColor,anchor=base,inner sep=0pt, outer sep=0pt, scale=  1.00] at ( 52.57, 26.40) {0.0};

\node[text=drawColor,anchor=base,inner sep=0pt, outer sep=0pt, scale=  1.00] at ( 76.06, 26.40) {0.2};

\node[text=drawColor,anchor=base,inner sep=0pt, outer sep=0pt, scale=  1.00] at ( 99.54, 26.40) {0.4};

\node[text=drawColor,anchor=base,inner sep=0pt, outer sep=0pt, scale=  1.00] at (123.02, 26.40) {0.6};

\node[text=drawColor,anchor=base,inner sep=0pt, outer sep=0pt, scale=  1.00] at (146.50, 26.40) {0.8};

\node[text=drawColor,anchor=base,inner sep=0pt, outer sep=0pt, scale=  1.00] at (169.98, 26.40) {1.0};

\path[draw=drawColor,line width= 0.4pt,line join=round,line cap=round] ( 48.00, 51.96) -- ( 48.00,152.74);

\path[draw=drawColor,line width= 0.4pt,line join=round,line cap=round] ( 48.00, 51.96) -- ( 42.00, 51.96);

\path[draw=drawColor,line width= 0.4pt,line join=round,line cap=round] ( 48.00, 72.11) -- ( 42.00, 72.11);

\path[draw=drawColor,line width= 0.4pt,line join=round,line cap=round] ( 48.00, 92.27) -- ( 42.00, 92.27);

\path[draw=drawColor,line width= 0.4pt,line join=round,line cap=round] ( 48.00,112.43) -- ( 42.00,112.43);

\path[draw=drawColor,line width= 0.4pt,line join=round,line cap=round] ( 48.00,132.58) -- ( 42.00,132.58);

\path[draw=drawColor,line width= 0.4pt,line join=round,line cap=round] ( 48.00,152.74) -- ( 42.00,152.74);

\node[text=drawColor,rotate= 90.00,anchor=base,inner sep=0pt, outer sep=0pt, scale=  1.00] at ( 33.60, 51.96) {0.0};

\node[text=drawColor,rotate= 90.00,anchor=base,inner sep=0pt, outer sep=0pt, scale=  1.00] at ( 33.60, 92.27) {0.4};

\node[text=drawColor,rotate= 90.00,anchor=base,inner sep=0pt, outer sep=0pt, scale=  1.00] at ( 33.60,132.58) {0.8};

\path[draw=drawColor,line width= 0.4pt,line join=round,line cap=round] ( 48.00, 48.00) --
	(174.67, 48.00) --
	(174.67,156.67) --
	( 48.00,156.67) --
	( 48.00, 48.00);
\end{scope}
\begin{scope}
\path[clip] (  0.00,  0.00) rectangle (180.67,180.67);
\definecolor{drawColor}{RGB}{0,0,0}

\node[text=drawColor,anchor=base,inner sep=0pt, outer sep=0pt, scale=  1.20] at (111.34,164.53) {\bfseries PP-plot amb una NIG};

\node[text=drawColor,anchor=base,inner sep=0pt, outer sep=0pt, scale=  1.00] at (111.34,  2.40) {$F_n(x)$};

\node[text=drawColor,rotate= 90.00,anchor=base,inner sep=0pt, outer sep=0pt, scale=  1.00] at (  9.60,102.34) {$F(x)$};
\end{scope}
\begin{scope}
\path[clip] ( 48.00, 48.00) rectangle (174.67,156.67);
\definecolor{drawColor}{RGB}{255,0,0}

\path[draw=drawColor,line width= 1.2pt,dash pattern=on 4pt off 4pt ,line join=round,line cap=round] ( 48.00, 48.03) -- (174.67,156.77);
\end{scope}
\begin{scope}
\path[clip] (228.67, 48.00) rectangle (355.35,156.67);
\definecolor{drawColor}{gray}{0.40}

\path[draw=drawColor,line width= 0.4pt,line join=round,line cap=round] (257.26, 65.67) -- (261.76, 70.17);

\path[draw=drawColor,line width= 0.4pt,line join=round,line cap=round] (257.26, 70.17) -- (261.76, 65.67);

\path[draw=drawColor,line width= 0.4pt,line join=round,line cap=round] (284.61, 96.83) -- (289.11,101.33);

\path[draw=drawColor,line width= 0.4pt,line join=round,line cap=round] (284.61,101.33) -- (289.11, 96.83);

\path[draw=drawColor,line width= 0.4pt,line join=round,line cap=round] (299.72,111.07) -- (304.22,115.57);

\path[draw=drawColor,line width= 0.4pt,line join=round,line cap=round] (299.72,115.57) -- (304.22,111.07);

\path[draw=drawColor,line width= 0.4pt,line join=round,line cap=round] (296.78,107.70) -- (301.28,112.20);

\path[draw=drawColor,line width= 0.4pt,line join=round,line cap=round] (296.78,112.20) -- (301.28,107.70);

\path[draw=drawColor,line width= 0.4pt,line join=round,line cap=round] (280.41, 93.20) -- (284.91, 97.70);

\path[draw=drawColor,line width= 0.4pt,line join=round,line cap=round] (280.41, 97.70) -- (284.91, 93.20);

\path[draw=drawColor,line width= 0.4pt,line join=round,line cap=round] (282.81, 95.36) -- (287.31, 99.86);

\path[draw=drawColor,line width= 0.4pt,line join=round,line cap=round] (282.81, 99.86) -- (287.31, 95.36);

\path[draw=drawColor,line width= 0.4pt,line join=round,line cap=round] (284.97, 97.13) -- (289.47,101.63);

\path[draw=drawColor,line width= 0.4pt,line join=round,line cap=round] (284.97,101.63) -- (289.47, 97.13);

\path[draw=drawColor,line width= 0.4pt,line join=round,line cap=round] (286.86, 98.78) -- (291.36,103.28);

\path[draw=drawColor,line width= 0.4pt,line join=round,line cap=round] (286.86,103.28) -- (291.36, 98.78);

\path[draw=drawColor,line width= 0.4pt,line join=round,line cap=round] (274.96, 87.65) -- (279.46, 92.15);

\path[draw=drawColor,line width= 0.4pt,line join=round,line cap=round] (274.96, 92.15) -- (279.46, 87.65);

\path[draw=drawColor,line width= 0.4pt,line join=round,line cap=round] (285.38, 97.48) -- (289.88,101.98);

\path[draw=drawColor,line width= 0.4pt,line join=round,line cap=round] (285.38,101.98) -- (289.88, 97.48);

\path[draw=drawColor,line width= 0.4pt,line join=round,line cap=round] (290.59,101.90) -- (295.09,106.40);

\path[draw=drawColor,line width= 0.4pt,line join=round,line cap=round] (290.59,106.40) -- (295.09,101.90);

\path[draw=drawColor,line width= 0.4pt,line join=round,line cap=round] (300.64,111.83) -- (305.14,116.33);

\path[draw=drawColor,line width= 0.4pt,line join=round,line cap=round] (300.64,116.33) -- (305.14,111.83);

\path[draw=drawColor,line width= 0.4pt,line join=round,line cap=round] (301.06,112.03) -- (305.56,116.53);

\path[draw=drawColor,line width= 0.4pt,line join=round,line cap=round] (301.06,116.53) -- (305.56,112.03);

\path[draw=drawColor,line width= 0.4pt,line join=round,line cap=round] (279.52, 91.89) -- (284.02, 96.39);

\path[draw=drawColor,line width= 0.4pt,line join=round,line cap=round] (279.52, 96.39) -- (284.02, 91.89);

\path[draw=drawColor,line width= 0.4pt,line join=round,line cap=round] (292.82,104.24) -- (297.32,108.74);

\path[draw=drawColor,line width= 0.4pt,line join=round,line cap=round] (292.82,108.74) -- (297.32,104.24);

\path[draw=drawColor,line width= 0.4pt,line join=round,line cap=round] (241.20, 52.42) -- (245.70, 56.92);

\path[draw=drawColor,line width= 0.4pt,line join=round,line cap=round] (241.20, 56.92) -- (245.70, 52.42);

\path[draw=drawColor,line width= 0.4pt,line join=round,line cap=round] (291.13,102.35) -- (295.63,106.85);

\path[draw=drawColor,line width= 0.4pt,line join=round,line cap=round] (291.13,106.85) -- (295.63,102.35);

\path[draw=drawColor,line width= 0.4pt,line join=round,line cap=round] (289.51,100.96) -- (294.01,105.46);

\path[draw=drawColor,line width= 0.4pt,line join=round,line cap=round] (289.51,105.46) -- (294.01,100.96);

\path[draw=drawColor,line width= 0.4pt,line join=round,line cap=round] (289.31,100.80) -- (293.81,105.30);

\path[draw=drawColor,line width= 0.4pt,line join=round,line cap=round] (289.31,105.30) -- (293.81,100.80);

\path[draw=drawColor,line width= 0.4pt,line join=round,line cap=round] (297.23,108.10) -- (301.73,112.60);

\path[draw=drawColor,line width= 0.4pt,line join=round,line cap=round] (297.23,112.60) -- (301.73,108.10);

\path[draw=drawColor,line width= 0.4pt,line join=round,line cap=round] (317.19,125.43) -- (321.69,129.93);

\path[draw=drawColor,line width= 0.4pt,line join=round,line cap=round] (317.19,129.93) -- (321.69,125.43);

\path[draw=drawColor,line width= 0.4pt,line join=round,line cap=round] (272.57, 85.55) -- (277.07, 90.05);

\path[draw=drawColor,line width= 0.4pt,line join=round,line cap=round] (272.57, 90.05) -- (277.07, 85.55);

\path[draw=drawColor,line width= 0.4pt,line join=round,line cap=round] (302.83,113.81) -- (307.33,118.31);

\path[draw=drawColor,line width= 0.4pt,line join=round,line cap=round] (302.83,118.31) -- (307.33,113.81);

\path[draw=drawColor,line width= 0.4pt,line join=round,line cap=round] (286.27, 98.23) -- (290.77,102.73);

\path[draw=drawColor,line width= 0.4pt,line join=round,line cap=round] (286.27,102.73) -- (290.77, 98.23);

\path[draw=drawColor,line width= 0.4pt,line join=round,line cap=round] (301.58,112.39) -- (306.08,116.89);

\path[draw=drawColor,line width= 0.4pt,line join=round,line cap=round] (301.58,116.89) -- (306.08,112.39);

\path[draw=drawColor,line width= 0.4pt,line join=round,line cap=round] (282.49, 95.09) -- (286.99, 99.59);

\path[draw=drawColor,line width= 0.4pt,line join=round,line cap=round] (282.49, 99.59) -- (286.99, 95.09);

\path[draw=drawColor,line width= 0.4pt,line join=round,line cap=round] (285.11, 97.22) -- (289.61,101.72);

\path[draw=drawColor,line width= 0.4pt,line join=round,line cap=round] (285.11,101.72) -- (289.61, 97.22);

\path[draw=drawColor,line width= 0.4pt,line join=round,line cap=round] (290.75,101.98) -- (295.25,106.48);

\path[draw=drawColor,line width= 0.4pt,line join=round,line cap=round] (290.75,106.48) -- (295.25,101.98);

\path[draw=drawColor,line width= 0.4pt,line join=round,line cap=round] (302.40,113.42) -- (306.90,117.92);

\path[draw=drawColor,line width= 0.4pt,line join=round,line cap=round] (302.40,117.92) -- (306.90,113.42);

\path[draw=drawColor,line width= 0.4pt,line join=round,line cap=round] (291.77,102.96) -- (296.27,107.46);

\path[draw=drawColor,line width= 0.4pt,line join=round,line cap=round] (291.77,107.46) -- (296.27,102.96);

\path[draw=drawColor,line width= 0.4pt,line join=round,line cap=round] (284.38, 96.60) -- (288.88,101.10);

\path[draw=drawColor,line width= 0.4pt,line join=round,line cap=round] (284.38,101.10) -- (288.88, 96.60);

\path[draw=drawColor,line width= 0.4pt,line join=round,line cap=round] (277.06, 89.40) -- (281.56, 93.90);

\path[draw=drawColor,line width= 0.4pt,line join=round,line cap=round] (277.06, 93.90) -- (281.56, 89.40);

\path[draw=drawColor,line width= 0.4pt,line join=round,line cap=round] (298.68,109.91) -- (303.18,114.41);

\path[draw=drawColor,line width= 0.4pt,line join=round,line cap=round] (298.68,114.41) -- (303.18,109.91);

\path[draw=drawColor,line width= 0.4pt,line join=round,line cap=round] (280.89, 93.67) -- (285.39, 98.17);

\path[draw=drawColor,line width= 0.4pt,line join=round,line cap=round] (280.89, 98.17) -- (285.39, 93.67);

\path[draw=drawColor,line width= 0.4pt,line join=round,line cap=round] (297.69,108.89) -- (302.19,113.39);

\path[draw=drawColor,line width= 0.4pt,line join=round,line cap=round] (297.69,113.39) -- (302.19,108.89);

\path[draw=drawColor,line width= 0.4pt,line join=round,line cap=round] (294.99,106.04) -- (299.49,110.54);

\path[draw=drawColor,line width= 0.4pt,line join=round,line cap=round] (294.99,110.54) -- (299.49,106.04);

\path[draw=drawColor,line width= 0.4pt,line join=round,line cap=round] (277.94, 90.62) -- (282.44, 95.12);

\path[draw=drawColor,line width= 0.4pt,line join=round,line cap=round] (277.94, 95.12) -- (282.44, 90.62);

\path[draw=drawColor,line width= 0.4pt,line join=round,line cap=round] (281.29, 94.00) -- (285.79, 98.50);

\path[draw=drawColor,line width= 0.4pt,line join=round,line cap=round] (281.29, 98.50) -- (285.79, 94.00);

\path[draw=drawColor,line width= 0.4pt,line join=round,line cap=round] (290.52,101.83) -- (295.02,106.33);

\path[draw=drawColor,line width= 0.4pt,line join=round,line cap=round] (290.52,106.33) -- (295.02,101.83);

\path[draw=drawColor,line width= 0.4pt,line join=round,line cap=round] (287.99, 99.79) -- (292.49,104.29);

\path[draw=drawColor,line width= 0.4pt,line join=round,line cap=round] (287.99,104.29) -- (292.49, 99.79);

\path[draw=drawColor,line width= 0.4pt,line join=round,line cap=round] (279.76, 92.25) -- (284.26, 96.75);

\path[draw=drawColor,line width= 0.4pt,line join=round,line cap=round] (279.76, 96.75) -- (284.26, 92.25);

\path[draw=drawColor,line width= 0.4pt,line join=round,line cap=round] (287.38, 99.29) -- (291.88,103.79);

\path[draw=drawColor,line width= 0.4pt,line join=round,line cap=round] (287.38,103.79) -- (291.88, 99.29);

\path[draw=drawColor,line width= 0.4pt,line join=round,line cap=round] (276.90, 89.34) -- (281.40, 93.84);

\path[draw=drawColor,line width= 0.4pt,line join=round,line cap=round] (276.90, 93.84) -- (281.40, 89.34);

\path[draw=drawColor,line width= 0.4pt,line join=round,line cap=round] (296.15,107.13) -- (300.65,111.63);

\path[draw=drawColor,line width= 0.4pt,line join=round,line cap=round] (296.15,111.63) -- (300.65,107.13);

\path[draw=drawColor,line width= 0.4pt,line join=round,line cap=round] (281.00, 93.74) -- (285.50, 98.24);

\path[draw=drawColor,line width= 0.4pt,line join=round,line cap=round] (281.00, 98.24) -- (285.50, 93.74);

\path[draw=drawColor,line width= 0.4pt,line join=round,line cap=round] (287.52, 99.33) -- (292.02,103.83);

\path[draw=drawColor,line width= 0.4pt,line join=round,line cap=round] (287.52,103.83) -- (292.02, 99.33);

\path[draw=drawColor,line width= 0.4pt,line join=round,line cap=round] (293.68,105.04) -- (298.18,109.54);

\path[draw=drawColor,line width= 0.4pt,line join=round,line cap=round] (293.68,109.54) -- (298.18,105.04);

\path[draw=drawColor,line width= 0.4pt,line join=round,line cap=round] (280.36, 93.19) -- (284.86, 97.69);

\path[draw=drawColor,line width= 0.4pt,line join=round,line cap=round] (280.36, 97.69) -- (284.86, 93.19);

\path[draw=drawColor,line width= 0.4pt,line join=round,line cap=round] (296.08,107.11) -- (300.58,111.61);

\path[draw=drawColor,line width= 0.4pt,line join=round,line cap=round] (296.08,111.61) -- (300.58,107.11);

\path[draw=drawColor,line width= 0.4pt,line join=round,line cap=round] (275.88, 88.53) -- (280.38, 93.03);

\path[draw=drawColor,line width= 0.4pt,line join=round,line cap=round] (275.88, 93.03) -- (280.38, 88.53);

\path[draw=drawColor,line width= 0.4pt,line join=round,line cap=round] (316.31,124.56) -- (320.81,129.06);

\path[draw=drawColor,line width= 0.4pt,line join=round,line cap=round] (316.31,129.06) -- (320.81,124.56);

\path[draw=drawColor,line width= 0.4pt,line join=round,line cap=round] (278.35, 91.02) -- (282.85, 95.52);

\path[draw=drawColor,line width= 0.4pt,line join=round,line cap=round] (278.35, 95.52) -- (282.85, 91.02);

\path[draw=drawColor,line width= 0.4pt,line join=round,line cap=round] (298.81,109.94) -- (303.31,114.44);

\path[draw=drawColor,line width= 0.4pt,line join=round,line cap=round] (298.81,114.44) -- (303.31,109.94);

\path[draw=drawColor,line width= 0.4pt,line join=round,line cap=round] (288.52,100.24) -- (293.02,104.74);

\path[draw=drawColor,line width= 0.4pt,line join=round,line cap=round] (288.52,104.74) -- (293.02,100.24);

\path[draw=drawColor,line width= 0.4pt,line join=round,line cap=round] (283.58, 96.12) -- (288.08,100.62);

\path[draw=drawColor,line width= 0.4pt,line join=round,line cap=round] (283.58,100.62) -- (288.08, 96.12);

\path[draw=drawColor,line width= 0.4pt,line join=round,line cap=round] (283.38, 96.00) -- (287.88,100.50);

\path[draw=drawColor,line width= 0.4pt,line join=round,line cap=round] (283.38,100.50) -- (287.88, 96.00);

\path[draw=drawColor,line width= 0.4pt,line join=round,line cap=round] (289.75,101.21) -- (294.25,105.71);

\path[draw=drawColor,line width= 0.4pt,line join=round,line cap=round] (289.75,105.71) -- (294.25,101.21);

\path[draw=drawColor,line width= 0.4pt,line join=round,line cap=round] (287.50, 99.33) -- (292.00,103.83);

\path[draw=drawColor,line width= 0.4pt,line join=round,line cap=round] (287.50,103.83) -- (292.00, 99.33);

\path[draw=drawColor,line width= 0.4pt,line join=round,line cap=round] (304.74,114.82) -- (309.24,119.32);

\path[draw=drawColor,line width= 0.4pt,line join=round,line cap=round] (304.74,119.32) -- (309.24,114.82);

\path[draw=drawColor,line width= 0.4pt,line join=round,line cap=round] (277.63, 90.21) -- (282.13, 94.71);

\path[draw=drawColor,line width= 0.4pt,line join=round,line cap=round] (277.63, 94.71) -- (282.13, 90.21);

\path[draw=drawColor,line width= 0.4pt,line join=round,line cap=round] (285.70, 97.86) -- (290.20,102.36);

\path[draw=drawColor,line width= 0.4pt,line join=round,line cap=round] (285.70,102.36) -- (290.20, 97.86);

\path[draw=drawColor,line width= 0.4pt,line join=round,line cap=round] (294.70,105.94) -- (299.20,110.44);

\path[draw=drawColor,line width= 0.4pt,line join=round,line cap=round] (294.70,110.44) -- (299.20,105.94);

\path[draw=drawColor,line width= 0.4pt,line join=round,line cap=round] (299.76,111.13) -- (304.26,115.63);

\path[draw=drawColor,line width= 0.4pt,line join=round,line cap=round] (299.76,115.63) -- (304.26,111.13);

\path[draw=drawColor,line width= 0.4pt,line join=round,line cap=round] (280.67, 93.40) -- (285.17, 97.90);

\path[draw=drawColor,line width= 0.4pt,line join=round,line cap=round] (280.67, 97.90) -- (285.17, 93.40);

\path[draw=drawColor,line width= 0.4pt,line join=round,line cap=round] (288.50,100.23) -- (293.00,104.73);

\path[draw=drawColor,line width= 0.4pt,line join=round,line cap=round] (288.50,104.73) -- (293.00,100.23);

\path[draw=drawColor,line width= 0.4pt,line join=round,line cap=round] (291.88,102.99) -- (296.38,107.49);

\path[draw=drawColor,line width= 0.4pt,line join=round,line cap=round] (291.88,107.49) -- (296.38,102.99);

\path[draw=drawColor,line width= 0.4pt,line join=round,line cap=round] (307.45,116.38) -- (311.95,120.88);

\path[draw=drawColor,line width= 0.4pt,line join=round,line cap=round] (307.45,120.88) -- (311.95,116.38);

\path[draw=drawColor,line width= 0.4pt,line join=round,line cap=round] (285.27, 97.40) -- (289.77,101.90);

\path[draw=drawColor,line width= 0.4pt,line join=round,line cap=round] (285.27,101.90) -- (289.77, 97.40);

\path[draw=drawColor,line width= 0.4pt,line join=round,line cap=round] (296.94,107.84) -- (301.44,112.34);

\path[draw=drawColor,line width= 0.4pt,line join=round,line cap=round] (296.94,112.34) -- (301.44,107.84);

\path[draw=drawColor,line width= 0.4pt,line join=round,line cap=round] (309.20,117.80) -- (313.70,122.30);

\path[draw=drawColor,line width= 0.4pt,line join=round,line cap=round] (309.20,122.30) -- (313.70,117.80);

\path[draw=drawColor,line width= 0.4pt,line join=round,line cap=round] (282.08, 94.73) -- (286.58, 99.23);

\path[draw=drawColor,line width= 0.4pt,line join=round,line cap=round] (282.08, 99.23) -- (286.58, 94.73);

\path[draw=drawColor,line width= 0.4pt,line join=round,line cap=round] (279.98, 92.80) -- (284.48, 97.30);

\path[draw=drawColor,line width= 0.4pt,line join=round,line cap=round] (279.98, 97.30) -- (284.48, 92.80);

\path[draw=drawColor,line width= 0.4pt,line join=round,line cap=round] (300.06,111.30) -- (304.56,115.80);

\path[draw=drawColor,line width= 0.4pt,line join=round,line cap=round] (300.06,115.80) -- (304.56,111.30);

\path[draw=drawColor,line width= 0.4pt,line join=round,line cap=round] (283.21, 95.78) -- (287.71,100.28);

\path[draw=drawColor,line width= 0.4pt,line join=round,line cap=round] (283.21,100.28) -- (287.71, 95.78);

\path[draw=drawColor,line width= 0.4pt,line join=round,line cap=round] (297.61,108.39) -- (302.11,112.89);

\path[draw=drawColor,line width= 0.4pt,line join=round,line cap=round] (297.61,112.89) -- (302.11,108.39);

\path[draw=drawColor,line width= 0.4pt,line join=round,line cap=round] (288.78,100.43) -- (293.28,104.93);

\path[draw=drawColor,line width= 0.4pt,line join=round,line cap=round] (288.78,104.93) -- (293.28,100.43);

\path[draw=drawColor,line width= 0.4pt,line join=round,line cap=round] (287.73, 99.58) -- (292.23,104.08);

\path[draw=drawColor,line width= 0.4pt,line join=round,line cap=round] (287.73,104.08) -- (292.23, 99.58);

\path[draw=drawColor,line width= 0.4pt,line join=round,line cap=round] (295.69,106.65) -- (300.19,111.15);

\path[draw=drawColor,line width= 0.4pt,line join=round,line cap=round] (295.69,111.15) -- (300.19,106.65);

\path[draw=drawColor,line width= 0.4pt,line join=round,line cap=round] (296.75,107.66) -- (301.25,112.16);

\path[draw=drawColor,line width= 0.4pt,line join=round,line cap=round] (296.75,112.16) -- (301.25,107.66);

\path[draw=drawColor,line width= 0.4pt,line join=round,line cap=round] (299.54,110.81) -- (304.04,115.31);

\path[draw=drawColor,line width= 0.4pt,line join=round,line cap=round] (299.54,115.31) -- (304.04,110.81);

\path[draw=drawColor,line width= 0.4pt,line join=round,line cap=round] (284.06, 96.41) -- (288.56,100.91);

\path[draw=drawColor,line width= 0.4pt,line join=round,line cap=round] (284.06,100.91) -- (288.56, 96.41);

\path[draw=drawColor,line width= 0.4pt,line join=round,line cap=round] (293.30,104.71) -- (297.80,109.21);

\path[draw=drawColor,line width= 0.4pt,line join=round,line cap=round] (293.30,109.21) -- (297.80,104.71);

\path[draw=drawColor,line width= 0.4pt,line join=round,line cap=round] (288.92,100.52) -- (293.42,105.02);

\path[draw=drawColor,line width= 0.4pt,line join=round,line cap=round] (288.92,105.02) -- (293.42,100.52);

\path[draw=drawColor,line width= 0.4pt,line join=round,line cap=round] (286.91, 98.80) -- (291.41,103.30);

\path[draw=drawColor,line width= 0.4pt,line join=round,line cap=round] (286.91,103.30) -- (291.41, 98.80);

\path[draw=drawColor,line width= 0.4pt,line join=round,line cap=round] (289.52,100.97) -- (294.02,105.47);

\path[draw=drawColor,line width= 0.4pt,line join=round,line cap=round] (289.52,105.47) -- (294.02,100.97);

\path[draw=drawColor,line width= 0.4pt,line join=round,line cap=round] (284.59, 96.83) -- (289.09,101.33);

\path[draw=drawColor,line width= 0.4pt,line join=round,line cap=round] (284.59,101.33) -- (289.09, 96.83);

\path[draw=drawColor,line width= 0.4pt,line join=round,line cap=round] (296.71,107.65) -- (301.21,112.15);

\path[draw=drawColor,line width= 0.4pt,line join=round,line cap=round] (296.71,112.15) -- (301.21,107.65);

\path[draw=drawColor,line width= 0.4pt,line join=round,line cap=round] (300.75,111.86) -- (305.25,116.36);

\path[draw=drawColor,line width= 0.4pt,line join=round,line cap=round] (300.75,116.36) -- (305.25,111.86);

\path[draw=drawColor,line width= 0.4pt,line join=round,line cap=round] (284.29, 96.55) -- (288.79,101.05);

\path[draw=drawColor,line width= 0.4pt,line join=round,line cap=round] (284.29,101.05) -- (288.79, 96.55);

\path[draw=drawColor,line width= 0.4pt,line join=round,line cap=round] (291.95,103.11) -- (296.45,107.61);

\path[draw=drawColor,line width= 0.4pt,line join=round,line cap=round] (291.95,107.61) -- (296.45,103.11);

\path[draw=drawColor,line width= 0.4pt,line join=round,line cap=round] (292.10,103.25) -- (296.60,107.75);

\path[draw=drawColor,line width= 0.4pt,line join=round,line cap=round] (292.10,107.75) -- (296.60,103.25);

\path[draw=drawColor,line width= 0.4pt,line join=round,line cap=round] (286.84, 98.76) -- (291.34,103.26);

\path[draw=drawColor,line width= 0.4pt,line join=round,line cap=round] (286.84,103.26) -- (291.34, 98.76);

\path[draw=drawColor,line width= 0.4pt,line join=round,line cap=round] (288.21,100.01) -- (292.71,104.51);

\path[draw=drawColor,line width= 0.4pt,line join=round,line cap=round] (288.21,104.51) -- (292.71,100.01);

\path[draw=drawColor,line width= 0.4pt,line join=round,line cap=round] (280.45, 93.24) -- (284.95, 97.74);

\path[draw=drawColor,line width= 0.4pt,line join=round,line cap=round] (280.45, 97.74) -- (284.95, 93.24);

\path[draw=drawColor,line width= 0.4pt,line join=round,line cap=round] (295.20,106.18) -- (299.70,110.68);

\path[draw=drawColor,line width= 0.4pt,line join=round,line cap=round] (295.20,110.68) -- (299.70,106.18);

\path[draw=drawColor,line width= 0.4pt,line join=round,line cap=round] (300.11,111.35) -- (304.61,115.85);

\path[draw=drawColor,line width= 0.4pt,line join=round,line cap=round] (300.11,115.85) -- (304.61,111.35);

\path[draw=drawColor,line width= 0.4pt,line join=round,line cap=round] (283.34, 95.91) -- (287.84,100.41);

\path[draw=drawColor,line width= 0.4pt,line join=round,line cap=round] (283.34,100.41) -- (287.84, 95.91);

\path[draw=drawColor,line width= 0.4pt,line join=round,line cap=round] (282.58, 95.21) -- (287.08, 99.71);

\path[draw=drawColor,line width= 0.4pt,line join=round,line cap=round] (282.58, 99.71) -- (287.08, 95.21);

\path[draw=drawColor,line width= 0.4pt,line join=round,line cap=round] (290.33,101.64) -- (294.83,106.14);

\path[draw=drawColor,line width= 0.4pt,line join=round,line cap=round] (290.33,106.14) -- (294.83,101.64);

\path[draw=drawColor,line width= 0.4pt,line join=round,line cap=round] (273.83, 86.72) -- (278.33, 91.22);

\path[draw=drawColor,line width= 0.4pt,line join=round,line cap=round] (273.83, 91.22) -- (278.33, 86.72);

\path[draw=drawColor,line width= 0.4pt,line join=round,line cap=round] (283.24, 95.79) -- (287.74,100.29);

\path[draw=drawColor,line width= 0.4pt,line join=round,line cap=round] (283.24,100.29) -- (287.74, 95.79);

\path[draw=drawColor,line width= 0.4pt,line join=round,line cap=round] (291.43,102.65) -- (295.93,107.15);

\path[draw=drawColor,line width= 0.4pt,line join=round,line cap=round] (291.43,107.15) -- (295.93,102.65);

\path[draw=drawColor,line width= 0.4pt,line join=round,line cap=round] (294.61,105.91) -- (299.11,110.41);

\path[draw=drawColor,line width= 0.4pt,line join=round,line cap=round] (294.61,110.41) -- (299.11,105.91);

\path[draw=drawColor,line width= 0.4pt,line join=round,line cap=round] (290.16,101.50) -- (294.66,106.00);

\path[draw=drawColor,line width= 0.4pt,line join=round,line cap=round] (290.16,106.00) -- (294.66,101.50);

\path[draw=drawColor,line width= 0.4pt,line join=round,line cap=round] (282.68, 95.27) -- (287.18, 99.77);

\path[draw=drawColor,line width= 0.4pt,line join=round,line cap=round] (282.68, 99.77) -- (287.18, 95.27);

\path[draw=drawColor,line width= 0.4pt,line join=round,line cap=round] (294.09,105.46) -- (298.59,109.96);

\path[draw=drawColor,line width= 0.4pt,line join=round,line cap=round] (294.09,109.96) -- (298.59,105.46);

\path[draw=drawColor,line width= 0.4pt,line join=round,line cap=round] (290.45,101.80) -- (294.95,106.30);

\path[draw=drawColor,line width= 0.4pt,line join=round,line cap=round] (290.45,106.30) -- (294.95,101.80);

\path[draw=drawColor,line width= 0.4pt,line join=round,line cap=round] (294.04,105.42) -- (298.54,109.92);

\path[draw=drawColor,line width= 0.4pt,line join=round,line cap=round] (294.04,109.92) -- (298.54,105.42);

\path[draw=drawColor,line width= 0.4pt,line join=round,line cap=round] (303.34,114.07) -- (307.84,118.57);

\path[draw=drawColor,line width= 0.4pt,line join=round,line cap=round] (303.34,118.57) -- (307.84,114.07);

\path[draw=drawColor,line width= 0.4pt,line join=round,line cap=round] (286.82, 98.76) -- (291.32,103.26);

\path[draw=drawColor,line width= 0.4pt,line join=round,line cap=round] (286.82,103.26) -- (291.32, 98.76);

\path[draw=drawColor,line width= 0.4pt,line join=round,line cap=round] (286.74, 98.67) -- (291.24,103.17);

\path[draw=drawColor,line width= 0.4pt,line join=round,line cap=round] (286.74,103.17) -- (291.24, 98.67);

\path[draw=drawColor,line width= 0.4pt,line join=round,line cap=round] (283.16, 95.68) -- (287.66,100.18);

\path[draw=drawColor,line width= 0.4pt,line join=round,line cap=round] (283.16,100.18) -- (287.66, 95.68);

\path[draw=drawColor,line width= 0.4pt,line join=round,line cap=round] (284.83, 96.99) -- (289.33,101.49);

\path[draw=drawColor,line width= 0.4pt,line join=round,line cap=round] (284.83,101.49) -- (289.33, 96.99);

\path[draw=drawColor,line width= 0.4pt,line join=round,line cap=round] (294.89,106.00) -- (299.39,110.50);

\path[draw=drawColor,line width= 0.4pt,line join=round,line cap=round] (294.89,110.50) -- (299.39,106.00);

\path[draw=drawColor,line width= 0.4pt,line join=round,line cap=round] (291.24,102.46) -- (295.74,106.96);

\path[draw=drawColor,line width= 0.4pt,line join=round,line cap=round] (291.24,106.96) -- (295.74,102.46);

\path[draw=drawColor,line width= 0.4pt,line join=round,line cap=round] (290.35,101.64) -- (294.85,106.14);

\path[draw=drawColor,line width= 0.4pt,line join=round,line cap=round] (290.35,106.14) -- (294.85,101.64);

\path[draw=drawColor,line width= 0.4pt,line join=round,line cap=round] (295.63,106.61) -- (300.13,111.11);

\path[draw=drawColor,line width= 0.4pt,line join=round,line cap=round] (295.63,111.11) -- (300.13,106.61);

\path[draw=drawColor,line width= 0.4pt,line join=round,line cap=round] (291.16,102.37) -- (295.66,106.87);

\path[draw=drawColor,line width= 0.4pt,line join=round,line cap=round] (291.16,106.87) -- (295.66,102.37);

\path[draw=drawColor,line width= 0.4pt,line join=round,line cap=round] (287.10, 98.98) -- (291.60,103.48);

\path[draw=drawColor,line width= 0.4pt,line join=round,line cap=round] (287.10,103.48) -- (291.60, 98.98);

\path[draw=drawColor,line width= 0.4pt,line join=round,line cap=round] (277.70, 90.23) -- (282.20, 94.73);

\path[draw=drawColor,line width= 0.4pt,line join=round,line cap=round] (277.70, 94.73) -- (282.20, 90.23);

\path[draw=drawColor,line width= 0.4pt,line join=round,line cap=round] (290.42,101.80) -- (294.92,106.30);

\path[draw=drawColor,line width= 0.4pt,line join=round,line cap=round] (290.42,106.30) -- (294.92,101.80);

\path[draw=drawColor,line width= 0.4pt,line join=round,line cap=round] (277.25, 89.52) -- (281.75, 94.02);

\path[draw=drawColor,line width= 0.4pt,line join=round,line cap=round] (277.25, 94.02) -- (281.75, 89.52);

\path[draw=drawColor,line width= 0.4pt,line join=round,line cap=round] (285.24, 97.36) -- (289.74,101.86);

\path[draw=drawColor,line width= 0.4pt,line join=round,line cap=round] (285.24,101.86) -- (289.74, 97.36);

\path[draw=drawColor,line width= 0.4pt,line join=round,line cap=round] (281.57, 94.32) -- (286.07, 98.82);

\path[draw=drawColor,line width= 0.4pt,line join=round,line cap=round] (281.57, 98.82) -- (286.07, 94.32);

\path[draw=drawColor,line width= 0.4pt,line join=round,line cap=round] (292.98,104.38) -- (297.48,108.88);

\path[draw=drawColor,line width= 0.4pt,line join=round,line cap=round] (292.98,108.88) -- (297.48,104.38);

\path[draw=drawColor,line width= 0.4pt,line join=round,line cap=round] (278.94, 91.47) -- (283.44, 95.97);

\path[draw=drawColor,line width= 0.4pt,line join=round,line cap=round] (278.94, 95.97) -- (283.44, 91.47);

\path[draw=drawColor,line width= 0.4pt,line join=round,line cap=round] (296.65,107.60) -- (301.15,112.10);

\path[draw=drawColor,line width= 0.4pt,line join=round,line cap=round] (296.65,112.10) -- (301.15,107.60);

\path[draw=drawColor,line width= 0.4pt,line join=round,line cap=round] (274.69, 87.55) -- (279.19, 92.05);

\path[draw=drawColor,line width= 0.4pt,line join=round,line cap=round] (274.69, 92.05) -- (279.19, 87.55);

\path[draw=drawColor,line width= 0.4pt,line join=round,line cap=round] (284.15, 96.43) -- (288.65,100.93);

\path[draw=drawColor,line width= 0.4pt,line join=round,line cap=round] (284.15,100.93) -- (288.65, 96.43);

\path[draw=drawColor,line width= 0.4pt,line join=round,line cap=round] (293.28,104.71) -- (297.78,109.21);

\path[draw=drawColor,line width= 0.4pt,line join=round,line cap=round] (293.28,109.21) -- (297.78,104.71);

\path[draw=drawColor,line width= 0.4pt,line join=round,line cap=round] (301.83,112.77) -- (306.33,117.27);

\path[draw=drawColor,line width= 0.4pt,line join=round,line cap=round] (301.83,117.27) -- (306.33,112.77);

\path[draw=drawColor,line width= 0.4pt,line join=round,line cap=round] (296.49,107.41) -- (300.99,111.91);

\path[draw=drawColor,line width= 0.4pt,line join=round,line cap=round] (296.49,111.91) -- (300.99,107.41);

\path[draw=drawColor,line width= 0.4pt,line join=round,line cap=round] (286.77, 98.72) -- (291.27,103.22);

\path[draw=drawColor,line width= 0.4pt,line join=round,line cap=round] (286.77,103.22) -- (291.27, 98.72);

\path[draw=drawColor,line width= 0.4pt,line join=round,line cap=round] (285.82, 97.94) -- (290.32,102.44);

\path[draw=drawColor,line width= 0.4pt,line join=round,line cap=round] (285.82,102.44) -- (290.32, 97.94);

\path[draw=drawColor,line width= 0.4pt,line join=round,line cap=round] (284.00, 96.36) -- (288.50,100.86);

\path[draw=drawColor,line width= 0.4pt,line join=round,line cap=round] (284.00,100.86) -- (288.50, 96.36);

\path[draw=drawColor,line width= 0.4pt,line join=round,line cap=round] (281.09, 93.87) -- (285.59, 98.37);

\path[draw=drawColor,line width= 0.4pt,line join=round,line cap=round] (281.09, 98.37) -- (285.59, 93.87);

\path[draw=drawColor,line width= 0.4pt,line join=round,line cap=round] (289.33,100.81) -- (293.83,105.31);

\path[draw=drawColor,line width= 0.4pt,line join=round,line cap=round] (289.33,105.31) -- (293.83,100.81);

\path[draw=drawColor,line width= 0.4pt,line join=round,line cap=round] (268.03, 81.23) -- (272.53, 85.73);

\path[draw=drawColor,line width= 0.4pt,line join=round,line cap=round] (268.03, 85.73) -- (272.53, 81.23);

\path[draw=drawColor,line width= 0.4pt,line join=round,line cap=round] (264.73, 77.06) -- (269.23, 81.56);

\path[draw=drawColor,line width= 0.4pt,line join=round,line cap=round] (264.73, 81.56) -- (269.23, 77.06);

\path[draw=drawColor,line width= 0.4pt,line join=round,line cap=round] (281.76, 94.39) -- (286.26, 98.89);

\path[draw=drawColor,line width= 0.4pt,line join=round,line cap=round] (281.76, 98.89) -- (286.26, 94.39);

\path[draw=drawColor,line width= 0.4pt,line join=round,line cap=round] (281.61, 94.33) -- (286.11, 98.83);

\path[draw=drawColor,line width= 0.4pt,line join=round,line cap=round] (281.61, 98.83) -- (286.11, 94.33);

\path[draw=drawColor,line width= 0.4pt,line join=round,line cap=round] (313.53,122.01) -- (318.03,126.51);

\path[draw=drawColor,line width= 0.4pt,line join=round,line cap=round] (313.53,126.51) -- (318.03,122.01);

\path[draw=drawColor,line width= 0.4pt,line join=round,line cap=round] (285.16, 97.26) -- (289.66,101.76);

\path[draw=drawColor,line width= 0.4pt,line join=round,line cap=round] (285.16,101.76) -- (289.66, 97.26);

\path[draw=drawColor,line width= 0.4pt,line join=round,line cap=round] (298.40,109.59) -- (302.90,114.09);

\path[draw=drawColor,line width= 0.4pt,line join=round,line cap=round] (298.40,114.09) -- (302.90,109.59);

\path[draw=drawColor,line width= 0.4pt,line join=round,line cap=round] (292.15,103.37) -- (296.65,107.87);

\path[draw=drawColor,line width= 0.4pt,line join=round,line cap=round] (292.15,107.87) -- (296.65,103.37);

\path[draw=drawColor,line width= 0.4pt,line join=round,line cap=round] (277.31, 89.58) -- (281.81, 94.08);

\path[draw=drawColor,line width= 0.4pt,line join=round,line cap=round] (277.31, 94.08) -- (281.81, 89.58);

\path[draw=drawColor,line width= 0.4pt,line join=round,line cap=round] (282.93, 95.45) -- (287.43, 99.95);

\path[draw=drawColor,line width= 0.4pt,line join=round,line cap=round] (282.93, 99.95) -- (287.43, 95.45);

\path[draw=drawColor,line width= 0.4pt,line join=round,line cap=round] (292.27,103.42) -- (296.77,107.92);

\path[draw=drawColor,line width= 0.4pt,line join=round,line cap=round] (292.27,107.92) -- (296.77,103.42);

\path[draw=drawColor,line width= 0.4pt,line join=round,line cap=round] (294.73,105.95) -- (299.23,110.45);

\path[draw=drawColor,line width= 0.4pt,line join=round,line cap=round] (294.73,110.45) -- (299.23,105.95);

\path[draw=drawColor,line width= 0.4pt,line join=round,line cap=round] (287.78, 99.61) -- (292.28,104.11);

\path[draw=drawColor,line width= 0.4pt,line join=round,line cap=round] (287.78,104.11) -- (292.28, 99.61);

\path[draw=drawColor,line width= 0.4pt,line join=round,line cap=round] (316.74,124.83) -- (321.24,129.33);

\path[draw=drawColor,line width= 0.4pt,line join=round,line cap=round] (316.74,129.33) -- (321.24,124.83);

\path[draw=drawColor,line width= 0.4pt,line join=round,line cap=round] (270.95, 83.56) -- (275.45, 88.06);

\path[draw=drawColor,line width= 0.4pt,line join=round,line cap=round] (270.95, 88.06) -- (275.45, 83.56);

\path[draw=drawColor,line width= 0.4pt,line join=round,line cap=round] (297.48,108.36) -- (301.98,112.86);

\path[draw=drawColor,line width= 0.4pt,line join=round,line cap=round] (297.48,112.86) -- (301.98,108.36);

\path[draw=drawColor,line width= 0.4pt,line join=round,line cap=round] (278.77, 91.39) -- (283.27, 95.89);

\path[draw=drawColor,line width= 0.4pt,line join=round,line cap=round] (278.77, 95.89) -- (283.27, 91.39);

\path[draw=drawColor,line width= 0.4pt,line join=round,line cap=round] (296.18,107.14) -- (300.68,111.64);

\path[draw=drawColor,line width= 0.4pt,line join=round,line cap=round] (296.18,111.64) -- (300.68,107.14);

\path[draw=drawColor,line width= 0.4pt,line join=round,line cap=round] (290.11,101.46) -- (294.61,105.96);

\path[draw=drawColor,line width= 0.4pt,line join=round,line cap=round] (290.11,105.96) -- (294.61,101.46);

\path[draw=drawColor,line width= 0.4pt,line join=round,line cap=round] (292.12,103.37) -- (296.62,107.87);

\path[draw=drawColor,line width= 0.4pt,line join=round,line cap=round] (292.12,107.87) -- (296.62,103.37);

\path[draw=drawColor,line width= 0.4pt,line join=round,line cap=round] (276.83, 89.27) -- (281.33, 93.77);

\path[draw=drawColor,line width= 0.4pt,line join=round,line cap=round] (276.83, 93.77) -- (281.33, 89.27);

\path[draw=drawColor,line width= 0.4pt,line join=round,line cap=round] (279.90, 92.56) -- (284.40, 97.06);

\path[draw=drawColor,line width= 0.4pt,line join=round,line cap=round] (279.90, 97.06) -- (284.40, 92.56);

\path[draw=drawColor,line width= 0.4pt,line join=round,line cap=round] (276.40, 88.90) -- (280.90, 93.40);

\path[draw=drawColor,line width= 0.4pt,line join=round,line cap=round] (276.40, 93.40) -- (280.90, 88.90);

\path[draw=drawColor,line width= 0.4pt,line join=round,line cap=round] (268.87, 81.81) -- (273.37, 86.31);

\path[draw=drawColor,line width= 0.4pt,line join=round,line cap=round] (268.87, 86.31) -- (273.37, 81.81);

\path[draw=drawColor,line width= 0.4pt,line join=round,line cap=round] (281.49, 94.20) -- (285.99, 98.70);

\path[draw=drawColor,line width= 0.4pt,line join=round,line cap=round] (281.49, 98.70) -- (285.99, 94.20);

\path[draw=drawColor,line width= 0.4pt,line join=round,line cap=round] (308.73,117.23) -- (313.23,121.73);

\path[draw=drawColor,line width= 0.4pt,line join=round,line cap=round] (308.73,121.73) -- (313.23,117.23);

\path[draw=drawColor,line width= 0.4pt,line join=round,line cap=round] (307.72,116.66) -- (312.22,121.16);

\path[draw=drawColor,line width= 0.4pt,line join=round,line cap=round] (307.72,121.16) -- (312.22,116.66);

\path[draw=drawColor,line width= 0.4pt,line join=round,line cap=round] (292.34,103.54) -- (296.84,108.04);

\path[draw=drawColor,line width= 0.4pt,line join=round,line cap=round] (292.34,108.04) -- (296.84,103.54);

\path[draw=drawColor,line width= 0.4pt,line join=round,line cap=round] (293.88,105.24) -- (298.38,109.74);

\path[draw=drawColor,line width= 0.4pt,line join=round,line cap=round] (293.88,109.74) -- (298.38,105.24);

\path[draw=drawColor,line width= 0.4pt,line join=round,line cap=round] (275.04, 87.77) -- (279.54, 92.27);

\path[draw=drawColor,line width= 0.4pt,line join=round,line cap=round] (275.04, 92.27) -- (279.54, 87.77);

\path[draw=drawColor,line width= 0.4pt,line join=round,line cap=round] (300.27,111.41) -- (304.77,115.91);

\path[draw=drawColor,line width= 0.4pt,line join=round,line cap=round] (300.27,115.91) -- (304.77,111.41);

\path[draw=drawColor,line width= 0.4pt,line join=round,line cap=round] (273.93, 86.75) -- (278.43, 91.25);

\path[draw=drawColor,line width= 0.4pt,line join=round,line cap=round] (273.93, 91.25) -- (278.43, 86.75);

\path[draw=drawColor,line width= 0.4pt,line join=round,line cap=round] (276.63, 89.11) -- (281.13, 93.61);

\path[draw=drawColor,line width= 0.4pt,line join=round,line cap=round] (276.63, 93.61) -- (281.13, 89.11);

\path[draw=drawColor,line width= 0.4pt,line join=round,line cap=round] (300.59,111.66) -- (305.09,116.16);

\path[draw=drawColor,line width= 0.4pt,line join=round,line cap=round] (300.59,116.16) -- (305.09,111.66);

\path[draw=drawColor,line width= 0.4pt,line join=round,line cap=round] (276.76, 89.17) -- (281.26, 93.67);

\path[draw=drawColor,line width= 0.4pt,line join=round,line cap=round] (276.76, 93.67) -- (281.26, 89.17);

\path[draw=drawColor,line width= 0.4pt,line join=round,line cap=round] (283.75, 96.20) -- (288.25,100.70);

\path[draw=drawColor,line width= 0.4pt,line join=round,line cap=round] (283.75,100.70) -- (288.25, 96.20);

\path[draw=drawColor,line width= 0.4pt,line join=round,line cap=round] (291.75,102.93) -- (296.25,107.43);

\path[draw=drawColor,line width= 0.4pt,line join=round,line cap=round] (291.75,107.43) -- (296.25,102.93);

\path[draw=drawColor,line width= 0.4pt,line join=round,line cap=round] (290.96,102.30) -- (295.46,106.80);

\path[draw=drawColor,line width= 0.4pt,line join=round,line cap=round] (290.96,106.80) -- (295.46,102.30);

\path[draw=drawColor,line width= 0.4pt,line join=round,line cap=round] (308.42,117.07) -- (312.92,121.57);

\path[draw=drawColor,line width= 0.4pt,line join=round,line cap=round] (308.42,121.57) -- (312.92,117.07);

\path[draw=drawColor,line width= 0.4pt,line join=round,line cap=round] (299.07,110.25) -- (303.57,114.75);

\path[draw=drawColor,line width= 0.4pt,line join=round,line cap=round] (299.07,114.75) -- (303.57,110.25);

\path[draw=drawColor,line width= 0.4pt,line join=round,line cap=round] (274.32, 87.28) -- (278.82, 91.78);

\path[draw=drawColor,line width= 0.4pt,line join=round,line cap=round] (274.32, 91.78) -- (278.82, 87.28);

\path[draw=drawColor,line width= 0.4pt,line join=round,line cap=round] (268.25, 81.38) -- (272.75, 85.88);

\path[draw=drawColor,line width= 0.4pt,line join=round,line cap=round] (268.25, 85.88) -- (272.75, 81.38);

\path[draw=drawColor,line width= 0.4pt,line join=round,line cap=round] (292.53,103.76) -- (297.03,108.26);

\path[draw=drawColor,line width= 0.4pt,line join=round,line cap=round] (292.53,108.26) -- (297.03,103.76);

\path[draw=drawColor,line width= 0.4pt,line join=round,line cap=round] (284.48, 96.69) -- (288.98,101.19);

\path[draw=drawColor,line width= 0.4pt,line join=round,line cap=round] (284.48,101.19) -- (288.98, 96.69);

\path[draw=drawColor,line width= 0.4pt,line join=round,line cap=round] (274.42, 87.44) -- (278.92, 91.94);

\path[draw=drawColor,line width= 0.4pt,line join=round,line cap=round] (274.42, 91.94) -- (278.92, 87.44);

\path[draw=drawColor,line width= 0.4pt,line join=round,line cap=round] (250.87, 58.83) -- (255.37, 63.33);

\path[draw=drawColor,line width= 0.4pt,line join=round,line cap=round] (250.87, 63.33) -- (255.37, 58.83);

\path[draw=drawColor,line width= 0.4pt,line join=round,line cap=round] (280.11, 92.87) -- (284.61, 97.37);

\path[draw=drawColor,line width= 0.4pt,line join=round,line cap=round] (280.11, 97.37) -- (284.61, 92.87);

\path[draw=drawColor,line width= 0.4pt,line join=round,line cap=round] (276.47, 88.99) -- (280.97, 93.49);

\path[draw=drawColor,line width= 0.4pt,line join=round,line cap=round] (276.47, 93.49) -- (280.97, 88.99);

\path[draw=drawColor,line width= 0.4pt,line join=round,line cap=round] (295.44,106.34) -- (299.94,110.84);

\path[draw=drawColor,line width= 0.4pt,line join=round,line cap=round] (295.44,110.84) -- (299.94,106.34);

\path[draw=drawColor,line width= 0.4pt,line join=round,line cap=round] (300.48,111.61) -- (304.98,116.11);

\path[draw=drawColor,line width= 0.4pt,line join=round,line cap=round] (300.48,116.11) -- (304.98,111.61);

\path[draw=drawColor,line width= 0.4pt,line join=round,line cap=round] (278.62, 91.29) -- (283.12, 95.79);

\path[draw=drawColor,line width= 0.4pt,line join=round,line cap=round] (278.62, 95.79) -- (283.12, 91.29);

\path[draw=drawColor,line width= 0.4pt,line join=round,line cap=round] (303.72,114.18) -- (308.22,118.68);

\path[draw=drawColor,line width= 0.4pt,line join=round,line cap=round] (303.72,118.68) -- (308.22,114.18);

\path[draw=drawColor,line width= 0.4pt,line join=round,line cap=round] (328.65,134.71) -- (333.15,139.21);

\path[draw=drawColor,line width= 0.4pt,line join=round,line cap=round] (328.65,139.21) -- (333.15,134.71);

\path[draw=drawColor,line width= 0.4pt,line join=round,line cap=round] (307.32,116.33) -- (311.82,120.83);

\path[draw=drawColor,line width= 0.4pt,line join=round,line cap=round] (307.32,120.83) -- (311.82,116.33);

\path[draw=drawColor,line width= 0.4pt,line join=round,line cap=round] (284.77, 96.98) -- (289.27,101.48);

\path[draw=drawColor,line width= 0.4pt,line join=round,line cap=round] (284.77,101.48) -- (289.27, 96.98);

\path[draw=drawColor,line width= 0.4pt,line join=round,line cap=round] (304.12,114.44) -- (308.62,118.94);

\path[draw=drawColor,line width= 0.4pt,line join=round,line cap=round] (304.12,118.94) -- (308.62,114.44);

\path[draw=drawColor,line width= 0.4pt,line join=round,line cap=round] (283.62, 96.12) -- (288.12,100.62);

\path[draw=drawColor,line width= 0.4pt,line join=round,line cap=round] (283.62,100.62) -- (288.12, 96.12);

\path[draw=drawColor,line width= 0.4pt,line join=round,line cap=round] (280.93, 93.69) -- (285.43, 98.19);

\path[draw=drawColor,line width= 0.4pt,line join=round,line cap=round] (280.93, 98.19) -- (285.43, 93.69);

\path[draw=drawColor,line width= 0.4pt,line join=round,line cap=round] (285.91, 97.98) -- (290.41,102.48);

\path[draw=drawColor,line width= 0.4pt,line join=round,line cap=round] (285.91,102.48) -- (290.41, 97.98);

\path[draw=drawColor,line width= 0.4pt,line join=round,line cap=round] (295.01,106.04) -- (299.51,110.54);

\path[draw=drawColor,line width= 0.4pt,line join=round,line cap=round] (295.01,110.54) -- (299.51,106.04);

\path[draw=drawColor,line width= 0.4pt,line join=round,line cap=round] (290.91,102.15) -- (295.41,106.65);

\path[draw=drawColor,line width= 0.4pt,line join=round,line cap=round] (290.91,106.65) -- (295.41,102.15);

\path[draw=drawColor,line width= 0.4pt,line join=round,line cap=round] (277.88, 90.37) -- (282.38, 94.87);

\path[draw=drawColor,line width= 0.4pt,line join=round,line cap=round] (277.88, 94.87) -- (282.38, 90.37);

\path[draw=drawColor,line width= 0.4pt,line join=round,line cap=round] (285.08, 97.20) -- (289.58,101.70);

\path[draw=drawColor,line width= 0.4pt,line join=round,line cap=round] (285.08,101.70) -- (289.58, 97.20);

\path[draw=drawColor,line width= 0.4pt,line join=round,line cap=round] (299.16,110.39) -- (303.66,114.89);

\path[draw=drawColor,line width= 0.4pt,line join=round,line cap=round] (299.16,114.89) -- (303.66,110.39);

\path[draw=drawColor,line width= 0.4pt,line join=round,line cap=round] (282.98, 95.45) -- (287.48, 99.95);

\path[draw=drawColor,line width= 0.4pt,line join=round,line cap=round] (282.98, 99.95) -- (287.48, 95.45);

\path[draw=drawColor,line width= 0.4pt,line join=round,line cap=round] (301.33,112.28) -- (305.83,116.78);

\path[draw=drawColor,line width= 0.4pt,line join=round,line cap=round] (301.33,116.78) -- (305.83,112.28);

\path[draw=drawColor,line width= 0.4pt,line join=round,line cap=round] (297.80,109.00) -- (302.30,113.50);

\path[draw=drawColor,line width= 0.4pt,line join=round,line cap=round] (297.80,113.50) -- (302.30,109.00);

\path[draw=drawColor,line width= 0.4pt,line join=round,line cap=round] (287.05, 98.95) -- (291.55,103.45);

\path[draw=drawColor,line width= 0.4pt,line join=round,line cap=round] (287.05,103.45) -- (291.55, 98.95);

\path[draw=drawColor,line width= 0.4pt,line join=round,line cap=round] (278.01, 90.66) -- (282.51, 95.16);

\path[draw=drawColor,line width= 0.4pt,line join=round,line cap=round] (278.01, 95.16) -- (282.51, 90.66);

\path[draw=drawColor,line width= 0.4pt,line join=round,line cap=round] (306.04,115.52) -- (310.54,120.02);

\path[draw=drawColor,line width= 0.4pt,line join=round,line cap=round] (306.04,120.02) -- (310.54,115.52);

\path[draw=drawColor,line width= 0.4pt,line join=round,line cap=round] (289.56,101.04) -- (294.06,105.54);

\path[draw=drawColor,line width= 0.4pt,line join=round,line cap=round] (289.56,105.54) -- (294.06,101.04);

\path[draw=drawColor,line width= 0.4pt,line join=round,line cap=round] (299.03,110.22) -- (303.53,114.72);

\path[draw=drawColor,line width= 0.4pt,line join=round,line cap=round] (299.03,114.72) -- (303.53,110.22);

\path[draw=drawColor,line width= 0.4pt,line join=round,line cap=round] (302.15,113.00) -- (306.65,117.50);

\path[draw=drawColor,line width= 0.4pt,line join=round,line cap=round] (302.15,117.50) -- (306.65,113.00);

\path[draw=drawColor,line width= 0.4pt,line join=round,line cap=round] (293.15,104.63) -- (297.65,109.13);

\path[draw=drawColor,line width= 0.4pt,line join=round,line cap=round] (293.15,109.13) -- (297.65,104.63);

\path[draw=drawColor,line width= 0.4pt,line join=round,line cap=round] (287.42, 99.30) -- (291.92,103.80);

\path[draw=drawColor,line width= 0.4pt,line join=round,line cap=round] (287.42,103.80) -- (291.92, 99.30);

\path[draw=drawColor,line width= 0.4pt,line join=round,line cap=round] (286.12, 98.15) -- (290.62,102.65);

\path[draw=drawColor,line width= 0.4pt,line join=round,line cap=round] (286.12,102.65) -- (290.62, 98.15);

\path[draw=drawColor,line width= 0.4pt,line join=round,line cap=round] (285.52, 97.70) -- (290.02,102.20);

\path[draw=drawColor,line width= 0.4pt,line join=round,line cap=round] (285.52,102.20) -- (290.02, 97.70);

\path[draw=drawColor,line width= 0.4pt,line join=round,line cap=round] (271.10, 83.59) -- (275.60, 88.09);

\path[draw=drawColor,line width= 0.4pt,line join=round,line cap=round] (271.10, 88.09) -- (275.60, 83.59);

\path[draw=drawColor,line width= 0.4pt,line join=round,line cap=round] (294.67,105.93) -- (299.17,110.43);

\path[draw=drawColor,line width= 0.4pt,line join=round,line cap=round] (294.67,110.43) -- (299.17,105.93);

\path[draw=drawColor,line width= 0.4pt,line join=round,line cap=round] (283.97, 96.35) -- (288.47,100.85);

\path[draw=drawColor,line width= 0.4pt,line join=round,line cap=round] (283.97,100.85) -- (288.47, 96.35);

\path[draw=drawColor,line width= 0.4pt,line join=round,line cap=round] (258.87, 70.88) -- (263.37, 75.38);

\path[draw=drawColor,line width= 0.4pt,line join=round,line cap=round] (258.87, 75.38) -- (263.37, 70.88);

\path[draw=drawColor,line width= 0.4pt,line join=round,line cap=round] (298.98,110.19) -- (303.48,114.69);

\path[draw=drawColor,line width= 0.4pt,line join=round,line cap=round] (298.98,114.69) -- (303.48,110.19);

\path[draw=drawColor,line width= 0.4pt,line join=round,line cap=round] (297.26,108.13) -- (301.76,112.63);

\path[draw=drawColor,line width= 0.4pt,line join=round,line cap=round] (297.26,112.63) -- (301.76,108.13);

\path[draw=drawColor,line width= 0.4pt,line join=round,line cap=round] (279.15, 91.71) -- (283.65, 96.21);

\path[draw=drawColor,line width= 0.4pt,line join=round,line cap=round] (279.15, 96.21) -- (283.65, 91.71);

\path[draw=drawColor,line width= 0.4pt,line join=round,line cap=round] (298.56,109.73) -- (303.06,114.23);

\path[draw=drawColor,line width= 0.4pt,line join=round,line cap=round] (298.56,114.23) -- (303.06,109.73);

\path[draw=drawColor,line width= 0.4pt,line join=round,line cap=round] (286.33, 98.26) -- (290.83,102.76);

\path[draw=drawColor,line width= 0.4pt,line join=round,line cap=round] (286.33,102.76) -- (290.83, 98.26);

\path[draw=drawColor,line width= 0.4pt,line join=round,line cap=round] (300.21,111.40) -- (304.71,115.90);

\path[draw=drawColor,line width= 0.4pt,line join=round,line cap=round] (300.21,115.90) -- (304.71,111.40);

\path[draw=drawColor,line width= 0.4pt,line join=round,line cap=round] (288.07, 99.94) -- (292.57,104.44);

\path[draw=drawColor,line width= 0.4pt,line join=round,line cap=round] (288.07,104.44) -- (292.57, 99.94);

\path[draw=drawColor,line width= 0.4pt,line join=round,line cap=round] (293.12,104.53) -- (297.62,109.03);

\path[draw=drawColor,line width= 0.4pt,line join=round,line cap=round] (293.12,109.03) -- (297.62,104.53);

\path[draw=drawColor,line width= 0.4pt,line join=round,line cap=round] (290.56,101.87) -- (295.06,106.37);

\path[draw=drawColor,line width= 0.4pt,line join=round,line cap=round] (290.56,106.37) -- (295.06,101.87);

\path[draw=drawColor,line width= 0.4pt,line join=round,line cap=round] (286.66, 98.58) -- (291.16,103.08);

\path[draw=drawColor,line width= 0.4pt,line join=round,line cap=round] (286.66,103.08) -- (291.16, 98.58);

\path[draw=drawColor,line width= 0.4pt,line join=round,line cap=round] (303.19,114.01) -- (307.69,118.51);

\path[draw=drawColor,line width= 0.4pt,line join=round,line cap=round] (303.19,118.51) -- (307.69,114.01);

\path[draw=drawColor,line width= 0.4pt,line join=round,line cap=round] (278.52, 91.18) -- (283.02, 95.68);

\path[draw=drawColor,line width= 0.4pt,line join=round,line cap=round] (278.52, 95.68) -- (283.02, 91.18);

\path[draw=drawColor,line width= 0.4pt,line join=round,line cap=round] (283.72, 96.20) -- (288.22,100.70);

\path[draw=drawColor,line width= 0.4pt,line join=round,line cap=round] (283.72,100.70) -- (288.22, 96.20);

\path[draw=drawColor,line width= 0.4pt,line join=round,line cap=round] (275.48, 88.15) -- (279.98, 92.65);

\path[draw=drawColor,line width= 0.4pt,line join=round,line cap=round] (275.48, 92.65) -- (279.98, 88.15);

\path[draw=drawColor,line width= 0.4pt,line join=round,line cap=round] (292.79,104.22) -- (297.29,108.72);

\path[draw=drawColor,line width= 0.4pt,line join=round,line cap=round] (292.79,108.72) -- (297.29,104.22);

\path[draw=drawColor,line width= 0.4pt,line join=round,line cap=round] (272.81, 85.77) -- (277.31, 90.27);

\path[draw=drawColor,line width= 0.4pt,line join=round,line cap=round] (272.81, 90.27) -- (277.31, 85.77);

\path[draw=drawColor,line width= 0.4pt,line join=round,line cap=round] (270.80, 83.44) -- (275.30, 87.94);

\path[draw=drawColor,line width= 0.4pt,line join=round,line cap=round] (270.80, 87.94) -- (275.30, 83.44);

\path[draw=drawColor,line width= 0.4pt,line join=round,line cap=round] (282.30, 94.84) -- (286.80, 99.34);

\path[draw=drawColor,line width= 0.4pt,line join=round,line cap=round] (282.30, 99.34) -- (286.80, 94.84);

\path[draw=drawColor,line width= 0.4pt,line join=round,line cap=round] (294.37,105.68) -- (298.87,110.18);

\path[draw=drawColor,line width= 0.4pt,line join=round,line cap=round] (294.37,110.18) -- (298.87,105.68);

\path[draw=drawColor,line width= 0.4pt,line join=round,line cap=round] (279.05, 91.64) -- (283.55, 96.14);

\path[draw=drawColor,line width= 0.4pt,line join=round,line cap=round] (279.05, 96.14) -- (283.55, 91.64);

\path[draw=drawColor,line width= 0.4pt,line join=round,line cap=round] (296.99,107.89) -- (301.49,112.39);

\path[draw=drawColor,line width= 0.4pt,line join=round,line cap=round] (296.99,112.39) -- (301.49,107.89);

\path[draw=drawColor,line width= 0.4pt,line join=round,line cap=round] (296.28,107.23) -- (300.78,111.73);

\path[draw=drawColor,line width= 0.4pt,line join=round,line cap=round] (296.28,111.73) -- (300.78,107.23);

\path[draw=drawColor,line width= 0.4pt,line join=round,line cap=round] (299.25,110.45) -- (303.75,114.95);

\path[draw=drawColor,line width= 0.4pt,line join=round,line cap=round] (299.25,114.95) -- (303.75,110.45);

\path[draw=drawColor,line width= 0.4pt,line join=round,line cap=round] (260.20, 73.30) -- (264.70, 77.80);

\path[draw=drawColor,line width= 0.4pt,line join=round,line cap=round] (260.20, 77.80) -- (264.70, 73.30);

\path[draw=drawColor,line width= 0.4pt,line join=round,line cap=round] (292.03,103.15) -- (296.53,107.65);

\path[draw=drawColor,line width= 0.4pt,line join=round,line cap=round] (292.03,107.65) -- (296.53,103.15);

\path[draw=drawColor,line width= 0.4pt,line join=round,line cap=round] (291.11,102.35) -- (295.61,106.85);

\path[draw=drawColor,line width= 0.4pt,line join=round,line cap=round] (291.11,106.85) -- (295.61,102.35);

\path[draw=drawColor,line width= 0.4pt,line join=round,line cap=round] (323.20,131.99) -- (327.70,136.49);

\path[draw=drawColor,line width= 0.4pt,line join=round,line cap=round] (323.20,136.49) -- (327.70,131.99);

\path[draw=drawColor,line width= 0.4pt,line join=round,line cap=round] (305.02,115.00) -- (309.52,119.50);

\path[draw=drawColor,line width= 0.4pt,line join=round,line cap=round] (305.02,119.50) -- (309.52,115.00);

\path[draw=drawColor,line width= 0.4pt,line join=round,line cap=round] (283.84, 96.27) -- (288.34,100.77);

\path[draw=drawColor,line width= 0.4pt,line join=round,line cap=round] (283.84,100.77) -- (288.34, 96.27);

\path[draw=drawColor,line width= 0.4pt,line join=round,line cap=round] (280.76, 93.48) -- (285.26, 97.98);

\path[draw=drawColor,line width= 0.4pt,line join=round,line cap=round] (280.76, 97.98) -- (285.26, 93.48);

\path[draw=drawColor,line width= 0.4pt,line join=round,line cap=round] (295.84,106.82) -- (300.34,111.32);

\path[draw=drawColor,line width= 0.4pt,line join=round,line cap=round] (295.84,111.32) -- (300.34,106.82);

\path[draw=drawColor,line width= 0.4pt,line join=round,line cap=round] (283.89, 96.31) -- (288.39,100.81);

\path[draw=drawColor,line width= 0.4pt,line join=round,line cap=round] (283.89,100.81) -- (288.39, 96.31);

\path[draw=drawColor,line width= 0.4pt,line join=round,line cap=round] (281.33, 94.08) -- (285.83, 98.58);

\path[draw=drawColor,line width= 0.4pt,line join=round,line cap=round] (281.33, 98.58) -- (285.83, 94.08);

\path[draw=drawColor,line width= 0.4pt,line join=round,line cap=round] (289.99,101.33) -- (294.49,105.83);

\path[draw=drawColor,line width= 0.4pt,line join=round,line cap=round] (289.99,105.83) -- (294.49,101.33);

\path[draw=drawColor,line width= 0.4pt,line join=round,line cap=round] (281.45, 94.20) -- (285.95, 98.70);

\path[draw=drawColor,line width= 0.4pt,line join=round,line cap=round] (281.45, 98.70) -- (285.95, 94.20);

\path[draw=drawColor,line width= 0.4pt,line join=round,line cap=round] (278.04, 90.69) -- (282.54, 95.19);

\path[draw=drawColor,line width= 0.4pt,line join=round,line cap=round] (278.04, 95.19) -- (282.54, 90.69);

\path[draw=drawColor,line width= 0.4pt,line join=round,line cap=round] (265.70, 78.19) -- (270.20, 82.69);

\path[draw=drawColor,line width= 0.4pt,line join=round,line cap=round] (265.70, 82.69) -- (270.20, 78.19);

\path[draw=drawColor,line width= 0.4pt,line join=round,line cap=round] (281.68, 94.38) -- (286.18, 98.88);

\path[draw=drawColor,line width= 0.4pt,line join=round,line cap=round] (281.68, 98.88) -- (286.18, 94.38);

\path[draw=drawColor,line width= 0.4pt,line join=round,line cap=round] (269.07, 82.01) -- (273.57, 86.51);

\path[draw=drawColor,line width= 0.4pt,line join=round,line cap=round] (269.07, 86.51) -- (273.57, 82.01);

\path[draw=drawColor,line width= 0.4pt,line join=round,line cap=round] (291.18,102.42) -- (295.68,106.92);

\path[draw=drawColor,line width= 0.4pt,line join=round,line cap=round] (291.18,106.92) -- (295.68,102.42);

\path[draw=drawColor,line width= 0.4pt,line join=round,line cap=round] (280.80, 93.48) -- (285.30, 97.98);

\path[draw=drawColor,line width= 0.4pt,line join=round,line cap=round] (280.80, 97.98) -- (285.30, 93.48);

\path[draw=drawColor,line width= 0.4pt,line join=round,line cap=round] (306.26,115.66) -- (310.76,120.16);

\path[draw=drawColor,line width= 0.4pt,line join=round,line cap=round] (306.26,120.16) -- (310.76,115.66);

\path[draw=drawColor,line width= 0.4pt,line join=round,line cap=round] (262.79, 75.41) -- (267.29, 79.91);

\path[draw=drawColor,line width= 0.4pt,line join=round,line cap=round] (262.79, 79.91) -- (267.29, 75.41);

\path[draw=drawColor,line width= 0.4pt,line join=round,line cap=round] (279.42, 91.83) -- (283.92, 96.33);

\path[draw=drawColor,line width= 0.4pt,line join=round,line cap=round] (279.42, 96.33) -- (283.92, 91.83);

\path[draw=drawColor,line width= 0.4pt,line join=round,line cap=round] (286.70, 98.61) -- (291.20,103.11);

\path[draw=drawColor,line width= 0.4pt,line join=round,line cap=round] (286.70,103.11) -- (291.20, 98.61);

\path[draw=drawColor,line width= 0.4pt,line join=round,line cap=round] (299.50,110.73) -- (304.00,115.23);

\path[draw=drawColor,line width= 0.4pt,line join=round,line cap=round] (299.50,115.23) -- (304.00,110.73);

\path[draw=drawColor,line width= 0.4pt,line join=round,line cap=round] (304.92,114.99) -- (309.42,119.49);

\path[draw=drawColor,line width= 0.4pt,line join=round,line cap=round] (304.92,119.49) -- (309.42,114.99);

\path[draw=drawColor,line width= 0.4pt,line join=round,line cap=round] (258.11, 67.35) -- (262.61, 71.85);

\path[draw=drawColor,line width= 0.4pt,line join=round,line cap=round] (258.11, 71.85) -- (262.61, 67.35);

\path[draw=drawColor,line width= 0.4pt,line join=round,line cap=round] (304.83,114.97) -- (309.33,119.47);

\path[draw=drawColor,line width= 0.4pt,line join=round,line cap=round] (304.83,119.47) -- (309.33,114.97);

\path[draw=drawColor,line width= 0.4pt,line join=round,line cap=round] (279.26, 91.72) -- (283.76, 96.22);

\path[draw=drawColor,line width= 0.4pt,line join=round,line cap=round] (279.26, 96.22) -- (283.76, 91.72);

\path[draw=drawColor,line width= 0.4pt,line join=round,line cap=round] (322.26,131.20) -- (326.76,135.70);

\path[draw=drawColor,line width= 0.4pt,line join=round,line cap=round] (322.26,135.70) -- (326.76,131.20);

\path[draw=drawColor,line width= 0.4pt,line join=round,line cap=round] (302.34,113.32) -- (306.84,117.82);

\path[draw=drawColor,line width= 0.4pt,line join=round,line cap=round] (302.34,117.82) -- (306.84,113.32);

\path[draw=drawColor,line width= 0.4pt,line join=round,line cap=round] (294.56,105.85) -- (299.06,110.35);

\path[draw=drawColor,line width= 0.4pt,line join=round,line cap=round] (294.56,110.35) -- (299.06,105.85);

\path[draw=drawColor,line width= 0.4pt,line join=round,line cap=round] (267.08, 79.58) -- (271.58, 84.08);

\path[draw=drawColor,line width= 0.4pt,line join=round,line cap=round] (267.08, 84.08) -- (271.58, 79.58);

\path[draw=drawColor,line width= 0.4pt,line join=round,line cap=round] (319.96,128.41) -- (324.46,132.91);

\path[draw=drawColor,line width= 0.4pt,line join=round,line cap=round] (319.96,132.91) -- (324.46,128.41);

\path[draw=drawColor,line width= 0.4pt,line join=round,line cap=round] (319.32,128.20) -- (323.82,132.70);

\path[draw=drawColor,line width= 0.4pt,line join=round,line cap=round] (319.32,132.70) -- (323.82,128.20);

\path[draw=drawColor,line width= 0.4pt,line join=round,line cap=round] (291.31,102.63) -- (295.81,107.13);

\path[draw=drawColor,line width= 0.4pt,line join=round,line cap=round] (291.31,107.13) -- (295.81,102.63);

\path[draw=drawColor,line width= 0.4pt,line join=round,line cap=round] (272.69, 85.67) -- (277.19, 90.17);

\path[draw=drawColor,line width= 0.4pt,line join=round,line cap=round] (272.69, 90.17) -- (277.19, 85.67);

\path[draw=drawColor,line width= 0.4pt,line join=round,line cap=round] (281.88, 94.45) -- (286.38, 98.95);

\path[draw=drawColor,line width= 0.4pt,line join=round,line cap=round] (281.88, 98.95) -- (286.38, 94.45);

\path[draw=drawColor,line width= 0.4pt,line join=round,line cap=round] (269.63, 82.52) -- (274.13, 87.02);

\path[draw=drawColor,line width= 0.4pt,line join=round,line cap=round] (269.63, 87.02) -- (274.13, 82.52);

\path[draw=drawColor,line width= 0.4pt,line join=round,line cap=round] (280.54, 93.34) -- (285.04, 97.84);

\path[draw=drawColor,line width= 0.4pt,line join=round,line cap=round] (280.54, 97.84) -- (285.04, 93.34);

\path[draw=drawColor,line width= 0.4pt,line join=round,line cap=round] (310.86,119.27) -- (315.36,123.77);

\path[draw=drawColor,line width= 0.4pt,line join=round,line cap=round] (310.86,123.77) -- (315.36,119.27);

\path[draw=drawColor,line width= 0.4pt,line join=round,line cap=round] (313.25,121.98) -- (317.75,126.48);

\path[draw=drawColor,line width= 0.4pt,line join=round,line cap=round] (313.25,126.48) -- (317.75,121.98);

\path[draw=drawColor,line width= 0.4pt,line join=round,line cap=round] (312.45,121.23) -- (316.95,125.73);

\path[draw=drawColor,line width= 0.4pt,line join=round,line cap=round] (312.45,125.73) -- (316.95,121.23);

\path[draw=drawColor,line width= 0.4pt,line join=round,line cap=round] (277.51, 89.90) -- (282.01, 94.40);

\path[draw=drawColor,line width= 0.4pt,line join=round,line cap=round] (277.51, 94.40) -- (282.01, 89.90);

\path[draw=drawColor,line width= 0.4pt,line join=round,line cap=round] (281.37, 94.10) -- (285.87, 98.60);

\path[draw=drawColor,line width= 0.4pt,line join=round,line cap=round] (281.37, 98.60) -- (285.87, 94.10);

\path[draw=drawColor,line width= 0.4pt,line join=round,line cap=round] (311.06,119.62) -- (315.56,124.12);

\path[draw=drawColor,line width= 0.4pt,line join=round,line cap=round] (311.06,124.12) -- (315.56,119.62);

\path[draw=drawColor,line width= 0.4pt,line join=round,line cap=round] (273.38, 86.05) -- (277.88, 90.55);

\path[draw=drawColor,line width= 0.4pt,line join=round,line cap=round] (273.38, 90.55) -- (277.88, 86.05);

\path[draw=drawColor,line width= 0.4pt,line join=round,line cap=round] (290.82,102.04) -- (295.32,106.54);

\path[draw=drawColor,line width= 0.4pt,line join=round,line cap=round] (290.82,106.54) -- (295.32,102.04);

\path[draw=drawColor,line width= 0.4pt,line join=round,line cap=round] (295.94,106.86) -- (300.44,111.36);

\path[draw=drawColor,line width= 0.4pt,line join=round,line cap=round] (295.94,111.36) -- (300.44,106.86);

\path[draw=drawColor,line width= 0.4pt,line join=round,line cap=round] (293.53,104.96) -- (298.03,109.46);

\path[draw=drawColor,line width= 0.4pt,line join=round,line cap=round] (293.53,109.46) -- (298.03,104.96);

\path[draw=drawColor,line width= 0.4pt,line join=round,line cap=round] (297.88,109.09) -- (302.38,113.59);

\path[draw=drawColor,line width= 0.4pt,line join=round,line cap=round] (297.88,113.59) -- (302.38,109.09);

\path[draw=drawColor,line width= 0.4pt,line join=round,line cap=round] (303.88,114.30) -- (308.38,118.80);

\path[draw=drawColor,line width= 0.4pt,line join=round,line cap=round] (303.88,118.80) -- (308.38,114.30);

\path[draw=drawColor,line width= 0.4pt,line join=round,line cap=round] (292.58,103.78) -- (297.08,108.28);

\path[draw=drawColor,line width= 0.4pt,line join=round,line cap=round] (292.58,108.28) -- (297.08,103.78);

\path[draw=drawColor,line width= 0.4pt,line join=round,line cap=round] (314.46,122.95) -- (318.96,127.45);

\path[draw=drawColor,line width= 0.4pt,line join=round,line cap=round] (314.46,127.45) -- (318.96,122.95);

\path[draw=drawColor,line width= 0.4pt,line join=round,line cap=round] (299.58,110.88) -- (304.08,115.38);

\path[draw=drawColor,line width= 0.4pt,line join=round,line cap=round] (299.58,115.38) -- (304.08,110.88);

\path[draw=drawColor,line width= 0.4pt,line join=round,line cap=round] (309.72,118.10) -- (314.22,122.60);

\path[draw=drawColor,line width= 0.4pt,line join=round,line cap=round] (309.72,122.60) -- (314.22,118.10);

\path[draw=drawColor,line width= 0.4pt,line join=round,line cap=round] (261.84, 73.97) -- (266.34, 78.47);

\path[draw=drawColor,line width= 0.4pt,line join=round,line cap=round] (261.84, 78.47) -- (266.34, 73.97);

\path[draw=drawColor,line width= 0.4pt,line join=round,line cap=round] (290.98,102.32) -- (295.48,106.82);

\path[draw=drawColor,line width= 0.4pt,line join=round,line cap=round] (290.98,106.82) -- (295.48,102.32);

\path[draw=drawColor,line width= 0.4pt,line join=round,line cap=round] (284.03, 96.40) -- (288.53,100.90);

\path[draw=drawColor,line width= 0.4pt,line join=round,line cap=round] (284.03,100.90) -- (288.53, 96.40);

\path[draw=drawColor,line width= 0.4pt,line join=round,line cap=round] (303.42,114.08) -- (307.92,118.58);

\path[draw=drawColor,line width= 0.4pt,line join=round,line cap=round] (303.42,118.58) -- (307.92,114.08);

\path[draw=drawColor,line width= 0.4pt,line join=round,line cap=round] (288.75,100.42) -- (293.25,104.92);

\path[draw=drawColor,line width= 0.4pt,line join=round,line cap=round] (288.75,104.92) -- (293.25,100.42);

\path[draw=drawColor,line width= 0.4pt,line join=round,line cap=round] (286.53, 98.44) -- (291.03,102.94);

\path[draw=drawColor,line width= 0.4pt,line join=round,line cap=round] (286.53,102.94) -- (291.03, 98.44);

\path[draw=drawColor,line width= 0.4pt,line join=round,line cap=round] (309.37,117.93) -- (313.87,122.43);

\path[draw=drawColor,line width= 0.4pt,line join=round,line cap=round] (309.37,122.43) -- (313.87,117.93);

\path[draw=drawColor,line width= 0.4pt,line join=round,line cap=round] (304.37,114.59) -- (308.87,119.09);

\path[draw=drawColor,line width= 0.4pt,line join=round,line cap=round] (304.37,119.09) -- (308.87,114.59);

\path[draw=drawColor,line width= 0.4pt,line join=round,line cap=round] (285.21, 97.34) -- (289.71,101.84);

\path[draw=drawColor,line width= 0.4pt,line join=round,line cap=round] (285.21,101.84) -- (289.71, 97.34);

\path[draw=drawColor,line width= 0.4pt,line join=round,line cap=round] (293.25,104.66) -- (297.75,109.16);

\path[draw=drawColor,line width= 0.4pt,line join=round,line cap=round] (293.25,109.16) -- (297.75,104.66);

\path[draw=drawColor,line width= 0.4pt,line join=round,line cap=round] (286.07, 98.12) -- (290.57,102.62);

\path[draw=drawColor,line width= 0.4pt,line join=round,line cap=round] (286.07,102.62) -- (290.57, 98.12);

\path[draw=drawColor,line width= 0.4pt,line join=round,line cap=round] (271.94, 84.70) -- (276.44, 89.20);

\path[draw=drawColor,line width= 0.4pt,line join=round,line cap=round] (271.94, 89.20) -- (276.44, 84.70);

\path[draw=drawColor,line width= 0.4pt,line join=round,line cap=round] (305.30,115.14) -- (309.80,119.64);

\path[draw=drawColor,line width= 0.4pt,line join=round,line cap=round] (305.30,119.64) -- (309.80,115.14);

\path[draw=drawColor,line width= 0.4pt,line join=round,line cap=round] (289.89,101.30) -- (294.39,105.80);

\path[draw=drawColor,line width= 0.4pt,line join=round,line cap=round] (289.89,105.80) -- (294.39,101.30);

\path[draw=drawColor,line width= 0.4pt,line join=round,line cap=round] (296.05,106.99) -- (300.55,111.49);

\path[draw=drawColor,line width= 0.4pt,line join=round,line cap=round] (296.05,111.49) -- (300.55,106.99);

\path[draw=drawColor,line width= 0.4pt,line join=round,line cap=round] (293.94,105.32) -- (298.44,109.82);

\path[draw=drawColor,line width= 0.4pt,line join=round,line cap=round] (293.94,109.82) -- (298.44,105.32);

\path[draw=drawColor,line width= 0.4pt,line join=round,line cap=round] (289.87,101.30) -- (294.37,105.80);

\path[draw=drawColor,line width= 0.4pt,line join=round,line cap=round] (289.87,105.80) -- (294.37,101.30);

\path[draw=drawColor,line width= 0.4pt,line join=round,line cap=round] (292.05,103.17) -- (296.55,107.67);

\path[draw=drawColor,line width= 0.4pt,line join=round,line cap=round] (292.05,107.67) -- (296.55,103.17);

\path[draw=drawColor,line width= 0.4pt,line join=round,line cap=round] (279.57, 91.89) -- (284.07, 96.39);

\path[draw=drawColor,line width= 0.4pt,line join=round,line cap=round] (279.57, 96.39) -- (284.07, 91.89);

\path[draw=drawColor,line width= 0.4pt,line join=round,line cap=round] (284.74, 96.90) -- (289.24,101.40);

\path[draw=drawColor,line width= 0.4pt,line join=round,line cap=round] (284.74,101.40) -- (289.24, 96.90);

\path[draw=drawColor,line width= 0.4pt,line join=round,line cap=round] (289.40,100.87) -- (293.90,105.37);

\path[draw=drawColor,line width= 0.4pt,line join=round,line cap=round] (289.40,105.37) -- (293.90,100.87);

\path[draw=drawColor,line width= 0.4pt,line join=round,line cap=round] (279.47, 91.88) -- (283.97, 96.38);

\path[draw=drawColor,line width= 0.4pt,line join=round,line cap=round] (279.47, 96.38) -- (283.97, 91.88);

\path[draw=drawColor,line width= 0.4pt,line join=round,line cap=round] (291.33,102.63) -- (295.83,107.13);

\path[draw=drawColor,line width= 0.4pt,line join=round,line cap=round] (291.33,107.13) -- (295.83,102.63);

\path[draw=drawColor,line width= 0.4pt,line join=round,line cap=round] (292.31,103.50) -- (296.81,108.00);

\path[draw=drawColor,line width= 0.4pt,line join=round,line cap=round] (292.31,108.00) -- (296.81,103.50);

\path[draw=drawColor,line width= 0.4pt,line join=round,line cap=round] (302.76,113.77) -- (307.26,118.27);

\path[draw=drawColor,line width= 0.4pt,line join=round,line cap=round] (302.76,118.27) -- (307.26,113.77);

\path[draw=drawColor,line width= 0.4pt,line join=round,line cap=round] (317.69,126.51) -- (322.19,131.01);

\path[draw=drawColor,line width= 0.4pt,line join=round,line cap=round] (317.69,131.01) -- (322.19,126.51);

\path[draw=drawColor,line width= 0.4pt,line join=round,line cap=round] (288.66,100.39) -- (293.16,104.89);

\path[draw=drawColor,line width= 0.4pt,line join=round,line cap=round] (288.66,104.89) -- (293.16,100.39);

\path[draw=drawColor,line width= 0.4pt,line join=round,line cap=round] (290.80,102.03) -- (295.30,106.53);

\path[draw=drawColor,line width= 0.4pt,line join=round,line cap=round] (290.80,106.53) -- (295.30,102.03);

\path[draw=drawColor,line width= 0.4pt,line join=round,line cap=round] (290.38,101.70) -- (294.88,106.20);

\path[draw=drawColor,line width= 0.4pt,line join=round,line cap=round] (290.38,106.20) -- (294.88,101.70);

\path[draw=drawColor,line width= 0.4pt,line join=round,line cap=round] (299.12,110.36) -- (303.62,114.86);

\path[draw=drawColor,line width= 0.4pt,line join=round,line cap=round] (299.12,114.86) -- (303.62,110.36);

\path[draw=drawColor,line width= 0.4pt,line join=round,line cap=round] (283.18, 95.75) -- (287.68,100.25);

\path[draw=drawColor,line width= 0.4pt,line join=round,line cap=round] (283.18,100.25) -- (287.68, 95.75);

\path[draw=drawColor,line width= 0.4pt,line join=round,line cap=round] (279.71, 92.15) -- (284.21, 96.65);

\path[draw=drawColor,line width= 0.4pt,line join=round,line cap=round] (279.71, 96.65) -- (284.21, 92.15);

\path[draw=drawColor,line width= 0.4pt,line join=round,line cap=round] (286.25, 98.22) -- (290.75,102.72);

\path[draw=drawColor,line width= 0.4pt,line join=round,line cap=round] (286.25,102.72) -- (290.75, 98.22);

\path[draw=drawColor,line width= 0.4pt,line join=round,line cap=round] (275.56, 88.18) -- (280.06, 92.68);

\path[draw=drawColor,line width= 0.4pt,line join=round,line cap=round] (275.56, 92.68) -- (280.06, 88.18);

\path[draw=drawColor,line width= 0.4pt,line join=round,line cap=round] (297.44,108.31) -- (301.94,112.81);

\path[draw=drawColor,line width= 0.4pt,line join=round,line cap=round] (297.44,112.81) -- (301.94,108.31);

\path[draw=drawColor,line width= 0.4pt,line join=round,line cap=round] (298.44,109.64) -- (302.94,114.14);

\path[draw=drawColor,line width= 0.4pt,line join=round,line cap=round] (298.44,114.14) -- (302.94,109.64);

\path[draw=drawColor,line width= 0.4pt,line join=round,line cap=round] (280.50, 93.29) -- (285.00, 97.79);

\path[draw=drawColor,line width= 0.4pt,line join=round,line cap=round] (280.50, 97.79) -- (285.00, 93.29);

\path[draw=drawColor,line width= 0.4pt,line join=round,line cap=round] (289.27,100.79) -- (293.77,105.29);

\path[draw=drawColor,line width= 0.4pt,line join=round,line cap=round] (289.27,105.29) -- (293.77,100.79);

\path[draw=drawColor,line width= 0.4pt,line join=round,line cap=round] (290.89,102.05) -- (295.39,106.55);

\path[draw=drawColor,line width= 0.4pt,line join=round,line cap=round] (290.89,106.55) -- (295.39,102.05);

\path[draw=drawColor,line width= 0.4pt,line join=round,line cap=round] (275.80, 88.52) -- (280.30, 93.02);

\path[draw=drawColor,line width= 0.4pt,line join=round,line cap=round] (275.80, 93.02) -- (280.30, 88.52);

\path[draw=drawColor,line width= 0.4pt,line join=round,line cap=round] (269.81, 82.68) -- (274.31, 87.18);

\path[draw=drawColor,line width= 0.4pt,line join=round,line cap=round] (269.81, 87.18) -- (274.31, 82.68);

\path[draw=drawColor,line width= 0.4pt,line join=round,line cap=round] (275.64, 88.18) -- (280.14, 92.68);

\path[draw=drawColor,line width= 0.4pt,line join=round,line cap=round] (275.64, 92.68) -- (280.14, 88.18);

\path[draw=drawColor,line width= 0.4pt,line join=round,line cap=round] (295.14,106.13) -- (299.64,110.63);

\path[draw=drawColor,line width= 0.4pt,line join=round,line cap=round] (295.14,110.63) -- (299.64,106.13);

\path[draw=drawColor,line width= 0.4pt,line join=round,line cap=round] (265.07, 77.69) -- (269.57, 82.19);

\path[draw=drawColor,line width= 0.4pt,line join=round,line cap=round] (265.07, 82.19) -- (269.57, 77.69);

\path[draw=drawColor,line width= 0.4pt,line join=round,line cap=round] (307.99,116.80) -- (312.49,121.30);

\path[draw=drawColor,line width= 0.4pt,line join=round,line cap=round] (307.99,121.30) -- (312.49,116.80);

\path[draw=drawColor,line width= 0.4pt,line join=round,line cap=round] (288.61,100.32) -- (293.11,104.82);

\path[draw=drawColor,line width= 0.4pt,line join=round,line cap=round] (288.61,104.82) -- (293.11,100.32);

\path[draw=drawColor,line width= 0.4pt,line join=round,line cap=round] (281.84, 94.44) -- (286.34, 98.94);

\path[draw=drawColor,line width= 0.4pt,line join=round,line cap=round] (281.84, 98.94) -- (286.34, 94.44);

\path[draw=drawColor,line width= 0.4pt,line join=round,line cap=round] (273.27, 86.05) -- (277.77, 90.55);

\path[draw=drawColor,line width= 0.4pt,line join=round,line cap=round] (273.27, 90.55) -- (277.77, 86.05);

\path[draw=drawColor,line width= 0.4pt,line join=round,line cap=round] (295.98,106.88) -- (300.48,111.38);

\path[draw=drawColor,line width= 0.4pt,line join=round,line cap=round] (295.98,111.38) -- (300.48,106.88);

\path[draw=drawColor,line width= 0.4pt,line join=round,line cap=round] (292.86,104.27) -- (297.36,108.77);

\path[draw=drawColor,line width= 0.4pt,line join=round,line cap=round] (292.86,108.77) -- (297.36,104.27);

\path[draw=drawColor,line width= 0.4pt,line join=round,line cap=round] (284.56, 96.81) -- (289.06,101.31);

\path[draw=drawColor,line width= 0.4pt,line join=round,line cap=round] (284.56,101.31) -- (289.06, 96.81);

\path[draw=drawColor,line width= 0.4pt,line join=round,line cap=round] (277.76, 90.24) -- (282.26, 94.74);

\path[draw=drawColor,line width= 0.4pt,line join=round,line cap=round] (277.76, 94.74) -- (282.26, 90.24);

\path[draw=drawColor,line width= 0.4pt,line join=round,line cap=round] (293.82,105.17) -- (298.32,109.67);

\path[draw=drawColor,line width= 0.4pt,line join=round,line cap=round] (293.82,109.67) -- (298.32,105.17);

\path[draw=drawColor,line width= 0.4pt,line join=round,line cap=round] (294.76,105.97) -- (299.26,110.47);

\path[draw=drawColor,line width= 0.4pt,line join=round,line cap=round] (294.76,110.47) -- (299.26,105.97);

\path[draw=drawColor,line width= 0.4pt,line join=round,line cap=round] (296.46,107.41) -- (300.96,111.91);

\path[draw=drawColor,line width= 0.4pt,line join=round,line cap=round] (296.46,111.91) -- (300.96,107.41);

\path[draw=drawColor,line width= 0.4pt,line join=round,line cap=round] (304.28,114.48) -- (308.78,118.98);

\path[draw=drawColor,line width= 0.4pt,line join=round,line cap=round] (304.28,118.98) -- (308.78,114.48);

\path[draw=drawColor,line width= 0.4pt,line join=round,line cap=round] (284.42, 96.63) -- (288.92,101.13);

\path[draw=drawColor,line width= 0.4pt,line join=round,line cap=round] (284.42,101.13) -- (288.92, 96.63);

\path[draw=drawColor,line width= 0.4pt,line join=round,line cap=round] (288.12, 99.96) -- (292.62,104.46);

\path[draw=drawColor,line width= 0.4pt,line join=round,line cap=round] (288.12,104.46) -- (292.62, 99.96);

\path[draw=drawColor,line width= 0.4pt,line join=round,line cap=round] (294.91,106.01) -- (299.41,110.51);

\path[draw=drawColor,line width= 0.4pt,line join=round,line cap=round] (294.91,110.51) -- (299.41,106.01);

\path[draw=drawColor,line width= 0.4pt,line join=round,line cap=round] (286.99, 98.88) -- (291.49,103.38);

\path[draw=drawColor,line width= 0.4pt,line join=round,line cap=round] (286.99,103.38) -- (291.49, 98.88);

\path[draw=drawColor,line width= 0.4pt,line join=round,line cap=round] (302.21,113.12) -- (306.71,117.62);

\path[draw=drawColor,line width= 0.4pt,line join=round,line cap=round] (302.21,117.62) -- (306.71,113.12);

\path[draw=drawColor,line width= 0.4pt,line join=round,line cap=round] (302.46,113.51) -- (306.96,118.01);

\path[draw=drawColor,line width= 0.4pt,line join=round,line cap=round] (302.46,118.01) -- (306.96,113.51);

\path[draw=drawColor,line width= 0.4pt,line join=round,line cap=round] (302.08,112.93) -- (306.58,117.43);

\path[draw=drawColor,line width= 0.4pt,line join=round,line cap=round] (302.08,117.43) -- (306.58,112.93);

\path[draw=drawColor,line width= 0.4pt,line join=round,line cap=round] (292.96,104.36) -- (297.46,108.86);

\path[draw=drawColor,line width= 0.4pt,line join=round,line cap=round] (292.96,108.86) -- (297.46,104.36);

\path[draw=drawColor,line width= 0.4pt,line join=round,line cap=round] (300.91,111.99) -- (305.41,116.49);

\path[draw=drawColor,line width= 0.4pt,line join=round,line cap=round] (300.91,116.49) -- (305.41,111.99);

\path[draw=drawColor,line width= 0.4pt,line join=round,line cap=round] (280.27, 93.06) -- (284.77, 97.56);

\path[draw=drawColor,line width= 0.4pt,line join=round,line cap=round] (280.27, 97.56) -- (284.77, 93.06);

\path[draw=drawColor,line width= 0.4pt,line join=round,line cap=round] (278.83, 91.42) -- (283.33, 95.92);

\path[draw=drawColor,line width= 0.4pt,line join=round,line cap=round] (278.83, 95.92) -- (283.33, 91.42);

\path[draw=drawColor,line width= 0.4pt,line join=round,line cap=round] (281.41, 94.15) -- (285.91, 98.65);

\path[draw=drawColor,line width= 0.4pt,line join=round,line cap=round] (281.41, 98.65) -- (285.91, 94.15);

\path[draw=drawColor,line width= 0.4pt,line join=round,line cap=round] (285.65, 97.83) -- (290.15,102.33);

\path[draw=drawColor,line width= 0.4pt,line join=round,line cap=round] (285.65,102.33) -- (290.15, 97.83);

\path[draw=drawColor,line width= 0.4pt,line join=round,line cap=round] (291.63,102.87) -- (296.13,107.37);

\path[draw=drawColor,line width= 0.4pt,line join=round,line cap=round] (291.63,107.37) -- (296.13,102.87);

\path[draw=drawColor,line width= 0.4pt,line join=round,line cap=round] (288.61,100.35) -- (293.11,104.85);

\path[draw=drawColor,line width= 0.4pt,line join=round,line cap=round] (288.61,104.85) -- (293.11,100.35);

\path[draw=drawColor,line width= 0.4pt,line join=round,line cap=round] (298.00,109.17) -- (302.50,113.67);

\path[draw=drawColor,line width= 0.4pt,line join=round,line cap=round] (298.00,113.67) -- (302.50,109.17);

\path[draw=drawColor,line width= 0.4pt,line join=round,line cap=round] (293.96,105.33) -- (298.46,109.83);

\path[draw=drawColor,line width= 0.4pt,line join=round,line cap=round] (293.96,109.83) -- (298.46,105.33);

\path[draw=drawColor,line width= 0.4pt,line join=round,line cap=round] (290.24,101.56) -- (294.74,106.06);

\path[draw=drawColor,line width= 0.4pt,line join=round,line cap=round] (290.24,106.06) -- (294.74,101.56);

\path[draw=drawColor,line width= 0.4pt,line join=round,line cap=round] (288.94,100.53) -- (293.44,105.03);

\path[draw=drawColor,line width= 0.4pt,line join=round,line cap=round] (288.94,105.03) -- (293.44,100.53);

\path[draw=drawColor,line width= 0.4pt,line join=round,line cap=round] (306.95,116.14) -- (311.45,120.64);

\path[draw=drawColor,line width= 0.4pt,line join=round,line cap=round] (306.95,120.64) -- (311.45,116.14);

\path[draw=drawColor,line width= 0.4pt,line join=round,line cap=round] (248.67, 56.27) -- (253.17, 60.77);

\path[draw=drawColor,line width= 0.4pt,line join=round,line cap=round] (248.67, 60.77) -- (253.17, 56.27);

\path[draw=drawColor,line width= 0.4pt,line join=round,line cap=round] (273.48, 86.14) -- (277.98, 90.64);

\path[draw=drawColor,line width= 0.4pt,line join=round,line cap=round] (273.48, 90.64) -- (277.98, 86.14);

\path[draw=drawColor,line width= 0.4pt,line join=round,line cap=round] (300.16,111.38) -- (304.66,115.88);

\path[draw=drawColor,line width= 0.4pt,line join=round,line cap=round] (300.16,115.88) -- (304.66,111.38);

\path[draw=drawColor,line width= 0.4pt,line join=round,line cap=round] (307.20,116.25) -- (311.70,120.75);

\path[draw=drawColor,line width= 0.4pt,line join=round,line cap=round] (307.20,120.75) -- (311.70,116.25);

\path[draw=drawColor,line width= 0.4pt,line join=round,line cap=round] (304.65,114.78) -- (309.15,119.28);

\path[draw=drawColor,line width= 0.4pt,line join=round,line cap=round] (304.65,119.28) -- (309.15,114.78);

\path[draw=drawColor,line width= 0.4pt,line join=round,line cap=round] (295.59,106.53) -- (300.09,111.03);

\path[draw=drawColor,line width= 0.4pt,line join=round,line cap=round] (295.59,111.03) -- (300.09,106.53);

\path[draw=drawColor,line width= 0.4pt,line join=round,line cap=round] (270.32, 83.14) -- (274.82, 87.64);

\path[draw=drawColor,line width= 0.4pt,line join=round,line cap=round] (270.32, 87.64) -- (274.82, 83.14);

\path[draw=drawColor,line width= 0.4pt,line join=round,line cap=round] (293.56,104.98) -- (298.06,109.48);

\path[draw=drawColor,line width= 0.4pt,line join=round,line cap=round] (293.56,109.48) -- (298.06,104.98);

\path[draw=drawColor,line width= 0.4pt,line join=round,line cap=round] (291.81,102.96) -- (296.31,107.46);

\path[draw=drawColor,line width= 0.4pt,line join=round,line cap=round] (291.81,107.46) -- (296.31,102.96);

\path[draw=drawColor,line width= 0.4pt,line join=round,line cap=round] (312.20,120.90) -- (316.70,125.40);

\path[draw=drawColor,line width= 0.4pt,line join=round,line cap=round] (312.20,125.40) -- (316.70,120.90);

\path[draw=drawColor,line width= 0.4pt,line join=round,line cap=round] (293.22,104.66) -- (297.72,109.16);

\path[draw=drawColor,line width= 0.4pt,line join=round,line cap=round] (293.22,109.16) -- (297.72,104.66);

\path[draw=drawColor,line width= 0.4pt,line join=round,line cap=round] (309.04,117.60) -- (313.54,122.10);

\path[draw=drawColor,line width= 0.4pt,line join=round,line cap=round] (309.04,122.10) -- (313.54,117.60);

\path[draw=drawColor,line width= 0.4pt,line join=round,line cap=round] (288.38,100.15) -- (292.88,104.65);

\path[draw=drawColor,line width= 0.4pt,line join=round,line cap=round] (288.38,104.65) -- (292.88,100.15);

\path[draw=drawColor,line width= 0.4pt,line join=round,line cap=round] (292.67,104.07) -- (297.17,108.57);

\path[draw=drawColor,line width= 0.4pt,line join=round,line cap=round] (292.67,108.57) -- (297.17,104.07);

\path[draw=drawColor,line width= 0.4pt,line join=round,line cap=round] (295.31,106.24) -- (299.81,110.74);

\path[draw=drawColor,line width= 0.4pt,line join=round,line cap=round] (295.31,110.74) -- (299.81,106.24);

\path[draw=drawColor,line width= 0.4pt,line join=round,line cap=round] (281.13, 93.89) -- (285.63, 98.39);

\path[draw=drawColor,line width= 0.4pt,line join=round,line cap=round] (281.13, 98.39) -- (285.63, 93.89);

\path[draw=drawColor,line width= 0.4pt,line join=round,line cap=round] (288.68,100.39) -- (293.18,104.89);

\path[draw=drawColor,line width= 0.4pt,line join=round,line cap=round] (288.68,104.89) -- (293.18,100.39);

\path[draw=drawColor,line width= 0.4pt,line join=round,line cap=round] (289.06,100.64) -- (293.56,105.14);

\path[draw=drawColor,line width= 0.4pt,line join=round,line cap=round] (289.06,105.14) -- (293.56,100.64);

\path[draw=drawColor,line width= 0.4pt,line join=round,line cap=round] (292.41,103.66) -- (296.91,108.16);

\path[draw=drawColor,line width= 0.4pt,line join=round,line cap=round] (292.41,108.16) -- (296.91,103.66);

\path[draw=drawColor,line width= 0.4pt,line join=round,line cap=round] (283.31, 95.89) -- (287.81,100.39);

\path[draw=drawColor,line width= 0.4pt,line join=round,line cap=round] (283.31,100.39) -- (287.81, 95.89);

\path[draw=drawColor,line width= 0.4pt,line join=round,line cap=round] (282.78, 95.35) -- (287.28, 99.85);

\path[draw=drawColor,line width= 0.4pt,line join=round,line cap=round] (282.78, 99.85) -- (287.28, 95.35);

\path[draw=drawColor,line width= 0.4pt,line join=round,line cap=round] (326.90,133.86) -- (331.40,138.36);

\path[draw=drawColor,line width= 0.4pt,line join=round,line cap=round] (326.90,138.36) -- (331.40,133.86);

\path[draw=drawColor,line width= 0.4pt,line join=round,line cap=round] (300.42,111.60) -- (304.92,116.10);

\path[draw=drawColor,line width= 0.4pt,line join=round,line cap=round] (300.42,116.10) -- (304.92,111.60);

\path[draw=drawColor,line width= 0.4pt,line join=round,line cap=round] (279.36, 91.83) -- (283.86, 96.33);

\path[draw=drawColor,line width= 0.4pt,line join=round,line cap=round] (279.36, 96.33) -- (283.86, 91.83);

\path[draw=drawColor,line width= 0.4pt,line join=round,line cap=round] (287.97, 99.79) -- (292.47,104.29);

\path[draw=drawColor,line width= 0.4pt,line join=round,line cap=round] (287.97,104.29) -- (292.47, 99.79);

\path[draw=drawColor,line width= 0.4pt,line join=round,line cap=round] (281.98, 94.63) -- (286.48, 99.13);

\path[draw=drawColor,line width= 0.4pt,line join=round,line cap=round] (281.98, 99.13) -- (286.48, 94.63);

\path[draw=drawColor,line width= 0.4pt,line join=round,line cap=round] (282.87, 95.37) -- (287.37, 99.87);

\path[draw=drawColor,line width= 0.4pt,line join=round,line cap=round] (282.87, 99.87) -- (287.37, 95.37);

\path[draw=drawColor,line width= 0.4pt,line join=round,line cap=round] (296.37,107.36) -- (300.87,111.86);

\path[draw=drawColor,line width= 0.4pt,line join=round,line cap=round] (296.37,111.86) -- (300.87,107.36);

\path[draw=drawColor,line width= 0.4pt,line join=round,line cap=round] (287.08, 98.97) -- (291.58,103.47);

\path[draw=drawColor,line width= 0.4pt,line join=round,line cap=round] (287.08,103.47) -- (291.58, 98.97);

\path[draw=drawColor,line width= 0.4pt,line join=round,line cap=round] (296.35,107.32) -- (300.85,111.82);

\path[draw=drawColor,line width= 0.4pt,line join=round,line cap=round] (296.35,111.82) -- (300.85,107.32);

\path[draw=drawColor,line width= 0.4pt,line join=round,line cap=round] (297.09,107.96) -- (301.59,112.46);

\path[draw=drawColor,line width= 0.4pt,line join=round,line cap=round] (297.09,112.46) -- (301.59,107.96);

\path[draw=drawColor,line width= 0.4pt,line join=round,line cap=round] (283.81, 96.26) -- (288.31,100.76);

\path[draw=drawColor,line width= 0.4pt,line join=round,line cap=round] (283.81,100.76) -- (288.31, 96.26);

\path[draw=drawColor,line width= 0.4pt,line join=round,line cap=round] (287.80, 99.63) -- (292.30,104.13);

\path[draw=drawColor,line width= 0.4pt,line join=round,line cap=round] (287.80,104.13) -- (292.30, 99.63);

\path[draw=drawColor,line width= 0.4pt,line join=round,line cap=round] (294.40,105.68) -- (298.90,110.18);

\path[draw=drawColor,line width= 0.4pt,line join=round,line cap=round] (294.40,110.18) -- (298.90,105.68);

\path[draw=drawColor,line width= 0.4pt,line join=round,line cap=round] (287.25, 99.16) -- (291.75,103.66);

\path[draw=drawColor,line width= 0.4pt,line join=round,line cap=round] (287.25,103.66) -- (291.75, 99.16);

\path[draw=drawColor,line width= 0.4pt,line join=round,line cap=round] (298.85,110.02) -- (303.35,114.52);

\path[draw=drawColor,line width= 0.4pt,line join=round,line cap=round] (298.85,114.52) -- (303.35,110.02);

\path[draw=drawColor,line width= 0.4pt,line join=round,line cap=round] (290.91,102.13) -- (295.41,106.63);

\path[draw=drawColor,line width= 0.4pt,line join=round,line cap=round] (290.91,106.63) -- (295.41,102.13);

\path[draw=drawColor,line width= 0.4pt,line join=round,line cap=round] (290.70,101.97) -- (295.20,106.47);

\path[draw=drawColor,line width= 0.4pt,line join=round,line cap=round] (290.70,106.47) -- (295.20,101.97);

\path[draw=drawColor,line width= 0.4pt,line join=round,line cap=round] (280.19, 93.04) -- (284.69, 97.54);

\path[draw=drawColor,line width= 0.4pt,line join=round,line cap=round] (280.19, 97.54) -- (284.69, 93.04);

\path[draw=drawColor,line width= 0.4pt,line join=round,line cap=round] (293.17,104.63) -- (297.67,109.13);

\path[draw=drawColor,line width= 0.4pt,line join=round,line cap=round] (293.17,109.13) -- (297.67,104.63);

\path[draw=drawColor,line width= 0.4pt,line join=round,line cap=round] (287.19, 99.10) -- (291.69,103.60);

\path[draw=drawColor,line width= 0.4pt,line join=round,line cap=round] (287.19,103.60) -- (291.69, 99.10);

\path[draw=drawColor,line width= 0.4pt,line join=round,line cap=round] (292.91,104.31) -- (297.41,108.81);

\path[draw=drawColor,line width= 0.4pt,line join=round,line cap=round] (292.91,108.81) -- (297.41,104.31);

\path[draw=drawColor,line width= 0.4pt,line join=round,line cap=round] (281.65, 94.33) -- (286.15, 98.83);

\path[draw=drawColor,line width= 0.4pt,line join=round,line cap=round] (281.65, 98.83) -- (286.15, 94.33);

\path[draw=drawColor,line width= 0.4pt,line join=round,line cap=round] (286.70, 98.61) -- (291.20,103.11);

\path[draw=drawColor,line width= 0.4pt,line join=round,line cap=round] (286.70,103.11) -- (291.20, 98.61);

\path[draw=drawColor,line width= 0.4pt,line join=round,line cap=round] (287.83, 99.63) -- (292.33,104.13);

\path[draw=drawColor,line width= 0.4pt,line join=round,line cap=round] (287.83,104.13) -- (292.33, 99.63);

\path[draw=drawColor,line width= 0.4pt,line join=round,line cap=round] (293.10,104.51) -- (297.60,109.01);

\path[draw=drawColor,line width= 0.4pt,line join=round,line cap=round] (293.10,109.01) -- (297.60,104.51);

\path[draw=drawColor,line width= 0.4pt,line join=round,line cap=round] (283.65, 96.16) -- (288.15,100.66);

\path[draw=drawColor,line width= 0.4pt,line join=round,line cap=round] (283.65,100.66) -- (288.15, 96.16);

\path[draw=drawColor,line width= 0.4pt,line join=round,line cap=round] (284.70, 96.84) -- (289.20,101.34);

\path[draw=drawColor,line width= 0.4pt,line join=round,line cap=round] (284.70,101.34) -- (289.20, 96.84);

\path[draw=drawColor,line width= 0.4pt,line join=round,line cap=round] (284.91, 97.05) -- (289.41,101.55);

\path[draw=drawColor,line width= 0.4pt,line join=round,line cap=round] (284.91,101.55) -- (289.41, 97.05);

\path[draw=drawColor,line width= 0.4pt,line join=round,line cap=round] (291.83,102.96) -- (296.33,107.46);

\path[draw=drawColor,line width= 0.4pt,line join=round,line cap=round] (291.83,107.46) -- (296.33,102.96);

\path[draw=drawColor,line width= 0.4pt,line join=round,line cap=round] (292.70,104.15) -- (297.20,108.65);

\path[draw=drawColor,line width= 0.4pt,line join=round,line cap=round] (292.70,108.65) -- (297.20,104.15);

\path[draw=drawColor,line width= 0.4pt,line join=round,line cap=round] (303.06,113.95) -- (307.56,118.45);

\path[draw=drawColor,line width= 0.4pt,line join=round,line cap=round] (303.06,118.45) -- (307.56,113.95);

\path[draw=drawColor,line width= 0.4pt,line join=round,line cap=round] (287.88, 99.64) -- (292.38,104.14);

\path[draw=drawColor,line width= 0.4pt,line join=round,line cap=round] (287.88,104.14) -- (292.38, 99.64);

\path[draw=drawColor,line width= 0.4pt,line join=round,line cap=round] (266.82, 79.21) -- (271.32, 83.71);

\path[draw=drawColor,line width= 0.4pt,line join=round,line cap=round] (266.82, 83.71) -- (271.32, 79.21);

\path[draw=drawColor,line width= 0.4pt,line join=round,line cap=round] (298.94,110.10) -- (303.44,114.60);

\path[draw=drawColor,line width= 0.4pt,line join=round,line cap=round] (298.94,114.60) -- (303.44,110.10);

\path[draw=drawColor,line width= 0.4pt,line join=round,line cap=round] (274.51, 87.45) -- (279.01, 91.95);

\path[draw=drawColor,line width= 0.4pt,line join=round,line cap=round] (274.51, 91.95) -- (279.01, 87.45);

\path[draw=drawColor,line width= 0.4pt,line join=round,line cap=round] (296.52,107.44) -- (301.02,111.94);

\path[draw=drawColor,line width= 0.4pt,line join=round,line cap=round] (296.52,111.94) -- (301.02,107.44);

\path[draw=drawColor,line width= 0.4pt,line join=round,line cap=round] (297.92,109.13) -- (302.42,113.63);

\path[draw=drawColor,line width= 0.4pt,line join=round,line cap=round] (297.92,113.63) -- (302.42,109.13);

\path[draw=drawColor,line width= 0.4pt,line join=round,line cap=round] (289.70,101.20) -- (294.20,105.70);

\path[draw=drawColor,line width= 0.4pt,line join=round,line cap=round] (289.70,105.70) -- (294.20,101.20);

\path[draw=drawColor,line width= 0.4pt,line join=round,line cap=round] (287.92, 99.69) -- (292.42,104.19);

\path[draw=drawColor,line width= 0.4pt,line join=round,line cap=round] (287.92,104.19) -- (292.42, 99.69);

\path[draw=drawColor,line width= 0.4pt,line join=round,line cap=round] (287.75, 99.61) -- (292.25,104.11);

\path[draw=drawColor,line width= 0.4pt,line join=round,line cap=round] (287.75,104.11) -- (292.25, 99.61);

\path[draw=drawColor,line width= 0.4pt,line join=round,line cap=round] (305.11,115.02) -- (309.61,119.52);

\path[draw=drawColor,line width= 0.4pt,line join=round,line cap=round] (305.11,119.52) -- (309.61,115.02);

\path[draw=drawColor,line width= 0.4pt,line join=round,line cap=round] (289.82,101.28) -- (294.32,105.78);

\path[draw=drawColor,line width= 0.4pt,line join=round,line cap=round] (289.82,105.78) -- (294.32,101.28);

\path[draw=drawColor,line width= 0.4pt,line join=round,line cap=round] (279.94, 92.71) -- (284.44, 97.21);

\path[draw=drawColor,line width= 0.4pt,line join=round,line cap=round] (279.94, 97.21) -- (284.44, 92.71);

\path[draw=drawColor,line width= 0.4pt,line join=round,line cap=round] (288.85,100.49) -- (293.35,104.99);

\path[draw=drawColor,line width= 0.4pt,line join=round,line cap=round] (288.85,104.99) -- (293.35,100.49);

\path[draw=drawColor,line width= 0.4pt,line join=round,line cap=round] (293.67,105.07) -- (298.17,109.57);

\path[draw=drawColor,line width= 0.4pt,line join=round,line cap=round] (293.67,109.57) -- (298.17,105.07);

\path[draw=drawColor,line width= 0.4pt,line join=round,line cap=round] (321.42,129.61) -- (325.92,134.11);

\path[draw=drawColor,line width= 0.4pt,line join=round,line cap=round] (321.42,134.11) -- (325.92,129.61);

\path[draw=drawColor,line width= 0.4pt,line join=round,line cap=round] (296.66,107.64) -- (301.16,112.14);

\path[draw=drawColor,line width= 0.4pt,line join=round,line cap=round] (296.66,112.14) -- (301.16,107.64);

\path[draw=drawColor,line width= 0.4pt,line join=round,line cap=round] (299.33,110.67) -- (303.83,115.17);

\path[draw=drawColor,line width= 0.4pt,line join=round,line cap=round] (299.33,115.17) -- (303.83,110.67);

\path[draw=drawColor,line width= 0.4pt,line join=round,line cap=round] (285.30, 97.43) -- (289.80,101.93);

\path[draw=drawColor,line width= 0.4pt,line join=round,line cap=round] (285.30,101.93) -- (289.80, 97.43);

\path[draw=drawColor,line width= 0.4pt,line join=round,line cap=round] (282.74, 95.33) -- (287.24, 99.83);

\path[draw=drawColor,line width= 0.4pt,line join=round,line cap=round] (282.74, 99.83) -- (287.24, 95.33);

\path[draw=drawColor,line width= 0.4pt,line join=round,line cap=round] (305.20,115.13) -- (309.70,119.63);

\path[draw=drawColor,line width= 0.4pt,line join=round,line cap=round] (305.20,119.63) -- (309.70,115.13);

\path[draw=drawColor,line width= 0.4pt,line join=round,line cap=round] (288.29,100.10) -- (292.79,104.60);

\path[draw=drawColor,line width= 0.4pt,line join=round,line cap=round] (288.29,104.60) -- (292.79,100.10);

\path[draw=drawColor,line width= 0.4pt,line join=round,line cap=round] (282.43, 95.03) -- (286.93, 99.53);

\path[draw=drawColor,line width= 0.4pt,line join=round,line cap=round] (282.43, 99.53) -- (286.93, 95.03);

\path[draw=drawColor,line width= 0.4pt,line join=round,line cap=round] (287.28, 99.20) -- (291.78,103.70);

\path[draw=drawColor,line width= 0.4pt,line join=round,line cap=round] (287.28,103.70) -- (291.78, 99.20);

\path[draw=drawColor,line width= 0.4pt,line join=round,line cap=round] (283.04, 95.48) -- (287.54, 99.98);

\path[draw=drawColor,line width= 0.4pt,line join=round,line cap=round] (283.04, 99.98) -- (287.54, 95.48);

\path[draw=drawColor,line width= 0.4pt,line join=round,line cap=round] (290.54,101.86) -- (295.04,106.36);

\path[draw=drawColor,line width= 0.4pt,line join=round,line cap=round] (290.54,106.36) -- (295.04,101.86);

\path[draw=drawColor,line width= 0.4pt,line join=round,line cap=round] (284.21, 96.49) -- (288.71,100.99);

\path[draw=drawColor,line width= 0.4pt,line join=round,line cap=round] (284.21,100.99) -- (288.71, 96.49);

\path[draw=drawColor,line width= 0.4pt,line join=round,line cap=round] (295.34,106.27) -- (299.84,110.77);

\path[draw=drawColor,line width= 0.4pt,line join=round,line cap=round] (295.34,110.77) -- (299.84,106.27);

\path[draw=drawColor,line width= 0.4pt,line join=round,line cap=round] (286.22, 98.21) -- (290.72,102.71);

\path[draw=drawColor,line width= 0.4pt,line join=round,line cap=round] (286.22,102.71) -- (290.72, 98.21);

\path[draw=drawColor,line width= 0.4pt,line join=round,line cap=round] (288.89,100.52) -- (293.39,105.02);

\path[draw=drawColor,line width= 0.4pt,line join=round,line cap=round] (288.89,105.02) -- (293.39,100.52);

\path[draw=drawColor,line width= 0.4pt,line join=round,line cap=round] (295.91,106.85) -- (300.41,111.35);

\path[draw=drawColor,line width= 0.4pt,line join=round,line cap=round] (295.91,111.35) -- (300.41,106.85);

\path[draw=drawColor,line width= 0.4pt,line join=round,line cap=round] (278.24, 90.81) -- (282.74, 95.31);

\path[draw=drawColor,line width= 0.4pt,line join=round,line cap=round] (278.24, 95.31) -- (282.74, 90.81);

\path[draw=drawColor,line width= 0.4pt,line join=round,line cap=round] (286.88, 98.78) -- (291.38,103.28);

\path[draw=drawColor,line width= 0.4pt,line join=round,line cap=round] (286.88,103.28) -- (291.38, 98.78);

\path[draw=drawColor,line width= 0.4pt,line join=round,line cap=round] (286.57, 98.47) -- (291.07,102.97);

\path[draw=drawColor,line width= 0.4pt,line join=round,line cap=round] (286.57,102.97) -- (291.07, 98.47);

\path[draw=drawColor,line width= 0.4pt,line join=round,line cap=round] (277.18, 89.49) -- (281.68, 93.99);

\path[draw=drawColor,line width= 0.4pt,line join=round,line cap=round] (277.18, 93.99) -- (281.68, 89.49);

\path[draw=drawColor,line width= 0.4pt,line join=round,line cap=round] (286.42, 98.34) -- (290.92,102.84);

\path[draw=drawColor,line width= 0.4pt,line join=round,line cap=round] (286.42,102.84) -- (290.92, 98.34);

\path[draw=drawColor,line width= 0.4pt,line join=round,line cap=round] (279.81, 92.30) -- (284.31, 96.80);

\path[draw=drawColor,line width= 0.4pt,line join=round,line cap=round] (279.81, 96.80) -- (284.31, 92.30);

\path[draw=drawColor,line width= 0.4pt,line join=round,line cap=round] (296.96,107.85) -- (301.46,112.35);

\path[draw=drawColor,line width= 0.4pt,line join=round,line cap=round] (296.96,112.35) -- (301.46,107.85);

\path[draw=drawColor,line width= 0.4pt,line join=round,line cap=round] (285.59, 97.77) -- (290.09,102.27);

\path[draw=drawColor,line width= 0.4pt,line join=round,line cap=round] (285.59,102.27) -- (290.09, 97.77);

\path[draw=drawColor,line width= 0.4pt,line join=round,line cap=round] (278.88, 91.44) -- (283.38, 95.94);

\path[draw=drawColor,line width= 0.4pt,line join=round,line cap=round] (278.88, 95.94) -- (283.38, 91.44);

\path[draw=drawColor,line width= 0.4pt,line join=round,line cap=round] (279.60, 92.02) -- (284.10, 96.52);

\path[draw=drawColor,line width= 0.4pt,line join=round,line cap=round] (279.60, 96.52) -- (284.10, 92.02);

\path[draw=drawColor,line width= 0.4pt,line join=round,line cap=round] (291.03,102.33) -- (295.53,106.83);

\path[draw=drawColor,line width= 0.4pt,line join=round,line cap=round] (291.03,106.83) -- (295.53,102.33);

\path[draw=drawColor,line width= 0.4pt,line join=round,line cap=round] (298.16,109.34) -- (302.66,113.84);

\path[draw=drawColor,line width= 0.4pt,line join=round,line cap=round] (298.16,113.84) -- (302.66,109.34);

\path[draw=drawColor,line width= 0.4pt,line join=round,line cap=round] (301.52,112.36) -- (306.02,116.86);

\path[draw=drawColor,line width= 0.4pt,line join=round,line cap=round] (301.52,116.86) -- (306.02,112.36);

\path[draw=drawColor,line width= 0.4pt,line join=round,line cap=round] (297.33,108.25) -- (301.83,112.75);

\path[draw=drawColor,line width= 0.4pt,line join=round,line cap=round] (297.33,112.75) -- (301.83,108.25);

\path[draw=drawColor,line width= 0.4pt,line join=round,line cap=round] (338.33,149.42) -- (342.83,153.92);

\path[draw=drawColor,line width= 0.4pt,line join=round,line cap=round] (338.33,153.92) -- (342.83,149.42);

\path[draw=drawColor,line width= 0.4pt,line join=round,line cap=round] (308.13,116.94) -- (312.63,121.44);

\path[draw=drawColor,line width= 0.4pt,line join=round,line cap=round] (308.13,121.44) -- (312.63,116.94);

\path[draw=drawColor,line width= 0.4pt,line join=round,line cap=round] (285.88, 97.97) -- (290.38,102.47);

\path[draw=drawColor,line width= 0.4pt,line join=round,line cap=round] (285.88,102.47) -- (290.38, 97.97);

\path[draw=drawColor,line width= 0.4pt,line join=round,line cap=round] (297.96,109.13) -- (302.46,113.63);

\path[draw=drawColor,line width= 0.4pt,line join=round,line cap=round] (297.96,113.63) -- (302.46,109.13);

\path[draw=drawColor,line width= 0.4pt,line join=round,line cap=round] (291.72,102.91) -- (296.22,107.41);

\path[draw=drawColor,line width= 0.4pt,line join=round,line cap=round] (291.72,107.41) -- (296.22,102.91);

\path[draw=drawColor,line width= 0.4pt,line join=round,line cap=round] (280.72, 93.47) -- (285.22, 97.97);

\path[draw=drawColor,line width= 0.4pt,line join=round,line cap=round] (280.72, 97.97) -- (285.22, 93.47);

\path[draw=drawColor,line width= 0.4pt,line join=round,line cap=round] (284.88, 97.05) -- (289.38,101.55);

\path[draw=drawColor,line width= 0.4pt,line join=round,line cap=round] (284.88,101.55) -- (289.38, 97.05);

\path[draw=drawColor,line width= 0.4pt,line join=round,line cap=round] (285.86, 97.94) -- (290.36,102.44);

\path[draw=drawColor,line width= 0.4pt,line join=round,line cap=round] (285.86,102.44) -- (290.36, 97.94);

\path[draw=drawColor,line width= 0.4pt,line join=round,line cap=round] (293.59,104.99) -- (298.09,109.49);

\path[draw=drawColor,line width= 0.4pt,line join=round,line cap=round] (293.59,109.49) -- (298.09,104.99);

\path[draw=drawColor,line width= 0.4pt,line join=round,line cap=round] (295.08,106.08) -- (299.58,110.58);

\path[draw=drawColor,line width= 0.4pt,line join=round,line cap=round] (295.08,110.58) -- (299.58,106.08);

\path[draw=drawColor,line width= 0.4pt,line join=round,line cap=round] (310.26,118.64) -- (314.76,123.14);

\path[draw=drawColor,line width= 0.4pt,line join=round,line cap=round] (310.26,123.14) -- (314.76,118.64);

\path[draw=drawColor,line width= 0.4pt,line join=round,line cap=round] (286.61, 98.53) -- (291.11,103.03);

\path[draw=drawColor,line width= 0.4pt,line join=round,line cap=round] (286.61,103.03) -- (291.11, 98.53);

\path[draw=drawColor,line width= 0.4pt,line join=round,line cap=round] (282.23, 94.82) -- (286.73, 99.32);

\path[draw=drawColor,line width= 0.4pt,line join=round,line cap=round] (282.23, 99.32) -- (286.73, 94.82);

\path[draw=drawColor,line width= 0.4pt,line join=round,line cap=round] (282.53, 95.11) -- (287.03, 99.61);

\path[draw=drawColor,line width= 0.4pt,line join=round,line cap=round] (282.53, 99.61) -- (287.03, 95.11);

\path[draw=drawColor,line width= 0.4pt,line join=round,line cap=round] (300.37,111.50) -- (304.87,116.00);

\path[draw=drawColor,line width= 0.4pt,line join=round,line cap=round] (300.37,116.00) -- (304.87,111.50);

\path[draw=drawColor,line width= 0.4pt,line join=round,line cap=round] (294.34,105.67) -- (298.84,110.17);

\path[draw=drawColor,line width= 0.4pt,line join=round,line cap=round] (294.34,110.17) -- (298.84,105.67);

\path[draw=drawColor,line width= 0.4pt,line join=round,line cap=round] (280.84, 93.55) -- (285.34, 98.05);

\path[draw=drawColor,line width= 0.4pt,line join=round,line cap=round] (280.84, 98.05) -- (285.34, 93.55);

\path[draw=drawColor,line width= 0.4pt,line join=round,line cap=round] (283.86, 96.30) -- (288.36,100.80);

\path[draw=drawColor,line width= 0.4pt,line join=round,line cap=round] (283.86,100.80) -- (288.36, 96.30);

\path[draw=drawColor,line width= 0.4pt,line join=round,line cap=round] (295.47,106.34) -- (299.97,110.84);

\path[draw=drawColor,line width= 0.4pt,line join=round,line cap=round] (295.47,110.84) -- (299.97,106.34);

\path[draw=drawColor,line width= 0.4pt,line join=round,line cap=round] (305.69,115.34) -- (310.19,119.84);

\path[draw=drawColor,line width= 0.4pt,line join=round,line cap=round] (305.69,119.84) -- (310.19,115.34);

\path[draw=drawColor,line width= 0.4pt,line join=round,line cap=round] (281.17, 93.90) -- (285.67, 98.40);

\path[draw=drawColor,line width= 0.4pt,line join=round,line cap=round] (281.17, 98.40) -- (285.67, 93.90);

\path[draw=drawColor,line width= 0.4pt,line join=round,line cap=round] (283.41, 96.01) -- (287.91,100.51);

\path[draw=drawColor,line width= 0.4pt,line join=round,line cap=round] (283.41,100.51) -- (287.91, 96.01);

\path[draw=drawColor,line width= 0.4pt,line join=round,line cap=round] (285.85, 97.97) -- (290.35,102.47);

\path[draw=drawColor,line width= 0.4pt,line join=round,line cap=round] (285.85,102.47) -- (290.35, 97.97);

\path[draw=drawColor,line width= 0.4pt,line join=round,line cap=round] (296.59,107.46) -- (301.09,111.96);

\path[draw=drawColor,line width= 0.4pt,line join=round,line cap=round] (296.59,111.96) -- (301.09,107.46);

\path[draw=drawColor,line width= 0.4pt,line join=round,line cap=round] (272.20, 85.37) -- (276.70, 89.87);

\path[draw=drawColor,line width= 0.4pt,line join=round,line cap=round] (272.20, 89.87) -- (276.70, 85.37);

\path[draw=drawColor,line width= 0.4pt,line join=round,line cap=round] (295.56,106.52) -- (300.06,111.02);

\path[draw=drawColor,line width= 0.4pt,line join=round,line cap=round] (295.56,111.02) -- (300.06,106.52);

\path[draw=drawColor,line width= 0.4pt,line join=round,line cap=round] (274.60, 87.45) -- (279.10, 91.95);

\path[draw=drawColor,line width= 0.4pt,line join=round,line cap=round] (274.60, 91.95) -- (279.10, 87.45);

\path[draw=drawColor,line width= 0.4pt,line join=round,line cap=round] (290.66,101.96) -- (295.16,106.46);

\path[draw=drawColor,line width= 0.4pt,line join=round,line cap=round] (290.66,106.46) -- (295.16,101.96);

\path[draw=drawColor,line width= 0.4pt,line join=round,line cap=round] (301.64,112.41) -- (306.14,116.91);

\path[draw=drawColor,line width= 0.4pt,line join=round,line cap=round] (301.64,116.91) -- (306.14,112.41);

\path[draw=drawColor,line width= 0.4pt,line join=round,line cap=round] (285.32, 97.44) -- (289.82,101.94);

\path[draw=drawColor,line width= 0.4pt,line join=round,line cap=round] (285.32,101.94) -- (289.82, 97.44);

\path[draw=drawColor,line width= 0.4pt,line join=round,line cap=round] (291.98,103.13) -- (296.48,107.63);

\path[draw=drawColor,line width= 0.4pt,line join=round,line cap=round] (291.98,107.63) -- (296.48,103.13);

\path[draw=drawColor,line width= 0.4pt,line join=round,line cap=round] (291.21,102.43) -- (295.71,106.93);

\path[draw=drawColor,line width= 0.4pt,line join=round,line cap=round] (291.21,106.93) -- (295.71,102.43);

\path[draw=drawColor,line width= 0.4pt,line join=round,line cap=round] (291.41,102.63) -- (295.91,107.13);

\path[draw=drawColor,line width= 0.4pt,line join=round,line cap=round] (291.41,107.13) -- (295.91,102.63);

\path[draw=drawColor,line width= 0.4pt,line join=round,line cap=round] (278.67, 91.31) -- (283.17, 95.81);

\path[draw=drawColor,line width= 0.4pt,line join=round,line cap=round] (278.67, 95.81) -- (283.17, 91.31);

\path[draw=drawColor,line width= 0.4pt,line join=round,line cap=round] (290.77,102.02) -- (295.27,106.52);

\path[draw=drawColor,line width= 0.4pt,line join=round,line cap=round] (290.77,106.52) -- (295.27,102.02);

\path[draw=drawColor,line width= 0.4pt,line join=round,line cap=round] (285.05, 97.19) -- (289.55,101.69);

\path[draw=drawColor,line width= 0.4pt,line join=round,line cap=round] (285.05,101.69) -- (289.55, 97.19);

\path[draw=drawColor,line width= 0.4pt,line join=round,line cap=round] (298.64,109.78) -- (303.14,114.28);

\path[draw=drawColor,line width= 0.4pt,line join=round,line cap=round] (298.64,114.28) -- (303.14,109.78);

\path[draw=drawColor,line width= 0.4pt,line join=round,line cap=round] (286.10, 98.12) -- (290.60,102.62);

\path[draw=drawColor,line width= 0.4pt,line join=round,line cap=round] (286.10,102.62) -- (290.60, 98.12);

\path[draw=drawColor,line width= 0.4pt,line join=round,line cap=round] (284.94, 97.06) -- (289.44,101.56);

\path[draw=drawColor,line width= 0.4pt,line join=round,line cap=round] (284.94,101.56) -- (289.44, 97.06);

\path[draw=drawColor,line width= 0.4pt,line join=round,line cap=round] (289.97,101.33) -- (294.47,105.83);

\path[draw=drawColor,line width= 0.4pt,line join=round,line cap=round] (289.97,105.83) -- (294.47,101.33);

\path[draw=drawColor,line width= 0.4pt,line join=round,line cap=round] (283.93, 96.33) -- (288.43,100.83);

\path[draw=drawColor,line width= 0.4pt,line join=round,line cap=round] (283.93,100.83) -- (288.43, 96.33);

\path[draw=drawColor,line width= 0.4pt,line join=round,line cap=round] (299.96,111.27) -- (304.46,115.77);

\path[draw=drawColor,line width= 0.4pt,line join=round,line cap=round] (299.96,115.77) -- (304.46,111.27);

\path[draw=drawColor,line width= 0.4pt,line join=round,line cap=round] (295.78,106.72) -- (300.28,111.22);

\path[draw=drawColor,line width= 0.4pt,line join=round,line cap=round] (295.78,111.22) -- (300.28,106.72);

\path[draw=drawColor,line width= 0.4pt,line join=round,line cap=round] (284.18, 96.48) -- (288.68,100.98);

\path[draw=drawColor,line width= 0.4pt,line join=round,line cap=round] (284.18,100.98) -- (288.68, 96.48);

\path[draw=drawColor,line width= 0.4pt,line join=round,line cap=round] (293.08,104.44) -- (297.58,108.94);

\path[draw=drawColor,line width= 0.4pt,line join=round,line cap=round] (293.08,108.94) -- (297.58,104.44);

\path[draw=drawColor,line width= 0.4pt,line join=round,line cap=round] (282.05, 94.73) -- (286.55, 99.23);

\path[draw=drawColor,line width= 0.4pt,line join=round,line cap=round] (282.05, 99.23) -- (286.55, 94.73);

\path[draw=drawColor,line width= 0.4pt,line join=round,line cap=round] (287.59, 99.37) -- (292.09,103.87);

\path[draw=drawColor,line width= 0.4pt,line join=round,line cap=round] (287.59,103.87) -- (292.09, 99.37);

\path[draw=drawColor,line width= 0.4pt,line join=round,line cap=round] (290.00,101.39) -- (294.50,105.89);

\path[draw=drawColor,line width= 0.4pt,line join=round,line cap=round] (290.00,105.89) -- (294.50,101.39);

\path[draw=drawColor,line width= 0.4pt,line join=round,line cap=round] (298.20,109.37) -- (302.70,113.87);

\path[draw=drawColor,line width= 0.4pt,line join=round,line cap=round] (298.20,113.87) -- (302.70,109.37);

\path[draw=drawColor,line width= 0.4pt,line join=round,line cap=round] (288.47,100.22) -- (292.97,104.72);

\path[draw=drawColor,line width= 0.4pt,line join=round,line cap=round] (288.47,104.72) -- (292.97,100.22);

\path[draw=drawColor,line width= 0.4pt,line join=round,line cap=round] (302.28,113.22) -- (306.78,117.72);

\path[draw=drawColor,line width= 0.4pt,line join=round,line cap=round] (302.28,117.72) -- (306.78,113.22);

\path[draw=drawColor,line width= 0.4pt,line join=round,line cap=round] (302.90,113.86) -- (307.40,118.36);

\path[draw=drawColor,line width= 0.4pt,line join=round,line cap=round] (302.90,118.36) -- (307.40,113.86);

\path[draw=drawColor,line width= 0.4pt,line join=round,line cap=round] (282.26, 94.84) -- (286.76, 99.34);

\path[draw=drawColor,line width= 0.4pt,line join=round,line cap=round] (282.26, 99.34) -- (286.76, 94.84);

\path[draw=drawColor,line width= 0.4pt,line join=round,line cap=round] (302.02,112.81) -- (306.52,117.31);

\path[draw=drawColor,line width= 0.4pt,line join=round,line cap=round] (302.02,117.31) -- (306.52,112.81);

\path[draw=drawColor,line width= 0.4pt,line join=round,line cap=round] (273.59, 86.39) -- (278.09, 90.89);

\path[draw=drawColor,line width= 0.4pt,line join=round,line cap=round] (273.59, 90.89) -- (278.09, 86.39);

\path[draw=drawColor,line width= 0.4pt,line join=round,line cap=round] (280.63, 93.39) -- (285.13, 97.89);

\path[draw=drawColor,line width= 0.4pt,line join=round,line cap=round] (280.63, 97.89) -- (285.13, 93.39);

\path[draw=drawColor,line width= 0.4pt,line join=round,line cap=round] (285.18, 97.28) -- (289.68,101.78);

\path[draw=drawColor,line width= 0.4pt,line join=round,line cap=round] (285.18,101.78) -- (289.68, 97.28);

\path[draw=drawColor,line width= 0.4pt,line join=round,line cap=round] (279.31, 91.79) -- (283.81, 96.29);

\path[draw=drawColor,line width= 0.4pt,line join=round,line cap=round] (279.31, 96.29) -- (283.81, 91.79);

\path[draw=drawColor,line width= 0.4pt,line join=round,line cap=round] (285.49, 97.63) -- (289.99,102.13);

\path[draw=drawColor,line width= 0.4pt,line join=round,line cap=round] (285.49,102.13) -- (289.99, 97.63);

\path[draw=drawColor,line width= 0.4pt,line join=round,line cap=round] (286.17, 98.17) -- (290.67,102.67);

\path[draw=drawColor,line width= 0.4pt,line join=round,line cap=round] (286.17,102.67) -- (290.67, 98.17);

\path[draw=drawColor,line width= 0.4pt,line join=round,line cap=round] (292.69,104.13) -- (297.19,108.63);

\path[draw=drawColor,line width= 0.4pt,line join=round,line cap=round] (292.69,108.63) -- (297.19,104.13);

\path[draw=drawColor,line width= 0.4pt,line join=round,line cap=round] (286.30, 98.24) -- (290.80,102.74);

\path[draw=drawColor,line width= 0.4pt,line join=round,line cap=round] (286.30,102.74) -- (290.80, 98.24);

\path[draw=drawColor,line width= 0.4pt,line join=round,line cap=round] (297.58,108.37) -- (302.08,112.87);

\path[draw=drawColor,line width= 0.4pt,line join=round,line cap=round] (297.58,112.87) -- (302.08,108.37);

\path[draw=drawColor,line width= 0.4pt,line join=round,line cap=round] (308.57,117.16) -- (313.07,121.66);

\path[draw=drawColor,line width= 0.4pt,line join=round,line cap=round] (308.57,121.66) -- (313.07,117.16);

\path[draw=drawColor,line width= 0.4pt,line join=round,line cap=round] (306.37,115.71) -- (310.87,120.21);

\path[draw=drawColor,line width= 0.4pt,line join=round,line cap=round] (306.37,120.21) -- (310.87,115.71);

\path[draw=drawColor,line width= 0.4pt,line join=round,line cap=round] (286.38, 98.32) -- (290.88,102.82);

\path[draw=drawColor,line width= 0.4pt,line join=round,line cap=round] (286.38,102.82) -- (290.88, 98.32);

\path[draw=drawColor,line width= 0.4pt,line join=round,line cap=round] (296.21,107.20) -- (300.71,111.70);

\path[draw=drawColor,line width= 0.4pt,line join=round,line cap=round] (296.21,111.70) -- (300.71,107.20);

\path[draw=drawColor,line width= 0.4pt,line join=round,line cap=round] (283.44, 96.02) -- (287.94,100.52);

\path[draw=drawColor,line width= 0.4pt,line join=round,line cap=round] (283.44,100.52) -- (287.94, 96.02);

\path[draw=drawColor,line width= 0.4pt,line join=round,line cap=round] (293.52,104.95) -- (298.02,109.45);

\path[draw=drawColor,line width= 0.4pt,line join=round,line cap=round] (293.52,109.45) -- (298.02,104.95);

\path[draw=drawColor,line width= 0.4pt,line join=round,line cap=round] (283.55, 96.08) -- (288.05,100.58);

\path[draw=drawColor,line width= 0.4pt,line join=round,line cap=round] (283.55,100.58) -- (288.05, 96.08);

\path[draw=drawColor,line width= 0.4pt,line join=round,line cap=round] (290.40,101.71) -- (294.90,106.21);

\path[draw=drawColor,line width= 0.4pt,line join=round,line cap=round] (290.40,106.21) -- (294.90,101.71);

\path[draw=drawColor,line width= 0.4pt,line join=round,line cap=round] (270.48, 83.20) -- (274.98, 87.70);

\path[draw=drawColor,line width= 0.4pt,line join=round,line cap=round] (270.48, 87.70) -- (274.98, 83.20);

\path[draw=drawColor,line width= 0.4pt,line join=round,line cap=round] (289.01,100.57) -- (293.51,105.07);

\path[draw=drawColor,line width= 0.4pt,line join=round,line cap=round] (289.01,105.07) -- (293.51,100.57);

\path[draw=drawColor,line width= 0.4pt,line join=round,line cap=round] (303.80,114.30) -- (308.30,118.80);

\path[draw=drawColor,line width= 0.4pt,line join=round,line cap=round] (303.80,118.80) -- (308.30,114.30);

\path[draw=drawColor,line width= 0.4pt,line join=round,line cap=round] (283.28, 95.80) -- (287.78,100.30);

\path[draw=drawColor,line width= 0.4pt,line join=round,line cap=round] (283.28,100.30) -- (287.78, 95.80);

\path[draw=drawColor,line width= 0.4pt,line join=round,line cap=round] (302.96,113.87) -- (307.46,118.37);

\path[draw=drawColor,line width= 0.4pt,line join=round,line cap=round] (302.96,118.37) -- (307.46,113.87);

\path[draw=drawColor,line width= 0.4pt,line join=round,line cap=round] (259.56, 72.51) -- (264.06, 77.01);

\path[draw=drawColor,line width= 0.4pt,line join=round,line cap=round] (259.56, 77.01) -- (264.06, 72.51);

\path[draw=drawColor,line width= 0.4pt,line join=round,line cap=round] (295.37,106.28) -- (299.87,110.78);

\path[draw=drawColor,line width= 0.4pt,line join=round,line cap=round] (295.37,110.78) -- (299.87,106.28);

\path[draw=drawColor,line width= 0.4pt,line join=round,line cap=round] (293.33,104.74) -- (297.83,109.24);

\path[draw=drawColor,line width= 0.4pt,line join=round,line cap=round] (293.33,109.24) -- (297.83,104.74);

\path[draw=drawColor,line width= 0.4pt,line join=round,line cap=round] (291.67,102.88) -- (296.17,107.38);

\path[draw=drawColor,line width= 0.4pt,line join=round,line cap=round] (291.67,107.38) -- (296.17,102.88);

\path[draw=drawColor,line width= 0.4pt,line join=round,line cap=round] (269.45, 82.42) -- (273.95, 86.92);

\path[draw=drawColor,line width= 0.4pt,line join=round,line cap=round] (269.45, 86.92) -- (273.95, 82.42);

\path[draw=drawColor,line width= 0.4pt,line join=round,line cap=round] (275.32, 88.04) -- (279.82, 92.54);

\path[draw=drawColor,line width= 0.4pt,line join=round,line cap=round] (275.32, 92.54) -- (279.82, 88.04);

\path[draw=drawColor,line width= 0.4pt,line join=round,line cap=round] (296.61,107.58) -- (301.11,112.08);

\path[draw=drawColor,line width= 0.4pt,line join=round,line cap=round] (296.61,112.08) -- (301.11,107.58);

\path[draw=drawColor,line width= 0.4pt,line join=round,line cap=round] (292.17,103.37) -- (296.67,107.87);

\path[draw=drawColor,line width= 0.4pt,line join=round,line cap=round] (292.17,107.87) -- (296.67,103.37);

\path[draw=drawColor,line width= 0.4pt,line join=round,line cap=round] (283.07, 95.53) -- (287.57,100.03);

\path[draw=drawColor,line width= 0.4pt,line join=round,line cap=round] (283.07,100.03) -- (287.57, 95.53);

\path[draw=drawColor,line width= 0.4pt,line join=round,line cap=round] (300.53,111.61) -- (305.03,116.11);

\path[draw=drawColor,line width= 0.4pt,line join=round,line cap=round] (300.53,116.11) -- (305.03,111.61);

\path[draw=drawColor,line width= 0.4pt,line join=round,line cap=round] (285.35, 97.47) -- (289.85,101.97);

\path[draw=drawColor,line width= 0.4pt,line join=round,line cap=round] (285.35,101.97) -- (289.85, 97.47);

\path[draw=drawColor,line width= 0.4pt,line join=round,line cap=round] (290.12,101.47) -- (294.62,105.97);

\path[draw=drawColor,line width= 0.4pt,line join=round,line cap=round] (290.12,105.97) -- (294.62,101.47);

\path[draw=drawColor,line width= 0.4pt,line join=round,line cap=round] (310.08,118.29) -- (314.58,122.79);

\path[draw=drawColor,line width= 0.4pt,line join=round,line cap=round] (310.08,122.79) -- (314.58,118.29);

\path[draw=drawColor,line width= 0.4pt,line join=round,line cap=round] (267.57, 80.42) -- (272.07, 84.92);

\path[draw=drawColor,line width= 0.4pt,line join=round,line cap=round] (267.57, 84.92) -- (272.07, 80.42);

\path[draw=drawColor,line width= 0.4pt,line join=round,line cap=round] (289.58,101.04) -- (294.08,105.54);

\path[draw=drawColor,line width= 0.4pt,line join=round,line cap=round] (289.58,105.54) -- (294.08,101.04);

\path[draw=drawColor,line width= 0.4pt,line join=round,line cap=round] (287.85, 99.63) -- (292.35,104.13);

\path[draw=drawColor,line width= 0.4pt,line join=round,line cap=round] (287.85,104.13) -- (292.35, 99.63);

\path[draw=drawColor,line width= 0.4pt,line join=round,line cap=round] (288.82,100.48) -- (293.32,104.98);

\path[draw=drawColor,line width= 0.4pt,line join=round,line cap=round] (288.82,104.98) -- (293.32,100.48);

\path[draw=drawColor,line width= 0.4pt,line join=round,line cap=round] (284.25, 96.50) -- (288.75,101.00);

\path[draw=drawColor,line width= 0.4pt,line join=round,line cap=round] (284.25,101.00) -- (288.75, 96.50);

\path[draw=drawColor,line width= 0.4pt,line join=round,line cap=round] (299.85,111.23) -- (304.35,115.73);

\path[draw=drawColor,line width= 0.4pt,line join=round,line cap=round] (299.85,115.73) -- (304.35,111.23);

\path[draw=drawColor,line width= 0.4pt,line join=round,line cap=round] (288.73,100.41) -- (293.23,104.91);

\path[draw=drawColor,line width= 0.4pt,line join=round,line cap=round] (288.73,104.91) -- (293.23,100.41);

\path[draw=drawColor,line width= 0.4pt,line join=round,line cap=round] (292.38,103.57) -- (296.88,108.07);

\path[draw=drawColor,line width= 0.4pt,line join=round,line cap=round] (292.38,108.07) -- (296.88,103.57);

\path[draw=drawColor,line width= 0.4pt,line join=round,line cap=round] (285.62, 97.82) -- (290.12,102.32);

\path[draw=drawColor,line width= 0.4pt,line join=round,line cap=round] (285.62,102.32) -- (290.12, 97.82);

\path[draw=drawColor,line width= 0.4pt,line join=round,line cap=round] (285.43, 97.59) -- (289.93,102.09);

\path[draw=drawColor,line width= 0.4pt,line join=round,line cap=round] (285.43,102.09) -- (289.93, 97.59);

\path[draw=drawColor,line width= 0.4pt,line join=round,line cap=round] (303.24,114.06) -- (307.74,118.56);

\path[draw=drawColor,line width= 0.4pt,line join=round,line cap=round] (303.24,118.56) -- (307.74,114.06);

\path[draw=drawColor,line width= 0.4pt,line join=round,line cap=round] (290.47,101.81) -- (294.97,106.31);

\path[draw=drawColor,line width= 0.4pt,line join=round,line cap=round] (290.47,106.31) -- (294.97,101.81);

\path[draw=drawColor,line width= 0.4pt,line join=round,line cap=round] (300.80,111.89) -- (305.30,116.39);

\path[draw=drawColor,line width= 0.4pt,line join=round,line cap=round] (300.80,116.39) -- (305.30,111.89);

\path[draw=drawColor,line width= 0.4pt,line join=round,line cap=round] (287.12, 99.00) -- (291.62,103.50);

\path[draw=drawColor,line width= 0.4pt,line join=round,line cap=round] (287.12,103.50) -- (291.62, 99.00);

\path[draw=drawColor,line width= 0.4pt,line join=round,line cap=round] (300.96,112.01) -- (305.46,116.51);

\path[draw=drawColor,line width= 0.4pt,line join=round,line cap=round] (300.96,116.51) -- (305.46,112.01);

\path[draw=drawColor,line width= 0.4pt,line join=round,line cap=round] (289.22,100.78) -- (293.72,105.28);

\path[draw=drawColor,line width= 0.4pt,line join=round,line cap=round] (289.22,105.28) -- (293.72,100.78);

\path[draw=drawColor,line width= 0.4pt,line join=round,line cap=round] (272.93, 85.82) -- (277.43, 90.32);

\path[draw=drawColor,line width= 0.4pt,line join=round,line cap=round] (272.93, 90.32) -- (277.43, 85.82);

\path[draw=drawColor,line width= 0.4pt,line join=round,line cap=round] (276.34, 88.83) -- (280.84, 93.33);

\path[draw=drawColor,line width= 0.4pt,line join=round,line cap=round] (276.34, 93.33) -- (280.84, 88.83);

\path[draw=drawColor,line width= 0.4pt,line join=round,line cap=round] (297.41,108.29) -- (301.91,112.79);

\path[draw=drawColor,line width= 0.4pt,line join=round,line cap=round] (297.41,112.79) -- (301.91,108.29);

\path[draw=drawColor,line width= 0.4pt,line join=round,line cap=round] (291.53,102.77) -- (296.03,107.27);

\path[draw=drawColor,line width= 0.4pt,line join=round,line cap=round] (291.53,107.27) -- (296.03,102.77);

\path[draw=drawColor,line width= 0.4pt,line join=round,line cap=round] (282.02, 94.66) -- (286.52, 99.16);

\path[draw=drawColor,line width= 0.4pt,line join=round,line cap=round] (282.02, 99.16) -- (286.52, 94.66);

\path[draw=drawColor,line width= 0.4pt,line join=round,line cap=round] (298.12,109.30) -- (302.62,113.80);

\path[draw=drawColor,line width= 0.4pt,line join=round,line cap=round] (298.12,113.80) -- (302.62,109.30);

\path[draw=drawColor,line width= 0.4pt,line join=round,line cap=round] (286.76, 98.70) -- (291.26,103.20);

\path[draw=drawColor,line width= 0.4pt,line join=round,line cap=round] (286.76,103.20) -- (291.26, 98.70);

\path[draw=drawColor,line width= 0.4pt,line join=round,line cap=round] (309.54,118.04) -- (314.04,122.54);

\path[draw=drawColor,line width= 0.4pt,line join=round,line cap=round] (309.54,122.54) -- (314.04,118.04);

\path[draw=drawColor,line width= 0.4pt,line join=round,line cap=round] (348.41,150.40) -- (352.91,154.90);

\path[draw=drawColor,line width= 0.4pt,line join=round,line cap=round] (348.41,154.90) -- (352.91,150.40);

\path[draw=drawColor,line width= 0.4pt,line join=round,line cap=round] (291.26,102.54) -- (295.76,107.04);

\path[draw=drawColor,line width= 0.4pt,line join=round,line cap=round] (291.26,107.04) -- (295.76,102.54);

\path[draw=drawColor,line width= 0.4pt,line join=round,line cap=round] (276.06, 88.59) -- (280.56, 93.09);

\path[draw=drawColor,line width= 0.4pt,line join=round,line cap=round] (276.06, 93.09) -- (280.56, 88.59);

\path[draw=drawColor,line width= 0.4pt,line join=round,line cap=round] (286.51, 98.37) -- (291.01,102.87);

\path[draw=drawColor,line width= 0.4pt,line join=round,line cap=round] (286.51,102.87) -- (291.01, 98.37);

\path[draw=drawColor,line width= 0.4pt,line join=round,line cap=round] (287.62, 99.38) -- (292.12,103.88);

\path[draw=drawColor,line width= 0.4pt,line join=round,line cap=round] (287.62,103.88) -- (292.12, 99.38);

\path[draw=drawColor,line width= 0.4pt,line join=round,line cap=round] (285.76, 97.90) -- (290.26,102.40);

\path[draw=drawColor,line width= 0.4pt,line join=round,line cap=round] (285.76,102.40) -- (290.26, 97.90);

\path[draw=drawColor,line width= 0.4pt,line join=round,line cap=round] (289.78,101.24) -- (294.28,105.74);

\path[draw=drawColor,line width= 0.4pt,line join=round,line cap=round] (289.78,105.74) -- (294.28,101.24);

\path[draw=drawColor,line width= 0.4pt,line join=round,line cap=round] (274.23, 87.09) -- (278.73, 91.59);

\path[draw=drawColor,line width= 0.4pt,line join=round,line cap=round] (274.23, 91.59) -- (278.73, 87.09);

\path[draw=drawColor,line width= 0.4pt,line join=round,line cap=round] (291.45,102.68) -- (295.95,107.18);

\path[draw=drawColor,line width= 0.4pt,line join=round,line cap=round] (291.45,107.18) -- (295.95,102.68);

\path[draw=drawColor,line width= 0.4pt,line join=round,line cap=round] (287.35, 99.28) -- (291.85,103.78);

\path[draw=drawColor,line width= 0.4pt,line join=round,line cap=round] (287.35,103.78) -- (291.85, 99.28);

\path[draw=drawColor,line width= 0.4pt,line join=round,line cap=round] (288.87,100.52) -- (293.37,105.02);

\path[draw=drawColor,line width= 0.4pt,line join=round,line cap=round] (288.87,105.02) -- (293.37,100.52);

\path[draw=drawColor,line width= 0.4pt,line join=round,line cap=round] (297.16,108.04) -- (301.66,112.54);

\path[draw=drawColor,line width= 0.4pt,line join=round,line cap=round] (297.16,112.54) -- (301.66,108.04);

\path[draw=drawColor,line width= 0.4pt,line join=round,line cap=round] (293.70,105.13) -- (298.20,109.63);

\path[draw=drawColor,line width= 0.4pt,line join=round,line cap=round] (293.70,109.63) -- (298.20,105.13);

\path[draw=drawColor,line width= 0.4pt,line join=round,line cap=round] (287.33, 99.23) -- (291.83,103.73);

\path[draw=drawColor,line width= 0.4pt,line join=round,line cap=round] (287.33,103.73) -- (291.83, 99.23);

\path[draw=drawColor,line width= 0.4pt,line join=round,line cap=round] (301.01,112.03) -- (305.51,116.53);

\path[draw=drawColor,line width= 0.4pt,line join=round,line cap=round] (301.01,116.53) -- (305.51,112.03);

\path[draw=drawColor,line width= 0.4pt,line join=round,line cap=round] (288.96,100.54) -- (293.46,105.04);

\path[draw=drawColor,line width= 0.4pt,line join=round,line cap=round] (288.96,105.04) -- (293.46,100.54);

\path[draw=drawColor,line width= 0.4pt,line join=round,line cap=round] (274.03, 86.87) -- (278.53, 91.37);

\path[draw=drawColor,line width= 0.4pt,line join=round,line cap=round] (274.03, 91.37) -- (278.53, 86.87);

\path[draw=drawColor,line width= 0.4pt,line join=round,line cap=round] (290.26,101.57) -- (294.76,106.07);

\path[draw=drawColor,line width= 0.4pt,line join=round,line cap=round] (290.26,106.07) -- (294.76,101.57);

\path[draw=drawColor,line width= 0.4pt,line join=round,line cap=round] (305.49,115.27) -- (309.99,119.77);

\path[draw=drawColor,line width= 0.4pt,line join=round,line cap=round] (305.49,119.77) -- (309.99,115.27);

\path[draw=drawColor,line width= 0.4pt,line join=round,line cap=round] (288.39,100.15) -- (292.89,104.65);

\path[draw=drawColor,line width= 0.4pt,line join=round,line cap=round] (288.39,104.65) -- (292.89,100.15);

\path[draw=drawColor,line width= 0.4pt,line join=round,line cap=round] (298.28,109.53) -- (302.78,114.03);

\path[draw=drawColor,line width= 0.4pt,line join=round,line cap=round] (298.28,114.03) -- (302.78,109.53);

\path[draw=drawColor,line width= 0.4pt,line join=round,line cap=round] (291.91,103.06) -- (296.41,107.56);

\path[draw=drawColor,line width= 0.4pt,line join=round,line cap=round] (291.91,107.56) -- (296.41,103.06);

\path[draw=drawColor,line width= 0.4pt,line join=round,line cap=round] (296.81,107.72) -- (301.31,112.22);

\path[draw=drawColor,line width= 0.4pt,line join=round,line cap=round] (296.81,112.22) -- (301.31,107.72);

\path[draw=drawColor,line width= 0.4pt,line join=round,line cap=round] (293.62,105.01) -- (298.12,109.51);

\path[draw=drawColor,line width= 0.4pt,line join=round,line cap=round] (293.62,109.51) -- (298.12,105.01);

\path[draw=drawColor,line width= 0.4pt,line join=round,line cap=round] (288.02, 99.87) -- (292.52,104.37);

\path[draw=drawColor,line width= 0.4pt,line join=round,line cap=round] (288.02,104.37) -- (292.52, 99.87);

\path[draw=drawColor,line width= 0.4pt,line join=round,line cap=round] (289.29,100.80) -- (293.79,105.30);

\path[draw=drawColor,line width= 0.4pt,line join=round,line cap=round] (289.29,105.30) -- (293.79,100.80);

\path[draw=drawColor,line width= 0.4pt,line join=round,line cap=round] (291.15,102.40) -- (295.65,106.90);

\path[draw=drawColor,line width= 0.4pt,line join=round,line cap=round] (291.15,106.90) -- (295.65,102.40);

\path[draw=drawColor,line width= 0.4pt,line join=round,line cap=round] (291.70,102.89) -- (296.20,107.39);

\path[draw=drawColor,line width= 0.4pt,line join=round,line cap=round] (291.70,107.39) -- (296.20,102.89);

\path[draw=drawColor,line width= 0.4pt,line join=round,line cap=round] (284.09, 96.41) -- (288.59,100.91);

\path[draw=drawColor,line width= 0.4pt,line join=round,line cap=round] (284.09,100.91) -- (288.59, 96.41);

\path[draw=drawColor,line width= 0.4pt,line join=round,line cap=round] (284.85, 97.00) -- (289.35,101.50);

\path[draw=drawColor,line width= 0.4pt,line join=round,line cap=round] (284.85,101.50) -- (289.35, 97.00);

\path[draw=drawColor,line width= 0.4pt,line join=round,line cap=round] (281.91, 94.52) -- (286.41, 99.02);

\path[draw=drawColor,line width= 0.4pt,line join=round,line cap=round] (281.91, 99.02) -- (286.41, 94.52);

\path[draw=drawColor,line width= 0.4pt,line join=round,line cap=round] (299.21,110.43) -- (303.71,114.93);

\path[draw=drawColor,line width= 0.4pt,line join=round,line cap=round] (299.21,114.93) -- (303.71,110.43);

\path[draw=drawColor,line width= 0.4pt,line join=round,line cap=round] (295.72,106.67) -- (300.22,111.17);

\path[draw=drawColor,line width= 0.4pt,line join=round,line cap=round] (295.72,111.17) -- (300.22,106.67);

\path[draw=drawColor,line width= 0.4pt,line join=round,line cap=round] (290.68,101.97) -- (295.18,106.47);

\path[draw=drawColor,line width= 0.4pt,line join=round,line cap=round] (290.68,106.47) -- (295.18,101.97);

\path[draw=drawColor,line width= 0.4pt,line join=round,line cap=round] (287.16, 99.07) -- (291.66,103.57);

\path[draw=drawColor,line width= 0.4pt,line join=round,line cap=round] (287.16,103.57) -- (291.66, 99.07);

\path[draw=drawColor,line width= 0.4pt,line join=round,line cap=round] (281.95, 94.60) -- (286.45, 99.10);

\path[draw=drawColor,line width= 0.4pt,line join=round,line cap=round] (281.95, 99.10) -- (286.45, 94.60);

\path[draw=drawColor,line width= 0.4pt,line join=round,line cap=round] (300.32,111.49) -- (304.82,115.99);

\path[draw=drawColor,line width= 0.4pt,line join=round,line cap=round] (300.32,115.99) -- (304.82,111.49);

\path[draw=drawColor,line width= 0.4pt,line join=round,line cap=round] (297.51,108.36) -- (302.01,112.86);

\path[draw=drawColor,line width= 0.4pt,line join=round,line cap=round] (297.51,112.86) -- (302.01,108.36);

\path[draw=drawColor,line width= 0.4pt,line join=round,line cap=round] (288.64,100.38) -- (293.14,104.88);

\path[draw=drawColor,line width= 0.4pt,line join=round,line cap=round] (288.64,104.88) -- (293.14,100.38);

\path[draw=drawColor,line width= 0.4pt,line join=round,line cap=round] (286.35, 98.31) -- (290.85,102.81);

\path[draw=drawColor,line width= 0.4pt,line join=round,line cap=round] (286.35,102.81) -- (290.85, 98.31);

\path[draw=drawColor,line width= 0.4pt,line join=round,line cap=round] (266.28, 78.92) -- (270.78, 83.42);

\path[draw=drawColor,line width= 0.4pt,line join=round,line cap=round] (266.28, 83.42) -- (270.78, 78.92);

\path[draw=drawColor,line width= 0.4pt,line join=round,line cap=round] (292.20,103.37) -- (296.70,107.87);

\path[draw=drawColor,line width= 0.4pt,line join=round,line cap=round] (292.20,107.87) -- (296.70,103.37);

\path[draw=drawColor,line width= 0.4pt,line join=round,line cap=round] (291.48,102.74) -- (295.98,107.24);

\path[draw=drawColor,line width= 0.4pt,line join=round,line cap=round] (291.48,107.24) -- (295.98,102.74);

\path[draw=drawColor,line width= 0.4pt,line join=round,line cap=round] (290.63,101.92) -- (295.13,106.42);

\path[draw=drawColor,line width= 0.4pt,line join=round,line cap=round] (290.63,106.42) -- (295.13,101.92);

\path[draw=drawColor,line width= 0.4pt,line join=round,line cap=round] (287.14, 99.04) -- (291.64,103.54);

\path[draw=drawColor,line width= 0.4pt,line join=round,line cap=round] (287.14,103.54) -- (291.64, 99.04);

\path[draw=drawColor,line width= 0.4pt,line join=round,line cap=round] (307.85,116.67) -- (312.35,121.17);

\path[draw=drawColor,line width= 0.4pt,line join=round,line cap=round] (307.85,121.17) -- (312.35,116.67);

\path[draw=drawColor,line width= 0.4pt,line join=round,line cap=round] (283.51, 96.06) -- (288.01,100.56);

\path[draw=drawColor,line width= 0.4pt,line join=round,line cap=round] (283.51,100.56) -- (288.01, 96.06);

\path[draw=drawColor,line width= 0.4pt,line join=round,line cap=round] (301.23,112.19) -- (305.73,116.69);

\path[draw=drawColor,line width= 0.4pt,line join=round,line cap=round] (301.23,116.69) -- (305.73,112.19);

\path[draw=drawColor,line width= 0.4pt,line join=round,line cap=round] (286.00, 98.06) -- (290.50,102.56);

\path[draw=drawColor,line width= 0.4pt,line join=round,line cap=round] (286.00,102.56) -- (290.50, 98.06);

\path[draw=drawColor,line width= 0.4pt,line join=round,line cap=round] (292.76,104.21) -- (297.26,108.71);

\path[draw=drawColor,line width= 0.4pt,line join=round,line cap=round] (292.76,108.71) -- (297.26,104.21);

\path[draw=drawColor,line width= 0.4pt,line join=round,line cap=round] (282.47, 95.07) -- (286.97, 99.57);

\path[draw=drawColor,line width= 0.4pt,line join=round,line cap=round] (282.47, 99.57) -- (286.97, 95.07);

\path[draw=drawColor,line width= 0.4pt,line join=round,line cap=round] (292.36,103.55) -- (296.86,108.05);

\path[draw=drawColor,line width= 0.4pt,line join=round,line cap=round] (292.36,108.05) -- (296.86,103.55);

\path[draw=drawColor,line width= 0.4pt,line join=round,line cap=round] (295.66,106.63) -- (300.16,111.13);

\path[draw=drawColor,line width= 0.4pt,line join=round,line cap=round] (295.66,111.13) -- (300.16,106.63);

\path[draw=drawColor,line width= 0.4pt,line join=round,line cap=round] (292.29,103.46) -- (296.79,107.96);

\path[draw=drawColor,line width= 0.4pt,line join=round,line cap=round] (292.29,107.96) -- (296.79,103.46);

\path[draw=drawColor,line width= 0.4pt,line join=round,line cap=round] (293.98,105.37) -- (298.48,109.87);

\path[draw=drawColor,line width= 0.4pt,line join=round,line cap=round] (293.98,109.87) -- (298.48,105.37);

\path[draw=drawColor,line width= 0.4pt,line join=round,line cap=round] (286.01, 98.11) -- (290.51,102.61);

\path[draw=drawColor,line width= 0.4pt,line join=round,line cap=round] (286.01,102.61) -- (290.51, 98.11);

\path[draw=drawColor,line width= 0.4pt,line join=round,line cap=round] (291.55,102.85) -- (296.05,107.35);

\path[draw=drawColor,line width= 0.4pt,line join=round,line cap=round] (291.55,107.35) -- (296.05,102.85);

\path[draw=drawColor,line width= 0.4pt,line join=round,line cap=round] (290.13,101.46) -- (294.63,105.96);

\path[draw=drawColor,line width= 0.4pt,line join=round,line cap=round] (290.13,105.96) -- (294.63,101.46);

\path[draw=drawColor,line width= 0.4pt,line join=round,line cap=round] (277.82, 90.35) -- (282.32, 94.85);

\path[draw=drawColor,line width= 0.4pt,line join=round,line cap=round] (277.82, 94.85) -- (282.32, 90.35);

\path[draw=drawColor,line width= 0.4pt,line join=round,line cap=round] (294.47,105.78) -- (298.97,110.28);

\path[draw=drawColor,line width= 0.4pt,line join=round,line cap=round] (294.47,110.28) -- (298.97,105.78);

\path[draw=drawColor,line width= 0.4pt,line join=round,line cap=round] (285.67, 97.83) -- (290.17,102.33);

\path[draw=drawColor,line width= 0.4pt,line join=round,line cap=round] (285.67,102.33) -- (290.17, 97.83);

\path[draw=drawColor,line width= 0.4pt,line join=round,line cap=round] (284.35, 96.59) -- (288.85,101.09);

\path[draw=drawColor,line width= 0.4pt,line join=round,line cap=round] (284.35,101.09) -- (288.85, 96.59);

\path[draw=drawColor,line width= 0.4pt,line join=round,line cap=round] (299.37,110.70) -- (303.87,115.20);

\path[draw=drawColor,line width= 0.4pt,line join=round,line cap=round] (299.37,115.20) -- (303.87,110.70);

\path[draw=drawColor,line width= 0.4pt,line join=round,line cap=round] (330.85,138.31) -- (335.35,142.81);

\path[draw=drawColor,line width= 0.4pt,line join=round,line cap=round] (330.85,142.81) -- (335.35,138.31);

\path[draw=drawColor,line width= 0.4pt,line join=round,line cap=round] (282.65, 95.21) -- (287.15, 99.71);

\path[draw=drawColor,line width= 0.4pt,line join=round,line cap=round] (282.65, 99.71) -- (287.15, 95.21);

\path[draw=drawColor,line width= 0.4pt,line join=round,line cap=round] (294.64,105.91) -- (299.14,110.41);

\path[draw=drawColor,line width= 0.4pt,line join=round,line cap=round] (294.64,110.41) -- (299.14,105.91);

\path[draw=drawColor,line width= 0.4pt,line join=round,line cap=round] (287.72, 99.50) -- (292.22,104.00);

\path[draw=drawColor,line width= 0.4pt,line join=round,line cap=round] (287.72,104.00) -- (292.22, 99.50);

\path[draw=drawColor,line width= 0.4pt,line join=round,line cap=round] (288.42,100.19) -- (292.92,104.69);

\path[draw=drawColor,line width= 0.4pt,line join=round,line cap=round] (288.42,104.69) -- (292.92,100.19);

\path[draw=drawColor,line width= 0.4pt,line join=round,line cap=round] (283.68, 96.16) -- (288.18,100.66);

\path[draw=drawColor,line width= 0.4pt,line join=round,line cap=round] (283.68,100.66) -- (288.18, 96.16);

\path[draw=drawColor,line width= 0.4pt,line join=round,line cap=round] (289.54,100.99) -- (294.04,105.49);

\path[draw=drawColor,line width= 0.4pt,line join=round,line cap=round] (289.54,105.49) -- (294.04,100.99);

\path[draw=drawColor,line width= 0.4pt,line join=round,line cap=round] (292.45,103.75) -- (296.95,108.25);

\path[draw=drawColor,line width= 0.4pt,line join=round,line cap=round] (292.45,108.25) -- (296.95,103.75);

\path[draw=drawColor,line width= 0.4pt,line join=round,line cap=round] (294.23,105.57) -- (298.73,110.07);

\path[draw=drawColor,line width= 0.4pt,line join=round,line cap=round] (294.23,110.07) -- (298.73,105.57);

\path[draw=drawColor,line width= 0.4pt,line join=round,line cap=round] (291.93,103.09) -- (296.43,107.59);

\path[draw=drawColor,line width= 0.4pt,line join=round,line cap=round] (291.93,107.59) -- (296.43,103.09);

\path[draw=drawColor,line width= 0.4pt,line join=round,line cap=round] (287.55, 99.34) -- (292.05,103.84);

\path[draw=drawColor,line width= 0.4pt,line join=round,line cap=round] (287.55,103.84) -- (292.05, 99.34);

\path[draw=drawColor,line width= 0.4pt,line join=round,line cap=round] (289.10,100.65) -- (293.60,105.15);

\path[draw=drawColor,line width= 0.4pt,line join=round,line cap=round] (289.10,105.15) -- (293.60,100.65);

\path[draw=drawColor,line width= 0.4pt,line join=round,line cap=round] (281.53, 94.22) -- (286.03, 98.72);

\path[draw=drawColor,line width= 0.4pt,line join=round,line cap=round] (281.53, 98.72) -- (286.03, 94.22);

\path[draw=drawColor,line width= 0.4pt,line join=round,line cap=round] (287.45, 99.31) -- (291.95,103.81);

\path[draw=drawColor,line width= 0.4pt,line join=round,line cap=round] (287.45,103.81) -- (291.95, 99.31);

\path[draw=drawColor,line width= 0.4pt,line join=round,line cap=round] (294.83,106.00) -- (299.33,110.50);

\path[draw=drawColor,line width= 0.4pt,line join=round,line cap=round] (294.83,110.50) -- (299.33,106.00);

\path[draw=drawColor,line width= 0.4pt,line join=round,line cap=round] (289.15,100.70) -- (293.65,105.20);

\path[draw=drawColor,line width= 0.4pt,line join=round,line cap=round] (289.15,105.20) -- (293.65,100.70);

\path[draw=drawColor,line width= 0.4pt,line join=round,line cap=round] (292.43,103.74) -- (296.93,108.24);

\path[draw=drawColor,line width= 0.4pt,line join=round,line cap=round] (292.43,108.24) -- (296.93,103.74);

\path[draw=drawColor,line width= 0.4pt,line join=round,line cap=round] (282.87, 95.42) -- (287.37, 99.92);

\path[draw=drawColor,line width= 0.4pt,line join=round,line cap=round] (282.87, 99.92) -- (287.37, 95.42);

\path[draw=drawColor,line width= 0.4pt,line join=round,line cap=round] (294.15,105.53) -- (298.65,110.03);

\path[draw=drawColor,line width= 0.4pt,line join=round,line cap=round] (294.15,110.03) -- (298.65,105.53);

\path[draw=drawColor,line width= 0.4pt,line join=round,line cap=round] (287.30, 99.21) -- (291.80,103.71);

\path[draw=drawColor,line width= 0.4pt,line join=round,line cap=round] (287.30,103.71) -- (291.80, 99.21);

\path[draw=drawColor,line width= 0.4pt,line join=round,line cap=round] (302.69,113.67) -- (307.19,118.17);

\path[draw=drawColor,line width= 0.4pt,line join=round,line cap=round] (302.69,118.17) -- (307.19,113.67);

\path[draw=drawColor,line width= 0.4pt,line join=round,line cap=round] (282.40, 94.99) -- (286.90, 99.49);

\path[draw=drawColor,line width= 0.4pt,line join=round,line cap=round] (282.40, 99.49) -- (286.90, 94.99);

\path[draw=drawColor,line width= 0.4pt,line join=round,line cap=round] (289.77,101.22) -- (294.27,105.72);

\path[draw=drawColor,line width= 0.4pt,line join=round,line cap=round] (289.77,105.72) -- (294.27,101.22);

\path[draw=drawColor,line width= 0.4pt,line join=round,line cap=round] (292.89,104.28) -- (297.39,108.78);

\path[draw=drawColor,line width= 0.4pt,line join=round,line cap=round] (292.89,108.78) -- (297.39,104.28);

\path[draw=drawColor,line width= 0.4pt,line join=round,line cap=round] (287.69, 99.46) -- (292.19,103.96);

\path[draw=drawColor,line width= 0.4pt,line join=round,line cap=round] (287.69,103.96) -- (292.19, 99.46);

\path[draw=drawColor,line width= 0.4pt,line join=round,line cap=round] (295.50,106.35) -- (300.00,110.85);

\path[draw=drawColor,line width= 0.4pt,line join=round,line cap=round] (295.50,110.85) -- (300.00,106.35);

\path[draw=drawColor,line width= 0.4pt,line join=round,line cap=round] (298.48,109.65) -- (302.98,114.15);

\path[draw=drawColor,line width= 0.4pt,line join=round,line cap=round] (298.48,114.15) -- (302.98,109.65);

\path[draw=drawColor,line width= 0.4pt,line join=round,line cap=round] (289.17,100.71) -- (293.67,105.21);

\path[draw=drawColor,line width= 0.4pt,line join=round,line cap=round] (289.17,105.21) -- (293.67,100.71);

\path[draw=drawColor,line width= 0.4pt,line join=round,line cap=round] (294.01,105.42) -- (298.51,109.92);

\path[draw=drawColor,line width= 0.4pt,line join=round,line cap=round] (294.01,109.92) -- (298.51,105.42);

\path[draw=drawColor,line width= 0.4pt,line join=round,line cap=round] (289.68,101.16) -- (294.18,105.66);

\path[draw=drawColor,line width= 0.4pt,line join=round,line cap=round] (289.68,105.66) -- (294.18,101.16);

\path[draw=drawColor,line width= 0.4pt,line join=round,line cap=round] (280.96, 93.73) -- (285.46, 98.23);

\path[draw=drawColor,line width= 0.4pt,line join=round,line cap=round] (280.96, 98.23) -- (285.46, 93.73);

\path[draw=drawColor,line width= 0.4pt,line join=round,line cap=round] (294.06,105.43) -- (298.56,109.93);

\path[draw=drawColor,line width= 0.4pt,line join=round,line cap=round] (294.06,109.93) -- (298.56,105.43);

\path[draw=drawColor,line width= 0.4pt,line join=round,line cap=round] (289.94,101.32) -- (294.44,105.82);

\path[draw=drawColor,line width= 0.4pt,line join=round,line cap=round] (289.94,105.82) -- (294.44,101.32);

\path[draw=drawColor,line width= 0.4pt,line join=round,line cap=round] (295.91,106.83) -- (300.41,111.33);

\path[draw=drawColor,line width= 0.4pt,line join=round,line cap=round] (295.91,111.33) -- (300.41,106.83);

\path[draw=drawColor,line width= 0.4pt,line join=round,line cap=round] (301.49,112.31) -- (305.99,116.81);

\path[draw=drawColor,line width= 0.4pt,line join=round,line cap=round] (301.49,116.81) -- (305.99,112.31);

\path[draw=drawColor,line width= 0.4pt,line join=round,line cap=round] (281.25, 93.94) -- (285.75, 98.44);

\path[draw=drawColor,line width= 0.4pt,line join=round,line cap=round] (281.25, 98.44) -- (285.75, 93.94);

\path[draw=drawColor,line width= 0.4pt,line join=round,line cap=round] (294.31,105.66) -- (298.81,110.16);

\path[draw=drawColor,line width= 0.4pt,line join=round,line cap=round] (294.31,110.16) -- (298.81,105.66);

\path[draw=drawColor,line width= 0.4pt,line join=round,line cap=round] (294.53,105.85) -- (299.03,110.35);

\path[draw=drawColor,line width= 0.4pt,line join=round,line cap=round] (294.53,110.35) -- (299.03,105.85);

\path[draw=drawColor,line width= 0.4pt,line join=round,line cap=round] (303.47,114.13) -- (307.97,118.63);

\path[draw=drawColor,line width= 0.4pt,line join=round,line cap=round] (303.47,118.63) -- (307.97,114.13);

\path[draw=drawColor,line width= 0.4pt,line join=round,line cap=round] (294.50,105.82) -- (299.00,110.32);

\path[draw=drawColor,line width= 0.4pt,line join=round,line cap=round] (294.50,110.32) -- (299.00,105.82);

\path[draw=drawColor,line width= 0.4pt,line join=round,line cap=round] (290.31,101.62) -- (294.81,106.12);

\path[draw=drawColor,line width= 0.4pt,line join=round,line cap=round] (290.31,106.12) -- (294.81,101.62);

\path[draw=drawColor,line width= 0.4pt,line join=round,line cap=round] (289.20,100.73) -- (293.70,105.23);

\path[draw=drawColor,line width= 0.4pt,line join=round,line cap=round] (289.20,105.23) -- (293.70,100.73);

\path[draw=drawColor,line width= 0.4pt,line join=round,line cap=round] (287.23, 99.16) -- (291.73,103.66);

\path[draw=drawColor,line width= 0.4pt,line join=round,line cap=round] (287.23,103.66) -- (291.73, 99.16);

\path[draw=drawColor,line width= 0.4pt,line join=round,line cap=round] (288.80,100.47) -- (293.30,104.97);

\path[draw=drawColor,line width= 0.4pt,line join=round,line cap=round] (288.80,104.97) -- (293.30,100.47);

\path[draw=drawColor,line width= 0.4pt,line join=round,line cap=round] (290.86,102.04) -- (295.36,106.54);

\path[draw=drawColor,line width= 0.4pt,line join=round,line cap=round] (290.86,106.54) -- (295.36,102.04);

\path[draw=drawColor,line width= 0.4pt,line join=round,line cap=round] (289.63,101.10) -- (294.13,105.60);

\path[draw=drawColor,line width= 0.4pt,line join=round,line cap=round] (289.63,105.60) -- (294.13,101.10);

\path[draw=drawColor,line width= 0.4pt,line join=round,line cap=round] (293.35,104.75) -- (297.85,109.25);

\path[draw=drawColor,line width= 0.4pt,line join=round,line cap=round] (293.35,109.25) -- (297.85,104.75);

\path[draw=drawColor,line width= 0.4pt,line join=round,line cap=round] (286.40, 98.33) -- (290.90,102.83);

\path[draw=drawColor,line width= 0.4pt,line join=round,line cap=round] (286.40,102.83) -- (290.90, 98.33);

\path[draw=drawColor,line width= 0.4pt,line join=round,line cap=round] (289.61,101.10) -- (294.11,105.60);

\path[draw=drawColor,line width= 0.4pt,line join=round,line cap=round] (289.61,105.60) -- (294.11,101.10);

\path[draw=drawColor,line width= 0.4pt,line join=round,line cap=round] (293.80,105.16) -- (298.30,109.66);

\path[draw=drawColor,line width= 0.4pt,line join=round,line cap=round] (293.80,109.66) -- (298.30,105.16);

\path[draw=drawColor,line width= 0.4pt,line join=round,line cap=round] (296.43,107.41) -- (300.93,111.91);

\path[draw=drawColor,line width= 0.4pt,line join=round,line cap=round] (296.43,111.91) -- (300.93,107.41);

\path[draw=drawColor,line width= 0.4pt,line join=round,line cap=round] (287.40, 99.29) -- (291.90,103.79);

\path[draw=drawColor,line width= 0.4pt,line join=round,line cap=round] (287.40,103.79) -- (291.90, 99.29);

\path[draw=drawColor,line width= 0.4pt,line join=round,line cap=round] (292.56,103.77) -- (297.06,108.27);

\path[draw=drawColor,line width= 0.4pt,line join=round,line cap=round] (292.56,108.27) -- (297.06,103.77);

\path[draw=drawColor,line width= 0.4pt,line join=round,line cap=round] (289.24,100.79) -- (293.74,105.29);

\path[draw=drawColor,line width= 0.4pt,line join=round,line cap=round] (289.24,105.29) -- (293.74,100.79);

\path[draw=drawColor,line width= 0.4pt,line join=round,line cap=round] (296.84,107.74) -- (301.34,112.24);

\path[draw=drawColor,line width= 0.4pt,line join=round,line cap=round] (296.84,112.24) -- (301.34,107.74);

\path[draw=drawColor,line width= 0.4pt,line join=round,line cap=round] (292.00,103.13) -- (296.50,107.63);

\path[draw=drawColor,line width= 0.4pt,line join=round,line cap=round] (292.00,107.63) -- (296.50,103.13);

\path[draw=drawColor,line width= 0.4pt,line join=round,line cap=round] (297.20,108.06) -- (301.70,112.56);

\path[draw=drawColor,line width= 0.4pt,line join=round,line cap=round] (297.20,112.56) -- (301.70,108.06);

\path[draw=drawColor,line width= 0.4pt,line join=round,line cap=round] (293.20,104.65) -- (297.70,109.15);

\path[draw=drawColor,line width= 0.4pt,line join=round,line cap=round] (293.20,109.15) -- (297.70,104.65);

\path[draw=drawColor,line width= 0.4pt,line join=round,line cap=round] (293.46,104.83) -- (297.96,109.33);

\path[draw=drawColor,line width= 0.4pt,line join=round,line cap=round] (293.46,109.33) -- (297.96,104.83);

\path[draw=drawColor,line width= 0.4pt,line join=round,line cap=round] (284.12, 96.41) -- (288.62,100.91);

\path[draw=drawColor,line width= 0.4pt,line join=round,line cap=round] (284.12,100.91) -- (288.62, 96.41);

\path[draw=drawColor,line width= 0.4pt,line join=round,line cap=round] (293.43,104.78) -- (297.93,109.28);

\path[draw=drawColor,line width= 0.4pt,line join=round,line cap=round] (293.43,109.28) -- (297.93,104.78);

\path[draw=drawColor,line width= 0.4pt,line join=round,line cap=round] (291.05,102.34) -- (295.55,106.84);

\path[draw=drawColor,line width= 0.4pt,line join=round,line cap=round] (291.05,106.84) -- (295.55,102.34);

\path[draw=drawColor,line width= 0.4pt,line join=round,line cap=round] (288.32,100.10) -- (292.82,104.60);

\path[draw=drawColor,line width= 0.4pt,line join=round,line cap=round] (288.32,104.60) -- (292.82,100.10);

\path[draw=drawColor,line width= 0.4pt,line join=round,line cap=round] (291.38,102.63) -- (295.88,107.13);

\path[draw=drawColor,line width= 0.4pt,line join=round,line cap=round] (291.38,107.13) -- (295.88,102.63);

\path[draw=drawColor,line width= 0.4pt,line join=round,line cap=round] (325.48,132.75) -- (329.98,137.25);

\path[draw=drawColor,line width= 0.4pt,line join=round,line cap=round] (325.48,137.25) -- (329.98,132.75);

\path[draw=drawColor,line width= 0.4pt,line join=round,line cap=round] (281.80, 94.41) -- (286.30, 98.91);

\path[draw=drawColor,line width= 0.4pt,line join=round,line cap=round] (281.80, 98.91) -- (286.30, 94.41);

\path[draw=drawColor,line width= 0.4pt,line join=round,line cap=round] (290.19,101.52) -- (294.69,106.02);

\path[draw=drawColor,line width= 0.4pt,line join=round,line cap=round] (290.19,106.02) -- (294.69,101.52);

\path[draw=drawColor,line width= 0.4pt,line join=round,line cap=round] (293.85,105.18) -- (298.35,109.68);

\path[draw=drawColor,line width= 0.4pt,line join=round,line cap=round] (293.85,109.68) -- (298.35,105.18);

\path[draw=drawColor,line width= 0.4pt,line join=round,line cap=round] (282.19, 94.77) -- (286.69, 99.27);

\path[draw=drawColor,line width= 0.4pt,line join=round,line cap=round] (282.19, 99.27) -- (286.69, 94.77);

\path[draw=drawColor,line width= 0.4pt,line join=round,line cap=round] (285.54, 97.74) -- (290.04,102.24);

\path[draw=drawColor,line width= 0.4pt,line join=round,line cap=round] (285.54,102.24) -- (290.04, 97.74);

\path[draw=drawColor,line width= 0.4pt,line join=round,line cap=round] (292.22,103.39) -- (296.72,107.89);

\path[draw=drawColor,line width= 0.4pt,line join=round,line cap=round] (292.22,107.89) -- (296.72,103.39);

\path[draw=drawColor,line width= 0.4pt,line join=round,line cap=round] (290.61,101.91) -- (295.11,106.41);

\path[draw=drawColor,line width= 0.4pt,line join=round,line cap=round] (290.61,106.41) -- (295.11,101.91);

\path[draw=drawColor,line width= 0.4pt,line join=round,line cap=round] (290.49,101.81) -- (294.99,106.31);

\path[draw=drawColor,line width= 0.4pt,line join=round,line cap=round] (290.49,106.31) -- (294.99,101.81);

\path[draw=drawColor,line width= 0.4pt,line join=round,line cap=round] (295.21,106.18) -- (299.71,110.68);

\path[draw=drawColor,line width= 0.4pt,line join=round,line cap=round] (295.21,110.68) -- (299.71,106.18);

\path[draw=drawColor,line width= 0.4pt,line join=round,line cap=round] (295.28,106.23) -- (299.78,110.73);

\path[draw=drawColor,line width= 0.4pt,line join=round,line cap=round] (295.28,110.73) -- (299.78,106.23);

\path[draw=drawColor,line width= 0.4pt,line join=round,line cap=round] (281.04, 93.86) -- (285.54, 98.36);

\path[draw=drawColor,line width= 0.4pt,line join=round,line cap=round] (281.04, 98.36) -- (285.54, 93.86);

\path[draw=drawColor,line width= 0.4pt,line join=round,line cap=round] (279.21, 91.71) -- (283.71, 96.21);

\path[draw=drawColor,line width= 0.4pt,line join=round,line cap=round] (279.21, 96.21) -- (283.71, 91.71);

\path[draw=drawColor,line width= 0.4pt,line join=round,line cap=round] (295.53,106.52) -- (300.03,111.02);

\path[draw=drawColor,line width= 0.4pt,line join=round,line cap=round] (295.53,111.02) -- (300.03,106.52);

\path[draw=drawColor,line width= 0.4pt,line join=round,line cap=round] (289.84,101.28) -- (294.34,105.78);

\path[draw=drawColor,line width= 0.4pt,line join=round,line cap=round] (289.84,105.78) -- (294.34,101.28);

\path[draw=drawColor,line width= 0.4pt,line join=round,line cap=round] (294.59,105.85) -- (299.09,110.35);

\path[draw=drawColor,line width= 0.4pt,line join=round,line cap=round] (294.59,110.35) -- (299.09,105.85);

\path[draw=drawColor,line width= 0.4pt,line join=round,line cap=round] (271.81, 84.52) -- (276.31, 89.02);

\path[draw=drawColor,line width= 0.4pt,line join=round,line cap=round] (271.81, 89.02) -- (276.31, 84.52);

\path[draw=drawColor,line width= 0.4pt,line join=round,line cap=round] (300.86,111.95) -- (305.36,116.45);

\path[draw=drawColor,line width= 0.4pt,line join=round,line cap=round] (300.86,116.45) -- (305.36,111.95);

\path[draw=drawColor,line width= 0.4pt,line join=round,line cap=round] (288.09, 99.94) -- (292.59,104.44);

\path[draw=drawColor,line width= 0.4pt,line join=round,line cap=round] (288.09,104.44) -- (292.59, 99.94);

\path[draw=drawColor,line width= 0.4pt,line join=round,line cap=round] (293.40,104.77) -- (297.90,109.27);

\path[draw=drawColor,line width= 0.4pt,line join=round,line cap=round] (293.40,109.27) -- (297.90,104.77);

\path[draw=drawColor,line width= 0.4pt,line join=round,line cap=round] (292.93,104.32) -- (297.43,108.82);

\path[draw=drawColor,line width= 0.4pt,line join=round,line cap=round] (292.93,108.82) -- (297.43,104.32);

\path[draw=drawColor,line width= 0.4pt,line join=round,line cap=round] (291.86,102.99) -- (296.36,107.49);

\path[draw=drawColor,line width= 0.4pt,line join=round,line cap=round] (291.86,107.49) -- (296.36,102.99);

\path[draw=drawColor,line width= 0.4pt,line join=round,line cap=round] (286.06, 98.11) -- (290.56,102.61);

\path[draw=drawColor,line width= 0.4pt,line join=round,line cap=round] (286.06,102.61) -- (290.56, 98.11);

\path[draw=drawColor,line width= 0.4pt,line join=round,line cap=round] (288.26,100.09) -- (292.76,104.59);

\path[draw=drawColor,line width= 0.4pt,line join=round,line cap=round] (288.26,104.59) -- (292.76,100.09);

\path[draw=drawColor,line width= 0.4pt,line join=round,line cap=round] (298.32,109.55) -- (302.82,114.05);

\path[draw=drawColor,line width= 0.4pt,line join=round,line cap=round] (298.32,114.05) -- (302.82,109.55);

\path[draw=drawColor,line width= 0.4pt,line join=round,line cap=round] (284.45, 96.67) -- (288.95,101.17);

\path[draw=drawColor,line width= 0.4pt,line join=round,line cap=round] (284.45,101.17) -- (288.95, 96.67);

\path[draw=drawColor,line width= 0.4pt,line join=round,line cap=round] (305.94,115.48) -- (310.44,119.98);

\path[draw=drawColor,line width= 0.4pt,line join=round,line cap=round] (305.94,119.98) -- (310.44,115.48);

\path[draw=drawColor,line width= 0.4pt,line join=round,line cap=round] (291.50,102.75) -- (296.00,107.25);

\path[draw=drawColor,line width= 0.4pt,line join=round,line cap=round] (291.50,107.25) -- (296.00,102.75);

\path[draw=drawColor,line width= 0.4pt,line join=round,line cap=round] (289.42,100.95) -- (293.92,105.45);

\path[draw=drawColor,line width= 0.4pt,line join=round,line cap=round] (289.42,105.45) -- (293.92,100.95);

\path[draw=drawColor,line width= 0.4pt,line join=round,line cap=round] (282.12, 94.75) -- (286.62, 99.25);

\path[draw=drawColor,line width= 0.4pt,line join=round,line cap=round] (282.12, 99.25) -- (286.62, 94.75);

\path[draw=drawColor,line width= 0.4pt,line join=round,line cap=round] (290.73,101.97) -- (295.23,106.47);

\path[draw=drawColor,line width= 0.4pt,line join=round,line cap=round] (290.73,106.47) -- (295.23,101.97);

\path[draw=drawColor,line width= 0.4pt,line join=round,line cap=round] (301.89,112.79) -- (306.39,117.29);

\path[draw=drawColor,line width= 0.4pt,line join=round,line cap=round] (301.89,117.29) -- (306.39,112.79);

\path[draw=drawColor,line width= 0.4pt,line join=round,line cap=round] (294.12,105.46) -- (298.62,109.96);

\path[draw=drawColor,line width= 0.4pt,line join=round,line cap=round] (294.12,109.96) -- (298.62,105.46);

\path[draw=drawColor,line width= 0.4pt,line join=round,line cap=round] (289.08,100.64) -- (293.58,105.14);

\path[draw=drawColor,line width= 0.4pt,line join=round,line cap=round] (289.08,105.14) -- (293.58,100.64);

\path[draw=drawColor,line width= 0.4pt,line join=round,line cap=round] (287.67, 99.46) -- (292.17,103.96);

\path[draw=drawColor,line width= 0.4pt,line join=round,line cap=round] (287.67,103.96) -- (292.17, 99.46);

\path[draw=drawColor,line width= 0.4pt,line join=round,line cap=round] (315.90,124.05) -- (320.40,128.55);

\path[draw=drawColor,line width= 0.4pt,line join=round,line cap=round] (315.90,128.55) -- (320.40,124.05);

\path[draw=drawColor,line width= 0.4pt,line join=round,line cap=round] (278.30, 90.91) -- (282.80, 95.41);

\path[draw=drawColor,line width= 0.4pt,line join=round,line cap=round] (278.30, 95.41) -- (282.80, 90.91);

\path[draw=drawColor,line width= 0.4pt,line join=round,line cap=round] (291.60,102.87) -- (296.10,107.37);

\path[draw=drawColor,line width= 0.4pt,line join=round,line cap=round] (291.60,107.37) -- (296.10,102.87);

\path[draw=drawColor,line width= 0.4pt,line join=round,line cap=round] (304.49,114.60) -- (308.99,119.10);

\path[draw=drawColor,line width= 0.4pt,line join=round,line cap=round] (304.49,119.10) -- (308.99,114.60);

\path[draw=drawColor,line width= 0.4pt,line join=round,line cap=round] (294.79,105.99) -- (299.29,110.49);

\path[draw=drawColor,line width= 0.4pt,line join=round,line cap=round] (294.79,110.49) -- (299.29,105.99);

\path[draw=drawColor,line width= 0.4pt,line join=round,line cap=round] (297.13,107.97) -- (301.63,112.47);

\path[draw=drawColor,line width= 0.4pt,line join=round,line cap=round] (297.13,112.47) -- (301.63,107.97);

\path[draw=drawColor,line width= 0.4pt,line join=round,line cap=round] (298.90,110.05) -- (303.40,114.55);

\path[draw=drawColor,line width= 0.4pt,line join=round,line cap=round] (298.90,114.55) -- (303.40,110.05);

\path[draw=drawColor,line width= 0.4pt,line join=round,line cap=round] (294.42,105.71) -- (298.92,110.21);

\path[draw=drawColor,line width= 0.4pt,line join=round,line cap=round] (294.42,110.21) -- (298.92,105.71);

\path[draw=drawColor,line width= 0.4pt,line join=round,line cap=round] (292.62,103.93) -- (297.12,108.43);

\path[draw=drawColor,line width= 0.4pt,line join=round,line cap=round] (292.62,108.43) -- (297.12,103.93);

\path[draw=drawColor,line width= 0.4pt,line join=round,line cap=round] (294.17,105.53) -- (298.67,110.03);

\path[draw=drawColor,line width= 0.4pt,line join=round,line cap=round] (294.17,110.03) -- (298.67,105.53);

\path[draw=drawColor,line width= 0.4pt,line join=round,line cap=round] (293.38,104.76) -- (297.88,109.26);

\path[draw=drawColor,line width= 0.4pt,line join=round,line cap=round] (293.38,109.26) -- (297.88,104.76);

\path[draw=drawColor,line width= 0.4pt,line join=round,line cap=round] (285.41, 97.58) -- (289.91,102.08);

\path[draw=drawColor,line width= 0.4pt,line join=round,line cap=round] (285.41,102.08) -- (289.91, 97.58);

\path[draw=drawColor,line width= 0.4pt,line join=round,line cap=round] (286.93, 98.81) -- (291.43,103.31);

\path[draw=drawColor,line width= 0.4pt,line join=round,line cap=round] (286.93,103.31) -- (291.43, 98.81);

\path[draw=drawColor,line width= 0.4pt,line join=round,line cap=round] (290.84,102.04) -- (295.34,106.54);

\path[draw=drawColor,line width= 0.4pt,line join=round,line cap=round] (290.84,106.54) -- (295.34,102.04);

\path[draw=drawColor,line width= 0.4pt,line join=round,line cap=round] (298.77,109.94) -- (303.27,114.44);

\path[draw=drawColor,line width= 0.4pt,line join=round,line cap=round] (298.77,114.44) -- (303.27,109.94);

\path[draw=drawColor,line width= 0.4pt,line join=round,line cap=round] (299.93,111.27) -- (304.43,115.77);

\path[draw=drawColor,line width= 0.4pt,line join=round,line cap=round] (299.93,115.77) -- (304.43,111.27);

\path[draw=drawColor,line width= 0.4pt,line join=round,line cap=round] (308.88,117.23) -- (313.38,121.73);

\path[draw=drawColor,line width= 0.4pt,line join=round,line cap=round] (308.88,121.73) -- (313.38,117.23);

\path[draw=drawColor,line width= 0.4pt,line join=round,line cap=round] (291.29,102.61) -- (295.79,107.11);

\path[draw=drawColor,line width= 0.4pt,line join=round,line cap=round] (291.29,107.11) -- (295.79,102.61);

\path[draw=drawColor,line width= 0.4pt,line join=round,line cap=round] (286.95, 98.85) -- (291.45,103.35);

\path[draw=drawColor,line width= 0.4pt,line join=round,line cap=round] (286.95,103.35) -- (291.45, 98.85);

\path[draw=drawColor,line width= 0.4pt,line join=round,line cap=round] (303.12,114.01) -- (307.62,118.51);

\path[draw=drawColor,line width= 0.4pt,line join=round,line cap=round] (303.12,118.51) -- (307.62,114.01);

\path[draw=drawColor,line width= 0.4pt,line join=round,line cap=round] (293.91,105.27) -- (298.41,109.77);

\path[draw=drawColor,line width= 0.4pt,line join=round,line cap=round] (293.91,109.77) -- (298.41,105.27);

\path[draw=drawColor,line width= 0.4pt,line join=round,line cap=round] (285.13, 97.24) -- (289.63,101.74);

\path[draw=drawColor,line width= 0.4pt,line join=round,line cap=round] (285.13,101.74) -- (289.63, 97.24);

\path[draw=drawColor,line width= 0.4pt,line join=round,line cap=round] (285.78, 97.92) -- (290.28,102.42);

\path[draw=drawColor,line width= 0.4pt,line join=round,line cap=round] (285.78,102.42) -- (290.28, 97.92);

\path[draw=drawColor,line width= 0.4pt,line join=round,line cap=round] (291.01,102.32) -- (295.51,106.82);

\path[draw=drawColor,line width= 0.4pt,line join=round,line cap=round] (291.01,106.82) -- (295.51,102.32);

\path[draw=drawColor,line width= 0.4pt,line join=round,line cap=round] (296.11,107.12) -- (300.61,111.62);

\path[draw=drawColor,line width= 0.4pt,line join=round,line cap=round] (296.11,111.62) -- (300.61,107.12);

\path[draw=drawColor,line width= 0.4pt,line join=round,line cap=round] (291.65,102.88) -- (296.15,107.38);

\path[draw=drawColor,line width= 0.4pt,line join=round,line cap=round] (291.65,107.38) -- (296.15,102.88);

\path[draw=drawColor,line width= 0.4pt,line join=round,line cap=round] (282.91, 95.45) -- (287.41, 99.95);

\path[draw=drawColor,line width= 0.4pt,line join=round,line cap=round] (282.91, 99.95) -- (287.41, 95.45);

\path[draw=drawColor,line width= 0.4pt,line join=round,line cap=round] (284.63, 96.84) -- (289.13,101.34);

\path[draw=drawColor,line width= 0.4pt,line join=round,line cap=round] (284.63,101.34) -- (289.13, 96.84);

\path[draw=drawColor,line width= 0.4pt,line join=round,line cap=round] (293.65,105.02) -- (298.15,109.52);

\path[draw=drawColor,line width= 0.4pt,line join=round,line cap=round] (293.65,109.52) -- (298.15,105.02);

\path[draw=drawColor,line width= 0.4pt,line join=round,line cap=round] (279.67, 92.09) -- (284.17, 96.59);

\path[draw=drawColor,line width= 0.4pt,line join=round,line cap=round] (279.67, 96.59) -- (284.17, 92.09);

\path[draw=drawColor,line width= 0.4pt,line join=round,line cap=round] (270.64, 83.32) -- (275.14, 87.82);

\path[draw=drawColor,line width= 0.4pt,line join=round,line cap=round] (270.64, 87.82) -- (275.14, 83.32);

\path[draw=drawColor,line width= 0.4pt,line join=round,line cap=round] (294.44,105.76) -- (298.94,110.26);

\path[draw=drawColor,line width= 0.4pt,line join=round,line cap=round] (294.44,110.26) -- (298.94,105.76);

\path[draw=drawColor,line width= 0.4pt,line join=round,line cap=round] (285.00, 97.13) -- (289.50,101.63);

\path[draw=drawColor,line width= 0.4pt,line join=round,line cap=round] (285.00,101.63) -- (289.50, 97.13);

\path[draw=drawColor,line width= 0.4pt,line join=round,line cap=round] (286.97, 98.88) -- (291.47,103.38);

\path[draw=drawColor,line width= 0.4pt,line join=round,line cap=round] (286.97,103.38) -- (291.47, 98.88);

\path[draw=drawColor,line width= 0.4pt,line join=round,line cap=round] (268.67, 81.66) -- (273.17, 86.16);

\path[draw=drawColor,line width= 0.4pt,line join=round,line cap=round] (268.67, 86.16) -- (273.17, 81.66);

\path[draw=drawColor,line width= 0.4pt,line join=round,line cap=round] (261.35, 73.94) -- (265.85, 78.44);

\path[draw=drawColor,line width= 0.4pt,line join=round,line cap=round] (261.35, 78.44) -- (265.85, 73.94);

\path[draw=drawColor,line width= 0.4pt,line join=round,line cap=round] (315.51,123.66) -- (320.01,128.16);

\path[draw=drawColor,line width= 0.4pt,line join=round,line cap=round] (315.51,128.16) -- (320.01,123.66);

\path[draw=drawColor,line width= 0.4pt,line join=round,line cap=round] (279.85, 92.36) -- (284.35, 96.86);

\path[draw=drawColor,line width= 0.4pt,line join=round,line cap=round] (279.85, 96.86) -- (284.35, 92.36);

\path[draw=drawColor,line width= 0.4pt,line join=round,line cap=round] (256.33, 64.77) -- (260.83, 69.27);

\path[draw=drawColor,line width= 0.4pt,line join=round,line cap=round] (256.33, 69.27) -- (260.83, 64.77);

\path[draw=drawColor,line width= 0.4pt,line join=round,line cap=round] (305.59,115.27) -- (310.09,119.77);

\path[draw=drawColor,line width= 0.4pt,line join=round,line cap=round] (305.59,119.77) -- (310.09,115.27);

\path[draw=drawColor,line width= 0.4pt,line join=round,line cap=round] (304.04,114.37) -- (308.54,118.87);

\path[draw=drawColor,line width= 0.4pt,line join=round,line cap=round] (304.04,118.87) -- (308.54,114.37);

\path[draw=drawColor,line width= 0.4pt,line join=round,line cap=round] (297.30,108.20) -- (301.80,112.70);

\path[draw=drawColor,line width= 0.4pt,line join=round,line cap=round] (297.30,112.70) -- (301.80,108.20);

\path[draw=drawColor,line width= 0.4pt,line join=round,line cap=round] (309.89,118.14) -- (314.39,122.64);

\path[draw=drawColor,line width= 0.4pt,line join=round,line cap=round] (309.89,122.64) -- (314.39,118.14);

\path[draw=drawColor,line width= 0.4pt,line join=round,line cap=round] (297.72,108.92) -- (302.22,113.42);

\path[draw=drawColor,line width= 0.4pt,line join=round,line cap=round] (297.72,113.42) -- (302.22,108.92);

\path[draw=drawColor,line width= 0.4pt,line join=round,line cap=round] (273.04, 85.82) -- (277.54, 90.32);

\path[draw=drawColor,line width= 0.4pt,line join=round,line cap=round] (273.04, 90.32) -- (277.54, 85.82);

\path[draw=drawColor,line width= 0.4pt,line join=round,line cap=round] (285.94, 97.99) -- (290.44,102.49);

\path[draw=drawColor,line width= 0.4pt,line join=round,line cap=round] (285.94,102.49) -- (290.44, 97.99);

\path[draw=drawColor,line width= 0.4pt,line join=round,line cap=round] (287.90, 99.66) -- (292.40,104.16);

\path[draw=drawColor,line width= 0.4pt,line join=round,line cap=round] (287.90,104.16) -- (292.40, 99.66);

\path[draw=drawColor,line width= 0.4pt,line join=round,line cap=round] (306.83,116.12) -- (311.33,120.62);

\path[draw=drawColor,line width= 0.4pt,line join=round,line cap=round] (306.83,120.62) -- (311.33,116.12);

\path[draw=drawColor,line width= 0.4pt,line join=round,line cap=round] (305.83,115.46) -- (310.33,119.96);

\path[draw=drawColor,line width= 0.4pt,line join=round,line cap=round] (305.83,119.96) -- (310.33,115.46);

\path[draw=drawColor,line width= 0.4pt,line join=round,line cap=round] (295.17,106.17) -- (299.67,110.67);

\path[draw=drawColor,line width= 0.4pt,line join=round,line cap=round] (295.17,110.67) -- (299.67,106.17);

\path[draw=drawColor,line width= 0.4pt,line join=round,line cap=round] (278.56, 91.29) -- (283.06, 95.79);

\path[draw=drawColor,line width= 0.4pt,line join=round,line cap=round] (278.56, 95.79) -- (283.06, 91.29);

\path[draw=drawColor,line width= 0.4pt,line join=round,line cap=round] (260.80, 73.66) -- (265.30, 78.16);

\path[draw=drawColor,line width= 0.4pt,line join=round,line cap=round] (260.80, 78.16) -- (265.30, 73.66);

\path[draw=drawColor,line width= 0.4pt,line join=round,line cap=round] (269.26, 82.22) -- (273.76, 86.72);

\path[draw=drawColor,line width= 0.4pt,line join=round,line cap=round] (269.26, 86.72) -- (273.76, 82.22);

\path[draw=drawColor,line width= 0.4pt,line join=round,line cap=round] (271.39, 84.11) -- (275.89, 88.61);

\path[draw=drawColor,line width= 0.4pt,line join=round,line cap=round] (271.39, 88.61) -- (275.89, 84.11);

\path[draw=drawColor,line width= 0.4pt,line join=round,line cap=round] (312.97,121.83) -- (317.47,126.33);

\path[draw=drawColor,line width= 0.4pt,line join=round,line cap=round] (312.97,126.33) -- (317.47,121.83);

\path[draw=drawColor,line width= 0.4pt,line join=round,line cap=round] (301.39,112.30) -- (305.89,116.80);

\path[draw=drawColor,line width= 0.4pt,line join=round,line cap=round] (301.39,116.80) -- (305.89,112.30);

\path[draw=drawColor,line width= 0.4pt,line join=round,line cap=round] (278.99, 91.47) -- (283.49, 95.97);

\path[draw=drawColor,line width= 0.4pt,line join=round,line cap=round] (278.99, 95.97) -- (283.49, 91.47);

\path[draw=drawColor,line width= 0.4pt,line join=round,line cap=round] (298.72,109.94) -- (303.22,114.44);

\path[draw=drawColor,line width= 0.4pt,line join=round,line cap=round] (298.72,114.44) -- (303.22,109.94);

\path[draw=drawColor,line width= 0.4pt,line join=round,line cap=round] (311.72,120.54) -- (316.22,125.04);

\path[draw=drawColor,line width= 0.4pt,line join=round,line cap=round] (311.72,125.04) -- (316.22,120.54);

\path[draw=drawColor,line width= 0.4pt,line join=round,line cap=round] (272.33, 85.40) -- (276.83, 89.90);

\path[draw=drawColor,line width= 0.4pt,line join=round,line cap=round] (272.33, 89.90) -- (276.83, 85.40);

\path[draw=drawColor,line width= 0.4pt,line join=round,line cap=round] (285.46, 97.62) -- (289.96,102.12);

\path[draw=drawColor,line width= 0.4pt,line join=round,line cap=round] (285.46,102.12) -- (289.96, 97.62);

\path[draw=drawColor,line width= 0.4pt,line join=round,line cap=round] (283.61, 96.15) -- (288.11,100.65);

\path[draw=drawColor,line width= 0.4pt,line join=round,line cap=round] (283.61,100.65) -- (288.11, 96.15);

\path[draw=drawColor,line width= 0.4pt,line join=round,line cap=round] (298.52,109.67) -- (303.02,114.17);

\path[draw=drawColor,line width= 0.4pt,line join=round,line cap=round] (298.52,114.17) -- (303.02,109.67);

\path[draw=drawColor,line width= 0.4pt,line join=round,line cap=round] (297.76,108.97) -- (302.26,113.47);

\path[draw=drawColor,line width= 0.4pt,line join=round,line cap=round] (297.76,113.47) -- (302.26,108.97);

\path[draw=drawColor,line width= 0.4pt,line join=round,line cap=round] (269.98, 82.83) -- (274.48, 87.33);

\path[draw=drawColor,line width= 0.4pt,line join=round,line cap=round] (269.98, 87.33) -- (274.48, 82.83);

\path[draw=drawColor,line width= 0.4pt,line join=round,line cap=round] (298.04,109.18) -- (302.54,113.68);

\path[draw=drawColor,line width= 0.4pt,line join=round,line cap=round] (298.04,113.68) -- (302.54,109.18);

\path[draw=drawColor,line width= 0.4pt,line join=round,line cap=round] (297.54,108.37) -- (302.04,112.87);

\path[draw=drawColor,line width= 0.4pt,line join=round,line cap=round] (297.54,112.87) -- (302.04,108.37);

\path[draw=drawColor,line width= 0.4pt,line join=round,line cap=round] (254.05, 62.68) -- (258.55, 67.18);

\path[draw=drawColor,line width= 0.4pt,line join=round,line cap=round] (254.05, 67.18) -- (258.55, 62.68);

\path[draw=drawColor,line width= 0.4pt,line join=round,line cap=round] (276.96, 89.35) -- (281.46, 93.85);

\path[draw=drawColor,line width= 0.4pt,line join=round,line cap=round] (276.96, 93.85) -- (281.46, 89.35);

\path[draw=drawColor,line width= 0.4pt,line join=round,line cap=round] (314.14,122.83) -- (318.64,127.33);

\path[draw=drawColor,line width= 0.4pt,line join=round,line cap=round] (314.14,127.33) -- (318.64,122.83);

\path[draw=drawColor,line width= 0.4pt,line join=round,line cap=round] (278.72, 91.36) -- (283.22, 95.86);

\path[draw=drawColor,line width= 0.4pt,line join=round,line cap=round] (278.72, 95.86) -- (283.22, 91.36);

\path[draw=drawColor,line width= 0.4pt,line join=round,line cap=round] (280.15, 92.87) -- (284.65, 97.37);

\path[draw=drawColor,line width= 0.4pt,line join=round,line cap=round] (280.15, 97.37) -- (284.65, 92.87);

\path[draw=drawColor,line width= 0.4pt,line join=round,line cap=round] (299.63,110.94) -- (304.13,115.44);

\path[draw=drawColor,line width= 0.4pt,line join=round,line cap=round] (299.63,115.44) -- (304.13,110.94);

\path[draw=drawColor,line width= 0.4pt,line join=round,line cap=round] (272.45, 85.41) -- (276.95, 89.91);

\path[draw=drawColor,line width= 0.4pt,line join=round,line cap=round] (272.45, 89.91) -- (276.95, 85.41);

\path[draw=drawColor,line width= 0.4pt,line join=round,line cap=round] (295.40,106.33) -- (299.90,110.83);

\path[draw=drawColor,line width= 0.4pt,line join=round,line cap=round] (295.40,110.83) -- (299.90,106.33);

\path[draw=drawColor,line width= 0.4pt,line join=round,line cap=round] (290.28,101.58) -- (294.78,106.08);

\path[draw=drawColor,line width= 0.4pt,line join=round,line cap=round] (290.28,106.08) -- (294.78,101.58);

\path[draw=drawColor,line width= 0.4pt,line join=round,line cap=round] (303.64,114.17) -- (308.14,118.67);

\path[draw=drawColor,line width= 0.4pt,line join=round,line cap=round] (303.64,118.67) -- (308.14,114.17);

\path[draw=drawColor,line width= 0.4pt,line join=round,line cap=round] (265.39, 78.01) -- (269.89, 82.51);

\path[draw=drawColor,line width= 0.4pt,line join=round,line cap=round] (265.39, 82.51) -- (269.89, 78.01);

\path[draw=drawColor,line width= 0.4pt,line join=round,line cap=round] (290.08,101.46) -- (294.58,105.96);

\path[draw=drawColor,line width= 0.4pt,line join=round,line cap=round] (290.08,105.96) -- (294.58,101.46);

\path[draw=drawColor,line width= 0.4pt,line join=round,line cap=round] (315.15,123.49) -- (319.65,127.99);

\path[draw=drawColor,line width= 0.4pt,line join=round,line cap=round] (315.15,127.99) -- (319.65,123.49);

\path[draw=drawColor,line width= 0.4pt,line join=round,line cap=round] (282.71, 95.29) -- (287.21, 99.79);

\path[draw=drawColor,line width= 0.4pt,line join=round,line cap=round] (282.71, 99.79) -- (287.21, 95.29);

\path[draw=drawColor,line width= 0.4pt,line join=round,line cap=round] (306.48,115.76) -- (310.98,120.26);

\path[draw=drawColor,line width= 0.4pt,line join=round,line cap=round] (306.48,120.26) -- (310.98,115.76);

\path[draw=drawColor,line width= 0.4pt,line join=round,line cap=round] (290.05,101.44) -- (294.55,105.94);

\path[draw=drawColor,line width= 0.4pt,line join=round,line cap=round] (290.05,105.94) -- (294.55,101.44);

\path[draw=drawColor,line width= 0.4pt,line join=round,line cap=round] (298.08,109.28) -- (302.58,113.78);

\path[draw=drawColor,line width= 0.4pt,line join=round,line cap=round] (298.08,113.78) -- (302.58,109.28);

\path[draw=drawColor,line width= 0.4pt,line join=round,line cap=round] (294.96,106.04) -- (299.46,110.54);

\path[draw=drawColor,line width= 0.4pt,line join=round,line cap=round] (294.96,110.54) -- (299.46,106.04);

\path[draw=drawColor,line width= 0.4pt,line join=round,line cap=round] (301.18,112.18) -- (305.68,116.68);

\path[draw=drawColor,line width= 0.4pt,line join=round,line cap=round] (301.18,116.68) -- (305.68,112.18);

\path[draw=drawColor,line width= 0.4pt,line join=round,line cap=round] (288.17,100.00) -- (292.67,104.50);

\path[draw=drawColor,line width= 0.4pt,line join=round,line cap=round] (288.17,104.50) -- (292.67,100.00);

\path[draw=drawColor,line width= 0.4pt,line join=round,line cap=round] (278.14, 90.70) -- (282.64, 95.20);

\path[draw=drawColor,line width= 0.4pt,line join=round,line cap=round] (278.14, 95.20) -- (282.64, 90.70);

\path[draw=drawColor,line width= 0.4pt,line join=round,line cap=round] (292.75,104.19) -- (297.25,108.69);

\path[draw=drawColor,line width= 0.4pt,line join=round,line cap=round] (292.75,108.69) -- (297.25,104.19);

\path[draw=drawColor,line width= 0.4pt,line join=round,line cap=round] (252.63, 59.29) -- (257.13, 63.79);

\path[draw=drawColor,line width= 0.4pt,line join=round,line cap=round] (252.63, 63.79) -- (257.13, 59.29);

\path[draw=drawColor,line width= 0.4pt,line join=round,line cap=round] (293.02,104.42) -- (297.52,108.92);

\path[draw=drawColor,line width= 0.4pt,line join=round,line cap=round] (293.02,108.92) -- (297.52,104.42);

\path[draw=drawColor,line width= 0.4pt,line join=round,line cap=round] (295.05,106.06) -- (299.55,110.56);

\path[draw=drawColor,line width= 0.4pt,line join=round,line cap=round] (295.05,110.56) -- (299.55,106.06);

\path[draw=drawColor,line width= 0.4pt,line join=round,line cap=round] (281.21, 93.92) -- (285.71, 98.42);

\path[draw=drawColor,line width= 0.4pt,line join=round,line cap=round] (281.21, 98.42) -- (285.71, 93.92);

\path[draw=drawColor,line width= 0.4pt,line join=round,line cap=round] (287.64, 99.42) -- (292.14,103.92);

\path[draw=drawColor,line width= 0.4pt,line join=round,line cap=round] (287.64,103.92) -- (292.14, 99.42);

\path[draw=drawColor,line width= 0.4pt,line join=round,line cap=round] (289.92,101.32) -- (294.42,105.82);

\path[draw=drawColor,line width= 0.4pt,line join=round,line cap=round] (289.92,105.82) -- (294.42,101.32);

\path[draw=drawColor,line width= 0.4pt,line join=round,line cap=round] (286.59, 98.51) -- (291.09,103.01);

\path[draw=drawColor,line width= 0.4pt,line join=round,line cap=round] (286.59,103.01) -- (291.09, 98.51);

\path[draw=drawColor,line width= 0.4pt,line join=round,line cap=round] (289.13,100.68) -- (293.63,105.18);

\path[draw=drawColor,line width= 0.4pt,line join=round,line cap=round] (289.13,105.18) -- (293.63,100.68);

\path[draw=drawColor,line width= 0.4pt,line join=round,line cap=round] (301.77,112.69) -- (306.27,117.19);

\path[draw=drawColor,line width= 0.4pt,line join=round,line cap=round] (301.77,117.19) -- (306.27,112.69);

\path[draw=drawColor,line width= 0.4pt,line join=round,line cap=round] (307.07,116.24) -- (311.57,120.74);

\path[draw=drawColor,line width= 0.4pt,line join=round,line cap=round] (307.07,120.74) -- (311.57,116.24);

\path[draw=drawColor,line width= 0.4pt,line join=round,line cap=round] (290.21,101.52) -- (294.71,106.02);

\path[draw=drawColor,line width= 0.4pt,line join=round,line cap=round] (290.21,106.02) -- (294.71,101.52);

\path[draw=drawColor,line width= 0.4pt,line join=round,line cap=round] (301.70,112.60) -- (306.20,117.10);

\path[draw=drawColor,line width= 0.4pt,line join=round,line cap=round] (301.70,117.10) -- (306.20,112.60);

\path[draw=drawColor,line width= 0.4pt,line join=round,line cap=round] (295.88,106.83) -- (300.38,111.33);

\path[draw=drawColor,line width= 0.4pt,line join=round,line cap=round] (295.88,111.33) -- (300.38,106.83);

\path[draw=drawColor,line width= 0.4pt,line join=round,line cap=round] (293.45,104.80) -- (297.95,109.30);

\path[draw=drawColor,line width= 0.4pt,line join=round,line cap=round] (293.45,109.30) -- (297.95,104.80);

\path[draw=drawColor,line width= 0.4pt,line join=round,line cap=round] (291.53,102.81) -- (296.03,107.31);

\path[draw=drawColor,line width= 0.4pt,line join=round,line cap=round] (291.53,107.31) -- (296.03,102.81);

\path[draw=drawColor,line width= 0.4pt,line join=round,line cap=round] (277.38, 89.65) -- (281.88, 94.15);

\path[draw=drawColor,line width= 0.4pt,line join=round,line cap=round] (277.38, 94.15) -- (281.88, 89.65);

\path[draw=drawColor,line width= 0.4pt,line join=round,line cap=round] (295.24,106.20) -- (299.74,110.70);

\path[draw=drawColor,line width= 0.4pt,line join=round,line cap=round] (295.24,110.70) -- (299.74,106.20);

\path[draw=drawColor,line width= 0.4pt,line join=round,line cap=round] (289.03,100.58) -- (293.53,105.08);

\path[draw=drawColor,line width= 0.4pt,line join=round,line cap=round] (289.03,105.08) -- (293.53,100.58);

\path[draw=drawColor,line width= 0.4pt,line join=round,line cap=round] (298.36,109.56) -- (302.86,114.06);

\path[draw=drawColor,line width= 0.4pt,line join=round,line cap=round] (298.36,114.06) -- (302.86,109.56);

\path[draw=drawColor,line width= 0.4pt,line join=round,line cap=round] (303.96,114.35) -- (308.46,118.85);

\path[draw=drawColor,line width= 0.4pt,line join=round,line cap=round] (303.96,118.85) -- (308.46,114.35);

\path[draw=drawColor,line width= 0.4pt,line join=round,line cap=round] (277.44, 89.78) -- (281.94, 94.28);

\path[draw=drawColor,line width= 0.4pt,line join=round,line cap=round] (277.44, 94.28) -- (281.94, 89.78);

\path[draw=drawColor,line width= 0.4pt,line join=round,line cap=round] (296.55,107.44) -- (301.05,111.94);

\path[draw=drawColor,line width= 0.4pt,line join=round,line cap=round] (296.55,111.94) -- (301.05,107.44);

\path[draw=drawColor,line width= 0.4pt,line join=round,line cap=round] (289.39,100.88) -- (293.89,105.38);

\path[draw=drawColor,line width= 0.4pt,line join=round,line cap=round] (289.39,105.38) -- (293.89,100.88);

\path[draw=drawColor,line width= 0.4pt,line join=round,line cap=round] (280.59, 93.35) -- (285.09, 97.85);

\path[draw=drawColor,line width= 0.4pt,line join=round,line cap=round] (280.59, 97.85) -- (285.09, 93.35);

\path[draw=drawColor,line width= 0.4pt,line join=round,line cap=round] (279.10, 91.64) -- (283.60, 96.14);

\path[draw=drawColor,line width= 0.4pt,line join=round,line cap=round] (279.10, 96.14) -- (283.60, 91.64);

\path[draw=drawColor,line width= 0.4pt,line join=round,line cap=round] (297.65,108.81) -- (302.15,113.31);

\path[draw=drawColor,line width= 0.4pt,line join=round,line cap=round] (297.65,113.31) -- (302.15,108.81);

\path[draw=drawColor,line width= 0.4pt,line join=round,line cap=round] (273.15, 85.98) -- (277.65, 90.48);

\path[draw=drawColor,line width= 0.4pt,line join=round,line cap=round] (273.15, 90.48) -- (277.65, 85.98);

\path[draw=drawColor,line width= 0.4pt,line join=round,line cap=round] (297.06,107.92) -- (301.56,112.42);

\path[draw=drawColor,line width= 0.4pt,line join=round,line cap=round] (297.06,112.42) -- (301.56,107.92);

\path[draw=drawColor,line width= 0.4pt,line join=round,line cap=round] (284.51, 96.70) -- (289.01,101.20);

\path[draw=drawColor,line width= 0.4pt,line join=round,line cap=round] (284.51,101.20) -- (289.01, 96.70);

\path[draw=drawColor,line width= 0.4pt,line join=round,line cap=round] (292.08,103.22) -- (296.58,107.72);

\path[draw=drawColor,line width= 0.4pt,line join=round,line cap=round] (292.08,107.72) -- (296.58,103.22);

\path[draw=drawColor,line width= 0.4pt,line join=round,line cap=round] (306.60,115.80) -- (311.10,120.30);

\path[draw=drawColor,line width= 0.4pt,line join=round,line cap=round] (306.60,120.30) -- (311.10,115.80);

\path[draw=drawColor,line width= 0.4pt,line join=round,line cap=round] (288.35,100.13) -- (292.85,104.63);

\path[draw=drawColor,line width= 0.4pt,line join=round,line cap=round] (288.35,104.63) -- (292.85,100.13);

\path[draw=drawColor,line width= 0.4pt,line join=round,line cap=round] (284.80, 96.98) -- (289.30,101.48);

\path[draw=drawColor,line width= 0.4pt,line join=round,line cap=round] (284.80,101.48) -- (289.30, 96.98);

\path[draw=drawColor,line width= 0.4pt,line join=round,line cap=round] (293.77,105.15) -- (298.27,109.65);

\path[draw=drawColor,line width= 0.4pt,line join=round,line cap=round] (293.77,109.65) -- (298.27,105.15);

\path[draw=drawColor,line width= 0.4pt,line join=round,line cap=round] (294.81,105.99) -- (299.31,110.49);

\path[draw=drawColor,line width= 0.4pt,line join=round,line cap=round] (294.81,110.49) -- (299.31,105.99);

\path[draw=drawColor,line width= 0.4pt,line join=round,line cap=round] (277.12, 89.43) -- (281.62, 93.93);

\path[draw=drawColor,line width= 0.4pt,line join=round,line cap=round] (277.12, 93.93) -- (281.62, 89.43);

\path[draw=drawColor,line width= 0.4pt,line join=round,line cap=round] (281.72, 94.39) -- (286.22, 98.89);

\path[draw=drawColor,line width= 0.4pt,line join=round,line cap=round] (281.72, 98.89) -- (286.22, 94.39);

\path[draw=drawColor,line width= 0.4pt,line join=round,line cap=round] (274.78, 87.62) -- (279.28, 92.12);

\path[draw=drawColor,line width= 0.4pt,line join=round,line cap=round] (274.78, 92.12) -- (279.28, 87.62);

\path[draw=drawColor,line width= 0.4pt,line join=round,line cap=round] (282.56, 95.19) -- (287.06, 99.69);

\path[draw=drawColor,line width= 0.4pt,line join=round,line cap=round] (282.56, 99.69) -- (287.06, 95.19);

\path[draw=drawColor,line width= 0.4pt,line join=round,line cap=round] (264.38, 77.01) -- (268.88, 81.51);

\path[draw=drawColor,line width= 0.4pt,line join=round,line cap=round] (264.38, 81.51) -- (268.88, 77.01);

\path[draw=drawColor,line width= 0.4pt,line join=round,line cap=round] (288.57,100.27) -- (293.07,104.77);

\path[draw=drawColor,line width= 0.4pt,line join=round,line cap=round] (288.57,104.77) -- (293.07,100.27);

\path[draw=drawColor,line width= 0.4pt,line join=round,line cap=round] (271.67, 84.43) -- (276.17, 88.93);

\path[draw=drawColor,line width= 0.4pt,line join=round,line cap=round] (271.67, 88.93) -- (276.17, 84.43);

\path[draw=drawColor,line width= 0.4pt,line join=round,line cap=round] (291.79,102.96) -- (296.29,107.46);

\path[draw=drawColor,line width= 0.4pt,line join=round,line cap=round] (291.79,107.46) -- (296.29,102.96);

\path[draw=drawColor,line width= 0.4pt,line join=round,line cap=round] (272.07, 85.19) -- (276.57, 89.69);

\path[draw=drawColor,line width= 0.4pt,line join=round,line cap=round] (272.07, 89.69) -- (276.57, 85.19);

\path[draw=drawColor,line width= 0.4pt,line join=round,line cap=round] (283.48, 96.06) -- (287.98,100.56);

\path[draw=drawColor,line width= 0.4pt,line join=round,line cap=round] (283.48,100.56) -- (287.98, 96.06);

\path[draw=drawColor,line width= 0.4pt,line join=round,line cap=round] (284.32, 96.55) -- (288.82,101.05);

\path[draw=drawColor,line width= 0.4pt,line join=round,line cap=round] (284.32,101.05) -- (288.82, 96.55);

\path[draw=drawColor,line width= 0.4pt,line join=round,line cap=round] (297.03,107.89) -- (301.53,112.39);

\path[draw=drawColor,line width= 0.4pt,line join=round,line cap=round] (297.03,112.39) -- (301.53,107.89);

\path[draw=drawColor,line width= 0.4pt,line join=round,line cap=round] (287.47, 99.32) -- (291.97,103.82);

\path[draw=drawColor,line width= 0.4pt,line join=round,line cap=round] (287.47,103.82) -- (291.97, 99.32);

\path[draw=drawColor,line width= 0.4pt,line join=round,line cap=round] (292.64,104.02) -- (297.14,108.52);

\path[draw=drawColor,line width= 0.4pt,line join=round,line cap=round] (292.64,108.52) -- (297.14,104.02);

\path[draw=drawColor,line width= 0.4pt,line join=round,line cap=round] (276.70, 89.16) -- (281.20, 93.66);

\path[draw=drawColor,line width= 0.4pt,line join=round,line cap=round] (276.70, 93.66) -- (281.20, 89.16);

\path[draw=drawColor,line width= 0.4pt,line join=round,line cap=round] (299.30,110.62) -- (303.80,115.12);

\path[draw=drawColor,line width= 0.4pt,line join=round,line cap=round] (299.30,115.12) -- (303.80,110.62);

\path[draw=drawColor,line width= 0.4pt,line join=round,line cap=round] (287.21, 99.12) -- (291.71,103.62);

\path[draw=drawColor,line width= 0.4pt,line join=round,line cap=round] (287.21,103.62) -- (291.71, 99.12);

\path[draw=drawColor,line width= 0.4pt,line join=round,line cap=round] (318.17,127.29) -- (322.67,131.79);

\path[draw=drawColor,line width= 0.4pt,line join=round,line cap=round] (318.17,131.79) -- (322.67,127.29);

\path[draw=drawColor,line width= 0.4pt,line join=round,line cap=round] (291.08,102.34) -- (295.58,106.84);

\path[draw=drawColor,line width= 0.4pt,line join=round,line cap=round] (291.08,106.84) -- (295.58,102.34);

\path[draw=drawColor,line width= 0.4pt,line join=round,line cap=round] (267.33, 79.98) -- (271.83, 84.48);

\path[draw=drawColor,line width= 0.4pt,line join=round,line cap=round] (267.33, 84.48) -- (271.83, 79.98);

\path[draw=drawColor,line width= 0.4pt,line join=round,line cap=round] (302.57,113.61) -- (307.07,118.11);

\path[draw=drawColor,line width= 0.4pt,line join=round,line cap=round] (302.57,118.11) -- (307.07,113.61);

\path[draw=drawColor,line width= 0.4pt,line join=round,line cap=round] (296.31,107.24) -- (300.81,111.74);

\path[draw=drawColor,line width= 0.4pt,line join=round,line cap=round] (296.31,111.74) -- (300.81,107.24);

\path[draw=drawColor,line width= 0.4pt,line join=round,line cap=round] (275.96, 88.55) -- (280.46, 93.05);

\path[draw=drawColor,line width= 0.4pt,line join=round,line cap=round] (275.96, 93.05) -- (280.46, 88.55);

\path[draw=drawColor,line width= 0.4pt,line join=round,line cap=round] (293.74,105.14) -- (298.24,109.64);

\path[draw=drawColor,line width= 0.4pt,line join=round,line cap=round] (293.74,109.64) -- (298.24,105.14);

\path[draw=drawColor,line width= 0.4pt,line join=round,line cap=round] (299.81,111.15) -- (304.31,115.65);

\path[draw=drawColor,line width= 0.4pt,line join=round,line cap=round] (299.81,115.65) -- (304.31,111.15);

\path[draw=drawColor,line width= 0.4pt,line join=round,line cap=round] (280.32, 93.14) -- (284.82, 97.64);

\path[draw=drawColor,line width= 0.4pt,line join=round,line cap=round] (280.32, 97.64) -- (284.82, 93.14);

\path[draw=drawColor,line width= 0.4pt,line join=round,line cap=round] (276.28, 88.77) -- (280.78, 93.27);

\path[draw=drawColor,line width= 0.4pt,line join=round,line cap=round] (276.28, 93.27) -- (280.78, 88.77);

\path[draw=drawColor,line width= 0.4pt,line join=round,line cap=round] (286.55, 98.47) -- (291.05,102.97);

\path[draw=drawColor,line width= 0.4pt,line join=round,line cap=round] (286.55,102.97) -- (291.05, 98.47);

\path[draw=drawColor,line width= 0.4pt,line join=round,line cap=round] (297.84,109.08) -- (302.34,113.58);

\path[draw=drawColor,line width= 0.4pt,line join=round,line cap=round] (297.84,113.58) -- (302.34,109.08);

\path[draw=drawColor,line width= 0.4pt,line join=round,line cap=round] (301.96,112.80) -- (306.46,117.30);

\path[draw=drawColor,line width= 0.4pt,line join=round,line cap=round] (301.96,117.30) -- (306.46,112.80);

\path[draw=drawColor,line width= 0.4pt,line join=round,line cap=round] (299.41,110.72) -- (303.91,115.22);

\path[draw=drawColor,line width= 0.4pt,line join=round,line cap=round] (299.41,115.22) -- (303.91,110.72);

\path[draw=drawColor,line width= 0.4pt,line join=round,line cap=round] (289.47,100.95) -- (293.97,105.45);

\path[draw=drawColor,line width= 0.4pt,line join=round,line cap=round] (289.47,105.45) -- (293.97,100.95);

\path[draw=drawColor,line width= 0.4pt,line join=round,line cap=round] (275.40, 88.06) -- (279.90, 92.56);

\path[draw=drawColor,line width= 0.4pt,line join=round,line cap=round] (275.40, 92.56) -- (279.90, 88.06);

\path[draw=drawColor,line width= 0.4pt,line join=round,line cap=round] (311.28,119.68) -- (315.78,124.18);

\path[draw=drawColor,line width= 0.4pt,line join=round,line cap=round] (311.28,124.18) -- (315.78,119.68);

\path[draw=drawColor,line width= 0.4pt,line join=round,line cap=round] (264.01, 76.74) -- (268.51, 81.24);

\path[draw=drawColor,line width= 0.4pt,line join=round,line cap=round] (264.01, 81.24) -- (268.51, 76.74);

\path[draw=drawColor,line width= 0.4pt,line join=round,line cap=round] (286.64, 98.54) -- (291.14,103.04);

\path[draw=drawColor,line width= 0.4pt,line join=round,line cap=round] (286.64,103.04) -- (291.14, 98.54);

\path[draw=drawColor,line width= 0.4pt,line join=round,line cap=round] (292.25,103.40) -- (296.75,107.90);

\path[draw=drawColor,line width= 0.4pt,line join=round,line cap=round] (292.25,107.90) -- (296.75,103.40);

\path[draw=drawColor,line width= 0.4pt,line join=round,line cap=round] (294.28,105.60) -- (298.78,110.10);

\path[draw=drawColor,line width= 0.4pt,line join=round,line cap=round] (294.28,110.10) -- (298.78,105.60);

\path[draw=drawColor,line width= 0.4pt,line join=round,line cap=round] (304.20,114.47) -- (308.70,118.97);

\path[draw=drawColor,line width= 0.4pt,line join=round,line cap=round] (304.20,118.97) -- (308.70,114.47);

\path[draw=drawColor,line width= 0.4pt,line join=round,line cap=round] (276.19, 88.70) -- (280.69, 93.20);

\path[draw=drawColor,line width= 0.4pt,line join=round,line cap=round] (276.19, 93.20) -- (280.69, 88.70);

\path[draw=drawColor,line width= 0.4pt,line join=round,line cap=round] (283.01, 95.47) -- (287.51, 99.97);

\path[draw=drawColor,line width= 0.4pt,line join=round,line cap=round] (283.01, 99.97) -- (287.51, 95.47);

\path[draw=drawColor,line width= 0.4pt,line join=round,line cap=round] (276.10, 88.69) -- (280.60, 93.19);

\path[draw=drawColor,line width= 0.4pt,line join=round,line cap=round] (276.10, 93.19) -- (280.60, 88.69);

\path[draw=drawColor,line width= 0.4pt,line join=round,line cap=round] (284.54, 96.71) -- (289.04,101.21);

\path[draw=drawColor,line width= 0.4pt,line join=round,line cap=round] (284.54,101.21) -- (289.04, 96.71);

\path[draw=drawColor,line width= 0.4pt,line join=round,line cap=round] (288.99,100.56) -- (293.49,105.06);

\path[draw=drawColor,line width= 0.4pt,line join=round,line cap=round] (288.99,105.06) -- (293.49,100.56);

\path[draw=drawColor,line width= 0.4pt,line join=round,line cap=round] (274.13, 86.87) -- (278.63, 91.37);

\path[draw=drawColor,line width= 0.4pt,line join=round,line cap=round] (274.13, 91.37) -- (278.63, 86.87);

\path[draw=drawColor,line width= 0.4pt,line join=round,line cap=round] (312.70,121.66) -- (317.20,126.16);

\path[draw=drawColor,line width= 0.4pt,line join=round,line cap=round] (312.70,126.16) -- (317.20,121.66);

\path[draw=drawColor,line width= 0.4pt,line join=round,line cap=round] (306.15,115.62) -- (310.65,120.12);

\path[draw=drawColor,line width= 0.4pt,line join=round,line cap=round] (306.15,120.12) -- (310.65,115.62);

\path[draw=drawColor,line width= 0.4pt,line join=round,line cap=round] (285.73, 97.90) -- (290.23,102.40);

\path[draw=drawColor,line width= 0.4pt,line join=round,line cap=round] (285.73,102.40) -- (290.23, 97.90);

\path[draw=drawColor,line width= 0.4pt,line join=round,line cap=round] (296.25,107.22) -- (300.75,111.72);

\path[draw=drawColor,line width= 0.4pt,line join=round,line cap=round] (296.25,111.72) -- (300.75,107.22);

\path[draw=drawColor,line width= 0.4pt,line join=round,line cap=round] (287.57, 99.37) -- (292.07,103.87);

\path[draw=drawColor,line width= 0.4pt,line join=round,line cap=round] (287.57,103.87) -- (292.07, 99.37);

\path[draw=drawColor,line width= 0.4pt,line join=round,line cap=round] (274.87, 87.64) -- (279.37, 92.14);

\path[draw=drawColor,line width= 0.4pt,line join=round,line cap=round] (274.87, 92.14) -- (279.37, 87.64);

\path[draw=drawColor,line width= 0.4pt,line join=round,line cap=round] (296.87,107.78) -- (301.37,112.28);

\path[draw=drawColor,line width= 0.4pt,line join=round,line cap=round] (296.87,112.28) -- (301.37,107.78);

\path[draw=drawColor,line width= 0.4pt,line join=round,line cap=round] (310.46,118.87) -- (314.96,123.37);

\path[draw=drawColor,line width= 0.4pt,line join=round,line cap=round] (310.46,123.37) -- (314.96,118.87);

\path[draw=drawColor,line width= 0.4pt,line join=round,line cap=round] (289.37,100.86) -- (293.87,105.36);

\path[draw=drawColor,line width= 0.4pt,line join=round,line cap=round] (289.37,105.36) -- (293.87,100.86);

\path[draw=drawColor,line width= 0.4pt,line join=round,line cap=round] (288.71,100.40) -- (293.21,104.90);

\path[draw=drawColor,line width= 0.4pt,line join=round,line cap=round] (288.71,104.90) -- (293.21,100.40);

\path[draw=drawColor,line width= 0.4pt,line join=round,line cap=round] (275.24, 88.04) -- (279.74, 92.54);

\path[draw=drawColor,line width= 0.4pt,line join=round,line cap=round] (275.24, 92.54) -- (279.74, 88.04);

\path[draw=drawColor,line width= 0.4pt,line join=round,line cap=round] (288.24,100.06) -- (292.74,104.56);

\path[draw=drawColor,line width= 0.4pt,line join=round,line cap=round] (288.24,104.56) -- (292.74,100.06);

\path[draw=drawColor,line width= 0.4pt,line join=round,line cap=round] (286.20, 98.17) -- (290.70,102.67);

\path[draw=drawColor,line width= 0.4pt,line join=round,line cap=round] (286.20,102.67) -- (290.70, 98.17);

\path[draw=drawColor,line width= 0.4pt,line join=round,line cap=round] (301.12,112.04) -- (305.62,116.54);

\path[draw=drawColor,line width= 0.4pt,line join=round,line cap=round] (301.12,116.54) -- (305.62,112.04);

\path[draw=drawColor,line width= 0.4pt,line join=round,line cap=round] (294.26,105.58) -- (298.76,110.08);

\path[draw=drawColor,line width= 0.4pt,line join=round,line cap=round] (294.26,110.08) -- (298.76,105.58);

\path[draw=drawColor,line width= 0.4pt,line join=round,line cap=round] (288.05, 99.89) -- (292.55,104.39);

\path[draw=drawColor,line width= 0.4pt,line join=round,line cap=round] (288.05,104.39) -- (292.55, 99.89);

\path[draw=drawColor,line width= 0.4pt,line join=round,line cap=round] (293.00,104.41) -- (297.50,108.91);

\path[draw=drawColor,line width= 0.4pt,line join=round,line cap=round] (293.00,108.91) -- (297.50,104.41);

\path[draw=drawColor,line width= 0.4pt,line join=round,line cap=round] (289.45,100.95) -- (293.95,105.45);

\path[draw=drawColor,line width= 0.4pt,line join=round,line cap=round] (289.45,105.45) -- (293.95,100.95);

\path[draw=drawColor,line width= 0.4pt,line join=round,line cap=round] (300.01,111.30) -- (304.51,115.80);

\path[draw=drawColor,line width= 0.4pt,line join=round,line cap=round] (300.01,115.80) -- (304.51,111.30);

\path[draw=drawColor,line width= 0.4pt,line join=round,line cap=round] (295.75,106.72) -- (300.25,111.22);

\path[draw=drawColor,line width= 0.4pt,line join=round,line cap=round] (295.75,111.22) -- (300.25,106.72);

\path[draw=drawColor,line width= 0.4pt,line join=round,line cap=round] (286.15, 98.17) -- (290.65,102.67);

\path[draw=drawColor,line width= 0.4pt,line join=round,line cap=round] (286.15,102.67) -- (290.65, 98.17);

\path[draw=drawColor,line width= 0.4pt,line join=round,line cap=round] (296.01,106.96) -- (300.51,111.46);

\path[draw=drawColor,line width= 0.4pt,line join=round,line cap=round] (296.01,111.46) -- (300.51,106.96);

\path[draw=drawColor,line width= 0.4pt,line join=round,line cap=round] (302.63,113.64) -- (307.13,118.14);

\path[draw=drawColor,line width= 0.4pt,line join=round,line cap=round] (302.63,118.14) -- (307.13,113.64);

\path[draw=drawColor,line width= 0.4pt,line join=round,line cap=round] (308.28,116.99) -- (312.78,121.49);

\path[draw=drawColor,line width= 0.4pt,line join=round,line cap=round] (308.28,121.49) -- (312.78,116.99);

\path[draw=drawColor,line width= 0.4pt,line join=round,line cap=round] (304.56,114.73) -- (309.06,119.23);

\path[draw=drawColor,line width= 0.4pt,line join=round,line cap=round] (304.56,119.23) -- (309.06,114.73);

\path[draw=drawColor,line width= 0.4pt,line join=round,line cap=round] (291.36,102.63) -- (295.86,107.13);

\path[draw=drawColor,line width= 0.4pt,line join=round,line cap=round] (291.36,107.13) -- (295.86,102.63);

\path[draw=drawColor,line width= 0.4pt,line join=round,line cap=round] (278.41, 91.04) -- (282.91, 95.54);

\path[draw=drawColor,line width= 0.4pt,line join=round,line cap=round] (278.41, 95.54) -- (282.91, 91.04);

\path[draw=drawColor,line width= 0.4pt,line join=round,line cap=round] (290.01,101.43) -- (294.51,105.93);

\path[draw=drawColor,line width= 0.4pt,line join=round,line cap=round] (290.01,105.93) -- (294.51,101.43);

\path[draw=drawColor,line width= 0.4pt,line join=round,line cap=round] (285.02, 97.14) -- (289.52,101.64);

\path[draw=drawColor,line width= 0.4pt,line join=round,line cap=round] (285.02,101.64) -- (289.52, 97.14);

\path[draw=drawColor,line width= 0.4pt,line join=round,line cap=round] (285.97, 98.06) -- (290.47,102.56);

\path[draw=drawColor,line width= 0.4pt,line join=round,line cap=round] (285.97,102.56) -- (290.47, 98.06);

\path[draw=drawColor,line width= 0.4pt,line join=round,line cap=round] (289.66,101.13) -- (294.16,105.63);

\path[draw=drawColor,line width= 0.4pt,line join=round,line cap=round] (289.66,105.63) -- (294.16,101.13);

\path[draw=drawColor,line width= 0.4pt,line join=round,line cap=round] (288.36,100.15) -- (292.86,104.65);

\path[draw=drawColor,line width= 0.4pt,line join=round,line cap=round] (288.36,104.65) -- (292.86,100.15);

\path[draw=drawColor,line width= 0.4pt,line join=round,line cap=round] (300.70,111.83) -- (305.20,116.33);

\path[draw=drawColor,line width= 0.4pt,line join=round,line cap=round] (300.70,116.33) -- (305.20,111.83);

\path[draw=drawColor,line width= 0.4pt,line join=round,line cap=round] (305.40,115.21) -- (309.90,119.71);

\path[draw=drawColor,line width= 0.4pt,line join=round,line cap=round] (305.40,119.71) -- (309.90,115.21);

\path[draw=drawColor,line width= 0.4pt,line join=round,line cap=round] (282.36, 94.96) -- (286.86, 99.46);

\path[draw=drawColor,line width= 0.4pt,line join=round,line cap=round] (282.36, 99.46) -- (286.86, 94.96);

\path[draw=drawColor,line width= 0.4pt,line join=round,line cap=round] (273.69, 86.57) -- (278.19, 91.07);

\path[draw=drawColor,line width= 0.4pt,line join=round,line cap=round] (273.69, 91.07) -- (278.19, 86.57);

\path[draw=drawColor,line width= 0.4pt,line join=round,line cap=round] (289.75,101.22) -- (294.25,105.72);

\path[draw=drawColor,line width= 0.4pt,line join=round,line cap=round] (289.75,105.72) -- (294.25,101.22);

\path[draw=drawColor,line width= 0.4pt,line join=round,line cap=round] (271.53, 84.29) -- (276.03, 88.79);

\path[draw=drawColor,line width= 0.4pt,line join=round,line cap=round] (271.53, 88.79) -- (276.03, 84.29);

\path[draw=drawColor,line width= 0.4pt,line join=round,line cap=round] (263.22, 75.50) -- (267.72, 80.00);

\path[draw=drawColor,line width= 0.4pt,line join=round,line cap=round] (263.22, 80.00) -- (267.72, 75.50);

\path[draw=drawColor,line width= 0.4pt,line join=round,line cap=round] (275.72, 88.44) -- (280.22, 92.94);

\path[draw=drawColor,line width= 0.4pt,line join=round,line cap=round] (275.72, 92.94) -- (280.22, 88.44);

\path[draw=drawColor,line width= 0.4pt,line join=round,line cap=round] (295.11,106.11) -- (299.61,110.61);

\path[draw=drawColor,line width= 0.4pt,line join=round,line cap=round] (295.11,110.61) -- (299.61,106.11);

\path[draw=drawColor,line width= 0.4pt,line join=round,line cap=round] (288.00, 99.80) -- (292.50,104.30);

\path[draw=drawColor,line width= 0.4pt,line join=round,line cap=round] (288.00,104.30) -- (292.50, 99.80);

\path[draw=drawColor,line width= 0.4pt,line join=round,line cap=round] (298.60,109.77) -- (303.10,114.27);

\path[draw=drawColor,line width= 0.4pt,line join=round,line cap=round] (298.60,114.27) -- (303.10,109.77);

\path[draw=drawColor,line width= 0.4pt,line join=round,line cap=round] (277.57, 90.20) -- (282.07, 94.70);

\path[draw=drawColor,line width= 0.4pt,line join=round,line cap=round] (277.57, 94.70) -- (282.07, 90.20);

\path[draw=drawColor,line width= 0.4pt,line join=round,line cap=round] (262.34, 74.97) -- (266.84, 79.47);

\path[draw=drawColor,line width= 0.4pt,line join=round,line cap=round] (262.34, 79.47) -- (266.84, 74.97);

\path[draw=drawColor,line width= 0.4pt,line join=round,line cap=round] (266.56, 79.16) -- (271.06, 83.66);

\path[draw=drawColor,line width= 0.4pt,line join=round,line cap=round] (266.56, 83.66) -- (271.06, 79.16);

\path[draw=drawColor,line width= 0.4pt,line join=round,line cap=round] (275.16, 87.95) -- (279.66, 92.45);

\path[draw=drawColor,line width= 0.4pt,line join=round,line cap=round] (275.16, 92.45) -- (279.66, 87.95);

\path[draw=drawColor,line width= 0.4pt,line join=round,line cap=round] (231.12, 49.78) -- (235.62, 54.28);

\path[draw=drawColor,line width= 0.4pt,line join=round,line cap=round] (231.12, 54.28) -- (235.62, 49.78);

\path[draw=drawColor,line width= 0.4pt,line join=round,line cap=round] (255.28, 63.59) -- (259.78, 68.09);

\path[draw=drawColor,line width= 0.4pt,line join=round,line cap=round] (255.28, 68.09) -- (259.78, 63.59);

\path[draw=drawColor,line width= 0.4pt,line join=round,line cap=round] (311.50,119.87) -- (316.00,124.37);

\path[draw=drawColor,line width= 0.4pt,line join=round,line cap=round] (311.50,124.37) -- (316.00,119.87);

\path[draw=drawColor,line width= 0.4pt,line join=round,line cap=round] (285.57, 97.77) -- (290.07,102.27);

\path[draw=drawColor,line width= 0.4pt,line join=round,line cap=round] (285.57,102.27) -- (290.07, 97.77);

\path[draw=drawColor,line width= 0.4pt,line join=round,line cap=round] (282.33, 94.95) -- (286.83, 99.45);

\path[draw=drawColor,line width= 0.4pt,line join=round,line cap=round] (282.33, 99.45) -- (286.83, 94.95);

\path[draw=drawColor,line width= 0.4pt,line join=round,line cap=round] (311.96,120.66) -- (316.46,125.16);

\path[draw=drawColor,line width= 0.4pt,line join=round,line cap=round] (311.96,125.16) -- (316.46,120.66);

\path[draw=drawColor,line width= 0.4pt,line join=round,line cap=round] (313.83,122.13) -- (318.33,126.63);

\path[draw=drawColor,line width= 0.4pt,line join=round,line cap=round] (313.83,126.63) -- (318.33,122.13);

\path[draw=drawColor,line width= 0.4pt,line join=round,line cap=round] (310.65,118.97) -- (315.15,123.47);

\path[draw=drawColor,line width= 0.4pt,line join=round,line cap=round] (310.65,123.47) -- (315.15,118.97);

\path[draw=drawColor,line width= 0.4pt,line join=round,line cap=round] (295.81,106.72) -- (300.31,111.22);

\path[draw=drawColor,line width= 0.4pt,line join=round,line cap=round] (295.81,111.22) -- (300.31,106.72);

\path[draw=drawColor,line width= 0.4pt,line join=round,line cap=round] (292.82,104.25) -- (297.32,108.75);

\path[draw=drawColor,line width= 0.4pt,line join=round,line cap=round] (292.82,108.75) -- (297.32,104.25);

\path[draw=drawColor,line width= 0.4pt,line join=round,line cap=round] (306.71,115.94) -- (311.21,120.44);

\path[draw=drawColor,line width= 0.4pt,line join=round,line cap=round] (306.71,120.44) -- (311.21,115.94);

\path[draw=drawColor,line width= 0.4pt,line join=round,line cap=round] (320.65,129.51) -- (325.15,134.01);

\path[draw=drawColor,line width= 0.4pt,line join=round,line cap=round] (320.65,134.01) -- (325.15,129.51);

\path[draw=drawColor,line width= 0.4pt,line join=round,line cap=round] (270.15, 82.88) -- (274.65, 87.38);

\path[draw=drawColor,line width= 0.4pt,line join=round,line cap=round] (270.15, 87.38) -- (274.65, 82.88);

\path[draw=drawColor,line width= 0.4pt,line join=round,line cap=round] (265.99, 78.42) -- (270.49, 82.92);

\path[draw=drawColor,line width= 0.4pt,line join=round,line cap=round] (265.99, 82.92) -- (270.49, 78.42);

\path[draw=drawColor,line width= 0.4pt,line join=round,line cap=round] (278.46, 91.15) -- (282.96, 95.65);

\path[draw=drawColor,line width= 0.4pt,line join=round,line cap=round] (278.46, 95.65) -- (282.96, 91.15);

\path[draw=drawColor,line width= 0.4pt,line join=round,line cap=round] (294.94,106.03) -- (299.44,110.53);

\path[draw=drawColor,line width= 0.4pt,line join=round,line cap=round] (294.94,110.53) -- (299.44,106.03);

\path[draw=drawColor,line width= 0.4pt,line join=round,line cap=round] (284.32, 96.58) -- (288.82,101.08);

\path[draw=drawColor,line width= 0.4pt,line join=round,line cap=round] (284.32,101.08) -- (288.82, 96.58);

\path[draw=drawColor,line width= 0.4pt,line join=round,line cap=round] (314.80,123.19) -- (319.30,127.69);

\path[draw=drawColor,line width= 0.4pt,line join=round,line cap=round] (314.80,127.69) -- (319.30,123.19);

\path[draw=drawColor,line width= 0.4pt,line join=round,line cap=round] (292.47,103.76) -- (296.97,108.26);

\path[draw=drawColor,line width= 0.4pt,line join=round,line cap=round] (292.47,108.26) -- (296.97,103.76);

\path[draw=drawColor,line width= 0.4pt,line join=round,line cap=round] (267.80, 80.76) -- (272.30, 85.26);

\path[draw=drawColor,line width= 0.4pt,line join=round,line cap=round] (267.80, 85.26) -- (272.30, 80.76);

\path[draw=drawColor,line width= 0.4pt,line join=round,line cap=round] (268.46, 81.55) -- (272.96, 86.05);

\path[draw=drawColor,line width= 0.4pt,line join=round,line cap=round] (268.46, 86.05) -- (272.96, 81.55);

\path[draw=drawColor,line width= 0.4pt,line join=round,line cap=round] (298.24,109.52) -- (302.74,114.02);

\path[draw=drawColor,line width= 0.4pt,line join=round,line cap=round] (298.24,114.02) -- (302.74,109.52);

\path[draw=drawColor,line width= 0.4pt,line join=round,line cap=round] (280.03, 92.83) -- (284.53, 97.33);

\path[draw=drawColor,line width= 0.4pt,line join=round,line cap=round] (280.03, 97.33) -- (284.53, 92.83);

\path[draw=drawColor,line width= 0.4pt,line join=round,line cap=round] (303.57,114.14) -- (308.07,118.64);

\path[draw=drawColor,line width= 0.4pt,line join=round,line cap=round] (303.57,118.64) -- (308.07,114.14);

\path[draw=drawColor,line width= 0.4pt,line join=round,line cap=round] (288.19,100.00) -- (292.69,104.50);

\path[draw=drawColor,line width= 0.4pt,line join=round,line cap=round] (288.19,104.50) -- (292.69,100.00);

\path[draw=drawColor,line width= 0.4pt,line join=round,line cap=round] (324.25,132.57) -- (328.75,137.07);

\path[draw=drawColor,line width= 0.4pt,line join=round,line cap=round] (324.25,137.07) -- (328.75,132.57);

\path[draw=drawColor,line width= 0.4pt,line join=round,line cap=round] (286.44, 98.35) -- (290.94,102.85);

\path[draw=drawColor,line width= 0.4pt,line join=round,line cap=round] (286.44,102.85) -- (290.94, 98.35);

\path[draw=drawColor,line width= 0.4pt,line join=round,line cap=round] (318.73,127.74) -- (323.23,132.24);

\path[draw=drawColor,line width= 0.4pt,line join=round,line cap=round] (318.73,132.24) -- (323.23,127.74);

\path[draw=drawColor,line width= 0.4pt,line join=round,line cap=round] (307.58,116.63) -- (312.08,121.13);

\path[draw=drawColor,line width= 0.4pt,line join=round,line cap=round] (307.58,121.13) -- (312.08,116.63);

\path[draw=drawColor,line width= 0.4pt,line join=round,line cap=round] (245.71, 54.20) -- (250.21, 58.70);

\path[draw=drawColor,line width= 0.4pt,line join=round,line cap=round] (245.71, 58.70) -- (250.21, 54.20);

\path[draw=drawColor,line width= 0.4pt,line join=round,line cap=round] (283.78, 96.24) -- (288.28,100.74);

\path[draw=drawColor,line width= 0.4pt,line join=round,line cap=round] (283.78,100.74) -- (288.28, 96.24);

\path[draw=drawColor,line width= 0.4pt,line join=round,line cap=round] (263.62, 76.08) -- (268.12, 80.58);

\path[draw=drawColor,line width= 0.4pt,line join=round,line cap=round] (263.62, 80.58) -- (268.12, 76.08);

\path[draw=drawColor,line width= 0.4pt,line join=round,line cap=round] (288.44,100.20) -- (292.94,104.70);

\path[draw=drawColor,line width= 0.4pt,line join=round,line cap=round] (288.44,104.70) -- (292.94,100.20);

\path[draw=drawColor,line width= 0.4pt,line join=round,line cap=round] (288.54,100.25) -- (293.04,104.75);

\path[draw=drawColor,line width= 0.4pt,line join=round,line cap=round] (288.54,104.75) -- (293.04,100.25);

\path[draw=drawColor,line width= 0.4pt,line join=round,line cap=round] (301.28,112.22) -- (305.78,116.72);

\path[draw=drawColor,line width= 0.4pt,line join=round,line cap=round] (301.28,116.72) -- (305.78,112.22);

\path[draw=drawColor,line width= 0.4pt,line join=round,line cap=round] (292.60,103.80) -- (297.10,108.30);

\path[draw=drawColor,line width= 0.4pt,line join=round,line cap=round] (292.60,108.30) -- (297.10,103.80);

\path[draw=drawColor,line width= 0.4pt,line join=round,line cap=round] (294.20,105.54) -- (298.70,110.04);

\path[draw=drawColor,line width= 0.4pt,line join=round,line cap=round] (294.20,110.04) -- (298.70,105.54);

\path[draw=drawColor,line width= 0.4pt,line join=round,line cap=round] (276.56, 89.10) -- (281.06, 93.60);

\path[draw=drawColor,line width= 0.4pt,line join=round,line cap=round] (276.56, 93.60) -- (281.06, 89.10);

\path[draw=drawColor,line width= 0.4pt,line join=round,line cap=round] (288.14,100.00) -- (292.64,104.50);

\path[draw=drawColor,line width= 0.4pt,line join=round,line cap=round] (288.14,104.50) -- (292.64,100.00);

\path[draw=drawColor,line width= 0.4pt,line join=round,line cap=round] (278.19, 90.78) -- (282.69, 95.28);

\path[draw=drawColor,line width= 0.4pt,line join=round,line cap=round] (278.19, 95.28) -- (282.69, 90.78);

\path[draw=drawColor,line width= 0.4pt,line join=round,line cap=round] (284.72, 96.87) -- (289.22,101.37);

\path[draw=drawColor,line width= 0.4pt,line join=round,line cap=round] (284.72,101.37) -- (289.22, 96.87);

\path[draw=drawColor,line width= 0.4pt,line join=round,line cap=round] (280.23, 93.05) -- (284.73, 97.55);

\path[draw=drawColor,line width= 0.4pt,line join=round,line cap=round] (280.23, 97.55) -- (284.73, 93.05);

\path[draw=drawColor,line width= 0.4pt,line join=round,line cap=round] (271.25, 83.87) -- (275.75, 88.37);

\path[draw=drawColor,line width= 0.4pt,line join=round,line cap=round] (271.25, 88.37) -- (275.75, 83.87);

\path[draw=drawColor,line width= 0.4pt,line join=round,line cap=round] (333.81,143.21) -- (338.31,147.71);

\path[draw=drawColor,line width= 0.4pt,line join=round,line cap=round] (333.81,147.71) -- (338.31,143.21);

\path[draw=drawColor,line width= 0.4pt,line join=round,line cap=round] (299.67,110.95) -- (304.17,115.45);

\path[draw=drawColor,line width= 0.4pt,line join=round,line cap=round] (299.67,115.45) -- (304.17,110.95);

\path[draw=drawColor,line width= 0.4pt,line join=round,line cap=round] (283.09, 95.56) -- (287.59,100.06);

\path[draw=drawColor,line width= 0.4pt,line join=round,line cap=round] (283.09,100.06) -- (287.59, 95.56);
\end{scope}
\begin{scope}
\path[clip] (  0.00,  0.00) rectangle (361.35,180.67);
\definecolor{drawColor}{RGB}{0,0,0}

\path[draw=drawColor,line width= 0.4pt,line join=round,line cap=round] (256.58, 48.00) -- (327.00, 48.00);

\path[draw=drawColor,line width= 0.4pt,line join=round,line cap=round] (256.58, 48.00) -- (256.58, 42.00);

\path[draw=drawColor,line width= 0.4pt,line join=round,line cap=round] (291.79, 48.00) -- (291.79, 42.00);

\path[draw=drawColor,line width= 0.4pt,line join=round,line cap=round] (327.00, 48.00) -- (327.00, 42.00);

\node[text=drawColor,anchor=base,inner sep=0pt, outer sep=0pt, scale=  1.00] at (256.58, 26.40) {-0.05};

\node[text=drawColor,anchor=base,inner sep=0pt, outer sep=0pt, scale=  1.00] at (291.79, 26.40) {0.00};

\node[text=drawColor,anchor=base,inner sep=0pt, outer sep=0pt, scale=  1.00] at (327.00, 26.40) {0.05};

\path[draw=drawColor,line width= 0.4pt,line join=round,line cap=round] (228.67, 70.68) -- (228.67,135.91);

\path[draw=drawColor,line width= 0.4pt,line join=round,line cap=round] (228.67, 70.68) -- (222.67, 70.68);

\path[draw=drawColor,line width= 0.4pt,line join=round,line cap=round] (228.67,103.29) -- (222.67,103.29);

\path[draw=drawColor,line width= 0.4pt,line join=round,line cap=round] (228.67,135.91) -- (222.67,135.91);

\node[text=drawColor,rotate= 90.00,anchor=base,inner sep=0pt, outer sep=0pt, scale=  1.00] at (214.28, 70.68) {-0.05};

\node[text=drawColor,rotate= 90.00,anchor=base,inner sep=0pt, outer sep=0pt, scale=  1.00] at (214.28,103.29) {0.00};

\node[text=drawColor,rotate= 90.00,anchor=base,inner sep=0pt, outer sep=0pt, scale=  1.00] at (214.28,135.91) {0.05};

\path[draw=drawColor,line width= 0.4pt,line join=round,line cap=round] (228.67, 48.00) --
	(355.35, 48.00) --
	(355.35,156.67) --
	(228.67,156.67) --
	(228.67, 48.00);
\end{scope}
\begin{scope}
\path[clip] (180.67,  0.00) rectangle (361.35,180.67);
\definecolor{drawColor}{RGB}{0,0,0}

\node[text=drawColor,anchor=base,inner sep=0pt, outer sep=0pt, scale=  1.20] at (292.01,164.53) {\bfseries QQ-plot amb una NIG};

\node[text=drawColor,anchor=base,inner sep=0pt, outer sep=0pt, scale=  1.00] at (292.01,  2.40) {Quantils NIG};

\node[text=drawColor,rotate= 90.00,anchor=base,inner sep=0pt, outer sep=0pt, scale=  1.00] at (190.27,102.34) {Quantils mostrals};
\end{scope}
\begin{scope}
\path[clip] (228.67, 48.00) rectangle (355.35,156.67);
\definecolor{drawColor}{RGB}{255,0,0}

\path[draw=drawColor,line width= 1.2pt,dash pattern=on 4pt off 4pt ,line join=round,line cap=round] (228.67, 44.83) -- (355.35,162.16);
\end{scope}
\end{tikzpicture}

	\caption{Comparació amb una distribució NIG}
	\label{fig:comparacio nig}
\end{figure}

\begin{lstlisting}[language=R]
# Comentari

Codi R
\end{lstlisting}

\bibliographystyle{plain}
\bibliography{fonts}

\end{document}
